\documentclass[11pt]{article}
\usepackage[T1]{fontenc}
\usepackage{amsmath,amssymb,amsthm}
\usepackage[margin=1in]{geometry}
\usepackage{enumitem}
\usepackage[colorlinks=true,linkcolor=blue,urlcolor=blue]{hyperref}

\newtheorem{theorem}{Theorem}[section]
\newtheorem{proposition}[theorem]{Proposition}
\newtheorem{lemma}[theorem]{Lemma}
\newtheorem{corollary}[theorem]{Corollary}
\newtheorem{definition}[theorem]{Definition}
\newtheorem{warning}[theorem]{Warning}
\newtheorem{firewall}[theorem]{Claim firewall}
\newtheorem{decision}[theorem]{Next decision}
\newtheorem{assessment}[theorem]{Assessment}
\newtheorem{ledger}[theorem]{Acquisition ledger}
\theoremstyle{remark}
\newtheorem{remark}[theorem]{Remark}

\title{Renewal Standard Model--Einstein Continuum Research Notes\\
Fourteenth Cycle, V1}
\author{}
\date{9 September 2026}

\begin{document}
\date{10 September 2026}
\maketitle
\tableofcontents

\section{Context, lineage, and purpose of the fourteenth cycle}
\label{sec:v14v1}

These research notes continue a long-form attempt to construct a
rigorous continuum limit of a regulated Standard Model coupled to
Einstein gravity.  The intended endpoint is one compatible continuum
state carrying the gauge, fermion, Higgs, and gravitational sectors,
with the physical gauge and diffeomorphism identities, a controlled
metric/coframe sector, Osterwalder--Schrader reconstruction, local
covariance, and the required massive and massless long-distance
structure.

That endpoint has not been reached.  The archives contain proved
conditional implications, exact finite-stage identities, compactness
and landing theorems, and explicit no-go results.  They do not yet
populate every hypothesis with one common regulator.  This compact
restart records the major acquired results, states the current claim
boundary, and then continues from the uniform
metric-compatible-connection problem reached at the end of the
thirteenth cycle.

The completed thirteenth cycle is frozen as
\[
 \texttt{Renewal\_SM\_Einstein\_Continuum\_Research\_Notes\_V100\_f13.tex}.
                                                                  \tag{1.1}
\]
Its exact record is
\[
\begin{array}{c|c}
 \text{lines}&29827\\
 \text{bytes}&1047050\\
 \text{SHA--256}&
 \texttt{585D42BB92D495BE9FCD043103784FC619F970E5D27B303A0F49529BF944129D}.
\end{array}                                                     \tag{1.2}
\]
Earlier frozen cycles remain independent proof archives.  Future
versions append to this V1 and replace only the immediate active
predecessor.  The newest Yang--Mills and Navier--Stokes companions are
audited at every fifth version.  At V100 the active file is frozen
with the next suffix and a new compact V1 begins.

\section{Major mathematical results acquired so far}
\label{sec:v14v1-summary}

This section is a proof-status map.  Each item retains the hypotheses
of its archived theorem; the summary does not turn a conditional
result into an unconditional construction.

\subsection{Finite regulators, determinants, topology, and symmetry}

\begin{enumerate}[leftmargin=2.2em]
\item Finite-cutoff renewal maps, local slice carriers, Schur--Feshbach
      reductions, determinant regularizations, and shell
      normalizations have been separated into explicit algebraic,
      spectral, and stability hypotheses.
\item Exact orbit--stabilizer weights and full-density coordinate
      covariance show why a partially compressed diagram cannot be
      assigned its final symmetry factor before the full orbit has
      been retained.
\item Relative trace-class or Hilbert--Schmidt increments control
      Fredholm determinant moduli and phases on zero-free paths.
      Exact local phase one-forms admit counterterms; cycle holonomy
      and determinant-line topology remain separate.
\item Gapped Hermitian carriers yield integer sector labels stable
      under certified homotopies.  A verified action cost exceeding
      sector entropy produces exponential sector tails.
\item A quantitative Faddeev--Popov gap gives a local
      diffeomorphism slice and relative determinant control.  Global
      quotient construction still requires overlap Jacobians,
      stabilizer and Gribov multiplicities, and equivariant
      refinement.
\end{enumerate}

\subsection{Anomalies, BRST structure, and physical Ward identities}

\begin{enumerate}[leftmargin=2.2em]
\item Finite cochain Hodge decompositions split a regulated anomaly
      into a harmonic obstruction and an exact removable part.
      Controlled chain homotopies transport the cohomology class
      through refinement.
\item Standard Model representation arithmetic cancels the universal
      local anomaly polynomial under the declared charge assignments.
      Regulator remainders, torsion classes, and global
      determinant-line anomalies require independent control.
\item A harmonic anomaly lying in a proved integral lattice and below
      its systole vanishes exactly.  The lattice membership, systole,
      and torsion checks are physical inputs rather than consequences
      of formal power counting.
\item Exact or vanishing-defect Ward identities pass through
      normalized complex shell martingales and through convergent
      joint states on declared local cores.
\item Closed unbounded BRST derivations admit a graph-domain passage
      when predual anomaly functionals and graph-norm operator errors
      converge.  The resulting anomaly functional is projectively
      natural under BRST-intertwining inclusions.
\end{enumerate}

\subsection{Gravity, contours, normalization, and continuum laws}

\begin{enumerate}[leftmargin=2.2em]
\item The bare Euclidean Einstein--Hilbert action is unbounded below
      along bounded fixed-volume, fixed-topology conformal
      oscillations.  It cannot be used without qualification as a
      positive amplitude base.
\item A coercive real Gaussian conformal reference has a
      Sobolev-valued projective continuum.  A naive fourth-order
      regulator cannot simultaneously retain nondegenerate fixed
      modes and the uniform reweighting moment needed for its removal.
\item Independent negative quadratic conformal modes have an exact
      imaginary-axis Gaussian contour.  Orientation normalization
      yields projective consistency and an infinite Sobolev-valued
      amplitude under a trace condition.
\item General complex contours require real-part coercivity,
      divisor avoidance, relative homology, normalized new-mode
      factors, and a denominator anchor.  Antilinear contour
      invariance proves reality; positivity requires a separate Gram
      factorization.
\item Projectively consistent Gaussian reference laws exist under
      covariance consistency and spectral-tail tightness.  Summable
      adjacent covariance, density, or conditional-state defects
      produce cofinal limits with quantitative tails.
\item Uniform gravitational regularity moments and a coercive action
      floor yield compact containment and Wasserstein subsequences on
      a separable state space.  Quotient tightness can be proved
      without a global gauge section when the metric quotient is
      Polish.
\item A countable distance-coordinate ledger reconstructs a bounded
      Polish quotient metric, its topology, and its Borel laws.  A
      dense observable ledger upgrades subsequential agreement to
      full-sequence convergence.
\end{enumerate}

\subsection{Quantum reconstruction, locality, and the joint state}

\begin{enumerate}[leftmargin=2.2em]
\item A joint classical--quantum state on
      \(C(K)\otimes_{\min}\widetilde{\mathcal K}({\cal H})\)
      has a trace-class conditional-density disintegration when a
      compact-resolvent energy moment prevents quantum mass from
      escaping to infinity.
\item Compatible local restrictions define a quasi-local state.
      Regionwise energy bounds suffice; no volume-uniform bound on
      the total infinite-volume energy is required.
\item An all-order positive compatible moment ledger is necessary:
      explicit moment examples prove that no fixed finite correlation
      order determines an interacting state.
\item Osterwalder--Schrader positivity on the physical algebra yields
      a Hilbert space, semigroup, Hamiltonian, and analytic
      continuation under the recorded domain and energy hypotheses.
      Dense subgroup covariance plus a uniform identity modulus
      extends to the continuum group; without the modulus, factorial
      characters give a counterexample.
\item Uniform ADM bounds yield a common causal cone and a countable
      rational region ledger.  Locality extends from the determining
      regions to all compact region pairs with a strict common causal
      margin.
\item Time-slice norm-density reconstructs the local net from a common
      slab.  Resolvent affiliation places unbounded local observables
      in the net, and spacelike affiliated observables strongly
      commute.
\item Dense test-function control and all-order product seminorms
      construct graph-continuous operator-valued distributions and
      their Wightman hierarchy on a common invariant domain.
\item Averaged clustering plus a cofinal Gram condition gives vacuum
      uniqueness.  Purity is distinct from uniqueness of the
      invariant vector; in a pure covariant classical-center sector,
      the metric center selects a fixed background class.
\end{enumerate}

\subsection{Recent noncentral Einstein-interface results}

The thirteenth cycle moved beyond a classical central metric by using
an invertible noncentral coframe coefficient \(E\) and
\[
 G=E^*\eta E.
\]
It established the following results.

\begin{enumerate}[leftmargin=2.2em]
\item Two-sided invertibility floors and inverse-closed smoothness give
      ordered first and second inverse derivatives, quantitative
      bounds, and cofinal landing for \(E^{-1}\) and the metric jets.
\item A differential graded connection has curvature and torsion with
      exact ordered Cartan identities.  At imperfect finite stages,
      the Bianchi failures equal the nilpotence and Leibniz defects
      evaluated on the actual connection and coframe.
\item The universal graded cyclic trace closes the curvature Chern
      forms modulo commutators.  Explicit \(2\times2\) examples prove
      that an ordinary coefficient-valued matrix trace is not cyclic
      and that left and right trace reversals need not agree.
\item The symmetric two-word continuation
      \[
       {\cal G}_{\mathrm J}=P-\frac14(gs+sg)
      \]
      is selected by classical normalization and reversal symmetry.
      Its divergence defect factors exactly into Ricci, scalar,
      metricity, and uncontracted-Bianchi defects.
\item For \(G=E^*\eta E\), the coframe-induced connection
      \(\Gamma_0=E^{-1}dE\) is metric compatible but flat.  Every
      metric-compatible connection is
      \[
       \Gamma=\Gamma_0+G^{-1}H,\qquad H^*=-H.
      \]
      Torsion-freeness is the linear Cartan selector equation
      \(G^{-1}H\theta=-{\cal T}_0\).
\item At a central orthonormal anchor the selector is an isomorphism,
      with an explicit Cartan coefficient inverse.  A Neumann margin
      gives unique metric-compatible torsion-free connections and
      cofinal landing in a local noncommutative cell.
\end{enumerate}

\section{Current claim boundary}
\label{sec:v14v1-boundary}

The full SM--Einstein continuum remains open.  The principal live
gates are
\[
\begin{array}{p{0.31\linewidth}|p{0.57\linewidth}}
\textbf{gate}&\textbf{missing evidence}\\ \hline
\text{uniform connection selector}
 &a finite or cofinal physical-region cover with a positive inverse
  margin\\
\text{ordered Ricci and scalar maps}
 &frame covariance, adjoint compatibility, classical reduction,
  differentiated landing, and the V97 defect relation\\
\text{Einstein/source matching}
 &support of the joint state on an ordered Einstein equation whose
  source is the already landed all-order stress tensor\\
\text{single-regulator population}
 &one regulator satisfying positivity, anomaly, determinant,
  quotient, tightness, locality, and connection hypotheses together\\
\text{infrared sector}
 &the Standard Model mass pattern, gravitational massless sector,
  clustering statements, and phase decomposition in the same
  continuum representation
\end{array}                                                     \tag{1.3}
\]
No claim below assumes those gates.

\section{Relative signature cells for the Cartan selector}
\label{sec:v14v1-relative}

The previous cycle bounded selector perturbations by
\(\|G^{-1}-G_0^{-1}\|\).  For an indefinite Einstein metric, a
positive Loewner comparison is unavailable.  The correct replacement
is a relative cell built from the absolute value and sign of the
self-adjoint anchor.

Let \(G_0=G_0^*\) be invertible in a matrix \(C^*\)-algebra.  Define
\[
 |G_0|:=(G_0^2)^{1/2},\qquad
 J_0:=G_0|G_0|^{-1}.                                        \tag{1.4}
\]
Then \(J_0=J_0^*=J_0^{-1}\), it commutes with \(|G_0|\), and
\[
 G_0=|G_0|^{1/2}J_0|G_0|^{1/2}.                             \tag{1.5}
\]
For another self-adjoint \(G\), set
\[
 Z_0(G):=
 |G_0|^{-1/2}(G-G_0)|G_0|^{-1/2},\qquad
 r_0(G):=\|Z_0(G)\|.                                        \tag{1.6}
\]

\begin{theorem}[Relative signature-preserving inverse cell]
\label{thm:v14v1-relativeinverse}
If
\[
 r_0(G)<1,                                                   \tag{1.7}
\]
then \(G\) is invertible and
\[
 G^{-1}
 =|G_0|^{-1/2}(J_0+Z_0(G))^{-1}|G_0|^{-1/2}.                \tag{1.8}
\]
Moreover,
\[
 \|G^{-1}\|
 \leq\frac{\|G_0^{-1}\|}{1-r_0(G)},                         \tag{1.9}
\]
and
\[
 \|G^{-1}-G_0^{-1}\|
 \leq
 \|G_0^{-1}\|\frac{r_0(G)}{1-r_0(G)}.                       \tag{1.10}
\]
The path \(G_t=G_0+t(G-G_0)\), \(0\leq t\leq1\), consists of
invertible self-adjoint elements.  In every finite-dimensional
representation, \(G\) and \(G_0\) therefore have the same inertia.
\end{theorem}

\begin{proof}
Equations (1.5)--(1.6) give
\[
 G=|G_0|^{1/2}(J_0+Z_0(G))|G_0|^{1/2}.                      \tag{1.11}
\]
Since
\[
 J_0+Z_0(G)=J_0\bigl(I+J_0Z_0(G)\bigr)
\]
and \(\|J_0Z_0(G)\|=r_0(G)<1\), the Neumann series proves
invertibility and
\[
 \|(J_0+Z_0(G))^{-1}\|\leq\frac1{1-r_0(G)}.
\]
This yields (1.8)--(1.9).  Also
\[
 (J_0+Z_0(G))^{-1}-J_0
 =
 \left[(I+J_0Z_0(G))^{-1}-I\right]J_0,
\]
whose norm is at most \(r_0(G)/(1-r_0(G))\).  Conjugating by
\(|G_0|^{-1/2}\) proves (1.10).

For \(G_t\), replace \(Z_0(G)\) by \(tZ_0(G)\); its norm remains
strictly below one.  Thus the path is self-adjoint and invertible.
Finite-dimensional inertia is constant on such a path.
\end{proof}

\begin{proposition}[Relative control is scale adapted]
\label{prop:v14v1-scale}
Let \(J=J^*=J^{-1}\), \(0<\varepsilon<1\), and for
\(\lambda>0\) set
\[
 G_{0,\lambda}=\lambda J,\qquad
 G_\lambda=\lambda(J+\varepsilon I).
\]
Then
\[
 r_{0,\lambda}(G_\lambda)=\varepsilon
 \quad\hbox{for every }\lambda,
 \qquad
 \|G_\lambda-G_{0,\lambda}\|=\lambda\varepsilon.             \tag{1.12}
\]
Thus the relative cell is uniform under an overall metric scale
where any fixed absolute-radius cell eventually fails.
\end{proposition}

\begin{proof}
Here \(|G_{0,\lambda}|=\lambda I\), so (1.6) gives
\(Z_{0,\lambda}(G_\lambda)=\varepsilon I\).  The absolute difference
is \(\lambda\varepsilon I\).
\end{proof}

\subsection{Relative Cartan-selector margin}

Fix an anchor pair \((G_0,\theta_0)\) for which
\[
 L_0(H):=G_0^{-1}H\theta_0
\]
is invertible on the real space of anti-self-adjoint one-form
matrices, and put
\[
 \kappa_0:=\|L_0^{-1}\|.                                    \tag{1.13}
\]
For \(G,\theta\), let \(L_{G,\theta}(H)=G^{-1}H\theta\).

\begin{lemma}[Relative selector perturbation]
\label{lem:v14v1-selectorbound}
If \(r_0(G)<1\), then
\[
 \|L_{G,\theta}-L_0\|
 \leq
 \|G_0^{-1}\|
 \left[
  \frac{r_0(G)}{1-r_0(G)}\,\|\theta\|
  +\|\theta-\theta_0\|
 \right]
 =:\Delta_0^{\mathrm{rel}}(G,\theta).                        \tag{1.14}
\]
\end{lemma}

\begin{proof}
For \(\|H\|=1\),
\[
\begin{aligned}
 G^{-1}H\theta-G_0^{-1}H\theta_0
 &=(G^{-1}-G_0^{-1})H\theta
   +G_0^{-1}H(\theta-\theta_0).
\end{aligned}
\]
Apply (1.10), submultiplicativity, and take the supremum over \(H\).
\end{proof}

\begin{theorem}[Relative noncentral Levi--Civita cell]
\label{thm:v14v1-selectorcell}
If
\[
 r_0(G)<1,\qquad
 q_0(G,\theta):=
 \kappa_0\Delta_0^{\mathrm{rel}}(G,\theta)<1,                \tag{1.15}
\]
then \(L_{G,\theta}\) is invertible and
\[
 \|L_{G,\theta}^{-1}\|
 \leq\frac{\kappa_0}{1-q_0(G,\theta)}.                       \tag{1.16}
\]
For every torsion source \({\cal T}_0\), the equations
\[
 H^*=-H,\qquad G^{-1}H\theta=-{\cal T}_0                    \tag{1.17}
\]
have the unique solution
\[
 H=-L_{G,\theta}^{-1}{\cal T}_0,\qquad
 \|H\|\leq
 \frac{\kappa_0\|{\cal T}_0\|}
 {1-q_0(G,\theta)}.                                         \tag{1.18}
\]
\end{theorem}

\begin{proof}
The bound (1.14) implies
\[
 \|L_0^{-1}(L_{G,\theta}-L_0)\|
 \leq q_0(G,\theta)<1.
\]
Factor \(L_{G,\theta}\) through \(L_0\) and apply the Neumann series.
This proves (1.16).  Equations (1.17)--(1.18) follow by applying the
inverse to the torsion source.
\end{proof}

\subsection{Finite-cover and cofinal persistence theorem}

\begin{theorem}[Uniform selector cover]
\label{thm:v14v1-cover}
Let \(X\) be a compact physical parameter space and suppose
\(x\mapsto(G(x),\theta(x),{\cal T}_0(x))\) is norm continuous.
Assume there are finitely many certified anchors
\((G_i,\theta_i)\), open sets \(U_i\) covering \(X\), and numbers
\(q_i^*<1\) such that, for every \(x\in U_i\),
\[
 r_i(G(x))<1,\qquad
 q_i(G(x),\theta(x))\leq q_i^*.                              \tag{1.19}
\]
Then every selector \(L_{G(x),\theta(x)}\) is invertible and
\[
 \sup_{x\in X}\|L_{G(x),\theta(x)}^{-1}\|
 \leq
 \max_i\frac{\kappa_i}{1-q_i^*}<\infty.                     \tag{1.20}
\]
The unique metric-compatible torsion-free correction \(H(x)\) is
continuous on \(X\).

If cofinal data \(G_n,\theta_n,{\cal T}_{0,n}\) converge uniformly on
\(X\) to the displayed data, and the inequalities (1.19) have a
strict uniform slack, then the same finite cover is valid for all
large \(n\), the inverse bound is uniform in \(n,x\), and
\[
 \sup_{x\in X}\|H_n(x)-H(x)\|\longrightarrow0.               \tag{1.21}
\]
\end{theorem}

\begin{proof}
For \(x\in U_i\), apply
Theorem~\ref{thm:v14v1-selectorcell} at anchor \(i\).  The finite
maximum gives (1.20).  Inversion is continuous on the open set of
invertible operators, so the local formulas agree on overlaps by
uniqueness and define a continuous \(H\).

Uniform convergence of \(G_n\) implies uniform convergence of every
fixed-anchor relative coordinate \(Z_i(G_n)\).  Uniform convergence
of \(\theta_n\) then gives uniform convergence of the bounds in
(1.14).  Strict slack and finiteness of the cover preserve (1.19)
for all large \(n\).  The resolvent identity
\[
 L_n^{-1}-L^{-1}=L_n^{-1}(L-L_n)L^{-1}
\]
gives uniform inverse convergence.  Combining it with uniform
convergence of the torsion source proves (1.21).
\end{proof}

\begin{remark}[What the cover theorem requires]
Theorem~\ref{thm:v14v1-cover} converts finite certified anchor data
into a global uniform selector bound on a declared compact physical
region.  It does not prove that the regulator's physical metric range
is compact, that a proposed list of anchors covers it, or that every
anchor selector is invertible.  Those are the next population tasks.
\end{remark}

\section{Acquisition ledger and next step}
\label{sec:v14v1-ledger}

\[
\begin{array}{p{0.29\linewidth}|p{0.23\linewidth}|p{0.38\linewidth}}
\textbf{gate}&\textbf{status at V1}&\textbf{content}\\ \hline
\text{indefinite relative metric cell}
 &\text{acquired}
 &absolute-value/sign factorization and inverse bounds\\
\text{signature preservation}
 &\text{acquired inside the cell}
 &invertible self-adjoint homotopy\\
\text{relative selector cell}
 &\text{acquired}
 &explicit margin \(q_i<1\)\\
\text{finite uniform cover theorem}
 &\text{acquired conditionally}
 &certified anchors plus strict slack\\
\text{physical cover population}
 &\text{open}
 &bounds for the actual cofinal metric/coframe range\\
\text{Ricci/scalar contraction}
 &\text{open}
 &ordered contraction on the landed connection curvature
\end{array}                                                     \tag{1.22}
\]

The next version should construct the ordered Ricci contraction first
at one certified selector anchor.  It must retain frame covariance
and the adjoint operation, and its differentiated form must expose the
intertwining defect needed by the symmetric Einstein trace reversal.
If the contraction cannot be defined on the present bimodule of
forms, the programme must return upstream to the choice of
differential calculus rather than use a classical index permutation
that the coefficient algebra does not permit.

\begin{assessment}[Status after fourteenth-cycle V1]
The archive and compact summary are complete, and the connection
selector now has a relative indefinite-metric cover theorem.  No
physical anchor cover, ordered Ricci/scalar map, Einstein support
identity, or full SM--Einstein continuum has been proved.
\end{assessment}
% =============================================================================
% BEGIN APPENDED FOURTEENTH-CYCLE V2 MATERIAL
% =============================================================================

\section{Fourteenth-cycle V2: anchor Ricci contraction and the
noncentral dual-frame obstruction}
\label{sec:v14v2}

\subsection{A central reference calculus}

The relative selector theorem in Section~\ref{sec:v14v1-relative}
constructs a metric-compatible torsion-free connection inside each
certified cell.  Curvature still has to be contracted.  We begin at a
reference calculus where coefficient extraction is unambiguous.

Let \(\mathfrak A\) be a unital \(C^*\)-algebra and suppose
\(\Omega^1\) is freely generated, as both a left and a right
\(\mathfrak A\)-module, by central one-forms
\[
 e^1,\ldots,e^m,\qquad ae^\mu=e^\mu a
 \quad(a\in\mathfrak A).                                    \tag{2.1}
\]
Assume the usual antisymmetric exterior basis.  Every curvature entry
then has a unique expansion
\[
 {\cal R}^a{}_b
 =\frac12 R^a{}_{b\mu\nu}e^\mu e^\nu,\qquad
 R^a{}_{b\mu\nu}=-R^a{}_{b\nu\mu}.                          \tag{2.2}
\]
Define the central-frame interior maps by
\[
 \iota_\rho(a)=0,\qquad
 \iota_\rho(e^\mu)=\delta^\mu_\rho,\qquad
 \iota_\rho(\alpha\beta)
 =\iota_\rho(\alpha)\beta
  +(-1)^{|\alpha|}\alpha\iota_\rho(\beta).                  \tag{2.3}
\]
Centrality makes (2.3) compatible with both coefficient actions.

\begin{lemma}[Coefficient recovery]
\label{lem:v14v2-recovery}
For the curvature (2.2),
\[
 \iota_\nu\iota_\rho{\cal R}^a{}_b
 =R^a{}_{b\rho\nu}.                                         \tag{2.4}
\]
\end{lemma}

\begin{proof}
Formula (2.3) gives
\[
 \iota_\rho(e^\mu e^\lambda)
 =\delta^\mu_\rho e^\lambda
  -\delta^\lambda_\rho e^\mu.
\]
Applying \(\iota_\nu\), inserting this in (2.2), and using
antisymmetry in \(\mu,\lambda\) yields (2.4).
\end{proof}

\begin{definition}[Anchor Ricci contraction]
At the central reference frame define
\[
 P_{b\nu}({\cal R})
 :=\sum_{a=1}^m
   \iota_\nu\iota_a{\cal R}^a{}_b
 =\sum_{a=1}^m R^a{}_{ba\nu}.                               \tag{2.5}
\]
This is an ordered linear coefficient contraction.  It introduces no
coefficient product and no cyclic trace.
\end{definition}

\begin{proposition}[Classical agreement]
\label{prop:v14v2-classicalricci}
When \(\mathfrak A\) is commutative and (2.2) is the curvature of a
classical metric connection, (2.5) is the usual Ricci tensor in the
sign convention fixed by (2.2).
\end{proposition}

\begin{proof}
The classical component definition in this convention is
\(\operatorname{Ric}_{b\nu}=R^a{}_{ba\nu}\), summed over \(a\).
This is exactly (2.5).
\end{proof}

\subsection{Symmetric scalar contraction and Einstein component}

Let \(G=(G_{\mu\nu})=G^*\) be an invertible metric coefficient matrix
and write \(g=G^{-1}=(g^{\mu\nu})\).  Define
\[
\begin{aligned}
 s_{\mathrm L}({\cal R})
 &:=\sum_{b,\nu}g^{\nu b}P_{b\nu}({\cal R}),\\
 s_{\mathrm R}({\cal R})
 &:=\sum_{b,\nu}P_{b\nu}({\cal R})g^{\nu b},\\
 s_{\mathrm J}({\cal R})
 &:=\frac12\bigl(s_{\mathrm L}({\cal R})
                 +s_{\mathrm R}({\cal R})\bigr).             \tag{2.6}
\end{aligned}
\]
The anchor Einstein component is
\[
 {\cal G}_{b\nu}^{\mathrm J}({\cal R})
 :=
 P_{b\nu}({\cal R})
 -\frac14\left(
   G_{b\nu}s_{\mathrm J}({\cal R})
   +s_{\mathrm J}({\cal R})G_{b\nu}
 \right).                                                     \tag{2.7}
\]

\begin{theorem}[Reality and classical reduction]
\label{thm:v14v2-reality}
Assume
\[
 P_{b\nu}({\cal R})^*=P_{\nu b}({\cal R}).                   \tag{2.8}
\]
Then
\[
 s_{\mathrm L}({\cal R})^*=s_{\mathrm R}({\cal R}),\qquad
 s_{\mathrm J}({\cal R})^*=s_{\mathrm J}({\cal R}),          \tag{2.9}
\]
and
\[
 \bigl({\cal G}_{b\nu}^{\mathrm J}({\cal R})\bigr)^*
 ={\cal G}_{\nu b}^{\mathrm J}({\cal R}).                    \tag{2.10}
\]
If the metric and Ricci coefficients commute, then
\[
 s_{\mathrm J}=g^{\nu b}P_{b\nu},\qquad
 {\cal G}_{b\nu}^{\mathrm J}
 =P_{b\nu}-\frac12G_{b\nu}s_{\mathrm J},                    \tag{2.11}
\]
the classical scalar and Einstein trace reversal.
\end{theorem}

\begin{proof}
Because \(g=g^*\), its entries satisfy
\((g^{\nu b})^*=g^{b\nu}\).  Hence
\[
\begin{aligned}
 s_{\mathrm L}^*
 &=\sum_{b,\nu}P_{b\nu}^* (g^{\nu b})^*
 =\sum_{b,\nu}P_{\nu b}g^{b\nu}
 =s_{\mathrm R},
\end{aligned}
\]
after relabelling \(b,\nu\).  This proves (2.9).  Taking the adjoint
of (2.7), using \(G_{b\nu}^*=G_{\nu b}\), proves (2.10).
Under commutativity, the two scalar orderings agree and the Jordan
product in (2.7) collapses to the usual one, proving (2.11).
\end{proof}

\subsection{Covariance under central changes of frame}

Let \(u=(u^a{}_r)\in GL_m(Z(\mathfrak A))\).  Suppose curvature,
metric, and inverse metric components transform by
\[
\begin{aligned}
 R'^a{}_{bcd}
 &=u^a{}_r (u^{-1})^s{}_b
   (u^{-1})^p{}_c (u^{-1})^q{}_d
   R^r{}_{spq},\\
 G'_{bd}
 &=(u^{-1})^s{}_b (u^{-1})^q{}_d G_{sq},\\
 g'^{db}
 &=u^d{}_p u^b{}_r g^{pr}.                                  \tag{2.12}
\end{aligned}
\]

\begin{theorem}[Central-frame covariance]
\label{thm:v14v2-centralcovariance}
Under (2.12),
\[
 P'_{bd}
 =(u^{-1})^s{}_b (u^{-1})^q{}_d P_{sq},\qquad
 s'_{\mathrm J}=s_{\mathrm J},                              \tag{2.13}
\]
and
\[
 ({\cal G}^{\mathrm J})'_{bd}
 =(u^{-1})^s{}_b (u^{-1})^q{}_d
   {\cal G}^{\mathrm J}_{sq}.                               \tag{2.14}
\]
\end{theorem}

\begin{proof}
Contract the first and third indices of the first formula in (2.12).
Centrality permits the factors to be grouped, and
\[
 \sum_a u^a{}_r (u^{-1})^p{}_a=\delta^p_r.
\]
This gives the Ricci formula in (2.13).  Insert that formula and the
inverse-metric formula into (2.6).  The two remaining pairs
\(u^{-1}u\) cancel, separately in the left and right orderings, so
both scalar contractions are invariant.  Finally insert these
relations into (2.7); centrality of the entries of \(u\) gives
(2.14).
\end{proof}

Theorem~\ref{thm:v14v2-centralcovariance} is a genuine covariance
result on the declared central frame group.  It is not yet covariance
under a frame matrix with noncentral entries.

\subsection{Differentiating the two interior contractions}

The contraction-intertwining defect can be made explicit before a
noncentral dual calculus is chosen.

\begin{lemma}[Double-contraction commutator identity]
\label{lem:v14v2-doublecommutator}
Let \(D\) be a degree-one operator on a graded space and let
\(I,J\) be degree-minus-one operators.  Put
\[
 K_I:=DI+ID,\qquad K_J:=DJ+JD.                               \tag{2.15}
\]
Then
\[
 DIJ-IJD=K_IJ-IK_J.                                         \tag{2.16}
\]
\end{lemma}

\begin{proof}
Use \(DI=K_I-ID\) and \(DJ=K_J-JD\):
\[
 DIJ=K_IJ-IDJ
 =K_IJ-IK_J+IJD.
\]
Rearrangement proves (2.16).
\end{proof}

Let \(D\) now denote the covariant differential on curvature forms
and their coefficient tensors, with frame-index contraction
compatible with \(D\).  Define the Cartan commutator defects
\[
 K_\rho:=D\iota_\rho+\iota_\rho D.                           \tag{2.17}
\]

\begin{theorem}[Exact anchor Ricci differentiation defect]
\label{thm:v14v2-riccidefect}
For the contraction (2.5),
\[
\boxed{
 {\cal J}_{P,b\nu}({\cal R})
 :=
 D P_{b\nu}({\cal R})-P_{b\nu}(D{\cal R})
 =
 \sum_a
 \left(K_\nu\iota_a-\iota_\nu K_a\right)
 {\cal R}^a{}_b.}                                           \tag{2.18}
\]
Thus differentiating Ricci is controlled by the actual
connection--interior commutators; it is not obtained by a formal
index move.
\end{theorem}

\begin{proof}
Apply Lemma~\ref{lem:v14v2-doublecommutator} with
\(I=\iota_\nu\), \(J=\iota_a\) to each entry
\({\cal R}^a{}_b\), then sum over \(a\).  Compatibility of \(D\)
with the contracted frame index introduces no additional term.
\end{proof}

For the scalar map, write
\[
 Q^{\nu b}:=Dg^{\nu b}.                                     \tag{2.19}
\]

\begin{corollary}[Exact scalar differentiation defect]
\label{cor:v14v2-scalardefect}
Let \({\cal J}_{s_{\mathrm J}}({\cal R})
 :=Ds_{\mathrm J}({\cal R})-s_{\mathrm J}(D{\cal R})\).
Then
\[
\boxed{
\begin{aligned}
 {\cal J}_{s_{\mathrm J}}({\cal R})
 =\frac12\sum_{b,\nu}\Bigl[
 &Q^{\nu b}P_{b\nu}
 +g^{\nu b}{\cal J}_{P,b\nu}\\
 &+{\cal J}_{P,b\nu}g^{\nu b}
 +P_{b\nu}Q^{\nu b}
 \Bigr].                                                     \tag{2.20}
\end{aligned}}
\]
In particular, inverse-metric compatibility \(Q=0\) leaves only the
symmetrized Ricci-intertwining defect.
\end{corollary}

\begin{proof}
Apply the degree-zero Leibniz rule to the two products in (2.6),
then subtract the same scalar contraction applied to \(D{\cal R}\).
The four displayed terms are the remaining differences.
\end{proof}

Equations (2.18)--(2.20) populate the abstract \(P\)- and
\(s\)-defects in the symmetric trace-reversal gate summarized in V1.
They do not assert that the right-hand side vanishes; the contracted
Bianchi cancellation still has to be proved in the selected calculus.

\subsection{A rank-one obstruction to ordinary noncentral duality}

The following elementary result explains why
Theorem~\ref{thm:v14v2-centralcovariance} cannot simply be extended by
dropping the word central.

\begin{theorem}[Noncentral dual-frame no-go]
\label{thm:v14v2-dualnogo}
Let \(\Omega^1=\mathfrak A\theta=\theta\mathfrak A\) be a free
rank-one bimodule with central generator \(\theta\).  Let
\(u\in GL_1(\mathfrak A)=\mathfrak A^\times\) and set
\(\theta'=u\theta\).  If there exists an
\(\mathfrak A\)-bimodule map
\[
 \iota':\Omega^1\longrightarrow\mathfrak A
 \quad\hbox{such that}\quad
 \iota'(\theta')=1,                                         \tag{2.21}
\]
then \(u\) is central.  Consequently, a genuinely noncentral frame
change has no ordinary bimodule-linear dual evaluation of this form.
\end{theorem}

\begin{proof}
Left linearity and (2.21) give
\[
 1=\iota'(u\theta)=u\,\iota'(\theta),
\]
so \(\iota'(\theta)=u^{-1}\).  For arbitrary \(a\in\mathfrak A\),
centrality of \(\theta\), left linearity, and right linearity give
\[
 a u^{-1}
 =\iota'(a\theta)
 =\iota'(\theta a)
 =u^{-1}a.
\]
Thus \(u^{-1}\), and hence \(u\), lies in \(Z(\mathfrak A)\).
\end{proof}

\begin{remark}[Required upstream datum]
Theorem~\ref{thm:v14v2-dualnogo} does not forbid noncommutative Ricci
geometry.  It proves that one must change at least one ingredient:
use a braided/twisted bimodule dual, restrict the admissible frame
group, or define curvature contraction through a separately specified
wedge splitting and metric pairing.  Ordinary bimodule duality cannot
be retained unchanged.
\end{remark}

\subsection{Norm landing at a certified anchor}

\begin{proposition}[Cofinal landing of the anchor contraction]
\label{prop:v14v2-landing}
Assume the coefficient-extraction maps in (2.2) are bounded.  If
\[
 {\cal R}_n\to{\cal R},\qquad
 G_n\to G,\qquad G_n^{-1}\to G^{-1}                         \tag{2.22}
\]
in the corresponding matrix coefficient norms, then
\[
 P_n\to P,\qquad
 s_{\mathrm J,n}\to s_{\mathrm J},\qquad
 {\cal G}^{\mathrm J}_n\to{\cal G}^{\mathrm J}.              \tag{2.23}
\]
\end{proposition}

\begin{proof}
The Ricci map (2.5) is a finite sum of bounded coefficient maps, so
\(P_n\to P\).  Each term of (2.6)--(2.7) is a two-factor ordered
product of convergent, hence uniformly bounded, sequences.  For such
products,
\[
 \|A_nB_n-AB\|
 \leq\|A_n-A\|\,\|B_n\|+\|A\|\,\|B_n-B\|.
\]
Finite summation proves (2.23).
\end{proof}

\subsection{Acquisition ledger and next upstream step}

\[
\begin{array}{p{0.29\linewidth}|p{0.23\linewidth}|p{0.38\linewidth}}
\textbf{gate}&\textbf{status after V2}&\textbf{content}\\ \hline
\text{anchor Ricci map}
 &\text{acquired}
 &P_{b\nu}=\sum_a\iota_\nu\iota_a{\cal R}^a{}_b\\
\text{scalar ordering}
 &\text{acquired}
 &s_{\mathrm J}=\frac12(s_{\mathrm L}+s_{\mathrm R})\\
\text{reality and classical limit}
 &\text{acquired conditionally}
 &Hermitian Ricci symmetry and commuting reduction\\
\text{central-frame covariance}
 &\text{acquired}
 &exact tensor transformation (2.13)--(2.14)\\
\text{differentiated contraction}
 &\text{reduced to explicit defects}
 &K_\rho=D\iota_\rho+\iota_\rho D\\
\text{noncentral ordinary duality}
 &\text{refuted}
 &rank-one bimodule obstruction\\
\text{noncentral Ricci covariance}
 &\text{open}
 &choose a braided dual or a wedge-splitting calculus
\end{array}                                                     \tag{2.24}
\]

The next version goes upstream to the missing bimodule datum.  It
should define a braided interior contraction, state the exact
relations needed for frame covariance and adjoint compatibility, and
show how two choices of braiding or wedge splitting change the Ricci
tensor.  A construction that is independent only modulo commutators
will not be promoted to an operator-valued Einstein equation without
an additional descent theorem.

\subsection{Parent-integrity record}

The installed parent was
\[
\begin{array}{c|c}
 \text{file}&
 \texttt{Renewal\_SM\_Einstein\_Continuum\_Research\_Notes\_V1.tex}\\
 \text{lines}&552\\
 \text{bytes}&21270\\
 \text{SHA--256}&
 \texttt{654399361B3F641F0E4382D27F67206B237D8BCB4B0FA4748030BBD2D49FAF74}.
\end{array}                                                     \tag{2.25}
\]

% =============================================================================
% END APPENDED FOURTEENTH-CYCLE V2 MATERIAL
% =============================================================================
% =============================================================================
% BEGIN APPENDED FOURTEENTH-CYCLE V3 MATERIAL
% =============================================================================

\section{Fourteenth-cycle V3: braided wedge splittings and the
Ricci-choice defect}
\label{sec:v14v3}

\subsection{Why a lift of two-forms is required}

The no-go theorem in Section~\ref{sec:v14v2} rules out an ordinary
bimodule dual after a genuinely noncentral frame change.  A more
intrinsic route starts from the exterior product
\[
 \wedge:\Omega^1\otimes_{\mathfrak A}\Omega^1
       \longrightarrow\Omega^2.                             \tag{3.1}
\]
Interior contraction of one curvature index requires a lift of a
two-form back to two ordered one-form slots.

\begin{definition}[Wedge splitting]
A wedge splitting is an \(\mathfrak A\)-bimodule map
\[
 j:\Omega^2\longrightarrow
   \Omega^1\otimes_{\mathfrak A}\Omega^1
 \quad\hbox{such that}\quad
 \wedge\circ j=I_{\Omega^2}.                                 \tag{3.2}
\]
Let
\[
 {\cal C}:
 M_m(\Omega^1\otimes_{\mathfrak A}\Omega^1)
 \longrightarrow M_m(\mathfrak A)                           \tag{3.3}
\]
denote a declared bounded contraction containing the frame-index
evaluation and the remaining one-form evaluation.  The corresponding
Ricci candidate is
\[
 P_j({\cal R}):={\cal C}(j{\cal R}),                         \tag{3.4}
\]
where \(j\) acts entrywise on the curvature matrix.
\end{definition}

At the central reference calculus of V2, \(j\) is the ordinary
antisymmetric lift and \({\cal C}\) is normalized so that (3.4)
equals (2.5).

\subsection{Choice dependence}

\begin{theorem}[Exact Ricci splitting defect]
\label{thm:v14v3-choice}
Let \(j_1,j_2\) be two wedge splittings.  Then
\[
 k:=j_2-j_1
 \quad\hbox{satisfies}\quad
 \operatorname{Ran}k\subseteq\ker\wedge,                    \tag{3.5}
\]
and
\[
 \boxed{
 P_{j_2}({\cal R})-P_{j_1}({\cal R})
 ={\cal C}\bigl(k{\cal R}\bigr).}                            \tag{3.6}
\]
Consequently the Ricci contraction is independent of the splitting
for all curvatures if and only if
\[
 {\cal C}\,k=0
 \quad\hbox{for every difference \(k\) of admissible splittings}.
                                                                  \tag{3.7}
\]
In norm,
\[
 \|P_{j_2}({\cal R})-P_{j_1}({\cal R})\|
 \leq\|{\cal C}\|\,\|j_2-j_1\|\,\|{\cal R}\|.                \tag{3.8}
\]
\end{theorem}

\begin{proof}
Subtract \(\wedge j_1=I\) from \(\wedge j_2=I\) to obtain
\(\wedge k=0\), which proves (3.5).  Subtract (3.4) for the two
splittings to obtain (3.6).  Formula (3.7) is exactly the condition
that the right-hand side vanish for every curvature.  Taking norms
gives (3.8).
\end{proof}

\begin{proposition}[The ambiguity is genuine]
\label{prop:v14v3-example}
Let \(\Omega^1\) be the two-dimensional vector space with basis
\(e^1,e^2\), with the ordinary exterior product.  On the generator
\(\xi=e^1\wedge e^2\), define
\[
\begin{aligned}
 j_0(\xi)
 &=\frac12(e^1\otimes e^2-e^2\otimes e^1),\\
 j_\lambda(\xi)
 &=j_0(\xi)+\lambda e^1\otimes e^1.                         \tag{3.9}
\end{aligned}
\]
Both are wedge splittings.  If
\[
 {\cal C}(e^1\otimes e^1)=e^1,\qquad
 {\cal C}(e^1\otimes e^2)=e^2,\qquad
 {\cal C}(e^2\otimes e^1)
 ={\cal C}(e^2\otimes e^2)=0,                               \tag{3.10}
\]
then
\[
 {\cal C}(j_\lambda\xi)-{\cal C}(j_0\xi)
 =\lambda e^1.                                               \tag{3.11}
\]
Thus the exterior two-form alone does not determine its Ricci
contraction.
\end{proposition}

\begin{proof}
Since \(e^1\wedge e^1=0\), (3.9) gives
\(\wedge j_\lambda(\xi)=\xi\).  Equation (3.11) follows directly
from (3.10).
\end{proof}

The example is intentionally elementary.  Classical geometry removes
this ambiguity by the canonical flip and antisymmetrizer.  In a
noncommutative bimodule there need not be a well-defined ordinary
flip, so its replacement must be part of the data.

\subsection{An involutive braided selector}

\begin{definition}[Admissible involutive braiding]
An admissible braiding is an \(\mathfrak A\)-bimodule automorphism
\[
 \sigma:\Omega^1\otimes_{\mathfrak A}\Omega^1
 \longrightarrow\Omega^1\otimes_{\mathfrak A}\Omega^1
\]
such that
\[
 \sigma^2=I,\qquad
 \wedge\sigma=-\wedge.                                      \tag{3.12}
\]
Set
\[
 A_\sigma:=\frac12(I-\sigma),\qquad
 {\cal T}_\sigma:=\operatorname{Ran}A_\sigma.                \tag{3.13}
\]
We call \(\sigma\) wedge-complete if
\[
 W_\sigma:=\wedge|_{{\cal T}_\sigma}:
 {\cal T}_\sigma\longrightarrow\Omega^2                     \tag{3.14}
\]
is a bimodule isomorphism.
\end{definition}

\begin{lemma}[Braided antisymmetric projection]
\label{lem:v14v3-projection}
Under (3.12),
\[
 A_\sigma^2=A_\sigma,\qquad
 \wedge A_\sigma=\wedge.                                    \tag{3.15}
\]
\end{lemma}

\begin{proof}
The first identity follows from \(\sigma^2=I\).  The second follows
from
\[
 \wedge A_\sigma
 =\frac12(\wedge-\wedge\sigma)
 =\frac12(\wedge+\wedge)=\wedge.
\]
\end{proof}

\begin{theorem}[Canonical splitting relative to a braiding]
\label{thm:v14v3-braidedsplitting}
If \(\sigma\) is wedge-complete, then
\[
 j_\sigma:=W_\sigma^{-1}:
 \Omega^2\longrightarrow{\cal T}_\sigma
 \subseteq\Omega^1\otimes_{\mathfrak A}\Omega^1             \tag{3.16}
\]
is the unique wedge splitting whose range lies in
\({\cal T}_\sigma\).  The associated Ricci candidate
\[
 P_\sigma({\cal R}):={\cal C}(j_\sigma{\cal R})              \tag{3.17}
\]
therefore has no remaining lift ambiguity after
\((\sigma,{\cal C})\) have been fixed.
\end{theorem}

\begin{proof}
By definition of \(W_\sigma\),
\(\wedge j_\sigma=I_{\Omega^2}\), so \(j_\sigma\) is a splitting.
If \(j\) is another splitting with range in \({\cal T}_\sigma\),
then
\[
 W_\sigma j=\wedge j=I.
\]
Uniqueness of the inverse of \(W_\sigma\) gives \(j=j_\sigma\).
Formula (3.17) is then determined by the declared contraction.
\end{proof}

\begin{remark}[Braid relation]
For tensor products of three or more one-forms, a full braided
exterior calculus normally also requires the Yang--Baxter relation
\[
 (\sigma\otimes I)(I\otimes\sigma)(\sigma\otimes I)
 =(I\otimes\sigma)(\sigma\otimes I)(I\otimes\sigma).
\]
The two-form splitting theorem uses only (3.12)--(3.14).  Higher-form
Bianchi and product identities must add the braid relation rather
than infer it from involutivity.
\end{remark}

\subsection{Covariance of the selected splitting}

Let \(U:\Omega^1\to\Omega^{1\prime}\) be a bimodule isomorphism and
let \(U^{(2)}:\Omega^2\to\Omega^{2\prime}\) satisfy
\[
 U^{(2)}\wedge
 =\wedge'(U\otimes U).                                      \tag{3.18}
\]
Assume the braidings obey
\[
 \sigma'(U\otimes U)=(U\otimes U)\sigma.                    \tag{3.19}
\]

\begin{theorem}[Braided frame covariance]
\label{thm:v14v3-covariance}
If both braidings are wedge-complete, then
\[
 j_{\sigma'}U^{(2)}
 =(U\otimes U)j_\sigma.                                     \tag{3.20}
\]
If the curvature and contraction transform by
\[
 {\cal R}'=U^{(2)}{\cal R},\qquad
 {\cal C}'(U\otimes U)={\cal U}_{\mathrm{Ric}}{\cal C},      \tag{3.21}
\]
then
\[
 P_{\sigma'}({\cal R}')
 ={\cal U}_{\mathrm{Ric}}P_\sigma({\cal R}).                \tag{3.22}
\]
\end{theorem}

\begin{proof}
Equation (3.19) sends
\({\cal T}_\sigma\) onto \({\cal T}_{\sigma'}\).  By (3.18),
\[
 \wedge'(U\otimes U)j_\sigma
 =U^{(2)}\wedge j_\sigma=U^{(2)}.
\]
Hence \((U\otimes U)j_\sigma\) is a splitting of \(U^{(2)}\) with
range in \({\cal T}_{\sigma'}\).  Uniqueness in
Theorem~\ref{thm:v14v3-braidedsplitting} proves (3.20).  Apply
\({\cal C}'\) and use (3.21) to obtain (3.22).
\end{proof}

This theorem states the exact covariance data.  It does not assert
that left multiplication by the noncentral coframe matrix is a
bimodule map; V2 proves that such an assertion can fail.

\subsection{Adjoint compatibility}

Assume the tensor product and two-form spaces carry involutions, and
write \(X^\dagger\) for the tensor involution.  Suppose
\[
 \sigma(X^\dagger)=\sigma(X)^\dagger,\qquad
 \wedge(X^\dagger)=\wedge(X)^*,                             \tag{3.23}
\]
and
\[
 {\cal C}(X^\dagger)={\cal C}(X)^*.                          \tag{3.24}
\]

\begin{proposition}[Reality of the braided contraction]
\label{prop:v14v3-reality}
Under (3.23), the selected splitting satisfies
\[
 j_\sigma(\alpha^*)=j_\sigma(\alpha)^\dagger.                \tag{3.25}
\]
Consequently, if the lifted curvature has the declared Hermitian
tensor symmetry, then \(P_\sigma\) has the Hermitian Ricci symmetry
required in (2.8).
\end{proposition}

\begin{proof}
Conditions (3.23) preserve \({\cal T}_\sigma\).  Both
\(j_\sigma(\alpha^*)\) and \(j_\sigma(\alpha)^\dagger\) lie in that
space, and applying \(\wedge\) gives \(\alpha^*\) in both cases.
Uniqueness in Theorem~\ref{thm:v14v3-braidedsplitting} proves
(3.25).  Equation (3.24) then transports the curvature symmetry to
the Ricci output.
\end{proof}

\subsection{Exact differentiated-splitting defect}

Let \(D_2\), \(D_\otimes\), and \(D_P\) be the covariant
differentials on \(\Omega^2\), the ordered two-one-form tensor space,
and the Ricci target.  Define
\[
\begin{aligned}
 {\cal J}_{j_\sigma}(\alpha)
 &:=D_\otimes j_\sigma(\alpha)-j_\sigma D_2\alpha,\\
 {\cal J}_{\cal C}(X)
 &:=D_P{\cal C}(X)-{\cal C}D_\otimes X.                     \tag{3.26}
\end{aligned}
\]

\begin{theorem}[Factorization of the Ricci intertwining defect]
\label{thm:v14v3-intertwining}
For \(P_\sigma={\cal C}j_\sigma\),
\[
\boxed{
 D_PP_\sigma({\cal R})-P_\sigma(D_2{\cal R})
 =
 {\cal J}_{\cal C}(j_\sigma{\cal R})
 +{\cal C}\bigl({\cal J}_{j_\sigma}({\cal R})\bigr).}        \tag{3.27}
\]
In norm,
\[
\begin{aligned}
 \|D_PP_\sigma({\cal R})-P_\sigma(D_2{\cal R})\|
 \leq{}&
 \|{\cal J}_{\cal C}(j_\sigma{\cal R})\|\\
 &+\|{\cal C}\|\,
   \|{\cal J}_{j_\sigma}({\cal R})\|.                        \tag{3.28}
\end{aligned}
\]
\end{theorem}

\begin{proof}
Add and subtract
\({\cal C}D_\otimes j_\sigma({\cal R})\):
\[
\begin{aligned}
 D_P{\cal C}j_\sigma({\cal R})
 -{\cal C}j_\sigma D_2{\cal R}
 ={}&
 \bigl(D_P{\cal C}-{\cal C}D_\otimes\bigr)
     j_\sigma({\cal R})\\
 &+{\cal C}
   \bigl(D_\otimes j_\sigma-j_\sigma D_2\bigr)({\cal R}).
\end{aligned}
\]
This is (3.27); the norm estimate follows by the triangle inequality.
\end{proof}

Equation (3.27) replaces one opaque Ricci defect by two geometric
ones: failure of the chosen contraction to intertwine the connection,
and failure of the braided lift to do so.  These are the quantities
that must enter the contracted Bianchi ledger.

\subsection{Cofinal landing}

\begin{theorem}[Landing under a uniform wedge inverse]
\label{thm:v14v3-landing}
Suppose finite-stage braided antisymmetric spaces have been identified
with common Banach spaces so that
\[
 W_{\sigma_n}\to W_\sigma,\qquad
 \sup_n\|W_{\sigma_n}^{-1}\|\leq M,\qquad
 W_\sigma^{-1}\ \hbox{exists}.                              \tag{3.29}
\]
Then
\[
 j_{\sigma_n}=W_{\sigma_n}^{-1}
 \longrightarrow W_\sigma^{-1}=j_\sigma.                   \tag{3.30}
\]
If also
\[
 {\cal C}_n\to{\cal C},\qquad
 {\cal R}_n\to{\cal R}                                      \tag{3.31}
\]
in operator and curvature norms, then
\[
 P_{\sigma_n}({\cal R}_n)
 \longrightarrow P_\sigma({\cal R}).                        \tag{3.32}
\]
\end{theorem}

\begin{proof}
The inverse identity
\[
 W_{\sigma_n}^{-1}-W_\sigma^{-1}
 =
 W_{\sigma_n}^{-1}(W_\sigma-W_{\sigma_n})W_\sigma^{-1}
\]
proves (3.30).  For (3.32), add and subtract the intermediate terms
\({\cal C}_nj_{\sigma_n}{\cal R}\) and
\({\cal C}_nj_\sigma{\cal R}\); uniform boundedness and
(3.30)--(3.31) make the three resulting differences tend to zero.
\end{proof}

\subsection{Acquisition ledger and next construction}

\[
\begin{array}{p{0.29\linewidth}|p{0.23\linewidth}|p{0.38\linewidth}}
\textbf{gate}&\textbf{status after V3}&\textbf{content}\\ \hline
\text{splitting dependence}
 &\text{exactly measured}
 &P_{j_2}-P_{j_1}={\cal C}(j_2-j_1){\cal R}\\
\text{canonical lift relative to \(\sigma\)}
 &\text{acquired}
 &j_\sigma=(\wedge|_{{\cal T}_\sigma})^{-1}\\
\text{braided frame covariance}
 &\text{acquired conditionally}
 &intertwining relations (3.18)--(3.21)\\
\text{adjoint compatibility}
 &\text{acquired conditionally}
 &star-preserving braiding, wedge, and contraction\\
\text{Ricci differentiation defect}
 &\text{factorized}
 &{\cal J}_{\cal C}+{\cal C}{\cal J}_{j_\sigma}\\
\text{cofinal Ricci landing}
 &\text{acquired conditionally}
 &uniform inverse margin for \(W_{\sigma_n}\)\\
\text{physical braiding}
 &\text{open}
 &construct from the regulator's differential calculus
\end{array}                                                     \tag{3.33}
\]

The next step is to test whether the noncentral coframe can transport
the central flip as a bimodule braiding.  If the transport fails
right-linearity, its defect must be computed and either absorbed into
a twisted tensor product or shown to vanish on the physical
observable subalgebra.  This upstream test precedes any claim of a
noncentral Ricci tensor.

\subsection{Parent-integrity record}

The installed parent was
\[
\begin{array}{c|c}
 \text{file}&
 \texttt{Renewal\_SM\_Einstein\_Continuum\_Research\_Notes\_V2.tex}\\
 \text{lines}&981\\
 \text{bytes}&34733\\
 \text{SHA--256}&
 \texttt{C2EC5BC03EF897112346D9C0E2B27DEB3B97FCF178231973B06CAB83180F9F2E}.
\end{array}                                                     \tag{3.34}
\]

% =============================================================================
% END APPENDED FOURTEENTH-CYCLE V3 MATERIAL
% =============================================================================
% =============================================================================
% BEGIN APPENDED FOURTEENTH-CYCLE V4 MATERIAL
% =============================================================================

\section{Fourteenth-cycle V4: coframe transport, the standard-action
no-go, and a twisted-bimodule repair}
\label{sec:v14v4}

\subsection{Failure on the unchanged bimodule}

Let
\[
 {\cal M}_0=\mathfrak A^m
\]
with componentwise right action and standard left action
\(\lambda_0(a)\xi=(aI_m)\xi\).  For
\(E\in GL_m(\mathfrak A)\), define
\[
 T_E:{\cal M}_0\longrightarrow{\cal M}_0,\qquad
 T_E\xi=E\xi.                                                \tag{4.1}
\]

\begin{theorem}[Standard-action transport obstruction]
\label{thm:v14v4-standardnogo}
The map \(T_E\) is always a right \(\mathfrak A\)-module
isomorphism.  It is a left \(\mathfrak A\)-module map for the
standard action if and only if every entry of \(E\) lies in
\(Z(\mathfrak A)\).  Therefore a genuinely noncentral coframe cannot
transport the central braiding by a tensor square
\(T_E\otimes_{\mathfrak A}T_E\) while the original bimodule structure
is kept unchanged.
\end{theorem}

\begin{proof}
For \(b\in\mathfrak A\),
\[
 T_E(\xi b)=E(\xi b)=(E\xi)b,
\]
so \(T_E\) is right-linear; its inverse is \(T_{E^{-1}}\).
Left-linearity is equivalent to
\[
 E(aI_m)\xi=(aI_m)E\xi
 \quad\hbox{for every \(a\) and \(\xi\)}.                    \tag{4.2}
\]
Testing (4.2) on the standard basis gives
\(E_{ij}a=aE_{ij}\) for every \(i,j,a\).  This is exactly
\(E_{ij}\in Z(\mathfrak A)\).

To define the tensor square over \(\mathfrak A\), the first copy must
respect the right action and the second copy must respect the left
action in the balancing relation.  The second property fails when
(4.2) fails.
\end{proof}

This theorem sharpens the rank-one dual-frame obstruction in V2: the
problem occurs already before contraction, at the balanced tensor
product used to define two-forms.

\subsection{Transporting the left representation}

The obstruction has an exact algebraic repair if the physical
calculus permits the left action to be transported.

\begin{definition}[Coframe-twisted bimodule]
Keep the underlying right module \(\mathfrak A^m\), and define
\[
 \lambda_E(a):=E(aI_m)E^{-1}\in M_m(\mathfrak A),\qquad
 a\cdot_E\xi:=\lambda_E(a)\xi.                               \tag{4.3}
\]
Denote the resulting bimodule by \({\cal M}_E\).
\end{definition}

\begin{theorem}[Exact bimodule transport]
\label{thm:v14v4-bimodule}
The map \(a\mapsto\lambda_E(a)\) is a unital algebra
representation by right-module endomorphisms.  The map
\[
 T_E:{\cal M}_0\longrightarrow{\cal M}_E                   \tag{4.4}
\]
is an \(\mathfrak A\)-bimodule isomorphism.
\end{theorem}

\begin{proof}
Using \(E^{-1}E=I_m\),
\[
 \lambda_E(a)\lambda_E(b)
 =E(aI_m)E^{-1}E(bI_m)E^{-1}
 =E(abI_m)E^{-1}
 =\lambda_E(ab),
\]
and \(\lambda_E(1)=I_m\).  A matrix over \(\mathfrak A\) acts by a
right-module endomorphism because matrix left multiplication
commutes with componentwise right multiplication.

Right-linearity of \(T_E\) was proved in
Theorem~\ref{thm:v14v4-standardnogo}.  For the left action,
\[
 T_E\lambda_0(a)\xi
 =E(aI_m)\xi
 =E(aI_m)E^{-1}E\xi
 =\lambda_E(a)T_E\xi.
\]
The inverse \(T_{E^{-1}}\) has the corresponding intertwining
property.
\end{proof}

\subsection{Transported braiding and wedge splitting}

Assume the reference bimodule \({\cal M}_0\) has an involutive
wedge-complete braiding \(\sigma_0\), a two-form module
\(\Omega_0^2\), a wedge map \(\wedge_0\), and selected splitting
\(j_0\).  Transport \(\Omega_0^2\) to a bimodule \(\Omega_E^2\) by a
chosen bimodule isomorphism
\[
 T_E^{(2)}:\Omega_0^2\longrightarrow\Omega_E^2.              \tag{4.5}
\]
Since (4.4) is a bimodule isomorphism, the balanced tensor map
\[
 S_E:=T_E\otimes_{\mathfrak A}T_E                           \tag{4.6}
\]
is well defined.  Define
\[
\begin{aligned}
 \sigma_E&:=S_E\sigma_0S_E^{-1},\\
 \wedge_E&:=T_E^{(2)}\wedge_0S_E^{-1},\\
 j_E&:=S_Ej_0(T_E^{(2)})^{-1}.                              \tag{4.7}
\end{aligned}
\]

\begin{theorem}[Transported braided calculus]
\label{thm:v14v4-transport}
The maps in (4.7) satisfy
\[
 \sigma_E^2=I,\qquad
 \wedge_E\sigma_E=-\wedge_E,\qquad
 \wedge_Ej_E=I.                                              \tag{4.8}
\]
If \(\sigma_0\) satisfies the braid relation, then so does
\(\sigma_E\).  If \(\sigma_0\) is wedge-complete, then
\(\sigma_E\) is wedge-complete and \(j_E\) is its unique selected
splitting.
\end{theorem}

\begin{proof}
Conjugation gives \(\sigma_E^2=I\).  The definitions and
\(\wedge_0\sigma_0=-\wedge_0\) give
\[
 \wedge_E\sigma_E
 =T_E^{(2)}\wedge_0\sigma_0S_E^{-1}
 =-T_E^{(2)}\wedge_0S_E^{-1}
 =-\wedge_E.
\]
Likewise,
\[
 \wedge_Ej_E
 =T_E^{(2)}\wedge_0j_0(T_E^{(2)})^{-1}=I.
\]
Conjugating each factor in the threefold tensor product transports
the braid relation.  Finally, \(S_E\) maps the reference
antisymmetric range bijectively onto the \(\sigma_E\)-antisymmetric
range, while (4.5) transports the restricted wedge isomorphism.
Wedge-completeness and uniqueness follow from
Theorem~\ref{thm:v14v3-braidedsplitting}.
\end{proof}

\begin{corollary}[Ricci covariance in the twisted category]
\label{cor:v14v4-ricci}
If the curvature and contraction are transported by
\[
 {\cal R}_E=T_E^{(2)}{\cal R}_0,\qquad
 {\cal C}_ES_E={\cal U}_{\mathrm{Ric}}{\cal C}_0,            \tag{4.9}
\]
then
\[
 P_E({\cal R}_E)
 :={\cal C}_Ej_E{\cal R}_E
 ={\cal U}_{\mathrm{Ric}}P_0({\cal R}_0).                   \tag{4.10}
\]
\end{corollary}

\begin{proof}
Insert (4.7) and (4.9):
\[
 {\cal C}_ES_Ej_0(T_E^{(2)})^{-1}T_E^{(2)}{\cal R}_0
 ={\cal U}_{\mathrm{Ric}}{\cal C}_0j_0{\cal R}_0.
\]
\end{proof}

Equation (4.10) proves exact covariance inside the transported
bimodule category.  It does not identify \({\cal M}_E\) with a
pre-existing physical one-form bimodule whose left action may differ
from (4.3).

\subsection{The transported differential}

Let \(d\) act entrywise and set
\[
 F:=E^{-1},\qquad
 \Gamma_0:=F\,dE.                                            \tag{4.11}
\]
For a degree-zero coefficient \(a\), define
\[
 \nabla_{\Gamma_0}^{\mathrm{ad}}a
 :=da+\Gamma_0(aI_m)-(aI_m)\Gamma_0.                         \tag{4.12}
\]

\begin{theorem}[Pure-gauge differential intertwining]
\label{thm:v14v4-differential}
The twisted representation satisfies
\[
 \boxed{
 d\lambda_E(a)
 =E\bigl(\nabla_{\Gamma_0}^{\mathrm{ad}}a\bigr)F.}           \tag{4.13}
\]
Thus the ordinary differential of the transported left action
corresponds exactly to the adjoint covariant differential at the
reference side.
\end{theorem}

\begin{proof}
From \(FE=I_m\),
\[
 dF=-F(dE)F=-\Gamma_0F,\qquad dE=E\Gamma_0.                 \tag{4.14}
\]
Differentiate \(\lambda_E(a)=E(aI_m)F\):
\[
\begin{aligned}
 d\lambda_E(a)
 &=(dE)(aI_m)F+E(da\,I_m)F+E(aI_m)dF\\
 &=E\left[
   \Gamma_0(aI_m)+da\,I_m-(aI_m)\Gamma_0
   \right]F,
\end{aligned}
\]
which is (4.13).
\end{proof}

The connection \(\Gamma_0\) in (4.11) is the flat
metric-compatible connection identified in the frozen V98 theorem.
Nonzero curvature is carried by the additional metric-skew correction
selected by torsion.  Equation (4.13) therefore supplies the
differential transport baseline, not the full gravitational
connection.

\subsection{Cofinal stability}

Let the module maps carry compatible operator norms and assume the
balanced tensor norm is a crossnorm.

\begin{lemma}[Tensor transport bounds]
\label{lem:v14v4-tensorbounds}
For two invertible coframes \(E,\widetilde E\),
\[
\begin{aligned}
 \|S_E-S_{\widetilde E}\|
 &\leq
 \|T_E-T_{\widetilde E}\|
 \bigl(\|T_E\|+\|T_{\widetilde E}\|\bigr),\\
 \|S_E^{-1}-S_{\widetilde E}^{-1}\|
 &\leq
 \|T_E^{-1}-T_{\widetilde E}^{-1}\|
 \bigl(\|T_E^{-1}\|+\|T_{\widetilde E}^{-1}\|\bigr).
                                                                  \tag{4.15}
\end{aligned}
\]
\end{lemma}

\begin{proof}
Use
\[
 T_E\otimes T_E-T_{\widetilde E}\otimes T_{\widetilde E}
 =(T_E-T_{\widetilde E})\otimes T_E
 +T_{\widetilde E}\otimes(T_E-T_{\widetilde E})
\]
and the crossnorm estimate.  Apply the same identity to the inverses.
\end{proof}

\begin{theorem}[Landing of the transported braiding]
\label{thm:v14v4-landing}
Suppose
\[
 E_n\to E,\qquad E_n^{-1}\to E^{-1}                         \tag{4.16}
\]
in matrix norm, and the maps \(T_{E_n}^{(2)}\) and their inverses
converge to \(T_E^{(2)}\) and its inverse.  Then
\[
 \lambda_{E_n}(a)\to\lambda_E(a)
 \quad\hbox{for every fixed \(a\)},                          \tag{4.17}
\]
and
\[
 \sigma_{E_n}\to\sigma_E,\qquad
 \wedge_{E_n}\to\wedge_E,\qquad
 j_{E_n}\to j_E                                             \tag{4.18}
\]
in operator norm.
\end{theorem}

\begin{proof}
Equation (4.17) follows from the three-factor product
\(E_n(aI_m)E_n^{-1}\).  Lemma~\ref{lem:v14v4-tensorbounds} gives
\(S_{E_n}\to S_E\) and convergence of the inverses.  Each expression
in (4.7) is then a finite product of convergent operator sequences,
which proves (4.18).
\end{proof}

\subsection{Physical-action defect and claim boundary}

Let \(\lambda_{\mathrm{phys}}\) be the left action supplied by the
actual regulator's one-form calculus.  Define
\[
 {\cal D}_E(a)
 :=\lambda_{\mathrm{phys}}(a)-E(aI_m)E^{-1}.                 \tag{4.19}
\]
The transported calculus is the physical calculus exactly when this
defect vanishes, together with the corresponding wedge and
differential identifications.  Smallness of \({\cal D}_E\) alone
does not prove that \(\lambda_{\mathrm{phys}}\) is conjugate to
\(\lambda_E\); representation stability is an additional theorem.

\[
\begin{array}{p{0.29\linewidth}|p{0.23\linewidth}|p{0.38\linewidth}}
\textbf{gate}&\textbf{status after V4}&\textbf{content}\\ \hline
\text{transport on unchanged bimodule}
 &\text{refuted for noncentral \(E\)}
 &T_E\text{ fails left-linearity}\\
\text{twisted left representation}
 &\text{acquired}
 &\lambda_E(a)=E(aI_m)E^{-1}\\
\text{transported braiding and split}
 &\text{acquired}
 &formulas (4.7)--(4.8)\\
\text{twisted-category Ricci covariance}
 &\text{acquired conditionally}
 &transported curvature and contraction\\
\text{differential transport}
 &\text{acquired for \(\Gamma_0\)}
 &identity (4.13)\\
\text{cofinal landing}
 &\text{acquired conditionally}
 &E_n,F_n,T_{E_n}^{(2)}\text{ converge}\\
\text{physical calculus identification}
 &\text{open}
 &vanishing or controlled descent of \({\cal D}_E\)
\end{array}                                                     \tag{4.20}
\]

The V5 boundary is the required companion audit.  After that audit,
the next substantive version should glue the local twisted modules
and their braided Ricci contractions across the finite selector cover.
The transition maps must satisfy an exact cocycle and preserve the
chosen wedge splittings; otherwise local Ricci tensors do not define
a global observable.

\subsection{Parent-integrity record}

The installed parent was
\[
\begin{array}{c|c}
 \text{file}&
 \texttt{Renewal\_SM\_Einstein\_Continuum\_Research\_Notes\_V3.tex}\\
 \text{lines}&1426\\
 \text{bytes}&48384\\
 \text{SHA--256}&
 \texttt{5330B6A6E384474DFCE4A4EE47EEF63D28CA3520B4502C76BC300D7A2A179DF8}.
\end{array}                                                     \tag{4.21}
\]

% =============================================================================
% END APPENDED FOURTEENTH-CYCLE V4 MATERIAL
% =============================================================================
% =============================================================================
% BEGIN APPENDED FOURTEENTH-CYCLE V5 MATERIAL
% =============================================================================

\section{Fourteenth-cycle V5: compulsory companion audit and the
constraint-before-gluing decision}
\label{sec:v14v5}

\subsection{Files inspected}

The mandatory five-version audit was performed on 10 September 2026
after installation of V4.  The newest active companion records at
that time were
\[
\begin{array}{p{0.39\linewidth}|r|r|p{0.27\linewidth}}
\textbf{file}&\textbf{lines}&\textbf{bytes}&\textbf{SHA--256}\\ \hline
\texttt{Renewal\_Yang\_Mills\_Mass\_Gap\_Research\_Notes\_V100.tex}
 &32206&1100952&
 \texttt{1FD1A60C0E15E0A1ECA28FD3F372398A3CFA2F5C74A39AD38DEC7C718B9DA36C}\\
\texttt{Renewal\_NS\_Stability\_Research\_Notes\_V26.tex}
 &5319&278283&
 \texttt{82C887F8C2C69E4F28FAD7C3DC8DE5B9099CB5A4BF431EBA1FFF3E91935978AD}.
\end{array}                                                     \tag{5.1}
\]
No companion file was modified.

\subsection{New Yang--Mills information}

The YM V98--V100 chain adds a result not present at the preceding SM
audit.

At V98, each uncompressed fine curl block is positive semidefinite
but has a momentum-dependent gauge kernel.  Differentiating that
kernel produces a nonzero cross term, and an exact rational vector
makes both signs of the independent derivative pencil indefinite.
The negative vector does not belong to the coupled physical range.
Thus imposing positivity before the vertical constraint is
mathematically false.

At V99, the full alias object is reconstructed on the fixed
\(45\)-dimensional constrained fibre, and the flat frame connection is
removed before differentiation.  Every one of the \(1024\) tested
covariant pencils has an exact positive rational margin.  This proves
a finite-grid constrained statement, not a torus theorem.

At V100, exact connection removal lowers the Laurent coefficient
stock and enlarges the certified local radii.  The improvement is
real, but the scalar common balls still miss the torus by roughly
three orders of magnitude.  The YM programme therefore moves to the
constraint-retaining relative Schur object rather than returning to
absolute coefficient radii.

\begin{proposition}[Transferable YM principle]
\label{prop:v14v5-ymtransfer}
For the SM--Einstein gluing step, the physical submodule and its
connection must be transported before any positivity, inverse-margin,
or contraction estimate is imposed.  Independent estimates on
ambient local modules cannot replace an estimate on the glued
physical range when that range moves with the coframe.
\end{proposition}

\begin{proof}
The YM moving-kernel lemma supplies an exact counterexample to the
opposite order of operations: each ambient block has a null direction,
its derivative couples that direction to the positive range, and both
signed derivative pencils become indefinite.  Compression to the
actual coupled range removes the inadmissible negative vector.
V4 of the present notes similarly has coframe-dependent twisted left
actions and braidings.  Hence an ambient estimate that ignores the
transition maps can test directions that do not define physical
sections.  Transporting the module and connection first is the
logically valid order.
\end{proof}

\subsection{Navier--Stokes status}

The NS note remains byte-identical to the source inspected in the
preceding audits.  Its current endpoint computes rank-one Oseen-kernel
derivative norms: the first-derivative restriction gives no gain and
the second gives a modest localized gain.  This continues to support
realized-input estimates, but supplies no gluing, braiding, Ricci, or
Einstein constant.

\subsection{Direction after the audit}

\begin{decision}[V6 glues the constrained twisted calculi]
\label{dec:v14v5-next}
Let \(\{U_i\}\) be a finite selector cover and let \(E_i\) be its
local coframes.  V6 shall:
\begin{enumerate}
\item construct the transition maps
      \(T_{ij}=T_{E_j}T_{E_i}^{-1}\) between the twisted bimodules;
\item prove the exact identity, inverse, and triple-overlap cocycle
      laws;
\item transport braidings, wedge splittings, connections, curvatures,
      and Ricci contractions through those transitions;
\item define the physical section module by the overlap equations
      before taking inverse or positivity bounds;
\item prove a uniform cofinal gluing theorem under transition and
      local-object convergence.
\end{enumerate}
Any overlap failure will be recorded as a cocycle or intertwining
defect.  It will not be hidden inside a componentwise norm.
\end{decision}

\begin{proof}
The exact coframe transport theorem of V4 supplies the local
bimodule isomorphisms.  Proposition~\ref{prop:v14v5-ymtransfer}
requires their constrained equalizer to be formed before estimates.
The listed items are therefore the next unresolved interface and do
not duplicate V1--V4.
\end{proof}

\subsection{Claim boundary and parent record}

This audit proves no new companion theorem and imports no companion
numerical constant.  It changes the SM proof order using an exact
structural counterexample.  The physical coframe cover, overlap
cocycle, global Ricci section, contracted Bianchi identity, joint
Einstein support, and full SM--Einstein continuum remain open.

The installed parent was
\[
\begin{array}{c|c}
 \text{file}&
 \texttt{Renewal\_SM\_Einstein\_Continuum\_Research\_Notes\_V4.tex}\\
 \text{lines}&1797\\
 \text{bytes}&59840\\
 \text{SHA--256}&
 \texttt{7E1EF6B6D12CD15208C839BB984A8DE565C2D21653A36B40C4C43C97756B1ACA}.
\end{array}                                                     \tag{5.2}
\]

% =============================================================================
% END APPENDED FOURTEENTH-CYCLE V5 MATERIAL
% =============================================================================
% =============================================================================
% BEGIN APPENDED FOURTEENTH-CYCLE V6 MATERIAL
% =============================================================================

\section{Fourteenth-cycle V6: descent of the twisted calculi and
global Ricci sections}
\label{sec:v14v6}

\subsection{Local coframes and transition maps}

Let \(X\) be the compact physical parameter region of the selector
cover in Theorem~\ref{thm:v14v1-cover}, and let
\(\{U_i\}_{i=1}^N\) be a finite open cover.  All identities below are
understood after restriction to the indicated overlaps.  Let
\[
 E_i\in GL_m(\mathfrak A(U_i)),\qquad
 T_i:=T_{E_i},\qquad
 \lambda_i(a)=E_i(aI_m)E_i^{-1}.                            \tag{6.1}
\]
Denote the corresponding twisted bimodule by \({\cal M}_i\).
On \(U_i\cap U_j\), define
\[
 \Phi_{ij}:=T_jT_i^{-1}:{\cal M}_i\longrightarrow{\cal M}_j,
 \qquad
 \Phi_{ij}\xi=E_jE_i^{-1}\xi.                               \tag{6.2}
\]

\begin{theorem}[Exact transition groupoid]
\label{thm:v14v6-cocycle}
Every \(\Phi_{ij}\) is a bimodule isomorphism and
\[
 \Phi_{ii}=I,\qquad
 \Phi_{ji}=\Phi_{ij}^{-1},\qquad
 \Phi_{jk}\Phi_{ij}=\Phi_{ik}                               \tag{6.3}
\]
on double and triple overlaps.  No commutativity of the coframe
entries is required.
\end{theorem}

\begin{proof}
Each \(T_i:{\cal M}_0\to{\cal M}_i\) is a bimodule isomorphism by
Theorem~\ref{thm:v14v4-bimodule}; hence so is \(T_jT_i^{-1}\).
The first two identities in (6.3) are immediate.  For the third,
\[
 \Phi_{jk}\Phi_{ij}
 =T_kT_j^{-1}T_jT_i^{-1}
 =T_kT_i^{-1}
 =\Phi_{ik}.
\]
The order of the factors is retained.
\end{proof}

\begin{definition}[Physical descended section module]
Define
\[
 \Gamma_{\Phi}({\cal M})
 :=
 \left\{
  (\xi_i)_{i=1}^N:
  \xi_j=\Phi_{ij}\xi_i
  \ \hbox{on every }U_i\cap U_j
 \right\}.                                                   \tag{6.4}
\]
This is the overlap equalizer on which subsequent operators must act.
\end{definition}

\begin{proposition}[Closedness and reference trivialization]
\label{prop:v14v6-sections}
With the maximum of the local Banach-module norms,
\(\Gamma_\Phi({\cal M})\) is a closed submodule of the finite product
of local modules.  Every reference section \(\zeta\) produces the
exactly glued family
\[
 \xi_i=T_i\zeta.                                             \tag{6.5}
\]
If all local coframes are expressed relative to the same reference
module, (6.5) identifies the descended module with that reference
section module.
\end{proposition}

\begin{proof}
For each overlap, the map
\[
 (\xi_\ell)_\ell\longmapsto \xi_j-\Phi_{ij}\xi_i
\]
is bounded and linear.  The equalizer (6.4) is the finite
intersection of its kernels and is therefore closed.  Equation (6.5)
glues because
\[
 \Phi_{ij}T_i\zeta=T_jT_i^{-1}T_i\zeta=T_j\zeta.
\]
The inverse identification is \(T_i^{-1}\xi_i\), independent of \(i\)
by (6.4).
\end{proof}

\subsection{Descent of braiding, wedge, and splitting}

Put
\[
 S_{ij}:=\Phi_{ij}\otimes_{\mathfrak A}\Phi_{ij},\qquad
 \Phi^{(2)}_{ij}:=T_j^{(2)}(T_i^{(2)})^{-1}.                 \tag{6.6}
\]
Let \(\sigma_i,\wedge_i,j_i\) be the transported data (4.7).

\begin{theorem}[Braided exterior descent]
\label{thm:v14v6-braideddescent}
On every overlap,
\[
\begin{aligned}
 S_{ij}\sigma_i&=\sigma_jS_{ij},\\
 \wedge_jS_{ij}&=\Phi^{(2)}_{ij}\wedge_i,\\
 j_j\Phi^{(2)}_{ij}&=S_{ij}j_i.                             \tag{6.7}
\end{aligned}
\]
The tensor and two-form transitions satisfy the same identity,
inverse, and triple-overlap cocycle laws as (6.3).  Hence the local
braided exterior calculi descend to the overlap equalizers.
\end{theorem}

\begin{proof}
Using \(S_i=T_i\otimes T_i\), one has
\(S_{ij}=S_jS_i^{-1}\).  Therefore
\[
 S_{ij}\sigma_i
 =S_jS_i^{-1}(S_i\sigma_0S_i^{-1})
 =S_j\sigma_0S_i^{-1}
 =\sigma_jS_{ij}.
\]
Insert the transported formulas
\[
 \wedge_i=T_i^{(2)}\wedge_0S_i^{-1},
 \qquad
 j_i=S_ij_0(T_i^{(2)})^{-1}
\]
to obtain the second and third identities in (6.7).
The cocycle statements follow by the cancellation proof used in
Theorem~\ref{thm:v14v6-cocycle}.
\end{proof}

\subsection{Connection and curvature overlap defects}

Let
\[
 \nabla_i:{\cal M}_i^\bullet\longrightarrow
           {\cal M}_i^{\bullet+1}
\]
be local degree-one connections, and use the induced transition on
the degree-shifted target without changing notation.  Define the
connection overlap defect
\[
 {\cal A}_{ij}
 :=\nabla_j\Phi_{ij}-\Phi_{ij}\nabla_i.                     \tag{6.8}
\]

\begin{theorem}[Connection-defect cocycle]
\label{thm:v14v6-connectiondefect}
On every triple overlap,
\[
 {\cal A}_{ik}
 ={\cal A}_{jk}\Phi_{ij}+\Phi_{jk}{\cal A}_{ij}.             \tag{6.9}
\]
In particular, the local connections define a connection on
\(\Gamma_\Phi({\cal M})\) if and only if
\[
 {\cal A}_{ij}=0
 \quad\hbox{on every overlap}.                              \tag{6.10}
\]
\end{theorem}

\begin{proof}
Use \(\Phi_{ik}=\Phi_{jk}\Phi_{ij}\):
\[
\begin{aligned}
 {\cal A}_{ik}
 &=\nabla_k\Phi_{jk}\Phi_{ij}
   -\Phi_{jk}\Phi_{ij}\nabla_i\\
 &=(\nabla_k\Phi_{jk}-\Phi_{jk}\nabla_j)\Phi_{ij}
   +\Phi_{jk}(\nabla_j\Phi_{ij}-\Phi_{ij}\nabla_i),
\end{aligned}
\]
which is (6.9).  If (6.10) holds and
\(\xi_j=\Phi_{ij}\xi_i\), then
\[
 \nabla_j\xi_j
 =\nabla_j\Phi_{ij}\xi_i
 =\Phi_{ij}\nabla_i\xi_i,
\]
so the differentiated family glues.  Conversely, testing all local
sections gives (6.10).
\end{proof}

Let \({\cal R}_i=\nabla_i^2\).

\begin{corollary}[Exact curvature overlap defect]
\label{cor:v14v6-curvaturedefect}
For every local section \(\xi_i\),
\[
 \boxed{
 ({\cal R}_j\Phi_{ij}-\Phi_{ij}{\cal R}_i)\xi_i
 =
 \nabla_j({\cal A}_{ij}\xi_i)
 +{\cal A}_{ij}(\nabla_i\xi_i).}                            \tag{6.11}
\]
Thus \({\cal A}_{ij}=0\) implies exact curvature descent.
\end{corollary}

\begin{proof}
As an operator identity,
\[
\begin{aligned}
 \nabla_j^2\Phi_{ij}-\Phi_{ij}\nabla_i^2
 &=
 \nabla_j(\nabla_j\Phi_{ij}-\Phi_{ij}\nabla_i)
 +(\nabla_j\Phi_{ij}-\Phi_{ij}\nabla_i)\nabla_i.
\end{aligned}
\]
Apply this identity to \(\xi_i\).
\end{proof}

Equation (6.11) is the gluing analogue of the Bianchi and
Ricci-intertwining defect formulas: a nonzero overlap error is kept as
a typed connection defect rather than absorbed into an ambient norm.

\subsection{Ricci and Einstein descent}

Let local curvature two-form transitions satisfy
\[
 {\cal R}_j=\Phi_{ij}^{(2)}{\cal R}_i,                       \tag{6.12}
\]
and suppose the local contraction maps obey
\[
 {\cal C}_jS_{ij}
 =\Psi_{ij}{\cal C}_i,                                      \tag{6.13}
\]
where \(\Psi_{ij}\) is the transition on the Ricci target and itself
satisfies the cocycle law.

\begin{theorem}[Global Ricci section]
\label{thm:v14v6-riccidescent}
For \(P_i={\cal C}_ij_i{\cal R}_i\),
\[
 P_j=\Psi_{ij}P_i                                           \tag{6.14}
\]
on every overlap.  Hence \((P_i)_i\) defines a section of the
descended Ricci module.
\end{theorem}

\begin{proof}
Use the third identity in (6.7), then (6.12)--(6.13):
\[
\begin{aligned}
 P_j
 &={\cal C}_jj_j\Phi_{ij}^{(2)}{\cal R}_i\\
 &={\cal C}_jS_{ij}j_i{\cal R}_i
 =\Psi_{ij}{\cal C}_ij_i{\cal R}_i
 =\Psi_{ij}P_i.
\end{aligned}
\]
\end{proof}

\begin{corollary}[Global symmetric Einstein section]
\label{cor:v14v6-einstein}
Suppose the metric \(G_i\), inverse metric, and symmetric scalar
\(s_{\mathrm J,i}\) transform compatibly with the same tensor
transition.  Then
\[
 {\cal G}^{\mathrm J}_{i}
 =P_i-\frac14(G_i s_{\mathrm J,i}
              +s_{\mathrm J,i}G_i)                          \tag{6.15}
\]
obeys the overlap law and defines a global descended Einstein
section.
\end{corollary}

\begin{proof}
The Ricci term glues by
Theorem~\ref{thm:v14v6-riccidescent}.  Apply the assumed metric and
scalar transformations to each ordered product in (6.15).
Associativity gives the same target transition for both orders.
\end{proof}

The hypotheses of this corollary are explicit.  In particular, local
existence of a Ricci formula does not imply (6.13), and a cyclic
scalar class is insufficient unless it descends to the declared
operator-valued scalar target.

\subsection{Constraint before positivity}

Let \({\cal M}_0\) be a Hilbert \(\mathfrak A\)-module.  Transport its
inner product to each local twisted module by
\[
 \langle\xi,\eta\rangle_i
 :=
 \left\langle T_i^{-1}\xi,T_i^{-1}\eta\right\rangle_0.       \tag{6.16}
\]

\begin{proposition}[Unitary transitions and descended positivity]
\label{prop:v14v6-positivity}
Every \(\Phi_{ij}\) is unitary for the transported inner products:
\[
 \langle\Phi_{ij}\xi,\Phi_{ij}\eta\rangle_j
 =\langle\xi,\eta\rangle_i.                                 \tag{6.17}
\]
If local adjointable operators \(A_i\) satisfy
\[
 A_j\Phi_{ij}=\Phi_{ij}A_i                                  \tag{6.18}
\]
and \(A_i\succeq0\) on one, hence every, trivialization, then
\[
 A(\xi_i)_i:=(A_i\xi_i)_i                                  \tag{6.19}
\]
is a positive operator on the descended physical section module.
\end{proposition}

\begin{proof}
Since \(T_j^{-1}\Phi_{ij}=T_i^{-1}\), equation (6.17) follows
immediately from (6.16).  Equation (6.18) makes (6.19) preserve the
overlap equalizer.  For a glued section, positivity is checked in any
local trivialization:
\[
 \langle\xi_i,A_i\xi_i\rangle_i\geq0.
\]
Equation (6.17) and the intertwining law make the value independent
of the chosen trivialization on overlaps.
\end{proof}

This is the required order of operations from the V5 audit: the
physical equalizer is formed first; positivity is then tested on it.
No independent ambient null direction is added.

\subsection{Cofinal descent}

\begin{theorem}[Cofinal preservation of exact gluing]
\label{thm:v14v6-cofinal}
For every \(n\), let a finite cover carry exact transitions
\(\Phi_{ij,n}\) and exactly glued local objects \(\xi_{i,n}\):
\[
 \xi_{j,n}=\Phi_{ij,n}\xi_{i,n}.                            \tag{6.20}
\]
Suppose, uniformly on overlaps,
\[
 \Phi_{ij,n}\to\Phi_{ij},\qquad
 \xi_{i,n}\to\xi_i,                                         \tag{6.21}
\]
and the transition norms are uniformly bounded.  Then
\[
 \xi_j=\Phi_{ij}\xi_i,                                      \tag{6.22}
\]
so the limit is an exact descended section.

The same conclusion applies simultaneously to the braided splittings,
connections, curvatures, Ricci tensors, and Einstein tensors whenever
their local convergence and exact finite-stage intertwining laws hold.
\end{theorem}

\begin{proof}
For every overlap,
\[
\begin{aligned}
 \|\xi_j-\Phi_{ij}\xi_i\|
 \leq{}&
 \|\xi_j-\xi_{j,n}\|
 +\|\Phi_{ij,n}\|\,\|\xi_{i,n}-\xi_i\|\\
 &+\|(\Phi_{ij,n}-\Phi_{ij})\xi_i\|.
\end{aligned}
\]
Each term tends to zero by (6.21) and uniform boundedness, proving
(6.22).  The remaining objects obey equations of the same form in
their corresponding finite products, so the identical argument
applies.
\end{proof}

\begin{corollary}[Uniform inverse margins descend]
\label{cor:v14v6-inversemargin}
Let \(L_{i,n}\) be local selector operators that intertwine the
transitions and are invertible with
\[
 \sup_{i,n}\|L_{i,n}^{-1}\|<\infty.                          \tag{6.23}
\]
If \(L_{i,n}\to L_i\) and every \(L_i\) is invertible, the local
solutions glue at each stage and land as the unique solution of the
global descended selector.
\end{corollary}

\begin{proof}
Intertwining and uniqueness make local inverse solutions agree on
overlaps.  The inverse identity proves local convergence, and
Theorem~\ref{thm:v14v6-cofinal} gives exact descent of the limit.
\end{proof}

\subsection{Acquisition ledger and remaining population}

\[
\begin{array}{p{0.29\linewidth}|p{0.23\linewidth}|p{0.38\linewidth}}
\textbf{gate}&\textbf{status after V6}&\textbf{content}\\ \hline
\text{twisted-module transition}
 &\text{acquired}
 &exact groupoid laws (6.3)\\
\text{physical section constraint}
 &\text{acquired}
 &closed overlap equalizer (6.4)\\
\text{braided exterior descent}
 &\text{acquired}
 &three intertwining laws (6.7)\\
\text{connection/curvature descent}
 &\text{reduced to explicit defect}
 &{\cal A}_{ij}\text{ and identity (6.11)}\\
\text{Ricci/Einstein descent}
 &\text{acquired conditionally}
 &curvature and contraction intertwiners\\
\text{positivity on physical range}
 &\text{acquired conditionally}
 &transported Hilbert-module norms\\
\text{cofinal gluing}
 &\text{acquired}
 &uniform transition and local-object convergence\\
\text{regulator population}
 &\text{open}
 &construct \(E_i,\Phi_{ij},{\cal A}_{ij}=0\), and
   \({\cal C}_i\) on the actual finite stages
\end{array}                                                     \tag{6.24}
\]

The next bottleneck is no longer algebraic descent.  It is the
physical overlap connection: the regulator must make
\({\cal A}_{ij}\) vanish, or provide summable defects with an exact
correction that restores the cocycle.  The next version will derive
that correction problem and its obstruction class rather than assume
connection compatibility.

\subsection{Parent-integrity record}

The installed parent was
\[
\begin{array}{c|c}
 \text{file}&
 \texttt{Renewal\_SM\_Einstein\_Continuum\_Research\_Notes\_V5.tex}\\
 \text{lines}&1933\\
 \text{bytes}&65727\\
 \text{SHA--256}&
 \texttt{D44B1CE72962516714F1B99933243E608881FE87EC5AA0C4C3432388DF2B5801}.
\end{array}                                                     \tag{6.25}
\]

% =============================================================================
% END APPENDED FOURTEENTH-CYCLE V6 MATERIAL
% =============================================================================
% =============================================================================
% BEGIN APPENDED FOURTEENTH-CYCLE V7 MATERIAL
% =============================================================================

\section{Fourteenth-cycle V7: connection correction cocycles and
uniqueness-forced Levi--Civita descent}
\label{sec:v14v7}

\subsection{Changing local connections}

Retain the transition groupoid \(\Phi_{ij}:{\cal M}_i\to{\cal M}_j\)
and the connection overlap defect
\[
 {\cal A}_{ij}
 =\nabla_j\Phi_{ij}-\Phi_{ij}\nabla_i                       \tag{7.1}
\]
from V6.  Let \(B_i\) be a degree-one local endomorphism and set
\[
 \widehat\nabla_i:=\nabla_i+B_i.                            \tag{7.2}
\]

\begin{proposition}[Exact correction law]
\label{prop:v14v7-correction}
The corrected overlap defect is
\[
 \boxed{
 \widehat{\cal A}_{ij}
 ={\cal A}_{ij}+B_j\Phi_{ij}-\Phi_{ij}B_i.}                 \tag{7.3}
\]
Hence exact connection descent after correction is equivalent to
\[
 {\cal A}_{ij}
 =\Phi_{ij}B_i-B_j\Phi_{ij}
 \quad\hbox{on every overlap}.                              \tag{7.4}
\]
\end{proposition}

\begin{proof}
Insert (7.2):
\[
\begin{aligned}
 \widehat\nabla_j\Phi_{ij}
 -\Phi_{ij}\widehat\nabla_i
 &=(\nabla_j+B_j)\Phi_{ij}
   -\Phi_{ij}(\nabla_i+B_i)\\
 &={\cal A}_{ij}+B_j\Phi_{ij}-\Phi_{ij}B_i.
\end{aligned}
\]
Setting this expression to zero gives (7.4).
\end{proof}

\subsection{The transition-valued Cech class}

Define the coboundary of a zero-cochain \(B=(B_i)\) by
\[
 (\delta_\Phi B)_{ij}:=
 B_j\Phi_{ij}-\Phi_{ij}B_i.                                 \tag{7.5}
\]
For a one-cochain \(A=(A_{ij})\), define
\[
 (\delta_\Phi A)_{ijk}
 :=
 A_{ik}-A_{jk}\Phi_{ij}-\Phi_{jk}A_{ij}.                   \tag{7.6}
\]

\begin{lemma}[Cocycle identity]
\label{lem:v14v7-cocycle}
For every zero-cochain \(B\),
\[
 \delta_\Phi(\delta_\Phi B)=0.                              \tag{7.7}
\]
The connection defect (7.1) satisfies
\[
 \delta_\Phi{\cal A}=0.                                     \tag{7.8}
\]
\end{lemma}

\begin{proof}
Insert (7.5) in (7.6) and use
\(\Phi_{ik}=\Phi_{jk}\Phi_{ij}\):
\[
\begin{aligned}
 &B_k\Phi_{ik}-\Phi_{ik}B_i\\
 &\quad-(B_k\Phi_{jk}-\Phi_{jk}B_j)\Phi_{ij}
       -\Phi_{jk}(B_j\Phi_{ij}-\Phi_{ij}B_i)=0.
\end{aligned}
\]
Equation (7.8) is exactly the connection-defect cocycle (6.9).
\end{proof}

\begin{definition}[Overlap obstruction class]
Let \(Z_\Phi^1:=\ker\delta_\Phi\) and
\(B_\Phi^1:=\operatorname{Ran}\delta_\Phi\).  The class
\[
 [{\cal A}]_\Phi\in Z_\Phi^1/B_\Phi^1                       \tag{7.9}
\]
is the connection overlap obstruction.  By
Proposition~\ref{prop:v14v7-correction}, an unrestricted exact
correction exists if and only if
\[
 [{\cal A}]_\Phi=0.                                         \tag{7.10}
\]
\end{definition}

This is a connection-valued descent class.  It is distinct from the
determinant-line, anomaly, and measure-overlap classes in the frozen
archives.

\subsection{Unrestricted partition correction}

Pull every overlap defect to the common reference trivialization:
\[
 a_{ij}:=T_j^{-1}{\cal A}_{ij}T_i.                          \tag{7.11}
\]
Then (7.8) becomes
\[
 a_{ik}=a_{ij}+a_{jk}.                                      \tag{7.12}
\]

\begin{theorem}[Fine-sheaf correction]
\label{thm:v14v7-partition}
Suppose \(X\) admits a central continuous partition of unity
\(\{\rho_k\}_{k=1}^N\) subordinate to the finite cover, and the full
space of local degree-one endomorphisms is allowed for corrections.
Define
\[
 b_i:=\sum_k\rho_k a_{ik},\qquad
 B_i:=T_i b_iT_i^{-1}.                                      \tag{7.13}
\]
Then
\[
 \widehat{\cal A}_{ij}=0.                                   \tag{7.14}
\]
If \(\|a_{ij}\|\leq\varepsilon\) on all overlaps, then
\[
 \|b_i\|\leq\varepsilon,\qquad
 \|B_i\|\leq\|T_i\|\,\|T_i^{-1}\|\,\varepsilon.              \tag{7.15}
\]
\end{theorem}

\begin{proof}
From (7.12),
\[
 a_{jk}-a_{ik}=-a_{ij}.
\]
Since \(\sum_k\rho_k=1\),
\[
 b_j-b_i
 =\sum_k\rho_k(a_{jk}-a_{ik})
 =-a_{ij}.                                                   \tag{7.16}
\]
Conjugating (7.16) by \(T_j\) on the left and \(T_i^{-1}\) on the
right gives
\[
 B_j\Phi_{ij}-\Phi_{ij}B_i=-{\cal A}_{ij}.
\]
Equation (7.3) proves (7.14).  Centrality of the \(\rho_k\) keeps the
correction module-linear, and convexity gives
\(\|b_i\|\leq\sum_k\rho_k\varepsilon=\varepsilon\).
The last bound follows by conjugation.
\end{proof}

Theorem~\ref{thm:v14v7-partition} does not yet solve the geometric
problem.  The correction \(B_i\) may destroy metric compatibility or
torsion-freeness.

\subsection{Admissible corrections}

Let \(G_i\) and \(\theta_i\) be the local metric and coframe, and
suppose \(\nabla_i\) is metric compatible and torsion free.  Adding
\(B_i\) changes the two defects by
\[
\begin{aligned}
 {\cal Q}_{\nabla_i+B_i}
 &=-\bigl(B_i^*G_i+G_iB_i\bigr),\\
 {\cal T}_{\nabla_i+B_i}
 &=B_i\theta_i.                                              \tag{7.17}
\end{aligned}
\]
Define the homogeneous admissible correction space
\[
 {\cal N}_i
 :=
 \left\{
 B:
 B^*G_i+G_iB=0,\quad B\theta_i=0
 \right\}.                                                   \tag{7.18}
\]

\begin{proposition}[Selector uniqueness kills admissible
corrections]
\label{prop:v14v7-rigidity}
If the metric/torsion selector at \(U_i\) is injective, then
\[
 {\cal N}_i=\{0\}.                                          \tag{7.19}
\]
\end{proposition}

\begin{proof}
If \(B\in{\cal N}_i\), then \(\nabla_i+B\) is a second
metric-compatible torsion-free connection for the same
\((G_i,\theta_i)\), by (7.17).  Injectivity of the Cartan selector,
equivalently uniqueness in the V1 selector cell, forces
\(\nabla_i+B=\nabla_i\), hence \(B=0\).
\end{proof}

\begin{corollary}[No partition repair inside a unique
Levi--Civita fibre]
\label{cor:v14v7-norepair}
Under the injectivity hypothesis on every chart, a correction
preserving both metricity and torsion exists if and only if
\[
 {\cal A}_{ij}=0
 \quad\hbox{already}.                                       \tag{7.20}
\]
\end{corollary}

\begin{proof}
Every admissible \(B_i\) vanishes by
Proposition~\ref{prop:v14v7-rigidity}.  Equation (7.3) then gives
\(\widehat{\cal A}_{ij}={\cal A}_{ij}\).
\end{proof}

This is why convex partition arguments valid for measures or Choi
operators cannot be copied into the Levi--Civita problem.

\subsection{Exact geometric transitions force descent}

Transport the connection from chart \(i\) to chart \(j\):
\[
 \nabla_i^{(j)}
 :=\Phi_{ij}\nabla_i\Phi_{ji}.                              \tag{7.21}
\]

\begin{theorem}[Uniqueness-forced Levi--Civita gluing]
\label{thm:v14v7-uniquenessgluing}
Assume on every overlap:
\begin{enumerate}
\item \(\Phi_{ij}\) transports \(G_i,\theta_i\), the involution,
      wedge, and differential to the corresponding chart-\(j\) data;
\item \(\nabla_i\) is metric compatible and torsion free;
\item the chart-\(j\) metric/torsion selector is injective.
\end{enumerate}
Then
\[
 \nabla_i^{(j)}=\nabla_j,\qquad
 {\cal A}_{ij}=0.                                           \tag{7.22}
\]
Consequently the local Levi--Civita connections, curvatures, braided
Ricci tensors, and symmetric Einstein tensors descend exactly.
\end{theorem}

\begin{proof}
Naturality of the transported involution and metric makes
\(\nabla_i^{(j)}\) metric compatible with \(G_j\).  Naturality of the
differential, wedge, and coframe makes its torsion equal to the
transport of the zero torsion of \(\nabla_i\).  Thus
\(\nabla_i^{(j)}\) and \(\nabla_j\) solve the same metricity and
torsion equations on chart \(j\).  Selector injectivity gives the
first equality in (7.22).

Multiplying that equality on the right by \(\Phi_{ij}\) gives
\(\nabla_j\Phi_{ij}=\Phi_{ij}\nabla_i\), which is the second
equality.  Curvature descent follows from (6.11); Ricci and Einstein
descent then follow from V6.
\end{proof}

\begin{remark}[What must be populated]
Theorem~\ref{thm:v14v7-uniquenessgluing} removes connection overlap
as an independent obstruction only after the transition has been
proved to transport the full geometric calculus.  A coframe matrix
cocycle alone is insufficient: differential, wedge, involution, and
metric compatibility must transform together.
\end{remark}

\subsection{Approximate class and cofinal use}

For finite stages, define the pulled defect norm
\[
 \varepsilon_n
 :=\max_{i,j}
 \|T_{j,n}^{-1}{\cal A}_{ij,n}T_{i,n}\|.                   \tag{7.23}
\]
The unrestricted partition correction has norm at most
\(\varepsilon_n\) in the reference trivialization.  If
\(\sum_n\varepsilon_n<\infty\), these corrections can enter a
summable diagonal.  They cannot be claimed to preserve both metricity
and torsion unless the corrected geometric data are also changed and
the selector is solved again.

The correct approximate route is therefore to transport
\((G_i,\theta_i,{\cal T}_{0,i})\), solve the unique selector in one
common fibre, and estimate the difference of the resulting solutions.
That quantitative stability theorem is the next version.

\subsection{Acquisition ledger}

\[
\begin{array}{p{0.29\linewidth}|p{0.23\linewidth}|p{0.38\linewidth}}
\textbf{gate}&\textbf{status after V7}&\textbf{content}\\ \hline
\text{connection correction law}
 &\text{acquired}
 &\widehat{\cal A}={\cal A}+\delta_\Phi B\\
\text{overlap obstruction class}
 &\text{acquired}
 &[{\cal A}]_\Phi\in Z_\Phi^1/B_\Phi^1\\
\text{unrestricted correction}
 &\text{acquired}
 &central partition formula (7.13)\\
\text{metric/torsion-preserving correction}
 &\text{rigid in selector cell}
 &{\cal N}_i=\{0\}\\
\text{exact Levi--Civita descent}
 &\text{acquired conditionally}
 &geometric naturality plus uniqueness\\
\text{approximate descent}
 &\text{open quantitatively}
 &compare transported selector solutions and residual data
\end{array}                                                     \tag{7.24}
\]

\subsection{Parent-integrity record}

The installed parent was
\[
\begin{array}{c|c}
 \text{file}&
 \texttt{Renewal\_SM\_Einstein\_Continuum\_Research\_Notes\_V6.tex}\\
 \text{lines}&2380\\
 \text{bytes}&79313\\
 \text{SHA--256}&
 \texttt{B891DE1A1E262B0A8F6F27DF4019B2F3BEE829A0B0F46AC64663DD13DBA9ACDD}.
\end{array}                                                     \tag{7.25}
\]

% =============================================================================
% END APPENDED FOURTEENTH-CYCLE V7 MATERIAL
% =============================================================================
% =============================================================================
% BEGIN APPENDED FOURTEENTH-CYCLE V8 MATERIAL
% =============================================================================

\section{Fourteenth-cycle V8: quantitative overlap stability from
selector data to Einstein sections}
\label{sec:v14v8}

\subsection{Two-chart selector comparison}

On an overlap \(U_i\cap U_j\), let
\[
 L_i:{\cal X}_i\to{\cal Y}_i,\qquad
 L_j:{\cal X}_j\to{\cal Y}_j                              \tag{8.1}
\]
be the metric/torsion selectors, and let
\[
 L_iH_i=-\tau_i,\qquad L_jH_j=-\tau_j.                      \tag{8.2}
\]
Here \(\tau_i,\tau_j\) denote the local torsion sources.  Let
\[
 U^X_{ij}:{\cal X}_i\to{\cal X}_j,\qquad
 U^Y_{ij}:{\cal Y}_i\to{\cal Y}_j                          \tag{8.3}
\]
be the maps induced by the geometric transition.  Transport the
chart-\(i\) data to chart \(j\):
\[
 \widetilde L_i
 :=U^Y_{ij}L_i(U^X_{ij})^{-1},\quad
 \widetilde H_i:=U^X_{ij}H_i,\quad
 \widetilde\tau_i:=U^Y_{ij}\tau_i.                         \tag{8.4}
\]
Then
\[
 \widetilde L_i\widetilde H_i=-\widetilde\tau_i.            \tag{8.5}
\]

\begin{theorem}[Selector overlap estimate]
\label{thm:v14v8-selector}
Assume \(L_j\) is invertible and
\[
 \|L_j^{-1}\|\leq M.                                        \tag{8.6}
\]
Define
\[
 \varepsilon_L:=\|L_j-\widetilde L_i\|,\qquad
 \varepsilon_\tau:=\|\tau_j-\widetilde\tau_i\|.             \tag{8.7}
\]
Then
\[
 \boxed{
 \|H_j-\widetilde H_i\|
 \leq M\left(
       \varepsilon_\tau
       +\varepsilon_L\|\widetilde H_i\|
      \right).}                                              \tag{8.8}
\]
If also
\[
 \|\widetilde L_i^{-1}\|\leq M,\qquad
 \|\widetilde\tau_i\|\leq T,                                \tag{8.9}
\]
then
\[
 \|H_j-\widetilde H_i\|
 \leq M\varepsilon_\tau+M^2T\varepsilon_L.                  \tag{8.10}
\]
\end{theorem}

\begin{proof}
Subtract the two equations in (8.2) and (8.5) after writing both in
the chart-\(j\) fibre:
\[
\begin{aligned}
 L_j(H_j-\widetilde H_i)
 &=-\tau_j-L_j\widetilde H_i\\
 &=(\widetilde\tau_i-\tau_j)
   +(\widetilde L_i-L_j)\widetilde H_i.
\end{aligned}
\]
Apply \(L_j^{-1}\), take norms, and use (8.6), proving (8.8).
Equation (8.5) and (8.9) give
\(\|\widetilde H_i\|\leq MT\), which yields (8.10).
\end{proof}

This comparison uses the selected physical correction \(H\), rather
than an ambient connection perturbation.  It therefore preserves the
constraint-before-estimate order fixed in V5.

\subsection{Connection overlap bound}

Write the selected connections as
\[
 \Gamma_j=\Gamma_{0,j}+G_j^{-1}H_j,\qquad
 \widetilde\Gamma_i
 =\widetilde\Gamma_{0,i}
  +\widetilde G_i^{-1}\widetilde H_i.                       \tag{8.11}
\]
Define
\[
\begin{aligned}
 \varepsilon_0
 &:=\|\Gamma_{0,j}-\widetilde\Gamma_{0,i}\|,\\
 \varepsilon_{G^{-1}}
 &:=\|G_j^{-1}-\widetilde G_i^{-1}\|.                       \tag{8.12}
\end{aligned}
\]

\begin{corollary}[Selected connection overlap estimate]
\label{cor:v14v8-connection}
Under the hypotheses of
Theorem~\ref{thm:v14v8-selector},
\[
\begin{aligned}
 \|\Gamma_j-\widetilde\Gamma_i\|
 \leq{}&
 \varepsilon_0
 +\|G_j^{-1}\|
   M\left(\varepsilon_\tau
          +\varepsilon_L\|\widetilde H_i\|\right)\\
 &+\varepsilon_{G^{-1}}\|\widetilde H_i\|.                  \tag{8.13}
\end{aligned}
\]
Under (8.9), the right-hand side is at most
\[
 \varepsilon_0
 +\|G_j^{-1}\|
   (M\varepsilon_\tau+M^2T\varepsilon_L)
 +MT\varepsilon_{G^{-1}}.                                  \tag{8.14}
\]
\end{corollary}

\begin{proof}
Subtract the two expressions in (8.11), then add and subtract
\(G_j^{-1}\widetilde H_i\):
\[
\begin{aligned}
 \Gamma_j-\widetilde\Gamma_i
 ={}&
 \Gamma_{0,j}-\widetilde\Gamma_{0,i}
 +G_j^{-1}(H_j-\widetilde H_i)\\
 &+(G_j^{-1}-\widetilde G_i^{-1})\widetilde H_i.
\end{aligned}
\]
Take norms and apply (8.8), then (8.9).
\end{proof}

In a common trivialization, the norm in (8.13) is the coefficient
form of the connection overlap defect \({\cal A}_{ij}\).  Exact
transport of the three selector inputs makes every epsilon vanish and
recovers the uniqueness theorem (7.22).

\subsection{Curvature propagation}

Let
\[
 {\cal R}(\Gamma):=d\Gamma+\Gamma^2.
\]
Assume
\[
 \|\Gamma_j\|,\|\widetilde\Gamma_i\|\leq C_\Gamma,\qquad
 \|d\Gamma_j-d\widetilde\Gamma_i\|\leq\varepsilon_d.        \tag{8.15}
\]

\begin{proposition}[Curvature overlap estimate]
\label{prop:v14v8-curvature}
If
\[
 \varepsilon_\Gamma
 :=\|\Gamma_j-\widetilde\Gamma_i\|,                          \tag{8.16}
\]
then
\[
 \boxed{
 \|{\cal R}_j-\widetilde{\cal R}_i\|
 \leq\varepsilon_d+2C_\Gamma\varepsilon_\Gamma.}            \tag{8.17}
\]
\end{proposition}

\begin{proof}
Use
\[
 \Gamma_j^2-\widetilde\Gamma_i^2
 =(\Gamma_j-\widetilde\Gamma_i)\Gamma_j
  +\widetilde\Gamma_i(\Gamma_j-\widetilde\Gamma_i).
\]
Submultiplicativity gives a bound
\(2C_\Gamma\varepsilon_\Gamma\); add the differential term.
\end{proof}

The derivative defect \(\varepsilon_d\) is indispensable.  Norm
closeness of connections alone does not control curvature.

\subsection{Braided Ricci propagation}

In the chart-\(j\) fibre, write the local and transported Ricci maps
as
\[
 P_j={\cal C}_jj_j{\cal R}_j,\qquad
 \widetilde P_i
 =\widetilde{\cal C}_i\widetilde j_i\widetilde{\cal R}_i.    \tag{8.18}
\]
Set
\[
\begin{aligned}
 \varepsilon_{\cal C}
 &:=\|{\cal C}_j-\widetilde{\cal C}_i\|,\\
 \varepsilon_j
 &:=\|j_j-\widetilde j_i\|,\\
 \varepsilon_{\cal R}
 &:=\|{\cal R}_j-\widetilde{\cal R}_i\|.                    \tag{8.19}
\end{aligned}
\]
Assume
\[
 \|{\cal C}_j\|,\|\widetilde{\cal C}_i\|\leq C_{\cal C},
 \quad
 \|j_j\|,\|\widetilde j_i\|\leq C_j,
 \quad
 \|{\cal R}_j\|,\|\widetilde{\cal R}_i\|\leq C_{\cal R}.     \tag{8.20}
\]

\begin{theorem}[Ricci overlap estimate]
\label{thm:v14v8-ricci}
Under (8.19)--(8.20),
\[
 \boxed{
 \|P_j-\widetilde P_i\|
 \leq
 \varepsilon_{\cal C}C_jC_{\cal R}
 +C_{\cal C}\varepsilon_jC_{\cal R}
 +C_{\cal C}C_j\varepsilon_{\cal R}.}                       \tag{8.21}
\]
\end{theorem}

\begin{proof}
Insert the two intermediate expressions
\(\widetilde{\cal C}_ij_j{\cal R}_j\) and
\(\widetilde{\cal C}_i\widetilde j_i{\cal R}_j\):
\[
\begin{aligned}
 P_j-\widetilde P_i
 ={}&
 ({\cal C}_j-\widetilde{\cal C}_i)j_j{\cal R}_j\\
 &+\widetilde{\cal C}_i(j_j-\widetilde j_i){\cal R}_j\\
 &+\widetilde{\cal C}_i\widetilde j_i
   ({\cal R}_j-\widetilde{\cal R}_i).
\end{aligned}
\]
Take norms and use (8.20).
\end{proof}

The three terms have different meanings: contraction mismatch,
braided-lift mismatch, and curvature mismatch.  They must remain
separate in the acquisition ledger.

\subsection{Scalar and Einstein propagation}

Let \(g_j=G_j^{-1}\) and let \(\widetilde g_i\) be the transported
inverse metric.  Suppose, in a matrix norm compatible with the finite
index contractions,
\[
 \|g_j\|,\|\widetilde g_i\|\leq C_g,\qquad
 \|P_j\|,\|\widetilde P_i\|\leq C_P.                         \tag{8.22}
\]
Let
\[
 \varepsilon_g:=\|g_j-\widetilde g_i\|,\qquad
 \varepsilon_P:=\|P_j-\widetilde P_i\|.                     \tag{8.23}
\]

\begin{proposition}[Symmetric scalar overlap estimate]
\label{prop:v14v8-scalar}
For the Jordan-ordered scalar contraction (2.6), there is a
finite index-norm constant \(c_m\) such that
\[
 \|s_{\mathrm J,j}-\widetilde s_{\mathrm J,i}\|
 \leq
 c_m\left(C_P\varepsilon_g+C_g\varepsilon_P\right).          \tag{8.24}
\]
For the entrywise maximum norm one may take \(c_m=m^2\).
\end{proposition}

\begin{proof}
For each of the two orderings, add and subtract the product with
one transported factor:
\[
 g_jP_j-\widetilde g_i\widetilde P_i
 =(g_j-\widetilde g_i)P_j
  +\widetilde g_i(P_j-\widetilde P_i).
\]
The right ordering is identical with the factors reversed.  Sum the
\(m^2\) contracted entries and average.
\end{proof}

Let
\[
 {\cal G}_j
 =P_j-\frac14(G_js_j+s_jG_j),\qquad
 \widetilde{\cal G}_i
 =\widetilde P_i
  -\frac14(\widetilde G_i\widetilde s_i
            +\widetilde s_i\widetilde G_i).                 \tag{8.25}
\]
Assume
\[
 \|G_j\|,\|\widetilde G_i\|\leq C_G,\quad
 \|s_j\|,\|\widetilde s_i\|\leq C_s,                        \tag{8.26}
\]
and put
\(\varepsilon_G=\|G_j-\widetilde G_i\|\),
\(\varepsilon_s=\|s_j-\widetilde s_i\|\).

\begin{corollary}[Einstein overlap estimate]
\label{cor:v14v8-einstein}
One has
\[
 \boxed{
 \|{\cal G}_j-\widetilde{\cal G}_i\|
 \leq
 \varepsilon_P
 +\frac12\left(C_s\varepsilon_G+C_G\varepsilon_s\right).}   \tag{8.27}
\]
\end{corollary}

\begin{proof}
Subtract (8.25), expand each ordered product difference, and use
(8.26).  The two left/right orders contribute the same upper bound
and carry the prefactor \(1/4\).
\end{proof}

\subsection{Cofinal approximate descent theorem}

\begin{theorem}[Vanishing overlap defects yield exact continuum
descent]
\label{thm:v14v8-cofinal}
Consider a fixed finite cover through a cofinal sequence.  Assume:
\begin{enumerate}
\item the selector inverse bounds \(M\), source bounds \(T\), and
      all constants in (8.15), (8.20), (8.22), and (8.26) are uniform;
\item the local geometric objects converge in the norms used above;
\item uniformly on overlaps,
\[
 \varepsilon_L,\varepsilon_\tau,\varepsilon_0,
 \varepsilon_{G^{-1}},\varepsilon_d,
 \varepsilon_{\cal C},\varepsilon_j,\varepsilon_G
 \longrightarrow0.                                         \tag{8.28}
\]
\end{enumerate}
Then the overlap defects of the selected connections, curvatures,
Ricci tensors, symmetric scalars, and symmetric Einstein tensors all
tend to zero.  Their local limits satisfy the exact overlap equations
and define global descended continuum sections.
\end{theorem}

\begin{proof}
Apply (8.14) to the connection, (8.17) to curvature, (8.21) to Ricci,
(8.24) to the scalar, and (8.27) to Einstein.  Each right-hand side
tends to zero in that order.  Local convergence and the vanishing
overlap differences give the exact limiting overlap equations by the
three-term argument in Theorem~\ref{thm:v14v6-cofinal}.
\end{proof}

\begin{corollary}[Summable diagonal version]
\label{cor:v14v8-summable}
If the right-hand sides of (8.14), (8.17), (8.21), (8.24), and
(8.27) are summable along a chosen diagonal, then every downstream
overlap defect is summable and can be inserted into the existing
cofinal Cauchy ledger without changing its convergence conclusion.
\end{corollary}

\begin{proof}
Each downstream bound is a finite linear combination of upstream
summable sequences with uniform coefficients.  Finite sums and
constant multiples of summable nonnegative sequences are summable.
\end{proof}

\begin{firewall}[Finite stages need not glue exactly]
\label{fire:v14v8-finite}
Theorem~\ref{thm:v14v8-cofinal} proves exact descent of the limit.
It does not turn a nonzero finite-stage defect into an exact
finite-stage global connection.  The unrestricted partition
correction of V7 generally breaks metricity or torsion.  Exact
finite-stage descent still requires geometric naturality or a new
selector solve after correcting the underlying data.
\end{firewall}

\subsection{Acquisition ledger and next target}

\[
\begin{array}{p{0.29\linewidth}|p{0.23\linewidth}|p{0.38\linewidth}}
\textbf{gate}&\textbf{status after V8}&\textbf{content}\\ \hline
\text{selector overlap}
 &\text{quantified}
 &M\varepsilon_\tau+M^2T\varepsilon_L\\
\text{connection overlap}
 &\text{quantified}
 &baseline, selector, and inverse-metric defects\\
\text{curvature overlap}
 &\text{quantified}
 &derivative plus quadratic connection terms\\
\text{Ricci overlap}
 &\text{quantified}
 &contraction, split, and curvature defects\\
\text{scalar/Einstein overlap}
 &\text{quantified}
 &ordered product bounds\\
\text{continuum exact descent}
 &\text{acquired conditionally}
 &uniform convergence and vanishing overlap ledger\\
\text{contracted Bianchi support}
 &\text{open}
 &combine descended Einstein section with the physical source
\end{array}                                                     \tag{8.29}
\]

The next version returns downstream.  It will place the descended
Einstein tensor and the already constructed stress tensor in one
common dual test space, formulate the exact finite-stage support
residual, and prove its passage to the joint continuum state.  This
must precede any assertion that the limiting law is supported on the
Einstein equation.

\subsection{Parent-integrity record}

The installed parent was
\[
\begin{array}{c|c}
 \text{file}&
 \texttt{Renewal\_SM\_Einstein\_Continuum\_Research\_Notes\_V7.tex}\\
 \text{lines}&2718\\
 \text{bytes}&89787\\
 \text{SHA--256}&
 \texttt{B5CD0E291A8618CD4E437D9B707F8F0BAEFACBD118BEF27DEA58E66EAD0978A3}.
\end{array}                                                     \tag{8.30}
\]

% =============================================================================
% END APPENDED FOURTEENTH-CYCLE V8 MATERIAL
% =============================================================================
% =============================================================================
% BEGIN APPENDED FOURTEENTH-CYCLE V9 MATERIAL
% =============================================================================

\section{Fourteenth-cycle V9: represented Einstein writers and
localized noncommutative residual Grams}
\label{sec:v14v9}

\subsection{Nonduplication boundary}

The frozen V69--V71 notes already prove:
\begin{enumerate}
\item the orthogonal split of a classical weak Einstein residual into
      diffeomorphism-tangent and transverse channels;
\item the distinction between a mean equation and an almost-sure weak
      equation;
\item positive finite-screen residual Grams and a countable test
      criterion;
\item Killing-safe stability of the metric-dependent Bianchi
      projections.
\end{enumerate}
The Grand Tensor manuscript separately proves a conditional local
weak Einstein--matter handoff once the common action, source, compact
screen, and finite residual hypotheses are populated.

This version does not reprove those statements.  It constructs the
missing interface from the braided, descended, operator-valued
Einstein tensor of V6--V8 to the represented local algebra containing
the already compiled stress tensor.

\subsection{A descent-compatible weak writer}

Let \({\cal T}\) be a separable locally convex space of compactly
supported symmetric test tensors, with a fixed countable dense family
\[
 h_1,h_2,\ldots.                                             \tag{9.1}
\]
Let \({\cal E}_i\) be the local target module containing the symmetric
Einstein tensor \({\cal G}_i^{\mathrm J}\), with overlap transition
\(\Psi_{ij}\).  Let \({\cal A}_{\mathrm{loc}}\) be a common local
observable \(*\)-algebra represented on an invariant core
\({\cal D}\).

\begin{definition}[Weak Einstein writer]
A chartwise weak writer is a family of linear maps
\[
 {\mathfrak W}_i:
 {\cal E}_i\times{\cal T}
 \longrightarrow \operatorname{Op}({\cal D}),               \tag{9.2}
\]
linear in both arguments, satisfying
\[
 {\mathfrak W}_j(\Psi_{ij}X,h)
 ={\mathfrak W}_i(X,h)                                      \tag{9.3}
\]
on overlaps and
\[
 {\mathfrak W}_i(X,h)^*
 ={\mathfrak W}_i(X^*,\overline h)                          \tag{9.4}
\]
as quadratic forms on \({\cal D}\).  A graph seminorm bound is
\[
 \|{\mathfrak W}_i(X,h)\psi\|
 \leq C_h\|X\|_{{\cal E}_i}\|\psi\|_{\cal D}.                \tag{9.5}
\]
\end{definition}

\begin{proposition}[The writer descends]
\label{prop:v14v9-writerdescent}
If \((X_i)_i\) is a descended section,
\(X_j=\Psi_{ij}X_i\), then
\[
 {\mathfrak W}(X,h):={\mathfrak W}_i(X_i,h)                 \tag{9.6}
\]
is independent of the chart.  If \(X=X^*\) and \(h\) is real, the
writer is symmetric on \({\cal D}\).
\end{proposition}

\begin{proof}
Chart independence is (9.3).  Symmetry follows from (9.4).
\end{proof}

For approximate overlap, define
\[
 \varepsilon_{{\mathfrak W},ij}(h)
 :=
 \|{\mathfrak W}_j\Psi_{ij}
   -{\mathfrak W}_i\|_h,                                    \tag{9.7}
\]
where the norm is the operator norm from \({\cal E}_i\) to the graph
seminorm associated with \(h\).

\begin{proposition}[Weak-writer overlap bound]
\label{prop:v14v9-writerbound}
If \(X_j\) and \(\widetilde X_i=\Psi_{ij}X_i\) satisfy
\(\|X_j-\widetilde X_i\|\leq\varepsilon_X\), then
\[
\begin{aligned}
 \|{\mathfrak W}_j(X_j,h)
   -{\mathfrak W}_i(X_i,h)\|_{\cal D}
 \leq
 C_h\varepsilon_X
 +\varepsilon_{{\mathfrak W},ij}(h)\|X_i\|.                 \tag{9.8}
\end{aligned}
\]
For \(X_i={\cal G}^{\mathrm J}_i\), the value
\(\varepsilon_X\) may be taken from the explicit Einstein overlap
bound (8.27).
\end{proposition}

\begin{proof}
Add and subtract
\({\mathfrak W}_j(\widetilde X_i,h)\), apply (9.5) to the first
difference, and (9.7) to the second.
\end{proof}

\subsection{The same-source residual}

Let
\[
 {\mathbb E}(h)
 :={\mathfrak W}({\cal G}^{\mathrm J},h)                    \tag{9.9}
\]
be the descended geometric writer.  Let
\({\mathbb T}^{\mathrm{ren}}(h)\) be the represented renormalized
stress writer already used in the Ward and Hamiltonian incidence
compilers.  Define
\[
 {\mathbb R}(h)
 :={\mathbb E}(h)-\kappa{\mathbb T}^{\mathrm{ren}}(h).       \tag{9.10}
\]

\begin{definition}[Same-source incidence]
The residual (9.10) has same-source incidence when the stress writer
is the image of the identical finite-stage stress insertion,
counterterm prescription, and local-algebra embedding used in:
\begin{enumerate}
\item the diffeomorphism Ward identity;
\item the Hamiltonian/energy incidence theorem;
\item the all-order local observable ledger.
\end{enumerate}
Equality merely of vacuum expectations is not same-source incidence.
\end{definition}

This condition prevents the Einstein equation from being closed with
a stress tensor different from the one whose conservation and
Hamiltonian meaning were proved.

\subsection{Why a scalar trace is too weak}

\begin{proposition}[Trace-zero residual no-go]
\label{prop:v14v9-tracenogo}
Let \(\tau=\frac12\operatorname{Tr}\) on \(M_2(\mathbb C)\) and
\[
 R=\begin{pmatrix}1&0\\0&-1\end{pmatrix}.
\]
Then
\[
 \tau(R)=0,\qquad R\neq0,\qquad \tau(R^*R)=1.                \tag{9.11}
\]
Thus a cyclic or vacuum-mean Einstein identity cannot be promoted to
an operator equation.  A positive square is required.
\end{proposition}

\begin{proof}
The three statements follow by direct matrix calculation.
\end{proof}

The cyclic trace constructed in the frozen V96 theorem is therefore
appropriate for Chern forms, but it is not a separating
operator-valued Einstein writer.

\subsection{Localized noncommutative residual Grams}

Let \((\pi_\omega,{\cal H}_\omega,\Omega_\omega)\) be the GNS
representation of the joint continuum state.  Choose a countable
unital \(*\)-subalgebra
\[
 {\cal A}_0=\operatorname{alg}^*(A_1,A_2,\ldots)
 \subset{\cal A}_{\mathrm{loc}}                             \tag{9.12}
\]
such that
\[
 {\cal D}_0:=\pi_\omega({\cal A}_0)\Omega_\omega             \tag{9.13}
\]
is dense in the declared local physical core.  Assume every
\({\mathbb R}(h_j)\) is defined on \({\cal D}_0\).

For finite \(q,L\), define the Gram matrix indexed by pairs
\((j,\ell)\), \(1\leq j\leq q\), \(1\leq\ell\leq L\):
\[
 {\mathbb G}_{q,L}^{\mathrm{loc}}
 :=
 \left(
 \left\langle
  {\mathbb R}(h_j)\pi_\omega(A_\ell)\Omega_\omega,\,
  {\mathbb R}(h_k)\pi_\omega(A_r)\Omega_\omega
 \right\rangle
 \right)_{(j,\ell),(k,r)}.                                  \tag{9.14}
\]

\begin{theorem}[Localized Gram criterion for the represented weak
Einstein equation]
\label{thm:v14v9-gram}
Every matrix (9.14) is positive semidefinite.  The following are
equivalent:
\begin{enumerate}
\item
\[
 {\mathbb R}(h_j)\psi=0
 \quad\hbox{for every \(j\) and every \(\psi\in{\cal D}_0\)}; \tag{9.15}
\]
\item
\[
 {\mathbb G}_{q,L}^{\mathrm{loc}}=0
 \quad\hbox{for every \(q,L\)}.                              \tag{9.16}
\]
\end{enumerate}
If the residuals are continuous in \(h\) on \({\cal T}\), then these
conditions imply
\[
 {\mathbb R}(h)\psi=0
 \quad(h\in{\cal T},\ \psi\in{\cal D}_0).                   \tag{9.17}
\]
\end{theorem}

\begin{proof}
The matrix in (9.14) is the Gram matrix of the displayed finite
family of Hilbert-space vectors and is therefore positive.

If (9.15) holds, every vector is zero and (9.16) follows.
Conversely, fix \(j,\ell\) and choose \(q\geq j\), \(L\geq\ell\).
The corresponding diagonal entry of the zero positive matrix is
\[
 \|{\mathbb R}(h_j)\pi_\omega(A_\ell)\Omega_\omega\|^2=0.
\]
Linearity gives zero on the algebraic span (9.13), proving (9.15).
Finally, density of \(\{h_j\}\) and continuity in \(h\) extend zero
to every \(h\in{\cal T}\).
\end{proof}

\begin{remark}[Vacuum-vector versus core equality]
The unlocalized condition
\(\|{\mathbb R}(h_j)\Omega_\omega\|=0\) proves only that the residual
annihilates the vacuum vector.  The multiplier index \(\ell\) in
(9.14) is what proves equality on the dense represented core.  It is
the noncommutative counterpart of the bounded cylinder multipliers
in the archived almost-sure weak-equation theorem.
\end{remark}

\subsection{Finite-stage landing of the localized Grams}

Let \(h_1,\ldots,h_q\) and \(A_1,\ldots,A_L\) be fixed before the
cutoff is sent to infinity.  At stage \(n\), let
\[
 {\mathbb R}_n(h_j)
 ={\mathbb E}_n(h_j)-\kappa{\mathbb T}^{\mathrm{ren}}_n(h_j) \tag{9.18}
\]
act on the corresponding transported core vectors, and define
\({\mathbb G}_{q,L,n}^{\mathrm{loc}}\) by the finite-stage analogue
of (9.14).

\begin{theorem}[Cofinal localized-Gram passage]
\label{thm:v14v9-cofinal}
Assume for every fixed \(q,L\):
\begin{enumerate}
\item the braided Einstein writers and the same-source stress writers
      occur on one regulator history;
\item the local products
\[
 \bigl({\mathbb R}_n(h_j)A_{\ell,n}\bigr)^*
 \bigl({\mathbb R}_n(h_k)A_{r,n}\bigr)                       \tag{9.19}
\]
belong to the all-order observable ledger;
\item their expectations converge to the entries in (9.14), with the
      required uniform integrability for unbounded writers;
\item
\[
 \|{\mathbb G}_{q,L,n}^{\mathrm{loc}}\|\longrightarrow0.     \tag{9.20}
\]
\end{enumerate}
Then
\[
 {\mathbb G}_{q,L}^{\mathrm{loc}}=0
 \quad\hbox{for every \(q,L\)},                              \tag{9.21}
\]
and the represented weak Einstein equation (9.17) holds on
\({\cal D}_0\).
\end{theorem}

\begin{proof}
Entrywise convergence on each fixed finite screen identifies
\({\mathbb G}_{q,L}^{\mathrm{loc}}\) as the limit of the regulated
Gram matrices.  Equation (9.20) makes that limit the zero matrix.
Apply Theorem~\ref{thm:v14v9-gram}, then the density and continuity
clause.
\end{proof}

\begin{corollary}[Defect budget for the regulated residual]
\label{cor:v14v9-budget}
Suppose on a fixed screen
\[
 \|{\mathbb R}_n(h_j)A_{\ell,n}\Omega_n\|
 \leq
 d_{n,j,\ell}^{\mathrm{geom}}
 +d_{n,j,\ell}^{\mathrm{source}}
 +d_{n,j,\ell}^{\mathrm{overlap}}
 +d_{n,j,\ell}^{\mathrm{writer}}.                           \tag{9.22}
\]
If every term on the right tends to zero and the associated products
land as in Theorem~\ref{thm:v14v9-cofinal}, then (9.20) holds.
A uniform summable bound permits the screen and cutoff to be enlarged
on the existing diagonal.
\end{corollary}

\begin{proof}
Every Gram entry is bounded by the product of the corresponding
vector norms by Cauchy--Schwarz.  The finite sum of vanishing defects
makes every vector norm, and hence the finite Gram norm, tend to zero.
Summability supports the usual diagonal tail estimate.
\end{proof}

The geometric term in (9.22) contains the local contracted-Bianchi
and Ricci/scalar defects.  The overlap term is controlled by V8.  The
source term is the finite Einstein equation residual for the same
stress insertion.  The writer term measures the passage from the
braided tensor module to the represented weak operator.

\subsection{Relation to the classical residual screen}

If the geometric and stress writers are central multiplication
operators on a classical random sector, and the multipliers
\(A_\ell\) range over a measure-determining cylinder algebra, then
(9.14) reduces to the localized positive residual tests in the
frozen V69 theorem.  Thus V9 extends that screen to the
noncommutative represented core; it does not replace the
diffeomorphism-tangent/transverse split or the Killing-safe projection
theorem.

\subsection{Acquisition ledger and next gate}

\[
\begin{array}{p{0.29\linewidth}|p{0.23\linewidth}|p{0.38\linewidth}}
\textbf{gate}&\textbf{status after V9}&\textbf{content}\\ \hline
\text{braided-to-weak writer}
 &\text{defined with exact descent conditions}
 &{\mathfrak W}_j\Psi_{ij}={\mathfrak W}_i\\
\text{same-source incidence}
 &\text{declared}
 &identical Ward, Hamiltonian, and Einstein stress writer\\
\text{scalar trace sufficiency}
 &\text{refuted}
 &explicit traceless matrix residual\\
\text{represented Einstein support}
 &\text{reduced to positive Grams}
 &localized multiplier matrices (9.14)\\
\text{cofinal support passage}
 &\text{acquired conditionally}
 &all-order occurrence, uniform integrability, and zero Gram limits\\
\text{physical residual population}
 &\text{open}
 &construct one regulator with the four defects in (9.22)
\end{array}                                                     \tag{9.23}
\]

The V10 boundary is the next compulsory companion audit.  After that
audit, the programme should return to the geometric term in (9.22):
derive the contracted Bianchi defect of the braided symmetric
Einstein tensor on the descended core and compare it with the same
stress Ward defect.  Only their joint cancellation can close the
diffeomorphism-tangent residual channel.

\subsection{Parent-integrity record}

The installed parent was
\[
\begin{array}{c|c}
 \text{file}&
 \texttt{Renewal\_SM\_Einstein\_Continuum\_Research\_Notes\_V8.tex}\\
 \text{lines}&3158\\
 \text{bytes}&102683\\
 \text{SHA--256}&
 \texttt{A9A9B7F3F735B4B67F637D454406D5BD94AA2F9AA4CFDD4CCF487E95BBCF8AEA}.
\end{array}                                                     \tag{9.24}
\]

% =============================================================================
% END APPENDED FOURTEENTH-CYCLE V9 MATERIAL
% =============================================================================
% =============================================================================
% BEGIN APPENDED FOURTEENTH-CYCLE V10 MATERIAL
% =============================================================================

\section{Fourteenth-cycle V10: compulsory companion audit and the
physical-constraint multiplier direction}
\label{sec:v14v10}

\subsection{Files inspected}

The mandatory audit was performed on 10 September 2026 after V9.
The newest active companion sources were
\[
\begin{array}{p{0.39\linewidth}|r|r|p{0.27\linewidth}}
\textbf{file}&\textbf{lines}&\textbf{bytes}&\textbf{SHA--256}\\ \hline
\texttt{Renewal\_Yang\_Mills\_Mass\_Gap\_Research\_Notes\_V1.tex}
 &486&18746&
 \texttt{F524719D4B48905C75E2B27695E634BA1D7824EB01B870ACBC708ACD677BB10E}\\
\texttt{Renewal\_NS\_Stability\_Research\_Notes\_V26.tex}
 &5319&278283&
 \texttt{82C887F8C2C69E4F28FAD7C3DC8DE5B9099CB5A4BF431EBA1FFF3E91935978AD}.
\end{array}                                                    \tag{10.1}
\]
No companion file was modified.

\subsection{Yang--Mills rollover and new theorem}

The YM programme has frozen its preceding hundred-version cycle and
restarted with a compact result summary.  Its substantive V1 theorem
works in a decomposition
\[
 {\cal E}=\operatorname{Ran}V\oplus\operatorname{Ran}W,
 \qquad C=W^*,\qquad\ker C=\operatorname{Ran}V.              \tag{10.2}
\]
For a Hermitian ambient pencil \(B\), it proves the exact equivalence
\[
 V^*BV\succeq\delta I
 \quad\Longleftrightarrow\quad
 B+ZC+C^*Z^*\succeq\delta I
 \quad\hbox{for some \(Z\)}.                                 \tag{10.3}
\]
The added multiplier vanishes on the constrained range and preserves
the full positive floor.  The YM note also records an essential
firewall: its existential formula for \(Z\) contains the compressed
block, so it cannot serve as a constructive proof of that same
block's positivity.  A sparse multiplier must be obtained from the
explicit constraints.

\begin{proposition}[Transfer to the Einstein residual programme]
\label{prop:v14v10-transfer}
Let \(C_n\) be the complete finite-stage physical constraint operator.
The contracted Bianchi and stress-Ward comparison should first be
formed on
\[
 {\cal H}_{\mathrm{phys},n}:=\ker C_n.                       \tag{10.4}
\]
A term discarded from the ambient identity is legitimate only when
it is written as an explicit right factor through \(C_n\), or when
its norm on \(\ker C_n\) is proved to vanish.  Ambient cancellation
that uses an inadmissible gauge, BRST, or diffeomorphism direction
does not close the physical residual.
\end{proposition}

\begin{proof}
Every operator of the form \(YC_n\) vanishes on (10.4).  Conversely,
in finite dimension, an operator that vanishes on \(\ker C_n\)
factors through \(C_n\); this factorization will be proved in V11.
The YM moving-kernel counterexample and multiplier theorem show why
the restriction is load-bearing and why a multiplier must be typed
by the constraint.  The same logic applies to the represented
Einstein writer of V9.
\end{proof}

\subsection{Navier--Stokes status}

The NS source is unchanged.  Its rank-one kernel calculation still
supports testing estimates on realized inputs, while its small
second-derivative gain supplies no physical-constraint, Bianchi,
stress, or Einstein coefficient.

\subsection{Direction after the audit}

\begin{decision}[V11 closes the tangent channel modulo constraints]
\label{dec:v14v10-next}
V11 shall:
\begin{enumerate}
\item prove the exact factorization
      \(D=DC^\dagger C+D(I-C^\dagger C)\) for a closed-range
      constraint;
\item identify \(D(I-C^\dagger C)\) as the unexplained physical
      residual and characterize exact factorization through \(C\);
\item derive the ordered geometric divergence defect from the
      braided Bianchi, Ricci, scalar, and metricity defects;
\item subtract the same-source stress Ward defect;
\item prove that equality of those defects modulo the constraint
      closes only the diffeomorphism-tangent channel;
\item retain the transverse localized Gram gate of V9 for the full
      Einstein equation.
\end{enumerate}
No constraint multiplier will be chosen using an already assumed
Einstein equation.
\end{decision}

\begin{proof}
Items 1--2 implement the exact finite-dimensional transfer from
(10.3).  Items 3--4 connect the geometric construction of V1--V8 to
the represented same-source residual of V9.  The archived V69 split
proves that tangent closure is weaker than the full field equation,
forcing item 5 and the retained transverse gate.
\end{proof}

\subsection{Claim boundary and parent record}

This audit proves no mass gap, Navier--Stokes estimate, or
SM--Einstein equation.  It selects the physical quotient on which the
next exact identity must be tested.

The installed parent was
\[
\begin{array}{c|c}
 \text{file}&
 \texttt{Renewal\_SM\_Einstein\_Continuum\_Research\_Notes\_V9.tex}\\
 \text{lines}&3547\\
 \text{bytes}&115951\\
 \text{SHA--256}&
 \texttt{E3736BF429D46CA72C78B413F19028F011F4805AEE820F4F7F756CFC09ADC339}.
\end{array}                                                    \tag{10.5}
\]

% =============================================================================
% END APPENDED FOURTEENTH-CYCLE V10 MATERIAL
% =============================================================================
% =============================================================================
% BEGIN APPENDED FOURTEENTH-CYCLE V11 MATERIAL
% =============================================================================

\section{Fourteenth-cycle V11: physical-constraint factorization and
the contracted Bianchi--Ward defect}
\label{sec:v14v11}

\subsection{The exact constraint factorization}

Let \(C:{\cal H}\to{\cal K}\) be a bounded operator between Hilbert
spaces with closed range.  Its Moore--Penrose inverse
\(C^\dagger:{\cal K}\to{\cal H}\) is bounded, and
\[
 Q_C:=C^\dagger C,\qquad P_C:=I-Q_C                         \tag{11.1}
\]
are the orthogonal projections onto
\((\ker C)^\perp\) and \(\ker C\), respectively.

\begin{theorem}[Constraint multiplier and physical remainder]
\label{thm:v14v11-factor}
For every bounded \(D:{\cal H}\to{\cal Y}\),
\[
 \boxed{
 D=(DC^\dagger)C+DP_C.}                                    \tag{11.2}
\]
The following conditions are equivalent:
\begin{enumerate}
\item \(D\) vanishes on \(\ker C\);
\item \(DP_C=0\);
\item there exists a bounded \(Y:{\cal K}\to{\cal Y}\) such that
      \(D=YC\).
\end{enumerate}
The canonical multiplier
\[
 Y_0:=DC^\dagger                                            \tag{11.3}
\]
satisfies \(Y_0=Y_0CC^\dagger\) and
\[
 \|Y_0\|\leq\|D\|\,\|C^\dagger\|.                           \tag{11.4}
\]
\end{theorem}

\begin{proof}
Since \(Q_C+P_C=I\),
\[
 D=DQ_C+DP_C=DC^\dagger C+DP_C,
\]
which proves (11.2).  Items 1 and 2 are equivalent because \(P_C\)
is the projection onto \(\ker C\).  Item 3 implies item 1 since
\(C\) vanishes there.  If item 2 holds, (11.2) gives item 3 with
\(Y=Y_0\).  The Moore--Penrose identity
\(C^\dagger= C^\dagger CC^\dagger\) proves the support statement for
\(Y_0\), and submultiplicativity gives (11.4).
\end{proof}

\begin{theorem}[Best multiplier residual]
\label{thm:v14v11-best}
For every bounded \(D:{\cal H}\to{\cal Y}\),
\[
 \boxed{
 \inf_{Y\in{\cal B}({\cal K},{\cal Y})}
 \|D-YC\|
 =\|DP_C\|
 =\|D|_{\ker C}\|.}                                        \tag{11.5}
\]
The canonical multiplier \(Y_0=DC^\dagger\) attains the infimum.
\end{theorem}

\begin{proof}
For every \(Y\),
\[
 (D-YC)P_C=DP_C
\]
because \(CP_C=0\).  Hence
\(\|D-YC\|\geq\|DP_C\|\).  With \(Y=DC^\dagger\), equation
(11.2) gives \(D-YC=DP_C\), so equality is attained.
The operator norm of \(DP_C\) is exactly the norm of the restriction
of \(D\) to the range of \(P_C=\ker C\).
\end{proof}

\begin{remark}[Noncircular use]
Formula (11.2) does not prove that \(D\) factors through \(C\).
It displays the physical remainder \(DP_C\).  Declaring
\(D=DC^\dagger C\) while omitting that remainder would assume the
desired constraint-kernel equation.  This is the analogue of the
constructive-multiplier firewall in the audited YM V1.
\end{remark}

\subsection{The ordered geometric divergence defect}

Let \({\cal R}\) be the braided curvature and put
\[
 {\cal B}:=\nabla^{\mathrm{ad}}{\cal R}                     \tag{11.6}
\]
for its uncontracted Bianchi residual.  Let
\({\mathsf{Ric}}\) and \({\mathsf s}\) be the selected braided Ricci
and symmetric scalar contraction maps.  Define their covariant
intertwining defects by
\[
\begin{aligned}
 {\cal J}_{\mathsf{Ric}}(A)
 &:=
 {\mathfrak D}{\mathsf{Ric}}(A)
 -{\mathsf{Ric}}(\nabla^{\mathrm{ad}}A),\\
 {\cal J}_{\mathsf s}(A)
 &:=
 d\,{\mathsf s}(A)
 -{\mathsf s}(\nabla^{\mathrm{ad}}A).                       \tag{11.7}
\end{aligned}
\]
Let
\[
 {\cal Q}:={\mathfrak D}G                                   \tag{11.8}
\]
be the metricity defect in the target tensor calculus.  The
Jordan-ordered geometric tensor is
\[
 {\mathscr G}
 :=
 {\mathsf{Ric}}({\cal R})
 -\frac14\left(
   G\,{\mathsf s}({\cal R})
   +{\mathsf s}({\cal R})\,G
 \right).                                                    \tag{11.9}
\]

\begin{theorem}[Exact braided geometric divergence identity]
\label{thm:v14v11-geometric}
With every product retained in the displayed order,
\[
\boxed{
\begin{aligned}
 {\mathfrak D}{\mathscr G}
 ={}&
 {\cal J}_{\mathsf{Ric}}({\cal R})
 +{\mathsf{Ric}}({\cal B})\\
 &-\frac14\Bigl[
 {\cal Q}\,{\mathsf s}({\cal R})
 +G\bigl(
   {\cal J}_{\mathsf s}({\cal R})
   +{\mathsf s}({\cal B})
  \bigr)\\
 &\hspace{19mm}
 +\bigl(
   {\cal J}_{\mathsf s}({\cal R})
   +{\mathsf s}({\cal B})
  \bigr)G
 +{\mathsf s}({\cal R})\,{\cal Q}
 \Bigr].
\end{aligned}}                                               \tag{11.10}
\]
\end{theorem}

\begin{proof}
Equations (11.6)--(11.7) give
\[
\begin{aligned}
 {\mathfrak D}{\mathsf{Ric}}({\cal R})
 &={\cal J}_{\mathsf{Ric}}({\cal R})
   +{\mathsf{Ric}}({\cal B}),\\
 d\,{\mathsf s}({\cal R})
 &={\cal J}_{\mathsf s}({\cal R})
   +{\mathsf s}({\cal B}).                                  \tag{11.11}
\end{aligned}
\]
Apply the covariant Leibniz rule to both ordered products in (11.9),
insert (11.8) and (11.11), and collect terms.  No cyclic permutation
is used.
\end{proof}

Define the right-hand side of (11.10) as
\[
 {\mathfrak b}_{\mathrm{geo}}
 :={\mathfrak D}{\mathscr G}.                               \tag{11.12}
\]
For an exact differential graded connection with metricity and
perfect contraction intertwining, every term in (11.10) vanishes.
At finite stage, (11.10) identifies its source: Bianchi, Ricci,
scalar, and metricity defects.

\begin{corollary}[Quantitative geometric bound]
\label{cor:v14v11-geometricbound}
In compatible submultiplicative norms,
\[
\begin{aligned}
 \|{\mathfrak b}_{\mathrm{geo}}\|
 \leq{}&
 \|{\cal J}_{\mathsf{Ric}}({\cal R})\|
 +\|{\mathsf{Ric}}({\cal B})\|\\
 &+\frac12\|{\cal Q}\|\,\|{\mathsf s}({\cal R})\|\\
 &+\frac12\|G\|
 \left(
  \|{\cal J}_{\mathsf s}({\cal R})\|
  +\|{\mathsf s}({\cal B})\|
 \right).                                                    \tag{11.13}
\end{aligned}
\]
\end{corollary}

\begin{proof}
Take norms in (11.10), pair the two metricity terms and the two
scalar terms, and use submultiplicativity.
\end{proof}

\subsection{Same-source Ward comparison}

Let \({\mathbb T}^{\mathrm{ren}}\) be the same represented
renormalized stress tensor specified in V9, and define its Ward
defect
\[
 {\mathfrak w}_{T}
 :={\mathfrak D}{\mathbb T}^{\mathrm{ren}}.                 \tag{11.14}
\]
The represented Einstein residual is
\[
 {\mathbb R}_{\mathrm{Ein}}
 :={\mathscr G}-\kappa{\mathbb T}^{\mathrm{ren}}.            \tag{11.15}
\]

\begin{theorem}[Contracted Bianchi--Ward identity]
\label{thm:v14v11-bianchiward}
One has the exact identity
\[
 \boxed{
 {\mathfrak D}{\mathbb R}_{\mathrm{Ein}}
 =
 {\mathfrak b}_{\mathrm{geo}}
 -\kappa{\mathfrak w}_{T}.}                                 \tag{11.16}
\]
No Einstein equation is used in deriving (11.16).
\end{theorem}

\begin{proof}
Apply \({\mathfrak D}\) to (11.15), use linearity, and insert
(11.12) and (11.14).
\end{proof}

This equality distinguishes two statements.  Vanishing of
\({\mathfrak w}_T\) is stress conservation.  Vanishing of
\({\mathfrak b}_{\mathrm{geo}}\) is the contracted geometric Bianchi
identity for the selected braided contraction.  Their equality
controls the divergence of the Einstein residual; it does not make
the residual itself zero.

\subsection{Closure on the physical constraint kernel}

At a finite test screen, let
\[
 C_{\mathrm{phys}}:{\cal H}_{\mathrm{kin}}\to{\cal K}_{\mathrm{con}}
                                                                  \tag{11.17}
\]
be the complete gauge, BRST, Gauss, and diffeomorphism constraint
operator with closed range.  Put
\[
 P_{\mathrm{phys}}
 :=I-C_{\mathrm{phys}}^\dagger C_{\mathrm{phys}}.            \tag{11.18}
\]
Treat the divergence defect in (11.16) as the operator
\[
 D_{\mathrm{BW}}
 :={\mathfrak b}_{\mathrm{geo}}
   -\kappa{\mathfrak w}_T.                                  \tag{11.19}
\]

\begin{theorem}[Physical tangent-channel gate]
\label{thm:v14v11-physical}
The following are equivalent:
\begin{enumerate}
\item the contracted Bianchi--Ward defect vanishes on the physical
      constraint kernel;
\item
\[
 D_{\mathrm{BW}}P_{\mathrm{phys}}=0;                        \tag{11.20}
\]
\item there exists a bounded constraint multiplier \(Y_{\mathrm{BW}}\)
      such that
\[
 D_{\mathrm{BW}}
 =Y_{\mathrm{BW}}C_{\mathrm{phys}}.                         \tag{11.21}
\]
\end{enumerate}
The exact unexplained physical defect is
\[
 {\cal U}_{\mathrm{BW}}
 :=D_{\mathrm{BW}}P_{\mathrm{phys}},                         \tag{11.22}
\]
and
\[
 \|{\cal U}_{\mathrm{BW}}\|
 =
 \inf_Y\|D_{\mathrm{BW}}-YC_{\mathrm{phys}}\|.               \tag{11.23}
\]
\end{theorem}

\begin{proof}
Apply Theorems~\ref{thm:v14v11-factor} and
\ref{thm:v14v11-best} to \(C=C_{\mathrm{phys}}\) and
\(D=D_{\mathrm{BW}}\).
\end{proof}

\begin{corollary}[Sufficient typed defect conditions]
\label{cor:v14v11-sufficient}
A sufficient condition for (11.20) is
\[
\begin{aligned}
 {\cal B}P_{\mathrm{phys}}&=0,&
 {\cal Q}P_{\mathrm{phys}}&=0,\\
 {\cal J}_{\mathsf{Ric}}({\cal R})P_{\mathrm{phys}}&=0,&
 {\cal J}_{\mathsf s}({\cal R})P_{\mathrm{phys}}&=0,\\
 {\mathfrak w}_TP_{\mathrm{phys}}&=0,                       \tag{11.24}
\end{aligned}
\]
with the contractions acting compatibly on the constrained core.
More generally, nonzero geometric and Ward terms may cancel in
(11.19), but the cancellation must occur on the same state,
regulator history, and represented core.
\end{corollary}

\begin{proof}
Insert (11.24) in (11.10), then in (11.19).  The last statement is
just (11.20) without requiring the sufficient terms to vanish
separately.
\end{proof}

\subsection{Tangent closure is not the Einstein equation}

Let \(K_gX={\cal L}_Xg\) be the infinitesimal diffeomorphism map on
the classical tensor screen, with
\[
 K_g^*=-2\operatorname{div}_g.
\]
When the represented weak writer intertwines
\({\mathfrak D}\) with \(\operatorname{div}_g\), equation (11.20)
annuls the diffeomorphism-tangent component of the Einstein residual
on the physical core.  The archived V69 split then gives
\[
 Q_g{\mathbb R}_{\mathrm{Ein}}=0.                           \tag{11.25}
\]
The transverse component
\[
 P_g{\mathbb R}_{\mathrm{Ein}}                              \tag{11.26}
\]
remains unconstrained by (11.16).

\begin{firewall}[No conservation-to-Einstein inference]
\label{fire:v14v11-noeinstein}
Equations (11.20)--(11.25) close only the
diffeomorphism-tangent channel.  The full represented weak Einstein
equation still requires the localized residual Grams of V9, including
the transverse screen.  Stress conservation and the geometric
Bianchi identity cannot replace those Grams.
\end{firewall}

\subsection{Cofinal physical-kernel passage}

Let \(C_n\) be finite-stage constraints with closed range and
\[
 P_n:=I-C_n^\dagger C_n.                                    \tag{11.27}
\]
Let
\[
 D_n:={\mathfrak b}_{\mathrm{geo},n}
      -\kappa{\mathfrak w}_{T,n}.                           \tag{11.28}
\]

\begin{theorem}[Cofinal tangent-channel landing]
\label{thm:v14v11-cofinal}
Suppose, after transport to common Hilbert screens,
\[
 C_n\to C,\qquad P_n\to P_C,\qquad D_n\to D                \tag{11.29}
\]
in operator norm.  If
\[
 \|D_nP_n\|\longrightarrow0,                                \tag{11.30}
\]
then
\[
 DP_C=0,                                                     \tag{11.31}
\]
so the limiting contracted Bianchi--Ward defect factors through the
limiting constraint and vanishes on the physical kernel.

If
\[
 \sum_n\|D_nP_n\|<\infty                                    \tag{11.32}
\]
along the chosen diagonal, this row is compatible with the existing
summable overlap and observable ledgers.
\end{theorem}

\begin{proof}
Use
\[
\begin{aligned}
 \|DP_C\|
 \leq{}&
 \|(D-D_n)P_C\|
 +\|D_n(P_C-P_n)\|
 +\|D_nP_n\|.
\end{aligned}
\]
The first two terms tend to zero by (11.29) and uniform boundedness
of the convergent sequence \(D_n\); the third tends to zero by
(11.30).  Equation (11.31) and
Theorem~\ref{thm:v14v11-factor} give the factorization.  Summability
in (11.32) is already in the form required by a diagonal tail
estimate.
\end{proof}

The convergence \(P_n\to P_C\) requires the Killing-safe or
constraint-kernel reduced singular-value floor.  It is not implied by
\(C_n\to C\) alone.

\subsection{Acquisition ledger and next target}

\[
\begin{array}{p{0.29\linewidth}|p{0.23\linewidth}|p{0.38\linewidth}}
\textbf{gate}&\textbf{status after V11}&\textbf{content}\\ \hline
\text{constraint factorization}
 &\text{acquired}
 &D=DC^\dagger C+DP_C\\
\text{best multiplier residual}
 &\text{acquired}
 &\inf_Y\|D-YC\|=\|DP_C\|\\
\text{braided geometric divergence}
 &\text{expanded exactly}
 &formula (11.10)\\
\text{same-source Bianchi--Ward comparison}
 &\text{acquired}
 &D_{\mathrm{BW}}={\mathfrak b}_{\mathrm{geo}}
   -\kappa{\mathfrak w}_T\\
\text{physical tangent channel}
 &\text{reduced to a multiplier residual}
 &{\cal U}_{\mathrm{BW}}=D_{\mathrm{BW}}P_{\mathrm{phys}}\\
\text{cofinal tangent landing}
 &\text{acquired conditionally}
 &constraint projection and defect convergence\\
\text{transverse Einstein equation}
 &\text{open}
 &localized positive residual Grams from V9
\end{array}                                                     \tag{11.33}
\]

The next version should merge the physical constraint projection
\(P_{\mathrm{phys}}\) with the metric-dependent transverse projection
\(P_g\) without assuming that they commute.  A two-projection
intersection or angle theorem is required before the transverse Gram
can be restricted to the genuine physical tensor core.

\subsection{Parent-integrity record}

The installed parent was
\[
\begin{array}{c|c}
 \text{file}&
 \texttt{Renewal\_SM\_Einstein\_Continuum\_Research\_Notes\_V10.tex}\\
 \text{lines}&3679\\
 \text{bytes}&121271\\
 \text{SHA--256}&
 \texttt{3A7ACC6AE7812F39A99CCDED3EC2E87C522E8CD9A8A8974DF3A118447EB2B565}.
\end{array}                                                     \tag{11.34}
\]

% =============================================================================
% END APPENDED FOURTEENTH-CYCLE V11 MATERIAL
% =============================================================================
% =============================================================================
% BEGIN APPENDED FOURTEENTH-CYCLE V12 MATERIAL
% =============================================================================

\section{Fourteenth-cycle V12: the physical--transverse intersection
projection and its angle gap}
\label{sec:v14v12}

\subsection{Two projections on the common residual screen}

Let \({\cal Z}\) be the common Hilbert screen in which represented
Einstein residual vectors live.  Let
\[
 P=P^*=P^2                                                   \tag{12.1}
\]
be the physical constraint-kernel projection, and let
\[
 S=S^*=S^2,\qquad T:=I-S                                    \tag{12.2}
\]
be the metric transverse and diffeomorphism-tangent projections.
No commutation between \(P\) and \(S\) is assumed.

\begin{proposition}[The naive product is not an intersection
projection]
\label{prop:v14v12-productnogo}
There exist orthogonal projections \(P,S\) such that
\[
 PS\neq SP,\qquad (PS)^2\neq PS,\qquad
 \operatorname{Ran}P\cap\operatorname{Ran}S=\{0\},
\]
while \(PS\neq0\).
\end{proposition}

\begin{proof}
In \(\mathbb R^2\), let \(P\) project onto \(e_1\) and let \(S\)
project onto
\[
 u=(\cos\theta,\sin\theta),\qquad 0<\theta<\frac{\pi}{2}.
\]
The two one-dimensional ranges are distinct, so their intersection is
zero.  Direct matrix multiplication gives
\[
 PS=
 \begin{pmatrix}
 \cos^2\theta&\cos\theta\sin\theta\\
 0&0
 \end{pmatrix}\neq0,
\]
and
\((PS)^2=\cos^2\theta\,PS\neq PS\).  Its adjoint is \(SP\neq PS\).
\end{proof}

Thus multiplying the two projectors would reintroduce nonphysical
directions into the Gram screen.

\subsection{Spectral construction of the intersection}

Define the positive contraction
\[
 A:=PSP.                                                     \tag{12.3}
\]

\begin{lemma}[The unit eigenspace is the intersection]
\label{lem:v14v12-unit}
One has
\[
 \ker(I-A)
 =\operatorname{Ran}P\cap\operatorname{Ran}S.               \tag{12.4}
\]
\end{lemma}

\begin{proof}
If \(x\) belongs to both ranges, then \(Px=Sx=x\), hence
\(Ax=x\).

Conversely, suppose \(Ax=x\).  Since \(A=PSP\), the vector \(x\)
lies in \(\operatorname{Ran}P\).  Therefore
\[
 \|x\|^2
 =\langle x,PSPx\rangle
 =\langle x,Sx\rangle
 =\|Sx\|^2.
\]
For an orthogonal projection,
\(\|x-Sx\|^2=\|x\|^2-\|Sx\|^2\), so \(Sx=x\).  Thus \(x\) lies in
both ranges.
\end{proof}

\begin{definition}[Physical--transverse projector]
Assume the point \(1\) is isolated in \(\sigma(A)\).  Define
\[
 R:=\mathbf 1_{\{1\}}(A).                                   \tag{12.5}
\]
By Lemma~\ref{lem:v14v12-unit}, \(R\) is the orthogonal projection
onto
\[
 {\cal Z}_{\mathrm{phys,tr}}
 :=\operatorname{Ran}P\cap\operatorname{Ran}S.              \tag{12.6}
\]
\end{definition}

If \(P\) and \(S\) commute, then \(R=PS\).  Formula (12.5) is the
replacement that remains correct without commutation.

\subsection{The protected angle gap}

Define
\[
 \gamma(P,S)
 :=
 1-\sup\bigl(\sigma(PSP)\setminus\{1\}\bigr),               \tag{12.7}
\]
with the supremum taken on the orthogonal complement of the unit
eigenspace.  The condition
\[
 \gamma(P,S)>0                                              \tag{12.8}
\]
is the protected angle gap between the two ranges after their exact
intersection has been removed.

\begin{theorem}[Norm reconstruction by alternating compression]
\label{thm:v14v12-alternating}
If \(\gamma:=\gamma(P,S)>0\), then
\[
 \|(PSP)^n-R\|\leq(1-\gamma)^n,\qquad n\geq1.                \tag{12.9}
\]
Hence the physical--transverse projector is obtained in operator norm
from the two original projections.
\end{theorem}

\begin{proof}
The self-adjoint positive contraction \(A=PSP\) acts as the identity
on \(\operatorname{Ran}R\).  On its orthogonal complement, its
spectrum lies in \([0,1-\gamma]\).  The spectral theorem therefore
gives
\[
 A^n=R+(A|_{\operatorname{Ran}R^\perp})^n
\]
and the second term has norm at most \((1-\gamma)^n\).
\end{proof}

\begin{remark}
The angle gap permits a nontrivial exact intersection.  It charges
only the next spectral value below one, just as the Killing-safe
reduced minimum modulus charges only the nonzero singular block.
\end{remark}

\subsection{Perturbation and cofinal stability}

Let \(\widetilde P,\widetilde S\) be another pair of orthogonal
projections and put
\[
 \widetilde A:=\widetilde P\widetilde S\widetilde P.         \tag{12.10}
\]

\begin{lemma}[Compression perturbation]
\label{lem:v14v12-compression}
One has
\[
 \|A-\widetilde A\|
 \leq2\|P-\widetilde P\|+\|S-\widetilde S\|.                \tag{12.11}
\]
\end{lemma}

\begin{proof}
Expand
\[
\begin{aligned}
 PSP-\widetilde P\widetilde S\widetilde P
 ={}&(P-\widetilde P)SP
 +\widetilde P(S-\widetilde S)P\\
 &+\widetilde P\widetilde S(P-\widetilde P).
\end{aligned}
\]
All projections have norm one, so the triangle inequality gives
(12.11).
\end{proof}

\begin{theorem}[Intersection-projection perturbation bound]
\label{thm:v14v12-perturbation}
Assume
\[
 \sigma(A),\sigma(\widetilde A)
 \subset[0,1-\gamma]\cup\{1\}                               \tag{12.12}
\]
for some \(\gamma>0\), and set
\[
 \delta:=\|A-\widetilde A\|<\frac{\gamma}{2}.                \tag{12.13}
\]
Let \(R\) and \(\widetilde R\) be the unit-eigenspace projections.
Then
\[
 \boxed{
 \|R-\widetilde R\|
 \leq\frac{2\delta}{\gamma-2\delta}.}                        \tag{12.14}
\]
\end{theorem}

\begin{proof}
Use the circle \(\Gamma\) centered at \(1\) with radius
\(\gamma/2\), oriented positively.  By (12.12), it separates the
unit eigenvalue from the rest of each spectrum.  The Riesz formulas
are
\[
 R=\frac1{2\pi i}\int_\Gamma(z-A)^{-1}\,dz,\qquad
 \widetilde R=\frac1{2\pi i}
 \int_\Gamma(z-\widetilde A)^{-1}\,dz.                      \tag{12.15}
\]
On \(\Gamma\),
\[
 \|(z-A)^{-1}\|\leq\frac2\gamma.
\]
The resolvent perturbation bound and (12.13) give
\[
 \|(z-\widetilde A)^{-1}\|
 \leq\frac1{\gamma/2-\delta}.
\]
Using
\[
 (z-A)^{-1}-(z-\widetilde A)^{-1}
 =(z-A)^{-1}(A-\widetilde A)(z-\widetilde A)^{-1},
\]
the length \(\pi\gamma\) of \(\Gamma\), and (12.15), yields
\[
 \|R-\widetilde R\|
 \leq
 \frac{\pi\gamma}{2\pi}
 \frac2\gamma
 \delta
 \frac1{\gamma/2-\delta}
 =
 \frac{2\delta}{\gamma-2\delta}.
\]
\end{proof}

\begin{corollary}[Cofinal intersection landing]
\label{cor:v14v12-cofinal}
Let \(P_n\to P\) and \(S_n\to S\) in operator norm.  Suppose all
pairs have a common angle gap \(\gamma>0\) in the sense of (12.12).
Then their physical--transverse projectors satisfy
\[
 R_n\longrightarrow R                                      \tag{12.16}
\]
in operator norm.
\end{corollary}

\begin{proof}
Lemma~\ref{lem:v14v12-compression} gives
\(P_nS_nP_n\to PSP\).  Apply
Theorem~\ref{thm:v14v12-perturbation} for all sufficiently large
\(n\).
\end{proof}

The common gap is a real acquisition row.  Norm convergence of the
two original projectors alone does not license a claim about a
changing intersection rank.

\subsection{Completion of the residual screen}

For the localized residual vectors of V9, write
\[
 v_{j\ell}
 :=
 {\mathbb R}_{\mathrm{Ein}}(h_j)
 \pi_\omega(A_\ell)\Omega_\omega\in{\cal Z}.                \tag{12.17}
\]
Assume the physical constraint has already been imposed:
\[
 Pv_{j\ell}=v_{j\ell}.                                      \tag{12.18}
\]
The tangent-channel result of V11 is
\[
 Tv_{j\ell}=0.                                              \tag{12.19}
\]

\begin{theorem}[Physical plus tangent closure reduces the full
residual to the intersection]
\label{thm:v14v12-residual}
Under (12.18)--(12.19),
\[
 v_{j\ell}=Rv_{j\ell}.                                      \tag{12.20}
\]
Consequently the following are equivalent:
\[
 v_{j\ell}=0\ \hbox{for all \(j,\ell\)}
 \quad\Longleftrightarrow\quad
 Rv_{j\ell}=0\ \hbox{for all \(j,\ell\)}.                   \tag{12.21}
\]
Thus, after physical and tangent closure, localized positive Grams
of the \(R\)-projected residuals are sufficient for the full
represented weak Einstein equation on the declared core.
\end{theorem}

\begin{proof}
Equation (12.19) and \(S=I-T\) give \(Sv_{j\ell}=v_{j\ell}\).
Together with (12.18), this places \(v_{j\ell}\) in the intersection
(12.6).  The orthogonal projection \(R\) acts as the identity there,
which proves (12.20)--(12.21).  The localized Gram criterion of V9
then gives the weak equation.
\end{proof}

\begin{remark}[No commutation assumption]
The proof of Theorem~\ref{thm:v14v12-residual} never uses \(PS=SP\).
It imposes the two constraints on the actual residual vector and then
uses the true intersection projection.  This is the
constraint-before-positivity order selected by the companion audits.
\end{remark}

\subsection{Finite-stage Gram passage}

Let \(R_n\) be the protected physical--transverse projection and let
\(v_{j\ell,n}\to v_{j\ell}\) on every fixed screen.

\begin{proposition}[Projected Gram landing]
\label{prop:v14v12-gramlanding}
If \(R_n\to R\), then
\[
 \langle R_nv_{j\ell,n},R_nv_{kr,n}\rangle
 \longrightarrow
 \langle Rv_{j\ell},Rv_{kr}\rangle.                         \tag{12.22}
\]
If every fixed finite projected Gram tends to zero and the limiting
vectors satisfy (12.18)--(12.19), then the represented weak Einstein
equation holds on the V9 core.
\end{proposition}

\begin{proof}
Operator-norm convergence of \(R_n\), vector convergence, and uniform
boundedness give \(R_nv_{j\ell,n}\to Rv_{j\ell}\).  Inner products
therefore converge.  Zero limiting Grams imply
\(Rv_{j\ell}=0\); Theorem~\ref{thm:v14v12-residual} and the V9
localized Gram theorem finish the proof.
\end{proof}

\subsection{Acquisition ledger and next population problem}

\[
\begin{array}{p{0.29\linewidth}|p{0.23\linewidth}|p{0.38\linewidth}}
\textbf{gate}&\textbf{status after V12}&\textbf{content}\\ \hline
\text{naive product projection}
 &\text{refuted}
 &explicit two-line counterexample\\
\text{physical--transverse projector}
 &\text{acquired}
 &R=\mathbf1_{\{1\}}(PSP)\\
\text{noncommuting projection stability}
 &\text{acquired conditionally}
 &protected angle gap \(\gamma>0\)\\
\text{cofinal intersection landing}
 &\text{acquired}
 &R_n\to R\text{ under uniform gap}\\
\text{full residual after tangent closure}
 &\text{reduced to \(R\)-Grams}
 &v=Rv\\
\text{physical angle-gap population}
 &\text{open}
 &uniform gap for the regulator constraint and Bianchi screens\\
\text{transverse Gram vanishing}
 &\text{open}
 &same-history finite-stage residual estimates
\end{array}                                                     \tag{12.23}
\]

The next version should obtain an angle-gap certificate from a
constraint incidence operator, rather than assume spectral
separation of \(PSP\).  A reduced singular-value bound for the map
from the physical range to the tangent quotient can provide such a
certificate and expose the exact failure vector when the angle
collapses.

\subsection{Parent-integrity record}

The installed parent was
\[
\begin{array}{c|c}
 \text{file}&
 \texttt{Renewal\_SM\_Einstein\_Continuum\_Research\_Notes\_V11.tex}\\
 \text{lines}&4146\\
 \text{bytes}&135330\\
 \text{SHA--256}&
 \texttt{30B137E520F529579F73CEE1CCD96B4901F9F551104A576F41D05570F11D940F}.
\end{array}                                                     \tag{12.24}
\]

% =============================================================================
% END APPENDED FOURTEENTH-CYCLE V12 MATERIAL
% =============================================================================
% =============================================================================
% BEGIN APPENDED FOURTEENTH-CYCLE V13 MATERIAL
% =============================================================================

\section{Fourteenth-cycle V13: incidence singular values and the
physical--transverse angle gap}
\label{sec:v14v13}

\subsection{Purpose and relation to V12}

V12 identified the correct physical--transverse projector
\[
 R=\mathbf 1_{\{1\}}(PSP),
\]
but its cofinal stability required a uniform spectral gap below the
unit eigenspace of \(PSP\).  The present note removes the mystery from
that hypothesis.  The optimal gap is exactly the square of a reduced
singular value of the incidence from the physical subspace into the
tangent subspace.  When the tangent subspace is generated by the
linearized diffeomorphism operator, the same gap is certified by the
reduced singular value of the adjoint gauge incidence.

No continuum claim is inferred from this identity alone.  A regulator
family must still populate a uniform lower bound.  What is acquired
here is the precise operator whose lower bound suffices, the constants
which transfer that bound to V12, and an exact failure vector when the
bound collapses.

\subsection{Compressed angle operator}

Let \({\cal Z}\) be a Hilbert space and let \(P,S\in{\cal B}({\cal Z})\)
be orthogonal projections.  Put
\[
 T:=I-S,\qquad
 {\cal P}:=\operatorname{Ran}P,\qquad
 {\cal M}:={\cal P}\cap\operatorname{Ran}S.                 \tag{13.1}
\]
Thus \(T\) is the tangent-channel projection, \({\cal P}\) is the
physical subspace, and \({\cal M}\) is the physical--transverse
intersection.  Regard
\[
 A_{\cal P}:=(PSP)|_{\cal P}:{\cal P}\longrightarrow{\cal P},
 \qquad
 B:=T|_{\cal P}:{\cal P}\longrightarrow\operatorname{Ran}T. \tag{13.2}
\]
The codomain in (13.2) is a Hilbert space because \(T\) is an
orthogonal projection.

For a bounded operator \(L:H_1\to H_2\) with
\((\ker L)^\perp\ne\{0\}\), write
\[
 \operatorname{rmin}(L)
 :=
 \inf_{\substack{x\in(\ker L)^\perp\\\|x\|=1}}\|Lx\|.       \tag{13.3}
\]
This is its reduced minimum modulus.  The degenerate case
\((\ker L)^\perp=\{0\}\) will always be separated explicitly.

\begin{lemma}[Tangent incidence factorization]
\label{lem:v14v13-factor}
The operator \(B\) satisfies
\[
 \ker B={\cal M},\qquad
 B^*B=I_{\cal P}-A_{\cal P}.                                \tag{13.4}
\]
In particular, both \(A_{\cal P}\) and \(B^*B\) reduce the
orthogonal decomposition
\[
 {\cal P}={\cal M}\oplus({\cal P}\ominus{\cal M}).           \tag{13.5}
\]
\end{lemma}

\begin{proof}
For \(x\in{\cal P}\), \(Bx=0\) is equivalent to \(Tx=0\), hence to
\(x\in\operatorname{Ran}S\).  This proves the kernel identity.
The adjoint of \(B\), with the domain and codomain in (13.2), is
\(B^*=P|_{\operatorname{Ran}T}\).  Therefore, for \(x\in{\cal P}\),
\[
 B^*Bx=PTx=P(I-S)Px=x-PSPx,
\]
which is (13.4).  A bounded self-adjoint operator and its kernel
projection commute, so (13.5) is reducing.
\end{proof}

Assume for the moment that
\({\cal P}\ominus{\cal M}\ne\{0\}\), and set
\[
 \beta(P,S):=\operatorname{rmin}(B).                         \tag{13.6}
\]
Define the optimal V12 gap by
\[
 \gamma(P,S)
 :=
 \sup\left\{\gamma\geq0:
 \sigma\!\left(A_{\cal P}|_{{\cal P}\ominus{\cal M}}\right)
 \subset[0,1-\gamma]\right\}.                               \tag{13.7}
\]
The value in (13.7) is allowed to vanish.

\begin{theorem}[Exact angle-gap identity]
\label{thm:v14v13-exactgap}
If \({\cal P}\ominus{\cal M}\ne\{0\}\), then
\[
 \boxed{\gamma(P,S)=\beta(P,S)^2.}                           \tag{13.8}
\]
Consequently, the following statements are equivalent:
\[
 \gamma(P,S)>0,\qquad
 \beta(P,S)>0,\qquad
 \operatorname{Ran}B\ \hbox{is closed}.                     \tag{13.9}
\]
If \({\cal P}={\cal M}\), then \(R=P\) and there is no complementary
angle channel to control.
\end{theorem}

\begin{proof}
On \({\cal P}\ominus{\cal M}\), Lemma~\ref{lem:v14v13-factor}
gives
\[
 A_{\cal P}=I-B^*B.                                         \tag{13.10}
\]
The spectral theorem for the positive operator \(B^*B\) gives
\[
 \inf\sigma\!\left(
 B^*B|_{{\cal P}\ominus{\cal M}}\right)
 =
 \inf_{\substack{x\in{\cal P}\ominus{\cal M}\\\|x\|=1}}
 \langle B^*Bx,x\rangle
 =
 \beta(P,S)^2.                                              \tag{13.11}
\]
Applying the spectral mapping theorem to (13.10) shows that the
right side of (13.11) is precisely the distance from the
complementary spectrum of \(A_{\cal P}\) to \(1\).  This proves
(13.8).

A bounded operator has closed range if and only if it is bounded
below on the orthogonal complement of its kernel.  For completeness,
if \(\beta>0\), a convergent sequence \(Bx_n\), with
\(x_n\perp\ker B\), makes \(x_n\) Cauchy and hence has a preimage.
Conversely, when the range is closed, the inverse of
\(B:(\ker B)^\perp\to\operatorname{Ran}B\) is bounded by the open
mapping theorem, which gives \(\beta>0\).  Together with (13.8), this
proves (13.9).
\end{proof}

\begin{corollary}[Alternating-compression rate with an exact constant]
\label{cor:v14v13-alternating}
Let \(R\) be the projection onto \({\cal M}\).  If
\(\beta(P,S)>0\), then
\[
 \|(PSP)^m-R\|
 \leq (1-\beta(P,S)^2)^m,\qquad m\geq1.                     \tag{13.12}
\]
The base in (13.12) is optimal whenever the top of the complementary
spectrum belongs to the spectrum.
\end{corollary}

\begin{proof}
On \({\cal M}\), \(PSP\) is the identity.  On its orthogonal
complement in \({\cal P}\), Theorem~\ref{thm:v14v13-exactgap} places
the spectrum in \([0,1-\beta^2]\).  On
\({\cal P}^{\perp}\), \(PSP=0\).  Functional calculus gives (13.12).
If \(1-\beta^2\) is spectral, the norm on the complementary reducing
subspace equals its \(m\)-th power.
\end{proof}

\subsection{Certificate from the diffeomorphism incidence}

The canonical operator \(B=TP\) already identifies the exact gap, but
a geometric regulator normally presents the tangent subspace as the
range of an infinitesimal gauge generator.  Let \({\cal X}\) be a
Hilbert space and let
\[
 K:{\cal X}\longrightarrow{\cal Z}                           \tag{13.13}
\]
be bounded with closed range
\[
 \operatorname{Ran}K=\operatorname{Ran}T.                    \tag{13.14}
\]
In an Einstein realization, \(K\) is the represented infinitesimal
diffeomorphism incidence on the chosen Sobolev or finite-stage
screen.  Define
\[
 \kappa_K:=\operatorname{rmin}(K)>0,\qquad M_K:=\|K\|,        \tag{13.15}
\]
and introduce the physical adjoint incidence
\[
 J:=K^*|_{\cal P}:{\cal P}\longrightarrow{\cal X}.           \tag{13.16}
\]

\begin{lemma}[Kernel of the adjoint incidence]
\label{lem:v14v13-jkernel}
Under (13.14),
\[
 \ker J={\cal M}.                                            \tag{13.17}
\]
\end{lemma}

\begin{proof}
For \(x\in{\cal P}\),
\[
 K^*x=0
 \quad\Longleftrightarrow\quad
 x\perp\operatorname{Ran}K
 \quad\Longleftrightarrow\quad
 x\perp\operatorname{Ran}T
 \quad\Longleftrightarrow\quad
 x\in\operatorname{Ran}S.
\]
Intersecting with \({\cal P}\) proves (13.17).
\end{proof}

Suppose again that \({\cal P}\ominus{\cal M}\ne\{0\}\), and put
\[
 \alpha(K,P):=\operatorname{rmin}(J).                        \tag{13.18}
\]

\begin{theorem}[Diffeomorphism-incidence angle certificate]
\label{thm:v14v13-incidence}
Under (13.13)--(13.18),
\[
 \boxed{
 \frac{\alpha(K,P)}{M_K}
 \leq\beta(P,S)
 \leq\frac{\alpha(K,P)}{\kappa_K}.}                          \tag{13.19}
\]
Hence
\[
 \boxed{
 \left(\frac{\alpha(K,P)}{M_K}\right)^2
 \leq\gamma(P,S)
 \leq
 \left(\frac{\alpha(K,P)}{\kappa_K}\right)^2.}               \tag{13.20}
\]
In particular, a positive reduced singular floor for \(K^*P\),
together with an upper norm bound for \(K\), populates the V12 angle
gap.
\end{theorem}

\begin{proof}
Because \(K^*S=0\), every \(x\in{\cal P}\) satisfies
\[
 Jx=K^*x=K^*Tx.                                              \tag{13.21}
\]
The upper estimate
\[
 \|Jx\|\leq M_K\|Tx\|                                       \tag{13.22}
\]
is immediate.

We also need the lower estimate for \(K^*\) on
\(\operatorname{Ran}K\).  In the polar decomposition
\(K=U|K|\), the restriction of \(|K|\) to
\((\ker K)^\perp\) is bounded below by \(\kappa_K\).
If \(y\in\operatorname{Ran}K\), then \(U^*y\in(\ker K)^\perp\),
\(\|U^*y\|=\|y\|\), and
\[
 \|K^*y\|=\||K|U^*y\|\geq\kappa_K\|y\|.                     \tag{13.23}
\]
Apply (13.22)--(13.23) to \(y=Tx\), and restrict \(x\) to the common
kernel complement
\({\cal P}\ominus{\cal M}\), identified by
Lemmas~\ref{lem:v14v13-factor} and \ref{lem:v14v13-jkernel}.  Taking
infima over its unit sphere gives
\[
 \kappa_K\beta(P,S)\leq\alpha(K,P)
 \leq M_K\beta(P,S),
\]
which is (13.19).  Square and use
Theorem~\ref{thm:v14v13-exactgap} to obtain (13.20).
\end{proof}

\begin{remark}[Normalized tangent generator]
If \(K\) is the inclusion
\(\operatorname{Ran}T\hookrightarrow{\cal Z}\), then
\(\kappa_K=M_K=1\) and \(J=TP\).  Both inequalities in (13.19)
become equalities.  A nonisometric geometric generator changes only
the conditioning constants in (13.19).
\end{remark}

\subsection{Sharp collapse witness}

\begin{proposition}[Loss of the singular floor permits intersection
jump]
\label{prop:v14v13-collapse}
Let \({\cal Z}=\mathbb R^2\), let \(P\) project onto
\(\mathbb Re_1\), and for \(0<\theta<\pi/2\) let \(S_\theta\) project
onto
\[
 u_\theta=(\cos\theta,\sin\theta).
\]
Then
\[
 {\cal M}_\theta=\{0\},\qquad
 \beta(P,S_\theta)=\sin\theta,\qquad
 \gamma(P,S_\theta)=\sin^2\theta,\qquad
 R_\theta=0.                                                \tag{13.24}
\]
Nevertheless \(S_\theta\to P\) as \(\theta\downarrow0\), while the
limiting intersection projector is \(R_0=P\).
\end{proposition}

\begin{proof}
For positive \(\theta\), the two one-dimensional ranges are distinct,
so their intersection is zero.  With
\(T_\theta=I-S_\theta\),
\[
 \|T_\theta e_1\|=\sin\theta,
 \qquad
 PS_\theta P|_{\mathbb Re_1}
 =\cos^2\theta\,I.                                          \tag{13.25}
\]
Equations (13.24) follow directly from (13.6)--(13.8).  The matrix
entries of \(S_\theta\) converge to those of \(P\), but at
\(\theta=0\) the ranges coincide.  Thus the exact failure vector
\(e_1\) has tangent incidence tending to zero and becomes a new
intersection vector at the limit.
\end{proof}

This example proves that the singular floor is structural.  It cannot
be removed from the cofinal projector argument and recovered from
norm convergence of \(P_n,S_n\) alone.

\subsection{Cofinal incidence theorem}

\begin{theorem}[Uniform incidence floor implies cofinal intersection
landing]
\label{thm:v14v13-cofinal}
For each \(n\), let \(P_n,S_n\) be orthogonal projections on
\({\cal Z}\), put \(T_n=I-S_n\), and let
\(K_n:{\cal X}_n\to{\cal Z}\) be bounded with closed range
\[
 \operatorname{Ran}K_n=\operatorname{Ran}T_n.                \tag{13.26}
\]
Assume
\[
 P_n\to P,\qquad S_n\to S                                   \tag{13.27}
\]
in operator norm.  Let
\[
 {\cal M}_n=\operatorname{Ran}P_n\cap\operatorname{Ran}S_n.
\]
Whenever \(\operatorname{Ran}P_n\ominus{\cal M}_n\ne\{0\}\),
assume the uniform bounds
\[
 \|K_n\|\leq M,\qquad
 \operatorname{rmin}
 \bigl(K_n^*|_{\operatorname{Ran}P_n}\bigr)\geq a            \tag{13.28}
\]
for constants \(0<a\leq M<\infty\).  In the degenerate case
\(\operatorname{Ran}P_n={\cal M}_n\), assume directly that it persists
under the declared screen identification.  Then the intersection
projections
\[
 R_n=\operatorname{proj}_{{\cal M}_n}
\]
converge in norm to
\[
 R=\operatorname{proj}_{\operatorname{Ran}P\cap
                            \operatorname{Ran}S}.            \tag{13.29}
\]
\end{theorem}

\begin{proof}
For every nondegenerate stage,
Theorem~\ref{thm:v14v13-incidence} gives
\[
 \gamma(P_n,S_n)\geq\gamma_0:=a^2/M^2>0.                    \tag{13.30}
\]
Thus
\[
 \sigma\!\left((P_nS_nP_n)|_{\operatorname{Ran}P_n}\right)
 \subset[0,1-\gamma_0]\cup\{1\}.                            \tag{13.31}
\]
The full-space compression has only the additional spectral point
\(0\).  By (13.27),
\[
 P_nS_nP_n\longrightarrow PSP                               \tag{13.32}
\]
in norm.  Spectral stability for self-adjoint operators shows that
the limiting spectrum also lies in
\([0,1-\gamma_0]\cup\{1\}\); in particular no complementary
spectral point can enter the open interval
\((1-\gamma_0,1)\).  The Riesz-projection estimate of V12 applied to
(13.32) now yields \(R_n\to R\).  The separately declared degenerate
case is immediate.
\end{proof}

\begin{corollary}[Cofinal projected-Gram landing from incidence data]
\label{cor:v14v13-gram}
Under the hypotheses of
Theorem~\ref{thm:v14v13-cofinal}, let
\(v_{j\ell,n}\to v_{j\ell}\) on every fixed V9 residual screen.
Then
\[
 \langle R_nv_{j\ell,n},R_nv_{kr,n}\rangle
 \longrightarrow
 \langle Rv_{j\ell},Rv_{kr}\rangle.                         \tag{13.33}
\]
If the limiting residuals satisfy the physical and tangent equations
of V12 and all limiting diagonal quantities in (13.33) vanish, the
represented weak Einstein equation holds on the declared core.
\end{corollary}

\begin{proof}
Theorem~\ref{thm:v14v13-cofinal} supplies \(R_n\to R\).
The projected-Gram proposition of V12 then applies verbatim.
Diagonal vanishing gives \(Rv_{j\ell}=0\), and the V12 intersection
reduction gives \(v_{j\ell}=0\).
\end{proof}

\subsection{Regulator acquisition rows}

For a concrete finite regulator, the certificate can be checked in
the following order:
\[
\begin{array}{c|p{0.69\linewidth}}
1&
construct the physical projector \(P_n\), for example from an exact
closed-range constraint kernel;\\
2&
construct the infinitesimal diffeomorphism incidence \(K_n\) and
verify \(\operatorname{Ran}K_n=\operatorname{Ran}T_n\);\\
3&
identify
\(\ker(K_n^*|_{\operatorname{Ran}P_n})={\cal M}_n\) exactly;\\
4&
bound \(K_n^*P_n\) below on
\(\operatorname{Ran}P_n\ominus{\cal M}_n\), and bound
\(\|K_n\|\) above;\\
5&
insert \(a^2/M^2\) into the V12 Riesz and projected-Gram estimates.
\end{array}                                                  \tag{13.34}
\]
In finite dimension, row 4 is the least positive singular value of
the displayed incidence matrix.  In the continuum it is a
closed-range or coercivity estimate on the declared Sobolev
completion.  A numerical singular value without an exact error
interval does not populate the analytic row.

\subsection{Acquisition ledger and next upstream move}

\[
\begin{array}{p{0.29\linewidth}|p{0.23\linewidth}|p{0.38\linewidth}}
\textbf{gate}&\textbf{status after V13}&\textbf{content}\\ \hline
\text{optimal physical--transverse gap}
 &\text{acquired}
 &\gamma(P,S)=\operatorname{rmin}(TP|_{\operatorname{Ran}P})^2\\
\text{geometric incidence certificate}
 &\text{acquired conditionally}
 &\gamma\geq
 \operatorname{rmin}(K^*|_{\operatorname{Ran}P})^2/\|K\|^2\\
\text{gap-collapse witness}
 &\text{acquired}
 &rotating lines and the vector \(e_1\)\\
\text{cofinal projector landing}
 &\text{reduced to a uniform incidence floor}
 &R_n\to R\\
\text{regulator singular floor}
 &\text{open}
 &same-history coercivity for \(K_n^*P_n\)\\
\text{transverse residual Grams}
 &\text{open}
 &same-history finite-stage estimates
\end{array}                                                  \tag{13.35}
\]

The next upstream step is to package the physical and tangent
conditions into one joint constraint operator.  Its kernel should be
exactly \({\cal M}\).  The angle certificate above should then prove
closed range, produce an explicit Moore--Penrose bound, and identify
the exact best multiplier remainder as \(R\) applied to the residual.
That construction will connect the constraint-multiplier lesson from
the V10 Yang--Mills audit directly to the Einstein residual screen.

\subsection{Parent-integrity record}

The installed parent was
\[
\begin{array}{c|c}
 \text{file}&
 \texttt{Renewal\_SM\_Einstein\_Continuum\_Research\_Notes\_V12.tex}\\
 \text{lines}&4533\\
 \text{bytes}&146768\\
 \text{SHA--256}&
 \texttt{BA93732BED2D4D6C7C03ADA3DB248238652EE2A9952D40FBEE31E9A7634045D0}.
\end{array}                                                  \tag{13.36}
\]

% =============================================================================
% END APPENDED FOURTEENTH-CYCLE V13 MATERIAL
% =============================================================================
% =============================================================================
% BEGIN APPENDED FOURTEENTH-CYCLE V14 MATERIAL
% =============================================================================

\section{Fourteenth-cycle V14: joint physical--tangent constraints,
exact multipliers, and the intersection remainder}
\label{sec:v14v14}

\subsection{Why a joint operator is needed}

The V10 Yang--Mills audit showed that positivity must be applied after
the physical constraint has been encoded, and V11 gave an exact
factorization for a single closed-range constraint.  V12--V13 added a
second, generally noncommuting, condition: removal of the
diffeomorphism-tangent channel.  Multiplying their two projections is
not legitimate.  This note instead stacks the two violations into one
operator.  Its kernel is the true physical--transverse intersection,
and the V13 angle floor proves that this stacked operator has closed
range.

The construction distinguishes two conclusions which must not be
conflated.  Vanishing of the joint constrained remainder says that a
residual is a sum of physical-constraint and tangent multipliers.
It says that the residual itself vanishes only after the residual is
also known to lie in the joint kernel.  This is the exact
constraint-before-positivity logic needed by the represented weak
Einstein screen.

\subsection{The stacked constraint}

Retain the V13 Hilbert space \({\cal Z}\) and orthogonal projections
\(P,S\).  Set
\[
 E:=I-P,\qquad T:=I-S,\qquad
 {\cal M}:=\operatorname{Ran}P\cap\operatorname{Ran}S.       \tag{14.1}
\]
Here \(E\) measures physical-constraint violation and \(T\) measures
the tangent component.  On the fixed Hilbert space
\({\cal W}:={\cal Z}\oplus{\cal Z}\), define
\[
 {\mathfrak C}:{\cal Z}\longrightarrow{\cal W},\qquad
 {\mathfrak C}x:=(Ex,Tx).                                   \tag{14.2}
\]

\begin{lemma}[Kernel and normal operator]
\label{lem:v14v14-kernel}
The stacked constraint satisfies
\[
 \ker{\mathfrak C}={\cal M},\qquad
 {\mathfrak C}^*(u,v)=Eu+Tv,\qquad
 {\mathfrak C}^*{\mathfrak C}=E+T.                          \tag{14.3}
\]
Moreover
\[
 {\mathfrak C}{\mathfrak C}^*
 =
 \begin{pmatrix}
 E&ET\\
 TE&T
 \end{pmatrix}.                                             \tag{14.4}
\]
The off-diagonal blocks in (14.4) retain the noncommutativity of the
two constraints.
\end{lemma}

\begin{proof}
Equation (14.2) vanishes precisely when \(Px=x\) and \(Sx=x\),
which gives the kernel.  For \(x,u,v\in{\cal Z}\),
\[
 \langle{\mathfrak C}x,(u,v)\rangle
 =\langle Ex,u\rangle+\langle Tx,v\rangle
 =\langle x,Eu+Tv\rangle.
\]
This proves the adjoint formula.  Since \(E^2=E\) and \(T^2=T\),
the two displayed normal-operator formulas follow by multiplication.
\end{proof}

Let
\[
 \beta:=\operatorname{rmin}
 \bigl(T|_{\operatorname{Ran}P}\bigr)                       \tag{14.5}
\]
when
\(\operatorname{Ran}P\ominus{\cal M}\ne\{0\}\).  By V13,
\(\beta^2\) is the optimal physical--transverse angle gap.

\begin{theorem}[Joint coercivity from the angle floor]
\label{thm:v14v14-coercivity}
Assume \(\beta>0\).  Then for every \(x\perp{\cal M}\),
\[
 \|{\mathfrak C}x\|
 \geq d(\beta)\|x\|,\qquad
 d(\beta):=
 \frac{\beta}{\sqrt{1+(1+\beta)^2}}.                        \tag{14.6}
\]
Consequently \(\operatorname{Ran}{\mathfrak C}\) is closed and
\[
 \|{\mathfrak C}^{\dagger}\|
 \leq\frac{\sqrt{1+(1+\beta)^2}}{\beta}.                    \tag{14.7}
\]
Conversely, if \({\mathfrak C}\) has closed range, then
\[
 \beta\geq\operatorname{rmin}({\mathfrak C})>0.              \tag{14.8}
\]
Thus, outside the trivial case
\(\operatorname{Ran}P={\cal M}\), the V13 angle gap is positive if
and only if the joint constraint has closed range.
\end{theorem}

\begin{proof}
Take \(x\perp{\cal M}\), and decompose it orthogonally as
\[
 x=p+q,\qquad p=Px,\qquad q=Ex.                              \tag{14.9}
\]
For \(m\in{\cal M}\), \(Pm=m\), and hence
\(\langle p,m\rangle=\langle x,m\rangle=0\).  Therefore
\(p\in\operatorname{Ran}P\ominus{\cal M}\), and (14.5) gives
\[
 \beta\|p\|\leq\|Tp\|
 \leq\|Tx\|+\|Tq\|
 \leq\|Tx\|+\|q\|.                                          \tag{14.10}
\]
It follows that
\[
 \begin{aligned}
 \|x\|
 &\leq\|p\|+\|q\|\\
 &\leq\beta^{-1}\|Tx\|+(1+\beta^{-1})\|q\|\\
 &\leq
 \frac{\sqrt{1+(1+\beta)^2}}{\beta}
 \bigl(\|Tx\|^2+\|q\|^2\bigr)^{1/2}\\
 &=d(\beta)^{-1}\|{\mathfrak C}x\|.
 \end{aligned}                                              \tag{14.11}
\]
This proves (14.6).  The closed-range theorem and the identity
\(\|{\mathfrak C}^{\dagger}\|
 =\operatorname{rmin}({\mathfrak C})^{-1}\)
give (14.7).

Conversely, let
\(x\in\operatorname{Ran}P\ominus{\cal M}\).  Then \(Ex=0\) and
\[
 \|Tx\|=\|{\mathfrak C}x\|
 \geq\operatorname{rmin}({\mathfrak C})\|x\|.                \tag{14.12}
\]
Taking the infimum proves (14.8).  The final equivalence follows from
V13.
\end{proof}

\begin{corollary}[Geometric incidence populates joint coercivity]
\label{cor:v14v14-incidence}
Let \(K:{\cal X}\to{\cal Z}\) be the V13 closed-range
diffeomorphism incidence with
\(\operatorname{Ran}K=\operatorname{Ran}T\).  If
\[
 \|K\|\leq M,\qquad
 \operatorname{rmin}
 \bigl(K^*|_{\operatorname{Ran}P}\bigr)\geq a>0,             \tag{14.13}
\]
then
\[
 \operatorname{rmin}({\mathfrak C})
 \geq d(a/M),\qquad
 \|{\mathfrak C}^{\dagger}\|
 \leq d(a/M)^{-1}.                                          \tag{14.14}
\]
\end{corollary}

\begin{proof}
The V13 incidence theorem gives \(\beta\geq a/M\).
The square of the function in (14.6) is
\[
 d(t)^2=\frac{t^2}{t^2+2t+2},
\]
whose derivative is positive for \(t>0\).  Apply
Theorem~\ref{thm:v14v14-coercivity}.
\end{proof}

\subsection{The exact intersection projector}

Let \(R\) denote the orthogonal projection onto \({\cal M}\), as in
V12.

\begin{theorem}[Moore--Penrose formula for the true intersection]
\label{thm:v14v14-mp}
Under the equivalent closed-range conditions of
Theorem~\ref{thm:v14v14-coercivity},
\[
 \boxed{R=I-{\mathfrak C}^{\dagger}{\mathfrak C}.}           \tag{14.15}
\]
Equivalently,
\[
 I-R
 ={\mathfrak C}^{\dagger}{\mathfrak C}
 ={\mathfrak C}^*
  ({\mathfrak C}{\mathfrak C}^*)^{\dagger}
  {\mathfrak C}.                                            \tag{14.16}
\]
Thus
\[
 {\cal M}^{\perp}
 =\operatorname{Ran}{\mathfrak C}^*
 =\operatorname{Ran}E+\operatorname{Ran}T,                  \tag{14.17}
\]
and the sum in (14.17) is closed.
\end{theorem}

\begin{proof}
For any closed-range operator \(L\), \(L^\dagger L\) is the
orthogonal projection onto \((\ker L)^\perp\).  Apply this to
\(L={\mathfrak C}\) and use
Lemma~\ref{lem:v14v14-kernel}.  This proves (14.15).  The standard
closed-range Moore--Penrose identity
\[
 L^\dagger L=L^*(LL^*)^\dagger L
\]
gives (14.16).  Also
\[
 \operatorname{Ran}{\mathfrak C}^*
 =\operatorname{Ran}E+\operatorname{Ran}T
\]
by the adjoint formula in (14.3).  A bounded operator and its adjoint
have closed range simultaneously, so this range is closed and equals
\((\ker{\mathfrak C})^\perp={\cal M}^\perp\).
\end{proof}

Formula (14.16), not a product such as \(PS\), is the normalized
noncommuting intersection compiler.  The mixed blocks \(ET\) and
\(TE\) in (14.4) are automatically included in its pseudoinverse.

\subsection{Exact vector multiplier remainder}

\begin{theorem}[Best joint multiplier approximation]
\label{thm:v14v14-vector}
For every \(v\in{\cal Z}\),
\[
 \boxed{
 v={\mathfrak C}^*\lambda_v+Rv,\qquad
 \lambda_v:=({\mathfrak C}^{\dagger})^*v.}                  \tag{14.18}
\]
The two terms in (14.18) are orthogonal, \(\lambda_v\) is the
minimum-norm multiplier producing \((I-R)v\), and
\[
 \boxed{
 \inf_{\lambda\in{\cal W}}
 \|v-{\mathfrak C}^*\lambda\|=\|Rv\|.}                      \tag{14.19}
\]
Moreover
\[
 \|\lambda_v\|
 \leq d(\beta)^{-1}\|v\|.                                  \tag{14.20}
\]
\end{theorem}

\begin{proof}
Taking adjoints in
\({\mathfrak C}^{\dagger}{\mathfrak C}=I-R\) gives
\[
 {\mathfrak C}^*({\mathfrak C}^{\dagger})^*=I-R,
\]
which proves (14.18).  Its first term lies in
\({\cal M}^{\perp}\) and its second in \({\cal M}\), so they are
orthogonal.

The Moore--Penrose solution
\(({\mathfrak C}^{\dagger})^*v\) is orthogonal to
\(\ker{\mathfrak C}^*\), and is therefore the unique solution of
minimum norm.  For any \(\lambda\),
\[
 R(v-{\mathfrak C}^*\lambda)=Rv,
\]
so the left side of (14.19) is at least \(\|Rv\|\).
The choice \(\lambda=\lambda_v\) attains equality by (14.18).
Finally (14.20) follows from (14.7).
\end{proof}

Writing
\(\lambda_v=(\lambda_v^{\rm phys},\lambda_v^{\rm tan})\),
equation (14.18) reads
\[
 v=E\lambda_v^{\rm phys}
   +T\lambda_v^{\rm tan}
   +Rv.                                                     \tag{14.21}
\]
The two multiplier terms in (14.21) need not be orthogonal to one
another.  Only their joint sum is the orthogonal
\({\cal M}^{\perp}\)-component.

\begin{corollary}[Augmented equation versus full equation]
\label{cor:v14v14-meaning}
For \(v\in{\cal Z}\), the following are equivalent:
\[
 Rv=0
 \quad\Longleftrightarrow\quad
 v={\mathfrak C}^*\lambda
 \ \hbox{for some }\lambda\in{\cal W}.                      \tag{14.22}
\]
If in addition
\[
 {\mathfrak C}v=0,                                          \tag{14.23}
\]
then (14.22) is equivalent to \(v=0\).
\end{corollary}

\begin{proof}
Equation (14.22) is (14.18)--(14.19).  Under (14.23),
\(v\in{\cal M}\), so \(Rv=v\).  Hence \(Rv=0\) is then exactly
\(v=0\).
\end{proof}

Condition (14.23) says simultaneously that the residual is physical
and has no tangent component.  Without it, a zero \(R\)-Gram proves
only the augmented multiplier equation in (14.22).

\subsection{Operator-valued residual factorization}

The same geometry controls a residual writer
\(D:{\cal Z}\to{\cal Y}\), where \({\cal Y}\) is any Hilbert space.

\begin{theorem}[Exact row factorization and optimal defect]
\label{thm:v14v14-row}
For every bounded \(D:{\cal Z}\to{\cal Y}\),
\[
 \boxed{
 D=L_D{\mathfrak C}+DR,\qquad
 L_D:=D{\mathfrak C}^{\dagger}.}                            \tag{14.24}
\]
Furthermore,
\[
 \boxed{
 \inf_{L\in{\cal B}({\cal W},{\cal Y})}
 \|D-L{\mathfrak C}\|=\|DR\|,}                              \tag{14.25}
\]
and
\[
 D=L{\mathfrak C}\ \hbox{for some bounded \(L\)}
 \quad\Longleftrightarrow\quad DR=0.                        \tag{14.26}
\]
The canonical factor obeys
\[
 \|L_D\|\leq d(\beta)^{-1}\|D\|.                            \tag{14.27}
\]
\end{theorem}

\begin{proof}
Insert \(I={\mathfrak C}^{\dagger}{\mathfrak C}+R\) from
(14.15) on the right of \(D\), which gives (14.24).  For arbitrary
\(L\),
\[
 (D-L{\mathfrak C})R=DR
\]
because \({\mathfrak C}R=0\).  Therefore
\(\|D-L{\mathfrak C}\|\geq\|DR\|\).  The choice
\(L=L_D\) attains equality, proving (14.25).  Equation (14.26)
follows, and (14.27) is (14.7).
\end{proof}

This is the two-constraint extension of the V11 factorization.  It
does not assume \(P\) and \(S\) commute and it supplies the exact
distance to joint constraint factorization.

\subsection{Localized positive Gram formulation}

Let \(v_a\in{\cal Z}\) denote any finite set of localized represented
Einstein residual vectors; the index \(a\) may stand for the V9 pair
\((j,\ell)\).  Define their optimally augmented residues
\[
 w_a:=v_a-{\mathfrak C}^*\lambda_{v_a}=Rv_a.                 \tag{14.28}
\]
Their Gram matrix is
\[
 G_{ab}^{\rm joint}
 :=\langle w_a,w_b\rangle
 =\langle Rv_a,Rv_b\rangle.                                \tag{14.29}
\]

\begin{proposition}[What the joint Gram proves]
\label{prop:v14v14-gram}
For a finite screen, the following are equivalent:
\[
 G_{aa}^{\rm joint}=0\ \hbox{for every \(a\)}
 \quad\Longleftrightarrow\quad
 v_a={\mathfrak C}^*\lambda_{v_a}\ \hbox{for every \(a\)}.   \tag{14.30}
\]
If also \({\mathfrak C}v_a=0\) for every \(a\), then these statements
are equivalent to
\[
 v_a=0\quad\hbox{for every \(a\)}.                          \tag{14.31}
\]
\end{proposition}

\begin{proof}
Since \(G_{aa}^{\rm joint}=\|Rv_a\|^2\), (14.30) follows from
Corollary~\ref{cor:v14v14-meaning}.  Under the additional joint-kernel
condition, the same corollary gives (14.31).
\end{proof}

Thus the correct positive screen is formed from the augmented
residuals after the joint constraint has been assembled.  The
unaugmented full equation requires the additional kernel membership;
this is exactly the physical-plus-tangent hypothesis used in the V12
residual theorem.

\subsection{Cofinal Moore--Penrose landing}

Let \(P_n\to P\) and \(S_n\to S\) in norm, and define on the fixed
spaces
\[
 E_n=I-P_n,\quad T_n=I-S_n,\quad
 {\mathfrak C}_nx=(E_nx,T_nx).                              \tag{14.32}
\]
Then
\[
 \|{\mathfrak C}_n-{\mathfrak C}\|
 \leq
 \bigl(\|P_n-P\|^2+\|S_n-S\|^2\bigr)^{1/2}.                 \tag{14.33}
\]

\begin{theorem}[Pseudoinverse and multiplier convergence]
\label{thm:v14v14-cofinal}
Assume the V13 uniform incidence certificate
\[
 \|K_n\|\leq M,\qquad
 \operatorname{rmin}
 \bigl(K_n^*|_{\operatorname{Ran}P_n}\bigr)\geq a>0.         \tag{14.34}
\]
Then, with \(d_0=d(a/M)>0\),
\[
 \sigma({\mathfrak C}_n^*{\mathfrak C}_n)
 \subset\{0\}\cup[d_0^2,2].                                \tag{14.35}
\]
Consequently
\[
 {\mathfrak C}_n^\dagger\longrightarrow
 {\mathfrak C}^\dagger,\qquad
 R_n=I-{\mathfrak C}_n^\dagger{\mathfrak C}_n
 \longrightarrow R                                         \tag{14.36}
\]
in operator norm.  If \(v_n\to v\), then the canonical multipliers
and remainders satisfy
\[
 ({\mathfrak C}_n^\dagger)^*v_n
 \longrightarrow({\mathfrak C}^\dagger)^*v,\qquad
 R_nv_n\longrightarrow Rv.                                 \tag{14.37}
\]
\end{theorem}

\begin{proof}
Corollary~\ref{cor:v14v14-incidence} gives the common reduced minimum
modulus \(d_0\).  Since
\(\|{\mathfrak C}_n\|^2
 \leq\|E_n\|^2+\|T_n\|^2=2\),
the spectral enclosure (14.35) follows.

Equation (14.33) gives norm convergence of \({\mathfrak C}_n\), hence
of
\(H_n={\mathfrak C}_n^*{\mathfrak C}_n\) to
\(H={\mathfrak C}^*{\mathfrak C}\).  On the common spectral set
\(\{0\}\cup[d_0^2,2]\), the function
\[
 f(0)=0,\qquad f(t)=t^{-1}\quad(t\geq d_0^2)
\]
is continuous.  Continuous functional calculus gives
\(H_n^\dagger=f(H_n)\to f(H)=H^\dagger\).  Since
\[
 {\mathfrak C}_n^\dagger
 =H_n^\dagger{\mathfrak C}_n^*,
\]
the first convergence in (14.36) follows.  The projector formula
(14.15) gives the second.  Equation (14.37) follows by applying the
convergent bounded operators to the convergent vectors.
\end{proof}

\begin{corollary}[Cofinal augmented-Gram landing]
\label{cor:v14v14-gramlanding}
If \(v_{a,n}\to v_a\) on every fixed screen under the hypotheses of
Theorem~\ref{thm:v14v14-cofinal}, then
\[
 \left\langle
 v_{a,n}-{\mathfrak C}_n^*
   ({\mathfrak C}_n^\dagger)^*v_{a,n},
 v_{b,n}-{\mathfrak C}_n^*
   ({\mathfrak C}_n^\dagger)^*v_{b,n}
 \right\rangle
 \longrightarrow
 \langle Rv_a,Rv_b\rangle.                                 \tag{14.38}
\]
If the limiting vectors also lie in \(\ker{\mathfrak C}\), zero
limiting diagonal Grams imply the full represented weak Einstein
equation on the declared V9 core.
\end{corollary}

\begin{proof}
Each augmented residual on the left is \(R_nv_{a,n}\) by (14.15).
Apply (14.36)--(14.37), continuity of the inner product, and
Proposition~\ref{prop:v14v14-gram}.
\end{proof}

\subsection{A noncommuting two-line check}

For the rotating-line example of V13, \(P\) projects onto \(e_1\) and
\(S_\theta\) onto \(u_\theta\), with \(0<\theta<\pi/2\).
Then \({\cal M}_\theta=0\), so \(R_\theta=0\) and
\({\mathfrak C}_\theta\) is injective.  In the standard basis,
\[
 {\mathfrak C}_\theta^*{\mathfrak C}_\theta
 =E+T_\theta
 =
 \begin{pmatrix}
 \sin^2\theta&-\sin\theta\cos\theta\\
 -\sin\theta\cos\theta&1+\cos^2\theta
 \end{pmatrix}.                                             \tag{14.39}
\]
Its determinant is \(\sin^2\theta\), so its smallest eigenvalue tends
to zero as \(\theta\downarrow0\).  At the limit the kernel becomes
\(\mathbb Re_1\).  This independently displays the pseudoinverse blowup
which accompanies the V13 intersection jump.

\subsection{Acquisition ledger and next bottleneck}

\[
\begin{array}{p{0.29\linewidth}|p{0.23\linewidth}|p{0.38\linewidth}}
\textbf{gate}&\textbf{status after V14}&\textbf{content}\\ \hline
\text{joint physical--tangent operator}
 &\text{acquired}
 &{\mathfrak C}x=((I-P)x,Tx)\\
\text{closed range}
 &\text{acquired from the V13 floor}
 &\operatorname{rmin}({\mathfrak C})\geq d(\beta)\\
\text{true intersection projector}
 &\text{acquired}
 &R=I-{\mathfrak C}^{\dagger}{\mathfrak C}\\
\text{best joint multiplier}
 &\text{acquired}
 &\inf_\lambda\|v-{\mathfrak C}^*\lambda\|=\|Rv\|\\
\text{operator factorization}
 &\text{acquired}
 &\inf_L\|D-L{\mathfrak C}\|=\|DR\|\\
\text{cofinal multiplier landing}
 &\text{acquired conditionally}
 &{\mathfrak C}_n^\dagger\to{\mathfrak C}^\dagger\\
\text{uniform regulator incidence floor}
 &\text{open}
 &same-history coercivity for \(K_n^*P_n\)\\
\text{full transverse Gram vanishing}
 &\text{open}
 &same-history residual estimate
\end{array}                                                  \tag{14.40}
\]

The next version is the mandatory fifth-version companion audit.
After that audit, the programme should go upstream to the source of
the open incidence floor.  The natural candidate is a gauge-fixed
elliptic or finite-stage normal operator for \(K_n^*P_n\), with its
kernel removed exactly.  Any coercivity claim must identify the
screen norm, kernel projection, boundary conditions, and constants
which remain uniform along the same regulator history.

\subsection{Parent-integrity record}

The installed parent was
\[
\begin{array}{c|c}
 \text{file}&
 \texttt{Renewal\_SM\_Einstein\_Continuum\_Research\_Notes\_V13.tex}\\
 \text{lines}&5026\\
 \text{bytes}&163054\\
 \text{SHA--256}&
 \texttt{C737E11A88CF9D0EB6F1D5B7E99D4078B6EDFE55C6E54019B52BB61F3ED53E4C}.
\end{array}                                                  \tag{14.41}
\]

% =============================================================================
% END APPENDED FOURTEENTH-CYCLE V14 MATERIAL
% =============================================================================
% =============================================================================
% BEGIN APPENDED FOURTEENTH-CYCLE V15 MATERIAL
% =============================================================================

\section{Fourteenth-cycle V15: compulsory companion audit and the
route to a populated incidence floor}
\label{sec:v14v15}

\subsection{Audit scope and immutable records}

This is the compulsory fifth-version audit requested for the shared
\texttt{latex\_notes} folder.  The newest numerical Yang--Mills and
Navier--Stokes notes were inspected read-only.  No companion file was
modified or removed.  The files selected by their active numerical
versions were
\[
\begin{array}{p{0.51\linewidth}|r|r}
\textbf{file}&\textbf{lines}&\textbf{bytes}\\ \hline
\texttt{Renewal\_Yang\_Mills\_Mass\_Gap\_Research\_Notes\_V2.tex}
 &872&32479\\
\texttt{Renewal\_NS\_Stability\_Research\_Notes\_V26.tex}
 &5319&278283
\end{array}                                                  \tag{15.1}
\]
with SHA--256 digests
\[
\begin{aligned}
 &\texttt{0FEF49FD192DF047585549D41C8EB904EDE64B51A29AE82AFB4824E21FC768B1},
                                                               \tag{15.2}\\
 &\texttt{82C887F8C2C69E4F28FAD7C3DC8DE5B9099CB5A4BF431EBA1FFF3E91935978AD},
                                                               \tag{15.3}
\end{aligned}
\]
respectively.  Line counts use the same UTF--8 split-line convention
as the SM version validator.

\subsection{Yang--Mills V2: exact coordinate elimination and exact
no-go witnesses}

The current Yang--Mills note proves that its reduced covariant pencil
has degree one in an active torus coordinate:
\[
 {\cal P}(z;\vartheta)
 =A(\vartheta)+zB(\vartheta)+\bar zB(\vartheta)^*,
 \qquad |z|=1.                                              \tag{15.4}
\]
It then gives the elementary sufficient certificate
\[
 \begin{pmatrix}
 A-X-\frac{\delta}{2}I&B\\
 B^*&X-\frac{\delta}{2}I
 \end{pmatrix}\succeq0
 \quad\Longrightarrow\quad
 {\cal P}(z;\vartheta)\succeq\delta I
 \quad\hbox{for every \(|z|=1\)}.                           \tag{15.5}
\]
The implication follows by testing the block matrix on
\((u,zu)\).  The note does not invoke an unproved matrix
Fej\'er--Riesz converse: a future use must exhibit \(X\), or exhibit
a two-term spectral factor.

Two attractive shortcuts are closed there by exact witnesses.  A
covariant factor-contraction bound fails, and every scalar splitting
\(X=tA\) fails at stated transverse corners, even though the direct
pencil has not been refuted.  The live target is a genuinely
matrix-valued \(X\), a spectral factor, or a return to the larger
constraint-multiplier system.  No Yang--Mills mass-gap conclusion is
claimed.

\subsection{Navier--Stokes V26: preserve the structured input before
taking norms}

The current Navier--Stokes note computes pointwise rank-one norms for
the first two derivatives of the Oseen kernel.  After contraction
with \(a\otimes a\), rotational symmetry reduces the squared norm to
a quadratic polynomial in
\[
 x=(n\cdot a)^2\in[0,1],                                    \tag{15.6}
\]
so the exact maximum is obtained at the endpoints or the vertex.
For the first derivative the rank-one and ambient symmetric-traceless
norms coincide.  For the second derivative the rank-one norm is
strictly smaller on a stated intermediate radial interval, with a
reported reduction reaching about eleven percent near one sample
radius.  The companion note records that this improves constants but
does not prove regularity.

The transferable lesson is structural: the special input incidence
must be retained through the contraction.  Replacing it prematurely
by the norm over the full ambient tensor space can discard a real
coercivity margin.

\subsection{Audit decision for the SM--Einstein bottleneck}

At the end of V14 the open local row was a uniform reduced singular
floor for
\[
 J_n:=K_n^*|_{\operatorname{Ran}P_n}.                       \tag{15.7}
\]
The two companion readings select the following route.

First, preserve the physical input projector inside the normal
operator:
\[
 J_n^*J_n
 =
 P_nK_nK_n^*P_n
 \quad\hbox{on }\operatorname{Ran}P_n.                      \tag{15.8}
\]
This is the SM analogue of retaining the rank-one input in the
Navier--Stokes contraction.  An estimate on \(K_nK_n^*\) over the
whole ambient space may be strictly weaker and can be zero for
irrelevant gauge reasons.

Second, remove the exact kernel before asking for strict positivity.
If
\[
 \Pi_n:=\operatorname{proj}_{\ker J_n},                     \tag{15.9}
\]
then the kernel-completed normal operator
\[
 \widehat H_n:=J_n^*J_n+\Pi_n                               \tag{15.10}
\]
is strictly positive precisely when \(J_n\) has a positive reduced
singular floor.  This converts a quotient coercivity question into a
full-space positive pencil without hiding the kernel.

Third, when the regulator symbol of \(\widehat H_n\) has degree one
in a phase, use the exact block implication (15.5).  The auxiliary
splitting must be matrix-valued unless the concrete symbol proves
otherwise; the Yang--Mills counterexamples forbid treating a scalar
ansatz as a general principle.

\begin{theorem}[Audit bridge]
\label{thm:v14v15-bridge}
Let \(J:H_1\to H_2\) be bounded and let
\(\Pi=\operatorname{proj}_{\ker J}\).  For every \(a>0\),
\[
 \operatorname{rmin}(J)\geq a
 \quad\Longleftrightarrow\quad
 J^*J+\Pi\succeq\min\{a^2,1\}I,                             \tag{15.11}
\]
provided the left side is interpreted on
\((\ker J)^\perp\).  Conversely, if
\[
 J^*J+\Pi\succeq\delta I                                   \tag{15.12}
\]
for \(\delta>0\), then
\[
 \operatorname{rmin}(J)\geq\sqrt{\delta}.                  \tag{15.13}
\]
\end{theorem}

\begin{proof}
The orthogonal decomposition
\(H_1=\ker J\oplus(\ker J)^\perp\) reduces \(J^*J\).
On the first summand \(J^*J+\Pi=I\).  On the second it equals
\(J^*J\), whose lower spectral bound is
\(\operatorname{rmin}(J)^2\).  This proves (15.11).
Restricting (15.12) to the second summand gives (15.13).
\end{proof}

The reverse implication in (15.13) uses the same \(\delta\) only when
\(\delta\leq1\), which is the range relevant to a normalized
certificate.  Formula (15.11) is the exact statement in all scales.

\begin{decision}[Direction after the V15 audit]
\label{dec:v14v15-direction}
V16 will develop a kernel-completed symbol compiler.  It will prove
that a displayed matrix-valued block splitting for one coordinate
gives a uniform reduced singular floor for \(K_n^*P_n\), and will
propagate the constant through V13--V14 to the intersection projector
and joint multiplier.  It will also state the exact failure witness
obtained from a negative block quadratic form.

No companion problem-specific estimate is imported.  In particular,
the Yang--Mills pencil and the Navier--Stokes kernel constants do not
prove an Einstein incidence bound.  They determine the proof
architecture for the SM regulator symbol.
\end{decision}

\subsection{Acquisition ledger after the audit}

\[
\begin{array}{p{0.29\linewidth}|p{0.23\linewidth}|p{0.38\linewidth}}
\textbf{gate}&\textbf{status after V15}&\textbf{content}\\ \hline
\text{companion checkpoint}
 &\text{acquired}
 &YM V2 and NS V26 read with hashes\\
\text{kernel completion}
 &\text{acquired abstractly}
 &J^*J+\Pi_{\ker J}\\
\text{structured physical incidence}
 &\text{selected}
 &retain \(P_nK_nK_n^*P_n\)\\
\text{one-coordinate certificate}
 &\text{selected conditionally}
 &matrix block splitting \(X\)\\
\text{concrete SM symbol}
 &\text{open}
 &derive from the same regulator history\\
\text{uniform completed-symbol floor}
 &\text{open}
 &exact formula or certified enclosure\\
\text{full transverse Gram vanishing}
 &\text{open}
 &same-history residual estimate
\end{array}                                                  \tag{15.14}
\]

The next compulsory companion audit is V20.

\subsection{Parent-integrity record}

The installed parent was
\[
\begin{array}{c|c}
 \text{file}&
 \texttt{Renewal\_SM\_Einstein\_Continuum\_Research\_Notes\_V14.tex}\\
 \text{lines}&5601\\
 \text{bytes}&181192\\
 \text{SHA--256}&
 \texttt{910A79EF4E215EB05472130D01474B5A2BD880B5FB03E8329D034F87CB215940}.
\end{array}                                                  \tag{15.15}
\]

% =============================================================================
% END APPENDED FOURTEENTH-CYCLE V15 MATERIAL
% =============================================================================
% =============================================================================
% BEGIN APPENDED FOURTEENTH-CYCLE V16 MATERIAL
% =============================================================================

\section{Fourteenth-cycle V16: a kernel-completed symbol compiler for
the Einstein incidence floor}
\label{sec:v14v16}

\subsection{Purpose}

V13 reduced stability of the physical--transverse intersection to a
uniform lower bound for the physical adjoint incidence
\[
 J=K^*|_{\operatorname{Ran}P}.
\]
V15 selected a kernel-completed normal symbol and a matrix-valued
one-coordinate certificate.  The present note makes that route
rigorous.

A useful noncircular feature is proved first.  One may begin with an
explicit projector onto the \emph{declared} gauge-zero modes without
assuming that all zero modes have been found.  If the completed normal
operator is strictly positive, that positivity proves that the
declared projector exhausts the kernel.  If a hidden zero mode exists,
it becomes an exact null witness and the certificate must fail.

\subsection{Noncircular completion of a declared kernel}

Let \(H_1,H_2\) be Hilbert spaces, let
\(J\in{\cal B}(H_1,H_2)\), and let
\(\Pi\in{\cal B}(H_1)\) satisfy
\[
 \Pi^2=\Pi=\Pi^*,\qquad J\Pi=0.                             \tag{16.1}
\]
Thus \(\operatorname{Ran}\Pi\) is a declared subspace of
\(\ker J\), but equality is not assumed.  Define
\[
 \widehat H:=J^*J+\Pi.                                      \tag{16.2}
\]

\begin{theorem}[Completed positivity proves kernel completeness]
\label{thm:v14v16-kernel}
If for some \(\delta>0\)
\[
 \widehat H\succeq\delta I,                                \tag{16.3}
\]
then
\[
 \boxed{\ker J=\operatorname{Ran}\Pi}                      \tag{16.4}
\]
and
\[
 \boxed{\operatorname{rmin}(J)\geq\sqrt{\delta}.}           \tag{16.5}
\]
If \(\operatorname{Ran}\Pi\ne\{0\}\), condition (16.3) also forces
\(\delta\leq1\).
\end{theorem}

\begin{proof}
The inclusion
\(\operatorname{Ran}\Pi\subset\ker J\) is (16.1).  Take
\(x\in\ker J\), and put \(y=(I-\Pi)x\).  Since \(J\Pi=0\),
\[
 Jy=Jx-J\Pi x=0,\qquad \Pi y=0.
\]
Therefore \(\langle\widehat Hy,y\rangle=0\).  Condition (16.3)
forces \(y=0\), so \(x=\Pi x\), proving (16.4).

If \(x\perp\ker J=\operatorname{Ran}\Pi\), then \(\Pi x=0\), and
(16.3) gives
\[
 \|Jx\|^2
 =\langle J^*Jx,x\rangle
 =\langle\widehat Hx,x\rangle
 \geq\delta\|x\|^2.
\]
Taking the infimum over the unit sphere in the kernel complement
proves (16.5).  Finally, for a unit
\(x\in\operatorname{Ran}\Pi\), one has
\(\widehat Hx=x\), hence \(\delta\leq1\).
\end{proof}

\begin{proposition}[Exact hidden-kernel obstruction]
\label{prop:v14v16-hidden}
If there exists \(0\ne y\in H_1\) such that
\[
 Jy=0,\qquad \Pi y=0,                                      \tag{16.6}
\]
then \(\widehat Hy=0\).  No positive \(\delta\) can satisfy (16.3).
\end{proposition}

\begin{proof}
Both summands of (16.2) annihilate \(y\).
\end{proof}

Thus kernel discovery and coercivity certification can be performed
together.  Defining \(\Pi\) through a pseudoinverse whose uniform
boundedness is the desired conclusion would be circular; verifying
(16.1) from exact gauge identities is not.

\subsection{One-coordinate Laurent compiler}

Now let \(H_1=\mathbb C^N\), let
\(\vartheta\) range over a compact parameter space \(\Theta\), and
let \(z\in\mathbb T\).  Suppose \(J(z,\vartheta)\) and an explicit
orthogonal projection \(\Pi(z,\vartheta)\) obey (16.1) pointwise.
Assume their completed normal symbol has active degree one:
\[
 \widehat H(z,\vartheta)
 =A(\vartheta)+zB(\vartheta)
  +\bar zB(\vartheta)^*,\qquad A(\vartheta)=A(\vartheta)^*.
                                                               \tag{16.7}
\]
No claim that a particular SM regulator has yet been shown to satisfy
(16.7) is made here; it is an acquisition row.

For a Hermitian matrix field \(X(\vartheta)\), define
\[
 {\cal Q}_{\delta,X}(\vartheta)
 :=
 \begin{pmatrix}
 A(\vartheta)-X(\vartheta)-\frac{\delta}{2}I&
 B(\vartheta)\\
 B(\vartheta)^*&
 X(\vartheta)-\frac{\delta}{2}I
 \end{pmatrix}.                                             \tag{16.8}
\]

\begin{theorem}[Kernel-completed one-coordinate certificate]
\label{thm:v14v16-block}
If
\[
 {\cal Q}_{\delta,X}(\vartheta)\succeq0
 \quad\hbox{for every }\vartheta\in\Theta                  \tag{16.9}
\]
for some \(\delta>0\), then, for every
\((z,\vartheta)\in\mathbb T\times\Theta\),
\[
 \widehat H(z,\vartheta)\succeq\delta I,                   \tag{16.10}
\]
\[
 \ker J(z,\vartheta)
 =\operatorname{Ran}\Pi(z,\vartheta),                       \tag{16.11}
\]
and
\[
 \operatorname{rmin}J(z,\vartheta)\geq\sqrt{\delta}.        \tag{16.12}
\]
\end{theorem}

\begin{proof}
For \(u\in\mathbb C^N\), set \(w_z=(u,zu)\).  Since
\(|z|=1\), direct multiplication gives
\[
 \begin{aligned}
 w_z^*{\cal Q}_{\delta,X}(\vartheta)w_z
 &=
 u^*\bigl(
 A(\vartheta)+zB(\vartheta)
 +\bar zB(\vartheta)^*-\delta I
 \bigr)u\\
 &=u^*(\widehat H(z,\vartheta)-\delta I)u.                 \tag{16.13}
 \end{aligned}
\]
Condition (16.9) proves (16.10).  Apply
Theorem~\ref{thm:v14v16-kernel} pointwise to obtain
(16.11)--(16.12).
\end{proof}

\begin{corollary}[Exact spectral-factor certificate]
\label{cor:v14v16-factor}
If an explicit matrix field \(G_0(\vartheta),G_1(\vartheta)\)
satisfies
\[
 \widehat H(z,\vartheta)-\delta I
 =
 \bigl(G_0(\vartheta)+zG_1(\vartheta)\bigr)^*
 \bigl(G_0(\vartheta)+zG_1(\vartheta)\bigr),                \tag{16.14}
\]
then Theorem~\ref{thm:v14v16-block} applies with
\[
 X(\vartheta)
 =G_1(\vartheta)^*G_1(\vartheta)+\frac{\delta}{2}I.          \tag{16.15}
\]
\end{corollary}

\begin{proof}
Expanding (16.14) gives
\[
 A-\delta I=G_0^*G_0+G_1^*G_1,\qquad B=G_0^*G_1.
\]
Substitution of (16.15) into (16.8) yields
\[
 {\cal Q}_{\delta,X}
 =
 \begin{pmatrix}G_0^*\\G_1^*\end{pmatrix}
 \begin{pmatrix}G_0&G_1\end{pmatrix}\succeq0.
\]
Use Theorem~\ref{thm:v14v16-block}.
\end{proof}

Only the forward implication is used.  No matrix
Fej\'er--Riesz converse is assumed.

\subsection{Exact and interval-verifiable positivity}

Condition (16.9) still contains the transverse parameter
\(\vartheta\).  The following two elementary mechanisms make clear
what counts as a proof.

\begin{proposition}[Exact algebraic certificate]
\label{prop:v14v16-exact}
If, for every \(\vartheta\), an explicitly verified factorization
\[
 {\cal Q}_{\delta,X}(\vartheta)
 =Y(\vartheta)^*Y(\vartheta)                               \tag{16.16}
\]
holds, then (16.9) holds.  The same conclusion follows from an exact
Hermitian \(LDL^*\) decomposition whose diagonal blocks are proved
positive semidefinite.
\end{proposition}

\begin{proof}
Both statements are direct quadratic-form identities.
\end{proof}

\begin{proposition}[Finite-cover enclosure]
\label{prop:v14v16-cover}
Let \(\Theta=\bigcup_{r=1}^R\Theta_r\).  Choose
\(\vartheta_r\in\Theta_r\).  Suppose exact or interval arithmetic
proves
\[
 {\cal Q}_{\delta,X}(\vartheta_r)\succeq\rho_r I,\qquad
 \sup_{\vartheta\in\Theta_r}
 \|{\cal Q}_{\delta,X}(\vartheta)
       -{\cal Q}_{\delta,X}(\vartheta_r)\|
 \leq\varepsilon_r                                        \tag{16.17}
\]
with
\[
 \rho_r-\varepsilon_r\geq0                                \tag{16.18}
\]
for every \(r\).  Then (16.9) holds.
If every difference in (16.18) is strictly positive, the block
certificate itself has that positive floor.
\end{proposition}

\begin{proof}
For a unit vector \(w\) and
\(\vartheta\in\Theta_r\),
\[
 \begin{aligned}
 w^*{\cal Q}_{\delta,X}(\vartheta)w
 &\geq
 w^*{\cal Q}_{\delta,X}(\vartheta_r)w
 -
 \|{\cal Q}_{\delta,X}(\vartheta)
       -{\cal Q}_{\delta,X}(\vartheta_r)\|\\
 &\geq\rho_r-\varepsilon_r.
 \end{aligned}
\]
Use (16.18).
\end{proof}

A floating mesh without the enclosure term
\(\varepsilon_r\) does not prove (16.9).  Conversely, the cover need
not be uniform: it may be refined only where the anchor floor is
small.

\begin{corollary}[Derivative-based box enclosure]
\label{cor:v14v16-derivative}
Suppose \(\Theta_r\) is a coordinate box centered at
\(\vartheta_r\), with side radii \(h_{r,j}\), and
\[
 \sup_{\vartheta\in\Theta_r}
 \|\partial_{\vartheta_j}{\cal Q}_{\delta,X}(\vartheta)\|
 \leq L_{r,j}.                                             \tag{16.19}
\]
Then (16.17) holds with
\[
 \varepsilon_r=\sum_jL_{r,j}h_{r,j}.                       \tag{16.20}
\]
\end{corollary}

\begin{proof}
Integrate the derivative along the coordinatewise polygonal path
from \(\vartheta_r\) to \(\vartheta\) and apply the triangle
inequality.
\end{proof}

\subsection{Witness hierarchy}

A failed auxiliary splitting and a failed incidence floor are
different statements.

\begin{proposition}[Three exact failure witnesses]
\label{prop:v14v16-witness}
Fix \(\delta>0\).

\begin{enumerate}
\item If \(w^*{\cal Q}_{\delta,X}(\vartheta)w<0\) for some
      \(w\in\mathbb C^{2N}\), the proposed splitting \(X\) fails.
      This alone need not refute (16.10).
\item If the negative witness has the structured form
      \(w=(u,zu)\), then
      \[
       u^*(\widehat H(z,\vartheta)-\delta I)u<0,             \tag{16.21}
      \]
      so the claimed floor \(\delta\) is false at that phase.
\item If \(u\ne0\), \(J(z,\vartheta)u=0\), and
      \(\Pi(z,\vartheta)u=0\), then \(u\) is a hidden-kernel
      witness and no positive completed-symbol floor exists.
\end{enumerate}
\end{proposition}

\begin{proof}
The first statement is the definition of failure of positive
semidefiniteness.  The identity (16.13) proves the second.
Proposition~\ref{prop:v14v16-hidden} proves the third.
\end{proof}

This hierarchy prevents an upstream no-go result for a convenient
certificate from being misreported as a no-go result for the direct
Einstein incidence symbol.

\subsection{Propagation to the V13--V14 continuum interface}

Return to the geometric notation at a fixed regulator stage.  Let
\[
 J=K^*|_{\operatorname{Ran}P},\qquad
 \operatorname{Ran}K=\operatorname{Ran}T,\qquad
 \|K\|\leq M.                                               \tag{16.22}
\]
Suppose Theorem~\ref{thm:v14v16-block} supplies the common symbol
floor \(\delta>0\).

\begin{theorem}[Constant propagation]
\label{thm:v14v16-propagation}
Under (16.22),
\[
 \alpha(K,P):=\operatorname{rmin}J\geq\sqrt{\delta},         \tag{16.23}
\]
\[
 \beta(P,S)\geq\frac{\sqrt{\delta}}{M},\qquad
 \gamma(P,S)\geq\frac{\delta}{M^2},                         \tag{16.24}
\]
and the joint constraint of V14 obeys
\[
 \operatorname{rmin}({\mathfrak C})
 \geq
 d\!\left(\frac{\sqrt{\delta}}{M}\right),                  \tag{16.25}
\]
\[
 \|{\mathfrak C}^{\dagger}\|
 \leq
 d\!\left(\frac{\sqrt{\delta}}{M}\right)^{-1},\qquad
 d(t)=\frac{t}{\sqrt{1+(1+t)^2}}.                          \tag{16.26}
\]
\end{theorem}

\begin{proof}
Equation (16.23) is (16.12).  The V13 incidence theorem gives
\(\beta\geq\alpha/\|K\|\), and its exact angle identity gives
\(\gamma=\beta^2\).  This proves (16.24).  The V14 coercivity theorem
then gives (16.25)--(16.26).
\end{proof}

\begin{corollary}[Uniform finite-stage compiler]
\label{cor:v14v16-cofinal}
Suppose every member of a same-history regulator family admits the
data (16.1), (16.7)--(16.9), and (16.22) with common constants
\[
 \delta_n\geq\delta_0>0,\qquad \|K_n\|\leq M.                \tag{16.27}
\]
If \(P_n\to P\), \(S_n\to S\), and the residual vectors converge on
each fixed screen, then:
\[
 R_n\to R,\qquad
 {\mathfrak C}_n^\dagger\to{\mathfrak C}^\dagger,           \tag{16.28}
\]
and the optimally augmented residual Grams converge to the V14
intersection Grams.  If the limiting residuals also satisfy the
physical and tangent equations, zero limiting diagonal Grams imply
the represented weak Einstein equation on the V9 core.
\end{corollary}

\begin{proof}
Equations (16.24)--(16.26) give common V13 angle and V14 joint
coercivity floors.  Apply the cofinal theorems of V13 and V14.
\end{proof}

All objects in (16.27) must come from the same regulator history.
A floor from one discretization cannot be combined with projectors or
residuals from another.

\subsection{Concrete acquisition protocol}

For each finite Einstein regulator, a complete certificate package
must contain
\[
\begin{array}{c|p{0.72\linewidth}}
\mathrm{S1}&
the exact physical projector \(P_n\), tangent incidence \(K_n\), and
the matrix of \(J_n=K_n^*|_{\operatorname{Ran}P_n}\);\\
\mathrm{S2}&
an explicit orthogonal projector \(\Pi_n\) satisfying
\(J_n\Pi_n=0\), derived from exact gauge identities;\\
\mathrm{S3}&
the Laurent coefficients \(A_n,B_n\) in one active phase and a proof
that no omitted Fourier modes occur;\\
\mathrm{S4}&
a genuinely matrix-valued \(X_n\), or an explicit spectral factor,
with an exact or interval-certified block inequality;\\
\mathrm{S5}&
a common \(\delta_0>0\), a common upper bound for \(\|K_n\|\), and
the resulting constants in (16.24)--(16.26);\\
\mathrm{S6}&
the same-history convergence of projectors, residuals, and localized
Grams required by the V13--V14 landing theorem.
\end{array}                                                  \tag{16.29}
\]

The V15 companion audit suggests the order S1--S4.  It does not
populate any one of those rows for the Einstein regulator.

\subsection{Claim boundary and next upstream task}

\begin{assessment}[What V16 acquires]
The kernel-completion step is now noncircular: completed positivity
both proves kernel completeness and supplies the singular floor.
A degree-one active phase can be removed by an exhibited
matrix-valued block splitting, with exact and interval-verifiable
versions.  The resulting constant propagates explicitly through the
physical--transverse projector and the joint multiplier construction.
\end{assessment}

\begin{firewall}[What remains open]
No concrete SM regulator symbol \(J_n\), declared kernel projector
\(\Pi_n\), Laurent coefficient bank, block splitting \(X_n\), or
positive uniform \(\delta_0\) has been populated here.  Consequently
no transverse Gram vanishing, regulator removal, Einstein continuum,
or full Standard Model--Einstein continuum follows from V16.
\end{firewall}

The next version must inspect the existing Einstein/Grand-Tensor
source for the most concrete available finite-stage physical
projector and infinitesimal diffeomorphism incidence.  If those data
do not determine the coefficient bank in S1--S3, the correct upstream
result is an interface theorem listing the missing typed maps and
showing which later claims are blocked by each omission.  Inventing a
symbol would be circular.

\subsection{Acquisition ledger}

\[
\begin{array}{p{0.29\linewidth}|p{0.23\linewidth}|p{0.38\linewidth}}
\textbf{gate}&\textbf{status after V16}&\textbf{content}\\ \hline
\text{kernel-completion logic}
 &\text{acquired}
 &positivity proves \(\ker J=\operatorname{Ran}\Pi\)\\
\text{one-coordinate compiler}
 &\text{acquired conditionally}
 &matrix block certificate\\
\text{certified transverse cover}
 &\text{acquired}
 &anchor floor minus enclosure radius\\
\text{failure classification}
 &\text{acquired}
 &splitting, direct-floor, and hidden-kernel witnesses\\
\text{constant propagation}
 &\text{acquired}
 &\delta\mapsto\gamma,\operatorname{rmin}{\mathfrak C}\\
\text{concrete Einstein coefficient bank}
 &\text{open}
 &source inspection required\\
\text{uniform block certificate}
 &\text{open}
 &same-history exact or interval proof\\
\text{full transverse Gram vanishing}
 &\text{open}
 &same-history residual estimate
\end{array}                                                  \tag{16.30}
\]

\subsection{Parent-integrity record}

The installed parent was
\[
\begin{array}{c|c}
 \text{file}&
 \texttt{Renewal\_SM\_Einstein\_Continuum\_Research\_Notes\_V15.tex}\\
 \text{lines}&5830\\
 \text{bytes}&189674\\
 \text{SHA--256}&
 \texttt{89D39DDA10E6A1B61BF5CAF21F3E3A3DA45FC2662DC9984B5BF8F892DB94E787}.
\end{array}                                                  \tag{16.31}
\]

% =============================================================================
% END APPENDED FOURTEENTH-CYCLE V16 MATERIAL
% =============================================================================
% =============================================================================
% BEGIN APPENDED FOURTEENTH-CYCLE V17 MATERIAL
% =============================================================================

\section{Fourteenth-cycle V17: typed source interface, gauge tangency,
and the nested-projection branch}
\label{sec:v14v17}

\subsection{Source provenance and purpose}

V16 required inspection of the attached sources before any concrete
Einstein incidence symbol was invented.  The two files inspected
read-only were
\[
\begin{array}{p{0.55\linewidth}|r|r}
\textbf{file}&\textbf{lines}&\textbf{bytes}\\ \hline
\texttt{Renewal\_Grand\_Tensor\_Monograph\_V067.tex}
 &68938&2804854\\
\texttt{Renewal\_Einstein\_Standard\_Model\_Physical\_Source\_Metric
\_Duhamel\_Ancestry\_and\_Quantum\_Continuum\_Programme\_V30.tex}
 &32301&1262823
\end{array}                                                  \tag{17.1}
\]
with SHA--256 digests
\[
\begin{aligned}
 &\texttt{FACD058EC8BB210AD8B296C9D7F4D78A229D08D7A2AE378808700655CB8BC0FE},
                                                               \tag{17.2}\\
 &\texttt{E0AE0E0990DE472021A1E90A9C10D41C3F747F28D97A5B1D8F1281852E3A0E69}.
                                                               \tag{17.3}
\end{aligned}
\]

The source inspection finds a more upstream branch than the V16
symbol search.  If infinitesimal diffeomorphisms preserve the
physical constraint surface, the tangent range is nested inside the
physical range.  The projections then commute, the angle gap is
exactly one, and no separate \(K^*P\) angle estimate is needed.  The
first missing row is therefore the typed constraint-propagation
identity which proves that nesting.

\subsection{What the sources already contain}

The Grand Tensor source contains the following finite and conditional
objects.

\begin{enumerate}
\item One finite common action
\[
 {\mathscr S}_{\rm tot}(q,\Phi)
 ={\mathscr S}_g(q)+{\mathscr S}_{\rm mat}(q,\Phi)
\]
and its geometric covector
\[
 {\cal E}^{\rm tot}_q
 =2D_q{\mathscr S}_{\rm tot}
 =2D_q{\mathscr S}_g-\mathsf T.                             \tag{17.4}
\]
\item A tangential relabeling map \({\cal L}_q\), a matter
relabeling map \({\cal L}_\Phi\), and the finite Noether relation
\[
 {\cal L}_q^*\mathsf T
 =2{\cal L}_\Phi^*E_\Phi.                                   \tag{17.5}
\]
On matter shell, common-action relabeling invariance supplies the
tangential Ward equation for the total geometric residual.
\item A finite HDA with self-adjoint tangential generators \(D_h[v]\)
on a physical carrier.  These are operators on the history carrier;
they are not, without another typed map, the geometric variation map
\({\cal L}_q\).
\item A positive coherent Ward Gram built from direct and relabelled
source syntheses.  Its \(o(\varepsilon_X^2)\) hypothesis kills the
infinitesimal Ward anomaly after the difference quotient has been
identified.
\item A determining metric--matter variation panel \(C_X\) with the
conditional floor
\[
 C_X^*C_X\succeq c_*G_{E,X}.                                \tag{17.6}
\]
Together with vanishing panel residual, (17.6) directly forces the
Riesz vector of the total first variation to zero.
\item A positive-screen compactness theorem and a conditional local
weak Einstein--matter handoff.
\end{enumerate}

The same source explicitly warns that ADM kinematics does not supply
a dynamical Einstein inverse: Hamiltonian and momentum constraints,
gauge-fixed evolution, and a constraint-propagation estimate are
additional inputs.  Neither inspected source gives one simultaneous
matrix tuple consisting of a residual Hilbert metric, a physical
constraint map on that residual screen, the geometric relabeling
matrix, its exact kernel projector, and a Laurent coefficient bank
for \(K^*P\).

\subsection{Metric whitening of the source residual}

The action residual in (17.4) is a covector, whereas V12--V16 use
Hilbert vectors.  The following conversion fixes the type.

Let \(Q\) be a finite-dimensional real or complex variation space
with a background inner product, and let
\[
 G=G^*>0                                                     \tag{17.7}
\]
be the faithful metric used on geometric variations.  Represent an
action covector by \(\epsilon\in Q\), so that its value on
\(k\in Q\) is \(\langle\epsilon,k\rangle\).  Its \(G\)-Riesz vector
and whitened residual are
\[
 e:=G^{-1}\epsilon,\qquad
 r:=G^{1/2}e=G^{-1/2}\epsilon.                              \tag{17.8}
\]
Let \(L:X\to Q\) be the infinitesimal geometric relabeling map and
define the whitened tangent incidence
\[
 K:=G^{1/2}L:X\longrightarrow Q.                            \tag{17.9}
\]

\begin{theorem}[Ward identity becomes transverse membership]
\label{thm:v14v17-whitening}
With the definitions (17.7)--(17.9),
\[
 L^*\epsilon=0
 \quad\Longleftrightarrow\quad
 K^*r=0.                                                    \tag{17.10}
\]
If \(K\) has closed range and
\[
 T:=KK^\dagger,\qquad S:=I-T,                               \tag{17.11}
\]
then (17.10) is equivalent to
\[
 Tr=0\qquad\hbox{and}\qquad Sr=r.                           \tag{17.12}
\]
\end{theorem}

\begin{proof}
Direct substitution gives
\[
 K^*r
 =L^*G^{1/2}G^{-1/2}\epsilon
 =L^*\epsilon,
\]
which proves (17.10).  The operator \(T=KK^\dagger\) is the
orthogonal projection onto \(\operatorname{Ran}K\), while
\(\ker K^*=(\operatorname{Ran}K)^\perp=\operatorname{Ran}S\).
This proves (17.12).
\end{proof}

Thus the common-action Ward identity supplies the transverse equation
only after the action covector, variation metric, and relabeling map
are taken from the same finite source card.

\subsection{Transport of a physical constraint}

Suppose a constraint on the \(G\)-Riesz residual is represented by a
closed-range map
\[
 C:Q\longrightarrow H_C.                                   \tag{17.13}
\]
In whitened coordinates put
\[
 \widetilde C:=CG^{-1/2},\qquad
 P:=I-\widetilde C^\dagger\widetilde C.                     \tag{17.14}
\]
Then \(P\) is the orthogonal projection onto the whitened physical
kernel
\[
 \ker\widetilde C
 =G^{1/2}\ker C.                                            \tag{17.15}
\]

\begin{proof}
The Moore--Penrose formula gives the projection statement.  Since
\(e=G^{-1/2}r\), one has
\(\widetilde Cr=Ce\), which gives (17.15).
\end{proof}

There is a necessary typing warning.  The V11 symbol
\(C_{\rm phys}:{\cal H}_{\rm kin}\to{\cal K}_{\rm con}\) and the V12
projection on the represented residual screen can be identified only
after an explicit isometry or closed embedding transports
\({\cal H}_{\rm kin}\) into that screen and intertwines the
constraints.  An ADM constraint on initial data is likewise not
automatically the map (17.13) on action-residual Riesz vectors.

\subsection{Exact gauge tangency removes the angle bottleneck}

The finite or linearized constraint-propagation identity is
\[
 \boxed{\widetilde C K=0.}                                  \tag{17.16}
\]
It says that every infinitesimal diffeomorphism direction lies in the
physical constraint kernel.

\begin{theorem}[Nested physical and tangent projections]
\label{thm:v14v17-nested}
Let \(\widetilde C\) and \(K\) have closed range, define \(P,T,S\) by
(17.11) and (17.14), and assume (17.16).  Then
\[
 \operatorname{Ran}T\subseteq\operatorname{Ran}P,           \tag{17.17}
\]
\[
 PT=TP=T,\qquad PS=SP=P-T,                                  \tag{17.18}
\]
and the physical--transverse projector of V12 is
\[
 \boxed{R=P-T.}                                             \tag{17.19}
\]
If \(\operatorname{Ran}T\ne\{0\}\), the V13 angle constants are
\[
 \boxed{\beta(P,S)=1,\qquad\gamma(P,S)=1.}                  \tag{17.20}
\]
\end{theorem}

\begin{proof}
Equation (17.16) gives
\(\operatorname{Ran}K\subseteq\ker\widetilde C\), which is
(17.17).  For orthogonal projections, range inclusion is equivalent
to \(PT=TP=T\).  Since \(S=I-T\), this proves (17.18).

The operator \(P-T\) is self-adjoint and
\[
 (P-T)^2=P-PT-TP+T=P-T,
\]
so it is an orthogonal projection.  Its range is
\[
 \operatorname{Ran}P\cap\ker T
 =\operatorname{Ran}P\cap\operatorname{Ran}S,
\]
which proves (17.19).

The orthogonal decomposition
\[
 \operatorname{Ran}P
 =
 \operatorname{Ran}T\oplus\operatorname{Ran}R               \tag{17.21}
\]
shows that the complement of the intersection inside
\(\operatorname{Ran}P\) is \(\operatorname{Ran}T\).  On that
subspace \(T\) is the identity.  Hence the reduced minimum modulus
of \(T|_{\operatorname{Ran}P}\) is one.  Apply the V13 exact
gap identity to obtain (17.20).
\end{proof}

\begin{corollary}[The V14 joint constraint becomes a partial
isometry]
\label{cor:v14v17-partial}
Under Theorem~\ref{thm:v14v17-nested}, let
\(E=I-P\) and
\[
 {\mathfrak C}x=(Ex,Tx).
\]
Then
\[
 ET=TE=0,\qquad
 {\mathfrak C}^*{\mathfrak C}=E+T=I-R,                      \tag{17.22}
\]
\[
 \boxed{{\mathfrak C}^{\dagger}={\mathfrak C}^*,\qquad
        \operatorname{rmin}({\mathfrak C})=1.}              \tag{17.23}
\]
In particular, all V14 multiplier estimates have constant one.
\end{corollary}

\begin{proof}
The ranges of \(E\) and \(T\) are orthogonal by (17.17), proving
\(ET=TE=0\).  Lemma V14 gives
\({\mathfrak C}^*{\mathfrak C}=E+T\), now an orthogonal projection.
Thus \({\mathfrak C}\) is a partial isometry with initial projection
\(I-R\); the Moore--Penrose inverse of a partial isometry is its
adjoint.  Its nonzero singular values are all one.
\end{proof}

This exact branch is strictly upstream of the V16 completed-symbol
search.  The latter remains useful when (17.16) is unavailable,
fails, or is known only approximately.

\subsection{Correct variational meaning of the intersection}

Let
\[
 {\cal A}:=\ker\widetilde C=\operatorname{Ran}P,\qquad
 {\cal G}:=\operatorname{Ran}K=\operatorname{Ran}T.         \tag{17.24}
\]
Under gauge tangency, \({\cal G}\subseteq{\cal A}\).  The orthogonal
representative of the gauge-reduced physical variation space is
\[
 {\cal A}\ominus{\cal G}
 ={\cal A}\cap{\cal G}^{\perp}
 =\operatorname{Ran}R.                                     \tag{17.25}
\]

\begin{theorem}[Constrained gauge-reduced Euler equation]
\label{thm:v14v17-quotient}
For a whitened action residual \(r\), the following are equivalent:
\[
 \langle r,h\rangle=0
 \quad\hbox{for every }h\in{\cal A}\ominus{\cal G},          \tag{17.26}
\]
\[
 \boxed{Rr=0,}                                              \tag{17.27}
\]
\[
 \boxed{r=E\lambda_{\rm con}+T\lambda_{\rm gau}}
 \quad\hbox{for some }\lambda_{\rm con},\lambda_{\rm gau}
 \in Q.                                                     \tag{17.28}
\]
The canonical orthogonal decomposition is
\[
 \boxed{r=Er+Tr+Rr.}                                        \tag{17.29}
\]
If the Ward equation \(K^*r=0\) also holds, then \(Tr=0\), and
gauge-reduced constrained stationarity is equivalent to
\[
 r=Er,                                                      \tag{17.30}
\]
a pure constraint multiplier.  It implies \(r=0\) only if that
multiplier is separately shown to vanish, or if \(r\in{\cal A}\).
\end{theorem}

\begin{proof}
The Riesz representation theorem and (17.25) make (17.26)
equivalent to (17.27).  By (17.22),
\[
 I-R=E+T
\]
is the orthogonal projection onto
\({\cal A}^{\perp}\oplus{\cal G}\).  This gives
(17.28)--(17.29).  The Ward equation is equivalent to \(Tr=0\) by
Theorem~\ref{thm:v14v17-whitening}.  The final statements follow
from the orthogonal sum.
\end{proof}

This theorem fixes the interpretation of the V14 projected Gram.
For a raw constrained action residual, \(Rr=0\) is the physical
Euler equation modulo constraint and gauge multipliers.  The stronger
V12 implication from \(Rr=0\) to \(r=0\) applies only after the
residual itself has independently been placed in
\(\operatorname{Ran}R\).

\subsection{Why the determining panel does not populate the angle}

The Grand Tensor determining panel (17.6) has a distinct and useful
role.

\begin{proposition}[Panel coercivity is direct residual control]
\label{prop:v14v17-panel}
If a panel \(D:Q\to Y\) satisfies
\[
 D^*D\succeq cI,\qquad c>0,                                 \tag{17.31}
\]
then
\[
 \|r\|\leq c^{-1/2}\|Dr\|.                                  \tag{17.32}
\]
However, (17.31) gives no positive lower bound for the
physical--tangent incidence \(K^*|_{\operatorname{Ran}P}\).
\end{proposition}

\begin{proof}
The first statement is immediate from
\(\|Dr\|^2\geq c\|r\|^2\).

For the second, take \(Q=\mathbb R^2\), \(D=I\), let \(K:\mathbb R
\to Q\) be \(Kt=te_1\), and let \(P_\theta\) project onto
\[
 u_\theta=(\sin\theta,\cos\theta),\qquad 0<\theta<\pi/2.
\]
Then \(D^*D=I\) for every \(\theta\), while
\[
 \operatorname{rmin}
 \bigl(K^*|_{\operatorname{Ran}P_\theta}\bigr)
 =\sin\theta\longrightarrow0.                              \tag{17.33}
\]
Thus the fixed determining-panel floor does not control the changing
physical--tangent angle.
\end{proof}

If the concrete panel residual \(Dr\) tends to zero with the common
floor (17.31), it directly proves \(r\to0\) and bypasses the angle
route.  It cannot be cited as proof of an incidence floor which it
does not contain.

The coherent Ward Gram is also an upper defect estimate: it proves
the Ward relation when its scaled difference tends to zero.  A lower
singular-value estimate is different data.  Likewise, the operator
HDA generator \(D_h[v]\) and the variation tangent \(L\) live in
different spaces until a same-history intertwiner is supplied.

\subsection{Typed source-interface card}

A concrete use of V12--V16 now requires one finite-stage card:
\[
\begin{array}{c|p{0.72\linewidth}}
\mathrm{I1}&
a geometric variation space \(Q_n\), faithful metric \(G_n\), and
total action covector \(\epsilon_n\);\\
\mathrm{I2}&
a relabeling parameter space \(X_n\) and the derivative
\(L_n:X_n\to Q_n\), giving \(K_n=G_n^{1/2}L_n\);\\
\mathrm{I3}&
a constraint map on the same Riesz-residual or admissible-variation
carrier, with the whitening and transport in (17.14);\\
\mathrm{I4}&
the exact gauge-tangency identity
\(\widetilde C_nK_n=0\), or a quantitative defect bound;\\
\mathrm{I5}&
the common-action Ward identity
\(L_n^*\epsilon_n=0\), or its controlled residual;\\
\mathrm{I6}&
a declaration whether the target is the raw equation \(r_n=0\) or
the constrained quotient equation \(R_nr_n=0\);\\
\mathrm{I7}&
same-history transport of all these maps and residual Grams to the
continuum screen.
\end{array}                                                  \tag{17.34}
\]

The inspected sources contain conditional or abstract versions of
I1, I2, I5, and the direct determining-panel alternative.  They do
not populate I3--I4 as a single matrix card on the represented
Einstein residual screen.  The attached V30 source supplies extensive
physical source-metric and tangent validation machinery but no
replacement for this typed gravity constraint card.

\subsection{Next upstream theorem}

When exact tangency (17.16) is unavailable, the first useful number is
\[
 \varepsilon_{\rm tan}:=\|(I-P)T\|,                         \tag{17.35}
\]
the sine of the largest constraint-escape angle of a tangent vector.
V18 will prove:

\begin{enumerate}
\item an explicit bound for (17.35) in terms of the same-history
constraint-propagation defect \(\|\widetilde C K\|\) and the
condition numbers of \(\widetilde C\) and \(K\);
\item a lower bound
\[
 \beta(P,S)\geq\sqrt{1-\varepsilon_{\rm tan}^2}
\]
under the precise closed-range hypotheses;
\item quantitative stability of \(R\), the joint pseudoinverse, and
the projected residual Grams;
\item an exact escape vector when the bound fails.
\end{enumerate}

This moves the open angle row from a large completed normal symbol to
the physically earlier constraint-propagation identity.  V16 remains
the fallback if the tangency defect is not small.

\subsection{Acquisition ledger}

\[
\begin{array}{p{0.29\linewidth}|p{0.23\linewidth}|p{0.38\linewidth}}
\textbf{gate}&\textbf{status after V17}&\textbf{content}\\ \hline
\text{source residual typing}
 &\text{acquired}
 &metric whitening of the action covector\\
\text{Ward-to-transverse interface}
 &\text{acquired}
 &L^*\epsilon=0\Longleftrightarrow K^*r=0\\
\text{exact gauge-tangent nesting}
 &\text{acquired conditionally}
 &\widetilde CK=0\Rightarrow R=P-T,\ \gamma=1\\
\text{joint multiplier on nested branch}
 &\text{acquired}
 &{\mathfrak C}^{\dagger}={\mathfrak C}^*\\
\text{variational quotient meaning}
 &\text{acquired}
 &Rr=0\text{ is stationarity modulo multipliers}\\
\text{determining-panel substitution}
 &\text{refuted}
 &fixed panel floor need not control the angle\\
\text{same-history constraint propagation}
 &\text{open}
 &populate \(\widetilde C_nK_n\)\\
\text{raw versus quotient target}
 &\text{open declaration}
 &depends on the selected gravity formulation
\end{array}                                                  \tag{17.36}
\]

\subsection{Parent-integrity record}

The installed parent was
\[
\begin{array}{c|c}
 \text{file}&
 \texttt{Renewal\_SM\_Einstein\_Continuum\_Research\_Notes\_V16.tex}\\
 \text{lines}&6331\\
 \text{bytes}&205772\\
 \text{SHA--256}&
 \texttt{BFA97971E5037A26F713B7B1CEBF66C343517EC188E860B8DC4C6097AE9EC06D}.
\end{array}                                                  \tag{17.37}
\]

% =============================================================================
% END APPENDED FOURTEENTH-CYCLE V17 MATERIAL
% =============================================================================
% =============================================================================
% BEGIN APPENDED FOURTEENTH-CYCLE V18 MATERIAL
% =============================================================================

\section{Fourteenth-cycle V18: quantitative constraint propagation,
sharp angle control, and three-defect Einstein closure}
\label{sec:v14v18}

\subsection{Setup}

Let \({\cal Z},{\cal H}_C,{\cal X}\) be Hilbert spaces.  Let
\[
 C:{\cal Z}\to{\cal H}_C,\qquad
 K:{\cal X}\to{\cal Z}                                     \tag{18.1}
\]
be bounded operators with closed range.  Define the physical and
tangent projections
\[
 P:=I-C^\dagger C,\qquad
 E:=I-P=C^\dagger C,\qquad
 T:=KK^\dagger,\qquad S:=I-T.                              \tag{18.2}
\]
Thus
\[
 \operatorname{Ran}P=\ker C,\qquad
 \operatorname{Ran}T=\operatorname{Ran}K.                  \tag{18.3}
\]
The exact gauge-tangency condition of V17 is \(CK=0\).  Its geometric
projection defect is
\[
 \varepsilon:=\|ET\|
 =\|(I-P)T\|.                                               \tag{18.4}
\]
This is the largest sine of constraint escape among unit tangent
vectors.

\subsection{The sharp two-projection identity}

Let
\[
 {\cal A}:=\operatorname{Ran}P,\qquad
 {\cal G}:=\operatorname{Ran}T,\qquad
 {\cal M}:={\cal A}\cap{\cal G}^{\perp}.                    \tag{18.5}
\]
The V13 incidence is
\[
 B:=T|_{\cal A}:{\cal A}\to{\cal G},                        \tag{18.6}
\]
with \(\ker B={\cal M}\).

\begin{theorem}[Constraint escape equals the complementary angle]
\label{thm:v14v18-sharp}
If
\[
 \varepsilon=\|(I-P)T\|<1,                                 \tag{18.7}
\]
then \(B\) has closed range equal to \({\cal G}\), and
\[
 \boxed{
 \beta(P,S)
 =\operatorname{rmin}B
 =\sqrt{1-\varepsilon^2}.}                                 \tag{18.8}
\]
Consequently the optimal V12 spectral gap is
\[
 \boxed{\gamma(P,S)=1-\varepsilon^2.}                      \tag{18.9}
\]
The exact tangency branch \(CK=0\) has \(\varepsilon=0\), recovering
\(\beta=\gamma=1\).
\end{theorem}

\begin{proof}
The adjoint of (18.6) is
\[
 B^*=P|_{\cal G}:{\cal G}\to{\cal A}.                       \tag{18.10}
\]
For every unit \(y\in{\cal G}\),
\[
 \|B^*y\|^2
 =\|Py\|^2
 =1-\|Ey\|^2.                                               \tag{18.11}
\]
Taking the infimum on the left and the supremum in the last term
gives
\[
 \operatorname{rmin}(B^*)
 =\sqrt{1-\|E|_{\cal G}\|^2}
 =\sqrt{1-\varepsilon^2}.                                  \tag{18.12}
\]
Under (18.7), \(B^*\) is bounded below.  Hence it is injective with
closed range, and the closed-range theorem gives
\[
 \operatorname{Ran}B=(\ker B^*)^\perp={\cal G}.             \tag{18.13}
\]
The restrictions of \(B\) and \(B^*\) to their kernel complements
have the same positive singular spectrum; equivalently, use the
polar decomposition of the surjection \(B\).  Their reduced minimum
moduli are therefore equal.  This proves (18.8).  The V13 exact
angle identity \(\gamma=\beta^2\) gives (18.9).  If \(CK=0\), then
\(ET=C^\dagger CKK^\dagger=0\).
\end{proof}

\begin{corollary}[Complementary compression estimate]
\label{cor:v14v18-compression}
Let \(R\) be the orthogonal projection onto \({\cal M}\).  Under
(18.7),
\[
 \boxed{\|PSP-R\|\leq\varepsilon^2,}                        \tag{18.14}
\]
and for \(m\geq1\),
\[
 \boxed{\|(PSP)^m-R\|\leq\varepsilon^{2m}.}                 \tag{18.15}
\]
\end{corollary}

\begin{proof}
On \({\cal M}\), \(PSP\) is the identity.  On
\({\cal A}\ominus{\cal M}\), V13 and (18.9) place its spectrum in
\([0,\varepsilon^2]\).  On \({\cal A}^{\perp}\), it is zero.
Functional calculus proves both estimates.
\end{proof}

Thus near gauge tangency makes alternating compression rapidly
convergent.  At exact tangency, one multiplication already gives
\(PSP=R=P-T\).

\subsection{The propagation defect controls the angle}

The operator identity
\[
 \boxed{
 ET=C^\dagger(CK)K^\dagger}                                \tag{18.16}
\]
is exact.

\begin{proof}
Use \(E=C^\dagger C\), \(T=KK^\dagger\), and associate the product.
\end{proof}

Define the same-history condition number
\[
 \eta:=\|C^\dagger\|\,\|CK\|\,\|K^\dagger\|.                \tag{18.17}
\]

\begin{theorem}[Constraint-propagation angle certificate]
\label{thm:v14v18-propagation}
One has
\[
 \varepsilon\leq\eta.                                      \tag{18.18}
\]
If \(\eta<1\), then
\[
 \boxed{
 \beta(P,S)\geq\sqrt{1-\eta^2},\qquad
 \gamma(P,S)\geq1-\eta^2.}                                 \tag{18.19}
\]
\end{theorem}

\begin{proof}
Take norms in (18.16) to obtain (18.18).  If \(\eta<1\), then
\(\varepsilon<1\).  Apply Theorem~\ref{thm:v14v18-sharp} and
\(\varepsilon\leq\eta\).
\end{proof}

All three factors in (18.17) are load-bearing.  A small raw defect
\(\|CK\|\) does not control the angle if either the constraint or
tangent range loses its reduced singular floor.

\begin{corollary}[Exact gauge tangency]
\label{cor:v14v18-exact}
If \(CK=0\), then
\[
 \operatorname{Ran}K\subseteq\ker C,\quad
 R=P-T,\quad
 \gamma=1,
\]
and the V16 completed-symbol route is unnecessary for the
physical--tangent intersection.
\end{corollary}

\begin{proof}
The first statement is immediate and the others follow from V17 or
Theorem~\ref{thm:v14v18-sharp}.
\end{proof}

\subsection{A sharper joint-constraint bound}

Let
\[
 {\mathfrak C}:{\cal Z}\to
 \operatorname{Ran}E\oplus\operatorname{Ran}T,\qquad
 {\mathfrak C}x=(Ex,Tx).                                   \tag{18.20}
\]
This is the V14 joint constraint with its codomain reduced to the
occurring two ranges.

\begin{theorem}[Joint pseudoinverse from near orthogonality]
\label{thm:v14v18-joint}
If \(\varepsilon<1\), then \({\mathfrak C}\) has closed range and
\[
 \boxed{
 \operatorname{rmin}({\mathfrak C})\geq\sqrt{1-\varepsilon},
 \qquad
 \|{\mathfrak C}^\dagger\|
 \leq(1-\varepsilon)^{-1/2}.}                              \tag{18.21}
\]
In terms of the computable propagation certificate,
\[
 \eta<1
 \quad\Longrightarrow\quad
 \|{\mathfrak C}^\dagger\|
 \leq(1-\eta)^{-1/2}.                                      \tag{18.22}
\]
\end{theorem}

\begin{proof}
On the reduced codomain,
\[
 {\mathfrak C}{\mathfrak C}^*
 =
 \begin{pmatrix}
 I_{\operatorname{Ran}E}&ET|_{\operatorname{Ran}T}\\
 TE|_{\operatorname{Ran}E}&I_{\operatorname{Ran}T}
 \end{pmatrix}.                                             \tag{18.23}
\]
For \(e\in\operatorname{Ran}E\) and
\(t\in\operatorname{Ran}T\),
\[
 \begin{aligned}
 \langle{\mathfrak C}{\mathfrak C}^*(e,t),(e,t)\rangle
 &=\|e\|^2+\|t\|^2+2\operatorname{Re}\langle e,t\rangle\\
 &\geq
 \|e\|^2+\|t\|^2-2\varepsilon\|e\|\|t\|\\
 &\geq(1-\varepsilon)(\|e\|^2+\|t\|^2).                    \tag{18.24}
 \end{aligned}
\]
Thus \({\mathfrak C}^*\) is bounded below by
\(\sqrt{1-\varepsilon}\).  It follows that
\({\mathfrak C}\) is onto its reduced codomain, has closed range,
and has the same positive singular lower bound.  This proves
(18.21).  Use \(\varepsilon\leq\eta\) for (18.22).
\end{proof}

This improves the general V14 bound when constraint and tangent
violations are nearly orthogonal.  It becomes the exact unit bound
at gauge tangency.

\subsection{Constraint and Ward defects control distance to the
intersection}

Let \(R\) again project onto
\({\cal M}=\ker C\cap(\operatorname{Ran}K)^\perp\).

\begin{lemma}[Individual defect estimates]
\label{lem:v14v18-defects}
For every \(r\in{\cal Z}\),
\[
 \|Er\|\leq\|C^\dagger\|\,\|Cr\|,                           \tag{18.25}
\]
\[
 \|Tr\|\leq\|K^\dagger\|\,\|K^*r\|.                        \tag{18.26}
\]
\end{lemma}

\begin{proof}
The first estimate follows from \(E=C^\dagger C\).
For the second, take adjoints in \(T=KK^\dagger\) to obtain
\[
 T=(K^\dagger)^*K^*,
\]
then take norms.
\end{proof}

\begin{theorem}[Three-defect residual estimate]
\label{thm:v14v18-three}
If \(\varepsilon<1\), then every \(r\in{\cal Z}\) satisfies
\[
\boxed{
\begin{aligned}
 \|r\|^2
 \leq{}&\|Rr\|^2\\
 &+\frac{
 \|C^\dagger\|^2\|Cr\|^2
 +\|K^\dagger\|^2\|K^*r\|^2}
 {1-\varepsilon}.
\end{aligned}}                                              \tag{18.27}
\]
If only the computable bound \(\eta<1\) is used, then
\[
\boxed{
\begin{aligned}
 \|r\|^2
 \leq{}&\|Rr\|^2\\
 &+\frac{
 \|C^\dagger\|^2\|Cr\|^2
 +\|K^\dagger\|^2\|K^*r\|^2}
 {1-\eta}.
\end{aligned}}                                              \tag{18.28}
\]
\end{theorem}

\begin{proof}
The V14 Moore--Penrose formula gives
\[
 I-R={\mathfrak C}^\dagger{\mathfrak C}.
\]
Hence Theorem~\ref{thm:v14v18-joint} and
Lemma~\ref{lem:v14v18-defects} give
\[
 \begin{aligned}
 \|(I-R)r\|^2
 &\leq\|{\mathfrak C}^\dagger\|^2
       \|{\mathfrak C}r\|^2\\
 &\leq
 \frac{\|Er\|^2+\|Tr\|^2}{1-\varepsilon}\\
 &\leq
 \frac{
 \|C^\dagger\|^2\|Cr\|^2
 +\|K^\dagger\|^2\|K^*r\|^2}
 {1-\varepsilon}.                                          \tag{18.29}
 \end{aligned}
\]
Since \(R\) is orthogonal,
\[
 \|r\|^2=\|Rr\|^2+\|(I-R)r\|^2,
\]
which proves (18.27).  Replacing
\(\varepsilon\) by its upper bound \(\eta\) gives (18.28).
\end{proof}

The three terms have distinct meanings:
\[
\begin{array}{c|c}
 \|Cr\|&\text{physical constraint defect}\\
 \|K^*r\|&\text{diffeomorphism Ward defect}\\
 \|Rr\|&\text{gauge-reduced physical Einstein residual}.
\end{array}                                                  \tag{18.30}
\]
No one of them can be omitted in the general constrained
formulation.

\begin{corollary}[Exact full-residual closure]
\label{cor:v14v18-exactclosure}
Under \(\varepsilon<1\),
\[
 Cr=0,\qquad K^*r=0,\qquad Rr=0
 \quad\Longrightarrow\quad r=0.                            \tag{18.31}
\]
\end{corollary}

\begin{proof}
Insert the three zero defects in (18.27).
\end{proof}

This recovers the strong V12 conclusion without assuming in advance
that \(r\) lies in both ranges.

\subsection{Cofinal three-defect theorem}

Let \(C_n,K_n\) be same-history finite-stage maps, with the associated
\(P_n,T_n,S_n,R_n\).  Let \(r_{a,n}\) be any fixed localized residual
screen, where \(a\) may encode the V9 test tensor and core multiplier.

\begin{theorem}[Uniform propagation closure]
\label{thm:v14v18-cofinal}
Assume constants \(c,k<\infty\) and \(0\leq\eta_0<1\) such that
\[
 \|C_n^\dagger\|\leq c,\qquad
 \|K_n^\dagger\|\leq k,\qquad
 \|C_nK_n\|\leq\frac{\eta_0}{ck}                           \tag{18.32}
\]
for all \(n\).  If, for every fixed index \(a\),
\[
 \|C_nr_{a,n}\|\longrightarrow0,\qquad
 \|K_n^*r_{a,n}\|\longrightarrow0,\qquad
 \|R_nr_{a,n}\|\longrightarrow0,                           \tag{18.33}
\]
then
\[
 \boxed{\|r_{a,n}\|\longrightarrow0.}                      \tag{18.34}
\]
If the localized products have the V9 all-order landing and uniform
integrability, the represented weak Einstein equation holds on the
declared limiting core.
\end{theorem}

\begin{proof}
The hypotheses imply
\[
 \eta_n
 =\|C_n^\dagger\|\,\|C_nK_n\|\,\|K_n^\dagger\|
 \leq\eta_0.
\]
Apply (18.28) at stage \(n\):
\[
 \|r_{a,n}\|^2
 \leq\|R_nr_{a,n}\|^2
 +\frac{
 c^2\|C_nr_{a,n}\|^2
 +k^2\|K_n^*r_{a,n}\|^2}
 {1-\eta_0}.
\]
Every term tends to zero by (18.33), proving (18.34).  The final
statement is the V9 localized-Gram landing theorem.
\end{proof}

\begin{corollary}[Uniform projector and multiplier landing]
\label{cor:v14v18-landing}
Under (18.32), if \(P_n\to P\) and \(S_n\to S\) in norm, then
\[
 R_n\to R,\qquad
 {\mathfrak C}_n^\dagger\to{\mathfrak C}^\dagger            \tag{18.35}
\]
in norm.
\end{corollary}

\begin{proof}
Theorem~\ref{thm:v14v18-propagation} gives the common angle gap
\[
 \gamma_n\geq1-\eta_0^2>0.
\]
Apply the V12 intersection-projection theorem and the V14
pseudoinverse theorem.  Alternatively, (18.21) gives the common
joint singular floor directly.
\end{proof}

\subsection{Escape witnesses and certificate failure}

\begin{proposition}[Tangent constraint-escape sequence]
\label{prop:v14v18-escape}
For every \(\rho<\varepsilon\), there is a unit vector
\(y_\rho\in\operatorname{Ran}K\) such that
\[
 \|(I-P)y_\rho\|>\rho,\qquad
 \|Py_\rho\|<\sqrt{1-\rho^2}.                              \tag{18.36}
\]
If \(\varepsilon=1\), there is a sequence of unit tangent vectors
\(y_m\) with
\[
 \|(I-P)y_m\|\to1,\qquad \|Py_m\|\to0.                     \tag{18.37}
\]
\end{proposition}

\begin{proof}
The definition
\(\varepsilon=\sup_{\|y\|=1,y\in\operatorname{Ran}K}\|Ey\|\)
gives \(y_\rho\).  Orthogonality of \(P\) and \(E=I-P\) gives
\[
 \|Py_\rho\|^2=1-\|Ey_\rho\|^2<1-\rho^2.
\]
Choose \(\rho\uparrow1\) for (18.37).
\end{proof}

The vectors in (18.37) are exact geometric diagnostics: tangent
directions escape almost completely from the declared physical
constraint kernel.  They do not by themselves refute the underlying
Einstein regulator; they refute the proposed constraint/tangent
incidence card.

If the upper estimate \(\eta\) in (18.17) is at least one while
\(\varepsilon<1\), only the condition-number certificate has failed.
One must compute \(ET\) more sharply before declaring a geometric
obstruction.

\subsection{Finite acquisition card}

The quantitative upstream card is now
\[
\begin{array}{c|p{0.72\linewidth}}
\mathrm{T1}&
exact matrices for \(C_n,K_n\) on one whitened same-history screen;\\
\mathrm{T2}&
certified reduced singular floors for \(C_n\) and \(K_n\), equivalently
upper bounds for their Moore--Penrose inverses;\\
\mathrm{T3}&
the propagation residual \(C_nK_n\), with an exact zero identity or
an outward norm bound;\\
\mathrm{T4}&
the resulting \(\eta_n<1\), angle gap, joint pseudoinverse bound, and
intersection projector;\\
\mathrm{T5}&
the three residual channels in (18.33), evaluated on the identical
localized Einstein vectors;\\
\mathrm{T6}&
cofinal product landing on the V9 represented core.
\end{array}                                                  \tag{18.38}
\]

This card is smaller and earlier than the V16 Laurent completed-symbol
card.  V16 remains available if T2--T3 cannot make \(\eta<1\).

\subsection{Acquisition ledger and next task}

\[
\begin{array}{p{0.29\linewidth}|p{0.23\linewidth}|p{0.38\linewidth}}
\textbf{gate}&\textbf{status after V18}&\textbf{content}\\ \hline
\text{constraint-escape angle}
 &\text{acquired sharply}
 &\gamma=1-\|(I-P)T\|^2\\
\text{propagation certificate}
 &\text{acquired}
 &\eta=\|C^\dagger\|\|CK\|\|K^\dagger\|\\
\text{joint pseudoinverse}
 &\text{sharpened}
 &\|{\mathfrak C}^\dagger\|\leq(1-\varepsilon)^{-1/2}\\
\text{full residual closure}
 &\text{acquired quantitatively}
 &constraint \(+\) Ward \(+\) \(R\)-Gram\\
\text{cofinal closure}
 &\text{acquired conditionally}
 &uniform \(\eta_0<1\) and three vanishing defects\\
\text{finite \(C_n,K_n\) matrices}
 &\text{open}
 &same-history source card T1\\
\text{uniform singular floors}
 &\text{open}
 &source card T2\\
\text{propagation residual values}
 &\text{open}
 &source card T3
\end{array}                                                  \tag{18.39}
\]

The next version should go one step further upstream and derive
\(CK=0\) from invariance of the constraint map under the finite
relabeling group.  If the constraint is equivariant, differentiation
on the constraint surface should give exact gauge tangency; away from
the surface it should give an explicit defect proportional to the
constraint residual.  This would replace T3 by a structural theorem
and leave only the conditioning and source-value rows to populate.

\subsection{Parent-integrity record}

The installed parent was
\[
\begin{array}{c|c}
 \text{file}&
 \texttt{Renewal\_SM\_Einstein\_Continuum\_Research\_Notes\_V17.tex}\\
 \text{lines}&6833\\
 \text{bytes}&223217\\
 \text{SHA--256}&
 \texttt{2DF15DFF9EF2F07E4C6484C0ABC418A50E69A98E669F4E7ED532E3A3A41663A2}.
\end{array}                                                  \tag{18.40}
\]

% =============================================================================
% END APPENDED FOURTEENTH-CYCLE V18 MATERIAL
% =============================================================================
% =============================================================================
% BEGIN APPENDED FOURTEENTH-CYCLE V19 MATERIAL
% =============================================================================

\section{Fourteenth-cycle V19: equivariant constraints generate gauge
tangency and asymptotic nesting}
\label{sec:v14v19}

\subsection{The nonlinear source of \(CK\)}

V18 treated the linearized constraint map \(C\) and gauge tangent
\(K\) as a same-history matrix pair.  Their product is not an
independent datum when both arise by differentiating an equivariant
nonlinear constraint system.  This note proves the exact identity,
includes an equivariance defect, and propagates the nonlinear
constraint residual through time.

Let \(Q\) and \(Y\) be finite-dimensional real or complex Hilbert
spaces, let \(U\subset Q\) be open, and let
\[
 \Phi:U\longrightarrow Y                                   \tag{19.1}
\]
be a \(C^1\) constraint map.  Let a finite-dimensional Lie group
\({\cal G}\) act \(C^1\)-smoothly on \(U\), and let
\[
 \rho:{\cal G}\longrightarrow{\rm GL}(Y)                   \tag{19.2}
\]
be a differentiable representation.  For
\(\xi\) in the Lie algebra \(\mathfrak g\), define
\[
 L_q\xi
 :=\left.\frac{d}{dt}\right|_{t=0}
       \exp(t\xi)\cdot q,\qquad
 A_\xi:=d\rho(\xi).                                         \tag{19.3}
\]
The linearized constraint is
\[
 C_q:=D\Phi_q:Q\longrightarrow Y.                           \tag{19.4}
\]

\subsection{Exact equivariance}

\begin{theorem}[Equivariance-to-tangency identity]
\label{thm:v14v19-equivariance}
Assume
\[
 \Phi(g\cdot q)=\rho(g)\Phi(q)                              \tag{19.5}
\]
whenever both sides are defined.  Then
\[
 \boxed{
 C_qL_q\xi=A_\xi\Phi(q)}
 \qquad(\xi\in\mathfrak g).                                \tag{19.6}
\]
In particular, on the constraint surface
\[
 \Phi(q)=0
 \quad\Longrightarrow\quad
 \boxed{C_qL_q=0.}                                         \tag{19.7}
\]
Hence the infinitesimal group orbit is tangent to
\(\Phi^{-1}(0)\).
\end{theorem}

\begin{proof}
Apply (19.5) to \(g=\exp(t\xi)\) and differentiate at \(t=0\).
The chain rule gives \(D\Phi_q[L_q\xi]\) on the left and
\(d\rho(\xi)\Phi(q)\) on the right, proving (19.6).
Equation (19.7) follows when \(\Phi(q)=0\).
\end{proof}

\begin{remark}[No regular-value hypothesis]
The identity (19.7) does not require \(0\) to be a regular value.
Regularity or a reduced singular floor for \(C_q\) is needed for a
uniform orthogonal projector, but not for the algebraic tangency
identity.
\end{remark}

\subsection{Metric whitening cancels from the propagation product}

Let \(G_q=G_q^*>0\) be the source metric on \(Q\).  Define the
whitened maps
\[
 \widetilde C_q:=C_qG_q^{-1/2},\qquad
 K_q:=G_q^{1/2}L_q.                                        \tag{19.8}
\]

\begin{proposition}[Metric-free middle product]
\label{prop:v14v19-whitening}
One has the exact identity
\[
 \boxed{\widetilde C_qK_q=C_qL_q.}                         \tag{19.9}
\]
Under (19.5),
\[
 \boxed{
 \widetilde C_qK_q\,\xi=A_\xi\Phi(q).}                     \tag{19.10}
\]
\end{proposition}

\begin{proof}
The positive square roots in (19.8) cancel:
\[
 \widetilde C_qK_q
 =C_qG_q^{-1/2}G_q^{1/2}L_q=C_qL_q.
\]
Insert (19.6).
\end{proof}

Thus the V18 angle certificate is independent of a choice of
coordinates in the middle variation space, provided the constraint
and tangent maps are whitened by the same metric.

Assume the infinitesimal target representation satisfies
\[
 \|A_\xi y\|_Y\leq a_q\|\xi\|_{\mathfrak g}\|y\|_Y.         \tag{19.11}
\]

\begin{corollary}[Off-shell propagation bound]
\label{cor:v14v19-offshell}
Under exact equivariance,
\[
 \boxed{
 \|\widetilde C_qK_q\|
 \leq a_q\|\Phi(q)\|.}                                     \tag{19.12}
\]
\end{corollary}

\begin{proof}
Take the supremum of (19.10) over unit \(\xi\) and use (19.11).
\end{proof}

\subsection{Approximate equivariance}

Define the finite equivariance defect
\[
 {\cal D}(g,q):=\Phi(g\cdot q)-\rho(g)\Phi(q).               \tag{19.13}
\]
Assume it is differentiable in the group direction at the identity,
and put
\[
 \delta_{\rm eq}(q)\xi
 :=
 \left.\frac{d}{dt}\right|_{t=0}
 {\cal D}(\exp(t\xi),q).                                    \tag{19.14}
\]

\begin{theorem}[Differentiated equivariance defect]
\label{thm:v14v19-defect}
One has
\[
 \boxed{
 \widetilde C_qK_q\,\xi
 =A_\xi\Phi(q)+\delta_{\rm eq}(q)\xi.}                      \tag{19.15}
\]
Consequently,
\[
 \boxed{
 \|\widetilde C_qK_q\|
 \leq a_q\|\Phi(q)\|+\|\delta_{\rm eq}(q)\|.}               \tag{19.16}
\]
\end{theorem}

\begin{proof}
Differentiate (19.13) at \(g=\exp(t\xi)\).  The left term gives
\(C_qL_q\xi\), the representation term gives \(A_\xi\Phi(q)\), and
the difference is (19.14).  Use (19.9) and take norms.
\end{proof}

A small finite equivariance defect without control of its derivative
does not imply (19.16).  The differentiated defect is the object
which enters the tangent angle.

\subsection{Population of the V18 angle}

Assume \(\widetilde C_q\) and \(K_q\) have closed range.  Put
\[
 P_q:=I-\widetilde C_q^\dagger\widetilde C_q,\qquad
 T_q:=K_qK_q^\dagger,\qquad S_q:=I-T_q.                    \tag{19.17}
\]
Let
\[
 c_q:=\|\widetilde C_q^\dagger\|,\qquad
 k_q:=\|K_q^\dagger\|,\qquad
 b_q:=a_q\|\Phi(q)\|+\|\delta_{\rm eq}(q)\|.                \tag{19.18}
\]

\begin{theorem}[Equivariance angle certificate]
\label{thm:v14v19-angle}
If
\[
 \zeta_q:=c_qk_qb_q<1,                                     \tag{19.19}
\]
then the physical--transverse angle satisfies
\[
 \boxed{
 \gamma(P_q,S_q)\geq1-\zeta_q^2,}                          \tag{19.20}
\]
and the V14 joint constraint obeys
\[
 \boxed{
 \|{\mathfrak C}_q^\dagger\|
 \leq(1-\zeta_q)^{-1/2}.}                                  \tag{19.21}
\]
If \(\Phi(q)=0\) and \(\delta_{\rm eq}(q)=0\), then
\[
 \zeta_q=0,\qquad
 R_q=P_q-T_q,\qquad
 \gamma(P_q,S_q)=1,\qquad
 {\mathfrak C}_q^\dagger={\mathfrak C}_q^*.                \tag{19.22}
\]
\end{theorem}

\begin{proof}
Theorem~\ref{thm:v14v19-defect} gives
\(\|\widetilde C_qK_q\|\leq b_q\).  Therefore the V18 propagation
condition number is bounded by
\[
 \|\widetilde C_q^\dagger\|
 \,\|\widetilde C_qK_q\|
 \,\|K_q^\dagger\|
 \leq\zeta_q.
\]
Apply the V18 angle and joint-pseudoinverse theorems.  The exact branch
is V17.
\end{proof}

This reduces the open angle row to four earlier quantities: the
nonlinear constraint residual, differentiated equivariance defect,
and the two reduced singular floors.

\subsection{Constraint propagation in physical time}

Let \(q:[0,\tau]\to U\) be an absolutely continuous finite-stage
trajectory and set
\[
 y(t):=\Phi(q(t)).                                          \tag{19.23}
\]
Assume the actual evolution gives, almost everywhere,
\[
 \dot y(t)=A(t)y(t)+e_{\rm prop}(t),                         \tag{19.24}
\]
where
\[
 \|A(t)\|\leq a(t),\qquad
 \|e_{\rm prop}(t)\|\leq d(t),                              \tag{19.25}
\]
with \(a,d\in L^1(0,\tau)\).

\begin{theorem}[Quantitative constraint propagation]
\label{thm:v14v19-time}
For \(0\leq t\leq\tau\),
\[
\boxed{
 \|y(t)\|
 \leq
 e^{\int_0^t a(s)\,ds}
 \left(
  \|y(0)\|+\int_0^t d(s)\,ds
 \right).}                                                  \tag{19.26}
\]
More sharply,
\[
 \|y(t)\|
 \leq e^{\int_0^t a}
 \|y(0)\|
 +\int_0^t
 e^{\int_s^t a(u)\,du}d(s)\,ds.                            \tag{19.27}
\]
In particular, exact propagation
\(e_{\rm prop}=0\) preserves the constraint surface.
\end{theorem}

\begin{proof}
Absolute continuity and (19.24)--(19.25) give
\[
 \frac{d}{dt}\|y(t)\|
 \leq a(t)\|y(t)\|+d(t)
\]
for almost every \(t\), using the upper Dini derivative at zeros.
Multiplication by
\(\exp(-\int_0^t a)\), integration, and multiplication back give
(19.27).  Bounding the exponential inside the integral by the full
factor gives (19.26).  If \(y(0)=0=d\), both bounds give \(y(t)=0\).
\end{proof}

\begin{corollary}[Time-uniform angle margin]
\label{cor:v14v19-timeangle}
Suppose on \(0\leq t\leq\tau\),
\[
 c_{q(t)}\leq c,\qquad
 k_{q(t)}\leq k,\qquad
 a_{q(t)}\leq a_*,\qquad
 \|\delta_{\rm eq}(q(t))\|\leq d_{\rm eq}.                 \tag{19.28}
\]
Let the right side of (19.26) be bounded by \(r_*\).  If
\[
 ck(a_*r_*+d_{\rm eq})\leq\zeta_*<1,                        \tag{19.29}
\]
then
\[
 \gamma(P_{q(t)},S_{q(t)})
 \geq1-\zeta_*^2
\]
uniformly for \(0\leq t\leq\tau\).
\end{corollary}

\begin{proof}
Use Theorem~\ref{thm:v14v19-time} in
Theorem~\ref{thm:v14v19-angle}.
\end{proof}

This is the missing logical bridge named by the Grand Tensor
constraint-propagation firewall: initial constraint control alone is
insufficient unless (19.24) has a controlled propagator and source
defect.

\subsection{Asymptotic nesting along a regulator history}

Let the preceding objects carry an index \(n\).  Assume
\[
 \sup_n c_n<\infty,\qquad
 \sup_n k_n<\infty,\qquad
 \sup_n a_n<\infty,                                        \tag{19.30}
\]
and
\[
 \|\Phi_n(q_n)\|\longrightarrow0,\qquad
 \|\delta_{{\rm eq},n}(q_n)\|\longrightarrow0.              \tag{19.31}
\]

\begin{theorem}[Asymptotically exact nesting]
\label{thm:v14v19-asymptotic}
Under (19.30)--(19.31), define
\[
 \varepsilon_n:=\|(I-P_n)T_n\|.
\]
Then
\[
 \varepsilon_n\longrightarrow0,\qquad
 \gamma(P_n,S_n)\longrightarrow1.                          \tag{19.32}
\]
Moreover
\[
 \|[P_n,T_n]\|\leq2\varepsilon_n\longrightarrow0,           \tag{19.33}
\]
\[
 \|P_nS_n-(P_n-T_n)\|
 =\|(I-P_n)T_n\|
 =\varepsilon_n\longrightarrow0,                           \tag{19.34}
\]
and
\[
 \|R_n-(P_n-T_n)\|
 \leq2\varepsilon_n+\varepsilon_n^2
 \longrightarrow0.                                         \tag{19.35}
\]
\end{theorem}

\begin{proof}
Equations (19.16), (19.30), and (19.31) make the V18 condition number
\(\eta_n\) tend to zero.  Hence
\(\varepsilon_n\leq\eta_n\to0\), and V18 gives
\(\gamma_n=1-\varepsilon_n^2\to1\).

Since
\[
 [P_n,T_n]=T_n(I-P_n)-(I-P_n)T_n,
\]
(19.33) follows by the triangle inequality and adjoint symmetry.
Also
\[
 P_nS_n-(P_n-T_n)
 =T_n-P_nT_n=(I-P_n)T_n,
\]
which proves (19.34).  Finally, V18 gives
\(\|P_nS_nP_n-R_n\|\leq\varepsilon_n^2\), while
\[
 T_n-P_nT_nP_n
 =(I-P_n)T_n+P_nT_n(I-P_n)
\]
has norm at most \(2\varepsilon_n\).  Combining these identities
proves (19.35).
\end{proof}

Thus the simple nested formula \(P_n-T_n\) becomes asymptotically
correct with a quantified error even before exact nesting is reached.

\subsection{A necessary counterexample}

\begin{countertheorem}[On-shell constraint without equivariance does
not imply tangency]
\label{cth:v14v19-nonequiv}
Let \(Q=\mathbb R^2\), let the additive group \(\mathbb R\) act by
\[
 t\cdot(x,y)=(x+t,y),
\]
and let
\[
 \Phi(x,y)=x.
\]
At \(q=(0,0)\),
\[
 \Phi(q)=0,
\]
but the infinitesimal tangent and linearized constraint are
\[
 L_q\xi=(\xi,0),\qquad C_q(u,v)=u,
\]
so
\[
 \boxed{C_qL_q=I_{\mathbb R}\ne0.}                          \tag{19.36}
\]
The group orbit leaves the constraint surface immediately.
\end{countertheorem}

\begin{proof}
All statements follow by differentiation.  The constraint is not
equivariant under a target representation which fixes zero, since
\(\Phi(t\cdot q)=t\ne0\).
\end{proof}

This example shows why a zero constraint value cannot replace the
equivariance or propagation row.

\subsection{Interaction with the three-defect Einstein estimate}

Let \(r_n\) be the whitened total Einstein residual on a fixed
localized screen.  Combining V18 with
Theorem~\ref{thm:v14v19-angle} gives
\[
\begin{aligned}
 \|r_n\|^2
 \leq{}&\|R_nr_n\|^2\\
 &+\frac{
 c_n^2\|\widetilde C_nr_n\|^2
 +k_n^2\|K_n^*r_n\|^2}
 {1-\zeta_n},                                               \tag{19.37}
\end{aligned}
\]
whenever \(\zeta_n<1\).  Here
\[
 \zeta_n
 =
 c_nk_n
 \bigl(
  a_n\|\Phi_n(q_n)\|
  +\|\delta_{{\rm eq},n}(q_n)\|
 \bigr).                                                    \tag{19.38}
\]
Therefore a complete same-history continuum route can populate the
denominator using nonlinear constraint propagation, rather than a
separate spectral search for \(K_n^*P_n\).

The three numerators in (19.37) remain independent.  Equivariance
controls the geometry of the decomposition; it does not make the
Einstein residual, Ward residual, or physical quotient residual
vanish.

\subsection{Acquisition card and next checkpoint}

The remaining finite data are now
\[
\begin{array}{c|p{0.72\linewidth}}
\mathrm{E1}&
an explicit nonlinear constraint map \(\Phi_n\) and relabeling
action on one finite gravity regulator;\\
\mathrm{E2}&
an exact equivariance identity or differentiated defect interval;\\
\mathrm{E3}&
uniform reduced singular floors for
\(\widetilde C_n=D\Phi_nG_n^{-1/2}\) and
\(K_n=G_n^{1/2}L_n\);\\
\mathrm{E4}&
initial constraint and time-propagation residual bounds;\\
\mathrm{E5}&
the three localized residual channels of (19.37);\\
\mathrm{E6}&
the V9 all-order landing and core-density hypotheses.
\end{array}                                                  \tag{19.39}
\]

V20 is the next compulsory companion audit.  After that audit the
next substantive version should use the source's finite graph
relabeling/HDA data to ask whether an explicit \(\Phi_n,L_n\) pair
can be extracted.  If not, E1 is the exact blocking row.

\subsection{Acquisition ledger}

\[
\begin{array}{p{0.29\linewidth}|p{0.23\linewidth}|p{0.38\linewidth}}
\textbf{gate}&\textbf{status after V19}&\textbf{content}\\ \hline
\text{equivariance-to-tangency}
 &\text{acquired}
 &D\Phi_qL_q=A\Phi(q)\\
\text{metric-whitened product}
 &\text{acquired}
 &\widetilde C_qK_q=C_qL_q\\
\text{approximate equivariance}
 &\text{acquired}
 &A\Phi+\delta_{\rm eq}\\
\text{constraint propagation}
 &\text{acquired}
 &Gronwall bound with source defect\\
\text{asymptotic nesting}
 &\text{acquired}
 &R_n-(P_n-T_n)\to0\\
\text{concrete nonlinear gravity constraint}
 &\text{open}
 &source card E1\\
\text{uniform linearized floors}
 &\text{open}
 &source card E3\\
\text{three residual numerators}
 &\text{open}
 &source card E5
\end{array}                                                  \tag{19.40}
\]

\subsection{Parent-integrity record}

The installed parent was
\[
\begin{array}{c|c}
 \text{file}&
 \texttt{Renewal\_SM\_Einstein\_Continuum\_Research\_Notes\_V18.tex}\\
 \text{lines}&7381\\
 \text{bytes}&239119\\
 \text{SHA--256}&
 \texttt{9F7D90F3B6D518815E8C73484E6BD1B3347A8C73D10D60C263212095A398E42A}.
\end{array}                                                  \tag{19.41}
\]

% =============================================================================
% END APPENDED FOURTEENTH-CYCLE V19 MATERIAL
% =============================================================================
% =============================================================================
% BEGIN APPENDED FOURTEENTH-CYCLE V20 MATERIAL
% =============================================================================

\section{Fourteenth-cycle V20: compulsory companion audit and the
first-class-constraint direction}
\label{sec:v14v20}

\subsection{Files inspected}

The required twentieth-version audit was performed read-only on the
newest numerical companion notes in the shared folder:
\[
\begin{array}{p{0.53\linewidth}|r|r}
\textbf{file}&\textbf{lines}&\textbf{bytes}\\ \hline
\texttt{Renewal\_Yang\_Mills\_Mass\_Gap\_Research\_Notes\_V6.tex}
 &1690&61699\\
\texttt{Renewal\_NS\_Stability\_Research\_Notes\_V26.tex}
 &5319&278283
\end{array}                                                  \tag{20.1}
\]
Their SHA--256 digests are
\[
\begin{aligned}
 &\texttt{03E0CF21CEFD3D5127133E6D39D9DED746BAEA9CC2BB9D6A3F86BA28AEE63ECB},
                                                               \tag{20.2}\\
 &\texttt{82C887F8C2C69E4F28FAD7C3DC8DE5B9099CB5A4BF431EBA1FFF3E91935978AD}.
                                                               \tag{20.3}
\end{aligned}
\]
No companion file was modified or removed.

\subsection{Yang--Mills V6}

Since the V15 audit, the Yang--Mills programme has advanced from V2
to V6.  It now contains:

\begin{enumerate}
\item exact noncommutative matrix splittings on active momentum
circles;
\item exact rational congruence certificates at all transverse
quarter-turn anchors;
\item a no-go theorem for transporting one frozen splitting across a
whole quarter cell by a global coefficientwise Frobenius stock;
\item a moving Hermitian splitting whose Laurent degree is at most one
in each remaining phase;
\item exact positive grid floors for that moving splitting and a
recursive plan which doubles the block dimension while eliminating
one phase at a time.
\end{enumerate}

The companion note carefully separates exact anchor certificates from
sampled positivity between anchors.  Its fixed-splitting no-go refutes
only a coarse proof device, not the direct pencil.  Its V6 recursive
construction has proved the first coordinate elimination and the
grid input for the second; it has not yet proved a full-torus
inequality or a mass gap.

The transferable result concerns the V16 fallback.  If an Einstein
kernel-completed normal symbol must be treated by Laurent positivity,
a globally frozen auxiliary matrix and a coarse ambient norm may
discard the cancellations needed for a useful cover.  A moving
matrix-valued splitting with controlled Fourier support can instead
be eliminated recursively.  No Yang--Mills floor or coefficient is
an Einstein input.

\subsection{Navier--Stokes V26}

The Navier--Stokes note is unchanged from V15.  Its exact rank-one
kernel-derivative calculation continues to support one rule: use a
restricted norm only after the actual perturbation is proved to lie
in the restricted input class.  It supplies no gravity constraint,
HDA, Ward, or Einstein residual estimate.

\subsection{Audit decision}

\begin{decision}[Primary and fallback routes after V20]
\label{dec:v14v20-route}
The primary route remains the V19 equivariant-constraint route:
derive \(C_qL_q=A\Phi(q)\) from the same finite gravitational
constraint algebra.  If exact first-class closure is available on the
constraint surface, the V17 nested branch gives angle gap one and
bypasses a completed-symbol search.

The V16 Laurent route remains a fallback when first-class tangency is
not exact or when a quantitative off-shell certificate is needed.
On that branch, use a moving matrix splitting and recursive
coordinate elimination; do not assume that a frozen splitting plus a
global Frobenius derivative stock will give a useful cover.
\end{decision}

\begin{proof}
The V19 theorem shows that equivariance controls the physical--tangent
angle at an earlier logical stage than positivity of \(K^*P\).
The Yang--Mills results concern how to certify the latter positivity
once a concrete Laurent pencil exists.  The routes are therefore
ordered rather than competing.  The Navier--Stokes result adds only
the structured-input warning.
\end{proof}

\subsection{The Grand Tensor HDA suggests the next source theorem}

The inspected Grand Tensor source has finite self-adjoint normal and
tangential generators satisfying
\[
\begin{aligned}
 -i[D_h[v],D_h[w]]&=D_h[[v,w]_h],\\
 -i[D_h[v],H_h[N]]&=H_h[\ell_h(v)N],\\
 -i[H_h[N],H_h[M]]
 &=D_h[G_h^{-1}(N\,d_hM-M\,d_hN)].
\end{aligned}                                               \tag{20.4}
\]
These are first-class closure identities: commutators of constraint
generators remain in the constraint span.  What has not yet been
written in the active SM notes is the exact bridge from (20.4) to the
V19 nonlinear constraint map.

V21 will prove that bridge in two settings:

\begin{enumerate}
\item a finite-dimensional Poisson phase space, where the constraint
map consists of the Hamiltonian and momentum constraint functions;
\item a finite density-state carrier, where the constraint map
consists of expectation values of the HDA generators and the gauge
orbit is their unitary conjugation.
\end{enumerate}

On either branch, first-class closure will imply
\[
 D\Phi_q\,L_q=A_q\Phi(q),
\]
hence exact tangency on the zero-constraint surface.  The theorem will
also identify the anomaly term when the HDA closes only up to a
residual.  Applying it to the Grand Tensor construction will still
require one declaration of the state/phase-space carrier and proof
that its constraint map is the physical one; operator bracket
identities alone do not choose that interpretation.

\subsection{Acquisition ledger}

\[
\begin{array}{p{0.29\linewidth}|p{0.23\linewidth}|p{0.38\linewidth}}
\textbf{gate}&\textbf{status after V20}&\textbf{content}\\ \hline
\text{companion checkpoint}
 &\text{acquired}
 &YM V6 and NS V26 read with hashes\\
\text{moving Laurent splitting}
 &\text{retained for fallback}
 &recursive matrix-valued coordinate elimination\\
\text{fixed-splitting global stock}
 &\text{rejected as general route}
 &exact YM no-cover witness\\
\text{primary angle route}
 &\text{unchanged}
 &equivariant constraint propagation\\
\text{HDA-to-equivariance bridge}
 &\text{next}
 &first-class constraint theorem\\
\text{physical carrier declaration}
 &\text{open}
 &phase-space or density-state branch\\
\text{three Einstein residual defects}
 &\text{open}
 &same-history source values
\end{array}                                                  \tag{20.5}
\]

The next compulsory companion audit is V25.

\subsection{Parent-integrity record}

The installed parent was
\[
\begin{array}{c|c}
 \text{file}&
 \texttt{Renewal\_SM\_Einstein\_Continuum\_Research\_Notes\_V19.tex}\\
 \text{lines}&7917\\
 \text{bytes}&253802\\
 \text{SHA--256}&
 \texttt{F9EE0F812EB7FBA3ED548720FD4423A7BCD4532D158BDA0A7146E58781AB9592}.
\end{array}                                                  \tag{20.6}
\]

% =============================================================================
% END APPENDED FOURTEENTH-CYCLE V20 MATERIAL
% =============================================================================
% =============================================================================
% BEGIN APPENDED FOURTEENTH-CYCLE V21 MATERIAL
% =============================================================================

\section{Fourteenth-cycle V21: first-class HDA closure implies
constraint-surface tangency}
\label{sec:v14v21}

\subsection{Purpose}

V19 proved that equivariance of a nonlinear constraint map implies
gauge tangency.  The Grand Tensor source presents the relevant
symmetry primarily as a finite hypersurface-deformation algebra
(HDA), namely closure of normal and tangential generators under
brackets.  This note supplies the missing bridge.  First-class
closure makes the constraint map equivariant along the Hamiltonian
gauge flow, and any bracket anomaly becomes exactly the differentiated
equivariance defect used in V19.

Two realizations are kept separate.  The first is a classical
finite-dimensional Poisson regulator.  The second uses expectation
constraints on a finite density-state carrier.  The second branch is
not silently identified with the stronger Dirac condition that every
constraint operator annihilates a state vector.

\subsection{Classical first-class constraints}

Let \(({\cal M},\{\cdot,\cdot\})\) be a finite-dimensional Poisson
manifold.  Let
\[
 \phi_1,\ldots,\phi_m\in C^1({\cal M};\mathbb R)             \tag{21.1}
\]
be constraint functions and define
\[
 \Phi(q):=(\phi_1(q),\ldots,\phi_m(q))\in\mathbb R^m.        \tag{21.2}
\]
For \(\xi=(\xi^1,\ldots,\xi^m)\), put
\[
 G_\xi:=\sum_{b=1}^m\xi^b\phi_b,\qquad
 K_q\xi:=X_{G_\xi}(q),                                      \tag{21.3}
\]
using the convention
\[
 Df[X_g]=\{f,g\}.                                           \tag{21.4}
\]
Assume the brackets have the decomposition
\[
 \{\phi_a,\phi_b\}(q)
 =
 \sum_{c=1}^m f_{ab}^{\,c}(q)\phi_c(q)
 +{\mathfrak a}_{ab}(q).                                   \tag{21.5}
\]
The final term is the finite HDA anomaly.

Define
\[
 [A_q(\xi)y]_a
 :=
 \sum_{b,c}\xi^b f_{ab}^{\,c}(q)y_c,                        \tag{21.6}
\]
\[
 [\delta_{{\rm HDA},q}\xi]_a
 :=
 \sum_b\xi^b{\mathfrak a}_{ab}(q).                         \tag{21.7}
\]

\begin{theorem}[First-class bracket-to-equivariance identity]
\label{thm:v14v21-classical}
The derivative \(C_q:=D\Phi_q\) satisfies
\[
 \boxed{
 C_qK_q\xi
 =A_q(\xi)\Phi(q)+\delta_{{\rm HDA},q}\xi.}                 \tag{21.8}
\]
If the algebra is first class without anomaly,
\({\mathfrak a}_{ab}=0\), then
\[
 C_qK_q=A_q\Phi(q).                                         \tag{21.9}
\]
On the constraint surface,
\[
 \Phi(q)=0,\qquad {\mathfrak a}_{ab}(q)=0
 \quad\Longrightarrow\quad
 \boxed{C_qK_q=0.}                                         \tag{21.10}
\]
\end{theorem}

\begin{proof}
For each component,
\[
 \begin{aligned}
 [C_qK_q\xi]_a
 &=D\phi_a(q)[X_{G_\xi}(q)]\\
 &=\{\phi_a,G_\xi\}(q)\\
 &=\sum_b\xi^b\{\phi_a,\phi_b\}(q).
 \end{aligned}
\]
Insert (21.5), then use (21.6)--(21.7).  This proves (21.8).
The remaining statements are its two specializations.
\end{proof}

\begin{corollary}[First-class constraints satisfy the V19
equivariance gate]
\label{cor:v14v21-v19}
Let \(G_q>0\) be a metric on \(T_q{\cal M}\), and whiten
\[
 \widetilde C_q=C_qG_q^{-1/2},\qquad
 \widetilde K_q=G_q^{1/2}K_q.                              \tag{21.11}
\]
Then
\[
 \boxed{
 \widetilde C_q\widetilde K_q
 =A_q\Phi(q)+\delta_{{\rm HDA},q}.}                         \tag{21.12}
\]
Hence the differentiated equivariance defect in V19 can be taken to
be \(\delta_{{\rm HDA},q}\).
\end{corollary}

\begin{proof}
The middle metric cancels as in V19:
\(\widetilde C_q\widetilde K_q=C_qK_q\).  Apply (21.8).
\end{proof}

\subsection{Coordinate-free finite HDA form}

Let \(V\) be the finite lapse-plus-shift label space and write the
constraint paired with \(\xi\in V\) as
\[
 {\cal C}[\xi](q):=\langle\Phi(q),\xi\rangle.                \tag{21.13}
\]
Suppose
\[
 \{{\cal C}[\eta],{\cal C}[\xi]\}(q)
 =
 {\cal C}[F_q(\eta,\xi)](q)
 +a_q(\eta,\xi),                                           \tag{21.14}
\]
where \(F_q:V\times V\to V\) is bilinear and \(a_q\) is the anomaly
form.

\begin{proposition}[Coadjoint form]
\label{prop:v14v21-coadjoint}
Let \(K_q\xi=X_{{\cal C}[\xi]}(q)\).  Then
\[
 \boxed{
 D\Phi_q[K_q\xi]
 =F_q(\,\cdot\,,\xi)^*\Phi(q)+a_q(\,\cdot\,,\xi).}          \tag{21.15}
\]
\end{proposition}

\begin{proof}
Pair the left side with arbitrary \(\eta\in V\):
\[
 \begin{aligned}
 \langle D\Phi_q[K_q\xi],\eta\rangle
 &=D{\cal C}[\eta](q)[X_{{\cal C}[\xi]}]\\
 &=\{{\cal C}[\eta],{\cal C}[\xi]\}(q).
 \end{aligned}
\]
Insert (21.14), then move \(F_q(\eta,\xi)\) to the dual side.
\end{proof}

For the HDA, \(V\) is the direct sum of lapse and shift labels.
The \(DD,DH,HH\) brackets give the three blocks of \(F_q\).

\subsection{Finite density-state realization}

Let \({\cal H}\) be a finite-dimensional Hilbert space and let
\[
 {\sf C}_1,\ldots,{\sf C}_m={\sf C}_m^*
 \in{\cal B}({\cal H})                                     \tag{21.16}
\]
be self-adjoint constraint generators.  On the affine space of
density matrices, define
\[
 \Phi_a(\rho):=\operatorname{Tr}(\rho{\sf C}_a).            \tag{21.17}
\]
For \(\xi\in\mathbb R^m\), put
\[
 {\sf G}_\xi:=\sum_b\xi^b{\sf C}_b,\qquad
 \rho_t:=e^{it{\sf G}_\xi}\rho e^{-it{\sf G}_\xi},          \tag{21.18}
\]
and
\[
 K_\rho\xi
 :=\left.\frac{d}{dt}\right|_{0}\rho_t
 =i[{\sf G}_\xi,\rho].                                     \tag{21.19}
\]
Assume
\[
 i[{\sf C}_a,{\sf C}_b]
 =
 \sum_c f_{ab}^{\,c}{\sf C}_c+{\sf A}_{ab},                \tag{21.20}
\]
where \({\sf A}_{ab}={\sf A}_{ab}^*\) is the closure anomaly.

\begin{theorem}[Expectation-constraint tangency]
\label{thm:v14v21-density}
For the constraint map (21.17),
\[
 \boxed{
 [D\Phi_\rho(K_\rho\xi)]_a
 =
 \sum_{b,c}\xi^bf_{ab}^{\,c}\Phi_c(\rho)
 +\sum_b\xi^b\operatorname{Tr}(\rho{\sf A}_{ab}).}          \tag{21.21}
\]
If \({\sf A}_{ab}=0\) and \(\Phi(\rho)=0\), then
\[
 \boxed{D\Phi_\rho K_\rho=0.}                              \tag{21.22}
\]
\end{theorem}

\begin{proof}
By cyclicity of the finite trace,
\[
 \begin{aligned}
 D\Phi_a(\rho)[K_\rho\xi]
 &=\operatorname{Tr}\bigl(i[{\sf G}_\xi,\rho]{\sf C}_a\bigr)\\
 &=\operatorname{Tr}\bigl(\rho\,i[{\sf C}_a,{\sf G}_\xi]\bigr)\\
 &=\sum_b\xi^b
   \operatorname{Tr}\bigl(\rho\,i[{\sf C}_a,{\sf C}_b]\bigr).
 \end{aligned}
\]
Insert (21.20), which gives (21.21).  The exact zero branch follows.
\end{proof}

\begin{corollary}[Operator anomaly bound]
\label{cor:v14v21-anomaly}
For a density matrix \(\rho\), define
\[
 \alpha_{\rm op}
 :=
 \left(\sum_{a,b}\|{\sf A}_{ab}\|_{\rm op}^2\right)^{1/2}.
                                                               \tag{21.23}
\]
With Euclidean norms on the finite constraint labels,
\[
 \left\|
 \left(
  \sum_b\xi^b\operatorname{Tr}(\rho{\sf A}_{ab})
 \right)_a
 \right\|
 \leq\alpha_{\rm op}\|\xi\|.                               \tag{21.24}
\]
\end{corollary}

\begin{proof}
Since \(\rho\geq0\) and \(\operatorname{Tr}\rho=1\),
\[
 |\operatorname{Tr}(\rho{\sf A}_{ab})|
 \leq\|{\sf A}_{ab}\|_{\rm op}.
\]
Apply Cauchy--Schwarz first in \(b\), then sum in \(a\).
\end{proof}

A central term
\({\sf A}_{ab}=z_{ab}I\) remains visible in (21.21); it cannot be
removed by imposing zero expectation of the noncentral constraints.

\begin{firewall}[Expectation constraints are a declared branch]
The condition \(\operatorname{Tr}(\rho{\sf C}_a)=0\) for every \(a\)
does not imply \({\sf C}_a\psi=0\) on a state vector, nor does it
construct the BRST cohomology.  The density-state theorem applies
only when expectation constraints are the declared finite physical
constraint map.  A Dirac or BRST realization needs its own carrier
and quotient.
\end{firewall}

\subsection{Application to the finite HDA identities}

At a fixed cutoff \(h\), choose bases
\[
 N_1,\ldots,N_p\quad\hbox{and}\quad
 v_1,\ldots,v_q                                             \tag{21.25}
\]
for the retained lapse and tangential label spaces, and form the
constraint-generator list
\[
 {\sf C}_{h,A}
 =
 \bigl(
  H_h[N_1],\ldots,H_h[N_p],
  D_h[v_1],\ldots,D_h[v_q]
 \bigr).                                                    \tag{21.26}
\]
The Grand Tensor finite HDA gives
\[
\begin{aligned}
 -i[D_h[v],D_h[w]]&=D_h[[v,w]_h],\\
 -i[D_h[v],H_h[N]]&=H_h[\ell_h(v)N],\\
 -i[H_h[N],H_h[M]]
 &=D_h[G_h^{-1}(N\,d_hM-M\,d_hN)].                         \tag{21.27}
\end{aligned}
\]

\begin{theorem}[Finite HDA tangency compiler]
\label{thm:v14v21-hda}
Assume:

\begin{enumerate}
\item every label on the right side of (21.27) lies in the retained
lapse-plus-shift space of (21.25);
\item all generators act on the same reducing finite carrier;
\item the signs in (21.26) are fixed once so that (21.27) has the
commutator convention (21.20);
\item any interface, Feshbach, or finite closure residual is collected
in an explicit anomaly bank \({\sf A}_{h,AB}\).
\end{enumerate}

Then the expectation constraint map
\[
 \Phi_{h,A}(\rho)
 :=\operatorname{Tr}(\rho{\sf C}_{h,A})                    \tag{21.28}
\]
obeys
\[
 \boxed{
 D\Phi_{h,\rho}K_{h,\rho}
 =
 A_{h,\rho}\Phi_h(\rho)
 +\delta_{{\rm HDA},h}(\rho),}                             \tag{21.29}
\]
where the last term consists exactly of the anomaly expectations.
If the exact HDA (21.27) holds without residual and
\(\Phi_h(\rho)=0\), then
\[
 \boxed{D\Phi_{h,\rho}K_{h,\rho}=0.}                       \tag{21.30}
\]
After common-metric whitening, (21.30) is the exact V17 nesting
condition.
\end{theorem}

\begin{proof}
Expand every right side of (21.27) in the bases (21.25).  This gives
the structure coefficients in (21.20), with all unabsorbed closure
terms placed in \({\sf A}_{h,AB}\).  Apply
Theorem~\ref{thm:v14v21-density}.  Metric whitening cancels in the
middle product by V19.
\end{proof}

\begin{assessment}[What the Grand Tensor HDA discharges]
On the literal branch of its finite-HDA theorem, the Grand Tensor
source asserts the same reducing carrier and vanishing interface or
Feshbach terms.  Once the retained label spaces are finite and closed
under the displayed brackets, the theorem above discharges the
algebraic tangency row for the expectation-constraint realization.
It does not select that realization as the physical ADM constraint
map, provide its matrices, or prove uniform singular floors.
\end{assessment}

\subsection{Quantitative anomalous HDA branch}

Suppose, in either the classical or density setting,
\[
 \|A_q(\xi)y\|
 \leq F_q\|\xi\|\,\|y\|,\qquad
 \|\delta_{{\rm HDA},q}\|\leq a_q.                         \tag{21.31}
\]
Then
\[
 \|C_qK_q\|
 \leq F_q\|\Phi(q)\|+a_q.                                  \tag{21.32}
\]

\begin{corollary}[HDA-to-angle estimate]
\label{cor:v14v21-angle}
If
\[
 \|C_q^\dagger\|\,\|K_q^\dagger\|
 \bigl(F_q\|\Phi(q)\|+a_q\bigr)
 \leq\zeta_q<1,                                            \tag{21.33}
\]
then
\[
 \boxed{\gamma(P_q,S_q)\geq1-\zeta_q^2}                   \tag{21.34}
\]
and
\[
 \boxed{\|{\mathfrak C}_q^\dagger\|
 \leq(1-\zeta_q)^{-1/2}.}                                  \tag{21.35}
\]
\end{corollary}

\begin{proof}
Equation (21.32) is the norm form of (21.8) or (21.21).
Apply the V18 propagation theorem.
\end{proof}

This identifies a finite HDA anomaly as a quantitative angle debit,
rather than merely a formal failure of covariance.

\subsection{Exact witnesses}

\begin{proposition}[Three first-class failure witnesses]
\label{prop:v14v21-witness}
Each of the following is an exact obstruction to the declared HDA
constraint card:

\begin{enumerate}
\item a bracket component outside the span of the retained generators;
\item a nonzero anomaly expectation on a zero-constraint density
state;
\item a vector \(\xi\) for which
\(D\Phi_qK_q\xi\ne0\) at \(\Phi(q)=0\).
\end{enumerate}

The second and third witnesses refute tangency.  The first refutes
closure of the chosen finite constraint bank but may be repaired by
enlarging that bank.
\end{proposition}

\begin{proof}
The first contradicts the first-class form (21.5) or (21.20).
At \(\Phi=0\), equations (21.8) and (21.21) reduce the tangent defect
to the anomaly.  This proves the remaining assertions.
\end{proof}

\subsection{Next upstream conditioning step}

On a symplectic regulator the gauge map is not independent of the
constraint derivative.  If the Poisson tensor is represented by a
compatible skew-adjoint isometry \({\mathbb J}\), then
\[
 K_q={\mathbb J}C_q^*                                      \tag{21.36}
\]
up to the fixed sign convention.  V22 should prove:

\begin{enumerate}
\item equality of the nonzero singular values of \(K_q\) and \(C_q\);
\item the coisotropy identity
\(C_q{\mathbb J}C_q^*=0\) on the first-class constraint surface;
\item reduction of the two conditioning factors in V18--V21 to one
constraint Gram \(C_qC_q^*\);
\item a kernel-completed positive Gram criterion which detects
reducibility among the finite HDA constraints.
\end{enumerate}

This attacks the remaining singular-floor row without introducing an
unrelated Einstein symbol.

\subsection{Acquisition ledger}

\[
\begin{array}{p{0.29\linewidth}|p{0.23\linewidth}|p{0.38\linewidth}}
\textbf{gate}&\textbf{status after V21}&\textbf{content}\\ \hline
\text{classical HDA tangency}
 &\text{acquired}
 &\{\phi_a,\phi_b\}=f_{ab}^{c}\phi_c+\mathfrak a_{ab}\\
\text{density-state HDA tangency}
 &\text{acquired}
 &expectation constraint branch\\
\text{anomaly-to-angle debit}
 &\text{acquired}
 &\zeta=\|C^\dagger\|\|K^\dagger\|(F\|\Phi\|+a)\\
\text{Grand Tensor algebraic bridge}
 &\text{acquired conditionally}
 &finite label closure on one carrier\\
\text{physical constraint realization}
 &\text{open declaration}
 &Poisson, expectation, Dirac, or BRST branch\\
\text{uniform constraint conditioning}
 &\text{open}
 &least positive eigenvalue of the constraint Gram\\
\text{three Einstein residual defects}
 &\text{open}
 &same-history source values
\end{array}                                                  \tag{21.37}
\]

\subsection{Parent-integrity record}

The installed parent was
\[
\begin{array}{c|c}
 \text{file}&
 \texttt{Renewal\_SM\_Einstein\_Continuum\_Research\_Notes\_V20.tex}\\
 \text{lines}&8103\\
 \text{bytes}&260979\\
 \text{SHA--256}&
 \texttt{6904A162AB1929A607A8C8C034AE976CE912B7FA6C470F2C3B78BEB39F262202}.
\end{array}                                                  \tag{21.38}
\]

% =============================================================================
% END APPENDED FOURTEENTH-CYCLE V21 MATERIAL
% =============================================================================
% =============================================================================
% BEGIN APPENDED FOURTEENTH-CYCLE V22 MATERIAL
% =============================================================================

\section{Fourteenth-cycle V22: symplectic constraint Grams, coisotropy,
and a single conditioning floor}
\label{sec:v14v22}

\subsection{Compatible finite Hamiltonian screen}

Let \(Z\) and \(Y\) be finite-dimensional Hilbert spaces.  Assume
\(Z\) carries a compatible complex structure
\[
 {\mathbb J}^*=-{\mathbb J},\qquad
 {\mathbb J}^*{\mathbb J}=I,\qquad
 {\mathbb J}^2=-I.                                         \tag{22.1}
\]
Let
\[
 C:Z\longrightarrow Y                                      \tag{22.2}
\]
be the derivative of the finite constraint map.  The Hamiltonian
gauge tangent generated by a multiplier \(\lambda\in Y\) is
\[
 K:={\mathbb J}C^*:Y\longrightarrow Z.                     \tag{22.3}
\]
A fixed sign change in (22.3) has no effect on any projection or
singular-value statement below.

\subsection{One singular spectrum controls both maps}

\begin{theorem}[Constraint--gauge singular-value identity]
\label{thm:v14v22-singular}
The maps \(C\), \(C^*\), and \(K\) have the same nonzero singular
values.  More precisely,
\[
 K^*K=CC^*,\qquad
 \ker K=\ker C^*,\qquad
 \operatorname{Ran}K={\mathbb J}\operatorname{Ran}C^*.      \tag{22.4}
\]
Consequently,
\[
 \boxed{
 \operatorname{rmin}K=\operatorname{rmin}C,}
 \qquad
 \boxed{\|K^\dagger\|=\|C^\dagger\|}                        \tag{22.5}
\]
whenever the nonzero range is present, and
\[
 \|K\|=\|C\|.                                               \tag{22.6}
\]
\end{theorem}

\begin{proof}
Using (22.1),
\[
 K^*K
 =C{\mathbb J}^*{\mathbb J}C^*
 =CC^*.
\]
This proves equality of the squared nonzero singular values and the
kernel identity.  The range identity follows from (22.3) and
unitarity of \({\mathbb J}\).  The reduced minimum moduli, pseudoinverse
norms, and operator norms are determined by the nonzero singular
values, proving (22.5)--(22.6).
\end{proof}

Thus the two conditioning factors in V18 are not independent on a
compatible Hamiltonian branch.

\subsection{The Poisson-bracket matrix is the propagation product}

Define
\[
 {\mathbb B}:=CK=C{\mathbb J}C^*:Y\longrightarrow Y.        \tag{22.7}
\]
It is skew-adjoint:
\[
 {\mathbb B}^*=-{\mathbb B}.                                \tag{22.8}
\]

\begin{proposition}[Linearized bracket identity]
\label{prop:v14v22-bracket}
If the components of the nonlinear constraint map are
\(\phi_1,\ldots,\phi_m\) in an orthonormal basis of \(Y\), then the
matrix entries of \({\mathbb B}\), up to the fixed Hamiltonian sign
convention, are
\[
 {\mathbb B}_{ab}=\{\phi_a,\phi_b\}.                        \tag{22.9}
\]
If the HDA is first class,
\[
 \{\phi_a,\phi_b\}
 =\sum_cf_{ab}^{\,c}\phi_c+{\mathfrak a}_{ab},              \tag{22.10}
\]
then
\[
 \boxed{
 {\mathbb B}=F(\Phi)+{\mathfrak A}_{\rm HDA}.}              \tag{22.11}
\]
\end{proposition}

\begin{proof}
The gradient of \(\phi_b\) is \(C^*e_b\), and its Hamiltonian vector
is \({\mathbb J}C^*e_b\).  Applying \(D\phi_a=e_a^*C\) gives
\(e_a^*C{\mathbb J}C^*e_b\), which is (22.9) under the convention of
V21.  Equation (22.11) is the matrix form of (22.10).
Skew-adjointness follows from (22.1).
\end{proof}

\begin{theorem}[Coisotropy is exact gauge tangency]
\label{thm:v14v22-coisotropic}
The following are equivalent:
\[
 CK=0,                                                      \tag{22.12}
\]
\[
 C{\mathbb J}C^*=0,                                        \tag{22.13}
\]
\[
 {\mathbb J}\operatorname{Ran}C^*\subseteq\ker C,           \tag{22.14}
\]
\[
 (\ker C)^{\perp_{\omega}}\subseteq\ker C,                  \tag{22.15}
\]
where
\[
 \omega(u,v):=\langle-{\mathbb J}u,v\rangle
\]
and \(\perp_\omega\) denotes symplectic orthogonality.  Thus exact
first-class closure on the zero-constraint surface makes the
linearized constraint surface coisotropic.
\end{theorem}

\begin{proof}
Equations (22.12)--(22.14) are the same statement after inserting
\(K={\mathbb J}C^*\).  Since
\[
 (\ker C)^\perp=\operatorname{Ran}C^*
\]
in finite dimension, its symplectic orthogonal is
\({\mathbb J}\operatorname{Ran}C^*\), up to the irrelevant sign.
This gives (22.15).
\end{proof}

\subsection{Constraint Gram and reducibility}

The multiplier-space constraint Gram is
\[
 G_C:=CC^*\succeq0.                                        \tag{22.16}
\]
Its kernel
\[
 {\cal R}_C:=\ker C^*                                      \tag{22.17}
\]
is the space of linear relations among the chosen constraint
components.  Let \(\Pi_C\) be an explicitly declared orthogonal
projection with
\[
 \operatorname{Ran}\Pi_C\subseteq\ker C^*.                 \tag{22.18}
\]
Define the completed constraint Gram
\[
 \widehat G_C:=CC^*+\Pi_C.                                  \tag{22.19}
\]

\begin{theorem}[Completed constraint Gram proves both completeness
and conditioning]
\label{thm:v14v22-completed}
If
\[
 \widehat G_C\succeq\mu I_Y                                \tag{22.20}
\]
for some \(\mu>0\), then
\[
 \boxed{\ker C^*=\operatorname{Ran}\Pi_C}                  \tag{22.21}
\]
and
\[
 \boxed{
 \operatorname{rmin}C
 =\operatorname{rmin}K
 \geq\sqrt{\mu},}
 \qquad
 \boxed{
 \|C^\dagger\|=\|K^\dagger\|
 \leq\mu^{-1/2}.}                                          \tag{22.22}
\]
If \(\operatorname{Ran}\Pi_C\ne\{0\}\), then necessarily
\(\mu\leq1\).
\end{theorem}

\begin{proof}
Apply the V16 kernel-completion theorem to \(J=C^*\) and
\(\Pi=\Pi_C\).  Its completed normal operator is
\[
 J^*J+\Pi_C=CC^*+\Pi_C.
\]
This gives (22.21) and the reduced floor for \(C^*\), hence for
\(C\).  Theorem~\ref{thm:v14v22-singular} transfers it to \(K\).
The pseudoinverse estimates and the final normalization statement
also follow from V16.
\end{proof}

A declared relation projector built from already known lapse/shift
dependencies is therefore checked, rather than assumed, by the same
positive Gram which supplies the regulator conditioning.

\subsection{One ratio controls the physical--tangent angle}

Let
\[
 b:=\|C{\mathbb J}C^*\|.                                   \tag{22.23}
\]

\begin{theorem}[Bracket-to-angle certificate]
\label{thm:v14v22-angle}
Assume (22.20) and
\[
 \theta:=\frac b\mu<1.                                     \tag{22.24}
\]
Let
\[
 P:=I-C^\dagger C,\qquad
 T:=KK^\dagger,\qquad S:=I-T.                              \tag{22.25}
\]
Then
\[
 \boxed{\|(I-P)T\|\leq\theta,}                             \tag{22.26}
\]
\[
 \boxed{\gamma(P,S)\geq1-\theta^2,}                        \tag{22.27}
\]
and the V14 joint constraint satisfies
\[
 \boxed{
 \|{\mathfrak C}^\dagger\|
 \leq(1-\theta)^{-1/2}.}                                   \tag{22.28}
\]
If \(C{\mathbb J}C^*=0\), then
\[
 R=P-T,\qquad\gamma=1,\qquad
 {\mathfrak C}^\dagger={\mathfrak C}^*.                    \tag{22.29}
\]
\end{theorem}

\begin{proof}
The V18 propagation condition number is
\[
 \eta
 =\|C^\dagger\|\,\|CK\|\,\|K^\dagger\|.
\]
By (22.7) and (22.22),
\[
 \eta\leq\mu^{-1}b=\theta.
\]
Apply the V18 angle and joint-pseudoinverse theorems.  The exact branch
is V17.
\end{proof}

Combining (22.11) with (22.23), a concrete bound is
\[
 b
 \leq
 \|F(\Phi)\|+\|{\mathfrak A}_{\rm HDA}\|.                  \tag{22.30}
\]
Thus the completed constraint-Gram floor must dominate the off-shell
first-class and anomaly matrix.

\subsection{Physical dimension count}

Let
\[
 n:=\dim Z,\qquad r:=\operatorname{rank}C.                  \tag{22.31}
\]

\begin{theorem}[Finite first-class degree count]
\label{thm:v14v22-dimension}
If \(C{\mathbb J}C^*=0\), then
\[
 2r\leq n.                                                  \tag{22.32}
\]
The gauge-reduced physical variation space
\[
 {\cal M}_{\rm phys}
 :=\ker C\cap(\operatorname{Ran}K)^\perp                   \tag{22.33}
\]
has dimension
\[
 \boxed{\dim{\cal M}_{\rm phys}=n-2r.}                     \tag{22.34}
\]
\end{theorem}

\begin{proof}
Theorem~\ref{thm:v14v22-coisotropic} gives
\(\operatorname{Ran}K\subseteq\ker C\).  Theorem
\ref{thm:v14v22-singular} gives
\(\dim\operatorname{Ran}K=r\), while
\(\dim\ker C=n-r\).  Hence \(r\leq n-r\), proving (22.32).
The orthogonal decomposition
\[
 \ker C
 =\operatorname{Ran}K\oplus{\cal M}_{\rm phys}
\]
then gives
\[
 \dim{\cal M}_{\rm phys}
 =(n-r)-r=n-2r.
\]
\end{proof}

This count uses the rank of the constraint derivative, not the raw
number of constraint labels.  Reducibility is removed by the kernel
projection \(\Pi_C\).

\subsection{Two exact examples}

\begin{example}[One first-class constraint]
Let \(Z=\mathbb R^2\) with
\[
 {\mathbb J}=
 \begin{pmatrix}0&-1\\1&0\end{pmatrix},
 \qquad C=(0\ \ 1).
\]
Then
\[
 C{\mathbb J}C^*=0,\qquad
 CC^*=1,\qquad
 K={\mathbb J}C^*=(-1,0)^{\mathsf T}.                      \tag{22.35}
\]
Thus \(r=1\), \(n-2r=0\), and the constraint surface consists only of
one gauge direction at the linearized level.
\end{example}

\begin{example}[A second-class pair]
Let \(Z=\mathbb R^2\) with the same \({\mathbb J}\), and take
\(C=I_2\).  Then
\[
 C{\mathbb J}C^*={\mathbb J}\ne0.                          \tag{22.36}
\]
Here \(r=2\) violates \(2r\leq n\), and the gauge range cannot lie in
\(\ker C=\{0\}\).  The bracket matrix is an exact witness that this
pair is second class rather than a first-class gauge system.
\end{example}

\begin{example}[Reducible labels]
With the first example, duplicate the constraint:
\[
 C=
 \begin{pmatrix}0&1\\0&1\end{pmatrix}.
\]
Then
\[
 CC^*=
 \begin{pmatrix}1&1\\1&1\end{pmatrix},
\quad
 \ker C^*=\operatorname{span}(1,-1).                       \tag{22.37}
\]
The raw Gram is singular because the labels are redundant.  Adding
the orthogonal projector onto \(\operatorname{span}(1,-1)\) gives a
strictly positive completed Gram and proves that this is the complete
relation space.
\end{example}

\subsection{Specialized three-defect Einstein estimate}

\begin{theorem}[Symplectic three-defect estimate]
\label{thm:v14v22-three}
Let \(R\) be the physical quotient projection and let \(r_{\rm Ein}\)
be a whitened finite Einstein residual.  Under (22.20)--(22.24), V18
becomes
\[
\boxed{
\begin{aligned}
 \|r_{\rm Ein}\|^2
 \leq{}&\|Rr_{\rm Ein}\|^2\\
 &+
 \frac{
 \mu^{-1}\bigl(
  \|Cr_{\rm Ein}\|^2+\|K^*r_{\rm Ein}\|^2
 \bigr)}
 {1-b/\mu}.
\end{aligned}}                                              \tag{22.38}
\]
\end{theorem}

\begin{proof}
By (22.22),
\[
 \|C^\dagger\|^2,\ \|K^\dagger\|^2\leq\mu^{-1}.
\]
Insert these and \(\eta\leq b/\mu\) into the V18 estimate.
\end{proof}

\begin{corollary}[Exact first-class closure]
\label{cor:v14v22-residual}
If \(b=0\), then
\[
 \|r_{\rm Ein}\|^2
 \leq
 \|Rr_{\rm Ein}\|^2
 +\mu^{-1}
 \bigl(\|Cr_{\rm Ein}\|^2+\|K^*r_{\rm Ein}\|^2\bigr).       \tag{22.39}
\]
Zero physical quotient, constraint, and Ward residuals imply the full
whitened residual is zero.
\end{corollary}

\begin{proof}
Set \(b=0\) in (22.38).
\end{proof}

\subsection{Cofinal stability from one completed Gram}

Let \(C_n:Z\to Y\) and compatible
\({\mathbb J}_n:Z\to Z\) act on fixed transported finite screens.
Let \(\Pi_n\) be the declared relation projections.  Assume
\[
 C_n\to C,\qquad
 {\mathbb J}_n\to{\mathbb J},\qquad
 \Pi_n\to\Pi                                             \tag{22.40}
\]
in norm and
\[
 C_nC_n^*+\Pi_n\succeq\mu I                                \tag{22.41}
\]
with a common \(\mu>0\).

\begin{theorem}[Constraint-Gram landing]
\label{thm:v14v22-cofinal}
Under (22.40)--(22.41),
\[
 C_n^\dagger\to C^\dagger,\qquad
 K_n^\dagger\to K^\dagger,                                 \tag{22.42}
\]
\[
 P_n\to P,\qquad T_n\to T.                                 \tag{22.43}
\]
If additionally
\[
 \|C_n{\mathbb J}_nC_n^*\|\leq\theta\mu,\qquad\theta<1,     \tag{22.44}
\]
then
\[
 R_n\to R,\qquad
 {\mathfrak C}_n^\dagger\to{\mathfrak C}^\dagger.           \tag{22.45}
\]
\end{theorem}

\begin{proof}
The completed-Gram floor separates zero from the positive spectrum of
\(C_nC_n^*\).  Continuous functional calculus, using a continuous
extension of \(t\mapsto t^{-1}\) from
\(\{0\}\cup[\mu,\infty)\), gives
\[
 (C_nC_n^*)^\dagger\to(CC^*)^\dagger.
\]
Since
\[
 C_n^\dagger=C_n^*(C_nC_n^*)^\dagger,
\]
the first convergence in (22.42) follows.  The identity
\[
 K_n^\dagger=(C_n^\dagger)^*{\mathbb J}_n^*
\]
gives the second.  The projection formulas
\[
 P_n=I-C_n^\dagger C_n,\qquad
 T_n={\mathbb J}_n(C_n^\dagger C_n){\mathbb J}_n^*
\]
give (22.43).  Finally, (22.44) and
Theorem~\ref{thm:v14v22-angle} give a common V12 angle gap and V14
joint floor, proving (22.45).
\end{proof}

This removes the need to assume projector convergence separately on
the compatible symplectic branch.

\subsection{Next coupled SM--Einstein step}

A full SM--Einstein regulator contains gravitational HDA constraints
and internal Gauss/BRST constraints on one common action carrier.
Their combined derivative need not be block diagonal after matter and
geometry are coupled.  V23 should study a block constraint map
\[
 C_{\rm tot}
 =
 \begin{pmatrix}
 C_{\rm grav}&M\\
 N&C_{\rm int}
 \end{pmatrix}                                             \tag{22.46}
\]
and prove a Schur-complement or graph-subspace floor for
\(C_{\rm tot}C_{\rm tot}^*\) from:

\begin{enumerate}
\item completed floors in the gravity and internal sectors;
\item exact relation projectors;
\item a quantitative mixed-incidence bound for \(M,N\);
\item the common symplectic first-class bracket.
\end{enumerate}

That is the first point in this cycle where the internal Standard
Model constraint bank and the Einstein HDA conditioning are assembled
rather than treated in parallel.

\subsection{Acquisition ledger}

\[
\begin{array}{p{0.29\linewidth}|p{0.23\linewidth}|p{0.38\linewidth}}
\textbf{gate}&\textbf{status after V22}&\textbf{content}\\ \hline
\text{constraint/gauge conditioning}
 &\text{unified}
 &\operatorname{rmin}K=\operatorname{rmin}C\\
\text{first-class propagation product}
 &\text{identified}
 &CK=C{\mathbb J}C^*\\
\text{reducibility completion}
 &\text{acquired}
 &CC^*+\Pi_C\succeq\mu I\\
\text{angle certificate}
 &\text{acquired}
 &b/\mu<1\\
\text{finite physical dimension}
 &\text{acquired}
 &n-2\operatorname{rank}C\\
\text{cofinal projector landing}
 &\text{acquired conditionally}
 &one uniform completed-Gram floor\\
\text{coupled SM--gravity constraint Gram}
 &\text{next}
 &mixed block floor\\
\text{concrete regulator values}
 &\text{open}
 &\mu,b,\text{ and residual channels}
\end{array}                                                  \tag{22.47}
\]

\subsection{Parent-integrity record}

The installed parent was
\[
\begin{array}{c|c}
 \text{file}&
 \texttt{Renewal\_SM\_Einstein\_Continuum\_Research\_Notes\_V21.tex}\\
 \text{lines}&8590\\
 \text{bytes}&275544\\
 \text{SHA--256}&
 \texttt{FB20AA24235F114D5E2CFC34E949D97BC80B497C98A0C94C3EEDC2DCF2DCA466}.
\end{array}                                                  \tag{22.48}
\]

% =============================================================================
% END APPENDED FOURTEENTH-CYCLE V22 MATERIAL
% =============================================================================

% =============================================================================
% BEGIN APPENDED FOURTEENTH-CYCLE V23 MATERIAL
% =============================================================================

\section{Fourteenth-cycle V23: coupled Standard-Model--Einstein
constraint Grams and the mixed-incidence Schur gate}
\label{sec:v14v23}

\subsection{Coupled constraint derivative}

Let
\[
 Z=Z_g\oplus Z_i,\qquad Y=Y_g\oplus Y_i                 \tag{23.1}
\]
be the finite phase-variation and constraint-label spaces for the
gravity and internal sectors.  Write the derivative of the combined
constraint map as
\[
 \boxed{
 C_{\rm tot}
 =
 \begin{pmatrix}
 C_g&M\\
 N&C_i
 \end{pmatrix}:Z_g\oplus Z_i\longrightarrow Y_g\oplus Y_i.}
                                                               \tag{23.2}
\]
Here
\[
 C_g:Z_g\to Y_g,\quad C_i:Z_i\to Y_i,\quad
 M:Z_i\to Y_g,\quad N:Z_g\to Y_i.                          \tag{23.3}
\]
The blocks \(M,N\) contain the linearized mixed matter--geometry
incidence.  They need not vanish for a common interacting action.

Let \(\Pi_g\) and \(\Pi_i\) be declared orthogonal projectors onto
known relation spaces among the gravity and internal constraint
labels.  Assume
\[
 C_g^*\Pi_g=0,\qquad M^*\Pi_g=0,                            \tag{23.4}
\]
\[
 C_i^*\Pi_i=0,\qquad N^*\Pi_i=0.                            \tag{23.5}
\]
Then
\[
 \Pi_0:=\Pi_g\oplus\Pi_i                                   \tag{23.6}
\]
has
\[
 C_{\rm tot}^*\Pi_0=0.                                     \tag{23.7}
\]
Conditions (23.4)--(23.5) are mixed-relation incidence rows; a
relation in an isolated sector need not remain a relation after
coupling unless these rows hold.

\subsection{The exact completed Gram}

Direct multiplication gives
\[
 C_{\rm tot}C_{\rm tot}^*
 =
 \begin{pmatrix}
 C_gC_g^*+MM^*&C_gN^*+MC_i^*\\
 NC_g^*+C_iM^*&NN^*+C_iC_i^*
 \end{pmatrix}.                                             \tag{23.8}
\]
Define
\[
 A:=C_gC_g^*+MM^*+\Pi_g,                                   \tag{23.9}
\]
\[
 D:=NN^*+C_iC_i^*+\Pi_i,                                   \tag{23.10}
\]
\[
 X:=C_gN^*+MC_i^*.                                         \tag{23.11}
\]
The completed coupled Gram is therefore
\[
 \boxed{
 \widehat G_{\rm tot}
 :=C_{\rm tot}C_{\rm tot}^*+\Pi_0
 =
 \begin{pmatrix}A&X\\X^*&D\end{pmatrix}.}                  \tag{23.12}
\]

Assume sector floors
\[
 A\succeq\mu_gI_{Y_g},\qquad
 D\succeq\mu_iI_{Y_i},\qquad
 \mu_g,\mu_i>0.                                            \tag{23.13}
\]
These may be proved directly or inherited from
\[
 C_gC_g^*+\Pi_g\succeq\mu_gI,\qquad
 C_iC_i^*+\Pi_i\succeq\mu_iI,                              \tag{23.14}
\]
because \(MM^*,NN^*\succeq0\).

\subsection{Two mixed-incidence floor certificates}

\begin{theorem}[Scalar cross-block floor]
\label{thm:v14v23-scalar}
Suppose
\[
 \|X\|\leq\rho,\qquad \rho^2<\mu_g\mu_i.                   \tag{23.15}
\]
Then
\[
 \widehat G_{\rm tot}\succeq\delta_{\rm sc}I_Y,             \tag{23.16}
\]
where
\[
 \boxed{
 \delta_{\rm sc}
 =
 \frac{\mu_g+\mu_i-
 \sqrt{(\mu_g-\mu_i)^2+4\rho^2}}{2}>0.}                    \tag{23.17}
\]
\end{theorem}

\begin{proof}
For \(u\in Y_g\), \(v\in Y_i\),
\[
 \begin{aligned}
 \left\langle
 \widehat G_{\rm tot}\binom uv,\binom uv
 \right\rangle
 &=
 \langle Au,u\rangle+\langle Dv,v\rangle
 +2\operatorname{Re}\langle Xv,u\rangle\\
 &\geq
 \mu_g\|u\|^2+\mu_i\|v\|^2-2\rho\|u\|\|v\|.
 \end{aligned}                                              \tag{23.18}
\]
The smallest eigenvalue of the scalar matrix
\[
 \begin{pmatrix}\mu_g&-\rho\\-\rho&\mu_i\end{pmatrix}
\]
is (23.17), positive exactly when (23.15) holds.  This proves
(23.16).
\end{proof}

The scalar bound may be wasteful because it ignores the anisotropy in
\(A,D\).  The normalized mixed incidence is
\[
 Y_{\rm mix}:=A^{-1/2}XD^{-1/2},\qquad
 \chi:=\|Y_{\rm mix}\|.                                    \tag{23.19}
\]

\begin{theorem}[Normalized Schur contraction]
\label{thm:v14v23-schur}
If
\[
 \chi<1,                                                    \tag{23.20}
\]
then
\[
 \boxed{
 \widehat G_{\rm tot}
 \succeq
 (1-\chi)
 \begin{pmatrix}A&0\\0&D\end{pmatrix}
 \succeq
 \delta_{\rm mix}I_Y,}                                     \tag{23.21}
\]
where
\[
 \boxed{
 \delta_{\rm mix}:=(1-\chi)\min\{\mu_g,\mu_i\}>0.}          \tag{23.22}
\]
Equivalently, the two Schur complements are strictly positive:
\[
 A-XD^{-1}X^*>0,\qquad
 D-X^*A^{-1}X>0.                                           \tag{23.23}
\]
\end{theorem}

\begin{proof}
Congruence by
\(\operatorname{diag}(A^{-1/2},D^{-1/2})\) transforms
\(\widehat G_{\rm tot}\) into
\[
 \begin{pmatrix}I&Y_{\rm mix}\\Y_{\rm mix}^*&I\end{pmatrix}.
                                                               \tag{23.24}
\]
For arbitrary \(a,b\),
\[
 2|\langle Y_{\rm mix}b,a\rangle|
 \leq2\chi\|a\|\|b\|
 \leq\chi(\|a\|^2+\|b\|^2).
\]
Hence (23.24) is bounded below by
\((1-\chi)I\).  Undoing the congruence gives (23.21)--(23.22).
The standard block Gaussian factorization gives
\[
 \begin{pmatrix}A&X\\X^*&D\end{pmatrix}
 =
 \begin{pmatrix}I&XD^{-1}\\0&I\end{pmatrix}
 \begin{pmatrix}A-XD^{-1}X^*&0\\0&D\end{pmatrix}
 \begin{pmatrix}I&0\\D^{-1}X^*&I\end{pmatrix},
\]
and the analogous factorization with \(A\) first.  Condition
\(\chi<1\) makes both Schur complements strictly positive, proving
(23.23).
\end{proof}

The normalized condition can succeed even when the coarse scalar
estimate (23.15) is inconclusive.

\subsection{Completeness of the coupled relation bank}

\begin{theorem}[Coupled kernel completion]
\label{thm:v14v23-completion}
Assume (23.4)--(23.5) and either
Theorem~\ref{thm:v14v23-scalar} or
Theorem~\ref{thm:v14v23-schur}.  Let \(\delta>0\) denote the
corresponding floor.  Then
\[
 \boxed{
 \ker C_{\rm tot}^*=\operatorname{Ran}\Pi_0}               \tag{23.25}
\]
and
\[
 \boxed{
 \operatorname{rmin}C_{\rm tot}\geq\sqrt{\delta},\qquad
 \|C_{\rm tot}^\dagger\|\leq\delta^{-1/2}.}                 \tag{23.26}
\]
\end{theorem}

\begin{proof}
Equation (23.7) places the declared relation range in
\(\ker C_{\rm tot}^*\).  The completed operator in (23.12) is exactly
\[
 C_{\rm tot}C_{\rm tot}^*+\Pi_0.
\]
Apply the V16 kernel-completion theorem with
\(J=C_{\rm tot}^*\).  Its positive floor proves both (23.25) and
(23.26).
\end{proof}

Thus a strict Schur contraction simultaneously proves that coupling
has introduced no hidden linear relation among the full constraint
bank.

\subsection{Separate floors do not suffice}

\begin{countertheorem}[Mixed coupling can create a new relation]
\label{cth:v14v23-newrelation}
Let all four sector spaces be one-dimensional and take
\[
 C_g=C_i=M=N=1.
\]
Then each isolated sector has constraint Gram \(1\), but
\[
 C_{\rm tot}
 =
 \begin{pmatrix}1&1\\1&1\end{pmatrix}                      \tag{23.27}
\]
has the new relation
\[
 (1,-1)\in\ker C_{\rm tot}^*.                              \tag{23.28}
\]
For this example,
\[
 A=D=2,\qquad X=2,\qquad\chi=1,                            \tag{23.29}
\]
so the strict mixed-incidence gate fails exactly.
\end{countertheorem}

\begin{proof}
The matrix in (23.27) has rank one and annihilates
\((1,-1)\).  Equations (23.29) follow from
(23.9)--(23.11).
\end{proof}

This witness shows why gravity and internal sector floors cannot be
combined by taking their minimum while ignoring the mixed incidence.

\subsection{Symplectic propagation of the coupled floor}

Let \(Z\) carry a compatible complex structure
\({\mathbb J}_{\rm tot}\), not necessarily block diagonal, and define
\[
 K_{\rm tot}:={\mathbb J}_{\rm tot}C_{\rm tot}^*.           \tag{23.30}
\]
By V22,
\[
 \operatorname{rmin}K_{\rm tot}
 =\operatorname{rmin}C_{\rm tot},\qquad
 \|K_{\rm tot}^\dagger\|
 =\|C_{\rm tot}^\dagger\|.                                 \tag{23.31}
\]
The combined first-class bracket matrix is
\[
 {\mathbb B}_{\rm tot}
 :=C_{\rm tot}{\mathbb J}_{\rm tot}C_{\rm tot}^*.           \tag{23.32}
\]

\begin{theorem}[Coupled SM--Einstein angle certificate]
\label{thm:v14v23-angle}
Under Theorem~\ref{thm:v14v23-completion}, suppose
\[
 b_{\rm tot}:=\|{\mathbb B}_{\rm tot}\|<\delta.             \tag{23.33}
\]
Let
\[
 P_{\rm tot}=I-C_{\rm tot}^\dagger C_{\rm tot},\qquad
 T_{\rm tot}=K_{\rm tot}K_{\rm tot}^\dagger,\qquad
 S_{\rm tot}=I-T_{\rm tot}.                                \tag{23.34}
\]
Then
\[
 \boxed{
 \gamma(P_{\rm tot},S_{\rm tot})
 \geq1-\left(\frac{b_{\rm tot}}{\delta}\right)^2,}          \tag{23.35}
\]
and
\[
 \boxed{
 \|{\mathfrak C}_{\rm tot}^\dagger\|
 \leq
 \left(1-\frac{b_{\rm tot}}{\delta}\right)^{-1/2}.}         \tag{23.36}
\]
On the exact first-class constraint surface,
\[
 {\mathbb B}_{\rm tot}=0
 \quad\Longrightarrow\quad
 R_{\rm tot}=P_{\rm tot}-T_{\rm tot},\qquad\gamma=1.        \tag{23.37}
\]
\end{theorem}

\begin{proof}
The V18 propagation condition number is at most
\[
 \|C_{\rm tot}^\dagger\|^2
 \|C_{\rm tot}{\mathbb J}_{\rm tot}C_{\rm tot}^*\|
 \leq b_{\rm tot}/\delta
\]
by (23.26) and (23.31).  Apply V18 or V22.
\end{proof}

This is the first coupled SM--Einstein constraint-angle theorem in
the present cycle.  It is conditional on a single common symplectic
carrier and the mixed block certificate.

\subsection{Coupled physical dimension}

Let
\[
 n_{\rm tot}:=\dim Z,\qquad
 m_{\rm tot}:=\dim Y,\qquad
 r_{\rm rel}:=\operatorname{rank}\Pi_0.                    \tag{23.38}
\]
Under Theorem~\ref{thm:v14v23-completion},
\[
 \operatorname{rank}C_{\rm tot}=m_{\rm tot}-r_{\rm rel}.    \tag{23.39}
\]

\begin{corollary}[Coupled first-class degree count]
\label{cor:v14v23-dimension}
If \({\mathbb B}_{\rm tot}=0\), then
\[
 2(m_{\rm tot}-r_{\rm rel})\leq n_{\rm tot},                \tag{23.40}
\]
and the gauge-reduced physical variation space has dimension
\[
 \boxed{
 n_{\rm tot}-2(m_{\rm tot}-r_{\rm rel}).}                  \tag{23.41}
\]
\end{corollary}

\begin{proof}
Equation (23.25) gives (23.39).  Apply the V22 first-class degree
count to \(C_{\rm tot}\).
\end{proof}

This count includes mixed coupling and reducibility.  It does not
identify the surviving modes with observed Standard Model particles;
that requires the representation, chirality, mass, and continuum
incidence results on the same carrier.

\subsection{Coupled residual estimate}

Let \(r_{\rm tot}\in Z\) be the whitened combined finite
SM--Einstein residual and let \(R_{\rm tot}\) be the physical quotient
projector.  Under (23.33),
\[
\boxed{
\begin{aligned}
 \|r_{\rm tot}\|^2
 \leq{}&\|R_{\rm tot}r_{\rm tot}\|^2\\
 &+\frac{\delta^{-1}
 \bigl(
  \|C_{\rm tot}r_{\rm tot}\|^2
  +\|K_{\rm tot}^*r_{\rm tot}\|^2
 \bigr)}
 {1-b_{\rm tot}/\delta}.
\end{aligned}}                                              \tag{23.42}
\]

\begin{proof}
Use the V22 symplectic three-defect estimate with the coupled floor
\(\delta\).
\end{proof}

Thus full finite residual closure follows from three same-history
coupled channels:
\[
 C_{\rm tot}r_{\rm tot}\to0,\qquad
 K_{\rm tot}^*r_{\rm tot}\to0,\qquad
 R_{\rm tot}r_{\rm tot}\to0,                               \tag{23.43}
\]
provided \(\delta\) stays uniformly above the bracket debit.

\subsection{Cofinal coupled landing}

Let every object carry an index \(n\) and be transported to fixed
finite screens.  Assume
\[
 A_n\succeq\mu_gI,\qquad D_n\succeq\mu_iI,\qquad
 \|A_n^{-1/2}X_nD_n^{-1/2}\|\leq\chi_0<1.                 \tag{23.44}
\]
Put
\[
 \delta_0:=(1-\chi_0)\min\{\mu_g,\mu_i\}.                  \tag{23.45}
\]
Assume also
\[
 \|{\mathbb B}_{{\rm tot},n}\|
 \leq\theta_0\delta_0,\qquad\theta_0<1.                    \tag{23.46}
\]

\begin{theorem}[Uniform coupled landing]
\label{thm:v14v23-cofinal}
If the block maps, compatible symplectic structures, and declared
relation projectors converge in norm on each fixed screen, then
\[
 C_{{\rm tot},n}^\dagger\to C_{\rm tot}^\dagger,\qquad
 P_{{\rm tot},n}\to P_{\rm tot},\qquad
 T_{{\rm tot},n}\to T_{\rm tot},                           \tag{23.47}
\]
\[
 R_{{\rm tot},n}\to R_{\rm tot},\qquad
 {\mathfrak C}_{{\rm tot},n}^\dagger
 \to{\mathfrak C}_{\rm tot}^\dagger.                       \tag{23.48}
\]
If the three defects (23.43) vanish on every fixed localized screen,
then the full finite residual tends to zero there.  With the V9
all-order product landing, the represented weak coupled equation
holds on the limiting declared core.
\end{theorem}

\begin{proof}
The Schur theorem gives the common completed-Gram floor
\(\delta_0\).  The V22 constraint-Gram landing theorem gives
(23.47).  The bracket ratio (23.46) gives a common angle and joint
floor, yielding (23.48).  Apply (23.42) and then the V9 landing
theorem.
\end{proof}

\subsection{Relation to the attached common-action source}

The Grand Tensor source already insists that gravitational, matter,
gauge, potential, counterterm, and boundary rewards occur on one
common action record, and it reconstructs mixed action Grams from a
single pressure jet.  That common provenance is necessary for
(23.2), but it does not itself populate the constraint derivative
blocks \(C_g,C_i,M,N\).  A mixed action Hessian and a mixed
constraint Gram are different derivatives and must not be
interchanged without an explicit identity.

A concrete regulator card must now provide
\[
\begin{array}{c|p{0.72\linewidth}}
\mathrm{C1}&
the gravity HDA and internal Gauss constraint components on one
finite symplectic carrier;\\
\mathrm{C2}&
the four derivative blocks in (23.2), including matter--geometry
cross derivatives;\\
\mathrm{C3}&
sector relation projectors and the mixed relation incidences
(23.4)--(23.5);\\
\mathrm{C4}&
the diagonal completed floors and normalized cross contraction
\(\chi<1\);\\
\mathrm{C5}&
the full first-class bracket matrix and anomaly debit
\(b_{\rm tot}<\delta\);\\
\mathrm{C6}&
the coupled residual channels and cofinal product landing.
\end{array}                                                  \tag{23.49}
\]

The classical internal Gauss bank fits this Hamiltonian compiler.
A quantum BRST cohomology is an additional quotient and is not
silently replaced by the classical Gauss constraint.

\subsection{Next task}

V24 should inspect the Grand Tensor common-action, matter--Legendre,
finite HDA, and Ward sections for an identity connecting their
measured mixed Hessian blocks to the four constraint derivatives in
(23.2).  If no such identity is present, V24 should prove a
non-substitution theorem and give the minimal second-variation
interface needed to derive \(M,N\).  This goes upstream of any
numerical Schur test.

\subsection{Acquisition ledger}

\[
\begin{array}{p{0.29\linewidth}|p{0.23\linewidth}|p{0.38\linewidth}}
\textbf{gate}&\textbf{status after V23}&\textbf{content}\\ \hline
\text{coupled constraint derivative}
 &\text{defined}
 &four gravity/internal blocks\\
\text{mixed relation incidence}
 &\text{identified}
 &M^*\Pi_g=N^*\Pi_i=0\\
\text{coupled Gram floor}
 &\text{acquired conditionally}
 &normalized Schur contraction \(\chi<1\)\\
\text{hidden coupled relations}
 &\text{excluded by the floor}
 &\ker C_{\rm tot}^*=\operatorname{Ran}\Pi_0\\
\text{coupled first-class angle}
 &\text{acquired}
 &b_{\rm tot}/\delta<1\\
\text{coupled cofinal closure}
 &\text{acquired conditionally}
 &three residual channels\\
\text{concrete mixed derivative blocks}
 &\text{open}
 &source card C2\\
\text{uniform numerical margins}
 &\text{open}
 &source cards C4--C5
\end{array}                                                  \tag{23.50}
\]

\subsection{Parent-integrity record}

The installed parent was
\[
\begin{array}{c|c}
 \text{file}&
 \texttt{Renewal\_SM\_Einstein\_Continuum\_Research\_Notes\_V22.tex}\\
 \text{lines}&9146\\
 \text{bytes}&290667\\
 \text{SHA--256}&
 \texttt{8469BEBAE4A259CA3F4B2AB53245A3A5AA9D34E3D45F6D77C74A7FC250B74526}.
\end{array}                                                  \tag{23.51}
\]

% =============================================================================
% END APPENDED FOURTEENTH-CYCLE V23 MATERIAL
% =============================================================================
% =============================================================================
% BEGIN APPENDED FOURTEENTH-CYCLE V24 MATERIAL
% =============================================================================

\section{Fourteenth-cycle V24: multiplier--phase Hessians as the
missing coupled-constraint interface}
\label{sec:v14v24}

\subsection{Source-audit conclusion}

The Grand Tensor common-action material contains several distinct
second derivatives:

\begin{enumerate}
\item a positive source Hessian and its Schur/Ward router;
\item pressure-jet reconstruction of marginal and mixed action-source
Grams;
\item matter velocity and Hamiltonian Hessians related by Legendre
reciprocity;
\item a measured gravity--matter HDA cross anomaly;
\item a finite total geometric first variation and determining
variation panel.
\end{enumerate}

None of the inspected theorems identifies these matrices with the
derivative of the full Hamiltonian, momentum, and internal Gauss
constraint map.  That derivative is a mixed Hessian only after the
actual constraint multipliers are included as typed action
coordinates.  This note proves the precise statement and its
converse limitation.

\subsection{Phase action curvature does not determine constraints}

\begin{countertheorem}[Same phase Hessian, arbitrary mixed constraint
conditioning]
\label{cth:v14v24-nosub}
For \(a\in\mathbb R\), define
\[
\begin{aligned}
 {\mathscr S}_a(x,y,\lambda_g,\lambda_i)
 :={}&\frac12(x^2+y^2)\\
 &+\lambda_g(x+ay)+\lambda_i(y+ax).
\end{aligned}                                               \tag{24.1}
\]
At \(\lambda_g=\lambda_i=0\), every member has the identical phase
action, stationary point, and phase Hessian
\[
 D_{(x,y)}^2{\mathscr S}_a=I_2.                             \tag{24.2}
\]
Nevertheless its constraint derivative is
\[
 C_{{\rm tot},a}
 =
 \begin{pmatrix}1&a\\a&1\end{pmatrix},                     \tag{24.3}
\]
whose singular values are
\[
 |1+a|,\qquad |1-a|.                                       \tag{24.4}
\]
In particular the reduced floor can be made arbitrarily small and a
new relation appears at \(a=1\), while (24.2) remains unchanged.
\end{countertheorem}

\begin{proof}
Differentiate (24.1) twice in the phase variables to obtain (24.2).
The multiplier equations are
\[
 \partial_{\lambda_g}{\mathscr S}_a=x+ay,\qquad
 \partial_{\lambda_i}{\mathscr S}_a=y+ax.
\]
Differentiating them in \((x,y)\) gives (24.3).  Its eigenvalues are
\(1+a\) and \(1-a\), proving (24.4).
\end{proof}

Thus even an exactly known positive phase Hessian does not populate
the V23 constraint Gram.  A common-action source Gram also does not
do so unless its coordinate bank explicitly contains the constraint
multipliers and their cross derivatives with the phase variables.

\subsection{Affine multiplier extraction theorem}

Let
\[
 Z=Z_g\oplus Z_i,\qquad
 Y=Y_g\oplus Y_i                                           \tag{24.5}
\]
be finite-dimensional phase and multiplier spaces.  Write
\[
 z=(z_g,z_i),\qquad
 \lambda=(\lambda_g,\lambda_i).
\]
Let the common finite action have the affine multiplier form
\[
 \boxed{
 {\mathscr S}(z,\lambda)
 =
 {\mathscr S}_0(z)
 +\langle\lambda_g,\Phi_g(z)\rangle
 +\langle\lambda_i,\Phi_i(z)\rangle.}                       \tag{24.6}
\]
Here \(\Phi_g\) is the gravity Hamiltonian/momentum constraint bank
and \(\Phi_i\) is the internal Gauss constraint bank.  Define
\[
 \Phi_{\rm tot}:=(\Phi_g,\Phi_i):Z\to Y.                    \tag{24.7}
\]

\begin{theorem}[Constraint derivative is the multiplier--phase mixed
Hessian]
\label{thm:v14v24-extraction}
At every \(z\),
\[
 \boxed{
 C_{\rm tot}:=D_z\Phi_{\rm tot}(z)
 =D_zD_\lambda{\mathscr S}(z,\lambda),}                    \tag{24.8}
\]
independently of \(\lambda\).  Relative to (24.5),
\[
 \boxed{
 C_{\rm tot}
 =
 \begin{pmatrix}
 D_{z_g}D_{\lambda_g}{\mathscr S}&
 D_{z_i}D_{\lambda_g}{\mathscr S}\\
 D_{z_g}D_{\lambda_i}{\mathscr S}&
 D_{z_i}D_{\lambda_i}{\mathscr S}
 \end{pmatrix}
 =
 \begin{pmatrix}C_g&M\\N&C_i\end{pmatrix}.}                \tag{24.9}
\]
Thus the four V23 blocks are four typed panels of one common action
Hessian.
\end{theorem}

\begin{proof}
Differentiating (24.6) in the multiplier variables gives
\[
 D_{\lambda_g}{\mathscr S}=\Phi_g(z),\qquad
 D_{\lambda_i}{\mathscr S}=\Phi_i(z).
\]
Differentiate these identities in \(z_g,z_i\).  Equality of the
displayed mixed derivatives follows directly from the affine form;
no inversion or field equation is used.
\end{proof}

\begin{corollary}[Mixed relation incidence is an action-Hessian
identity]
\label{cor:v14v24-relations}
Let \(\Pi_g,\Pi_i\) be declared relation projectors.  The V23
conditions
\[
 C_g^*\Pi_g=M^*\Pi_g=0,\qquad
 C_i^*\Pi_i=N^*\Pi_i=0                                    \tag{24.10}
\]
are equivalent to
\[
 (D_zD_{\lambda_g}{\mathscr S})^*\Pi_g=0,\qquad
 (D_zD_{\lambda_i}{\mathscr S})^*\Pi_i=0.                  \tag{24.11}
\]
\end{corollary}

\begin{proof}
The first multiplier row of (24.9) is \((C_g\ M)\), and the second is
\((N\ C_i)\).  Taking adjoints gives the equivalence.
\end{proof}

This shows exactly which common-action Hessian rows must be measured
to verify that isolated sector relations survive the coupling.

\subsection{Gauge-fixed and nonlinear multiplier warning}

If an action contains a nonlinear multiplier term,
\[
 {\mathscr S}(z,\lambda)
 =
 {\mathscr S}_0(z)+\langle\lambda,\Phi(z)\rangle
 +{\mathscr S}_{\rm gf}(z,\lambda),                        \tag{24.12}
\]
then
\[
 D_zD_\lambda{\mathscr S}
 =
 D\Phi(z)+D_zD_\lambda{\mathscr S}_{\rm gf}.               \tag{24.13}
\]
The second term belongs to the declared gauge-fixed system.  It
cannot be included in the physical constraint derivative without
changing the constraint map.

\begin{proposition}[Exact gauge-fixing subtraction]
\label{prop:v14v24-gf}
If \({\mathscr S}_{\rm gf}\) is known, then
\[
 \boxed{
 D\Phi(z)
 =
 D_zD_\lambda{\mathscr S}
 -D_zD_\lambda{\mathscr S}_{\rm gf}.}                      \tag{24.14}
\]
If the second term is not separately identified, the mixed action
Hessian determines only the gauge-fixed derivative, not the physical
constraint derivative.
\end{proposition}

\begin{proof}
Rearrange (24.13).  The non-identification statement follows because
arbitrary mixed terms can be shifted between
\(\langle\lambda,\Phi\rangle\) and
\({\mathscr S}_{\rm gf}\) while keeping their sum fixed.
\end{proof}

This is the Hessian version of the Grand Tensor firewall that a
positive lapse--shift Hessian belongs to a declared reduced or
gauge-fixed system.

\subsection{Normalization of multiplier coordinates}

Let \(U:Y\to Y\) be invertible and set
\[
 \lambda'=U\lambda.
\]
Writing the same affine action in the new coordinate gives
\[
 \Phi'=U^{-*}\Phi,\qquad
 C'=U^{-*}C.                                                \tag{24.15}
\]

\begin{proposition}[Invariant quotient, noninvariant numerical floor]
\label{prop:v14v24-normalization}
Under (24.15),
\[
 \ker C'=\ker C,\qquad
 \operatorname{Ran}(C')^*=\operatorname{Ran}C^*.           \tag{24.16}
\]
Hence the physical projection \(P\), Hamiltonian gauge range
\({\mathbb J}\operatorname{Ran}C^*\), tangent projection \(T\), and
quotient projection \(R\) are unchanged.  The numerical singular
floor of \(C'\) need not equal that of \(C\) unless \(U\) is unitary
for the declared multiplier metric.
\end{proposition}

\begin{proof}
Left multiplication by an invertible map does not change the kernel,
and
\[
 (C')^*=C^*U^{-1}
\]
has the same range as \(C^*\).  The projection statements follow.
Singular values are preserved only by unitary left multiplication.
\end{proof}

Therefore every quantitative V22--V23 floor must record the physical
multiplier metric and units.  Rescaling a constraint can improve or
degrade the displayed number without changing the quotient geometry.

\subsection{Finite-difference acquisition with an outward error}

Suppose the action is available as an exact or interval-valued
finite-stage oracle.  For unit directions \(u\in Z\), \(v\in Y\),
let
\[
 f(s,t):={\mathscr S}(z+su,\lambda+tv).                    \tag{24.17}
\]
Define the symmetric mixed difference
\[
 {\cal D}_{h,k}f
 :=
 \frac{
 f(h,k)-f(h,-k)-f(-h,k)+f(-h,-k)}
 {4hk}.                                                     \tag{24.18}
\]

\begin{theorem}[Certified mixed-derivative enclosure]
\label{thm:v14v24-difference}
Assume \(f_{st}\) exists on
\([-h,h]\times[-k,k]\) and
\[
 |\partial_sf_{st}(s,t)|\leq L_s,\qquad
 |\partial_tf_{st}(s,t)|\leq L_t.                          \tag{24.19}
\]
Then
\[
 \boxed{
 \left|
 {\cal D}_{h,k}f
 -D_zD_\lambda{\mathscr S}(z,\lambda)[u,v]
 \right|
 \leq\frac{L_sh+L_tk}{2}.}                                 \tag{24.20}
\]
\end{theorem}

\begin{proof}
Two applications of the fundamental theorem of calculus give
\[
 {\cal D}_{h,k}f
 =
 \frac1{4hk}
 \int_{-h}^{h}\int_{-k}^{k}
 f_{st}(s,t)\,dt\,ds.                                      \tag{24.21}
\]
By (24.19),
\[
 |f_{st}(s,t)-f_{st}(0,0)|
 \leq L_s|s|+L_t|t|.
\]
Average this inequality over the rectangle.  Since the averages of
\(|s|\) and \(|t|\) are \(h/2\) and \(k/2\), respectively,
(24.20) follows.
\end{proof}

Applying (24.20) to a declared orthonormal basis gives outward
intervals for every entry of \(C_g,M,N,C_i\).  An operator-norm
radius for the assembled matrix must then be obtained from a verified
matrix norm bound; entrywise midpoint values alone are not a Schur
certificate.

\subsection{Perturbation of the coupled completed Gram}

Let \(C\) be the exact constraint derivative and
\(\widetilde C=C+\Delta\) a certified approximation with
\[
 \|\Delta\|\leq d.                                         \tag{24.22}
\]
Let \(\Pi\) be a common declared relation projector satisfying
\(C^*\Pi=\widetilde C^*\Pi=0\).  Suppose
\[
 CC^*+\Pi\succeq\mu I.                                     \tag{24.23}
\]

\begin{theorem}[Hessian-panel error debit]
\label{thm:v14v24-perturb}
One has
\[
 \left\|
 \widetilde C\widetilde C^*-CC^*
 \right\|
 \leq2\|C\|d+d^2.                                         \tag{24.24}
\]
Consequently, if
\[
 \widetilde\mu
 :=\mu-2\|C\|d-d^2>0,                                     \tag{24.25}
\]
then
\[
 \widetilde C\widetilde C^*+\Pi
 \succeq\widetilde\mu I,                                  \tag{24.26}
\]
the declared relation projector is complete for
\(\widetilde C^*\), and
\[
 \operatorname{rmin}\widetilde C\geq\sqrt{\widetilde\mu}.
                                                               \tag{24.27}
\]
\end{theorem}

\begin{proof}
Expand
\[
 \widetilde C\widetilde C^*-CC^*
 =C\Delta^*+\Delta C^*+\Delta\Delta^*
\]
and take norms to obtain (24.24).  Subtract this radius from
(24.23), then apply the V16 completed-kernel theorem.
\end{proof}

The common relation incidence in the hypothesis is essential.  If
the approximation violates it, its relation defect must be included
as a separate block rather than hidden in (24.25).

\subsection{What the existing Grand Tensor theorems do and do not
populate}

The common-action Hessian/Ward-router theorem proves a Schur identity
among source maps when its graph symmetry lies in the Hessian kernel.
The matter--Legendre theorem proves reciprocity of velocity and
Hamiltonian Hessians and identifies an HDA cross anomaly.  The
pressure-jet theorem reconstructs all mixed blocks among its declared
reward coordinates.

These results populate (24.8) only if the declared coordinate bank
contains:

\begin{enumerate}
\item the actual lapse and shift multipliers;
\item the actual internal Gauss multipliers;
\item the phase variables with respect to which the constraint map is
linearized;
\item the direct multiplier--phase mixed derivatives, separated from
gauge-fixing terms;
\item one physical multiplier metric and normalization.
\end{enumerate}

The inspected text does not present that five-row matrix card.
Accordingly, V23's \(M,N\) remain unpopulated.  The source's measured
gravity--matter HDA cross anomaly contributes instead to the bracket
debit \(b_{\rm tot}\), after its norm and same-history incidence are
supplied.

\subsection{Cofinal interface}

Let \({\mathscr S}_n\) be same-history finite actions with affine
multiplier coordinates and let their certified mixed panels satisfy
\[
 \|D_zD_\lambda{\mathscr S}_n-C_{\rm tot}\|\to0.            \tag{24.28}
\]
If their relation projectors converge and a uniform completed-Gram
floor survives the error debit (24.25), then V22--V23 give convergence
of the physical, tangent, and quotient projections.  If the coupled
bracket debit stays below that floor and the three localized residual
channels vanish, the V23 cofinal weak-equation conclusion follows.

This statement remains conditional on actual actions and panels; it
does not promote the abstract source Hessian to a constraint
derivative.

\subsection{Next checkpoint}

V25 is the next compulsory companion audit.  After that audit, V26
should assemble a fully typed finite acquisition ledger for one
selected graph regulator: multiplier coordinates, phase variables,
constraint values, mixed Hessian panels, relation projectors, HDA
bracket anomaly, and the three Einstein residual channels.  If the
sources lack numerical or symbolic entries for that card, the ledger
should return an explicit finite list of unpopulated matrices rather
than another abstract positivity theorem.

\subsection{Acquisition ledger}

\[
\begin{array}{p{0.29\linewidth}|p{0.23\linewidth}|p{0.38\linewidth}}
\textbf{gate}&\textbf{status after V24}&\textbf{content}\\ \hline
\text{action-Hessian substitution}
 &\text{refuted}
 &same phase Hessian, arbitrary constraint floor\\
\text{constraint extraction}
 &\text{acquired}
 &C_{\rm tot}=D_zD_\lambda{\mathscr S}\\
\text{four coupled blocks}
 &\text{typed}
 &multiplier--phase panels\\
\text{gauge-fixing separation}
 &\text{acquired}
 &subtract \(D_zD_\lambda{\mathscr S}_{\rm gf}\)\\
\text{normalization invariance}
 &\text{acquired}
 &quotient fixed, numerical floor metric-dependent\\
\text{finite-difference enclosure}
 &\text{acquired}
 &mixed panel plus outward remainder\\
\text{concrete multiplier panel}
 &\text{open}
 &source values not present in one card\\
\text{coupled Schur margin}
 &\text{open}
 &requires populated \(C_g,M,N,C_i\)
\end{array}                                                  \tag{24.29}
\]

\subsection{Parent-integrity record}

The installed parent was
\[
\begin{array}{c|c}
 \text{file}&
 \texttt{Renewal\_SM\_Einstein\_Continuum\_Research\_Notes\_V23.tex}\\
 \text{lines}&9704\\
 \text{bytes}&306678\\
 \text{SHA--256}&
 \texttt{E36FF4251C00A4ABC9CC64879262024F22F5CBC2DA45D0FC76D8E1640DF7F711}.
\end{array}                                                  \tag{24.30}
\]

% =============================================================================
% END APPENDED FOURTEENTH-CYCLE V24 MATERIAL
% =============================================================================
% =============================================================================
% BEGIN APPENDED FOURTEENTH-CYCLE V25 MATERIAL
% =============================================================================

\section{Fourteenth-cycle V25: compulsory companion audit before the
finite regulator card}
\label{sec:v14v25}

\subsection{Files inspected}

The required twenty-fifth-version audit was performed read-only on
\[
\begin{array}{p{0.53\linewidth}|r|r}
\textbf{file}&\textbf{lines}&\textbf{bytes}\\ \hline
\texttt{Renewal\_Yang\_Mills\_Mass\_Gap\_Research\_Notes\_V7.tex}
 &1868&67814\\
\texttt{Renewal\_NS\_Stability\_Research\_Notes\_V26.tex}
 &5319&278283.
\end{array}                                                  \tag{25.1}
\]
Their SHA--256 digests are
\[
\begin{aligned}
 &\texttt{4D52B6E7FABB15A73D312694D1F598560A9AD36F2B0393D247481DB2E36AABBD},
                                                               \tag{25.2}\\
 &\texttt{82C887F8C2C69E4F28FAD7C3DC8DE5B9099CB5A4BF431EBA1FFF3E91935978AD}.
                                                               \tag{25.3}
\end{aligned}
\]
No companion file was modified or removed.

\subsection{Yang--Mills V7}

The Yang--Mills programme has advanced one version since the V20
audit.  V7 exactly regroups its moving degree-one \(90\times90\)
block in a second phase and constructs rational noncommutative
splittings for the resulting \(180\times180\) blocks at all sixteen
remaining transverse anchors.  Exact congruence certificates give
strict positive floors and thereby eliminate two momentum coordinates
on the remaining \(4^2\) grid.  The smaller Wilson block floor requires
higher rational precision.  Two recursive coordinate eliminations
remain, and no full-torus or mass-gap conclusion is claimed.

This is relevant only to the V16 fallback.  It demonstrates that a
moving degree-one splitting can be lifted recursively while every
acceptance step remains an exact rational matrix certificate.  It
also shows that the certified margin can shrink substantially after a
lift, so precision and floor accounting must be carried at every
level.  None of its matrices or floors applies to the coupled
SM--Einstein constraint card of V23--V24.

\subsection{Navier--Stokes V26}

The Navier--Stokes note remains unchanged.  Its rank-one Oseen-kernel
calculation supplies no multiplier--phase Hessian, constraint
relation, first-class bracket, or Einstein residual value.  Its
structured-input warning remains compatible with keeping the
physical constraint and gauge ranges explicit.

\subsection{Audit decision}

\begin{decision}[Direction after V25]
\label{dec:v14v25-next}
The primary route remains:
\[
 \text{common affine multiplier action}
 \Longrightarrow
 C_{\rm tot}=D_zD_\lambda{\mathscr S}
 \Longrightarrow
 \text{completed coupled constraint Gram}
 \Longrightarrow
 \text{first-class angle and three-defect closure}.        \tag{25.4}
\]
No companion result changes this order.

V26 will produce a finite acquisition ledger for one selected graph
regulator.  Each row will be classified as:

\begin{enumerate}
\item explicitly defined in the attached source;
\item conditionally assumed there;
\item derivable from a displayed identity once source values are
provided;
\item absent and therefore still requiring a new computation.
\end{enumerate}

The ledger will include dimensions, domains and codomains, action
coordinates, derivative blocks, relation projectors, completed floors,
HDA anomaly, Ward residual, physical quotient residual, and cofinal
transport.  An abstract theorem will not be counted as a populated
finite matrix.
\end{decision}

\begin{proof}
V21--V24 identify the minimal data needed by the primary route.
The companion audit contributes a method for a different Laurent
positivity branch but no entry in that primary data set.  The
classification therefore prevents a conditional source theorem from
being mistaken for a numerical or symbolic regulator card.
\end{proof}

\subsection{Acquisition ledger}

\[
\begin{array}{p{0.29\linewidth}|p{0.23\linewidth}|p{0.38\linewidth}}
\textbf{gate}&\textbf{status after V25}&\textbf{content}\\ \hline
\text{companion checkpoint}
 &\text{acquired}
 &YM V7 and NS V26 read with hashes\\
\text{recursive Laurent fallback}
 &\text{strengthened externally}
 &two exact YM coordinate eliminations\\
\text{primary coupled route}
 &\text{unchanged}
 &multiplier--phase Hessian first\\
\text{finite regulator ledger}
 &\text{next}
 &typed population audit\\
\text{coupled derivative values}
 &\text{open}
 &four V23 blocks\\
\text{uniform coupled floor}
 &\text{open}
 &relation and Schur certificates\\
\text{three residual channels}
 &\text{open}
 &same-history evaluations
\end{array}                                                  \tag{25.5}
\]

The next compulsory companion audit is V30.

\subsection{Parent-integrity record}

The installed parent was
\[
\begin{array}{c|c}
 \text{file}&
 \texttt{Renewal\_SM\_Einstein\_Continuum\_Research\_Notes\_V24.tex}\\
 \text{lines}&10178\\
 \text{bytes}&321578\\
 \text{SHA--256}&
 \texttt{58B4FA15EC573B00C6B4E96D5DBCC9508F0C2BD9280132BAEA5CBE532C557D0F}.
\end{array}                                                  \tag{25.6}
\]

% =============================================================================
% END APPENDED FOURTEENTH-CYCLE V25 MATERIAL
% =============================================================================
% =============================================================================
% BEGIN APPENDED FOURTEENTH-CYCLE V26 MATERIAL
% =============================================================================

\section{Fourteenth-cycle V26: a finite acquisition ledger for the
coupled constraint card}
\label{sec:v14v26}

\subsection{Purpose and status convention}

The preceding notes reduced the finite constraint problem to a
specific object,
\[
 C_{{\rm tot},h}=D_zD_\lambda\mathscr S_h
 =
 \begin{pmatrix}
 C_{g,h}&M_h\\
 N_h&C_{i,h}
 \end{pmatrix}:
 Z_{g,h}\oplus Z_{i,h}\longrightarrow
 Y_{g,h}\oplus Y_{i,h}.                         \tag{26.1}
\]
Here \(z\) is a phase variable and \(\lambda\) is an affine
multiplier variable.  Formula (26.1) is a typing requirement, not a
claim that the four matrices have already been extracted from the
attached source.  This version fixes one candidate finite spatial
cell and records, row by row, what is populated and what is not.

We use the following four status symbols:
\[
 \mathbf E=\hbox{explicitly populated},\qquad
 \mathbf C=\hbox{conditional source input},\qquad
 \mathbf D=\hbox{derived after named rows are populated},\qquad
 \mathbf M=\hbox{missing finite datum}.                         \tag{26.2}
\]
A row marked \(\mathbf C\) is not promoted to \(\mathbf E\) merely
because a source theorem proves consequences from it.  A row marked
\(\mathbf D\) carries no independent numerical content.

\subsection{The selected graph skeleton}

Let
\[
 X_4=\{1,2,3,4\},\qquad
 E_4=\operatorname{span}_{\mathbb R}
 \{e_{12},e_{13},e_{14},e_{23},e_{24},e_{34}\}
 \cong\mathbb R^6,                                      \tag{26.3}
\]
where every edge is oriented from its smaller to its larger endpoint.
Put
\[
 V_4=\mathbb R^{X_4},\qquad
 W_4=\left\{x\in V_4:\sum_{a=1}^4x_a=0\right\},
 \qquad P_{W_4}=I-\frac14\mathbf1\mathbf1^*.             \tag{26.4}
\]
The incidence map \(B_4:E_4\to W_4\) is
\[
 B_4e_{ab}=e_b-e_a.                                      \tag{26.5}
\]
In the displayed edge order its matrix, regarded first as a map into
\(V_4\), is
\[
 [B_4]=
 \begin{pmatrix}
 -1&-1&-1& 0& 0& 0\\
  1& 0& 0&-1&-1& 0\\
  0& 1& 0& 1& 0&-1\\
  0& 0& 1& 0& 1& 1
 \end{pmatrix}.                                          \tag{26.6}
\]

\begin{lemma}[Exact \(K_4\) cut--loop card]
\label{lem:v14v26-k4}
The matrix in (26.6) satisfies
\[
 B_4B_4^*=4P_{W_4},
 \qquad
 P_{\rm cut}=\frac14B_4^*B_4,
 \qquad
 P_{\rm loop}=I_{E_4}-P_{\rm cut}.                       \tag{26.7}
\]
Consequently
\[
 E_4=\operatorname{Ran}B_4^*\mathbin{\oplus}\ker B_4,
 \qquad
 \dim\operatorname{Ran}B_4^*=3,qquad
 \dim\ker B_4=3.                                        \tag{26.8}
\]
On \(W_4\), the nonzero singular values of \(B_4\) all equal
\(2\).
\end{lemma}

\begin{proof}
The diagonal entries of \(B_4B_4^*\) are the three edges incident on
each vertex, and every off-diagonal entry is \(-1\).  Hence
\[
 B_4B_4^*=4I_{V_4}-\mathbf1\mathbf1^*=4P_{W_4}.
\]
It follows that \(\frac14B_4^*B_4\) is self-adjoint and
\[
 \left(\frac14B_4^*B_4\right)^2
 =\frac1{16}B_4^*(4P_{W_4})B_4
 =\frac14B_4^*B_4,
\]
because \(P_{W_4}B_4=B_4\).  Its range is
\(\operatorname{Ran}B_4^*\), so it is the cut projector and its
orthogonal complement is \(\ker B_4\).  The rank is
\(\dim W_4=3\), which proves (26.8).  Restricting the first identity
in (26.7) to \(W_4\) gives the singular values.
\end{proof}

Thus the graph-algebra part of the candidate regulator is
\(\mathbf E\).  In the locked source normalization with a positive
weight \(\vartheta_h\), the associated endpoint and loop Grams are,
conditionally on that normalization,
\[
 G_{{\rm end},h}=4\vartheta_h I_3,qquad
 G_{{\rm loop},h}=\vartheta_h I_3,qquad
 G_{{\rm cross},h}=0.                                    \tag{26.9}
\]
The value and uniform lower bound of \(\vartheta_h\) belong to the
chosen source packet.  More importantly, neither (26.7) nor (26.9)
is the constraint Gram of (26.1).

For bookkeeping, the selected skeleton is the spatial cell \(K_4\)
together with one discrete time slab \([t,t+h]\).  The Grand Tensor
complete-cell selector supplies this branch only under its stated
operational hypotheses; retention of one common history carrier on
the slab is likewise conditional.  We denote such a carrier by
\(\mathcal H_h^{\rm ret}\) without assigning it a dimension that the
source does not assign.

\subsection{Typed acquisition ledger}

The complete input card needed at this cell is
\[
 \mathcal A_h=
 \bigl(
 \mathcal H_h^{\rm ret},Z_{g,h},Z_{i,h},Y_{g,h},Y_{i,h},
 \mathbb J_h,\mathscr S_h,\Pi_{g,h},\Pi_{i,h},
 \mathcal B_h,\mathcal R_h,\iota_h
 \bigr),                                                  \tag{26.10}
\]
where \(\mathcal B_h\) is the finite bracket card,
\(\mathcal R_h\) is the same-history residual card, and \(\iota_h\)
denotes the cofinal comparison maps.  The present population audit is
\[
\begin{array}{p{0.05\linewidth}|p{0.25\linewidth}|p{0.13\linewidth}|p{0.45\linewidth}}
\textbf{row}&\textbf{finite datum}&\textbf{status}&\textbf{precise content}\\ \hline
L1&graph chain spaces and incidence&\mathbf E&
\(E_4,W_4,B_4,P_{\rm cut},P_{\rm loop}\), with (26.3)--(26.8)\\
L2&selected retained history carrier&\mathbf C&
one common carrier is assumed by the finite HDA and weak-handoff theorems;
no concrete basis is selected\\
L3&gravitational phase space&\mathbf M&
finite Hilbert space \(Z_{g,h}\), including canonical coordinates and
inner product\\
L4&internal SM phase space&\mathbf M&
finite Hilbert space \(Z_{i,h}\), including gauge and matter canonical
coordinates on the same carrier\\
L5&multiplier spaces&\mathbf M&
based Hilbert spaces \(Y_{g,h}\) for lapse/shift constraints and
\(Y_{i,h}\) for internal Gauss constraints\\
L6&common symplectic operator&\mathbf M&
matrix \(\mathbb J_h:Z_h\to Z_h\),
\(\mathbb J_h^*=-\mathbb J_h\), \(\mathbb J_h^*\mathbb J_h=I\)\\
L7&affine multiplier action&\mathbf M&
\(\mathscr S_h(z,\lambda)=\mathscr S_{0,h}(z)+
\langle\lambda,\Phi_h(z)\rangle\) in the bases of L3--L6\\
L8&mixed Hessian panel&\mathbf D&
\(C_{{\rm tot},h}=D\Phi_h=D_zD_\lambda\mathscr S_h\), after L7\\
L9&four coupled blocks&\mathbf D&
\(C_{g,h},M_h,N_h,C_{i,h}\), obtained by compressing L8 to the
splittings in L3--L5\\
L10&declared label-relation projectors&\mathbf M&
exact bases for \(\operatorname{Ran}\Pi_{g,h}\subset Y_{g,h}\) and
\(\operatorname{Ran}\Pi_{i,h}\subset Y_{i,h}\), with the four mixed
incidences \(C_{g,h}^*\Pi_{g,h}=M_h^*\Pi_{g,h}=0\) and
\(C_{i,h}^*\Pi_{i,h}=N_h^*\Pi_{i,h}=0\)\\
L11&completed diagonal Grams&\mathbf D&
\(A_h=C_{g,h}C_{g,h}^*+M_hM_h^*+\Pi_{g,h}\) and
\(D_h=N_hN_h^*+C_{i,h}C_{i,h}^*+\Pi_{i,h}\)\\
L12&coupled Schur certificate&\mathbf D&
\(X_h=C_{g,h}N_h^*+M_hC_{i,h}^*\), floors
\(\mu_{g,h},\mu_{i,h}\), and
\(\chi_h=\|A_h^{-1/2}X_hD_h^{-1/2}\|\)\\
L13&finite HDA brackets&\mathbf C&
exact displayed HDA structure on a reducing retained carrier; the
gravity--matter cross anomaly is symbolically defined\\
L14&bracket anomaly norm&\mathbf M&
same-basis value or outward enclosure of
\(a_h=\|C_{{\rm tot},h}\mathbb J_hC_{{\rm tot},h}^*\|\)\\
L15&source frame and ADM--Cartan reconstruction&\mathbf C&
conditional on source floors, coherent history, horizontal gauge,
and Hodge--Sobolev data\\
L16&finite Einstein residual functional&\mathbf E&
defined by the common-action residual theorem\\
L17&represented residual vectors and Gram&\mathbf D&
Riesz representatives on the selected finite Hilbert structure,
after L2--L7 are fixed\\
L18&three residual values&\mathbf M&
constraint, tangent-quotient, and Einstein/Ward residuals evaluated
on one retained history\\
L19&positive cofinal screens and comparison maps&\mathbf C&
assumptions of the local weak Einstein handoff; no populated uniform
card on the selected branch
\end{array}                                                  \tag{26.11}
\]

Rows L3--L7 are deliberately separate.  A list of formal lapse,
shift, and gauge labels does not fix their Hilbert normalization; a
kinematic metric and extrinsic curvature do not by themselves fix a
canonical symplectic phase chart; and a common action does not by
itself display the affine multiplier coordinates required by L7.

\subsection{The compiler theorem for a completed ledger}

Put \(Z_h=Z_{g,h}\oplus Z_{i,h}\) and
\(Y_h=Y_{g,h}\oplus Y_{i,h}\).  For populated L7 and L10 set
\[
 C_h=C_{{\rm tot},h}=D\Phi_h,\qquad
 \Pi_{0,h}=\Pi_{g,h}\oplus\Pi_{i,h},\qquad
 \widehat G_h=C_hC_h^*+\Pi_{0,h}.                        \tag{26.12}
\]
The mixed incidences in L10 give \(C_h^*\Pi_{0,h}=0\).  Equality
\(\operatorname{Ran}\Pi_{0,h}=\ker C_h^*\) is not assumed: it is a
consequence of the completed-Gram floor below.  Without that floor,
coupling may create additional relations.

\begin{theorem}[Finite-ledger sufficiency]
\label{thm:v14v26-ledger}
Suppose rows L2--L10 are populated on a finite regulator, the four
mixed relation incidences in L10 hold, and the coupled completed Gram
satisfies
\[
 \widehat G_h\succeq\delta_h I_{Y_h},\qquad \delta_h>0.    \tag{26.13}
\]
Then
\[
 \boxed{\ker C_h^*=\operatorname{Ran}\Pi_{0,h}}.
\]
Define
\[
 P_h=I-C_h^\dagger C_h,\qquad
 K_h=\mathbb J_hC_h^*,\qquad
 T_h=K_hK_h^\dagger,\qquad S_h=I-T_h,
 \qquad
 R_h=P_{\ker C_h\cap(\operatorname{Ran}K_h)^\perp}.       \tag{26.14}
\]
Then all operators in (26.14) are determined by the ledger,
\[
 \|C_h^\dagger\|\le\delta_h^{-1/2},\qquad
 \|K_h^\dagger\|\le\delta_h^{-1/2}.                     \tag{26.15}
\]
If
\[
 a_h=\|C_h\mathbb J_hC_h^*\|<\delta_h,                   \tag{26.16}
\]
then, with \(\varepsilon_h=\|(I-P_h)T_h\|\),
\[
 \varepsilon_h\le\frac{a_h}{\delta_h}<1,\qquad
 \gamma(P_h,S_h)=1-\varepsilon_h^2
 \ge 1-\left(\frac{a_h}{\delta_h}\right)^2.             \tag{26.17}
\]
For every \(r\in Z_h\),
\[
 \|r\|^2\le \|R_hr\|^2+
 \frac{\delta_h^{-1}\|C_hr\|^2+
       \delta_h^{-1}\|K_h^*r\|^2}
      {1-a_h/\delta_h}.                                  \tag{26.18}
\]
If the finite brackets are first class at the background, then
\(a_h=0\), \(\operatorname{Ran}K_h\subseteq\ker C_h\), and
\[
 R_h=P_h-T_h,\qquad \gamma(P_h,S_h)=1.                    \tag{26.19}
\]
Thus L2--L14 compile the coupled physical projector, its angle
certificate, and the two constraint/gauge defect terms in the
same-history residual estimate.  Rows L15--L19 are still required for
the geometric and cofinal continuum handoff.
\end{theorem}

\begin{proof}
The mixed incidences give \(C_h^*\Pi_{0,h}=0\).  Apply the V16
kernel-completion theorem to (26.13).  It proves
\(\ker C_h^*=\operatorname{Ran}\Pi_{0,h}\) and says that every
nonzero singular value of \(C_h\) is at least
\(\sqrt{\delta_h}\), proving the first estimate in (26.15).  Since
\(\mathbb J_h\) is unitary, \(K_h=\mathbb J_hC_h^*\) has the same
nonzero singular values, proving the second.

The exact escape identity from V18 gives
\[
 (I-P_h)T_h
 =C_h^\dagger(C_hK_h)K_h^\dagger,
 \qquad C_hK_h=C_h\mathbb J_hC_h^*.
\]
Taking norms and using (26.15) proves the first inequality in
(26.17).  The sharp two-subspace identity from V18 then gives
\(\gamma(P_h,S_h)=1-\varepsilon_h^2\), proving the remainder of
(26.17).  Applying the V18 three-defect estimate and substituting
(26.15) and \(\varepsilon_h\le a_h/\delta_h\) proves (26.18).

At a first-class background, the Poisson bracket matrix of the
constraint map is \(C_h\mathbb J_hC_h^*\).  Its vanishing gives
\(C_hK_h=0\), hence
\(\operatorname{Ran}K_h\subseteq\ker C_h\).  Orthogonal nested
projectors then satisfy \(R_h=P_h-T_h\), and the exact gap between
\(P_h\) and \(S_h=I-T_h\) is one.  The last assertion follows by
tracing which ledger rows define each quantity; it does not assert the
geometric hypotheses L15--L19.
\end{proof}

\begin{corollary}[Diagonal-to-coupled certificate]
\label{cor:v14v26-schur}
Suppose the diagonal completed Grams in L11 obey
\[
 A_h\succeq\mu_{g,h}I,qquad
 D_h\succeq\mu_{i,h}I                                  \tag{26.20}
\]
and the exact coupled completed Gram has the block form
\[
 \widehat G_h=\begin{pmatrix}A_h&X_h\\X_h^*&D_h\end{pmatrix}.
                                                               \tag{26.21}
\]
If
\[
 \chi_h=\|A_h^{-1/2}X_hD_h^{-1/2}\|<1,                  \tag{26.22}
\]
then Theorem \ref{thm:v14v26-ledger} applies with
\[
 \delta_h=(1-\chi_h)\min\{\mu_{g,h},\mu_{i,h}\}.         \tag{26.23}
\]
\end{corollary}

\begin{proof}
For \(u\in Y_{g,h}\), \(v\in Y_{i,h}\), put
\(p=A_h^{1/2}u\) and \(q=D_h^{1/2}v\).  Then
\[
 2|\langle u,X_hv\rangle|
 \le 2\chi_h\|p\|\|q\|
 \le\chi_h(\|p\|^2+\|q\|^2).
\]
Consequently the quadratic form of (26.21) is at least
\((1-\chi_h)(\|p\|^2+\|q\|^2)\), which is bounded below by
the right side of (26.23) times \(\|u\|^2+\|v\|^2\).
\end{proof}

The corollary explains exactly what L12 must contain.  Merely listing
the four blocks of (26.1) is insufficient: their relation spaces and
the normalized off-diagonal coupling must also be certified.

\subsection{Non-substitution rules}

Several already available positive operators have tempting but
incorrect types.  The following table records the firewall.
\[
\begin{array}{p{0.22\linewidth}|p{0.28\linewidth}|p{0.40\linewidth}}
\textbf{available object}&\textbf{actual type}&\textbf{why it cannot fill L8--L12}\\ \hline
\(K_4\) incidence \(B_4\)&\(E_4\to W_4\)&
its singular value \(2\) measures graph boundary conditioning, not
phase-to-constraint conditioning\\
endpoint/loop source Gram&source-channel coefficient space&
its floor controls reconstruction of source modes after a source packet
is chosen, not the derivative of lapse/shift/Gauss constraints\\
finite Dirac graph regulator&fermion coefficient space to itself&
its spectral gap controls a doubled fermion Hamiltonian and determinant
phase; it has neither multiplier codomain nor gravitational rows\\
determining panel&evaluation map on a chosen observable/residual class&
injectivity on that class does not supply a canonical phase chart or a
constraint singular-value floor\\
finite HDA identity&bilinear bracket of constraint labels&
first-class closure controls tangency but not the nonzero singular values
of the constraint derivative\\
ADM source-frame theorem&source coefficients to reconstructed geometry&
it is conditional on a source-frame floor and begins after the finite
constraint card required here
\end{array}                                                  \tag{26.24}
\]

\begin{proposition}[Bracket closure has no hidden Gram floor]
\label{prop:v14v26-no-floor}
First-class closure alone cannot imply a positive lower bound for the
completed constraint Gram.
\end{proposition}

\begin{proof}
Let \(C:Z\to Y\) be any nonzero finite constraint derivative with
\(C\mathbb JC^*=0\), and for \(0<\rho\le1\) put
\(C_\rho=\rho C\).  Then
\[
 C_\rho\mathbb JC_\rho^*=0                              \tag{26.25}
\]
for every \(\rho\), so the same first-class statement holds.  The
nonzero eigenvalues of \(C_\rho C_\rho^*\), however, are \(\rho^2\)
times those of \(CC^*\) and tend to zero.  The relation projector is
unchanged, so the least eigenvalue of
\(C_\rho C_\rho^*+P_{\ker C^*}\) tends to zero whenever the constraint
range is nontrivial.  No scale-free bracket identity supplies L11.
\end{proof}

\begin{proposition}[Graph conditioning has no canonical constraint
interpretation]
\label{prop:v14v26-type}
The identities (26.7)--(26.9) do not determine any one of
\(C_{g,h},M_h,N_h,C_{i,h}\).
\end{proposition}

\begin{proof}
Their common domain, when defined, is the selected phase space
\(Z_{g,h}\oplus Z_{i,h}\), and their common codomain is the multiplier
label space \(Y_{g,h}\oplus Y_{i,h}\).  By contrast, \(B_4\) has the
fixed combinatorial type \(E_4\to W_4\).  The source provides no
canonical unitary identifications of these pairs of spaces.  Even
after arbitrary identifications in coincident dimensions, composing
\(B_4\) with arbitrary invertible coordinate changes on phase or
multiplier space changes the mixed Hessian and its Hilbert singular
values while preserving the graph identity (26.7).  Therefore the
graph card does not determine the constraint blocks.
\end{proof}

\subsection{The first missing row}

\begin{proposition}[Document-level acquisition result]
\label{prop:v14v26-first-missing}
On the selected \(K_4\), one-time-slab branch, the first object that
must be newly constructed before a coupled constraint floor can be
proved is a based affine multiplier chart
\[
 (z_g,z_i;N,\beta,A_0)longmapsto
 \mathscr S_h(z_g,z_i;N,\beta,A_0)                       \tag{26.26}
\]
on one retained carrier, together with the mixed derivative
\[
 D_{(z_g,z_i)}D_{(N,\beta,A_0)}\mathscr S_h.              \tag{26.27}
\]
Here \(N\) denotes lapse labels, \(\beta\) shift labels, and \(A_0\)
internal temporal-gauge or Gauss multipliers.  Their actual finite
bases and inner products are part of the required construction.
\end{proposition}

\begin{proof}
The attached material explicitly supplies the finite graph chain
algebra, conditional retained-carrier HDA formulae, a common-action
framework, matter Legendre data, source reconstruction statements,
and finite residual functionals.  These objects populate L1, support
L2, L13, L15, and L16, and motivate L3--L7.  They do not display, in
one finite coordinate card, the phase bases, multiplier bases,
symplectic matrix, and affine lapse/shift/internal-Gauss coupling.
Without these data (26.27) has no matrix representation and even its
Hilbert adjoint is undetermined.  Relations, completed Grams, and
Schur margins are functions of that adjoint, so L10--L12 cannot be
computed first.  V24 proves that once (26.26) is affine, (26.27) is
exactly the desired constraint derivative; hence (26.26)--(26.27),
rather than another abstract positivity lemma, is the first missing
row.
\end{proof}

This proposition is a statement about what has been acquired from the
attached finite programme, not a claim that no discretization of ADM
or Yang--Mills theory is known elsewhere.  Importing an external
discretization would require proving that it uses the same histories,
source normalization, action, and cofinal maps.

\subsection{A machine-checkable acquisition packet}

To prevent later ambiguity, a populated finite card at scale \(h\)
shall contain the following tuple:
\[
\begin{split}
\mathfrak P_h=\big(&
\mathcal B(Z_{g,h}),\mathcal B(Z_{i,h}),
\mathcal B(Y_{g,h}),\mathcal B(Y_{i,h}),
[\mathbb J_h],\\
&[D_zD_\lambda\mathscr S_h],
\mathcal B(\ker C_h^*),[\Pi_{0,h}],
\underline\delta_h,\overline a_h,
\mathbf r_h,\mathsf G_{\mathbf r,h},
[\iota_h]\big).                                         \tag{26.28}
\end{split}
\]
Here \(\mathcal B(U)\) is an ordered orthonormal basis, brackets denote
the resulting matrices, \(\underline\delta_h\) is a proved lower bound
for (26.13), \(\overline a_h\) is a proved upper bound for (26.16),
\(\mathbf r_h\) is the vector of the three same-history residual
channels, and \(\mathsf G_{\mathbf r,h}\) is their represented Gram.
Every floating-point field must be accompanied by an outward interval
or an exact certificate.  In particular, acceptance requires
\[
 0\le\overline a_h<\underline\delta_h,
 \qquad
 \underline\gamma_h=
 1-(\overline a_h/\underline\delta_h)^2>0.                \tag{26.29}
\]

\begin{lemma}[Basis covariance of the acceptance inequality]
\label{lem:v14v26-basis}
If the phase and label bases in (26.28) are changed by orthogonal
matrices, then \(\delta_h\), \(a_h\), and the inequality (26.29) are
unchanged.
\end{lemma}

\begin{proof}
For orthogonal \(U:Z_h\to Z_h\) and \(V:Y_h\to Y_h\), the matrices
transform as
\[
 C_h\mapsto VC_hU^*,\qquad
 \mathbb J_h\mapsto U\mathbb J_hU^*,\qquad
 \Pi_{0,h}\mapsto V\Pi_{0,h}V^*.
\]
Thus \(\widehat G_h\mapsto V\widehat G_hV^*\) and
\(C_h\mathbb J_hC_h^*\mapsto
V(C_h\mathbb J_hC_h^*)V^*\).  Eigenvalues and operator norms are
preserved.
\end{proof}

This invariance does not permit arbitrary nonunitary rescaling of
multiplier labels.  Such a rescaling changes the Hilbert problem and
must be recorded as part of L5.

\subsection{Upstream move selected for V27}

The source does contain a positive matter velocity Hessian and a
Legendre-transform theorem.  The next useful upstream question is
therefore whether a finite kinetic Hessian, together with an affine
lapse/shift/Gauss action, canonically supplies L3--L8 on the retained
carrier.  V27 will formulate and prove the precise finite Legendre
completion statement, including its failure modes, and then compare
its hypotheses against the Grand Tensor action blocks.  It will not
identify the graph incidence floor with a constraint floor.

The acquisition state after V26 is
\[
\begin{array}{p{0.31\linewidth}|p{0.20\linewidth}|p{0.38\linewidth}}
\textbf{gate}&\textbf{status}&\textbf{next operation}\\ \hline
\(K_4\) cut--loop algebra&closed&retain as exact spatial skeleton\\
retained common carrier&conditional&select a concrete carrier and basis\\
canonical phase chart&open&derive from a nondegenerate finite Legendre map\\
affine multiplier chart&open&exhibit lapse, shift, and internal Gauss rows\\
coupled mixed Hessian&open&differentiate the populated action\\
relation and Schur certificate&blocked downstream&compute only after the
mixed Hessian exists\\
angle certificate&compiler proved&apply (26.13)--(26.18)\\
geometric/cofinal handoff&conditional&populate L15--L19 on the same histories
\end{array}                                                  \tag{26.30}
\]

\subsection{Parent-integrity record}

The installed parent was
\[
\begin{array}{c|c}
 \text{file}&
 \texttt{Renewal\_SM\_Einstein\_Continuum\_Research\_Notes\_V25.tex}\\
 \text{lines}&10327\\
 \text{bytes}&327003\\
 \text{SHA--256}&
 \texttt{4E7DCAFA4009E4B9335C588E953998657E994075E8ABACC9DB860ACE37FCE486}.
\end{array}                                                  \tag{26.31}
\]

% =============================================================================
% END APPENDED FOURTEENTH-CYCLE V26 MATERIAL
% =============================================================================
% =============================================================================
% BEGIN APPENDED FOURTEENTH-CYCLE V27 MATERIAL
% =============================================================================

\section{Fourteenth-cycle V27: finite Legendre completion of an
ADM--Gauss multiplier chart}
\label{sec:v14v27}

\subsection{Why the ordinary Legendre theorem is not yet the chart}

V26 identified the first missing finite row as an affine lapse,
shift, and internal-Gauss multiplier action on one based phase space.
The Grand Tensor matter theorem proves, for a positive matter velocity
Hessian \(\mathbb K\), the local Legendre identities
\[
 p=D_v\mathscr L,
 \qquad D_p\mathscr H=v,
 \qquad D_p^2\mathscr H=\mathbb K^{-1}.                    \tag{27.1}
\]
It also proves the typed matter bracket, the
\(a_rb_r=1\) reciprocity relation, and a formula for the
gravity--matter cross anomaly.  Those are retained.  The new issue is
whether the lapse dependence of \(\mathbb K^{-1}\), and the shift and
Gauss dependence of the velocity, assemble into a Hamiltonian that is
linear in the multiplier labels.  Positivity of \(\mathbb K\) alone
does not imply that linearity.

The construction below is finite-dimensional and local in a selected
configuration chart.  It allows an indefinite gravitational kinetic
form: nondegeneracy, rather than positivity, is what the Legendre map
requires.  Positivity remains available on matter subblocks when the
source theorem supplies it.

\subsection{Inverse-lapse kinetic normal form}

Let \(Q_h\) be a finite-dimensional real Hilbert configuration space.
Using its inner product to identify \(Q_h\) with \(Q_h^*\), put
\[
 Z_h=Q_h\oplus Q_h,
 \qquad z=(q,p),
 \qquad
 \mathbb J_h=
 \begin{pmatrix}0&-I\\ I&0\end{pmatrix}.                  \tag{27.2}
\]
Let \(Y_{N,h},Y_{\beta,h},Y_{G,h}\) be based finite Hilbert spaces of
lapse, shift, and internal temporal-gauge labels.  Suppose that for
q in a configuration neighbourhood we are given smooth maps
\[
 \mathcal A_q:Y_{N,h}\longrightarrow\operatorname{Sym}(Q_h),
 \quad
 \mathcal V(q)\in Y_{N,h},                                \tag{27.3}
\]
\[
 \mathcal D_q:Y_{\beta,h}\longrightarrow Q_h,
 \qquad
 \mathcal R_q:Y_{G,h}\longrightarrow Q_h.                \tag{27.4}
\]
The first map is linear in \(N\).  On an open lapse cone
\(\mathcal U_N\subset Y_{N,h}\), assume
\[
 \mathcal A_q(N)=\mathcal A_q(N)^*,
 \qquad \operatorname{rmin}\mathcal A_q(N)>0.             \tag{27.5}
\]
No sign condition is imposed in (27.5).

For
\[
 w=\dot q-\mathcal D_q\beta-\mathcal R_q a_0             \tag{27.6}
\]
define the finite inverse-lapse Lagrangian
\[
 \mathscr L_h(q,\dot q;N,\beta,a_0)
 =\frac12\langle w,\mathcal A_q(N)^{-1}w\rangle
  -\langle N,\mathcal V(q)\rangle.                        \tag{27.7}
\]

\begin{theorem}[Finite ADM--Gauss Legendre compiler]
\label{thm:v14v27-compiler}
Under (27.3)--(27.7), the Legendre map in \(\dot q\) is invertible and
is given by
\[
 p=D_{\dot q}\mathscr L_h
   =\mathcal A_q(N)^{-1}w,
 \qquad
 \dot q=\mathcal A_q(N)p+mathcal D_q\beta+
              \mathcal R_q a_0.                           \tag{27.8}
\]
Its Hamiltonian is exactly affine in the three multiplier banks:
\[
 \mathscr H_h(q,p;N,\beta,a_0)
 =\langle N,\Phi_{N,h}(q,p)\rangle
  +\langle\beta,\Phi_{\beta,h}(q,p)\rangle
  +\langle a_0,\Phi_{G,h}(q,p)\rangle,                    \tag{27.9}
\]
where
\[
 \langle N,\Phi_{N,h}(q,p)\rangle
 =\frac12\langle p,\mathcal A_q(N)p\rangle
  +\langle N,\mathcal V(q)\rangle,                        \tag{27.10}
\]
\[
 \Phi_{\beta,h}(q,p)=\mathcal D_q^*p,
 \qquad
 \Phi_{G,h}(q,p)=\mathcal R_q^*p.                         \tag{27.11}
\]
Consequently the canonical phase action
\[
 \mathscr S_{H,h}=\int
 \bigl(\langle p,\dot q\rangle-
 \langle(N,\beta,a_0),\Phi_h(q,p)\rangle\bigr)\,dt       \tag{27.12}
\]
has the multiplier equations \(\Phi_h=0\).  After the harmless sign
coordinate \(\lambda=-(N,\beta,a_0)\), it has the V24 form
\[
 \mathscr S_{H,h}=\mathscr S_{0,h}+
 \int\langle\lambda,\Phi_h(z)\rangle\,dt,
 \qquad
 D_zD_\lambda\mathscr S_{H,h}=D\Phi_h.                   \tag{27.13}
\]
Thus (27.3)--(27.7), in declared bases, populate V26 rows L3--L8.
They do not by themselves prove first-class closure or a completed
constraint-Gram floor.
\end{theorem}

\begin{proof}
Differentiating (27.7) in \(\dot q\) gives the first identity in
(27.8).  Invertibility of \(\mathcal A_q(N)\) gives the second.  Now
\[
 \begin{aligned}
 \langle p,\dot q\rangle-\mathscr L_h
 &={}
 \langle p,\mathcal A_q(N)p+\mathcal D_q\beta+
                    \mathcal R_q a_0\rangle\\
 &\quad-\\frac12\langle\mathcal A_q(N)p,
                         \mathcal A_q(N)^{-1}
                         \mathcal A_q(N)p\rangle
       +\langle N,\mathcal V(q)\rangle\\
 &=\frac12\langle p,\mathcal A_q(N)p\rangle
   +\langle p,\mathcal D_q\beta\rangle
   +\langle p,\mathcal R_q a_0\rangle
   +\langle N,\mathcal V(q)\rangle.
 \end{aligned}                                             \tag{27.14}
\]
Linearity of \(N\mapsto\mathcal A_q(N)\) and the adjoint identities
in (27.11) give (27.9)--(27.11).  Varying (27.12) in the independent
multipliers gives the three constraints.  The sign change in (27.13)
is an orthogonal change on the multiplier label space, so it records
the V24 convention without changing any singular value or relation.
Equation (27.2) is skew-adjoint and unitary, hence is the compatible
canonical symplectic operator required by V26.  Affinity supplies the
mixed Hessian, while bracket closure and conditioning are separate
properties of that derivative.
\end{proof}

\subsection{The actual constraint derivative supplied by the compiler}

For \(u=(\xi,\pi)\in Z_h\), differentiation of
(27.10)--(27.11) gives the weak rows
\[
\boxed{
 \begin{aligned}
 \langle N,D\Phi_{N,h}(z)u\rangle
 ={}&\langle\pi,\mathcal A_q(N)p\rangle
 +\frac12\langle p,(D_q\mathcal A[\xi])(N)p\rangle\\
 &+\langle N,D_q\mathcal V[\xi]\rangle,
 \end{aligned}}                                            \tag{27.15}
\]
\[
\boxed{
 \langle\beta,D\Phi_{\beta,h}(z)u\rangle
 =\langle\pi,\mathcal D_q\beta\rangle
  +\langle p,(D_q\mathcal D[\xi])\beta\rangle,}          \tag{27.16}
\]
\[
\boxed{
 \langle a_0,D\Phi_{G,h}(z)u\rangle
 =\langle\pi,\mathcal R_q a_0\rangle
  +\langle p,(D_q\mathcal R[\xi])a_0\rangle.}            \tag{27.17}
\]
These formulas are immediately matrix-valued once bases for the four
spaces and coordinate formulae for
\(\mathcal A,\mathcal V,\mathcal D,\mathcal R\) are supplied.

\begin{corollary}[Four-block extraction]
\label{cor:v14v27-blocks}
Suppose
\[
 Q_h=Q_{g,h}\oplus Q_{i,h},
 \qquad
 Z_h=Z_{g,h}\oplus Z_{i,h},                               \tag{27.18}
\]
where \(Z_{s,h}=Q_{s,h}\oplus Q_{s,h}\), and suppose
\(Y_{N,h}\oplus Y_{\beta,h}\) is declared gravitational while
\(Y_{G,h}\) is internal.  Compressing (27.15)--(27.17) to these
splittings gives uniquely
\[
 D\Phi_h=
 \begin{pmatrix}C_{g,h}&M_h\\N_h&C_{i,h}\end{pmatrix}.   \tag{27.19}
\]
Every nonzero off-diagonal derivative of
\(\mathcal A,\mathcal V,\mathcal D\), or \(\mathcal R\) contributes
to \(M_h\) or \(N_h\); it may not be discarded before the V26 Schur
test.
\end{corollary}

\begin{proof}
The ordered orthogonal decompositions in (27.18) give four canonical
matrix compressions of the single linear map \(D\Phi_h:Z_h\to Y_h\).
Equations (27.15)--(27.17) list every derivative entering that map.
A term depending on a variation from the other sector lies, by
definition, in an off-diagonal compression.
\end{proof}

\begin{proposition}[Uniform Legendre transport]
\label{prop:v14v27-transport}
Let the maps in (27.3)--(27.4) carry a cofinal index \(n\).  On a fixed
finite screen suppose
\[
 \mathcal A_{q,n}(N)\longrightarrow\mathcal A_q(N),
 \qquad
 \operatorname{rmin}\mathcal A_{q,n}(N)\ge\kappa>0,       \tag{27.20}
\]
and suppose the maps and their first \(q\)-derivatives converge in
operator norm.  Then the Legendre inverses, the constraints
(27.10)--(27.11), and their derivatives (27.15)--(27.17) converge in
operator norm on bounded \((q,p)\)-sets.
\end{proposition}

\begin{proof}
The inverse identity
\[
 A_n^{-1}-A^{-1}=A_n^{-1}(A-A_n)A^{-1}                    \tag{27.21}
\]
and (27.20) give norm convergence of the inverse kinetic operators.
All remaining expressions are finite sums of products of convergent
maps with bounded \(q,p\).  The product rule applied to
(27.15)--(27.17) gives the derivative claim.
\end{proof}

This proposition transports a chart once it has been populated.  It
does not provide the uniform constraint floor \(\delta_h\); that is a
singular-value statement about the complete matrix (27.19).

\subsection{An exact criterion for multiplier affinity}

The inverse-lapse normal form is sufficient, but it is useful to know
what must be tested in a more general source action.

\begin{theorem}[Legendre affinity criterion]
\label{thm:v14v27-affinity}
Fix \(q\), and let \(\mathscr L(q,v;\lambda)\) have an invertible
velocity Hessian on a connected velocity chart.  Let
\(\mathscr H(q,p;\lambda)\) be its Legendre transform.  On a
star-shaped momentum chart, \(\mathscr H\) is affine in \(\lambda\)
if and only if all three functions
\[
 \mathscr H(q,0;\lambda),
 \qquad D_p\mathscr H(q,0;\lambda),
 \qquad D_p^2\mathscr H(q,p;\lambda)                       \tag{27.22}
\]
are affine in \(\lambda\) for every \(p\) in the chart.  Moreover,
\[
 D_p^2\mathscr H(q,p;\lambda)
 =\left[D_v^2\mathscr L(q,v(q,p,\lambda);\lambda)
  \right]^{-1}.                                            \tag{27.23}
\]
Thus inverse-Hessian affinity is a necessary multiplier-row test; for
a velocity-quadratic action it becomes sufficient together with the
affinity of the zero-momentum drift and potential.
\end{theorem}

\begin{proof}
If \(\mathscr H\) is affine in \(\lambda\), every \(p\)-derivative in
(27.22) is affine.  Conversely, Taylor's formula with integral
remainder gives
\[
 \begin{aligned}
 \mathscr H(q,p;\lambda)
 ={}&\mathscr H(q,0;\lambda)
 +D_p\mathscr H(q,0;\lambda)[p]\\
 &+\int_0^1(1-t)
 D_p^2\mathscr H(q,tp;\lambda)[p,p],dt .
 \end{aligned}                                             \tag{27.24}
\]
Every term on the right is affine in \(\lambda\), proving the
converse.  Differentiate
\(p=D_v\mathscr L(q,v;\lambda)\) at fixed \(q,\lambda\):
\[
 I=D_v^2\mathscr L\,D_pv.
\]
The envelope identity gives \(D_p\mathscr H=v\), hence
\(D_p^2\mathscr H=D_pv\), proving (27.23).  In the
velocity-quadratic case the Hessian is independent of \(v\), and the
constant and linear terms in momentum are exactly the potential and
drift terms.
\end{proof}

\begin{countertheorem}[Positive velocity Hessian need not yield an
affine lapse]
\label{cth:v14v27-positive-nonaffine}
For \(N\in\mathbb R\), let
\[
 \mathscr L_N(v)=\frac12e^Nv^2.                            \tag{27.25}
\]
Then \(D_v^2\mathscr L_N=e^N>0\) for every \(N\), but
\[
 \mathscr H_N(p)=\frac12e^{-N}p^2                         \tag{27.26}
\]
is not affine in \(N\) on any open interval.
\end{countertheorem}

\begin{proof}
The momentum is \(p=e^Nv\), so \(v=e^{-N}p\), which gives
(27.26).  Its second \(N\)-derivative is
\(\frac12e^{-N}p^2\), nonzero when \(p\ne0\).
\end{proof}

This counterexample is the precise reason that the positive matter
Hessian in (27.1) cannot, without a lapse-incidence identity, be
entered as V26 row L7.

\subsection{Affinity does not imply first-class closure}

\begin{countertheorem}[An affine multiplier action can be second
class]
\label{cth:v14v27-second-class}
On \(Z=\mathbb R_q\oplus\mathbb R_p\), take
\[
 \mathscr S=\mathscr S_0+lambda_1q+lambda_2p.             \tag{27.27}
\]
The constraint map is affine and its derivative is the identity, but
\[
 \{q,p\}=1.                                                \tag{27.28}
\]
Thus its bracket matrix is nonzero at the constraint surface.
\end{countertheorem}

\begin{proof}
The multiplier derivatives of (27.27) are \(q\) and \(p\).  Their
canonical Poisson bracket is one.  Equivalently, with the canonical
\(\mathbb J\), the matrix \(C\mathbb JC^*\) is \(\mathbb J\ne0\).
\end{proof}

Hence the compiler has a strict order:
\[
 \text{Legendre regularity}
 \Longrightarrow\text{canonical phase chart}
 \Longrightarrow\text{affine constraint derivative}
 \Longrightarrow\text{relation and Gram tests}
 \Longrightarrow\text{HDA/angle test}.                    \tag{27.29}
\]
No arrow in (27.29) is reversed in the acquisition ledger.

\subsection{Comparison with the attached source}

The exact source comparison is now
\[
\begin{array}{p{0.28\linewidth}|p{0.15\linewidth}|p{0.46\linewidth}}
\textbf{compiler datum}&\textbf{status}&\textbf{source comparison}\\ \hline
positive matter velocity Hessian&\mathbf C&
the typed matter Legendre theorem assumes \(\mathbb K\succ0\) and derives
\(D_p^2\mathscr H=\mathbb K^{-1}\)\\
matter reciprocity/HDA current&\mathbf C&
exact after conductance matching, cross-kernel symmetry, and a positive
species-current determining Gram\\
full configuration space \(Q_{g,h}\oplus Q_{i,h}\)&\mathbf M&
no ordered canonical basis is supplied on the selected \(K_4\) slab\\
gravitational Legendre block&\mathbf M&
the ADM metric/extrinsic-curvature reconstruction is kinematic and does not
display a nondegenerate finite DeWitt Legendre matrix in the required basis\\
lapse co-metric incidence \(N\mapsto\mathcal A_q(N)\)&\mathbf M&
the inverse matter Hessian is not displayed as a linear function of the
finite lapse labels\\
shift drift \(\mathcal D_q\)&\mathbf M&
endpoint/Cartan reconstruction supplies a shift under conditional source
hypotheses, but not its canonical action on every phase coordinate\\
internal gauge drift \(\mathcal R_q\)&\mathbf M&
classical Gauss and BRST labels occur, but no same-card temporal-gauge drift
matrix is displayed\\
potential row \(\mathcal V(q)\)&\mathbf M&
local scalar and plaquette actions are explicit, but their assembly with the
gravitational lapse bank on this phase chart is not populated\\
constraint derivative&\mathbf D&
equations (27.15)--(27.19) compute it once the preceding four missing maps
are supplied
\end{array}                                                  \tag{27.30}
\]


\begin{decision}[Direction after V27]
\label{dec:v14v27-next}
The Legendre step is no longer an abstract bottleneck: Theorem
\ref{thm:v14v27-compiler} gives a sufficient finite normal form and
Theorem \ref{thm:v14v27-affinity} gives an exact source test.  The
first missing source row has moved upstream to the maps
\[
 \boxed{\mathcal A_q,\quad\mathcal V(q),\quad
        \mathcal D_q,\quad\mathcal R_q}                    \tag{27.31}
\]
on a common based gravitational--internal configuration carrier.

V28 will attack the least geometric of these missing maps first.  It
will determine what the \(K_4\) endpoint representation can and cannot
supply for the shift drift \(\mathcal D_q\), construct the universal
infinitesimal pullback action on vertex and edge fields, and isolate
the additional field-representation data needed to extend that action
to coframes, connections, scalars, and fermions.  If the endpoint
representation does not determine the drift, V28 will give the exact
kernel witness and move further upstream.
\end{decision}

\subsection{Acquisition ledger}

\[
\begin{array}{p{0.30\linewidth}|p{0.20\linewidth}|p{0.39\linewidth}}
\textbf{gate}&\textbf{status after V27}&\textbf{content}\\ \hline
finite Legendre compiler&proved&inverse-lapse normal form\\
canonical symplectic chart&derived conditionally&standard \(Q_h\oplus Q_h\)
chart after a configuration basis is supplied\\
affine multiplier test&proved&zero jet plus inverse-Hessian affinity\\
matter Legendre block&partially populated&positive typed sector only\\
gravitational kinetic block&open&nondegenerate finite DeWitt row\\
lapse incidence&open&linear co-metric map on the common carrier\\
shift drift&next&finite field pullback representation\\
internal Gauss drift&open&temporal-gauge infinitesimal action\\
constraint Gram and angle&downstream&apply V26 only after (27.31)
\end{array}                                                  \tag{27.32}
\]

\subsection{Parent-integrity record}

The installed parent was
\[
\begin{array}{c|c}
 \text{file}&
 \texttt{Renewal\_SM\_Einstein\_Continuum\_Research\_Notes\_V26.tex}\\
 \text{lines}&10891\\
 \text{bytes}&349036\\
 \text{SHA--256}&
 \texttt{FF1342293DFF657E65980B94661B3CB9251F7EB1CD3626C80D77E439728D576B}.
\end{array}                                                  \tag{27.33}
\]

% =============================================================================
% END APPENDED FOURTEENTH-CYCLE V27 MATERIAL
% =============================================================================
% =============================================================================
% BEGIN APPENDED FOURTEENTH-CYCLE V28 MATERIAL
% =============================================================================

\section{Fourteenth-cycle V28: endpoint shift drift, cut lifts, and the
finite-algebra obstruction}
\label{sec:v14v28}

\subsection{Source interface and the question to be answered}

The source endpoint theorem conditionally reconstructs a shift
\(\beta_h\), an actual pullback
\(\mathscr P_{\beta,h}G_h(t+h)\), and the limiting expression
\[
 K=\frac1{2\mathcal N}
 \bigl(\partial_0G-\mathcal L_\beta G\bigr).                \tag{28.1}
\]
V27 requires a different, earlier object:
\[
 \mathcal D_q:Y_{\beta,h}\longrightarrow Q_h,              \tag{28.2}
\]
the infinitesimal action of a shift label on every finite
configuration coordinate.  This note asks how much of (28.2) follows
from the exact \(K_4\) endpoint algebra in V26.

There are two distinct issues.  First, a three-component endpoint
shift must be identified with the three-dimensional standard
\(K_4\) module.  Second, a vertex displacement must be lifted to an
edge flow before it can act on fields.  The first has a tetrahedral
solution after a frame is chosen.  The second has an unavoidable
three-dimensional circulation ambiguity.

\subsection{The tetrahedral vector identification}

Let
\[
 r_1=\frac1{\sqrt3}(1,1,1),\quad
 r_2=\frac1{\sqrt3}(1,-1,-1),\quad
 r_3=\frac1{\sqrt3}(-1,1,-1),\quad
 r_4=\frac1{\sqrt3}(-1,-1,1).                              \tag{28.3}
\]
Then
\[
 \sum_{a=1}^4r_a=0,\qquad
 r_a\cdot r_b=
 \begin{cases}1,&a=b,\\-1/3,&a\ne b.\end{cases}            \tag{28.4}
\]
Recall \(W_4=\mathbf1^\perp\subset\mathbb R^4\).

\begin{lemma}[Oriented tetrahedral frame]
\label{lem:v14v28-tetra}
The map
\[
 \tau:W_4\longrightarrow\mathbb R^3,\qquad
 \tau x=\frac{\sqrt3}{2}\sum_{a=1}^4x_ar_a                 \tag{28.5}
\]
is an isometry.  It identifies the standard \(K_4\) vertex module
with one spatial vector port after the ordered tetrahedral frame
(28.3) is declared.
\end{lemma}

\begin{proof}
For \(x,y\in W_4\), (28.4) gives
\[
 \begin{aligned}
 \left\langle\sum_ax_ar_a,\sum_by_br_b\right\rangle
 &=\sum_ax_ay_a-\frac13\sum_{a\ne b}x_ay_b\\
 &=\sum_ax_ay_a+\frac13\sum_ax_ay_a
 =\frac43\langle x,y\rangle,
 \end{aligned}
\]
because \(\sum_{a,b}x_ay_b=0\).  Multiplication by
\(\sqrt3/2\) therefore gives an isometry.  The last assertion records
the required frame choice; changing the ordered tetrahedron composes
\(\tau\) with an orthogonal transformation.
\end{proof}

The source's rank-one vector-isotypic statement provides one shift
port only conditionally.  Even on that branch, using (28.5) requires
the physical vector frame to be identified with the ordered
tetrahedral frame.  We therefore write
\[
 \widehat\beta=\tau^{-1}\beta\in W_4                     \tag{28.6}
\]
only after this interface is declared.

\subsection{No nonzero exact infinitesimal pullback on a finite set}

Let \(\mathcal F_4=\mathbb R^{X_4}\) carry pointwise multiplication.

\begin{theorem}[Finite commutative algebras have no shift derivation]
\label{thm:v14v28-no-derivation}
Every derivation \(D:\mathcal F_4\to\mathcal F_4\) is zero.  Moreover,
every algebra automorphism of \(\mathcal F_4\) is induced by a
permutation of \(X_4\).  Hence every continuous one-parameter group of
algebra automorphisms through the identity is constant.
\end{theorem}

\begin{proof}
Let \(e_a\) be the indicator of vertex \(a\).  Since \(e_a^2=e_a\),
the Leibniz rule gives
\[
 D(e_a)=2e_aD(e_a).                                       \tag{28.7}
\]
At coordinate \(a\), (28.7) says \(D(e_a)_a=2D(e_a)_a\);
at every other coordinate its right side is zero.  Thus
\(D(e_a)=0\).  The \(e_a\) form a basis, so \(D=0\).

An algebra automorphism preserves the nonzero primitive idempotents.
Those are exactly \(e_1,\ldots,e_4\), so the automorphism permutes
them.  Conversely every vertex permutation is an automorphism.
The automorphism group is therefore the finite discrete group
\(S_4\); a continuous path in it beginning at the identity is
constant.
\end{proof}

This theorem does not forbid a useful finite difference operator.  It
says that such an operator cannot simultaneously be nonzero and obey
the exact continuum Leibniz rule on the fixed vertex algebra.  The
source's actual endpoint pullback compares two time slices; it need
not be a continuous algebra automorphism of one fixed four-point set.

\subsection{A canonical edge-flux action on vertex scalars}

Represent an edge flow by an antisymmetric array
\(b_{xy}=-b_{yx}\), equivalently by its six components on the
oriented edges of V26.  Define
\[
 (\mathcal L_b^{(0)}f)_x
 :=\frac12\sum_{y\ne x}b_{xy}(f_y-f_x).                    \tag{28.8}
\]
Let \(J_4:E_4\to V_4\) be the unsigned endpoint map
\[
 J_4e_{ab}=e_a+e_b.
\]
In the fixed V26 orientation, (28.8) is
\[
 \mathcal L_b^{(0)}
 =\frac12J_4\operatorname{Diag}(b)B_4^*.                  \tag{28.9}
\]
Although (28.9) uses an oriented coordinate vector, simultaneous
reversal of an edge changes the signs of both \(b_e\) and the
corresponding row of \(B_4^*\); the operator is orientation
independent.

\begin{proposition}[Exact properties of the scalar flux action]
\label{prop:v14v28-flux}
For every \(b\in E_4\),
\[
 \mathcal L_b^{(0)}\mathbf1=0,\qquad
 \|\mathcal L_b^{(0)}\|
 \le\sqrt6\,\|b\|_2.                                      \tag{28.10}
\]
The assignment \(b\mapsto\mathcal L_b^{(0)}\) is injective and
equivariant under vertex permutations.  Its Leibniz defect is
\[
\boxed{
 \begin{aligned}
 \Delta_b(f,g)
 &:=
 \mathcal L_b^{(0)}(fg)
 -f\,\mathcal L_b^{(0)}g
 -g\,\mathcal L_b^{(0)}f\\
 &=-\frac12B_4\operatorname{Diag}(b)
 \bigl((B_4^*f)\odot(B_4^*g)\bigr).
 \end{aligned}}                                            \tag{28.11}
\]
In particular,
\[
 \|\Delta_b(f,g)\|_2
 \le \|b\|_\infty
 \|(B_4^*f)\odot(B_4^*g)\|_2.                              \tag{28.12}
\]
\end{proposition}

\begin{proof}
Every edge difference of \(\mathbf1\) vanishes, proving the first
identity.  Direct multiplication gives
\[
 J_4J_4^*=2I+\mathbf1\mathbf1^*,
\]
so \(\|J_4\|=\sqrt6\).  V26 gives \(\|B_4^*\|=2\), and
\(\|\operatorname{Diag}(b)\|\le\|b\|_2\).  Equation (28.9) therefore
gives (28.10).

For \(x\ne y\), the \((x,y)\) entry of
\(\mathcal L_b^{(0)}\) is \(b_{xy}/2\).  Thus the operator determines
every component of \(b\), proving injectivity.  Formula (28.8) makes
permutation equivariance immediate.

For one oriented edge \(a<b\), its contribution to the Leibniz defect
at \(a\) is
\[
 \frac{b_{ab}}2(f_b-f_a)(g_b-g_a),
\]
while its contribution at \(b\) is the negative of this expression.
Summing the signed endpoint contributions gives (28.11).  Since
\(\|B_4\|=2\), taking norms gives (28.12).
\end{proof}

Under a physical scaling \(h^{-1}\mathcal L_b^{(0)}\), the Leibniz
defect is bounded by
\[
 h^{-1}\|b\|_\infty
 \|(B_4^*f)\odot(B_4^*g)\|_2.                              \tag{28.13}
\]
It therefore vanishes on a smooth cofinal branch if both edge
differences are \(O(h)\) and \(b\) is uniformly bounded.  This is a
cofinal consistency statement, not an exact finite derivation.

\subsection{The cut lift and the circulation kernel}

An endpoint imbalance \(\widehat\beta\in W_4\) does not itself belong
to the edge-flow space.  The exact sequence is
\[
 0\longrightarrow\ker B_4
 \longrightarrow E_4
 \xrightarrow{\,B_4\,}W_4
 \longrightarrow0.                                        \tag{28.14}
\]

\begin{theorem}[Minimum-norm endpoint lift]
\label{thm:v14v28-cut-lift}
For \(\widehat\beta\in W_4\), define
\[
 b_{\rm cut}(\widehat\beta)
 :=B_4^\dagger\widehat\beta
 =\frac14B_4^*\widehat\beta.                               \tag{28.15}
\]
Then
\[
 B_4b_{\rm cut}(\widehat\beta)=\widehat\beta,\qquad
 \|b_{\rm cut}(\widehat\beta)\|_2=\frac12\|\widehat\beta\|_2.
                                                               \tag{28.16}
\]
It is the unique minimum-norm lift.  Every other lift is uniquely
\[
 b=b_{\rm cut}(\widehat\beta)+c,\qquad c\in\ker B_4.        \tag{28.17}
\]
Consequently
\[
 \mathcal D^{\rm sc}_{\widehat\beta}f
 :=\mathcal L_{b_{\rm cut}(\widehat\beta)}^{(0)}f          \tag{28.18}
\]
is a canonical \(K_4\) scalar drift only after the orthogonal
zero-circulation convention has been declared.
\end{theorem}

\begin{proof}
V26 gives \(B_4B_4^*=4I\) on \(W_4\), which proves the first identity
in (28.16).  Moreover,
\[
 \|b_{\rm cut}\|^2
 =\frac1{16}\langle\widehat\beta,B_4B_4^*
                    \widehat\beta\rangle
 =\frac14\|\widehat\beta\|^2.
\]
The orthogonal decomposition
\(E_4=\operatorname{Ran}B_4^*\oplus\ker B_4\) proves
(28.17) and the minimum-norm assertion.  Equation (28.18) is then
well-defined.  The last qualification is necessary because
orthogonality is extra Hilbert data and the physical endpoint theorem
does not state that circulation vanishes.
\end{proof}

In vertex coordinates, (28.18) has the explicit formula
\[
 \boxed{
 (\mathcal D^{\rm sc}_{\widehat\beta}f)_x
 =\frac18\sum_{y\ne x}
   (\widehat\beta_y-\widehat\beta_x)(f_y-f_x).}             \tag{28.19}
\]
Thus it can be entered as a literal bilinear matrix once
\(\widehat\beta=\tau^{-1}\beta\) is fixed.

\begin{countertheorem}[Endpoint data do not determine circulation]
\label{cth:v14v28-cycle}
Let
\[
 c=e_{12}+e_{23}-e_{13}.                                   \tag{28.20}
\]
Then \(B_4c=0\), but \(\mathcal L_c^{(0)}\ne0\).  Hence two edge flows
with identical endpoint imbalance can induce different scalar
drifts.
\end{countertheorem}

\begin{proof}
Using \(B_4e_{ab}=e_b-e_a\),
\[
 B_4c=(e_2-e_1)+(e_3-e_2)-(e_3-e_1)=0.
\]
For \(f=e_1\), formula (28.8) gives
\[
 \mathcal L_c^{(0)}f=(0,-1/2,1/2,0),
\]
which is nonzero.  By (28.17), adding \(c\) preserves the endpoint
imbalance while changing the drift.
\end{proof}

The witness (28.20) is the exact upstream obstruction requested in
V27.  The endpoint vector port fixes at most the cut component of an
edge flow.  Its loop component is new data unless a zero-circulation
principle is proved from the physical action.

\subsection{What incidence determines on edge fields}

Let \(A:V_4\to V_4\) be a scalar drift with \(A\mathbf1=0\), and put
\(\overline A=P_{W_4}AP_{W_4}\).  Define its cut-edge lift by
\[
 A_{\rm cut}^{(1)}
 :=\frac14B_4^*\overline A B_4:E_4\to E_4.                 \tag{28.21}
\]

\begin{theorem}[Exact-edge lift and loop nonuniqueness]
\label{thm:v14v28-edge-lift}
The map in (28.21) has range in
\(\operatorname{Ran}B_4^*\), vanishes on \(\ker B_4\), and satisfies
\[
 A_{\rm cut}^{(1)}B_4^*
 =B_4^*\overline A
 \quad\hbox{on }W_4.                                       \tag{28.22}
\]
It is the unique operator with those three properties.  Every
operator
\[
 A^{(1)}=A_{\rm cut}^{(1)}
 +P_{\rm loop}UP_{\rm cut}
 +P_{\rm cut}VP_{\rm loop}
 +P_{\rm loop}WP_{\rm loop}                               \tag{28.23}
\]
has the same prescribed action (28.22) on exact edge fields precisely
when \(U=0\); the blocks \(V,W\) remain invisible to (28.22).
Accordingly, vertex incidence determines transport of exact
one-forms but not transport of loop curvature or a general edge
connection.
\end{theorem}

\begin{proof}
Equation (28.21) visibly factors through \(B_4\) and \(B_4^*\), so it
kills \(\ker B_4\) and lands in \(\operatorname{Ran}B_4^*\).  On
\(W_4\),
\[
 A_{\rm cut}^{(1)}B_4^*
 =\frac14B_4^*\overline A B_4B_4^*
 =B_4^*\overline A.
\]
If another operator has the stated range, kernel, and intertwining
property, its restriction to
\(\operatorname{Ran}B_4^*=\operatorname{Ran}(B_4^*|_{W_4})\)
is fixed by (28.22), and its restriction to the loop space is zero;
hence it equals (28.21).

For (28.23), applying it to \(B_4^*W_4\) adds the loop-valued block
\(P_{\rm loop}UP_{\rm cut}\).  Thus (28.22) forces \(U=0\).
The other two added blocks have domain in the loop space and are not
tested on exact edge fields.
\end{proof}

For coframes, metrics, gauge connections, and spinors, even the cut
lift is not the whole representation.  One additionally needs the
tensor representation, local frame rotation, internal representation,
connection transport, and, for spinors, a spin lift.  These data are
not contained in the bare incidence matrix.

\subsection{The correct endpoint-family interface}

The preceding fixed-cell construction is useful as a fallback, but
the source itself uses an actual cross-time endpoint map.  Let
\[
 \mathscr P_{s\beta,h}^{Q}:
 Q_h(t+h)\longrightarrow Q_h(t),\qquad
 \mathscr P_{0,h}^{Q}=I,                                  \tag{28.24}
\]
be a differentiable family in a declared field representation.  The
finite shift drift required in V27 is
\[
 \boxed{
 \mathcal D_{q,h}\beta
 :=-\left.\frac{d}{ds}\right|_{s=0}
   \mathscr P_{s\beta,h}^{Q}q.}                            \tag{28.25}
\]
The sign agrees with
\(\partial_0-\mathcal L_\beta\) in (28.1).

\begin{proposition}[Endpoint derivative suffices for the V27 row]
\label{prop:v14v28-endpoint}
If (28.24) is differentiable in operator norm, linear to first order
in \(\beta\), and supplied in the based common configuration space,
then (28.25) defines the linear map
\(\mathcal D_q:Y_{\beta,h}\to Q_h\) required by V27.  Its
\(q\)-derivative is obtained by differentiating the same family in
the background coordinate.  If the family is instead an exact
algebra-automorphism group on the fixed vertex scalar algebra, then
(28.25) vanishes there.
\end{proposition}

\begin{proof}
Operator-norm differentiability makes (28.25) a vector of \(Q_h\);
first-order linearity gives linearity in \(\beta\).  Differentiation
in \(q\) gives the second term in V27 equation (27.16).  In the
exact fixed-algebra case the derivative is a derivation, so Theorem
\ref{thm:v14v28-no-derivation} makes it zero.
\end{proof}

Thus a nonzero finite shift must enter through at least one of:
\[
 \begin{array}{ll}
 \text{(a)}&\text{a cross-time map between moving sampling spaces},\\
 \text{(b)}&\text{a nonmultiplicative finite transport with a controlled
                 Leibniz defect},\\
 \text{(c)}&\text{an enlarged differential calculus carrying edge and
                 higher-cell data}.
 \end{array}                                                \tag{28.26}
\]
The Grand Tensor endpoint theorem conditionally supplies the idea of
(a), but it does not display (28.24) in bases on the full
gravity--internal configuration space.

\subsection{Source comparison and next acquisition}

The shift row now separates into the following entries:
\[
\begin{array}{p{0.29\linewidth}|p{0.16\linewidth}|p{0.44\linewidth}}
\textbf{datum}&\textbf{status}&\textbf{result}\\ \hline
one physical vector port&\mathbf C&
rank-one vector-isotypic source theorem\\
tetrahedral identification&\mathbf E&
explicit isometry (28.5), after an ordered frame is declared\\
endpoint-to-edge lift&\mathbf D&
minimum-norm cut lift (28.15) under zero circulation\\
circulation sector&\mathbf M&
three-dimensional \(\ker B_4\), with witness (28.20)\\
vertex-scalar drift&\mathbf D&
explicit formula (28.19) under the preceding choices\\
exact finite Leibniz action&\text{impossible}&
Theorem \ref{thm:v14v28-no-derivation}\\
exact-edge transport&\mathbf D&
cut lift (28.21); loop blocks remain missing\\
full field endpoint family&\mathbf M&
matrices (28.24) for metric/coframe, gauge, scalar, and fermion fields\\
cofinal consistency&\mathbf C&
requires the scaled Leibniz defect (28.13) and endpoint-family
transport to vanish on fixed screens
\end{array}                                                 \tag{28.27}
\]

\begin{decision}[Direction after V28]
\label{dec:v14v28-next}
The \(K_4\) graph can populate a scalar shift drift only after a
tetrahedral frame and a zero-circulation lift are declared.  It cannot
populate the full V27 map \(\mathcal D_q\), because general edge and
spin fields require loop and representation blocks.

V29 will attack the internal temporal-gauge drift
\(\mathcal R_q\), where the finite weak and gauge representations are
more explicit than the gravitational endpoint family.  It will
derive the Gauss constraint and its derivative from the vertex and
edge group actions, identify the stabilizer relations exactly, and
test whether this populates the internal diagonal block
\(C_{i,h}\) without importing a fermion spectral gap.  The unresolved
shift circulation and full endpoint-family rows remain recorded.
\end{decision}

\subsection{Acquisition ledger}

\[
\begin{array}{p{0.30\linewidth}|p{0.20\linewidth}|p{0.39\linewidth}}
\textbf{gate}&\textbf{status after V28}&\textbf{content}\\ \hline
endpoint vector identification&conditional/explicit&
rank-one source port plus tetrahedral frame\\
cut shift lift&constructed&
\(B_4^\dagger=\frac14B_4^*\)\\
loop shift lift&open&
three circulation coordinates\\
scalar drift&constructed conditionally&
orientation-independent flux action\\
edge/coframe drift&partial&
exact-edge block only\\
full endpoint pullback derivative&open&
common based field representation\\
internal Gauss drift&next&
differentiate the finite gauge action\\
lapse and gravitational kinetic rows&open&
unchanged from V27
\end{array}                                                 \tag{28.28}
\]

\subsection{Parent-integrity record}

The installed parent was
\[
\begin{array}{c|c}
 \text{file}&
 \texttt{Renewal\_SM\_Einstein\_Continuum\_Research\_Notes\_V27.tex}\\
 \text{lines}&11337\\
 \text{bytes}&365745\\
 \text{SHA--256}&
 \texttt{2E89B550FB3AB61A25845D06AE8FF8B823B01ABAAA087B9B41D3829CDE9AE510}.
\end{array}                                                 \tag{28.29}
\]

% =============================================================================
% END APPENDED FOURTEENTH-CYCLE V28 MATERIAL
% =============================================================================
% =============================================================================
% BEGIN APPENDED FOURTEENTH-CYCLE V29 MATERIAL
% =============================================================================

\section{Fourteenth-cycle V29: the finite internal Gauss momentum map
and its vacuum relation floor}
\label{sec:v14v29}

\subsection{Scope}

The active Grand Tensor Standard-Model action supplies the finite
gauge transformations
\[
 U_e\longmapsto g_{t(e)}U_eg_{s(e)}^{-1},\qquad
 H_v\longmapsto\rho_H(g_v)H_v,\qquad
 \Psi_v\longmapsto\rho_\Psi(g_v)\Psi_v,                    \tag{29.1}
\]
and their infinitesimal/BRST form
\[
 sU_e=c_{t(e)}U_e-U_ec_{s(e)},\qquad
 sH_v=c_vH_v,\qquad s\Psi_v=c_v\Psi_v.                    \tag{29.2}
\]
The source derives the Ward and classical BRST identities from gauge
invariance.  It does not display the canonical temporal-gauge
multiplier row in the V27 phase coordinates.  This note constructs
that row for the ordinary cotangent gauge--Higgs sector and computes
its exact relation space at the invariant vacuum.

Let \(G=G_{\rm SM}\) be a compact matrix group with Lie algebra
\(\mathfrak g\) and a fixed \(\operatorname{Ad}\)-invariant positive
inner product.  Let \(X=(V,E)\) be a connected finite oriented graph;
the selected application is \(X=K_4\).  Put
\[
 \mathfrak g_X=\mathfrak g^V.                              \tag{29.3}
\]
For clarity the first construction uses the bosonic configuration
manifold
\[
 \mathcal Q_X=G^E\times\prod_{v\in V}V_{H,v}.              \tag{29.4}
\]
Fermions and their independent dual amplitudes are returned to the
ledger below because their natural finite phase space is graded and
first order.

\subsection{Infinitesimal action and canonical momentum map}

Right-trivialize an edge tangent by
\[
 \delta U_e\longmapsto \delta U_eU_e^{-1}\in\mathfrak g.
\]
For \(\xi=(\xi_v)_{v\in V}\in\mathfrak g_X\), differentiating (29.1)
gives the fundamental configuration vector
\[
 \boxed{
 (\mathcal R_q\xi)_e
 =\xi_{t(e)}-\operatorname{Ad}_{U_e}\xi_{s(e)},\qquad
 (\mathcal R_q\xi)_{H,v}=\rho_{H,*}(\xi_v)H_v.}            \tag{29.5}
\]
This is the V27 internal drift candidate.

Let \(T^*\mathcal Q_X\) have its canonical symplectic form.  In the
same trivialization write a phase point as
\[
 z=(U,H;E,\pi_H),                                         \tag{29.6}
\]
where \(E_e\in\mathfrak g\) and \(\pi_{H,v}\in V_{H,v}^*\).

\begin{theorem}[Finite Gauss momentum map]
\label{thm:v14v29-momentum}
The cotangent lift of (29.1) is Hamiltonian with momentum map
\[
 \Phi_{G,X}:T^*\mathcal Q_X\longrightarrow\mathfrak g_X
\]
defined by
\[
 \boxed{\langle\Phi_{G,X}(z),\xi\rangle
 =\langle(E,\pi_H),\mathcal R_q\xi\rangle.}                \tag{29.7}
\]
In vertex components,
\[
 \boxed{
 \begin{aligned}
 (\Phi_{G,X})_v
 ={}&\sum_{e:t(e)=v}E_e
 -\sum_{e:s(e)=v}\operatorname{Ad}_{U_e^{-1}}E_e
 +j_{H,v},\\
 \langle j_{H,v},\xi_v\rangle
 :={}&\langle\pi_{H,v},\rho_{H,*}(\xi_v)H_v\rangle .
 \end{aligned}}                                            \tag{29.8}
\]
Adding a temporal gauge multiplier \(a_0\in\mathfrak g_X\) to the
Hamiltonian as
\[
 \langle a_0,\Phi_{G,X}\rangle                             \tag{29.9}
\]
reproduces exactly the V27 drift term
\(\langle p,\mathcal R_q a_0\rangle\) and gives the Gauss equation
\(\Phi_{G,X}=0\).
\end{theorem}

\begin{proof}
The canonical one-form on a cotangent bundle is invariant under every
cotangent-lifted configuration action.  Contracting it with the
fundamental vector generated by \(\xi\) gives (29.7), the standard
cotangent momentum-map identity.  To compute its components, insert
(29.5):
\[
 \sum_e\langle E_e,\xi_{t(e)}\rangle
 -\sum_e\langle E_e,\operatorname{Ad}_{U_e}\xi_{s(e)}\rangle
 +\sum_v\langle\pi_{H,v},\rho_{H,*}(\xi_v)H_v\rangle.
\]
Ad-invariance changes the second edge pairing to
\(\langle\operatorname{Ad}_{U_e^{-1}}E_e,\xi_{s(e)}\rangle\).
Collecting coefficients of each independent \(\xi_v\) proves (29.8).
Equation (29.7) makes (29.9) identical to the last term in the V27
Hamiltonian (27.9); variation in \(a_0\) gives the stated constraint.
\end{proof}

\subsection{Derivative, adjoint, and exact relation space}

Let
\[
 C_{G,z}:=D_z\Phi_{G,X}:T_z(T^*\mathcal Q_X)
                    \longrightarrow\mathfrak g_X.         \tag{29.10}
\]

\begin{proposition}[Weak derivative and relation criterion]
\label{prop:v14v29-derivative}
For a phase variation \(u=(\eta,\pi)\),
\[
 \boxed{
 \langle C_{G,z}u,\xi\rangle
 =\langle\pi,\mathcal R_q\xi\rangle
  +\langle p,(D_q\mathcal R[\eta])\xi\rangle.}             \tag{29.11}
\]
Consequently
\[
 \boxed{
 \ker C_{G,z}^*
 =\left\{\xi\in\mathfrak g_X:
 \mathcal R_q\xi=0,\quad
 \langle p,(D_q\mathcal R[\eta])\xi\rangle=0
 \ \hbox{for every }\eta\right\}.}                         \tag{29.12}
\]
This is the infinitesimal stabilizer of the full cotangent phase
point.  At zero momentum,
\[
 \ker C_{G,(q,0)}^*=\ker\mathcal R_q,\qquad
 C_{G,(q,0)}C_{G,(q,0)}^*=\mathcal R_q^*\mathcal R_q.      \tag{29.13}
\]
\end{proposition}

\begin{proof}
Differentiate (29.7) simultaneously in configuration and momentum.
The momentum variation gives the first term of (29.11), and the
configuration variation of the fundamental vector gives the second.
A label belongs to \(\ker C_{G,z}^*\) precisely when the right side
vanishes for every independently variable \(\pi\) and \(\eta\).
The \(\pi\) variations give \(\mathcal R_q\xi=0\); the \(\eta\)
variations give the second condition in (29.12).  If \(p=0\), the
second condition is automatic and
\[
 C_{G,(q,0)}(\eta,\pi)=\mathcal R_q^*\pi.
\]
Multiplying by its adjoint proves the Gram identity in (29.13).
\end{proof}

Thus the relation projector is a stabilizer projector at the chosen
phase point.  It is not determined by the abstract gauge group alone
and may jump across orbit types.

\subsection{First-class bracket}

\begin{theorem}[Exact finite Gauss closure]
\label{thm:v14v29-first-class}
For \(\xi,\eta\in\mathfrak g_X\), set
\(\Phi_\xi=\langle\Phi_{G,X},\xi\rangle\).  With the sign convention
in (29.7),
\[
 \boxed{\{\Phi_\xi,\Phi_\eta\}=\Phi_{[\xi,\eta]}}           \tag{29.14}
\]
where the bracket on \(\mathfrak g_X\) is vertexwise.  Hence the Gauss
constraint bank is first class and
\[
 C_{G,z}\mathbb J_zC_{G,z}^*=0                             \tag{29.15}
\]
at every point of \(\Phi_{G,X}^{-1}(0)\).
\end{theorem}

\begin{proof}
The Hamiltonian vector field of \(\Phi_\xi\) is the cotangent lift of
the fundamental vector field \(\mathcal R\xi\).  Fundamental vector
fields for a left action satisfy
\([(\mathcal R\xi)^\#,(\mathcal R\eta)^\#]
=(\mathcal R[\xi,\eta])^\#\).  Therefore the Hamiltonian vector fields
of the two sides of (29.14) agree.  Both sides vanish on the zero
section, so their difference is the zero constant.  At a zero of the
momentum map, (29.14) vanishes for all labels.  The matrix of these
brackets is exactly \(C_{G,z}\mathbb J_zC_{G,z}^*\), proving
(29.15).
\end{proof}

The source BRST identity is compatible with (29.14): replacing the
label by the ghost and using \(sc=-c^2\) is the graded encoding of the
same Lie bracket.  BRST nilpotence does not supply the Gram floor
computed next.

\subsection{The invariant \(K_4\) vacuum}

Take the source vacuum
\[
 U_e=I,\qquad H_v=v_Hh_0,\qquad E_e=0,\qquad\pi_{H,v}=0,   \tag{29.16}
\]
with \(v_H>0\).  Let
\[
 L_{h_0}:\mathfrak g\longrightarrow V_H,\qquad
 L_{h_0}\xi=\rho_{H,*}(\xi)(v_Hh_0),                       \tag{29.17}
\]
and
\[
 \mathfrak h:=\ker L_{h_0}.
\]
For the active Standard-Model representation the source identifies
the stabilizer group with
\[
 H\cong U(3),\qquad \dim\mathfrak h=9,\qquad
 \dim\mathfrak h^\perp=3.                                 \tag{29.18}
\]

\begin{theorem}[Exact vacuum relation projector on \(K_4\)]
\label{thm:v14v29-relations}
At (29.16),
\[
 \boxed{
 \ker C_{G,0}^*
 =\left\{(\xi,\xi,\xi,\xi):\xi\in\mathfrak h\right\}.}     \tag{29.19}
\]
With the product inner product on \(\mathfrak g^4\), its orthogonal
projector is
\[
 \boxed{
 (\Pi_{G,0}\zeta)_v
 =P_{\mathfrak h}\left(\frac14\sum_{x=1}^4\zeta_x\right)
 \quad(v=1,\ldots,4).}                                    \tag{29.20}
\]
Thus the local gauge label space has dimension \(48\), the declared
relation rank is \(9\), and the linearized bosonic Gauss bank has rank
\(39\).
\end{theorem}

\begin{proof}
At identity links, the edge component of (29.5) is
\[
 \xi_{t(e)}-\xi_{s(e)}.
\]
It vanishes on every edge of the connected graph \(K_4\) exactly when
all four vertex labels equal one common \(\xi\).  The Higgs component
then vanishes exactly when \(L_{h_0}\xi=0\), namely
\(\xi\in\mathfrak h\).  Equation (29.13) identifies this stabilizer
with \(\ker C_{G,0}^*\), proving (29.19).  Orthogonal projection first
takes the vertex average and then its \(\mathfrak h\)-component,
which proves (29.20).  Finally
\(\dim G_{\rm SM}=12\), so
\(\dim\mathfrak g^4=48\); rank--nullity gives \(39\).
\end{proof}

The nine relations are global unbroken gauge generators at this
specific vacuum.  They are not nine missing constraints, and they
need not persist at a generic Higgs and momentum background.

\subsection{A completed internal Gauss floor}

Let the edge tangent norm at (29.16) use weights \(w_e>0\), and the
Higgs tangent norm use weights \(m_v>0\).  Put
\[
 w_{\min}=\min_e w_e,\qquad m_{\min}=\min_v m_v.            \tag{29.21}
\]
Let
\[
 M_H=L_{h_0}^*L_{h_0}.
\]
Since \(\ker M_H=\mathfrak h\), define its positive orbit floor
\[
 \nu_H=\lambda_{\min}
 \left(M_H\big|_{\mathfrak h^\perp}\right)>0.              \tag{29.22}
\]

\begin{theorem}[Bosonic \(K_4\) Gauss completed-Gram certificate]
\label{thm:v14v29-floor}
At the vacuum (29.16), use the weighted version of
\(\mathcal R_{q_0}^*\mathcal R_{q_0}\) and the relation projector
(29.20).  Then
\[
 \boxed{
 C_{G,0}C_{G,0}^*+\Pi_{G,0}
 \succeq
 \delta_{G,0}I_{\mathfrak g^4},\qquad
 \delta_{G,0}:=
 \min\{4w_{\min},m_{\min}\nu_H,1\}>0.}                    \tag{29.23}
\]
In particular,
\[
 \operatorname{rmin}C_{G,0}\ge\sqrt{\delta_{G,0}},
 \qquad
 \|C_{G,0}^\dagger\|\le\delta_{G,0}^{-1/2}.                \tag{29.24}
\]
\end{theorem}

\begin{proof}
Write a label uniquely as
\[
 \zeta=\mathbf1\otimes\bar\zeta+\zeta_W,\qquad
 \bar\zeta=\frac14\sum_v\zeta_v,\qquad
 \sum_v(\zeta_W)_v=0.                                     \tag{29.25}
\]
The identity-link edge action is
\((B_4^*\otimes I_{\mathfrak g})\zeta\).  V26 gives
\[
 \sum_e w_e\|\zeta_{t(e)}-\zeta_{s(e)}\|^2
 \ge4w_{\min}\|\zeta_W\|^2.                               \tag{29.26}
\]
The Higgs action gives
\[
 \sum_vm_v\|L_{h_0}\zeta_v\|^2
 \ge m_{\min}\nu_H
 \sum_v\|P_{\mathfrak h^\perp}\zeta_v\|^2.                \tag{29.27}
\]
Finally,
\[
 \langle\zeta,\Pi_{G,0}\zeta\rangle
 =\|\mathbf1\otimes P_{\mathfrak h}\bar\zeta\|^2.          \tag{29.28}
\]
The three mutually orthogonal components
\[
 \zeta_W,\qquad
 \mathbf1\otimes P_{\mathfrak h^\perp}\bar\zeta,\qquad
 \mathbf1\otimes P_{\mathfrak h}\bar\zeta
\]
are controlled respectively by (29.26), (29.27), and (29.28).
Taking their least coefficient proves (29.23).  Apply the V16
kernel-completion theorem, or restrict (29.23) to
\((\ker C_{G,0}^*)^\perp\), to obtain (29.24).
\end{proof}

The formula is symbolic but finite and executable.  The graph
constant \(4\) is exact.  The physical values of \(w_{\min}\),
\(m_{\min}\), and \(\nu_H\), and their uniform cofinal behavior, must
still be read from the selected action and Hilbert normalization.

\begin{corollary}[Exact internal angle at the vacuum]
\label{cor:v14v29-angle}
For the bosonic Gauss block at (29.16), the Hamiltonian gauge
generator is \(K_{G,0}=\mathbb J C_{G,0}^*\).  Its range lies in
\(\ker C_{G,0}\), and the V26 physical--gauge angle is exactly one.
\end{corollary}

\begin{proof}
The vacuum satisfies \(\Phi_{G,X}=0\).  The first-class identity
(29.15) gives \(C_{G,0}K_{G,0}=0\).  Apply V26
Theorem \ref{thm:v14v26-ledger} with anomaly debit zero and the floor
(29.23).
\end{proof}

\subsection{What has and has not been populated}

The construction uses no eigenvalue of the finite Dirac operator.
That spectral gap has the wrong type for (29.23).  The exact internal
status is
\[
\begin{array}{p{0.29\linewidth}|p{0.17\linewidth}|p{0.43\linewidth}}
\textbf{datum}&\textbf{status}&\textbf{content}\\ \hline
finite gauge action on \(U,H\)&\mathbf E&
source transformations (29.1)--(29.2)\\
bosonic drift \(\mathcal R_q\)&\mathbf D&
explicit formula (29.5) after a spatial cotangent split\\
Gauss momentum map&\mathbf D&
equations (29.7)--(29.8)\\
Gauss first-class bracket&\mathbf D&
equivariant momentum-map identity (29.14)\\
vacuum relation projector&\mathbf D&
diagonal \(U(3)\) projector (29.20)\\
vacuum bosonic Gram floor&\mathbf D&
formula (29.23), with physical weights still to be inserted\\
graded fermion phase row&\mathbf M&
the source uses chiral fields and independent dual amplitudes, not the
ordinary real cotangent chart (29.4)\\
temporal/spatial action split&\mathbf M&
the Euclidean all-link action is explicit, but its time-link Legendre
reduction to (29.9) is not displayed in one based card\\
gravity-dependent whitening&\mathbf M&
the common Hodge metric can create mixed derivatives after Riesz
identification\\
generic-background relation floor&\mathbf M&
stabilizer and orbit Gram must be recomputed away from (29.16)
\end{array}                                                 \tag{29.29}
\]

Thus V29 populates a genuine internal diagonal constraint card on the
bosonic invariant-vacuum branch.  It does not yet populate the full
\(C_{i,h}\) for the chiral phase space, nor either mixed block
\(M_h,N_h\).

\begin{decision}[Direction after V29]
\label{dec:v14v29-next}
V30 is the compulsory companion audit.  After inspecting the newest
Yang--Mills and Navier--Stokes notes, V30 will compare their advances
against the new Gauss orbit-Gram card.  Unless a companion supplies a
same-carrier temporal link reduction, the primary route will remain:
\[
 \text{spacetime time-link split}
 \longrightarrow
 \text{graded gauge--fermion phase chart}
 \longrightarrow
 \text{gravity-dependent whitening}
 \longrightarrow
 \text{full coupled Schur card}.                           \tag{29.30}
\]
V30 will then select the next substantive upstream row; it will not
replace (29.23) by a Yang--Mills Wilson-shell or fermion gap.
\end{decision}

\subsection{Acquisition ledger}

\[
\begin{array}{p{0.30\linewidth}|p{0.20\linewidth}|p{0.39\linewidth}}
\textbf{gate}&\textbf{status after V29}&\textbf{content}\\ \hline
bosonic internal Gauss map&constructed&
cotangent momentum map\\
vacuum relation bank&proved on the bosonic vacuum&
diagonal unbroken \(U(3)\)\\
vacuum internal floor&proved conditionally&
\(\min\{4w_{\min},m_{\min}\nu_H,1\}\)\\
internal first-class angle&proved on the bosonic vacuum&
zero bracket debit on the constraint surface\\
fermion Gauss phase row&open&
graded symplectic/Berezin chart\\
full shift drift&open&
loop and endpoint-family blocks from V28\\
lapse/gravity kinetic row&open&
V27 inverse-lapse incidence\\
companion checkpoint&next&
YM and NS audit at V30
\end{array}                                                 \tag{29.31}
\]

\subsection{Parent-integrity record}

The installed parent was
\[
\begin{array}{c|c}
 \text{file}&
 \texttt{Renewal\_SM\_Einstein\_Continuum\_Research\_Notes\_V28.tex}\\
 \text{lines}&11841\\
 \text{bytes}&383556\\
 \text{SHA--256}&
 \texttt{E1F5A33B355210167FAE6183612236B5F267E1AE5EA8A7F58D817C2404F35D77}.
\end{array}                                                 \tag{29.32}
\]

% =============================================================================
% END APPENDED FOURTEENTH-CYCLE V29 MATERIAL
% =============================================================================

% =============================================================================
% BEGIN APPENDED FOURTEENTH-CYCLE V30 MATERIAL
% =============================================================================

\section{Fourteenth-cycle V30: compulsory companion audit and the
moving-Gauss-floor direction}
\label{sec:v14v30}

\subsection{Files inspected}

The required thirtieth-version audit was performed read-only on
\[
\begin{array}{p{0.51\linewidth}|r|r}
\textbf{file}&\textbf{lines}&\textbf{bytes}\\ \hline
\texttt{Renewal\_Yang\_Mills\_Mass\_Gap\_Research\_Notes\_V9.tex}
 &2260&81930\\
\texttt{Renewal\_NS\_Stability\_Research\_Notes\_V26.tex}
 &5319&278283.
\end{array}                                                  \tag{30.1}
\]
Their SHA--256 digests are
\[
\begin{aligned}
 &\texttt{E357F84CA31644B4C499CAD1235F3372713083D8631D4CAA72CAF590F1B30A84},
                                                               \tag{30.2}\\
 &\texttt{82C887F8C2C69E4F28FAD7C3DC8DE5B9099CB5A4BF431EBA1FFF3E91935978AD}.
                                                               \tag{30.3}
\end{aligned}
\]
No companion file was modified or removed.

\subsection{Yang--Mills V8--V9}

At the V25 audit, the Yang--Mills programme had reached V7 and had
eliminated two momentum coordinates exactly at the remaining
quarter-turn grid.  V8 has since proved, by exact Gaussian-integer
negative witnesses, that the naive degree-one interpolation of the
second family of Riccati splittings is indefinite for both the Wilson
and endpoint actions.  The witnesses refute that auxiliary
interpolation, not the original compressed pencil.  Pointwise
optimization and restoration of full cardinal functions do not give
one stable common repair.

V9 then goes upstream to a factor
\[
 H=\mathcal A^*\mathcal A,\qquad
 K=\mathcal A^*\mathcal D\mathcal A                         \tag{30.4}
\]
and defines
\[
 S=\frac{K+K^*}{2},\qquad
 \Omega=\frac12H^{-1}(K-K^*),\qquad
 \mathcal B=\mathcal D\mathcal A-\mathcal A\Omega.          \tag{30.5}
\]
It proves exactly that \(\Omega\) is \(H\)-skew,
\(\mathcal A^*\mathcal B=S\), and
\[
 2qH\pm\nabla H\succeq0
 \quad\Longleftrightarrow\quad
 -qH\preceq S\preceq qH
 \quad\Longleftrightarrow\quad
 SH^{-1}S\preceq q^2H.                                    \tag{30.6}
\]
Thus the normal factor derivative and its \(H\)-skew tangential part
do not control the Gram variation.  V9 also proves exact positive
quarter-turn-grid floors for the stronger \(q=2/3\) pencil:
approximately \(0.0070406973\) for Wilson and \(0.2075386118\) for
the endpoint action.  Positivity between grid points remains open.

\subsection{Import into the SM--Einstein route}

The algebraic cancellation in (30.5)--(30.6) applies to V29 with
\[
 \mathcal A=\mathcal R_q,\qquad
 H=\mathcal R_q^*\mathcal R_q
 =C_{G,(q,0)}C_{G,(q,0)}^*                                \tag{30.7}
\]
on the complement of the stabilizer relation space.  A moving gauge
frame can make \(\|\mathcal D\mathcal R_q\|\) large even when the
orbit Gram changes slowly.  Bounding the full derivative would
therefore charge a coordinate artifact.  The correct transport
quantity is the Hermitian part of
\(\mathcal R_q^*\mathcal D\mathcal R_q\) in the moving quotient
metric.

This import does not transfer either Yang--Mills numerical floor.
Its factor, fibre, momentum variables, and constraint compression are
different from the \(K_4\) Standard-Model Gauss orbit map.  It also
does not permit use of \(H^{-1}\) before a relation quotient and a
positive floor have been established.  V29 supplies such a floor at
one bosonic vacuum; V31 must transport it without circularity.

\subsection{Navier--Stokes V26}

The Navier--Stokes companion is unchanged.  It computes exact
rank-one refinements of first and second Oseen-kernel derivative
norms.  The first gives no improvement; the second gives a modest
localized improvement, with no global regularity conclusion.  Its
main lesson remains that a norm tailored to the actual structured
input can be strictly smaller than the unrestricted operator norm,
but the reduction must be proved on that input class.

There is no Navier--Stokes Gauss action, stabilizer relation, affine
multiplier chart, or constraint Gram to import.  The structured-input
lesson supports using the Hermitian quotient derivative in
(30.6), rather than a larger unrelated factor norm.

\begin{decision}[Direction after V30]
\label{dec:v14v30-next}
V31 will prove a moving-relation Gauss-floor transport theorem.  Along
a \(C^1\) background path with constant stabilizer dimension it will:

\begin{enumerate}
\item use Kato parallel transport to freeze the moving relation
projector;
\item restrict the Gauss factor to the fixed relation-free quotient;
\item remove the quotient-metric skew connection as in (30.5);
\item derive a two-sided Gronwall comparison for the orbit Gram from
a bound on its Hermitian logarithmic derivative; and
\item propagate the V29 vacuum floor over an explicit path length.
\end{enumerate}

The theorem will fail with a named witness when the stabilizer
dimension jumps.  It will not use the YM quarter-grid floors or the
NS rank-one constants as Standard-Model data.
\end{decision}

\subsection{Acquisition ledger}

\[
\begin{array}{p{0.30\linewidth}|p{0.20\linewidth}|p{0.39\linewidth}}
\textbf{gate}&\textbf{status after V30}&\textbf{content}\\ \hline
companion checkpoint&acquired&
YM V9 and NS V26 read with hashes\\
naive splitting interpolation&ruled out in YM&
exact auxiliary-block witnesses\\
skew/normal cancellation&imported structurally&
Hermitian Gram derivative only\\
vacuum Gauss floor&retained&
V29 symbolic \(K_4\) certificate\\
moving relation quotient&next&
Kato transport under constant rank\\
generic-background Gauss floor&next&
Gronwall propagation from the vacuum\\
graded fermion and gravitational rows&open&
unchanged
\end{array}                                                 \tag{30.8}
\]

The next compulsory companion audit is V35.

\subsection{Parent-integrity record}

The installed parent was
\[
\begin{array}{c|c}
 \text{file}&
 \texttt{Renewal\_SM\_Einstein\_Continuum\_Research\_Notes\_V29.tex}\\
 \text{lines}&12304\\
 \text{bytes}&399332\\
 \text{SHA--256}&
 \texttt{427A937C1C0509A11B6A6614B9B700ED6DDEF990B621CDF2CD3F201D94F6FF07}.
\end{array}                                                 \tag{30.9}
\]

% =============================================================================
% END APPENDED FOURTEENTH-CYCLE V30 MATERIAL
% =============================================================================
% =============================================================================
% BEGIN APPENDED FOURTEENTH-CYCLE V31 MATERIAL
% =============================================================================

\section{Fourteenth-cycle V31: transport of a Gauss floor on a
constant orbit-type stratum}
\label{sec:v14v31}

\subsection{Moving relations must be frozen before differentiating}

Let \(Y,Z\) be fixed finite-dimensional Hilbert spaces and let
\[
 C(s):Z\longrightarrow Y,\qquad 0\le s\le L,              \tag{31.1}
\]
be a \(C^1\) family of constraint derivatives.  Assume for now that
\(\operatorname{rank}C(s)\) is constant and that the orthogonal
relation projector
\[
 \Pi(s)=P_{\ker C(s)^*},\qquad Q(s)=I-\Pi(s)               \tag{31.2}
\]
is \(C^1\).  The completed Gram is
\[
 \widehat G(s)=C(s)C(s)^*+\Pi(s).                          \tag{31.3}
\]
Differentiating (31.3) directly mixes physical change of
\(C(s)C(s)^*\) with rotation of the relation subspace.  The first step
is therefore to identify all relation-free quotients with one fixed
space.

\begin{lemma}[Kato freezing of the relation projector]
\label{lem:v14v31-kato}
Let \(U(s)\) solve
\[
 \dot U(s)=[\dot\Pi(s),\Pi(s)]U(s),\qquad U(0)=I.           \tag{31.4}
\]
Then \(U(s)\) is orthogonal and
\[
 U(s)^*\Pi(s)U(s)=\Pi(0),\qquad
 U(s)^*Q(s)U(s)=Q(0).                                     \tag{31.5}
\]
\end{lemma}

\begin{proof}
Both \(\Pi\) and \(\dot\Pi\) are self-adjoint, so
\([\dot\Pi,\Pi]^*=-[\dot\Pi,\Pi]\).  Equation (31.4) therefore gives
\[
 \frac d{ds}(U^*U)=0,
\]
and \(U(0)=I\) makes \(U\) orthogonal.  Differentiating
\(\Pi^2=\Pi\) gives
\[
 \Pi\dot\Pi\Pi=0,\qquad
 \dot\Pi=\dot\Pi\Pi+\Pi\dot\Pi.
\]
Consequently
\[
 [[\dot\Pi,\Pi],\Pi]=\dot\Pi.                              \tag{31.6}
\]
Using \(\dot U=[\dot\Pi,\Pi]U\) and (31.6),
\[
 \frac d{ds}(U^*\Pi U)
 =U^*\bigl(\dot\Pi+[\Pi,[\dot\Pi,\Pi]]\bigr)U=0.
\]
This proves the first identity in (31.5); the second follows by
taking complements.
\end{proof}

Put \(Y_0=Q(0)Y\) and define the injective quotient factor
\[
 A(s):Y_0\longrightarrow Z,\qquad
 A(s)=C(s)^*U(s)\big|_{Y_0}.                               \tag{31.7}
\]
Its Gram is
\[
 H(s)=A(s)^*A(s)
 =Q(0)U(s)^*C(s)C(s)^*U(s)Q(0).                           \tag{31.8}
\]
The nonzero spectrum of \(C(s)C(s)^*\) is exactly the spectrum of
\(H(s)\).  On the frozen relation space \(\Pi(0)Y\), the transported
completed Gram equals the identity.

\subsection{Removal of the quotient-metric skew connection}

Set
\[
 K(s)=A(s)^*\dot A(s),\qquad
 S(s)=\frac{K(s)+K(s)^*}{2}.                              \tag{31.9}
\]
Then
\[
 \dot H(s)=K(s)+K(s)^*=2S(s).                             \tag{31.10}
\]
Where \(H(s)>0\), define
\[
 \Omega(s)=\frac12H(s)^{-1}\bigl(K(s)-K(s)^*\bigr),
 \qquad
 B(s)=\dot A(s)-A(s)\Omega(s).                            \tag{31.11}
\]

\begin{theorem}[Quotient skew-removal identity]
\label{thm:v14v31-skew}
For every \(s\) on the constant-rank stratum,
\[
 \Omega^*H+H\Omega=0,\qquad
 A^*B=S=S^*,\qquad
 B^*A+A^*B=\dot H.                                        \tag{31.12}
\]
For every \(q\ge0\), the following conditions are equivalent:
\[
\begin{array}{ll}
\textup{(i)}&-2qH\preceq\dot H\preceq2qH,\\
\textup{(ii)}&-qH\preceq S\preceq qH,\\
\textup{(iii)}&SH^{-1}S\preceq q^2H,\\
\textup{(iv)}&B^*P_AB\preceq q^2A^*A,
\end{array}                                                \tag{31.13}
\]
where \(P_A=A(A^*A)^{-1}A^*\) is the range projector.  Thus neither
the normal part of \(\dot A\) nor its \(H\)-skew tangential part enters
the relative Gram derivative.
\end{theorem}

\begin{proof}
Since \(H\Omega=(K-K^*)/2\),
\[
 \Omega^*H=(K^*-K)/2,
\]
which gives the first identity in (31.12).  Also
\[
 A^*B=K-H\Omega=(K+K^*)/2=S.
\]
The last identity in (31.12) follows from (31.10), or directly from
the preceding equation.

Conditions (i) and (ii) are identical by (31.10).  Congruence by
\(H^{-1/2}\) turns (ii) into
\[
 -qI\preceq T\preceq qI,\qquad
 T=H^{-1/2}SH^{-1/2}=T^*.
\]
By the spectral theorem this is equivalent to
\(T^2\preceq q^2I\), which is (iii) after undoing the congruence.
Finally
\[
 B^*P_AB
 =B^*A H^{-1}A^*B
 =SH^{-1}S,                                               \tag{31.14}
\]
so (iii) and (iv) agree.
\end{proof}

This is the V30 Yang--Mills cancellation applied to the actual
moving constraint quotient.  Formula (31.13) is structural; using
\(H^{-1}\) in (iii)--(iv) does not by itself prove a floor.  Conditions
(i)--(ii), by contrast, can be certified as inverse-free Loewner
inequalities.

\subsection{Relative and absolute floor propagation}

\begin{theorem}[Constant-stratum completed-Gram transport]
\label{thm:v14v31-transport}
Suppose
\[
 H(0)\succeq d_0I_{Y_0},\qquad d_0>0,                      \tag{31.15}
\]
and an integrable function \(q:[0,L]\to[0,\infty)\) satisfies the
inverse-free inequalities
\[
 -q(s)H(s)\preceq S(s)\preceq q(s)H(s)                    \tag{31.16}
\]
for every \(s\).  Put
\[
 Q_s=\int_0^s q(t)\,dt.
\]
Then
\[
 \boxed{
 e^{-2Q_s}H(0)\preceq H(s)\preceq e^{2Q_s}H(0)}            \tag{31.17}
\]
and
\[
 \boxed{
 \widehat G(s)\succeq
 \delta(s)I_Y,\qquad
 \delta(s)=\min\{1,e^{-2Q_s}d_0\}.}                        \tag{31.18}
\]
Consequently
\[
 \|C(s)^\dagger\|\le e^{Q_s}d_0^{-1/2}.                   \tag{31.19}
\]
\end{theorem}

\begin{proof}
For fixed \(y\in Y_0\), let
\(f_y(s)=\langle y,H(s)y\rangle\).  Equations
(31.10) and (31.16) give
\[
 -2q(s)f_y(s)\le\dot f_y(s)\le2q(s)f_y(s).                \tag{31.20}
\]
Scalar Gronwall applied to both inequalities gives
\[
 e^{-2Q_s}f_y(0)\le f_y(s)\le e^{2Q_s}f_y(0).
\]
Since this holds for every \(y\), it is exactly (31.17).  The lower
bound in (31.15) yields the quotient floor
\(e^{-2Q_s}d_0\).  Conjugation by the orthogonal Kato transport
preserves spectra, and the completed Gram is the identity on the
relation space.  This proves (31.18).  Its restriction to the
nonzero singular sector of \(C(s)\) proves (31.19).
\end{proof}

When a relative certificate is not yet available, there is a
noncircular additive fallback.

\begin{proposition}[Absolute symmetric-derivative debit]
\label{prop:v14v31-absolute}
If (31.15) holds and
\[
 \|S(s)\|\le\sigma(s),                                    \tag{31.21}
\]
then
\[
 H(s)\succeq
 \left(d_0-2\int_0^s\sigma(t)\,dt\right)I_{Y_0}.           \tag{31.22}
\]
In particular the right side is a positive quotient floor whenever
the parenthesis is positive.
\end{proposition}

\begin{proof}
For fixed \(y\),
\[
 \frac d{ds}\langle y,H(s)y\rangle
 =2\langle y,S(s)y\rangle
 \ge-2\sigma(s)\|y\|^2.
\]
Integrate and apply (31.15).
\end{proof}

Unlike a bound on \(\|\dot A\|\), (31.21) has already discarded
normal factor motion and quotient-metric skew rotation.  Unlike the
relative form (31.16), it requires no inverse and no bootstrap floor.

\subsection{Application to the \(K_4\) Gauss vacuum}

Let \(q(s)\) be a \(C^1\) path of bosonic gauge--Higgs configurations
and, initially, zero momenta.  Put
\[
 C_G(s)=D\Phi_G(q(s),0),\qquad
 \Pi_G(s)=P_{\ker C_G(s)^*}.                               \tag{31.23}
\]
At \(s=0\), take the V29 invariant \(K_4\) vacuum.  Its quotient floor
is
\[
 d_{G,0}=\min\{4w_{\min},m_{\min}\nu_H\}.                  \tag{31.24}
\]

\begin{corollary}[Gauss orbit-stratum radius]
\label{cor:v14v31-gauss}
Assume along \(0\le s\le L\):

\begin{enumerate}
\item the cotangent gauge stabilizer has constant dimension nine;
\item \(\Pi_G(s)\) is \(C^1\);
\item after Kato freezing, the Hermitian factor derivative obeys
\[
 -q_G(s)H_G(s)\preceq S_G(s)\preceq q_G(s)H_G(s).
                                                               \tag{31.25}
\]
\end{enumerate}

Then
\[
 \boxed{
 C_G(s)C_G(s)^*+\Pi_G(s)
 \succeq
 \min\left\{1,
 e^{-2\int_0^s q_G(t)\,dt}d_{G,0}\right\}I.}              \tag{31.26}
\]
The internal Gauss bracket remains exactly first class on the zero
set of its momentum map, so its physical--gauge angle remains one.
\end{corollary}

\begin{proof}
The V29 relation theorem supplies the initial relation space and
(31.24).  Apply Theorem \ref{thm:v14v31-transport} to the Gauss
constraint derivative.  Equivariance of the momentum map is
independent of the chosen orbit-type coordinates and gives zero
bracket debit on the constraint surface, as in V29.
\end{proof}

This corollary replaces an unstructured small-neighbourhood
continuity claim by an accumulated, basis-covariant path debit.  To
execute it, one must still populate either the Loewner pair (31.25)
or the absolute symmetric derivative in (31.21) from the same
gauge--Higgs action.

\subsection{Why a stabilizer jump cannot be crossed silently}

\begin{countertheorem}[Completed floors can collapse at a rank jump]
\label{cth:v14v31-rank-jump}
Let \(Y=Z=\mathbb R\) and
\[
 C(s)=s.
\]
Then
\[
 \Pi(0)=1,\qquad \Pi(s)=0\quad(s\ne0),
\]
and every individual completed Gram is positive:
\[
 \widehat G(0)=1,\qquad \widehat G(s)=s^2\quad(s\ne0).      \tag{31.27}
\]
Nevertheless
\[
 \inf_{0<|s|\le L}\lambda_{\min}\widehat G(s)=0.           \tag{31.28}
\]
\end{countertheorem}

\begin{proof}
At \(s=0\), \(C(0)^*=0\), so the relation space is all of
\(\mathbb R\).  At every nonzero \(s\), the adjoint is injective and
the relation space is zero.  Formula (31.27) follows directly, and
(31.28) follows by \(s\to0\).
\end{proof}

The V29 invariant vacuum is an enhanced-stabilizer point.  A generic
link holonomy can reduce its diagonal \(U(3)\) stabilizer.  The nine
vacuum relations then become newly active constraint directions, and
their singular values begin at zero.  Therefore (31.26) applies only
inside the constant orbit-type stratum.  It cannot establish a
uniform generic-background floor by continuity from the vacuum.

\subsection{Next upstream move}

\begin{decision}[Direction after V31]
\label{dec:v14v31-next}
The transport theorem closes the constant-rank analytic step, but the
vacuum is not a safe universal seed.  V32 will avoid crossing orbit
types.  For arbitrary link variables and Higgs values on connected
\(K_4\), it will:

\begin{enumerate}
\item identify the Gauss relation algebra as the intersection of a
root holonomy centralizer with the transported Higgs stabilizers;
\item express the orbit Gram as a weighted covariant graph Laplacian
plus a vertex Higgs orbit mass;
\item prove a rooted covariant Poincare inequality from a spanning
tree and one anchor mass; and
\item obtain a direct completed-Gram floor on each chosen background,
without transporting through the enhanced vacuum.
\end{enumerate}

This is the correct upstream response to the rank-jump witness.  The
fermion and gravity-dependent whitening rows will remain separate
until the bosonic covariant graph floor is explicit.
\end{decision}

\subsection{Acquisition ledger}

\[
\begin{array}{p{0.30\linewidth}|p{0.20\linewidth}|p{0.39\linewidth}}
\textbf{gate}&\textbf{status after V31}&\textbf{content}\\ \hline
moving relation frame&proved&
Kato unitary (31.4)\\
skew/normal removal&proved&
quotient identity (31.12)--(31.14)\\
relative floor transport&proved&
exponential bound (31.17)--(31.19)\\
inverse-free fallback&proved&
additive symmetric debit (31.22)\\
vacuum-stratum Gauss floor&conditional&
path certificate (31.26)\\
rank-changing transport&ruled out&
scalar witness (31.27)--(31.28)\\
generic bosonic Gauss floor&next&
covariant graph Laplacian plus anchor\\
graded fermion and coupled gravity rows&open&
unchanged
\end{array}                                                 \tag{31.29}
\]

\subsection{Parent-integrity record}

The installed parent was
\[
\begin{array}{c|c}
 \text{file}&
 \texttt{Renewal\_SM\_Einstein\_Continuum\_Research\_Notes\_V30.tex}\\
 \text{lines}&12472\\
 \text{bytes}&405909\\
 \text{SHA--256}&
 \texttt{C54741CCD5BD76F32A44391244F59AEEB8D74D0B3074E7BFB05967640087A2D3}.
\end{array}                                                 \tag{31.30}
\]

% =============================================================================
% END APPENDED FOURTEENTH-CYCLE V31 MATERIAL
% =============================================================================
% =============================================================================
% BEGIN APPENDED FOURTEENTH-CYCLE V32 MATERIAL
% =============================================================================

\section{Fourteenth-cycle V32: a rooted covariant \(K_4\) Gauss floor
at an arbitrary bosonic background}
\label{sec:v14v32}

\subsection{Covariant graph and Higgs orbit operator}

Retain the compact Standard-Model gauge group \(G\), its Lie algebra
\(\mathfrak g\), and the \(\operatorname{Ad}\)-invariant Hilbert metric
of V29.  On the oriented \(K_4\), let
\[
 q=(U_e,H_v),\qquad
 L_v:\mathfrak g\longrightarrow V_{H,v},\quad
 L_v\xi=\rho_{H,*}(\xi)H_v.                               \tag{32.1}
\]
For a vertex gauge label \(\xi=(\xi_v)_{v=1}^4\), define
\[
 (d_U\xi)_e
 =\xi_{t(e)}-\operatorname{Ad}_{U_e}\xi_{s(e)},\qquad
 (L_H\xi)_v=L_v\xi_v.                                     \tag{32.2}
\]
With positive edge and vertex weights \(w_e,m_v\), the bosonic orbit
energy is
\[
 \mathcal E_q(\xi)
 =\sum_{e\in E_4}w_e\|(d_U\xi)_e\|^2
  +\sum_{v=1}^4m_v\|L_v\xi_v\|^2.                         \tag{32.3}
\]
This is the quadratic form of
\(\mathcal R_q^*\mathcal R_q\) in the weighted tangent metric.

Choose the star spanning tree
\[
 T=\{12,13,14\}
\]
rooted at vertex \(1\), and put
\[
 P_1=I,\qquad P_v=U_{1v}\quad(v=2,3,4).                   \tag{32.4}
\]
For the three chords \(ab\in\{23,24,34\}\), define the rooted
holonomies
\[
 \mathcal H_{ab}=P_b^{-1}U_{ab}P_a.                       \tag{32.5}
\]

\subsection{Exact relation algebra}

Define the root anchor
\[
 \mathfrak A_q:\mathfrak g\longrightarrow
 \left(\bigoplus_{ab\in\{23,24,34\}}\mathfrak g\right)
 \oplus\left(\bigoplus_{v=1}^4V_{H,v}\right)               \tag{32.6}
\]
by
\[
 \boxed{
 \mathfrak A_qx=
 \left(
 \bigl(\sqrt{w_{ab}}(I-\operatorname{Ad}_{\mathcal H_{ab}})x
       \bigr)_{ab},
 \bigl(\sqrt{m_v}L_v\operatorname{Ad}_{P_v}x\bigr)_v
 \right).}                                                 \tag{32.7}
\]
Let
\[
 \mathfrak h_q:=\ker\mathfrak A_q.                         \tag{32.8}
\]

\begin{theorem}[Holonomy--Higgs description of all Gauss relations]
\label{thm:v14v32-relations}
The kernel of the infinitesimal gauge action (32.2) is
\[
 \boxed{
 \ker\mathcal R_q
 =\left\{
 \bigl(\operatorname{Ad}_{P_v}x\bigr)_{v=1}^4:
 x\in\mathfrak h_q
 \right\}.}                                                \tag{32.9}
\]
Equivalently,
\[
 \boxed{
 \mathfrak h_q=
 \bigcap_{ab\in\{23,24,34\}}
 \ker(I-\operatorname{Ad}_{\mathcal H_{ab}})
 \ \cap\
 \bigcap_{v=1}^4
 \operatorname{Ad}_{P_v^{-1}}\ker L_v.}                   \tag{32.10}
\]
Thus the relation algebra is the intersection of the rooted holonomy
centralizer with the transported Higgs stabilizers.
\end{theorem}

\begin{proof}
If \(d_U\xi=0\) on the three tree edges, then
\[
 \xi_v=\operatorname{Ad}_{P_v}\xi_1\qquad(v=1,2,3,4).
                                                               \tag{32.11}
\]
For a chord \(a\to b\), substitute (32.11) in
\(\xi_b=\operatorname{Ad}_{U_{ab}}\xi_a\) and conjugate by
\(P_b^{-1}\).  The result is
\[
 x=\operatorname{Ad}_{P_b^{-1}U_{ab}P_a}x
   =\operatorname{Ad}_{\mathcal H_{ab}}x,\qquad x=\xi_1.
\]
The Higgs component vanishes precisely when
\(L_v\operatorname{Ad}_{P_v}x=0\) for every \(v\).  These are exactly
the kernel conditions in (32.7), proving (32.9)--(32.10).  The
converse follows by reversing the same substitutions.
\end{proof}

Let \(P_{\mathfrak h_q}\) be the orthogonal projector in
\(\mathfrak g\).  For \(\xi\in\mathfrak g^4\), put
\[
 \zeta_v=\operatorname{Ad}_{P_v^{-1}}\xi_v,\qquad
 \overline\zeta=\frac14\sum_{v=1}^4\zeta_v.                \tag{32.12}
\]

\begin{corollary}[Explicit relation projector]
\label{cor:v14v32-projector}
The orthogonal projector onto (32.9) is
\[
 \boxed{
 (\Pi_q\xi)_v
 =\operatorname{Ad}_{P_v}
 P_{\mathfrak h_q}\overline\zeta.}                         \tag{32.13}
\]
\end{corollary}

\begin{proof}
Every \(\operatorname{Ad}_{P_v}\) is orthogonal.  After the unitary
coordinate change \(\xi_v\mapsto\zeta_v\), the relation space is the
diagonal copy
\(\{(x,x,x,x):x\in\mathfrak h_q\}\).  Orthogonal projection onto it
takes the vertex average followed by
\(P_{\mathfrak h_q}\).  Transforming back gives (32.13).
\end{proof}

Equations (32.7) and (32.13) give an exact finite relation card at
every chosen bosonic background.  No continuation from the enhanced
vacuum is used.

\subsection{Tree defects and the root-anchor error map}

For \(\xi\in\mathfrak g^4\), set
\[
 x=\zeta_1=\xi_1,\qquad
 d_v=\zeta_v-x\quad(v=2,3,4),\qquad
 d=(d_2,d_3,d_4),                                         \tag{32.14}
\]
and write \(d_1=0\).  The tree-edge energies are exactly
\[
 \|(d_U\xi)_{1v}\|=\|d_v\|.                               \tag{32.15}
\]
For a chord \(a\to b\),
\[
 \operatorname{Ad}_{P_b^{-1}}(d_U\xi)_{ab}
 =(I-\operatorname{Ad}_{\mathcal H_{ab}})x
   +d_b-\operatorname{Ad}_{\mathcal H_{ab}}d_a.           \tag{32.16}
\]
The Higgs rows are
\[
 L_v\xi_v
 =L_v\operatorname{Ad}_{P_v}x
  +L_v\operatorname{Ad}_{P_v}d_v.                         \tag{32.17}
\]

Let \(\mathfrak E_q:(\mathfrak g)^3\to\operatorname{codom}\mathfrak A_q\)
be the exact error map whose chord and Higgs components are
\[
 \bigl(\mathfrak E_qd\bigr)_{ab}
 =\sqrt{w_{ab}}
  \bigl(d_b-\operatorname{Ad}_{\mathcal H_{ab}}d_a\bigr),
                                                               \tag{32.18}
\]
\[
 \bigl(\mathfrak E_qd\bigr)_{H,v}
 =\sqrt{m_v}L_v\operatorname{Ad}_{P_v}d_v,                 \tag{32.19}
\]
with \(d_1=0\).  Put
\[
 \Lambda_q=\|\mathfrak E_q\|.                             \tag{32.20}
\]
If \(y_q(\xi)\) denotes the vector of all weighted chord and Higgs
residuals, then (32.16)--(32.19) give the exact factorization
\[
 \boxed{y_q(\xi)=\mathfrak A_qx+\mathfrak E_qd.}           \tag{32.21}
\]

\subsection{A direct rooted covariant Poincare floor}

Since \(\mathfrak A_q\) is finite-dimensional, its restriction to
\(\mathfrak h_q^\perp\) has the positive floor
\[
 \nu_q:=
 \lambda_{\min}\left(
 \mathfrak A_q^*\mathfrak A_q
 \big|_{\mathfrak h_q^\perp}\right)>0                     \tag{32.22}
\]
when \(\mathfrak h_q^\perp\ne0\).  If the complement is zero, set
\(\nu_q=+\infty\) and interpret \(\nu_q^{-1}=0\).
Let
\[
 w_T=\min\{w_{12},w_{13},w_{14}\},\qquad
 \omega_q=\min\{1,w_T\}.                                  \tag{32.23}
\]
Define
\[
 c_a(q)=\frac{2}{\sqrt{\nu_q}},\qquad
 c_d(q)=1+\frac{\sqrt3}{2}
              +\frac{2\Lambda_q}{\sqrt{\nu_q}},           \tag{32.24}
\]
and
\[
 \boxed{
 \delta_q^{\rm root}
 :=\frac{\omega_q}
 {1+c_a(q)^2+c_d(q)^2}.}                                  \tag{32.25}
\]

\begin{theorem}[Arbitrary-background bosonic Gauss floor]
\label{thm:v14v32-floor}
For every \(\xi\in\mathfrak g^4\),
\[
 \boxed{
 \mathcal E_q(\xi)+\|\Pi_q\xi\|^2
 \ge\delta_q^{\rm root}\|\xi\|^2.}                        \tag{32.26}
\]
At zero canonical momentum, where
\(C_{G,(q,0)}C_{G,(q,0)}^*=\mathcal R_q^*\mathcal R_q\),
this is the completed constraint-Gram certificate
\[
 \boxed{
 C_{G,(q,0)}C_{G,(q,0)}^*+\Pi_q
 \succeq\delta_q^{\rm root}I_{\mathfrak g^4}.}             \tag{32.27}
\]
Consequently the declared relation bank (32.13) is complete and
\[
 \|C_{G,(q,0)}^\dagger\|
 \le(\delta_q^{\rm root})^{-1/2}.                          \tag{32.28}
\]
\end{theorem}

\begin{proof}
Let
\[
 D=\left(\sum_{v=2}^4\|d_v\|^2\right)^{1/2},\qquad
 A=\|y_q(\xi)\|,\qquad R=\|\Pi_q\xi\|.                     \tag{32.29}
\]
By (32.21),
\[
 \|\mathfrak A_qx\|\le A+\Lambda_qD.
\]
Therefore
\[
 \|P_{\mathfrak h_q^\perp}x\|
 \le\nu_q^{-1/2}(A+\Lambda_qD).                            \tag{32.30}
\]
Now
\[
 \overline\zeta=x+\overline d,\qquad
 \overline d=\frac14(d_2+d_3+d_4),\qquad
 \|\overline d\|\le\frac{\sqrt3}{4}D.                     \tag{32.31}
\]
By (32.13),
\[
 R=2\|P_{\mathfrak h_q}\overline\zeta\|,
\]
and hence
\[
 \|P_{\mathfrak h_q}x\|
 \le\frac R2+\frac{\sqrt3}{4}D.                           \tag{32.32}
\]
The unitary change to \(\zeta\) and the triangle inequality give
\[
 \begin{aligned}
 \|\xi\|=\|\zeta\|
 &\le2\|x\|+D\\
 &\le R+\frac2{\sqrt{\nu_q}}A+
 \left(1+\frac{\sqrt3}{2}
       +\frac{2\Lambda_q}{\sqrt{\nu_q}}\right)D\\
 &=R+c_a(q)A+c_d(q)D.
 \end{aligned}                                             \tag{32.33}
\]
Cauchy--Schwarz now yields
\[
 \|\xi\|^2
 \le\bigl(1+c_a(q)^2+c_d(q)^2\bigr)
       (R^2+A^2+D^2).                                     \tag{32.34}
\]
The energy splits into tree and anchor rows, so
\[
 \mathcal E_q(\xi)+R^2
 \ge w_TD^2+A^2+R^2
 \ge\omega_q(D^2+A^2+R^2).                               \tag{32.35}
\]
Combining (32.34)--(32.35) proves (32.26).  The zero-momentum Gram
identity is V29 equation (29.13), which gives (32.27).  The V16
kernel-completion theorem then proves relation completeness and
(32.28).
\end{proof}

The bound is deliberately elementary and may be conservative.  Its
advantage is that every quantity is read from matrices of the chosen
background: three tree weights, a root anchor on the
twelve-dimensional algebra \(\mathfrak g\), and the finite error norm
\(\Lambda_q\).  It remains valid at every orbit type because the
relation projector is recomputed rather than transported across a
rank jump.

\begin{corollary}[Uniform cofinal criterion]
\label{cor:v14v32-cofinal}
For a regulator sequence \(q_h\), suppose one can choose star trees
so that
\[
 \inf_h w_{T,h}>0,\qquad
 \inf_h\nu_{q_h}>0,\qquad
 \sup_h\Lambda_{q_h}<\infty.                              \tag{32.36}
\]
Then
\[
 \inf_h\delta_{q_h}^{\rm root}>0.                         \tag{32.37}
\]
If the relation projectors and constraint derivatives converge on
each fixed screen, the internal bosonic Gauss projectors and
pseudoinverses converge there.
\end{corollary}

\begin{proof}
The right side of (32.25) is bounded below by a positive constant
under (32.36), proving (32.37).  Apply the V16/V22 completed-Gram
landing theorem on each fixed screen.
\end{proof}

Condition (32.36) is a physical acquisition, not a consequence of
connectedness.  Holonomies and Higgs fields may approach a more
symmetric stratum, making \(\nu_{q_h}\to0\).

\subsection{Consistency checks}

At the V29 invariant vacuum, all rooted holonomies equal the identity
and all Higgs values equal \(v_Hh_0\).  Then
\[
 \mathfrak h_q=\mathfrak h\cong\mathfrak u(3),             \tag{32.38}
\]
so (32.13) reduces to the diagonal projector (29.20).  The direct
V29 spectral decomposition gives the sharper floor
\(\min\{4w_{\min},m_{\min}\nu_H,1\}\).  The rooted bound (32.25) is a
generic-background fallback and is not substituted for that sharper
symmetric-vacuum calculation.

For a generic background, chord holonomies can reduce the root
centralizer and inhomogeneous Higgs values can further reduce the
intersection (32.10).  The corresponding formerly relational
directions are controlled directly by \(\nu_q\), so no continuity
argument through the rank change is needed.

\subsection{What remains before the full internal block}

The arbitrary-background status is
\[
\begin{array}{p{0.29\linewidth}|p{0.17\linewidth}|p{0.43\linewidth}}
\textbf{datum}&\textbf{status}&\textbf{content}\\ \hline
bosonic relation algebra&derived exactly&
holonomy--Higgs intersection (32.10)\\
bosonic relation projector&derived exactly&
transported root average (32.13)\\
bosonic completed floor&derived conditionally&
root certificate (32.25)--(32.27)\\
uniform regulator floor&open&
three quantitative rows (32.36)\\
nonzero-momentum phase point&open&
cotangent stabilizer includes the momentum condition from V29\\
graded fermion Gauss row&open&
independent chiral amplitudes require a graded phase treatment\\
gravity-dependent Riesz map&open&
common Hodge whitening can contribute the coupled blocks\\
temporal-link Legendre origin&open&
the Euclidean all-link action has not been reduced to one canonical
time slab
\end{array}                                                 \tag{32.39}
\]

\begin{decision}[Direction after V32]
\label{dec:v14v32-next}
V32 closes the finite bosonic relation and floor problem at an
arbitrary fixed \(K_4\) background, subject to explicit root-anchor
margins.  V33 will now examine the graded fermion row rather than
folding it into the positive bosonic Gram.  Starting from independent
\(\Psi,\overline\Psi\) and the source gauge action, it will construct
the even graded symplectic form and fermionic gauge momentum map,
derive its contribution to the Gauss identity, and prove precisely
what its nilpotent coefficients can and cannot contribute to a real
positive constraint floor.  This will decide whether the finite
Hilbert projector must be built on the bosonic body first and the CAR
sector represented afterward.
\end{decision}

\subsection{Acquisition ledger}

\[
\begin{array}{p{0.30\linewidth}|p{0.20\linewidth}|p{0.39\linewidth}}
\textbf{gate}&\textbf{status after V32}&\textbf{content}\\ \hline
generic Gauss relations&closed at fixed cutoff&
root holonomy and Higgs stabilizers\\
generic relation projector&closed at fixed cutoff&
formula (32.13)\\
generic bosonic floor&closed conditionally&
finite root-anchor card\\
uniform root-anchor margins&open&
cofinal condition (32.36)\\
graded fermion Gauss row&next&
even graded momentum map\\
gravity-dependent whitening&open&
full coupled blocks\\
temporal canonical reduction&open&
common one-slab action
\end{array}                                                 \tag{32.40}
\]

\subsection{Parent-integrity record}

The installed parent was
\[
\begin{array}{c|c}
 \text{file}&
 \texttt{Renewal\_SM\_Einstein\_Continuum\_Research\_Notes\_V31.tex}\\
 \text{lines}&12870\\
 \text{bytes}&418009\\
 \text{SHA--256}&
 \texttt{B90A68027AD450C1B6BE50660491420F60F89EB8C99C4A4200D815B2132DA8D2}.
\end{array}                                                 \tag{32.41}
\]

% =============================================================================
% END APPENDED FOURTEENTH-CYCLE V32 MATERIAL
% =============================================================================
% =============================================================================
% BEGIN APPENDED FOURTEENTH-CYCLE V33 MATERIAL
% =============================================================================

\section{Fourteenth-cycle V33: the graded fermion Gauss row and body
reduction of the positive constraint floor}
\label{sec:v14v33}

\subsection{Why a graded row must be typed separately}

The source Standard-Model action uses independent chiral amplitudes
\(\Psi\) and \(\overline\Psi\), with gauge transformations
\[
 \Psi_v\longmapsto\rho_\Psi(g_v)\Psi_v,\qquad
 \overline\Psi_v\longmapsto
 \overline\Psi_v\rho_\Psi(g_v)^{-1}.                       \tag{33.1}
\]
Classically these variables are odd.  They do not form an ordinary
real Hilbert cotangent space, and their nilpotent coefficients do not
carry an ordered positive cone.  The correct finite statement has
two layers:

\begin{enumerate}
\item an even graded momentum map supplies the fermionic charge in the
Gauss constraint and preserves exact first-class closure;
\item the real positive completed-Gram floor is a statement about the
body, while the nilpotent fermion correction is handled by finite
algebraic inversion.
\end{enumerate}

This separation retains the source Ward/BRST identities without
using a finite-Dirac spectral gap as a constraint floor.

\subsection{Finite Grassmann phase space}

Let \(\Lambda\) be a finite-dimensional real Grassmann algebra, let
\(\mathfrak n\subset\Lambda\) be its ideal of elements with zero body,
and choose \(r\) with
\[
 \mathfrak n^r=0.                                         \tag{33.2}
\]
Let \(F_v\) be the realification of the finite complex chiral
one-particle fibre at vertex \(v\).  The odd configuration space is
\[
 \mathcal Q_F=\bigoplus_{v\in V}\Pi F_v,                   \tag{33.3}
\]
and its even cotangent supermanifold has odd canonical coordinates
\[
 (\Psi_v,\overline\Psi_v)
 \in(\Lambda_{\overline1}\otimes F_v)
 \oplus(\Lambda_{\overline1}\otimes F_v^*).                \tag{33.4}
\]
Here \(\Pi\) denotes parity reversal.  The canonical cotangent
symplectic form is even; consequently its Poisson bracket is the
ordinary even graded bracket.  We fix the cotangent sign convention
in which a lifted configuration action has momentum map equal to the
momentum paired with the fundamental configuration vector.

For \(\xi=(\xi_v)\in\mathfrak g^V\), write
\[
 T_v(\xi_v)=\rho_{\Psi,*}(\xi_v).
\]
The infinitesimal action is
\[
 \delta_\xi\Psi_v=T_v(\xi_v)\Psi_v,\qquad
 \delta_\xi\overline\Psi_v=-\overline\Psi_vT_v(\xi_v).
                                                               \tag{33.5}
\]

\begin{theorem}[Graded fermion Gauss momentum map]
\label{thm:v14v33-momentum}
The cotangent lift of (33.5) is Hamiltonian.  Its even momentum map
\(\Phi_F\) is characterized by
\[
 \boxed{
 \langle\Phi_F(\Psi,\overline\Psi),\xi\rangle
 =\sum_{v\in V}
 \overline\Psi_vT_v(\xi_v)\Psi_v.}                         \tag{33.6}
\]
For phase variations
\((\dot\Psi,\dot{\overline\Psi})\),
\[
 \boxed{
 \begin{aligned}
 \left\langle D\Phi_F
  [\dot\Psi,\dot{\overline\Psi}],\xi\right\rangle
 =\sum_v\bigl(
 &\dot{\overline\Psi}_vT_v(\xi_v)\Psi_v\\
 &+\overline\Psi_vT_v(\xi_v)\dot\Psi_v
 \bigr).
 \end{aligned}}                                            \tag{33.7}
\]
The graded Poisson brackets obey
\[
 \boxed{
 \{\langle\Phi_F,\xi\rangle,
   \langle\Phi_F,\eta\rangle\}_{\rm gr}
 =\langle\Phi_F,[\xi,\eta]\rangle.}                        \tag{33.8}
\]
\end{theorem}

\begin{proof}
The canonical cotangent momentum-map construction applies to a
linear action on a finite super vector space.  Pairing the odd
cotangent coordinate \(\overline\Psi_v\) with the fundamental vector
\(T_v(\xi_v)\Psi_v\) gives (33.6); the product is even because it is
quadratic in odd variables.  Differentiation gives (33.7).

The Hamiltonian vector field of the left side of (33.6) is the
cotangent lift of (33.5).  Commutators of the lifted fundamental
fields reproduce
\[
 [T_v(\xi_v),T_v(\eta_v)]=T_v([\xi_v,\eta_v]).
\]
The canonical cotangent momentum map vanishes on the zero section,
so there is no additive central constant.  This proves equivariance,
which is exactly (33.8).
\end{proof}

Combining (33.6) with the bosonic momentum map of V29 gives
\[
 \Phi_G^{\rm tot}
 =\Phi_G^{\rm bos}+\Phi_F.                                \tag{33.9}
\]
The direct-sum graded symplectic form and (29.14), (33.8) imply
\[
 \{\langle\Phi_G^{\rm tot},\xi\rangle,
   \langle\Phi_G^{\rm tot},\eta\rangle\}_{\rm gr}
 =\langle\Phi_G^{\rm tot},[\xi,\eta]\rangle.               \tag{33.10}
\]
Thus the fermionic charge does not create a classical Gauss anomaly.

\subsection{BRST compatibility}

Let \(c\in\Lambda_{\overline1}\otimes\mathfrak g^V\) be the ghost and
set
\[
 sc=-c^2,\qquad s\Psi=c\Psi,\qquad
 s\overline\Psi=-\overline\Psi c.                         \tag{33.11}
\]

\begin{corollary}[Nilpotence on the fermion row]
\label{cor:v14v33-brst}
With the graded Leibniz rule,
\[
 s^2\Psi=0,\qquad s^2\overline\Psi=0.                     \tag{33.12}
\]
If the fermion action is gauge invariant, its BRST variation
vanishes.
\end{corollary}

\begin{proof}
Using that \(c\) is odd,
\[
 s^2\Psi=(sc)\Psi-c(s\Psi)=-c^2\Psi+c^2\Psi=0.
\]
The dual calculation is the same with the right action.  The BRST
variation of an invariant action is its infinitesimal gauge
variation with the label replaced by \(c\), hence is zero.
\end{proof}

This recovers the source's classical BRST row in the canonical
graded phase language.  It does not construct the quantum BRST
cohomology.

\subsection{Body of the fermionic constraint derivative}

Let
\[
 \operatorname{body}:\Lambda_{\overline0}\longrightarrow\mathbb R
\]
be the quotient by \(\mathfrak n\).  Every term in (33.6) is
quadratic in odd generators, so
\[
 \operatorname{body}\Phi_F=0.                             \tag{33.13}
\]
At the zero fermion configuration,
\[
 \Phi_F(0,0)=0,\qquad D\Phi_F(0,0)=0.                     \tag{33.14}
\]
At a general Grassmann configuration, the coefficients in (33.7)
belong to the nilpotent ideal.

\begin{proposition}[The fermion row has zero positive body]
\label{prop:v14v33-zero-body}
Let \(C_{\rm gr}\) be the full graded Gauss derivative and
\(C_{\rm bos}\) its bosonic restriction, extended by zero on the odd
columns.  In any based extension over \(\Lambda\),
\[
 \operatorname{body}C_{\rm gr}=C_{\rm bos},               \tag{33.15}
\]
and the body of every even Gram expression constructed polynomially
from \(C_{\rm gr}\) and its graded adjoint is the corresponding
bosonic Gram.  In particular the fermionic classical row cannot
increase, decrease, or create a strictly positive real
constraint-Gram floor.
\end{proposition}

\begin{proof}
Equation (33.7) shows that every fermionic entry of the derivative
has an odd coefficient and every even correction formed from two such
entries lies in \(\mathfrak n\).  The body homomorphism kills
\(\mathfrak n\) and preserves sums, products, and real adjoints.
Therefore applying it to the full derivative leaves the bosonic
matrix, and applying it to a Gram leaves the bosonic Gram.
A real positive spectral floor is determined by that body matrix;
nilpotent coefficients have zero body and no ordered real magnitude.
\end{proof}

There is also a representation-theoretic obstruction to treating a
nonzero nilpotent as a positive operator.

\begin{lemma}[Positive nilpotents vanish in Hilbert representations]
\label{lem:v14v33-positive-nilpotent}
If a bounded Hilbert-space operator \(N\) is positive and
\(N^k=0\) for some \(k\), then \(N=0\).
\end{lemma}

\begin{proof}
A positive operator is self-adjoint and
\(\|N\|\) equals its spectral radius.  Nilpotence makes its spectrum
\(\{0\}\), so \(\|N\|=0\).
\end{proof}

Thus no faithful positive Hilbert interpretation can turn a
classical nilpotent fermion correction into an additional positive
constraint direction.

\subsection{Finite algebraic completion over the Grassmann algebra}

Although nilpotent corrections carry no positive floor, they can be
inverted algebraically once the body floor is known.

\begin{theorem}[Body-floor inversion of the graded completed Gram]
\label{thm:v14v33-body-inverse}
Let
\[
 \widehat G_{\rm gr}
 =\widehat G_{\rm bos}+N                                  \tag{33.16}
\]
be a finite square matrix over \(\Lambda_{\overline0}\), where all
entries of \(N\) lie in \(\mathfrak n\).  Then
\(\widehat G_{\rm gr}\) is invertible over \(\Lambda_{\overline0}\)
if and only if \(\widehat G_{\rm bos}\) is invertible over
\(\mathbb R\).  If
\[
 \widehat G_{\rm bos}\succeq\delta I,\qquad\delta>0,        \tag{33.17}
\]
then
\[
 \boxed{
 \widehat G_{\rm gr}^{-1}
 =\sum_{j=0}^{r-1}
  \bigl(-\widehat G_{\rm bos}^{-1}N\bigr)^j
  \widehat G_{\rm bos}^{-1}.}                             \tag{33.18}
\]
\end{theorem}

\begin{proof}
If \(\widehat G_{\rm gr}\) is invertible, applying the body
homomorphism to
\(\widehat G_{\rm gr}\widehat G_{\rm gr}^{-1}=I\) proves that
\(\widehat G_{\rm bos}\) is invertible.  Conversely, set
\(M=\widehat G_{\rm bos}^{-1}N\).  Every entry of \(M\) lies in
\(\mathfrak n\), so \(M^r=0\) after increasing \(r\), if necessary,
to the nilpotence index of the finite matrix ideal.  Hence
\[
 (I+M)^{-1}=\sum_{j=0}^{r-1}(-M)^j.
\]
Since
\(\widehat G_{\rm gr}=\widehat G_{\rm bos}(I+M)\), this is
(33.18).  Condition (33.17) supplies the required real inverse and
the bound
\(\|\widehat G_{\rm bos}^{-1}\|\le\delta^{-1}\).
\end{proof}

The finite series (33.18) is an algebraic inverse, not a
Moore--Penrose inverse in an ordered super Hilbert space.  Its real
conditioning remains the V32 bosonic body floor.

\begin{corollary}[Graded extension of the rooted \(K_4\) card]
\label{cor:v14v33-rooted}
At a zero-momentum bosonic background satisfying V32
Theorem \ref{thm:v14v32-floor}, the classical fermion charge produces
a nilpotent extension of the completed Gauss Gram.  It is algebraically
invertible by (33.18), with body inverse bounded by
\[
 \|(\widehat G_{\rm gr}^{-1})_{\rm body}\|
 \le(\delta_q^{\rm root})^{-1}.                            \tag{33.19}
\]
The real relation projector and positive floor remain those of V32.
\end{corollary}

\begin{proof}
Proposition \ref{prop:v14v33-zero-body} identifies the body with the
V32 bosonic completed Gram, and V32 supplies its positive floor.
Apply Theorem \ref{thm:v14v33-body-inverse}.
\end{proof}

\subsection{Quantized CAR charge is a different positive card}

After CAR quantization on a finite fermion Fock space,
\(\mathcal F(F)\), the gauge charge is represented by second
quantization:
\[
 \widehat Q_\xi=\mathrm d\Gamma(iT(\xi))=\widehat Q_\xi^*.
\]
The source identity gives
\[
 [\widehat Q_\xi,\widehat Q_\eta]
 =i\widehat Q_{[\xi,\eta]}.                                \tag{33.20}
\]
For a faithful density matrix \(\rho\), one may form the positive
centered covariance
\[
 G_\rho(\xi,\eta)
 =\frac12\operatorname{Tr}\rho
 \left\{\widehat Q_\xi^\circ,\widehat Q_\eta^\circ\right\},
 \qquad
 \widehat Q_\xi^\circ
 =\widehat Q_\xi-\operatorname{Tr}(\rho\widehat Q_\xi)I.   \tag{33.21}
\]
This Gram measures visibility of represented charges in the chosen
state.  It is not the classical constraint derivative
\(D\Phi_G\), and its floor depends on the representation and state.
A finite-Dirac energy gap does not determine (33.21).

\begin{firewall}[Classical, represented, and physical Gauss rows]
\label{fire:v14v33-three-rows}
The graded momentum map (33.6), the bosonic body constraint projector
of V32, and the quantum covariance (33.21) are three distinct
objects.  Exact BRST or CAR closure does not identify their Grams.
A full quantum constraint handoff must specify which represented
charge space, state metric, and physical quotient are used.
\end{firewall}

\subsection{Acquisition result and next direction}

The internal constraint card now has the form
\[
\begin{array}{p{0.29\linewidth}|p{0.17\linewidth}|p{0.43\linewidth}}
\textbf{datum}&\textbf{status}&\textbf{content}\\ \hline
graded fermion gauge action&explicit&
source representation (33.1)\\
graded fermion momentum map&derived&
charge formula (33.6)\\
graded first-class closure&derived&
equation (33.8)\\
fermionic real floor contribution&exactly zero&
body theorem (33.15)\\
graded completed inverse&derived conditionally&
finite series (33.18) from the bosonic floor\\
quantum charge covariance&defined, unpopulated&
requires a CAR state and represented charge basis\\
full bosonic body card&conditional&
V32 root-anchor margins\\
common temporal action split&open&
still required to derive the canonical multiplier from the source
Euclidean slab
\end{array}                                                 \tag{33.22}
\]

\begin{decision}[Direction after V33]
\label{dec:v14v33-next}
The classical graded fermion row no longer blocks algebraic
completion: it is a nilpotent extension controlled by the bosonic
body floor.  The first unresolved provenance question is now whether
the source's group-valued temporal links yield the affine multiplier
\(a_0\) required by V27 at finite cutoff.

V34 will perform the one-slab temporal-link audit.  It will compare a
Wilson plaquette containing a temporal link with a first-order
Hamiltonian action, prove that the exact group-valued link dependence
is generally non-affine, and identify the additional transfer-matrix
or electric-field transform needed to obtain the exact Gauss
multiplier.  A small-\(h\) expansion will be kept distinct from an
exact finite constraint card.
\end{decision}

\subsection{Acquisition ledger}

\[
\begin{array}{p{0.30\linewidth}|p{0.20\linewidth}|p{0.39\linewidth}}
\textbf{gate}&\textbf{status after V33}&\textbf{content}\\ \hline
fermion classical Gauss row&closed algebraically&
graded cotangent momentum map\\
fermion bracket closure&closed&
equivariant graded momentum map\\
positive real conditioning&bosonic body only&
V32 floor unchanged\\
nilpotent correction inverse&closed&
finite Neumann series\\
represented CAR charge Gram&open&
state-dependent quantum card\\
temporal link to affine multiplier&next&
one-slab canonical reduction\\
gravity-dependent whitening&open&
coupled SM--Einstein blocks
\end{array}                                                 \tag{33.23}
\]

\subsection{Parent-integrity record}

The installed parent was
\[
\begin{array}{c|c}
 \text{file}&
 \texttt{Renewal\_SM\_Einstein\_Continuum\_Research\_Notes\_V32.tex}\\
 \text{lines}&13311\\
 \text{bytes}&432054\\
 \text{SHA--256}&
 \texttt{E6524319EBCAB3EF2E058C7499EE1FC763133881C7E3FBE44F9E6534C9F2FEB6}.
\end{array}                                                 \tag{33.24}
\]

% =============================================================================
% END APPENDED FOURTEENTH-CYCLE V33 MATERIAL
% =============================================================================



% =============================================================================
% BEGIN APPENDED FOURTEENTH-CYCLE V34 MATERIAL
% =============================================================================

\section{Fourteenth-cycle V34: temporal links, exact group averaging,
and the missing affine Gauss multiplier}
\label{sec:v14v34}

\subsection{The provenance question}

The source active Standard-Model action is a Euclidean cell action
built from link holonomies, Higgs covariant differences, plaquettes,
and the gauge-covariant Dirac--Yukawa operator.  It is exactly gauge
invariant.  V27, however, needs a canonical one-slab action of the
form
\[
 \mathscr S_{H,h}
 =\mathscr S_{0,h}
 +\langle a_{0,h},\Phi_{G,h}\rangle,                       \tag{34.1}
\]
with \(a_{0,h}\in\mathfrak g^{V_{\rm sp}}\) an affine multiplier.  The
fact that both objects encode gauge symmetry does not make their
finite dependence on temporal data identical.

Let \(U_{0,v}\in G\) denote a temporal link.  A plaquette containing
that link depends on products such as
\[
 U_{0,v}U_{i,v+\widehat0}
 U_{0,v+\widehat i}^{-1}U_{i,v}^{-1}.                     \tag{34.2}
\]
Thus the source variable is group-valued and enters through a
character of a product, not through a linear pairing with a Lie
algebra element.

\subsection{Finite nonaffinity}

\begin{countertheorem}[One \(U(1)\) plaquette is already nonaffine]
\label{cth:v14v34-u1}
Let a temporal \(U(1)\) link be
\[
 U_0(a)=e^{iha}
\]
and freeze the other three link phases of a plaquette into
\(\alpha\).  The Wilson plaquette action is
\[
 S_h(a)=w\bigl(1-\cos(ha+\alpha)\bigr),\qquad w>0.          \tag{34.3}
\]
On every open interval on which
\(\cos(ha+\alpha)\) is not identically zero,
\[
 S_h''(a)=wh^2\cos(ha+\alpha)\ne0.                         \tag{34.4}
\]
Hence \(S_h\) is not affine in the logarithmic temporal-link
coordinate.
\end{countertheorem}

\begin{proof}
Differentiate (34.3) twice.  An affine function has identically zero
second derivative, so (34.4) excludes affinity.
\end{proof}

This is not a nonabelian complication or a fermion effect.  It occurs
in the smallest compact gauge group and one plaquette.

\begin{theorem}[No global nonconstant affine coordinate on a compact
gauge group]
\label{thm:v14v34-no-global-affine}
Let \(G\) be compact and connected.  Suppose a continuous
\(F:G\to\mathbb R\) satisfies, for every \(\xi\in\mathfrak g\),
\[
 F(\exp(t\xi))=F(e)+t\ell(\xi)                             \tag{34.5}
\]
for all \(t\in\mathbb R\), with \(\ell:\mathfrak g\to\mathbb R\)
linear.  Then \(\ell=0\), and \(F\) is constant on every
one-parameter subgroup.  If these subgroups generate \(G\), then
\(F\) is constant.
\end{theorem}

\begin{proof}
For fixed \(\xi\), compactness of \(G\) makes
\(\{\exp(t\xi):t\in\mathbb R\}\) relatively compact.  Continuity makes
\(F\) bounded on its closure.  The right side of (34.5) is bounded in
\(t\) only when \(\ell(\xi)=0\).  This holds for every \(\xi\), so
\(\ell=0\).  Connected compact Lie groups are generated by their
one-parameter subgroups, proving the last assertion.
\end{proof}

Therefore a global group-valued temporal link cannot itself be the
nontrivial affine multiplier in (34.1).  One must either pass to a
local infinitesimal limit, introduce a first-order Hamiltonian
variable, or impose gauge invariance by exact group projection.

\subsection{What the logarithmic expansion actually proves}

Let \(F_h:G\to\mathbb R\) be \(C^2\) near the identity and put
\[
 f_h(a)=F_h(\exp(ha)).
\]
For a bounded Lie-algebra ball \(\|a\|\le R\), define
\[
 M_h(R)=
 \sup_{\substack{\|a\|\le R\\0\le t\le1}}
 \left\|D_a^2\bigl[F_h(\exp(tha))\bigr]\right\|.           \tag{34.6}
\]

\begin{proposition}[Controlled infinitesimal multiplier row]
\label{prop:v14v34-taylor}
For \(\|a\|\le R\),
\[
 \boxed{
 f_h(a)=F_h(e)+h\,dF_h(e)[a]+\mathcal R_h(a),\qquad
 |\mathcal R_h(a)|\le\frac12M_h(R)\|a\|^2.}                \tag{34.7}
\]
If the Hessian in (34.6) is written with respect to the unscaled
variable \(tha\), so that
\[
 M_h(R)\le h^2\widetilde M_h(R),                           \tag{34.8}
\]
then
\[
 |\mathcal R_h(a)|
 \le\frac{h^2}{2}\widetilde M_h(R)\|a\|^2.                \tag{34.9}
\]
The first derivative \(dF_h(e)\) gives a candidate Gauss row only if
the rescaled remainder required by the action normalization tends to
zero uniformly.
\end{proposition}

\begin{proof}
Apply one-variable Taylor's theorem with integral remainder to
\(t\mapsto F_h(\exp(tha))\) between \(0\) and \(1\).  Its first
derivative at zero is \(h\,dF_h(e)[a]\), and its second derivative is
bounded by (34.6), proving (34.7).  The chain-rule scaling in (34.8)
gives (34.9).
\end{proof}

For an action divided by one time step, (34.9) leaves an error of
order
\[
 h\,\widetilde M_h(R)\|a\|^2.                             \tag{34.10}
\]
A continuum multiplier follows only with a uniform bound making
(34.10) vanish.  At fixed \(h\), the higher terms remain and (34.1)
is not exact.

\subsection{Exact finite alternative: Haar projection}

Let a compact gauge group \(\mathcal G_X=G^{V_{\rm sp}}\) act
unitarily on a finite Hilbert space \(\mathcal H_h\) by
\(\mathcal U_h(g)\).  Define
\[
 P_{{\rm inv},h}
 :=\int_{\mathcal G_X}\mathcal U_h(g)\,dg,                 \tag{34.11}
\]
where Haar measure has total mass one.

\begin{theorem}[Temporal-link integration is an exact gauge
projector]
\label{thm:v14v34-haar}
The operator \(P_{{\rm inv},h}\) is the orthogonal projector onto
\[
 \mathcal H_h^{\mathcal G_X}
 =\{\psi:\mathcal U_h(g)\psi=\psi\ \hbox{for every }g\}.
                                                               \tag{34.12}
\]
If \(\mathcal G_X\) is connected and its derived representation has
self-adjoint generators \(\widehat\Phi_h(\xi)\), then
\[
 \mathcal H_h^{\mathcal G_X}
 =\bigcap_{\xi\in\mathfrak g_X}
   \ker\widehat\Phi_h(\xi).                                \tag{34.13}
\]
Thus integration over a temporal gauge link can impose Gauss
invariance exactly, but it does so by compact group averaging rather
than by an affine finite Lie-algebra multiplier.
\end{theorem}

\begin{proof}
Haar invariance gives
\[
 P_{{\rm inv},h}^*
 =P_{{\rm inv},h},\qquad
 P_{{\rm inv},h}^2
 =\int\!\!\int\mathcal U_h(gg')\,dg\,dg'
 =P_{{\rm inv},h}.                                       \tag{34.14}
\]
It also gives
\(\mathcal U_h(k)P_{{\rm inv},h}=P_{{\rm inv},h}\), so the
range is invariant.  Conversely an invariant vector is fixed by the
integral.  This proves (34.12).

For a connected group, a vector is invariant precisely when it is
annihilated by every infinitesimal generator: one direction follows
by differentiation, and the other by exponentiation and generation
of the identity component.  This proves (34.13).
\end{proof}

In a Peter--Weyl transform, (34.11) contracts at each vertex onto the
invariant tensor subspace.  This is the exact finite representation
of the Gauss law used by a transfer construction.  It provides a
quantum physical projector, not the classical mixed Hessian
\(D_zD_{a_0}\mathscr S_h\) required by V26.

\subsection{Exact Hamiltonian branch}

On the bosonic phase space of V29, let \(\Theta_h\) be the canonical
one-form, let \(H_{0,h}\) be gauge invariant, and let
\(\Phi_{G,h}\) be the momentum map.  The first-order action
\[
 \boxed{
 \mathscr S_h^{(1)}
 =\int\left[
 \Theta_h(\dot z)-H_{0,h}(z)
 -\langle a_0,\Phi_{G,h}(z)\rangle
 \right]dt}                                                \tag{34.15}
\]
is exactly affine in \(a_0\).

\begin{proposition}[First-order sufficiency]
\label{prop:v14v34-first-order}
If a selected regulator supplies (34.15) in based finite spaces, then
\[
 -D_zD_{a_0}\mathscr S_h^{(1)}
 =D_z\Phi_{G,h}=C_{G,h}.                                  \tag{34.16}
\]
The V29--V33 relation, floor, first-class, and graded-extension
theorems then apply to its internal Gauss block.  Exact equivalence
between (34.15) and a Euclidean temporal-link action requires an
additional transfer/Fourier--Legendre theorem.
\end{proposition}

\begin{proof}
Differentiate (34.15) first in \(a_0\) and then in \(z\) to obtain
(34.16).  The cited theorems use precisely the momentum map and its
derivative.  None of those differentiations proves equality of
(34.15) with the character-valued plaquette action, so that equality
is an independent hypothesis.
\end{proof}

A sufficient source bridge would be a same-carrier transfer identity
of the form
\[
 T_h=P_{{\rm inv},h}\,e^{-h\widehat H_h}\,
     P_{{\rm inv},h}                                      \tag{34.17}
\]
together with a classical-symbol/Legendre identification of
\(\widehat H_h\), the electric variables, and the V29 momentum map.
The source contains Peter--Weyl and finite action machinery, but the
inspected active action theorem does not display (34.17) on the
selected \(K_4\) slab.

\subsection{Acquisition verdict}

\[
\begin{array}{p{0.29\linewidth}|p{0.17\linewidth}|p{0.43\linewidth}}
\textbf{datum}&\textbf{status}&\textbf{content}\\ \hline
Euclidean temporal link&explicit&
group-valued link in plaquette and covariant matter terms\\
exact logarithmic affinity&ruled out&
\(U(1)\) witness (34.3)--(34.4)\\
first infinitesimal Gauss row&derived conditionally&
Taylor coefficient in (34.7)\\
cofinal affine limit&open&
uniform rescaled remainder (34.10)\\
exact finite quantum Gauss projector&derived&
Haar average (34.11)--(34.13)\\
exact classical affine multiplier&conditional&
first-order action (34.15)\\
Euclidean-to-Hamiltonian equivalence&missing&
same-carrier transfer/Fourier--Legendre card (34.17)\\
internal constraint floor&conditional&
V32 body floor once the Hamiltonian branch is selected
\end{array}                                                 \tag{34.18}
\]

\begin{decision}[Direction after V34]
\label{dec:v14v34-next}
V34 rules out reading an exact affine multiplier directly from a
finite Wilson temporal link.  There are now two honest branches:

\begin{enumerate}
\item a classical Hamiltonian branch, which must populate the
first-order action (34.15) and prove its relation to the common
source action;
\item a quantum transfer branch, which uses the exact Haar projector
(34.11) and must build its represented constraint/residual metric
separately.
\end{enumerate}

V35 is the compulsory companion audit.  It will inspect the newest
Yang--Mills and Navier--Stokes notes for a transfer, electric-field,
or quotient-metric construction relevant to this fork.  V36 will
then continue on the branch with the strongest same-carrier input.
\end{decision}

\subsection{Acquisition ledger}

\[
\begin{array}{p{0.30\linewidth}|p{0.20\linewidth}|p{0.39\linewidth}}
\textbf{gate}&\textbf{status after V34}&\textbf{content}\\ \hline
finite temporal-link affinity&ruled out&
compact-character obstruction\\
cofinal Lie-algebra row&conditional&
Taylor remainder gate\\
quantum Gauss projection&exact&
Haar/Peter--Weyl projection\\
classical Gauss multiplier&exact if supplied&
first-order canonical action\\
source equivalence&open&
transfer/Fourier--Legendre bridge\\
companion checkpoint&next&
YM and NS audit at V35\\
gravity-dependent whitening&open&
coupled SM--Einstein Schur card
\end{array}                                                 \tag{34.19}
\]

\subsection{Parent-integrity record}

The installed parent was
\[
\begin{array}{c|c}
 \text{file}&
 \texttt{Renewal\_SM\_Einstein\_Continuum\_Research\_Notes\_V33.tex}\\
 \text{lines}&13732\\
 \text{bytes}&446910\\
 \text{SHA--256}&
 \texttt{45BF618B7FC1A1CAB07186C122FB39C276F7617E4C29B6742D154CBE6C5049EF}.
\end{array}                                                 \tag{34.20}
\]

% =============================================================================
% END APPENDED FOURTEENTH-CYCLE V34 MATERIAL
% =============================================================================
% =============================================================================
% BEGIN APPENDED FOURTEENTH-CYCLE V35 MATERIAL
% =============================================================================

\section{Fourteenth-cycle V35: compulsory companion audit and the
temporal auxiliary--constraint split}
\label{sec:v14v35}

\subsection{Files inspected}

The required thirty-fifth-version audit was performed read-only on
\[
\begin{array}{p{0.51\linewidth}|r|r}
\textbf{file}&\textbf{lines}&\textbf{bytes}\\ \hline
\texttt{Renewal\_Yang\_Mills\_Mass\_Gap\_Research\_Notes\_V10.tex}
 &2563&93475\\
\texttt{Renewal\_NS\_Stability\_Research\_Notes\_V26.tex}
 &5319&278283.
\end{array}                                                  \tag{35.1}
\]
Their SHA--256 digests are
\[
\begin{aligned}
 &\texttt{864CD554F72FDB9E7F43F83532E0D6C6F5C1AE53506676DEBFB2B16A4AD1B342},
                                                               \tag{35.2}\\
 &\texttt{82C887F8C2C69E4F28FAD7C3DC8DE5B9099CB5A4BF431EBA1FFF3E91935978AD}.
                                                               \tag{35.3}
\end{aligned}
\]
No companion file was modified or removed.

\subsection{Yang--Mills V10}

Yang--Mills V10 first proves an exact multiplier anatomy.  Relative
to a retained/constraint decomposition, an ambient Hermitian matrix
and multiplier have blocks
\[
 M=\begin{pmatrix}A&E\\E^*&D\end{pmatrix},\qquad
 Z=VZ_R+WZ_0.                                               \tag{35.4}
\]
The augmented matrix is
\[
 M+ZC+C^*Z^*
 =
 \begin{pmatrix}
 A&E+Z_R\\
 E^*+Z_R^*&D+Z_0+Z_0^*
 \end{pmatrix}.                                            \tag{35.5}
\]
The minimum-norm lift has \(Z_R=0\) and cannot cancel a nonzero cross
block.  The unique cross-cancelling value is \(Z_R=-E\), but it leaves
\(A\) unchanged.  An ambient positive completion exists exactly when
the retained compression \(A\) was already positive.  Thus multiplier
completion cannot prove its own load-bearing compression.

V10 also advances the stronger four-thirds Yang--Mills pencil from
isolated quarter-turn points to all active circles over the
transverse \(4^3\) grid, with exact rational floors.  The remaining
transverse three-torus problem is open.  Neither this Laurent
certificate nor its numerical floors belongs to the Standard-Model
temporal-link card.

The multiplier theorem is relevant to V34.  Adding a temporal or
constraint-side variable can organize cross terms, but cannot decide
whether the retained phase equation is a genuine constraint.  That
decision must be made from the temporal-variable Hessian and the
reduced phase action before any positivity completion.

\subsection{Navier--Stokes V26}

The Navier--Stokes file is unchanged.  Its rank-one Oseen-kernel
refinement provides neither a temporal link, a gauge projector, nor
an auxiliary-variable Schur complement.  Its structured-input
lesson remains compatible with restricting a mixed Hessian to the
actual unsolved multiplier kernel.

\subsection{Direction selected}

A group-valued temporal link is nonlinear.  V34 showed that it cannot
be read as an exact affine Lie-algebra multiplier.  It would also be
incorrect to call every equation obtained by varying it a phase
constraint.  If the vertical temporal-link Hessian is invertible, the
implicit-function theorem solves for that link and produces a Schur
correction to the phase Hessian; no constraint remains.  Only
vertical kernel directions can survive as unsolved multiplier
equations.

\begin{decision}[Direction after V35]
\label{dec:v14v35-next}
V36 will prove the finite auxiliary--constraint splitting theorem for
an action
\[
 \mathscr S(z,\lambda),\qquad \lambda\in\mathcal M_\lambda,
                                                               \tag{35.6}
\]
with a manifold-valued temporal variable.  At a stationary point it
will:

\begin{enumerate}
\item define the intrinsic mixed Hessian and vertical Hessian;
\item split the temporal tangent into the vertical range and kernel;
\item eliminate the invertible auxiliary block by an exact Schur
complement;
\item identify the only candidate constraint derivative as the
compression of the mixed Hessian to the vertical kernel; and
\item recover V24's affine formula when the entire vertical Hessian
vanishes.
\end{enumerate}

An explicit scalar counterexample will show that a nonzero mixed
Hessian with invertible vertical Hessian is not a constraint.  The
theorem will determine exactly what must be computed from the source
temporal-link action before either the classical or quantum V34
branch can be accepted.
\end{decision}

\subsection{Acquisition ledger}

\[
\begin{array}{p{0.30\linewidth}|p{0.20\linewidth}|p{0.39\linewidth}}
\textbf{gate}&\textbf{status after V35}&\textbf{content}\\ \hline
companion checkpoint&acquired&
YM V10 and NS V26 read with hashes\\
multiplier completion firewall&imported structurally&
retained compression unchanged\\
YM active-circle floor&not imported&
different pencil and fibre\\
temporal vertical Hessian&next&
auxiliary versus constraint directions\\
nonlinear mixed derivative&next&
kernel-compressed phase map\\
Euclidean/Hamiltonian bridge&open&
depends on V36 classification\\
gravity-dependent whitening&open&
coupled SM--Einstein card
\end{array}                                                 \tag{35.7}
\]

The next compulsory companion audit is V40.

\subsection{Parent-integrity record}

The installed parent was
\[
\begin{array}{c|c}
 \text{file}&
 \texttt{Renewal\_SM\_Einstein\_Continuum\_Research\_Notes\_V34.tex}\\
 \text{lines}&14066\\
 \text{bytes}&458739\\
 \text{SHA--256}&
 \texttt{188F5812C8F103262045D536E34A07ACDE0BFD5A03480E97B076040AAB84A873}.
\end{array}                                                 \tag{35.8}
\]

% =============================================================================
% END APPENDED FOURTEENTH-CYCLE V35 MATERIAL
% =============================================================================
% =============================================================================
% BEGIN APPENDED FOURTEENTH-CYCLE V36 MATERIAL
% =============================================================================

\section{Fourteenth-cycle V36: the nonlinear temporal
auxiliary--constraint Hessian split}
\label{sec:v14v36}

\subsection{Intrinsic Hessian blocks at a temporal stationary point}

Let \(Z\) be a finite phase/history manifold and let
\(\mathcal M_\lambda\) be a finite manifold of temporal variables.
For a \(C^3\) action
\[
 \mathscr S:Z\times\mathcal M_\lambda\longrightarrow\mathbb R,
                                                               \tag{36.1}
\]
fix \((z_0,\lambda_0)\) with
\[
 d_\lambda\mathscr S(z_0,\lambda_0)=0.                    \tag{36.2}
\]
Choose Hilbert metrics on
\[
 \mathsf Z=T_{z_0}Z,\qquad
 \mathsf Y=T_{\lambda_0}\mathcal M_\lambda.
\]
At (36.2), define
\[
 H=D_z^2\mathscr S:\mathsf Z\to\mathsf Z,\qquad
 C=D_zD_\lambda\mathscr S:\mathsf Z\to\mathsf Y,\qquad
 A=D_\lambda^2\mathscr S:\mathsf Y\to\mathsf Y.           \tag{36.3}
\]
The order in \(C\) means
\[
 \langle Cx,y\rangle
 =D_zD_\lambda\mathscr S[x,y].
\]

\begin{lemma}[Coordinate invariance of the temporal blocks]
\label{lem:v14v36-intrinsic}
At a point satisfying (36.2), \(A\) and \(C\) are intrinsic bilinear
maps under changes of temporal coordinates.  In particular their
rank, the subspace \(\ker A\), and the singular values of a
metric-whitened compression of \(C\) do not depend on the chosen
coordinate chart.
\end{lemma}

\begin{proof}
The difference between two covariant Hessians is the first derivative
of the function contracted with the difference tensor of the two
connections.  In the temporal slots that first derivative is
\(d_\lambda\mathscr S\), which vanishes by (36.2).  Equivalently, the
nonlinear second-derivative term in the coordinate-change rule is
multiplied by the temporal gradient and vanishes.  The remaining
transformation is the ordinary tensorial change of bases.  Orthogonal
whitening preserves ranks and singular values.
\end{proof}

Put
\[
 \mathsf K=\ker A,\qquad
 \mathsf R=\mathsf K^\perp=\operatorname{Ran}A,\qquad
 P_K=P_{\mathsf K},\quad P_R=I-P_K.                       \tag{36.4}
\]
Since \(A=A^*\), its restriction
\[
 A_R=A|_{\mathsf R}:\mathsf R\longrightarrow\mathsf R     \tag{36.5}
\]
is invertible.  Split
\[
 C_K=P_KC,\qquad C_R=P_RC.                                \tag{36.6}
\]

\subsection{Linearized stationarity dichotomy}

The temporal stationarity equation is
\[
 \Phi(z,\lambda):=\nabla_\lambda\mathscr S(z,\lambda)=0.
                                                               \tag{36.7}
\]
Its linearization at (36.2) is
\[
 C\,\delta z+A\,\delta\lambda=0.                           \tag{36.8}
\]

\begin{theorem}[Auxiliary--constraint splitting]
\label{thm:v14v36-split}
Write
\[
 \delta\lambda=k+r,\qquad k\in\mathsf K,\quad r\in\mathsf R.
\]
Then (36.8) is equivalent to the two equations
\[
 \boxed{C_K\,\delta z=0,}                                  \tag{36.9}
\]
\[
 \boxed{r=-A_R^{-1}C_R\,\delta z.}                         \tag{36.10}
\]
Thus \(r\) is a solved auxiliary response.  The only linearized
constraint on the phase variation is (36.9), with constraint
derivative
\[
 \boxed{C_{\rm con}=C_K=P_{\ker A}D_zD_\lambda\mathscr S.} \tag{36.11}
\]
If \(A\) is invertible, \(C_{\rm con}=0\) and the temporal equation
imposes no linearized phase constraint.  If \(A=0\), every temporal
direction is a multiplier and \(C_{\rm con}=C\), recovering the V24
affine formula.
\end{theorem}

\begin{proof}
Apply \(P_K\) and \(P_R\) to (36.8).  Since \(A\) is self-adjoint,
\(P_KA=0\), \(AP_K=0\), and \(P_RA=AP_R=A_RP_R\).  The kernel
projection gives (36.9).  The range projection gives
\[
 C_R\delta z+A_Rr=0,
\]
whose unique solution is (36.10).  The two limiting cases follow
from \(\ker A=\{0\}\) and \(\ker A=\mathsf Y\), respectively.
\end{proof}

This theorem changes V26 row L8 for a nonlinear temporal-link branch:
the required matrix is \(C_K\), not the uncompressed mixed Hessian
\(C\).

\subsection{Schur reduction of the auxiliary directions}

The implicit-function theorem applied to the range part of (36.7)
gives, near the base point and with the kernel coordinate held fixed,
a function
\[
 r=r(z,k),\qquad r(z_0,0)=0,                              \tag{36.12}
\]
solving
\[
 P_R\Phi(z,k+r(z,k))=0.                                   \tag{36.13}
\]
Let
\[
 \mathscr S_{\rm red}(z,k)
 :=\mathscr S(z,k+r(z,k)).                                \tag{36.14}
\]

\begin{theorem}[Exact second-variation Schur complement]
\label{thm:v14v36-schur}
At the base point,
\[
 D_zr=-A_R^{-1}C_R,                                       \tag{36.15}
\]
and the reduced phase Hessian is
\[
 \boxed{
 D_z^2\mathscr S_{\rm red}
 =H-C_R^*A_R^{-1}C_R.}                                    \tag{36.16}
\]
The mixed kernel row remains
\[
 D_zD_k\mathscr S_{\rm red}=C_K.                          \tag{36.17}
\]
Consequently the auxiliary block changes the retained phase
curvature through (36.16), but it cannot be counted as part of the
constraint Gram \(C_KC_K^*\).
\end{theorem}

\begin{proof}
Differentiate (36.13) in \(z\) at the base point:
\[
 C_R+A_RD_zr=0,
\]
which proves (36.15).  The chain rule for (36.14) gives
\[
 \begin{aligned}
 D_z^2\mathscr S_{\rm red}
 ={}&H+C_R^*D_zr+(D_zr)^*C_R\\
 &+(D_zr)^*A_RD_zr.
 \end{aligned}
\]
Substitution of (36.15), using \(A_R=A_R^*\), cancels one of the two
cross terms and gives (36.16).

For the mixed derivative, the direct term is \(C_K\).  Every
correction contains a block \(P_KA P_R\), which is zero because
\(\mathsf K=\ker A\) and \(A\) is self-adjoint.  This proves (36.17).
The last assertion follows from the different codomains of the Schur
term and the constraint Gram.
\end{proof}

The formula remains valid for an indefinite \(A_R\).  Positivity is
not needed for elimination; nondegeneracy is.

\subsection{Two exact witnesses}

\begin{countertheorem}[A mixed Hessian need not be a constraint]
\label{cth:v14v36-auxiliary}
Let
\[
 \mathscr S(z,a)=\frac12(a-z)^2.                           \tag{36.18}
\]
At \((0,0)\),
\[
 A=1,\qquad C=-1.                                         \tag{36.19}
\]
The temporal equation is \(a-z=0\), which solves \(a=z\); it does not
impose \(z=0\).  The reduced action is identically zero.
\end{countertheorem}

\begin{proof}
All statements follow by direct differentiation.  Since \(A\) is
invertible, Theorem \ref{thm:v14v36-split} gives
\(\mathsf K=0\) and \(C_{\rm con}=0\), while (36.16) gives
\(1-(-1)^2=0\).
\end{proof}

\begin{example}[One constraint and one auxiliary variable]
\label{ex:v14v36-partial}
For
\[
 \mathscr S(z_1,z_2;k,r)
 =kz_1+\frac12(r-z_2)^2,                                  \tag{36.20}
\]
one has
\[
 A=\begin{pmatrix}0&0\\0&1\end{pmatrix},\qquad
 C=\begin{pmatrix}1&0\\0&-1\end{pmatrix}.                 \tag{36.21}
\]
The kernel equation is \(z_1=0\), while the range equation solves
\(r=z_2\).  Thus \(C_K=(1\ 0)\), not the full matrix in (36.21).
\end{example}

These witnesses show why the nonzero mixed derivative of a Wilson
temporal link cannot be entered directly as a Gauss constraint row.

\subsection{When the kernel equation is an exact nonlinear constraint}

A zero vertical Hessian at one point is only a second-order
statement.  Higher temporal derivatives can still solve a kernel
coordinate nonlinearly.  The following normal form is sufficient for
an exact constraint.

\begin{theorem}[Partial-affine normal form]
\label{thm:v14v36-normal}
Suppose a neighbourhood has coordinates
\((z,k,r)\in Z\times\mathsf K\times\mathsf R\) in which
\[
 \mathscr S(z,k,r)
 =\mathscr S_0(z)+\langle k,\Phi_K(z)\rangle+\Psi(z,r),
                                                               \tag{36.22}
\]
and \(D_r^2\Psi\) is invertible.  Then the exact temporal equations
are
\[
 \Phi_K(z)=0,\qquad D_r\Psi(z,r)=0.                        \tag{36.23}
\]
The second equation solves \(r=r(z)\), while the first is independent
of \(k\) and has derivative
\[
 D\Phi_K=D_zD_k\mathscr S=C_{\rm con}.                     \tag{36.24}
\]
Conversely, if a temporal kernel coordinate appears nonlinearly,
Theorem \ref{thm:v14v36-split} supplies only its linearized constraint
at the base point.
\end{theorem}

\begin{proof}
Differentiate (36.22) in \(k\) and \(r\) to obtain (36.23).
Invertibility of \(D_r^2\Psi\) gives the local solution by the
implicit-function theorem.  Differentiating the first equation in
\(z\) gives (36.24).  If higher \(k\)-derivatives occur, the
\(k\)-stationarity equation can depend on \(k\); the vertical Hessian
at one point does not exclude such terms.
\end{proof}

The globally affine V27 action is the special case
\(\mathsf R=0\).  The compact-link obstruction of V34 says that a
nontrivial Wilson temporal link does not generally admit (36.22) with
the entire temporal tangent in \(\mathsf K\).

\subsection{Completed constraint card after the split}

Let
\[
 \Pi_K=P_{\ker C_K^*},\qquad
 \widehat G_K=C_KC_K^*+\Pi_K.                             \tag{36.25}
\]

\begin{proposition}[Correct positivity target]
\label{prop:v14v36-target}
If
\[
 \widehat G_K\succeq\delta_KI_{\mathsf K},\qquad
 \delta_K>0,                                               \tag{36.26}
\]
then the V16--V23 projector and angle machinery applies to the
kernel constraint bank.  Neither
\[
 CC^*\quad\hbox{nor}\quad C_R C_R^*                       \tag{36.27}
\]
can replace (36.25).  The range block belongs instead to the Schur
term in (36.16).
\end{proposition}

\begin{proof}
Theorem \ref{thm:v14v36-split} identifies the codomain of the unsolved
linearized constraint equation as \(\mathsf K\) and its derivative as
\(C_K\).  The completed-Gram theorem therefore applies to (36.25).
The range equation has already been solved by (36.10), so its
derivative cannot be counted again as an independent constraint.
Equation (36.16) shows its actual effect.
\end{proof}

This is the temporal-variable analogue of the YM V10 warning:
completion and cross cancellation do not create positivity in the
retained object.

\subsection{Application to a group-valued temporal link}

For a temporal-link action, use a left- or right-invariant metric to
identify
\[
 T_{U_0}G^{V_{\rm sp}}\cong\mathfrak g^{V_{\rm sp}}.
\]
At a temporal stationary link \(U_{0,*}\), the finite acquisition is
now the triple
\[
 \boxed{
 A_{00}=D_{U_0}^2\mathscr S,\qquad
 C_{z0}=D_zD_{U_0}\mathscr S,\qquad
 P_{\ker A_{00}}.}                                        \tag{36.28}
\]
The candidate classical constraint derivative is
\[
 C_{\rm con}=P_{\ker A_{00}}C_{z0}.                       \tag{36.29}
\]

For the \(U(1)\) witness (34.3),
\[
 A_{00}=wh^2\cos(ha+\alpha).                              \tag{36.30}
\]
At a stationary point with nonzero cosine, \(A_{00}\) is invertible,
so that plaquette variable is auxiliary at the level of the local
second variation.  It is not an affine Gauss multiplier.  Gauge
invariance of the whole time slab may create collective kernel
directions only after all temporal links, boundary variables, and
matter rows are assembled.

\subsection{Source-acquisition verdict}

\[
\begin{array}{p{0.29\linewidth}|p{0.17\linewidth}|p{0.43\linewidth}}
\textbf{datum}&\textbf{status}&\textbf{content}\\ \hline
nonlinear temporal compiler&proved&
kernel compression (36.11)\\
auxiliary elimination&proved&
Schur complement (36.16)\\
exact nonlinear constraint&conditional&
partial-affine normal form (36.22)\\
single Wilson temporal link&generically auxiliary&
nonzero vertical Hessian (36.30)\\
assembled temporal kernel&unpopulated&
requires the complete one-slab Hessian (36.28)\\
correct constraint Gram&defined&
completed Gram of \(C_K\) only\\
quantum Haar branch&unchanged&
exact invariant projector from V34\\
gravity/internal coupling&open&
must be computed after temporal kernel extraction
\end{array}                                                 \tag{36.31}
\]

\begin{decision}[Direction after V36]
\label{dec:v14v36-next}
The first missing finite source row is now the assembled one-slab
vertical Hessian \(A_{00}\), not an assumed affine temporal
coordinate.  V37 will exploit exact gauge invariance before attempting
a numerical Hessian.  It will differentiate the simultaneous gauge
action on phase and temporal variables, derive the Noether Hessian
complex
\[
 \begin{pmatrix}H&C^*\\C&A\end{pmatrix}
 \binom{\mathcal R_z\xi}{\mathcal R_\lambda\xi}=0,          \tag{36.32}
\]
and determine which temporal zero modes are pure gauge, which produce
kernel constraints after auxiliary elimination, and which require
boundary gauge fixing.  This will prevent symmetry zero modes from
being counted as physical Gauss relations.
\end{decision}

\subsection{Acquisition ledger}

\[
\begin{array}{p{0.30\linewidth}|p{0.20\linewidth}|p{0.39\linewidth}}
\textbf{gate}&\textbf{status after V36}&\textbf{content}\\ \hline
temporal auxiliary split&closed at second order&
spectral kernel/range decomposition\\
linearized constraint derivative&closed formally&
\(P_{\ker A}C\)\\
exact affine constraint&conditional&
normal form (36.22)\\
one-slab vertical Hessian&next&
gauge-covariant assembled matrix\\
symmetry zero modes&next&
Noether Hessian complex\\
bosonic/graded internal floors&retained&
apply only after kernel extraction\\
full coupled Schur card&open&
gravity-dependent mixed blocks
\end{array}                                                 \tag{36.33}
\]

\subsection{Parent-integrity record}

The installed parent was
\[
\begin{array}{c|c}
 \text{file}&
 \texttt{Renewal\_SM\_Einstein\_Continuum\_Research\_Notes\_V35.tex}\\
 \text{lines}&14223\\
 \text{bytes}&464707\\
 \text{SHA--256}&
 \texttt{3D8F16B490DB0F8568191DEBBBC5CE8FD41DC4D2911E9EA4C55205B504D4D2F2}.
\end{array}                                                 \tag{36.34}
\]

% =============================================================================
% END APPENDED FOURTEENTH-CYCLE V36 MATERIAL
% =============================================================================
% =============================================================================
% BEGIN APPENDED FOURTEENTH-CYCLE V37 MATERIAL
% =============================================================================

\section{Fourteenth-cycle V37: the gauge Noether Hessian complex after
temporal auxiliary elimination}
\label{sec:v14v37}

\subsection{The full infinitesimal Ward identity}

Let a finite-dimensional Lie group \(\mathcal G\) act smoothly on
\(Z\times\mathcal M_\lambda\), and let
\[
 \mathcal R_z(x):\mathfrak g\to T_zZ,\qquad
 \mathcal R_\lambda(x):\mathfrak g\to T_\lambda\mathcal M_\lambda
                                                               \tag{37.1}
\]
be the two components of its infinitesimal action at
\(x=(z,\lambda)\).  Suppose
\[
 \mathscr S(g\cdot z,g\cdot\lambda)=\mathscr S(z,\lambda).
                                                               \tag{37.2}
\]
Write
\[
 E(x)=d\mathscr S(x)
\]
for the full Euler covector and
\[
 \mathbb H(x)=D^2\mathscr S(x)
 =\begin{pmatrix}H&C^*\\C&A\end{pmatrix}                  \tag{37.3}
\]
in the phase/temporal splitting.

\begin{theorem}[Off-shell Noether Hessian identity]
\label{thm:v14v37-offshell}
For every \(\xi\in\mathfrak g\),
\[
 \boxed{
 \mathbb H(x)
 \binom{\mathcal R_z(x)\xi}{\mathcal R_\lambda(x)\xi}
 =
 -\bigl(D_x\mathcal R(x)[\,\cdot\,]\xi\bigr)^*E(x).}       \tag{37.4}
\]
At a full stationary point \(x_0\), this becomes the Hessian complex
\[
 \boxed{
 \begin{pmatrix}H&C^*\\C&A\end{pmatrix}
 \binom{\mathcal R_z\xi}{\mathcal R_\lambda\xi}=0.}        \tag{37.5}
\]
Equivalently,
\[
 H\mathcal R_z\xi+C^*\mathcal R_\lambda\xi=0,              \tag{37.6}
\]
\[
 C\mathcal R_z\xi+A\mathcal R_\lambda\xi=0.                \tag{37.7}
\]
\end{theorem}

\begin{proof}
Differentiate (37.2) at the group identity:
\[
 E(x)[\mathcal R(x)\xi]=0.                                \tag{37.8}
\]
Differentiate (37.8) in an arbitrary direction \(u\):
\[
 \mathbb H(x)[u,\mathcal R(x)\xi]
 +E(x)[D_x\mathcal R(x)[u]\xi]=0.
\]
Riesz representation in the first slot gives (37.4).  At a stationary
point \(E(x_0)=0\), the right side vanishes, proving (37.5).
The two block rows are (37.6)--(37.7).
\end{proof}

The right side of (37.4) is the exact off-shell Ward correction.  A
small Euler residual controls a Hessian orbit defect only together
with a bound on \(D\mathcal R\); gauge invariance alone does not make
an off-shell Hessian singular.

\subsection{Descent through the V36 split}

At a stationary point, retain the notation
\[
 \mathsf K=\ker A,\qquad \mathsf R=\mathsf K^\perp,\qquad
 C_K=P_KC,\qquad C_R=P_RC.                                \tag{37.9}
\]
For a gauge label \(\xi\), write
\[
 k_\xi=P_K\mathcal R_\lambda\xi,\qquad
 r_\xi=P_R\mathcal R_\lambda\xi.                          \tag{37.10}
\]
Define the reduced phase Hessian
\[
 H_{\rm red}=H-C_R^*A_R^{-1}C_R.                          \tag{37.11}
\]

\begin{theorem}[Reduced Noether Hessian complex]
\label{thm:v14v37-reduced}
At a full stationary point,
\[
 \boxed{C_K\mathcal R_z\xi=0,}                             \tag{37.12}
\]
\[
 \boxed{r_\xi=-A_R^{-1}C_R\mathcal R_z\xi,}                \tag{37.13}
\]
and
\[
 \boxed{
 H_{\rm red}\mathcal R_z\xi+C_K^*k_\xi=0.}                \tag{37.14}
\]
Equivalently, the reduced Hessian
\[
 \mathbb H_{\rm red}
 =
 \begin{pmatrix}H_{\rm red}&C_K^*\\C_K&0\end{pmatrix}
 :\mathsf Z\oplus\mathsf K\to\mathsf Z\oplus\mathsf K     \tag{37.15}
\]
annihilates the reduced gauge generator:
\[
 \boxed{
 \mathbb H_{\rm red}
 \binom{\mathcal R_z\xi}{k_\xi}=0.}                        \tag{37.16}
\]
\end{theorem}

\begin{proof}
Apply \(P_K\) and \(P_R\) to (37.7).  The kernel projection gives
(37.12), and the range projection gives (37.13), exactly as in V36.
Now decompose (37.6):
\[
 H\mathcal R_z\xi+C_K^*k_\xi+C_R^*r_\xi=0.
\]
Substitution of (37.13) gives (37.14).  Equations
(37.12) and (37.14) are the two rows of (37.16).
\end{proof}

Equation (37.12) is the finite tangency identity: every infinitesimal
phase gauge orbit lies in the linearized kernel constraint surface.
Equation (37.14) says that the temporal kernel component supplies the
multiplier part of the same reduced Noether zero mode.

\subsection{Stabilizers yield constraint relations}

Let
\[
 \mathfrak s_z=\ker\mathcal R_z
\]
be the phase stabilizer algebra at the stationary point.

\begin{corollary}[Symmetry-generated relation bank]
\label{cor:v14v37-relations}
The map
\[
 B_K=P_K\mathcal R_\lambda:\mathfrak g\to\mathsf K         \tag{37.17}
\]
satisfies
\[
 \boxed{
 B_K(\mathfrak s_z)\subseteq\ker C_K^*.}                  \tag{37.18}
\]
If
\[
 \operatorname{Ran}\bigl(B_K|_{\mathfrak s_z}\bigr)
 =\ker C_K^*,                                              \tag{37.19}
\]
then the orthogonal projector onto that range is the complete
relation projector for the kernel constraint bank.
\end{corollary}

\begin{proof}
If \(\xi\in\mathfrak s_z\), then \(\mathcal R_z\xi=0\).
Equation (37.14) reduces to
\[
 C_K^*k_\xi=0,
\]
and \(k_\xi=B_K\xi\), proving (37.18).  Equality (37.19) is precisely
the assertion that no further linear relations exist.
\end{proof}

Condition (37.19) must be proved by a completed-Gram floor.  Gauge
symmetry supplies declared relations but does not prove their
completeness.

\subsection{Three types of temporal kernel direction}

The temporal kernel \(\mathsf K\) has the orthogonal decomposition
\[
 \mathsf K
 =\overline{\operatorname{Ran}B_K}
 \oplus\bigl(\operatorname{Ran}B_K\bigr)^\perp.            \tag{37.20}
\]
At finite dimension the closure bar may be omitted, but it emphasizes
the later cofinal interpretation.

\begin{definition}[Gauge, relational, and additional kernel rows]
\label{def:v14v37-kernel-types}
A direction \(k\in\mathsf K\) is:

\begin{enumerate}
\item \emph{gauge-generated} if \(k=B_K\xi\) for some \(\xi\);
\item \emph{relational} if \(k=B_K\xi\) with
\(\mathcal R_z\xi=0\), in which case (37.18) gives
\(k\in\ker C_K^*\);
\item \emph{additional} if
\(k\perp\operatorname{Ran}B_K\).
\end{enumerate}
\end{definition}

A gauge-generated direction need not be a relation: when
\(\mathcal R_z\xi\ne0\), equation (37.14) couples it to a nonzero
phase gauge tangent.  An additional kernel direction is a candidate
constraint multiplier but has no gauge provenance until another
symmetry or conservation law is supplied.

\begin{proposition}[Gauge zero modes are not constraint relations]
\label{prop:v14v37-not-relation}
For \(k_\xi=B_K\xi\),
\[
 C_K^*k_\xi=-H_{\rm red}\mathcal R_z\xi.                  \tag{37.21}
\]
Thus \(k_\xi\) is a relation only if the reduced phase gauge tangent
lies in \(\ker H_{\rm red}\).  In particular
\(\mathcal R_z\xi=0\) is sufficient, but not necessary.
\end{proposition}

\begin{proof}
This is equation (37.14).  The stated criterion follows immediately.
\end{proof}

Therefore the kernel of the temporal Hessian cannot be inserted
wholesale as a constraint relation bank.  One must first compute
\(C_K\) and its adjoint.

\subsection{Gauge fixing and retained compression}

Let
\[
 \mathcal G_{\rm red}:\mathfrak g\to\mathsf Z\oplus\mathsf K,\qquad
 \mathcal G_{\rm red}\xi=
 \binom{\mathcal R_z\xi}{B_K\xi}.                          \tag{37.22}
\]
Theorem \ref{thm:v14v37-reduced} gives
\[
 \operatorname{Ran}\mathcal G_{\rm red}
 \subseteq\ker\mathbb H_{\rm red}.                        \tag{37.23}
\]
Choose a gauge-fixing map
\[
 F:\mathsf Z\oplus\mathsf K\to\mathsf G_{\rm fix}
\]
such that \(F\mathcal G_{\rm red}\) is injective on the quotient of
\(\mathfrak g\) by the full stabilizer.

\begin{proposition}[Gauge completion does not alter physical
compression]
\label{prop:v14v37-gauge-completion}
Let \(P_{\rm phys}\) project orthogonally onto
\((\operatorname{Ran}\mathcal G_{\rm red})^\perp\).  For every
positive gauge-fixing metric \(W\),
\[
 P_{\rm phys}
 \bigl(\mathbb H_{\rm red}+F^*WF\bigr)
 P_{\rm phys}
 =
 P_{\rm phys}\mathbb H_{\rm red}P_{\rm phys}
 +P_{\rm phys}F^*WFP_{\rm phys}.                          \tag{37.24}
\]
If \(F\) is chosen to vanish on the physical complement, then
\[
 \boxed{
 P_{\rm phys}
 \bigl(\mathbb H_{\rm red}+F^*WF\bigr)
 P_{\rm phys}
 =
 P_{\rm phys}\mathbb H_{\rm red}P_{\rm phys}.}            \tag{37.25}
\]
Hence gauge completion can lift gauge zero modes but cannot create a
positive physical compression.
\end{proposition}

\begin{proof}
Equation (37.24) is direct compression.  Under the stated adapted
choice, \(FP_{\rm phys}=0\), so the second term vanishes, proving
(37.25).
\end{proof}

This is the Noether form of the YM V10 multiplier firewall.  A
gauge-fixing floor must not be counted as a dynamical or constraint
floor on the physical quotient.

\subsection{A minimal example}

Let
\[
 \mathscr S(z,a)=\frac12(z-a)^2.                           \tag{37.26}
\]
It is invariant under simultaneous translations
\[
 z\mapsto z+t,\qquad a\mapsto a+t.                        \tag{37.27}
\]
At \((0,0)\),
\[
 \mathbb H=
 \begin{pmatrix}1&-1\\-1&1\end{pmatrix},\qquad
 \mathbb H\binom11=0.                                     \tag{37.28}
\]
Yet \(A=1\) is invertible, so V36 eliminates \(a=z\) and leaves no
constraint.  The gauge zero mode of the full Hessian is therefore not
evidence for a multiplier kernel.

\begin{proof}
Direct differentiation gives (37.28).  The temporal equation is
\(a-z=0\), so the conclusion follows from V36.
\end{proof}

\subsection{Application to one gauge time slab}

For a one-slab lattice action, the gauge algebra contains labels on
both time boundaries.  Temporal links transform by
\[
 U_{0,v}\longmapsto
 g_{1,v}U_{0,v}g_{0,v}^{-1}.                              \tag{37.29}
\]
At a fixed temporal link, its infinitesimal action is
\[
 \mathcal R_\lambda(\xi_0,\xi_1)_v
 =\xi_{1,v}-\operatorname{Ad}_{U_{0,v}}\xi_{0,v}           \tag{37.30}
\]
in right trivialization.  Spatial boundary fields supply
\(\mathcal R_z(\xi_0,\xi_1)\).  Equations
(37.5)--(37.16) therefore become explicit matrices once the complete
one-slab action and boundary convention are fixed.

The source active action supplies the transformation laws and exact
Ward identity, but the inspected theorem does not specify:

\begin{enumerate}
\item which boundary gauge labels are fixed, integrated, or retained;
\item the based temporal-link Hessian \(A_{00}\);
\item the phase/temporal mixed block \(C_{z0}\); or
\item the physical gauge-fixing complement used for the reduced
Hessian.
\end{enumerate}

Without these rows, a Ward identity does not determine the dimension
of \(\ker A_{00}\) or the matrix \(C_K\).

\subsection{Acquisition verdict and next direction}

\[
\begin{array}{p{0.29\linewidth}|p{0.17\linewidth}|p{0.43\linewidth}}
\textbf{datum}&\textbf{status}&\textbf{content}\\ \hline
off-shell Hessian Ward identity&proved&
equation (37.4)\\
stationary Noether complex&proved&
block kernel (37.5)\\
auxiliary-reduced complex&proved&
equations (37.12)--(37.16)\\
symmetry relation inclusion&proved&
equation (37.18)\\
relation completeness&open&
requires a completed \(C_K\) Gram\\
gauge-fixing firewall&proved&
physical compression unchanged\\
one-slab generator maps&partly explicit&
temporal action (37.30), boundary action still to be based\\
one-slab Hessian matrices&missing&
source population required
\end{array}                                                 \tag{37.31}
\]

\begin{decision}[Direction after V37]
\label{dec:v14v37-next}
V38 will make the boundary convention explicit on
\(K_4\times\{0,1\}\).  It will compute the time-incidence
cut/diagonal decomposition of (37.30), compare fixed-boundary,
diagonal-boundary, and Haar-integrated choices, and identify which
temporal-link gauge directions can remain in the V36 kernel.  This
finite two-slice complex must be fixed before the source plaquette
Hessian can be requested or certified.
\end{decision}

\subsection{Acquisition ledger}

\[
\begin{array}{p{0.30\linewidth}|p{0.20\linewidth}|p{0.39\linewidth}}
\textbf{gate}&\textbf{status after V37}&\textbf{content}\\ \hline
Noether Hessian complex&closed&
full and reduced identities\\
gauge-generated relations&identified&
stabilizer image in \(\ker C_K^*\)\\
relation completeness&open&
kernel constraint Gram\\
gauge-fixing contribution&separated&
cannot supply physical floor\\
two-boundary gauge convention&next&
finite time-incidence complex\\
temporal Hessian population&open&
after boundary convention\\
coupled SM--Einstein card&open&
after temporal kernel extraction
\end{array}                                                 \tag{37.32}
\]

\subsection{Parent-integrity record}

The installed parent was
\[
\begin{array}{c|c}
 \text{file}&
 \texttt{Renewal\_SM\_Einstein\_Continuum\_Research\_Notes\_V36.tex}\\
 \text{lines}&14649\\
 \text{bytes}&478777\\
 \text{SHA--256}&
 \texttt{C83B84F043143B6CC3C023148CFCF2F455949D06FC7AD3E67932EE0585219C46}.
\end{array}                                                 \tag{37.33}
\]

% =============================================================================
% END APPENDED FOURTEENTH-CYCLE V37 MATERIAL
% =============================================================================
% =============================================================================
% BEGIN APPENDED FOURTEENTH-CYCLE V38 MATERIAL
% =============================================================================

\section{Fourteenth-cycle V38: the two-boundary gauge complex and
temporal-gauge/Haar equivalence}
\label{sec:v14v38}

\subsection{The time-incidence map}

Let
\[
 Y_0=Y_1=\mathfrak g^{V_4},\qquad
 Y_\partial=Y_0\oplus Y_1.                                \tag{38.1}
\]
For temporal links \(U_{0,v}\in G\), define the orthogonal block
operator
\[
 \mathcal U_0=\bigoplus_{v=1}^4\operatorname{Ad}_{U_{0,v}}
 :Y_0\longrightarrow Y_1.                                \tag{38.2}
\]
The infinitesimal action of the two boundary gauge labels on the
temporal links is
\[
 D_0:Y_\partial\longrightarrow\mathfrak g^{V_4},\qquad
 D_0(\xi_0,\xi_1)=\xi_1-\mathcal U_0\xi_0.                \tag{38.3}
\]
Thus
\[
 D_0=\begin{pmatrix}-\mathcal U_0&I\end{pmatrix}.          \tag{38.4}
\]

\begin{theorem}[Exact temporal cut/diagonal decomposition]
\label{thm:v14v38-time-cut}
The map \(D_0\) satisfies
\[
 D_0D_0^*=2I,\qquad D_0^\dagger=\frac12D_0^*.             \tag{38.5}
\]
Its kernel and orthogonal range are
\[
 \ker D_0
 =\{(\xi,\mathcal U_0\xi):\xi\in Y_0\},                   \tag{38.6}
\]
\[
 \operatorname{Ran}D_0^*
 =\{(-\mathcal U_0^*\eta,\eta):
       \eta\in\mathfrak g^{V_4}\}.                         \tag{38.7}
\]
The corresponding projectors are
\[
 \boxed{
 P_{\rm cut}^{\,t}
 =D_0^\dagger D_0
 =\frac12
 \begin{pmatrix}
 I&-\mathcal U_0^*\\
 -\mathcal U_0&I
 \end{pmatrix},}                                          \tag{38.8}
\]
\[
 \boxed{
 P_{\rm diag}^{\,t}
 =I-P_{\rm cut}^{\,t}
 =\frac12
 \begin{pmatrix}
 I&\mathcal U_0^*\\
 \mathcal U_0&I
 \end{pmatrix}.}                                          \tag{38.9}
\]
All nonzero singular values of \(D_0\) equal \(\sqrt2\).
\end{theorem}

\begin{proof}
Orthogonality of \(\mathcal U_0\) gives
\[
 D_0D_0^*
 =\mathcal U_0\mathcal U_0^*+I=2I,
\]
which proves (38.5).  Equation (38.6) is the equation
\(\xi_1=\mathcal U_0\xi_0\), and (38.7) follows from
\(D_0^*\eta=(-\mathcal U_0^*\eta,\eta)\).  Multiplication of
\(\frac12D_0^*D_0\) gives (38.8), and subtraction gives (38.9).
The singular-value assertion follows from (38.5).
\end{proof}

The temporal diagonal is covariant: it identifies the final boundary
label with parallel transport of the initial label.  At
\(U_{0,v}=I\), it is the ordinary diagonal
\(\xi_1=\xi_0\).

\subsection{Boundary stabilizers and pure temporal gauge modes}

Let \(q_0,q_1\) be the spatial gauge--Higgs data on the two
boundaries and define their infinitesimal stabilizer algebras
\[
 \mathfrak h_0=\ker\mathcal R_{q_0}\subset Y_0,\qquad
 \mathfrak h_1=\ker\mathcal R_{q_1}\subset Y_1.            \tag{38.10}
\]
Their pure temporal gauge image is
\[
 \mathsf K_{\rm pg}
 :=D_0(\mathfrak h_0\oplus\mathfrak h_1)
 =\mathfrak h_1+\mathcal U_0\mathfrak h_0
 \subset\mathfrak g^{V_4}.                                \tag{38.11}
\]

\begin{theorem}[Exact pure-gauge temporal kernel]
\label{thm:v14v38-pure-gauge}
At a full stationary point of a gauge-invariant one-slab action, with
vertical temporal Hessian \(A_{00}\) and mixed block \(C_{z0}\),
\[
 \boxed{
 \mathsf K_{\rm pg}
 \subseteq\ker A_{00}\cap\ker C_{z0}^*.}                  \tag{38.12}
\]
Moreover,
\[
 \boxed{
 \dim\mathsf K_{\rm pg}
 =\dim\mathfrak h_0+\dim\mathfrak h_1
 -\dim\bigl(\mathfrak h_1\cap
             \mathcal U_0\mathfrak h_0\bigr).}            \tag{38.13}
\]
The kernel of \(D_0\) inside
\(\mathfrak h_0\oplus\mathfrak h_1\) is isomorphic to
\[
 \mathfrak h_0\cap\mathcal U_0^{-1}\mathfrak h_1,          \tag{38.14}
\]
the stabilizer of the entire temporal edge together with both
boundary configurations.
\end{theorem}

\begin{proof}
For
\((\xi_0,\xi_1)\in\mathfrak h_0\oplus\mathfrak h_1\), the phase
component of the slab gauge generator is zero and its temporal
component is \(D_0(\xi_0,\xi_1)\).  Apply the two rows of the
stationary Noether Hessian identity (37.6)--(37.7):
\[
 C_{z0}^*D_0(\xi_0,\xi_1)=0,\qquad
 A_{00}D_0(\xi_0,\xi_1)=0.
\]
This proves (38.12).

The image of the map
\((x,y)\mapsto y-\mathcal U_0x\) from
\(\mathfrak h_0\oplus\mathfrak h_1\) is the sum in (38.11).
The dimension formula for a sum of finite-dimensional subspaces gives
(38.13).  Its kernel consists of
\(y=\mathcal U_0x\) with
\(x\in\mathfrak h_0\) and
\(\mathcal U_0x\in\mathfrak h_1\), which is (38.14).
\end{proof}

The directions in (38.12) are relations of the temporal kernel
constraint row.  They must be included in its declared relation
projector before a completed-Gram floor is tested.

\subsection{The time-incidence floor is not a constraint floor}

\begin{firewall}[Different domains and different questions]
\label{fire:v14v38-no-substitution}
The singular value \(\sqrt2\) in Theorem
\ref{thm:v14v38-time-cut} conditions the boundary-gauge action on the
temporal-link tangent:
\[
 D_0:Y_0\oplus Y_1\to T_{U_0}G^{V_4}.
\]
The V36 constraint derivative has type
\[
 C_K:T_zZ_{\rm slab}\to\ker A_{00}.
\]
There is no canonical identification of these maps.  In particular,
(38.5) cannot be used as a lower bound for
\(C_KC_K^*\).
\end{firewall}

The time cut is useful for gauge coordinates and relation counting,
not for physical constraint conditioning.

\subsection{Three boundary conventions}

The finite slab admits distinct, noninterchangeable choices:
\[
\begin{array}{p{0.24\linewidth}|p{0.29\linewidth}|p{0.36\linewidth}}
\textbf{convention}&\textbf{allowed labels}&\textbf{effect}\\ \hline
fixed boundaries&
\(\xi_0=\xi_1=0\)&
no boundary gauge zero mode; the transfer kernel has fixed gauge
frames\\
one boundary fixed&
\(\xi_0=0\), \(\xi_1\) free&
\(D_0(0,\xi_1)=\xi_1\); every temporal-link tangent can be put in
temporal gauge, while the final boundary fields transform\\
both boundaries quotient&
\((\xi_0,\xi_1)\) free&
the covariant diagonal (38.6) is the residual simultaneous gauge
action; temporal-link integration becomes a boundary Haar projection\\
periodic identification&
\(\xi_0=\xi_1\)&
the temporal action is
\((I-\mathcal U_0)\xi_0\), whose kernel is the
\(\operatorname{Ad}_{U_0}\)-centralizer
\end{array}                                                 \tag{38.15}
\]
A constraint card must state which row of (38.15) it uses.  Boundary
fixing, quotienting, and periodic holonomy have different relation
spaces.

\subsection{Exact temporal-gauge/Haar identity}

Let \(\mathcal Q_0,\mathcal Q_1\) be finite boundary configuration
spaces with invariant measures, and let
\[
 K_h(q_1,q_0;u),\qquad u\in\mathcal G_X:=G^{V_4},          \tag{38.16}
\]
be a one-slab kernel satisfying the exact two-boundary covariance
\[
 K_h(g_1q_1,g_0q_0;g_1ug_0^{-1})
 =K_h(q_1,q_0;u).                                         \tag{38.17}
\]

\begin{theorem}[Temporal-link integration equals boundary projection]
\label{thm:v14v38-haar-equivalence}
With normalized Haar measure,
\[
 \boxed{
 \int_{\mathcal G_X}K_h(q_1,q_0;u)\,du
 =
 \int_{\mathcal G_X}K_h(gq_1,q_0;e)\,dg.}                 \tag{38.18}
\]
If \(T_{h,e}\) is the temporal-gauge integral operator with kernel
\(K_h(q_1,q_0;e)\), then the operator obtained by integrating all
temporal links is
\[
 \boxed{T_h=P_{{\rm inv},1}T_{h,e},}                       \tag{38.19}
\]
where \(P_{{\rm inv},1}\) is the Haar projector on the final boundary.
By covariance it may equivalently be placed on the initial boundary,
and
\[
 T_h=P_{{\rm inv},1}T_hP_{{\rm inv},0}.                   \tag{38.20}
\]
\end{theorem}

\begin{proof}
In (38.17), choose \(g_0=e\) and \(g_1=u^{-1}\).  Then
\[
 K_h(q_1,q_0;u)=K_h(u^{-1}q_1,q_0;e).                     \tag{38.21}
\]
Integrate and use invariance of Haar measure under inversion to obtain
(38.18).  Applying the two sides to a boundary function gives
(38.19).  Using (38.17) with general \(g_0,g_1\) moves the group
average between the two boundaries, and idempotence of the Haar
projectors gives (38.20).
\end{proof}

This theorem is exact at finite cutoff.  It does not require a
logarithmic expansion of \(u\) or an affine \(a_0\).  It shows where
the Gauss law lives on the quantum transfer branch: temporal-link
integration projects boundary data to gauge invariants.

\begin{corollary}[Temporal gauge does not remove Gauss projection]
\label{cor:v14v38-temporal-gauge}
Under the hypotheses of Theorem
\ref{thm:v14v38-haar-equivalence}, setting \(u=e\) is equivalent to
the original temporal-link integral only when the missing boundary
Haar projector in (38.19) is retained.
\end{corollary}

\begin{proof}
Equation (38.19) is the exact comparison.  Dropping the projector
replaces \(P_{{\rm inv},1}T_{h,e}\) by \(T_{h,e}\), which acts also on
noninvariant boundary vectors.
\end{proof}

\subsection{Source implications}

The source finite action is exactly gauge invariant, so its Boltzmann
kernel has the covariance (38.17) once one spatial time slab,
boundary fields, temporal links, and the integration measure have
been declared.  The source also contains gauge-equivariant Gibbs and
Perron--Doob rows.  The inspected action theorem does not itself
supply one based transfer operator \(T_{h,e}\), its reflection
Hilbert space, or the equality of that operator with the source's
physical clock.

The new acquisition split is therefore
\[
\begin{array}{p{0.29\linewidth}|p{0.17\linewidth}|p{0.43\linewidth}}
\textbf{datum}&\textbf{status}&\textbf{content}\\ \hline
time-incidence complex&explicit&
equations (38.3)--(38.9)\\
pure temporal gauge relations&derived&
equations (38.11)--(38.14)\\
boundary convention&must be selected&
four inequivalent rows (38.15)\\
temporal-gauge/Haar equivalence&proved&
finite identity (38.18)--(38.20)\\
source slab kernel&structurally available&
requires a declared time slicing of the common action\\
reflection Hilbert transfer&unpopulated&
kernel, measure, quotient, and positivity card\\
classical kernel constraint&unpopulated&
still requires \(A_{00},C_{z0}\) and V36 compression
\end{array}                                                 \tag{38.22}
\]

\begin{decision}[Direction after V38]
\label{dec:v14v38-next}
The quantum branch now has an exact finite Gauss projection without
an affine temporal link.  V39 will build the minimal positive
transfer interface around (38.19): reflection positivity will give
a Hilbert-space contraction, strict positivity will define its
self-adjoint logarithm, and the Haar projector will be shown to
commute with the Hamiltonian.  The note will identify the precise
kernel, faithfulness, and small-eigenvalue conditions needed before
this transfer Hamiltonian can be matched to the source physical
clock.  The classical V36 Hessian branch remains open in parallel.
\end{decision}

\subsection{Acquisition ledger}

\[
\begin{array}{p{0.30\linewidth}|p{0.20\linewidth}|p{0.39\linewidth}}
\textbf{gate}&\textbf{status after V38}&\textbf{content}\\ \hline
two-boundary time complex&closed&
exact cut/diagonal projectors\\
pure temporal gauge relations&closed formally&
boundary stabilizer image\\
temporal gauge fixing&typed&
must retain transformed boundary data\\
exact quantum Gauss projection&closed&
temporal-link/Haar identity\\
positive transfer realization&next&
reflection quotient and logarithm\\
classical temporal Hessian&open&
V36 kernel compression\\
coupled SM--Einstein card&open&
common physical clock and whitening
\end{array}                                                 \tag{38.23}
\]

\subsection{Parent-integrity record}

The installed parent was
\[
\begin{array}{c|c}
 \text{file}&
 \texttt{Renewal\_SM\_Einstein\_Continuum\_Research\_Notes\_V37.tex}\\
 \text{lines}&15070\\
 \text{bytes}&491878\\
 \text{SHA--256}&
 \texttt{5F8288914F032713C59D73258D2E8122801D99427E961BA275CDE07F58A40E15}.
\end{array}                                                 \tag{38.24}
\]

% =============================================================================
% END APPENDED FOURTEENTH-CYCLE V38 MATERIAL
% =============================================================================
% =============================================================================
% BEGIN APPENDED FOURTEENTH-CYCLE V39 MATERIAL
% =============================================================================

\section{Fourteenth-cycle V39: a reflection-positive gauge transfer
compiler and its Hamiltonian logarithm}
\label{sec:v14v39}

\subsection{Finite Osterwalder--Schrader quotient data}

Let \(\mathcal V_{+,h}\) be a finite complex vector space of
positive-time observables on one selected slab.  Let
\[
 R_h=R_h^*\succeq0                                        \tag{39.1}
\]
be its reflection Gram, so the Osterwalder--Schrader pre-inner
product is
\[
 \langle x,y\rangle_{{\rm OS},h}=x^*R_hy.                 \tag{39.2}
\]
Put
\[
 \mathcal N_h=\ker R_h,\qquad
 \mathcal H_{{\rm OS},h}
 =\mathcal V_{+,h}/\mathcal N_h.                          \tag{39.3}
\]
Let \(L_h:\mathcal V_{+,h}\to\mathcal V_{+,h}\) be the one-step
positive-time translation before quotient.

\begin{theorem}[Finite OS transfer compiler]
\label{thm:v14v39-os}
Assume
\[
 L_h\mathcal N_h\subseteq\mathcal N_h,                    \tag{39.4}
\]
\[
 R_hL_h=L_h^*R_h,                                         \tag{39.5}
\]
\[
 L_h^*R_hL_h\preceq R_h,                                  \tag{39.6}
\]
and
\[
 R_hL_h\succeq0.                                          \tag{39.7}
\]
Then \(L_h\) descends to a well-defined operator
\(T_{0,h}\) on \(\mathcal H_{{\rm OS},h}\), and
\[
 \boxed{T_{0,h}=T_{0,h}^*,\qquad
        0\preceq T_{0,h}\preceq I.}                       \tag{39.8}
\]
Every assertion is equivalent to a finite matrix relation on the
source slab.
\end{theorem}

\begin{proof}
Condition (39.4) makes the quotient action
\[
 T_{0,h}[x]=[L_hx]
\]
well-defined.  Equation (39.5) gives
\[
 \langle[x],T_{0,h}[y]\rangle_{\rm OS}
 =x^*R_hL_hy
 =(L_hx)^*R_hy
 =\langle T_{0,h}[x],[y]\rangle_{\rm OS},
\]
so \(T_{0,h}\) is self-adjoint.  Equation (39.6) gives
\[
 \|T_{0,h}[x]\|_{\rm OS}^2
 =x^*L_h^*R_hL_hx
 \le x^*R_hx,
\]
hence \(\|T_{0,h}\|\le1\).  Finally (39.7) gives
\[
 \langle[x],T_{0,h}[x]\rangle_{\rm OS}=x^*R_hL_hx\ge0.
\]
A positive self-adjoint contraction satisfies (39.8).
\end{proof}

Reflection positivity alone gives (39.1).  The descent, reflection
symmetry, contraction, and one-step positivity rows
(39.4)--(39.7) remain separate finite checks.

\subsection{Gauge descent and temporal-link integration}

Suppose the boundary gauge group \(\mathcal G_h\) acts on
\(\mathcal V_{+,h}\) through \(U_h(g)\) and
\[
 U_h(g)^*R_hU_h(g)=R_h,\qquad
 U_h(g)\mathcal N_h\subseteq\mathcal N_h.                 \tag{39.9}
\]
Then the action descends unitarily to
\(\mathcal H_{{\rm OS},h}\).  Let
\[
 P_h=\int_{\mathcal G_h}U_h(g)\,dg                        \tag{39.10}
\]
denote the descended Haar projector.

\begin{theorem}[Gauge-projected positive transfer]
\label{thm:v14v39-gauge}
Assume Theorem \ref{thm:v14v39-os}, (39.9), and
\[
 L_hU_h(g)=U_h(g)L_h\qquad(g\in\mathcal G_h).              \tag{39.11}
\]
Then \(P_h\) is an orthogonal projector,
\[
 [P_h,T_{0,h}]=0,                                         \tag{39.12}
\]
and
\[
 T_h=P_hT_{0,h}P_h                                        \tag{39.13}
\]
is a positive self-adjoint contraction with
\[
 \operatorname{Ran}T_h\subseteq\mathcal H_{{\rm OS},h}^{\mathcal G_h}.
                                                               \tag{39.14}
\]
On the invariant Hilbert space
\[
 \mathcal H_{{\rm phys},h}=P_h\mathcal H_{{\rm OS},h},     \tag{39.15}
\]
the physical transfer is simply
\[
 T_{{\rm phys},h}
 =T_{0,h}|_{\mathcal H_{{\rm phys},h}}.                   \tag{39.16}
\]
Under the V38 slab covariance, (39.13) is exactly the operator
obtained by integrating the temporal links.
\end{theorem}

\begin{proof}
Equation (39.9) makes the descended representation unitary, so Haar
averaging gives an orthogonal projector.  Equation (39.11) passes to
the quotient and integration gives (39.12).  Commuting compression
of a positive self-adjoint contraction is again a positive
self-adjoint contraction, proving (39.13)--(39.14).  On the range of
\(P_h\), both projectors in (39.13) act as the identity, proving
(39.16).  The final assertion is V38
Theorem \ref{thm:v14v38-haar-equivalence}.
\end{proof}

The full-space operator \(T_h\) has the entire noninvariant subspace
in its kernel.  Its logarithm must therefore be taken on
\(\mathcal H_{{\rm phys},h}\), not on the unprojected Hilbert space.

\subsection{The Hamiltonian logarithm}

\begin{theorem}[Finite physical Hamiltonian]
\label{thm:v14v39-log}
Suppose, in addition, that
\[
 T_{{\rm phys},h}\succeq m_hI_{\mathcal H_{{\rm phys},h}},
 \qquad m_h>0.                                             \tag{39.17}
\]
Then
\[
 \boxed{
 H_{{\rm phys},h}
 :=-\frac1h\log T_{{\rm phys},h}}                          \tag{39.18}
\]
is a bounded positive self-adjoint operator and
\[
 T_{{\rm phys},h}=e^{-hH_{{\rm phys},h}},\qquad
 0\preceq H_{{\rm phys},h}
 \preceq\frac{|\log m_h|}{h}I.                            \tag{39.19}
\]
It commutes with every boundary gauge transformation and with every
infinitesimal gauge generator on the invariant carrier.
\end{theorem}

\begin{proof}
The spectrum of the positive contraction in (39.17) lies in
\([m_h,1]\).  Continuous functional calculus therefore defines
\(\log T_{{\rm phys},h}\), which is self-adjoint and has spectrum in
\([\log m_h,0]\).  Equations (39.18)--(39.19) follow.  Since the
transfer commutes with the gauge representation, so does every
continuous function of it.  Differentiating the group commutation
gives commutation with the infinitesimal generators.
\end{proof}

If \(T_{{\rm phys},h}\) is only semidefinite, functional calculus
defines a logarithm on its support
\[
 S_h=P_{(0,1]}(T_{{\rm phys},h}),                          \tag{39.20}
\]
but states in \(\ker T_{{\rm phys},h}\) are removed.  A source card
must state whether that support reduction is physical or an artifact
of a finite transfer kernel.

\subsection{Faithfulness is not a mass gap}

Let the transfer eigenvalues on the physical space be
\[
 1\ge\lambda_{0,h}\ge\lambda_{1,h}\ge\cdots>0.             \tag{39.21}
\]
After vacuum normalization, the finite excitation gap is
\[
 \Delta_h
 =-\frac1h\log\frac{\lambda_{1,h}}{\lambda_{0,h}}.         \tag{39.22}
\]

\begin{countertheorem}[A logarithm floor gives no excitation gap]
\label{cth:v14v39-no-gap}
For \(T=I_n\),
\[
 T\succeq I_n,\qquad -h^{-1}\log T=0,
\]
but every transfer eigenvalue equals one and \(\Delta_h=0\).  More
generally, a lower bound on the smallest transfer eigenvalue controls
the largest finite energy, not separation of the two largest
transfer eigenvalues.
\end{countertheorem}

\begin{proof}
Direct spectral calculation gives the first statement.  Formula
(39.18) reverses the spectral order: small transfer eigenvalues give
large energies, whereas (39.22) depends on the top two transfer
eigenvalues.
\end{proof}

Thus (39.17) is a faithfulness condition for the logarithm.  Vacuum
uniqueness, a positive excitation gap, particle poles, and the full
quantum continuum require separate spectral and cofinal arguments.

\subsection{Stable logarithms and clock matching}

\begin{lemma}[Operator-logarithm perturbation bound]
\label{lem:v14v39-log-bound}
If
\[
 mI\preceq A,B\preceq I,\qquad m>0,
\]
then
\[
 \boxed{\|\log A-\log B\|\le m^{-1}\|A-B\|.}              \tag{39.23}
\]
\end{lemma}

\begin{proof}
Use the resolvent representation
\[
 \log A-\log B
 =\int_0^\infty
 \bigl((B+tI)^{-1}-(A+tI)^{-1}\bigr)\,dt
\]
and the resolvent identity.  Its norm is at most
\[
 \|A-B\|\int_0^\infty(m+t)^{-2}\,dt
 =m^{-1}\|A-B\|.
\]
\end{proof}

Let \(T_{{\rm clock},h}\) be the transfer induced by the source's
declared physical clock after transport to the same invariant OS
screen.  Define
\[
 \epsilon_{{\rm clock},h}
 =\|T_{{\rm phys},h}-T_{{\rm clock},h}\|.                 \tag{39.24}
\]

\begin{corollary}[Generator matching criterion]
\label{cor:v14v39-clock}
If both transfers lie in \([mI,I]\) on a fixed screen, then
\[
 \left\|
 -h^{-1}\log T_{{\rm phys},h}
 +h^{-1}\log T_{{\rm clock},h}
 \right\|
 \le\frac{\epsilon_{{\rm clock},h}}{mh}.                  \tag{39.25}
\]
Hence
\[
 \epsilon_{{\rm clock},h}=o(h)                            \tag{39.26}
\]
is sufficient for equality of their limiting generators on that
screen.
\end{corollary}

\begin{proof}
Apply Lemma \ref{lem:v14v39-log-bound} and divide by \(h\).
\end{proof}

A zero accepted/action gap or equality of one stationary row does not
imply (39.26).  The physical-clock comparison must be made at the
operator level on the same quotient.

\subsection{Exact acquisition packet}

A finite quantum transfer card at scale \(h\) now consists of
\[
 \mathfrak T_h=
 \bigl(
 [R_h],[L_h],[U_h(g)],[\mathcal N_h],
 [P_h],[T_{{\rm phys},h}],m_h,
 \lambda_{0,h},\lambda_{1,h},
 \epsilon_{{\rm clock},h}
 \bigr).                                                   \tag{39.27}
\]
The acceptance rows are
\[
\begin{array}{p{0.30\linewidth}|p{0.20\linewidth}|p{0.39\linewidth}}
\textbf{row}&\textbf{certificate}&\textbf{purpose}\\ \hline
OS positivity&\(R_h\succeq0\)&reflection quotient\\
transfer descent&(39.4)--(39.7)&positive contraction\\
gauge covariance&(39.9)--(39.11)&Haar projection\\
physical faithfulness&\(T_{\rm phys}\succeq m_hI\)&
Hamiltonian logarithm\\
vacuum/gap spectrum&(39.21)--(39.22)&
physical excitation scale\\
clock match&\(\epsilon_{{\rm clock},h}=o(h)\)&
same continuum generator
\end{array}                                                 \tag{39.28}
\]

The source contains reflection, gauge, common-action, and physical
clock structures, but the inspected material does not populate all
matrices in (39.27) on one selected \(K_4\) time slab.

\begin{decision}[Direction after V39]
\label{dec:v14v39-next}
V40 is the compulsory companion audit.  It will check whether the
newest Yang--Mills or Navier--Stokes notes supply a transfer
faithfulness, quotient logarithm, or operator clock-matching
certificate.  V41 will then address the remaining interface between
the quantum transfer branch and the classical constraint branch:
the semiclassical principal symbol of
\(H_{{\rm phys},h}\) must reproduce the V29 Gauss momentum map and the
V27 Hamiltonian action on a common phase chart.
\end{decision}

\subsection{Acquisition ledger}

\[
\begin{array}{p{0.30\linewidth}|p{0.20\linewidth}|p{0.39\linewidth}}
\textbf{gate}&\textbf{status after V39}&\textbf{content}\\ \hline
finite OS quotient&compiler proved&
matrix rows (39.1)--(39.7)\\
gauge-projected transfer&compiler proved&
Haar compression\\
physical Hamiltonian logarithm&conditional&
strict transfer faithfulness\\
mass gap&separate and open&
top spectral ratio\\
physical clock match&defined&
operator error \(o(h)\)\\
literal source transfer card&unpopulated&
matrices (39.27)\\
classical temporal Hessian branch&open&
V36--V38
\end{array}                                                 \tag{39.29}
\]

\subsection{Parent-integrity record}

The installed parent was
\[
\begin{array}{c|c}
 \text{file}&
 \texttt{Renewal\_SM\_Einstein\_Continuum\_Research\_Notes\_V38.tex}\\
 \text{lines}&15430\\
 \text{bytes}&503943\\
 \text{SHA--256}&
 \texttt{E206F00D53C4CE00DE317AD546DC3A602E39AB24EE88ED46C320D7E9328B48F5}.
\end{array}                                                 \tag{39.30}
\]

% =============================================================================
% END APPENDED FOURTEENTH-CYCLE V39 MATERIAL
% =============================================================================
% =============================================================================
% BEGIN APPENDED FOURTEENTH-CYCLE V40 MATERIAL
% =============================================================================

\section{Fourteenth-cycle V40: compulsory companion audit and the
finite-atlas certification lesson}
\label{sec:v14v40}

\subsection{Files inspected}

The fortieth-version audit was performed read-only on the newest
active companion notes
\[
\begin{array}{p{0.51\linewidth}|r|r}
\textbf{file}&\textbf{lines}&\textbf{bytes}\\ \hline
\texttt{Renewal\_Yang\_Mills\_Mass\_Gap\_Research\_Notes\_V12.tex}
 &2964&108936\\
\texttt{Renewal\_NS\_Stability\_Research\_Notes\_V26.tex}
 &5319&278283.
\end{array}                                                  \tag{40.1}
\]
Their SHA--256 digests are
\[
\begin{aligned}
 &\texttt{4EF8F50A20C51A4B6BCB8E465A05412C83D1AD6B9A06578BD94FF8250B2A2EB4},
                                                               \tag{40.2}\\
 &\texttt{82C887F8C2C69E4F28FAD7C3DC8DE5B9099CB5A4BF431EBA1FFF3E91935978AD}.
                                                               \tag{40.3}
\end{aligned}
\]
The Yang--Mills file changed after the first inventory made during
this interrupted iteration: the live V12 is a substantive successor
of V11, not a duplicate.  The values in (40.1)--(40.3) are those of
the files actually audited for this installed version.  No companion
file was modified or removed.

\subsection{Yang--Mills V12: what transfers and what does not}

Yang--Mills V12 gives an exact rational-phase counterexample to the
full-cardinal interpolation of the V10 strong-pencil splittings.  Its
negative vector belongs to the doubled certificate block, so V12
correctly limits the conclusion: the selected global interpolation
fails, while the underlying Yang--Mills pencil is not thereby shown
negative.

The useful replacement is an independent local-splitting cover.
At each transverse point one may hold one positive splitting fixed
on a neighbourhood.  Compactness reduces a uniform positivity proof
to finitely many such neighbourhoods, and no agreement of the
splittings on overlaps is needed.  The quantitative debit on a patch
is the operator norm of the change of the diagonal and Laurent
off-diagonal blocks.

This method is adopted as a certification pattern for future
Standard-Model screen estimates.  It does not populate the present
transfer card.  The Yang--Mills block, its splitting variables, its
rational witness, and its numerical floors are different from
\[
 R_h,\quad L_h,\quad T_{{\rm phys},h},\quad
 T_{{\rm clock},h}                                         \tag{40.4}
\]
of V39.  In particular, failure of one auxiliary splitting
interpolation is neither failure of reflection positivity nor
failure of the physical transfer.

\begin{proposition}[Scope of a local certificate import]
\label{prop:v14v40-scope}
Let \(K\) be compact and let \(\mathcal Q(x;X)\) be a continuous
Hermitian certificate depending on a base point \(x\in K\) and an
auxiliary variable \(X\).  If for every \(x\) there are \(X_x\) and
\(\eta_x>0\) such that
\[
 \mathcal Q(x;X_x)\succeq 2\eta_x I,                       \tag{40.5}
\]
then a finite collection of fixed auxiliary values certifies
\[
 \mathcal Q(y;X_{x_j})\succeq\eta_{x_j}I
 \quad\text{on open sets }U_{x_j},\qquad
 K=\bigcup_{j=1}^J U_{x_j}.                                \tag{40.6}
\]
This conclusion supplies a positivity proof for \(\mathcal Q\) only.
It supplies no identity between \(\mathcal Q\) and a different
operator family \(\widetilde{\mathcal Q}\).
\end{proposition}

\begin{proof}
Continuity and the open character of strict positive definiteness
give, for each \(x\), a neighbourhood \(U_x\) on which
\(\|\mathcal Q(y;X_x)-\mathcal Q(x;X_x)\|<\eta_x\).
Weyl's inequality gives (40.6).  Compactness supplies a finite
subcover.  The hypotheses contain no relation to
\(\widetilde{\mathcal Q}\), so an assertion about that family cannot
follow.  This last observation is the precise firewall preventing
the Yang--Mills certificate from being substituted for (40.4).
\end{proof}

\subsection{Navier--Stokes V26}

The Navier--Stokes file is unchanged from the V35 audit.  Its final
active block computes rank-one norms of Oseen-kernel derivatives and
uses them to refine ancient-solution estimates.  It contains no
gauge Haar projector, OS reflection matrix, one-step transfer,
physical-transfer faithfulness floor, or comparison with the
source's Standard-Model clock.  Its older report of a
variable-clock result in a then-current Yang--Mills note is
provenance about that historical companion, not a transfer theorem
for the present \(K_4\) Standard-Model screen.

\subsection{Result of the transfer-card search}

The audit therefore leaves every literal acquisition in
\[
 \mathfrak T_h=
 \bigl(
 [R_h],[L_h],[U_h(g)],[\mathcal N_h],
 [P_h],[T_{{\rm phys},h}],m_h,
 \lambda_{0,h},\lambda_{1,h},
 \epsilon_{{\rm clock},h}
 \bigr)                                                    \tag{40.7}
\]
unpopulated.  The reusable result is methodological: once a common
SM operator family has been extracted, strict local certificates may
be assembled as a finite atlas without interpolating their auxiliary
variables.

The next obstruction lies before any cofinal use of (40.7).  V27 and
V29 produced a classical Hamiltonian action and a Gauss momentum map,
whereas V39 produced a conditional quantum transfer Hamiltonian.
They have not yet been shown to be quantizations of one another.

\begin{decision}[Direction after V40]
\label{dec:v14v40-next}
V41 will build the finite semiclassical handoff compiler.  It will
distinguish
\[
 -h^{-1}\log \langle T\rangle_z
 \quad\text{from}\quad
 \langle -h^{-1}\log T\rangle_z,                           \tag{40.8}
\]
quantify their Jensen gap by the transfer variance, and prove the
extra visibility condition required after Haar projection.  It will
then state the exact symbol, commutator, and constraint rows needed
for the V39 quantum generator to reproduce the V27--V29 classical
constrained dynamics.  No coherent-state or quantization data absent
from the source will be invented.
\end{decision}

\subsection{Acquisition ledger}

\[
\begin{array}{p{0.30\linewidth}|p{0.20\linewidth}|p{0.39\linewidth}}
\textbf{gate}&\textbf{status after V40}&\textbf{content}\\ \hline
companion checkpoint&acquired&
YM V12 and NS V26 read with live hashes\\
global YM interpolation&failed there&
exact auxiliary-certificate counterexample\\
independent local atlas&imported structurally&
finite-cover positivity method\\
literal SM transfer card&unpopulated&
no companion substitution\\
semiclassical Hamiltonian symbol&next&
operator logarithm versus scalar symbol\\
projected-state visibility&next&
Haar quotient localization\\
cofinal continuum and mass gap&open&
requires source cards and uniform estimates
\end{array}                                                 \tag{40.9}
\]

The next compulsory companion audit is V45.

\subsection{Parent-integrity record}

The installed parent was
\[
\begin{array}{c|c}
 \text{file}&
 \texttt{Renewal\_SM\_Einstein\_Continuum\_Research\_Notes\_V39.tex}\\
 \text{lines}&15800\\
 \text{bytes}&515494\\
 \text{SHA--256}&
 \texttt{B24EB09EF40AC789BF69FF8A96C8223425FB5B6DB111F5BE376A52510F37E412}.
\end{array}                                                 \tag{40.10}
\]

% =============================================================================
% END APPENDED FOURTEENTH-CYCLE V40 MATERIAL
% =============================================================================
% =============================================================================
% BEGIN APPENDED FOURTEENTH-CYCLE V41 MATERIAL
% =============================================================================

\section{Fourteenth-cycle V41: the finite semiclassical transfer--constraint
handoff}
\label{sec:v14v41}

\subsection{Purpose and two small parameters}

V27 produced a classical Hamiltonian action after the Legendre
transform, V29 produced the finite bosonic Gauss momentum map, and
V39 produced a conditional quantum Hamiltonian from a faithful
physical transfer.  These are three different constructions until a
symbol map identifies them on one phase chart.

Write \(h>0\) for the time-slab width and \(\hbar>0\) for the
semiclassical parameter.  They must not be identified silently.  At
a finite screen let
\[
 {\cal H}_{\hbar,h},\qquad U_{\hbar,h}(g),\qquad
 P_{\hbar,h}=\int_G U_{\hbar,h}(g)\,dg                       \tag{41.1}
\]
be a Hilbert space, a unitary representation of the compact gauge
group, and its normalized Haar projector.  Let
\[
 0\prec m_{\hbar,h}I
 \preceq T_{\hbar,h}\preceq I,\qquad
 [T_{\hbar,h},U_{\hbar,h}(g)]=0,                            \tag{41.2}
\]
and set, on the invariant carrier,
\[
 \widehat H_{\hbar,h}
 =-h^{-1}\log\!\left(
 T_{\hbar,h}|_{P_{\hbar,h}{\cal H}_{\hbar,h}}\right).       \tag{41.3}
\]
All results below are finite-dimensional or bounded-operator
statements.  They do not assert that (41.1)--(41.3) have been
extracted from the source.

For a normalized vector \(\psi_z\) labelled by a point \(z\) of a
classical phase chart, define the unprojected covariant symbol
\[
 \sigma_{\hbar,h}(A)(z)=\langle\psi_z,A\psi_z\rangle.       \tag{41.4}
\]
When
\[
 v_{\hbar,h}(z)
 :=\langle\psi_z,P_{\hbar,h}\psi_z\rangle
 =\|P_{\hbar,h}\psi_z\|^2>0,                               \tag{41.5}
\]
define
\[
 \psi_z^{P}
 =v_{\hbar,h}(z)^{-1/2}P_{\hbar,h}\psi_z,\qquad
 \sigma_{\hbar,h}^{P}(A)(z)
 =\langle\psi_z^{P},A\psi_z^{P}\rangle.                    \tag{41.6}
\]

\subsection{The scalar transfer logarithm is not the Hamiltonian symbol}

For a normalized physical vector \(\phi\), put
\[
 \tau_\phi=\langle\phi,T_{\hbar,h}\phi\rangle,\qquad
 E_\phi^{\rm tr}=-h^{-1}\log\tau_\phi,                     \tag{41.7}
\]
and
\[
 {\rm Var}_\phi(T_{\hbar,h})
 =\langle\phi,T_{\hbar,h}^2\phi\rangle-\tau_\phi^2.        \tag{41.8}
\]

\begin{theorem}[Jensen--variance transfer bound]
\label{thm:v14v41-jensen}
Under (41.2), for every normalized physical \(\phi\),
\[
 \boxed{
 0\le
 \langle\phi,\widehat H_{\hbar,h}\phi\rangle
 -E_\phi^{\rm tr}
 \le
 \frac{{\rm Var}_\phi(T_{\hbar,h})}
      {2h\,m_{\hbar,h}^{\,2}}.}                            \tag{41.9}
\]
Thus the logarithm of a transfer expectation and the expectation of
the operator logarithm agree only after a variance estimate.
\end{theorem}

\begin{proof}
Let \(\mu_\phi\) be the spectral probability measure of
\(T_{\hbar,h}\) in the state \(\phi\).  It is supported in
\([m_{\hbar,h},1]\), and
\[
 \tau_\phi=\int\lambda\,d\mu_\phi(\lambda),\qquad
 \langle\phi,\widehat H_{\hbar,h}\phi\rangle
 =h^{-1}\int-\log\lambda\,d\mu_\phi(\lambda).              \tag{41.10}
\]
Convexity of \(f(\lambda)=-\log\lambda\) gives the lower bound in
(41.9).  Since
\[
 0<f''(\lambda)=\lambda^{-2}\le m_{\hbar,h}^{-2},          \tag{41.11}
\]
Taylor's theorem at \(\tau_\phi\) gives
\[
 f(\lambda)
 \le f(\tau_\phi)+f'(\tau_\phi)(\lambda-\tau_\phi)
 +\frac{(\lambda-\tau_\phi)^2}{2m_{\hbar,h}^{\,2}}.        \tag{41.12}
\]
Integrate.  The linear term vanishes and the quadratic term is
(41.8).  Division by \(h\) proves (41.9).
\end{proof}

\begin{corollary}[Bounded-energy variance debit]
\label{cor:v14v41-energy-variance}
Suppose \(0\preceq\widehat H\preceq MI\) and
\(T=e^{-h\widehat H}\), so \(m=e^{-hM}\).  Then
\[
 {\rm Var}_{\phi}(T)
 \le h^2{\rm Var}_{\phi}(\widehat H),                      \tag{41.13}
\]
and consequently
\[
 0\le\langle\widehat H\rangle_\phi-E_\phi^{\rm tr}
 \le\frac{h e^{2hM}}2\,{\rm Var}_{\phi}(\widehat H).       \tag{41.14}
\]
\end{corollary}

\begin{proof}
For a real random variable \(X\) with the spectral law of
\(\widehat H\), and an independent copy \(X'\),
\[
 {\rm Var}(F(X))
 =\frac12{\mathbb E}|F(X)-F(X')|^2.                        \tag{41.15}
\]
On \([0,M]\), \(F(x)=e^{-hx}\) has Lipschitz constant at most \(h\).
Equation (41.13) follows from (41.15), and (41.14) follows from
Theorem~\ref{thm:v14v41-jensen}.
\end{proof}

The direct sufficient condition along a cofinal family is
\[
 \sup_{z\in K}
 \frac{{\rm Var}_{\psi_z^P}(T_{\hbar,h})}
      {h\,m_{\hbar,h}^{\,2}}\longrightarrow0              \tag{41.16}
\]
on every fixed compact physical screen \(K\).  A uniform energy
bound and a uniform energy-variance bound make (41.16) automatic as
\(h\to0\), by (41.14).  Transfer faithfulness alone does not.

\begin{countertheorem}[Strict faithfulness does not remove the Jensen gap]
\label{cth:v14v41-jensen-gap}
Let
\[
 T=\begin{pmatrix}1&0\\0&m\end{pmatrix},\qquad
 \phi=2^{-1/2}(1,1),\qquad 0<m<1.
\]
Then \(T\succeq mI\), but
\[
 \langle\widehat H\rangle_\phi
 =-\frac{\log m}{2h}
 >
 -\frac1h\log\frac{1+m}{2}
 =E_\phi^{\rm tr}.                                        \tag{41.17}
\]
\end{countertheorem}

\begin{proof}
The formulas are direct.  Strict inequality is strict convexity of
\(-\log\), because the spectral measure has two distinct atoms.
\end{proof}

\subsection{Haar visibility and stability of projected symbols}

\begin{lemma}[Projected-symbol ratio]
\label{lem:v14v41-projected-ratio}
Let \(P\) be an orthogonal projection, \(\|\psi\|=1\), and
\(v=\langle\psi,P\psi\rangle>0\).  For any bounded \(A\),
\[
 \langle\psi^P,A\psi^P\rangle
 =\frac{\langle\psi,PAP\psi\rangle}{v}.                    \tag{41.18}
\]
If \(v\ge v_0>0\) and, for a scalar \(a\),
\[
 \bigl|\langle\psi,P(A-aI)P\psi\rangle\bigr|\le\epsilon,   \tag{41.19}
\]
then
\[
 \bigl|\langle\psi^P,A\psi^P\rangle-a\bigr|
 \le\epsilon/v_0.                                         \tag{41.20}
\]
\end{lemma}

\begin{proof}
Insert \(\psi^P=P\psi/\sqrt v\).  This gives (41.18);
subtract \(a\) and use \(v\ge v_0\) to obtain (41.20).
\end{proof}

Thus an absolute numerator estimate is not a physical-symbol
estimate without a lower bound for the Haar visibility \(v\).

\begin{countertheorem}[Haar projection can destroy localization]
\label{cth:v14v41-visibility}
Let \(G=U(1)\) act on \(\mathbb C^2\) by
\[
 U(e^{i\theta})=\begin{pmatrix}1&0\\0&e^{i\theta}\end{pmatrix},
 \qquad P=\begin{pmatrix}1&0\\0&0\end{pmatrix}.            \tag{41.21}
\]
For \(0<\epsilon<1\), set
\[
 \psi_\epsilon=(\epsilon,\sqrt{1-\epsilon^2}).
\]
Then
\[
 v(\epsilon)=\epsilon^2,\qquad
 \psi_\epsilon^P=(1,0),                                   \tag{41.22}
\]
so every projected state is identical and the projection becomes
undefined at the limiting vector \((0,1)\).  For the invariant
observable \(A=\operatorname{diag}(0,1)\),
\[
 \langle\psi_\epsilon,A\psi_\epsilon\rangle=1-\epsilon^2,
 \qquad
 \langle\psi_\epsilon^P,A\psi_\epsilon^P\rangle=0.         \tag{41.23}
\]
\end{countertheorem}

\begin{proof}
Normalized Haar integration gives the matrix \(P\) in (41.21).
Equations (41.22)--(41.23) follow by multiplication.
\end{proof}

This finite example rules out an inference from the existence of a
Haar projector to preservation of a chosen classical localization.
A source acquisition must either prove a visibility floor on its
selected chart or construct coherent states directly on the reduced
phase space.

\subsection{Why the Gauss symbol must be taken before projection}

Choose the generator convention
\[
 U_{\hbar,h}(\exp t\xi)
 =\exp\!\left(-\frac{it}{\hbar}
              \widehat\Phi_{\hbar,h}(\xi)\right).          \tag{41.24}
\]

\begin{proposition}[Projected Gauss symbols vanish identically]
\label{prop:v14v41-projected-gauss}
For every \(\xi\) in the gauge Lie algebra,
\[
 \widehat\Phi_{\hbar,h}(\xi)P_{\hbar,h}
 =P_{\hbar,h}\widehat\Phi_{\hbar,h}(\xi)=0.                \tag{41.25}
\]
Consequently, whenever \(\psi_z^P\) exists,
\[
 \sigma_{\hbar,h}^P
   \bigl(\widehat\Phi_{\hbar,h}(\xi)\bigr)(z)=0.           \tag{41.26}
\]
Hence projected coherent states cannot recover the off-shell
classical momentum map
\(\langle\Phi_G(z),\xi\rangle\).
\end{proposition}

\begin{proof}
Haar invariance gives
\[
 U_{\hbar,h}(\exp t\xi)P_{\hbar,h}=P_{\hbar,h}
 =P_{\hbar,h}U_{\hbar,h}(\exp t\xi).                       \tag{41.27}
\]
Differentiate at \(t=0\) using (41.24), which proves (41.25).
Equation (41.26) follows.  The V29 momentum map is generally nonzero
away from its zero level, whereas the left side of (41.26) is zero
at every label.  Therefore equality cannot hold off shell.
\end{proof}

There are precisely two consistent routes:
\[
\begin{array}{p{0.27\linewidth}|p{0.29\linewidth}|p{0.34\linewidth}}
\textbf{route}&\textbf{symbol carrier}&\textbf{classical target}\\ \hline
Dirac before reduction&
unprojected covariant states&
full map \(\Phi_G\), then impose \(\Phi_G=0\)\\
quantize after reduction&
coherent states on reduced phase space&
physical Hamiltonian on \(\Phi_G^{-1}(0)/G\).
\end{array}                                                 \tag{41.28}
\]
Mixing the projected carrier of the second row with the off-shell
constraint target of the first row would force a false zero symbol.

\subsection{A conditional finite handoff theorem}

Let \(Z\) be the finite classical phase chart of V27--V29, let
\[
 \Phi_\xi(z)=\langle\Phi_G(z),\xi\rangle,\qquad
 Z_0=\Phi_G^{-1}(0),                                       \tag{41.29}
\]
and let \(H_{\rm cl}\) be the classical Hamiltonian obtained from the
V27 Legendre branch.  Let \(K\subset Z\) be compact and put
\(K_0=K\cap Z_0\).  Suppose a cofinal family
\(\alpha\mapsto(\hbar_\alpha,h_\alpha)\) supplies normalized
unprojected states \(\psi_{\alpha,z}\), the data
(41.1)--(41.3), and quantum constraints
\(\widehat\Phi_\alpha(\xi)\).

Define the four errors
\[
\begin{aligned}
 \epsilon_{H,\alpha}
 &=\sup_{z\in K_0}
   \left|
   \sigma_\alpha^P(\widehat H_\alpha)(z)-H_{\rm cl}(z)
   \right|,                                                \tag{41.30}\\
 \epsilon_{G,\alpha}
 &=\sup_{\substack{z\in K\\\|\xi\|\le1}}
   \left|
   \sigma_\alpha(\widehat\Phi_\alpha(\xi))(z)-\Phi_\xi(z)
   \right|,                                                \tag{41.31}\\
 \epsilon_{{\rm dyn},\alpha}(a)
 &=\sup_{z\in K}
   \left|
   \sigma_\alpha\!\left(
    \frac{1}{i\hbar_\alpha}
    [\operatorname{Op}_\alpha(a),\widehat H_\alpha]
   \right)(z)
   -\{a,H_{\rm cl}\}(z)
   \right|,                                                \tag{41.32}\\
 \delta_{{\rm tr},\alpha}
 &=\sup_{z\in K_0}
   \frac{{\rm Var}_{\psi_{\alpha,z}^P}(T_\alpha)}
        {2h_\alpha m_\alpha^2}.                            \tag{41.33}
\end{aligned}
\]
In (41.32), both operators must be realized on one carrier; if
\(\widehat H_\alpha\) is defined only after projection, then
\(a\) is a reduced observable and the supremum is over \(K_0/G\).

\begin{theorem}[Finite semiclassical transfer--constraint handoff]
\label{thm:v14v41-handoff}
Assume:

\begin{enumerate}
\item \(v_\alpha(z)\ge v_{0,\alpha}>0\) on \(K_0\);
\item the faithful transfer conditions (41.2) hold;
\item
\[
 \epsilon_{H,\alpha}+\epsilon_{G,\alpha}
 +\delta_{{\rm tr},\alpha}\longrightarrow0;               \tag{41.34}
\]
\item for every observable in a fixed separating finite family,
\(\epsilon_{{\rm dyn},\alpha}(a)\to0\).
\end{enumerate}

Then, uniformly on \(K_0\),
\[
 -h_\alpha^{-1}
 \log\langle\psi_{\alpha,z}^P,
             T_\alpha\psi_{\alpha,z}^P\rangle
 \longrightarrow H_{\rm cl}(z).                           \tag{41.35}
\]
Uniformly on \(K\), the unprojected quantum-constraint symbols
converge to the V29 Gauss momentum map.  The initial Heisenberg
derivatives of the separating observables converge to their
classical Hamiltonian derivatives.  Therefore these rows are
sufficient for a finite-screen principal-symbol identification of
the transfer generator, Gauss constraints, and classical flow.
\end{theorem}

\begin{proof}
Apply Theorem~\ref{thm:v14v41-jensen} to
\(\phi=\psi_{\alpha,z}^P\).  Equations (41.30) and (41.33) give
\[
 \left|
 -h_\alpha^{-1}\log
 \langle\psi_{\alpha,z}^P,T_\alpha\psi_{\alpha,z}^P\rangle
 -H_{\rm cl}(z)
 \right|
 \le\delta_{{\rm tr},\alpha}+\epsilon_{H,\alpha},          \tag{41.36}
\]
which proves (41.35).  Equation (41.31) proves the constraint-symbol
statement.  The Heisenberg derivative at time zero is
\[
 \left.\frac d{dt}\right|_{t=0}
 \sigma_\alpha\!\left(
 e^{it\widehat H_\alpha/\hbar_\alpha}
 \operatorname{Op}_\alpha(a)
 e^{-it\widehat H_\alpha/\hbar_\alpha}\right)
 =
 \sigma_\alpha\!\left(
 \frac1{i\hbar_\alpha}
 [\operatorname{Op}_\alpha(a),\widehat H_\alpha]\right).
                                                               \tag{41.37}
\]
Equation (41.32) gives convergence to
\(\{a,H_{\rm cl}\}\).  A separating finite family determines the
finite-screen tangent vector in the selected coordinates.  The
visibility hypothesis makes every projected state in these formulas
defined and prevents the normalization denominator from vanishing.
\end{proof}

The theorem is a sufficiency statement, not an existence theorem.
In particular, it does not derive a continuum Hilbert space,
essential self-adjointness of an unbounded limiting Hamiltonian,
vacuum uniqueness, or a positive excitation gap.

\subsection{Constraint-algebra and preservation rows}

If the quantum constraints satisfy the exact normalization
\[
 [\widehat\Phi_\alpha(\xi),\widehat\Phi_\alpha(\eta)]
 =i\hbar_\alpha\widehat\Phi_\alpha([\xi,\eta]),             \tag{41.38}
\]
then a symbol-commutator estimate is still needed to identify the
classical Poisson bracket.  Namely, if
\[
 \left|
 \sigma_\alpha\!\left(
 \frac{[\widehat\Phi_\alpha(\xi),
        \widehat\Phi_\alpha(\eta)]}{i\hbar_\alpha}\right)(z)
 -\{\Phi_\xi,\Phi_\eta\}(z)
 \right|\le\epsilon_{{\rm br},\alpha}\|\xi\|\|\eta\|,       \tag{41.39}
\]
then (41.31), (41.38), and (41.39) imply
\[
 \left|
 \{\Phi_\xi,\Phi_\eta\}(z)-\Phi_{[\xi,\eta]}(z)
 \right|
 \le
 \epsilon_{{\rm br},\alpha}\|\xi\|\|\eta\|
 +\epsilon_{G,\alpha}\|[\xi,\eta]\|.                       \tag{41.40}
\]
The proof is substitution followed by the triangle inequality.

Likewise exact quantum gauge covariance
\[
 [\widehat\Phi_\alpha(\xi),\widehat H_\alpha]=0            \tag{41.41}
\]
implies classical constraint preservation only when the corresponding
commutator-to-Poisson error is controlled:
\[
 \left|
 \sigma_\alpha\!\left(
 \frac{[\widehat\Phi_\alpha(\xi),\widehat H_\alpha]}
      {i\hbar_\alpha}\right)(z)
 -\{\Phi_\xi,H_{\rm cl}\}(z)
 \right|\le\epsilon_{{\rm inv},\alpha}\|\xi\|.             \tag{41.42}
\]
Then
\[
 |\{\Phi_\xi,H_{\rm cl}\}(z)|
 \le\epsilon_{{\rm inv},\alpha}\|\xi\|.                   \tag{41.43}
\]
Thus exact operator commutators do not replace the principal-symbol
row that connects them to the classical bracket.

\subsection{The source acquisition card}

The finite symbol card exposed by V41 is
\[
 \mathfrak S_{\hbar,h}(K)=
 \bigl(
 [\psi_z],[\operatorname{Op}_{\hbar,h}],
 v_0,m,
 \epsilon_H,\epsilon_G,\epsilon_{\rm dyn},
 \epsilon_{\rm br},\epsilon_{\rm inv},\delta_{\rm tr}
 \bigr).                                                   \tag{41.44}
\]
Its carriers must be tied to the same \(K_4\) phase variables,
boundary convention, temporal transfer, and physical clock used in
V27--V39.  A generic family of coherent states on an unrelated
lattice phase space cannot populate this card.

The first missing literal source row is now precise:
\[
 \boxed{\text{a source-defined quantization/coherent-state map on the
 selected \(K_4\) phase chart, compatible with Haar reduction and the
 physical clock, with estimates for (41.30)--(41.33).}}     \tag{41.45}
\]
No such map, visibility floor, or transfer-variance ledger has been
located in the inspected source passages.

\begin{decision}[Direction after V41]
\label{dec:v14v41-next}
V42 will go upstream on the hardest new row rather than assume a
coherent-state package.  It will derive a gauge-orbit Gram formula
for \(v_{\hbar,h}(z)\), prove a quantitative visibility floor from a
finite overlap matrix, and determine exactly which stabilizer and
orbit-separation data are needed on the \(K_4\) screen.  If that floor
cannot survive cofinally, the programme must switch from
group-averaged kinematical states to a direct reduced-phase
quantization.
\end{decision}

\subsection{Acquisition ledger}

\[
\begin{array}{p{0.30\linewidth}|p{0.20\linewidth}|p{0.39\linewidth}}
\textbf{gate}&\textbf{status after V41}&\textbf{content}\\ \hline
operator versus scalar log&proved&
Jensen--variance bound\\
Haar visibility&necessary&
ratio theorem and counterexample\\
off-shell Gauss symbol&carrier fixed&
must be taken before projection\\
finite symbol handoff&conditional theorem&
errors (41.30)--(41.33)\\
constraint algebra/preservation&typed&
commutator-to-Poisson rows\\
literal coherent-state card&missing&
source acquisition (41.45)\\
orbit visibility floor&next&
finite overlap Gram\\
full SM--Einstein continuum&open&
cofinal construction not yet supplied
\end{array}                                                 \tag{41.46}
\]

\subsection{Parent-integrity record}

The installed parent was
\[
\begin{array}{c|c}
 \text{file}&
 \texttt{Renewal\_SM\_Einstein\_Continuum\_Research\_Notes\_V40.tex}\\
 \text{lines}&15988\\
 \text{bytes}&523073\\
 \text{SHA--256}&
 \texttt{A5440482F7D883D915C3580107894B50A6E58BDD82507E8459910A63C5D7A50E}.
\end{array}                                                 \tag{41.47}
\]

% =============================================================================
% END APPENDED FOURTEENTH-CYCLE V41 MATERIAL
% =============================================================================
% =============================================================================
% BEGIN APPENDED FOURTEENTH-CYCLE V42 MATERIAL
% =============================================================================

\section{Fourteenth-cycle V42: exact Haar-orbit Gram certificates and
the visibility atlas}
\label{sec:v14v42}

\subsection{The finite representation problem}

Fix one finite screen and abbreviate
\[
 {\cal H}={\cal H}_{\hbar,h},\qquad
 U=U_{\hbar,h},\qquad
 P=\int_GU(g)\,dg.                                         \tag{42.1}
\]
Let \(d=\dim_{\mathbb C}{\cal H}<\infty\).  V41 showed that
\[
 v(z)=\langle\psi_z,P\psi_z\rangle=\|P\psi_z\|^2           \tag{42.2}
\]
is the normalization denominator for a Haar-projected coherent
state.  This version reduces (42.2) to two finite certificates: an
exact orbit Gram matrix and a finite atlas of invariant
coefficients.

\subsection{Exact finite Haar cubature for one representation}

\begin{theorem}[Finite operator cubature for Haar projection]
\label{thm:v14v42-cubature}
There are points \(g_1,\ldots,g_J\in G\) and weights
\[
 w_j\ge0,\qquad \sum_{j=1}^Jw_j=1,\qquad
 J\le 2d^2+1,                                              \tag{42.3}
\]
such that
\[
 \boxed{P=\sum_{j=1}^Jw_jU(g_j).}                          \tag{42.4}
\]
The nodes and weights depend on the finite representation, but not
on \(z\).
\end{theorem}

\begin{proof}
Regard the complex matrix algebra
\({\rm End}_{\mathbb C}({\cal H})\) as a real vector space of
dimension \(2d^2\).  The compact set
\[
 {\cal U}=\{U(g):g\in G\}
\]
has Haar barycenter \(P\), so \(P\) belongs to the convex hull of
\({\cal U}\).  Carathéodory's theorem represents any point of a
convex hull in a real \(2d^2\)-dimensional space by a convex
combination of at most \(2d^2+1\) points.  This is (42.3)--(42.4).
\end{proof}

This theorem does not claim an efficient cubature.  It proves that
the visibility question for a fixed finite representation always
has an exact finite witness.

For a normalized \(\psi_z\), define orbit vectors and their Gram
matrix by
\[
 u_j(z)=U(g_j)\psi_z,\qquad
 G_z=(\langle u_i(z),u_j(z)\rangle)_{i,j=1}^J.             \tag{42.5}
\]

\begin{corollary}[Exact orbit Gram formula]
\label{cor:v14v42-orbit-gram}
With \(w=(w_1,\ldots,w_J)^{\mathsf T}\),
\[
 \boxed{v(z)=w^*G_zw.}                                    \tag{42.6}
\]
Equivalently,
\[
 v(z)=\int_G\langle\psi_z,U(g)\psi_z\rangle\,dg.           \tag{42.7}
\]
Although the integrand in (42.7) may be complex, both expressions
are real and nonnegative.
\end{corollary}

\begin{proof}
By (42.4),
\[
 P\psi_z=\sum_jw_ju_j(z).
\]
Since \(P=P^*=P^2\),
\[
 v(z)=\|P\psi_z\|^2
 =\sum_{i,j}\overline{w_i}w_j
   \langle u_i(z),u_j(z)\rangle
 =w^*G_zw.
\]
The first equality in (42.7) follows directly from (42.1).
\end{proof}

A small eigenvalue of the full Gram matrix is not the required
quantity.  The exact visibility tests the distinguished coefficient
vector \(w\).  In particular, \(G_z\) may be singular while
\(w^*G_zw>0\).

\subsection{The invariant-coordinate formula}

Let
\[
 {\cal H}^G=P{\cal H}
\]
have dimension \(r\), and choose an orthonormal basis
\(e_1,\ldots,e_r\).  Define
\[
 c(z)=
 \bigl(\langle e_1,\psi_z\rangle,\ldots,
       \langle e_r,\psi_z\rangle\bigr)^{\mathsf T}
 \in\mathbb C^r.                                          \tag{42.8}
\]

\begin{proposition}[Trivial-isotypic visibility]
\label{prop:v14v42-trivial}
The following identities and equivalences hold:
\[
 v(z)=\|c(z)\|_{\mathbb C^r}^2
 =w^*G_zw,                                                 \tag{42.9}
\]
and
\[
 v(z)>0
 \quad\Longleftrightarrow\quad
 P\psi_z\ne0
 \quad\Longleftrightarrow\quad
 \psi_z\ \hbox{has a nonzero trivial-isotypic component}. \tag{42.10}
\]
\end{proposition}

\begin{proof}
The orthogonal projection onto
\(\operatorname{span}\{e_1,\ldots,e_r\}\) is
\[
 P=\sum_{a=1}^r|e_a\rangle\langle e_a|.
\]
Substitution in (42.2) gives the first expression in (42.9);
Corollary~\ref{cor:v14v42-orbit-gram} gives the second.  Each
equivalence in (42.10) is then immediate.
\end{proof}

Thus visibility is a representation-theoretic coefficient, not a
consequence of classical orbit dimension alone.

\subsection{A finite visibility atlas}

Let \(K\) be a compact metric subset of the selected classical phase
chart.  Suppose \(z_1,\ldots,z_N\in K\), radii \(r_j>0\), and constants
\(L_j\ge0\) satisfy
\[
 K\subset\bigcup_{j=1}^N B(z_j,r_j),                       \tag{42.11}
\]
and
\[
 \|\psi_z-\psi_{z_j}\|
 \le L_j\,d_K(z,z_j)
 \qquad(z\in K\cap B(z_j,r_j)).                            \tag{42.12}
\]
Put
\[
 \eta_j=\sqrt{v(z_j)}=\|P\psi_{z_j}\|.                     \tag{42.13}
\]

\begin{theorem}[Finite atlas visibility floor]
\label{thm:v14v42-atlas}
If
\[
 \eta_j>L_jr_j\qquad(1\le j\le N),                         \tag{42.14}
\]
then
\[
 \boxed{
 \inf_{z\in K}v(z)
 \ge
 \min_{1\le j\le N}(\eta_j-L_jr_j)^2>0.}                  \tag{42.15}
\]
Each centre value \(\eta_j^2\) can be supplied either by the
invariant-coordinate formula or by the exact orbit Gram formula.
\end{theorem}

\begin{proof}
The projection \(P\) is a contraction.  If
\(z\in B(z_j,r_j)\), then
\[
 \|P\psi_z-P\psi_{z_j}\|
 \le\|\psi_z-\psi_{z_j}\|
 \le L_jr_j.                                              \tag{42.16}
\]
The reverse triangle inequality and (42.13) give
\[
 \sqrt{v(z)}=\|P\psi_z\|
 \ge\eta_j-L_jr_j.                                        \tag{42.17}
\]
Square and minimize over the finite cover.
\end{proof}

This is the visibility analogue of the independent local-certificate
method imported structurally at V40.  The invariant basis and the
local state derivative may change between patches; they need not be
interpolated into one global auxiliary choice.

\begin{corollary}[Certified approximate Gram centres]
\label{cor:v14v42-approx-gram}
Suppose a Hermitian approximation \(\widetilde G_j\) obeys
\[
 \|G_{z_j}-\widetilde G_j\|_{\rm op}\le\rho_j.             \tag{42.18}
\]
Then
\[
 v(z_j)\ge
 \underline v_j:=
 w^*\widetilde G_jw-\rho_j\|w\|^2.                        \tag{42.19}
\]
If \(\underline v_j>0\) and
\[
 \sqrt{\underline v_j}>L_jr_j,                            \tag{42.20}
\]
then (42.15) holds with
\(\eta_j\) replaced by \(\sqrt{\underline v_j}\).
\end{corollary}

\begin{proof}
The quadratic-form error is bounded by
\[
 |w^*(G_{z_j}-\widetilde G_j)w|
 \le\rho_j\|w\|^2.
\]
This proves (42.19).  Repeat the proof of
Theorem~\ref{thm:v14v42-atlas}.
\end{proof}

\subsection{Stabilizers and orbit separation are insufficient}

\begin{countertheorem}[Maximal projective stabilizer with zero visibility]
\label{cth:v14v42-stabilizer}
Let \({\cal H}=\mathbb C\) and let \(U(1)\) act by the nontrivial
character
\[
 U(e^{i\theta})\psi=e^{in\theta}\psi,\qquad n\in\mathbb Z\setminus\{0\}.
                                                               \tag{42.21}
\]
Every nonzero vector has a single-point projective orbit and hence
maximal projective stabilizer, but
\[
 P=\int_0^{2\pi}e^{in\theta}\frac{d\theta}{2\pi}=0,
 \qquad v(\psi)=0.                                        \tag{42.22}
\]
\end{countertheorem}

\begin{proof}
The phase in (42.21) does not change the projective ray, so the
projective orbit has no separation at all.  Orthogonality of
nontrivial characters gives (42.22).
\end{proof}

The missing datum is the phase carried along the orbit, equivalently
the trivial-isotypic coefficient (42.8).  A classical stabilizer
ledger such as V32's root holonomy centralizer is valuable for rank
and reduction, but it cannot by itself prove quantum Haar
visibility.

\begin{countertheorem}[Finite visibility need not be cofinal]
\label{cth:v14v42-cofinal}
On
\[
 {\cal H}_N=(\mathbb C^2)^{\otimes N},
 \qquad
 U_N(e^{i\theta})
 =\operatorname{diag}(1,e^{i\theta})^{\otimes N},
\]
let
\[
 \psi_{\epsilon,N}
 =\bigl(\epsilon|0\rangle+
        \sqrt{1-\epsilon^2}|1\rangle\bigr)^{\otimes N},
 \qquad 0<\epsilon<1.
\]
Then
\[
 v_N=\epsilon^{2N}>0\quad\hbox{for every finite \(N\)},\qquad
 v_N\longrightarrow0.                                    \tag{42.23}
\]
\end{countertheorem}

\begin{proof}
All tensor factors have nonnegative integer charge.  The only
charge-zero basis vector is \(|0\rangle^{\otimes N}\), whose
coefficient in \(\psi_{\epsilon,N}\) is \(\epsilon^N\).  Haar
projection keeps precisely this vector, so (42.23) follows.
\end{proof}

Therefore a positive visibility value at each regulator is not a
uniform cofinal floor.

\subsection{Uniform floors versus relative numerator control}

A uniform visibility floor is sufficient for stable quotient
symbols, but it is not logically necessary when numerator errors
are relative to the same small denominator.

\begin{proposition}[Relative projected-symbol certificate]
\label{prop:v14v42-relative}
Let \(v_\alpha(z)>0\) and let
\[
 N_\alpha(z)
 =\langle\psi_{\alpha,z},
   P_\alpha A_\alpha P_\alpha\psi_{\alpha,z}\rangle.
\]
If for a target \(a(z)\)
\[
 |N_\alpha(z)-a(z)v_\alpha(z)|
 \le r_\alpha(z)v_\alpha(z),                              \tag{42.24}
\]
then
\[
 |\sigma_\alpha^P(A_\alpha)(z)-a(z)|
 \le r_\alpha(z),                                         \tag{42.25}
\]
regardless of the size of \(v_\alpha(z)\).
\end{proposition}

\begin{proof}
Divide (42.24) by the positive number \(v_\alpha(z)\) and use
(41.18).
\end{proof}

Hence the cofinal acceptance alternatives are:

\begin{enumerate}
\item prove \(\inf_{\alpha,z}v_\alpha(z)>0\) and control absolute
numerator errors;
\item allow \(v_\alpha\to0\), but prove every numerator error is
\(o(v_\alpha)\); or
\item quantize the reduced phase space directly, so no Haar
normalization denominator occurs.
\end{enumerate}

The second route can be badly conditioned and requires
source-specific cancellation.  It must not be inferred from
pointwise positivity.

\subsection{The \(K_4\) acquisition ledger}

For the selected \(K_4\) screen, an exact visibility card is
\[
 \mathfrak V_{\hbar,h}(K)=
 \bigl(
 [U_{\hbar,h}], [P_{\hbar,h}],
 [\psi_z],
 [g_j,w_j],
 [G_z],
 [e_a],
 [z_j,r_j,L_j,\underline v_j]
 \bigr).                                                   \tag{42.26}
\]
There are two equivalent centre certificates,
\[
 \underline v_j
 \le \sum_a|\langle e_a,\psi_{z_j}\rangle|^2
 =w^*G_{z_j}w.                                             \tag{42.27}
\]
To populate either side one must know the finite quantum
representation and the coherent vectors.  The classical K4
incidence, Gauss rank, stabilizer, and root-anchor floors of
V26--V32 do not determine those data.

\begin{assessment}[What V42 resolves]
\label{ass:v14v42-resolves}
The visibility row is now a finite, exact certification problem.
Haar integration itself is not an infinite-dimensional obstacle:
Theorem~\ref{thm:v14v42-cubature} replaces it by at most
\(2d^2+1\) representation matrices.  The unresolved obstruction is
source acquisition of the representation, the state map, and a
cofinal invariant-component estimate.  Stabilizer information alone
has been ruled out as a substitute.
\end{assessment}

\begin{decision}[Direction after V42]
\label{dec:v14v42-next}
V43 will construct the alternative local compiler on a regular
reduced symplectic slice.  For a finite linearized bosonic slice it
will build Gaussian coherent states, compute exactly the covariant
symbols of affine and quadratic observables, and prove the
commutator-to-Poisson row.  This will determine how much of the V41
symbol card follows canonically once the V29 constraint reduction
and a positive complex structure are supplied, while keeping the
source-specific existence of that slice open.
\end{decision}

\subsection{Acquisition ledger}

\[
\begin{array}{p{0.30\linewidth}|p{0.20\linewidth}|p{0.39\linewidth}}
\textbf{gate}&\textbf{status after V42}&\textbf{content}\\ \hline
finite Haar cubature&proved&
at most \(2d^2+1\) nodes\\
orbit Gram visibility&proved&
\(v=w^*G_zw\)\\
finite visibility atlas&proved&
centre floor minus state variation\\
stabilizer-only inference&refuted&
nontrivial-character counterexample\\
cofinal pointwise inference&refuted&
tensor-power counterexample\\
relative quotient route&proved&
error relative to visibility\\
literal \(K_4\) visibility card&missing&
quantum representation and states\\
reduced Gaussian compiler&next&
regular symplectic slice
\end{array}                                                 \tag{42.28}
\]

\subsection{Parent-integrity record}

The installed parent was
\[
\begin{array}{c|c}
 \text{file}&
 \texttt{Renewal\_SM\_Einstein\_Continuum\_Research\_Notes\_V41.tex}\\
 \text{lines}&16522\\
 \text{bytes}&540895\\
 \text{SHA--256}&
 \texttt{7C330FF0C541304346946AA52D25DE64D69D5B59F491FDE750B42CDA2BD22D8F}.
\end{array}                                                 \tag{42.29}
\]

% =============================================================================
% END APPENDED FOURTEENTH-CYCLE V42 MATERIAL
% =============================================================================
% =============================================================================
% BEGIN APPENDED FOURTEENTH-CYCLE V43 MATERIAL
% =============================================================================

\section{Fourteenth-cycle V43: reduced symplectic Gaussian quantization
and the exact quadratic symbol sector}
\label{sec:v14v43}

\subsection{Regular finite reduction}

Let \((Z,\Omega)\) be the finite classical phase space at one screen,
with Hamiltonian action of a compact gauge group \(G\) and momentum
map \(\Phi\).  Fix \(x\in\Phi^{-1}(0)\).  Write
\[
 C_x=\ker D\Phi_x,\qquad
 {\cal O}_x=\{X_{\Phi_\xi}(x):\xi\in\mathfrak g\}.          \tag{43.1}
\]

\begin{proposition}[Linear reduced symplectic space]
\label{prop:v14v43-reduction}
Assume \(x\) lies in a regular orbit-type stratum and
\[
 C_x^{\Omega}={\cal O}_x.                                  \tag{43.2}
\]
Then
\[
 V_x^{\rm red}=C_x/{\cal O}_x                              \tag{43.3}
\]
has a well-defined nondegenerate symplectic form
\[
 \Omega_x^{\rm red}([u],[v])=\Omega_x(u,v).                \tag{43.4}
\]
\end{proposition}

\begin{proof}
First-classness gives
\({\cal O}_x\subset C_x\).  If \(u\) is replaced by
\(u+o\), \(o\in{\cal O}_x=C_x^\Omega\), then
\(\Omega_x(o,v)=0\) for every \(v\in C_x\); hence (43.4) is
well defined.  If \([u]\) is in the radical of (43.4), then
\(\Omega_x(u,C_x)=0\), so \(u\in C_x^\Omega={\cal O}_x\).
Thus \([u]=0\), proving nondegeneracy.
\end{proof}

Condition (43.2) is the regularity row.  V29 proves a finite Gauss
momentum map and its first-class algebra, but those facts alone do
not prove constant orbit type or (43.2) on a cofinal family.

\begin{lemma}[Invariant compatible complex structure]
\label{lem:v14v43-complex}
Let a compact group \(K\) act symplectically on a finite-dimensional
real symplectic vector space \((V,\omega)\).  There is a
\(K\)-invariant complex structure \(J\) such that
\[
 J^2=-I,\qquad
 \omega(Ju,Jv)=\omega(u,v),\qquad
 g(u,v):=\omega(u,Jv)
\]
is a positive inner product.
\end{lemma}

\begin{proof}
Average any positive inner product over \(K\), obtaining an invariant
inner product \(g_0\).  Define the invertible \(g_0\)-skew operator
\(A\) by
\[
 \omega(u,v)=g_0(Au,v).                                    \tag{43.5}
\]
The positive operator \(-A^2\) commutes with \(A\) and with \(K\).
Functional calculus therefore defines
\[
 J=A(-A^2)^{-1/2}.                                         \tag{43.6}
\]
Then \(J^2=-I\), \(J\) is \(K\)-invariant, and it preserves
\(\omega\).  Moreover
\[
 \omega(u,Jv)
 =g_0\!\left((-A^2)^{1/2}u,v\right),                       \tag{43.7}
\]
which is positive definite.
\end{proof}

Applied to (43.3), this lemma supplies a Gaussian polarization once
the regular reduced space is known.  It does not select one
canonically from the source.

\subsection{Scaled bosonic coherent states}

Use \(J\) to identify \(V_x^{\rm red}\) with a complex one-particle
space \({\cal K}\), normalized so that
\[
 \omega(u,v)=2\,{\rm Im}\langle u,v\rangle_{\cal K},        \tag{43.8}
\]
where the Hermitian product is linear in its second argument.  On
the bosonic Fock space
\[
 {\cal F}_s({\cal K})
 =\bigoplus_{n=0}^{\infty}{\cal K}^{\widehat\otimes n},    \tag{43.9}
\]
let \(a_\hbar(f)=\sqrt\hbar\,a(f)\).  Then
\[
 [a_\hbar(f),a_\hbar^\dagger(g)]
 =\hbar\langle f,g\rangle I.                               \tag{43.10}
\]
For \(z\in{\cal K}\), the normalized exponential vector
\[
 |z;\hbar\rangle
 =
 e^{-\|z\|^2/(2\hbar)}
 \exp\!\left(\hbar^{-1}a_\hbar^\dagger(z)\right)|0\rangle  \tag{43.11}
\]
satisfies
\[
 a_\hbar(f)|z;\hbar\rangle
 =\langle f,z\rangle|z;\hbar\rangle.                       \tag{43.12}
\]

For a Hermitian one-particle operator \(A\), define
\[
 q_A(z)=\langle z,Az\rangle,\qquad
 \widehat q_{\hbar,A}:=\hbar\,d\Gamma(A).                  \tag{43.13}
\]

\begin{theorem}[Exact normal-ordered quadratic symbol]
\label{thm:v14v43-quadratic-symbol}
For every Hermitian \(A\),
\[
 \boxed{
 \langle z;\hbar|\widehat q_{\hbar,A}|z;\hbar\rangle
 =q_A(z).}                                                 \tag{43.14}
\]
For Hermitian \(A,B\),
\[
 \boxed{
 \frac1{i\hbar}
 [\widehat q_{\hbar,A},\widehat q_{\hbar,B}]
 =
 \widehat q_{\hbar,-i[A,B]}.}                              \tag{43.15}
\]
With the symplectic convention (43.8),
\[
 \{q_A,q_B\}=q_{-i[A,B]}.                                  \tag{43.16}
\]
Thus the covariant symbol and commutator-to-Poisson rows are exact
on the number-conserving quadratic algebra.
\end{theorem}

\begin{proof}
Choose an orthonormal basis \((e_j)\).  On the finite-particle core,
\[
 \hbar d\Gamma(A)
 =\sum_{j,k}\langle e_j,Ae_k\rangle
   a_\hbar^\dagger(e_j)a_\hbar(e_k).                       \tag{43.17}
\]
Equation (43.12) and its adjoint give (43.14).
The standard second-quantization identity
\[
 [d\Gamma(A),d\Gamma(B)]=d\Gamma([A,B])
\]
gives (43.15) after inserting the factors of \(\hbar\).

For completeness, differentiating \(q_A\) shows that its Hamiltonian
vector field is \(X_{q_A}(z)=-iAz\), under the convention
\(\omega(X_f,\cdot)=Df\).  Hence
\[
 \{q_A,q_B\}
 =\omega(-iAz,-iBz)
 =2\,{\rm Im}\langle Az,Bz\rangle
 =\langle z,-i[A,B]z\rangle,
\]
which is (43.16).
\end{proof}

\subsection{Exact linearized Gauss quantization}

Suppose a residual compact gauge algebra acts unitarily on
\({\cal K}\), with anti-Hermitian infinitesimal representation
\(T(\xi)\).  Put
\[
 A_\xi=iT(\xi)=A_\xi^*,\qquad
 \Phi_\xi^{\rm lin}(z)=q_{A_\xi}(z),\qquad
 \widehat\Phi_{\hbar}(\xi)
 =\widehat q_{\hbar,A_\xi}.                                \tag{43.18}
\]

\begin{corollary}[Exact Gauss symbol and algebra]
\label{cor:v14v43-gauss}
For all \(\xi,\eta\),
\[
 \langle z;\hbar|\widehat\Phi_\hbar(\xi)|z;\hbar\rangle
 =\Phi_\xi^{\rm lin}(z),                                   \tag{43.19}
\]
and
\[
 [\widehat\Phi_\hbar(\xi),\widehat\Phi_\hbar(\eta)]
 =i\hbar\widehat\Phi_\hbar([\xi,\eta]),                    \tag{43.20}
\]
provided \(T([\xi,\eta])=[T(\xi),T(\eta)]\).
\end{corollary}

\begin{proof}
Equation (43.19) is (43.14).  Since
\[
 [A_\xi,A_\eta]
 =[iT(\xi),iT(\eta)]
 =-T([\xi,\eta])
 =iA_{[\xi,\eta]},
\]
equation (43.15) gives (43.20).
\end{proof}

There is no zero-point shift in (43.19), because
\(\hbar d\Gamma(A)\) is normal ordered.  This exact result applies to
the linearized V29 momentum map only after a symplectic identification
between its reduced variables and \({\cal K}\) has been supplied.
Nonlinear gauge actions and nonquadratic momentum maps require
separate symbol remainders.

\subsection{Positive quadratic Hamiltonians and their transfer}

Assume the reduced quadratic Hamiltonian is compatible with the
chosen polarization:
\[
 H_{\rm cl}^{(2)}(z)=q_K(z),\qquad K=K^*\succeq0.           \tag{43.21}
\]
Define
\[
 \widehat H_\hbar^{(2)}=\hbar d\Gamma(K),\qquad
 T_{\hbar,h}^{(2)}=e^{-h\widehat H_\hbar^{(2)}}.            \tag{43.22}
\]

\begin{theorem}[Exact coherent transfer formula]
\label{thm:v14v43-transfer}
For every \(z\in{\cal K}\),
\[
 \boxed{
 \langle z;\hbar|T_{\hbar,h}^{(2)}|z;\hbar\rangle
 =
 \exp\!\left[
 -\frac1\hbar
 \langle z,(I-e^{-h\hbar K})z\rangle
 \right].}                                                 \tag{43.23}
\]
Consequently the scalar transfer energy is
\[
 E_{\hbar,h}^{\rm tr}(z)
 =
 \frac1{h\hbar}
 \langle z,(I-e^{-h\hbar K})z\rangle,                      \tag{43.24}
\]
and
\[
 \boxed{
 0\le H_{\rm cl}^{(2)}(z)-E_{\hbar,h}^{\rm tr}(z)
 \le\frac{h\hbar}{2}\langle z,K^2z\rangle.}                \tag{43.25}
\]
\end{theorem}

\begin{proof}
Let \(\alpha=z/\sqrt\hbar\).  Normalized exponential vectors obey
\[
 \langle E(\alpha),\Gamma(C)E(\alpha)\rangle
 =\exp\bigl(\langle\alpha,C\alpha\rangle-\|\alpha\|^2\bigr)
                                                               \tag{43.26}
\]
for a one-particle contraction \(C\).  Since
\[
 e^{-h\hbar d\Gamma(K)}=\Gamma(e^{-h\hbar K}),
\]
equation (43.23) follows from (43.26).

For \(s=h\hbar\) and \(x\ge0\),
\[
 0\le x-\frac{1-e^{-sx}}s\le\frac{s x^2}{2}.              \tag{43.27}
\]
The first inequality follows from \(e^{-sx}\ge1-sx\); the second
follows from
\(e^{-sx}\le1-sx+s^2x^2/2\).
Apply the spectral theorem for \(K\) to (43.27) and take the
quadratic form in \(z\).  This proves (43.25).
\end{proof}

The energy variance is also exact:
\[
 {\rm Var}_{z;\hbar}(\widehat H_\hbar^{(2)})
 =\hbar\langle z,K^2z\rangle.                              \tag{43.28}
\]
It follows either by normal ordering the square of
\(\hbar d\Gamma(K)\) or by differentiating (43.23) twice with respect
to \(h\) at zero.  Thus the V41 variance debit is naturally
\(O(h\hbar)\) on fixed compact subsets of \({\cal K}\).

\subsection{The unbounded-logarithm branch}

The Fock transfer in (43.22) is injective but, when \(K\) has a
positive mode and occupation is unbounded, it has no strictly
positive operator lower bound.  This does not prevent an unbounded
Hamiltonian reconstruction.

\begin{theorem}[Injective transfer gives an unbounded logarithm]
\label{thm:v14v43-unbounded-log}
Let \(T\) be a positive self-adjoint contraction on a Hilbert space
and assume \(\ker T=\{0\}\).  Then
\[
 H=-h^{-1}\log T                                           \tag{43.29}
\]
is a positive self-adjoint operator with dense domain
\[
 {\cal D}(H)=
 \left\{\psi:
 \int_{(0,1]}|\log\lambda|^2\,
 d\langle\psi,E_T(\lambda)\psi\rangle<\infty
 \right\},                                                 \tag{43.30}
\]
and \(T=e^{-hH}\).  A uniform floor \(T\succeq mI\) is needed for a
bounded logarithm, not for the existence of the self-adjoint
unbounded generator.
\end{theorem}

\begin{proof}
Because \(\ker T=E_T(\{0\}){\cal H}=\{0\}\), the spectral function
\(-h^{-1}\log\lambda\) is finite \(E_T\)-almost everywhere.  The
unbounded Borel functional calculus gives the self-adjoint operator
(43.29) on (43.30).  Truncating the spectral interval to
\([1/n,1]\) shows that the union of bounded spectral subspaces lies
in (43.30) and is dense.  Functional calculus then gives
\(e^{-hH}=T\).
\end{proof}

This theorem refines the V39 acquisition tree.  At a finite
occupation cutoff, a positive floor yields a bounded logarithm.  On
the full Fock carrier, injectivity and domain control replace that
floor.

\subsection{A hard-cutoff firewall}

Let \(a_\hbar,a_\hbar^\dagger\) be one-mode operators and let \(P_N\)
project onto occupations \(0,\ldots,N\).  Define
\[
 b_{\hbar,N}=P_Na_\hbar P_N.                               \tag{43.31}
\]

\begin{countertheorem}[Occupation truncation breaks the CCR at the cap]
\label{cth:v14v43-cutoff}
One has
\[
 [b_{\hbar,N},b_{\hbar,N}^\dagger]
 =\hbar\bigl(I_{N+1}-(N+1)|N\rangle\langle N|\bigr).       \tag{43.32}
\]
Thus a finite occupation cutoff does not preserve the exact
quadratic commutator compiler on states reaching the top sector.
\end{countertheorem}

\begin{proof}
For \(0\le n<N\),
\[
 b_{\hbar,N}b_{\hbar,N}^\dagger|n\rangle
 =\hbar(n+1)|n\rangle,\qquad
 b_{\hbar,N}^\dagger b_{\hbar,N}|n\rangle
 =\hbar n|n\rangle.
\]
For \(n=N\), the first term vanishes and the second equals
\(\hbar N|N\rangle\).  These eigenvalues are exactly those of
(43.32).
\end{proof}

A cutoff can still be used if the coherent-state probability of the
cap is bounded and entered as a defect.  It cannot be declared an
exact finite-dimensional representation of the CCR.

\subsection{What is filled and what remains}

On a regular reduced linear slice with a compatible complex
structure and a number-conserving positive quadratic Hamiltonian,
V43 fills the following rows without asymptotic error:
\[
 \epsilon_H=0,\qquad
 \epsilon_G=0,\qquad
 \epsilon_{\rm br}=0                                      \tag{43.33}
\]
for the corresponding normal-ordered quadratic observables.  The
scalar-transfer error has the explicit bound (43.25), and no Haar
visibility denominator appears after direct reduction.

The remaining source-dependent rows are:

\begin{enumerate}
\item a regular \(K_4\) reduced symplectic slice satisfying (43.2);
\item identification of the V27 Hamiltonian Hessian and V29 Gauss
linearization on that slice;
\item a polarization for which the quadratic Hamiltonian has the
number-conserving form (43.21), or a controlled Bogoliubov
replacement;
\item estimates for cubic and higher classical terms, fermionic
sectors, gravity dependence, and cofinal changes of screen; and
\item either the unbounded-domain branch of
Theorem~\ref{thm:v14v43-unbounded-log} or a cutoff-defect ledger
compatible with (43.32).
\end{enumerate}

\begin{decision}[Direction after V43]
\label{dec:v14v43-next}
V44 will remove the number-conserving restriction at quadratic
order.  It will treat a general positive bosonic quadratic form by a
finite Williamson/Bogoliubov transform, compute its vacuum and
coherent-state symbol correction, and state the exact stability
condition needed for the transformation to remain uniform along the
cofinal family.  This is the next upstream row before nonlinear
remainders can be estimated.
\end{decision}

\subsection{Acquisition ledger}

\[
\begin{array}{p{0.30\linewidth}|p{0.20\linewidth}|p{0.39\linewidth}}
\textbf{gate}&\textbf{status after V43}&\textbf{content}\\ \hline
regular linear reduction&conditional&
symplectic quotient theorem\\
compatible polarization&exists finitely&
compact-invariant complex structure\\
quadratic covariant symbol&exact&
normal-ordered Fock quantization\\
Gauss commutator algebra&exact&
linear unitary action\\
coherent transfer energy&proved&
\(O(h\hbar)\) error\\
unbounded transfer log&proved&
injectivity and spectral domain\\
hard occupation cutoff&defect exposed&
top-sector CCR term\\
general positive quadratic form&next&
Williamson/Bogoliubov compiler\\
full nonlinear continuum&open&
source and cofinal estimates missing
\end{array}                                                 \tag{43.34}
\]

\subsection{Parent-integrity record}

The installed parent was
\[
\begin{array}{c|c}
 \text{file}&
 \texttt{Renewal\_SM\_Einstein\_Continuum\_Research\_Notes\_V42.tex}\\
 \text{lines}&16954\\
 \text{bytes}&553849\\
 \text{SHA--256}&
 \texttt{56E382FE010ED15972F3BF78522DA228FDDF845CA2EE01B1EE10DD2506C33AFE}.
\end{array}                                                 \tag{43.35}
\]

% =============================================================================
% END APPENDED FOURTEENTH-CYCLE V43 MATERIAL
% =============================================================================
% =============================================================================
% BEGIN APPENDED FOURTEENTH-CYCLE V44 MATERIAL
% =============================================================================

\section{Fourteenth-cycle V44: Williamson--Bogoliubov reduction,
squeezing control, and the Einstein vacuum-energy row}
\label{sec:v14v44}

\subsection{General positive quadratic data}

Let \(V=\mathbb R^{2n}\) carry the standard symplectic matrix
\[
 \Omega=
 \begin{pmatrix}0&I_n\\-I_n&0\end{pmatrix},\qquad
 \omega(u,v)=u^{\mathsf T}\Omega v.                        \tag{44.1}
\]
After the constraint reduction of V43, consider the general
quadratic Hamiltonian
\[
 H_G(z)=\frac12z^{\mathsf T}Gz,\qquad
 G=G^{\mathsf T}\succ0.                                   \tag{44.2}
\]
No number-conserving polarization is assumed.

\begin{theorem}[Finite Williamson reduction with quantitative bounds]
\label{thm:v14v44-williamson}
There are a symplectic matrix \(S\) and positive numbers
\(\nu_1,\ldots,\nu_n\) such that
\[
 S^{\mathsf T}GS
 =D_\nu:=
 \operatorname{diag}(\nu_1,\ldots,\nu_n,
                     \nu_1,\ldots,\nu_n),                 \tag{44.3}
\]
\[
 S^{\mathsf T}\Omega S=\Omega.                             \tag{44.4}
\]
If
\[
 aI\preceq G\preceq bI,\qquad 0<a\le b,                   \tag{44.5}
\]
then
\[
 a\le\nu_j\le b,\qquad
 \|S\|\le\sqrt{b/a},\qquad
 \|S^{-1}\|\le\sqrt{b/a}.                                 \tag{44.6}
\]
\end{theorem}

\begin{proof}
Set
\[
 M=G^{1/2}\Omega G^{1/2}.
\]
This is a real invertible skew-symmetric matrix.  Its orthogonal
normal form gives an orthogonal \(O\), with orientations chosen in
each two-plane, such that
\[
 O^{\mathsf T}MO=\Omega D_\nu.                             \tag{44.7}
\]
Define
\[
 S=G^{-1/2}OD_\nu^{1/2}.                                   \tag{44.8}
\]
Then \(S^{\mathsf T}GS=D_\nu\).  Moreover,
\[
 G^{-1/2}\Omega G^{-1/2}=-M^{-1}.
\]
Using (44.7), the right side becomes
\(O D_\nu^{-1}\Omega O^{\mathsf T}\); substitution in
\(S^{\mathsf T}\Omega S\), and the fact that \(D_\nu\)
commutes with \(\Omega\), gives (44.4).

The \(\nu_j\) are the singular values of \(M\), each repeated twice.
Under (44.5),
\[
 \sigma_{\min}(M)\ge a,\qquad \|M\|\le b,                  \tag{44.9}
\]
because \(\Omega\) is orthogonal.  This proves the first part of
(44.6).  Formula (44.8) then gives
\[
 \|S\|\le\|G^{-1/2}\|\|D_\nu^{1/2}\|
 \le\sqrt{b/a}.
\]
The identity
\(S^{-1}=D_\nu^{-1/2}O^{\mathsf T}G^{1/2}\)
gives the same bound for \(S^{-1}\).
\end{proof}

The numbers \(\nu_j\) are the symplectic frequencies.  Uniform
ellipticity of the reduced Hessian controls both the frequencies and
the squeezing condition number.

\subsection{Squeezed quantization}

Put \(z=Sy\), where
\(y=(q_1,\ldots,q_n,p_1,\ldots,p_n)\).  Then
\[
 H_G(Sy)
 =\frac12\sum_{j=1}^n\nu_j(q_j^2+p_j^2).                  \tag{44.10}
\]
Let \(\mu_\hbar(S)\) be either metaplectic implementer of \(S\).
With
\[
 a_{\hbar,j}
 =\frac{\widehat q_j+i\widehat p_j}{\sqrt2},\qquad
 [a_{\hbar,j},a_{\hbar,k}^\dagger]
 =\hbar\delta_{jk},                                       \tag{44.11}
\]
the Weyl quantization of (44.10) is
\[
 \operatorname{Op}_\hbar^{W}(H_G)
 \simeq
 \sum_{j=1}^n\nu_j
 \left(a_{\hbar,j}^\dagger a_{\hbar,j}+\frac\hbar2\right).
                                                               \tag{44.12}
\]
Here \(\simeq\) denotes metaplectic conjugacy by
\(\mu_\hbar(S)\).

For \(x\in V\), set \(y=S^{-1}x\) and define the squeezed coherent
state
\[
 |x;\hbar,S\rangle
 =\mu_\hbar(S)|y;\hbar\rangle.                             \tag{44.13}
\]

\begin{theorem}[Exact squeezed quadratic symbols]
\label{thm:v14v44-squeezed-symbol}
Let
\[
 E_0(G,\hbar)=\frac\hbar2\sum_{j=1}^n\nu_j,                \tag{44.14}
\]
and define the normal-ordered quadratic Hamiltonian
\[
 \widehat H_{G,\hbar}^{N}
 =\operatorname{Op}_\hbar^{W}(H_G)-E_0(G,\hbar)I.          \tag{44.15}
\]
Then
\[
 \boxed{
 \langle x;\hbar,S|
 \widehat H_{G,\hbar}^{N}
 |x;\hbar,S\rangle=H_G(x),}                               \tag{44.16}
\]
whereas
\[
 \langle x;\hbar,S|
 \operatorname{Op}_\hbar^{W}(H_G)
 |x;\hbar,S\rangle
 =H_G(x)+E_0(G,\hbar).                                    \tag{44.17}
\]
\end{theorem}

\begin{proof}
Metaplectic covariance of Weyl quantization reduces the calculation
to (44.10).  In the ordinary coherent state centred at \(y\),
\[
 \langle a_{\hbar,j}^\dagger a_{\hbar,j}\rangle
 =\frac{q_j^2+p_j^2}{2}.                                  \tag{44.18}
\]
Substitution in (44.12) gives (44.17); subtracting (44.14) gives
(44.16).
\end{proof}

Thus the Bogoliubov transformation changes the coherent-state
covariance but introduces no phase-space-dependent symbol error at
quadratic order.  The only scalar correction is the vacuum energy
(44.14).

\subsection{Exact transfer formula after squeezing}

Let
\[
 K_\nu=\operatorname{diag}(\nu_1,\ldots,\nu_n)
\]
on the one-particle complex space and let
\[
 \widehat H_{G,\hbar}^{N}
 \simeq\hbar d\Gamma(K_\nu),\qquad
 T_{G,\hbar,h}^{N}=e^{-h\widehat H_{G,\hbar}^{N}}.          \tag{44.19}
\]
Writing
\[
 \alpha(y)=2^{-1/2}(q+ip)\in\mathbb C^n,                   \tag{44.20}
\]
Theorem~\ref{thm:v14v43-transfer} gives
\[
 \langle x;\hbar,S|T_{G,\hbar,h}^{N}|x;\hbar,S\rangle
 =
 \exp\!\left[
 -\frac1\hbar
 \langle\alpha(y),
 (I-e^{-h\hbar K_\nu})\alpha(y)\rangle
 \right].                                                  \tag{44.21}
\]
Consequently
\[
 0\le
 H_G(x)+h^{-1}\log
 \langle x;\hbar,S|T_{G,\hbar,h}^{N}|x;\hbar,S\rangle
 \le
 \frac{h\hbar}{2}
 \langle\alpha(y),K_\nu^2\alpha(y)\rangle.                 \tag{44.22}
\]
For the unrenormalized Weyl transfer,
\[
 T_{G,\hbar,h}^{W}
 =e^{-h\operatorname{Op}_\hbar^W(H_G)}
 =e^{-hE_0(G,\hbar)}T_{G,\hbar,h}^{N},                    \tag{44.23}
\]
so its scalar transfer energy is larger by exactly \(E_0(G,\hbar)\).

\subsection{Vacuum energy is an Einstein source row}

For a nongravitating finite system, subtracting (44.14) changes only
the energy origin.  In the coupled Einstein programme, a scalar
vacuum term contributes to the metric variation through the volume
density and therefore changes the cosmological source.  Hence
(44.15) is not an innocuous convention unless the source action
contains the corresponding counterterm.

\begin{proposition}[Cofinal alternatives for the quadratic vacuum term]
\label{prop:v14v44-vacuum}
For a cofinal family indexed by \(\alpha\), write
\[
 E_{0,\alpha}
 =\frac{\hbar_\alpha}{2}
  \sum_{j=1}^{n_\alpha}\nu_{j,\alpha}.                     \tag{44.24}
\]
One of the following source treatments must be recorded, with the chosen counterterm and volume normalization fixed:

\begin{enumerate}
\item \(E_{0,\alpha}\to0\), so the quadratic vacuum source vanishes;
\item \(E_{0,\alpha}\to E_0^{\rm ren}\) and the limit is retained as
a finite cosmological contribution;
\item a declared counterterm cancels a specified part of
\(E_{0,\alpha}\), with the residual shown to converge; or
\item the residual has no finite limit, in which case this
quadratic cofinal branch does not yet define an Einstein source.
\end{enumerate}
\end{proposition}

\begin{proof}
Equation (44.24) is the exact difference between the Weyl and
normal-ordered Hamiltonians.  Coupling a constant energy density to
the spacetime volume produces a term proportional to the metric in
the stress tensor.  Therefore the coupled limit requires a finite
declared limit of the retained scalar coefficient.  The four cases
exhaust vanishing, finite retention, finite renormalized retention,
and failure of convergence.
\end{proof}

In particular, fixed finite dimension gives
\(E_0=O(\hbar)\) when the frequencies remain bounded, but a growing
number of screen modes requires the separate trace condition
\[
 \hbar_\alpha\sum_{j=1}^{n_\alpha}\nu_{j,\alpha}
 \longrightarrow0                                        \tag{44.25}
\]
before the vacuum term may be discarded.

\subsection{Uniform squeezing and frequency stability}

Let \(G\) and \(\widetilde G\) obey the same ellipticity window
\[
 aI\preceq G,\widetilde G\preceq bI.                       \tag{44.26}
\]
Let \(\nu_j(G)\) and \(\nu_j(\widetilde G)\) be ordered symplectic
frequencies.

\begin{proposition}[Lipschitz control of symplectic frequencies]
\label{prop:v14v44-frequency}
After ordering,
\[
 \max_j|\nu_j(G)-\nu_j(\widetilde G)|
 \le\sqrt{\frac ba}\,\|G-\widetilde G\|.                   \tag{44.27}
\]
\end{proposition}

\begin{proof}
For positive matrices bounded below by \(aI\),
\[
 \|G^{1/2}-\widetilde G^{1/2}\|
 \le\frac{\|G-\widetilde G\|}{2\sqrt a}.                  \tag{44.28}
\]
This follows from the standard resolvent integral for the square
root.  Set
\(M=G^{1/2}\Omega G^{1/2}\) and define
\(\widetilde M\) similarly.  Since both square roots have norm at
most \(\sqrt b\),
\[
 \|M-\widetilde M\|
 \le\sqrt{\frac ba}\,\|G-\widetilde G\|.                   \tag{44.29}
\]
The symplectic frequencies are the singular values of \(M\), each
repeated twice.  The singular-value perturbation inequality gives
(44.27).
\end{proof}

Combining Theorem~\ref{thm:v14v44-williamson} and
Proposition~\ref{prop:v14v44-frequency}, a uniform window
\(0<a\le G_\alpha\le b<\infty\) supplies:

\begin{enumerate}
\item frequencies bounded away from zero and infinity;
\item uniformly conditioned squeezed covariances;
\item stable frequency sums whenever their multiplicity-weighted
trace is controlled; and
\item uniform constants in the quadratic transfer bound on compact
sets of the Williamson coordinates.
\end{enumerate}

The matrices \(S_\alpha\) themselves need not be continuous through
frequency degeneracies; the frequency multiset and the associated
spectral subspaces are the invariant data.  A cofinal proof should
use those invariant objects rather than a chosen ordered eigenbasis.

\subsection{Gauge-equivariant polarization}

Suppose a compact residual gauge group acts by symplectic matrices
\(R(g)\) and preserves \(G\):
\[
 R(g)^{\mathsf T}GR(g)=G.                                 \tag{44.30}
\]
Define \(A_G\) intrinsically by
\[
 \omega(u,A_Gv)=u^{\mathsf T}Gv.                           \tag{44.31}
\]
Then \(A_G\) commutes with every \(R(g)\).  Since \(G\succ0\), its
polar complex structure
\[
 J_G=A_G(-A_G^2)^{-1/2}                                   \tag{44.32}
\]
is compatible and also commutes with \(R(g)\), where the positive square root is taken in the \(G\)-inner product.  Thus the
Williamson polarization can be chosen gauge equivariantly; the
one-particle frequency operator then commutes with the residual
gauge representation.

\begin{proof}
Use (44.30), symplectic invariance, and (44.31):
\[
 \omega(Ru,A_GRv)
 =G(Ru,Rv)=G(u,v)=\omega(Ru,RA_Gv).
\]
Nondegeneracy gives \(A_GR=RA_G\).  Functional calculus gives the
same commutation for (44.32).  Positivity and compatibility follow
from the polar construction as in
Lemma~\ref{lem:v14v43-complex}.
\end{proof}

This removes a possible quadratic gauge anomaly at the polarization
step, provided the reduced Hessian is strictly positive and gauge
invariant.

\subsection{Zero modes are not squeezed oscillators}

\begin{countertheorem}[Semidefinite Hessian has no oscillator vacuum]
\label{cth:v14v44-zero-mode}
For one canonical pair,
\[
 H(q,p)=\frac12p^2                                         \tag{44.33}
\]
has a zero symplectic frequency.  Its Schrödinger quantization is
the free Hamiltonian \(-\hbar^2\partial_q^2/2\), whose spectral
bottom is zero but has no normalizable ground vector on
\(L^2(\mathbb R)\).  No finite invertible Williamson transform turns
(44.33) into \(\nu(q^2+p^2)/2\) with \(\nu>0\).
\end{countertheorem}

\begin{proof}
The Hamiltonian matrix of (44.33) is nilpotent and has only the
eigenvalue zero, whereas a positive oscillator has Hamiltonian
eigenvalues \(\pm i\nu\).  Similarity under a symplectic change of
coordinates preserves the spectrum.  The spectral statement is the
Fourier representation of the free Laplacian: multiplication by
\(\hbar^2k^2/2\) has no nonzero \(L^2\) eigenfunction at zero.
\end{proof}

All gauge, diffeomorphism, and relation zero modes must therefore be
removed or treated as constraints before applying the positive
Williamson compiler.  This is why the V32 relation projector and the
V43 regular reduction precede V44.

\subsection{Acquisition card and next direction}

The general quadratic source card is
\[
 \mathfrak W_\alpha=
 \bigl(
 [V_\alpha^{\rm red},\omega_\alpha],
 [G_\alpha],a_\alpha,b_\alpha,
 [\nu_{j,\alpha}],
 [S_\alpha],
 E_{0,\alpha},
 [\hbox{counterterm}_\alpha],
 [R_\alpha(g)]
 \bigr).                                                   \tag{44.34}
\]
No literal \(K_4\) evaluation of this card has yet been extracted
from the attached action.  Theorems V43--V44 show what would follow
once it is.

\begin{decision}[Direction after V44]
\label{dec:v14v44-next}
V45 is the compulsory companion audit.  It will check the newest
Yang--Mills and Navier--Stokes notes for a local positive-Hessian
atlas, a variable-clock spectral estimate, an OS transfer
construction, or a regulator trace/vacuum-energy control that can
populate (44.34).  V46 will then return to the first literal source
row: extraction of the coupled \(K_4\) reduced quadratic Hessian and
its metric-dependent ellipticity window.
\end{decision}

\subsection{Acquisition ledger}

\[
\begin{array}{p{0.30\linewidth}|p{0.20\linewidth}|p{0.39\linewidth}}
\textbf{gate}&\textbf{status after V44}&\textbf{content}\\ \hline
general positive quadratic form&proved&
Williamson reduction\\
squeezed covariant symbol&exact&
normal-ordered Hamiltonian\\
quadratic transfer error&proved&
\(O(h\hbar)\) after squeezing\\
vacuum energy&exposed&
Einstein source row (44.24)\\
uniform squeezing&conditional&
two-sided Hessian ellipticity\\
frequency perturbation&proved&
bound (44.27)\\
gauge-equivariant polarization&proved&
invariant positive Hessian\\
zero-mode substitution&refuted&
free-particle counterexample\\
literal \(K_4\) Hessian card&missing&
source extraction deferred to V46
\end{array}                                                 \tag{44.35}
\]

\subsection{Parent-integrity record}

The installed parent was
\[
\begin{array}{c|c}
 \text{file}&
 \texttt{Renewal\_SM\_Einstein\_Continuum\_Research\_Notes\_V43.tex}\\
 \text{lines}&17408\\
 \text{bytes}&568283\\
 \text{SHA--256}&
 \texttt{3249512A37E4217FF88101923351334EB61E9C66F3CD75F6C6DF7860CFDE1A25}.
\end{array}                                                 \tag{44.36}
\]

% =============================================================================
% END APPENDED FOURTEENTH-CYCLE V44 MATERIAL
% =============================================================================

% =============================================================================
% BEGIN APPENDED FOURTEENTH-CYCLE V45 MATERIAL
% =============================================================================

\section{Fourteenth-cycle V45: compulsory companion audit and the
relative-Hessian extraction rule}
\label{sec:v14v45}

\subsection{Live companion provenance}

The required forty-fifth-version audit was performed read-only on
\[
\begin{array}{p{0.51\linewidth}|r|r}
\textbf{file}&\textbf{lines}&\textbf{bytes}\\ \hline
\texttt{Renewal\_Yang\_Mills\_Mass\_Gap\_Research\_Notes\_V16.tex}
 &3504&128707\\
\texttt{Renewal\_NS\_Stability\_Research\_Notes\_V26.tex}
 &5319&278283.
\end{array}                                                  \tag{45.1}
\]
Their SHA--256 digests are
\[
\begin{aligned}
 &\texttt{18A4F49E25A5E8D74351BE92EDBA756B019F79BB1613E921F1C22166627FDFAF},
                                                               \tag{45.2}\\
 &\texttt{82C887F8C2C69E4F28FAD7C3DC8DE5B9099CB5A4BF431EBA1FFF3E91935978AD}.
                                                               \tag{45.3}
\end{aligned}
\]
The live Yang--Mills V15 and V16 files have the same byte count and
the same SHA--256 digest (45.2); V16 ends with the V15 section and
contains no V16 mathematical appendage.  It is therefore treated as
the highest-numbered live provenance file, but not as an additional
mathematical iteration.  No companion file was modified or removed.

\subsection{Yang--Mills V13--V15}

The new Yang--Mills block constructs the first exact local
three-dimensional patch for its stronger four-thirds pencil.  It
then replaces an ambient perturbation estimate by a congruence
metric adapted to the centre block, and finally transforms each
Fourier coefficient before taking its remainder norm.  In that
specific calculation the certified radius grows from \(2^{-18}\) to
\(2^{-10}\) and then to \(2^{-7}\).

The numerical matrices, radii, and Loewner floors belong to the
Yang--Mills Wilson pencil and are not imported.  The reusable rule is
the order of certification.  For a positive centre Hessian \(G_0\),
one should bound
\[
 \left\|R_0^*
 \bigl(G(q)-G_0\bigr)R_0\right\|                           \tag{45.4}
\]
with a rational or otherwise certified approximate whitening
\(R_0\), and transform the source coefficients before summing their
norms.  An ambient Frobenius stock can lose the anisotropy that
controls the true relative floor.

This rule is consistent with the whitening and Schur-complement
firewalls already proved earlier in the SM notes.  It sharpens the
planned literal extraction rather than adding a duplicate theorem.

\subsection{Navier--Stokes V26}

The Navier--Stokes file remains unchanged.  Its Oseen-kernel
rank-one norm refinement supplies no reduced Standard-Model
Hessian, symplectic frequency trace, temporal transfer, OS quotient,
or Einstein vacuum counterterm.

\subsection{No transfer or vacuum-energy substitution}

Neither companion provides any of
\[
 [G_{K_4}^{\rm red}],\quad
 [\nu_{j,K_4}],\quad
 \hbar\sum_j\nu_{j,K_4},\quad
 [R_h],[L_h],[T_{{\rm phys},h}],\quad
 \epsilon_{{\rm clock},h}.                                \tag{45.5}
\]
The local Yang--Mills block floor proves positivity of an auxiliary
Laurent certificate.  It is not an eigenvalue floor for the coupled
SM--Einstein physical Hessian, and it does not control the
Williamson trace in V44.

\begin{proposition}[Permitted structural import]
\label{prop:v14v45-import}
Suppose a source Hessian family has an exact finite coefficient
expansion
\[
 G(q)=G_0+\sum_{\alpha\ne0}c_\alpha(q)G_\alpha,            \tag{45.6}
\]
and a certified invertible \(R_0\) obeys
\[
 R_0^*G_0R_0\succeq\beta I,\qquad
 R_0^*R_0\preceq\gamma I.                                 \tag{45.7}
\]
If
\[
 \sum_{\alpha\ne0}|c_\alpha(q)|
 \,\|R_0^*G_\alpha R_0\|\le\rho<\beta,                    \tag{45.8}
\]
then
\[
 G(q)\succeq\frac{\beta-\rho}{\gamma}I.                   \tag{45.9}
\]
Only the proof pattern (45.6)--(45.9), not any companion coefficient,
is transferable.
\end{proposition}

\begin{proof}
The triangle inequality in the whitened metric gives
\[
 R_0^*G(q)R_0
 \succeq(\beta-\rho)I.
\]
For \(x=R_0y\),
\[
 x^*G(q)x\ge(\beta-\rho)\|y\|^2,\qquad
 \|x\|^2\le\gamma\|y\|^2,
\]
which proves (45.9).
\end{proof}

\subsection{Direction selected}

\begin{decision}[Direction after V45]
\label{dec:v14v45-next}
V46 will return to the attached Einstein--Standard-Model source and
the Grand Tensor monograph.  It will locate the literal finite
\(K_4\) action, its phase variables, its multiplier equations, and
the declared physical clock.  It will then determine whether a
reduced quadratic Hessian can actually be assembled from the stated
coefficients.  If the source gives only schematic operators, V46
will record the exact missing coefficient rows rather than assigning
them values.

Any local ellipticity test will use the centre-relative,
coefficientwise order (45.6)--(45.9).  V46 will not infer a
Williamson floor from the unrelated Yang--Mills patch.
\end{decision}

\subsection{Acquisition ledger}

\[
\begin{array}{p{0.30\linewidth}|p{0.20\linewidth}|p{0.39\linewidth}}
\textbf{gate}&\textbf{status after V45}&\textbf{content}\\ \hline
companion checkpoint&acquired&
YM V16 and NS V26 with live hashes\\
YM V15/V16 relation&duplicate verified&
same byte count and digest\\
relative coefficient method&imported structurally&
source Hessian rule (45.6)--(45.9)\\
YM numerical floors&not imported&
different operator family\\
literal reduced \(K_4\) Hessian&next&
attached-source extraction\\
Williamson trace/counterterm&unpopulated&
no companion data\\
OS transfer and clock rows&unpopulated&
no companion data\\
full continuum&open&
cofinal source estimates absent
\end{array}                                                 \tag{45.10}
\]

The next compulsory companion audit is V50.

\subsection{Parent-integrity record}

The installed parent was
\[
\begin{array}{c|c}
 \text{file}&
 \texttt{Renewal\_SM\_Einstein\_Continuum\_Research\_Notes\_V44.tex}\\
 \text{lines}&17866\\
 \text{bytes}&582988\\
 \text{SHA--256}&
 \texttt{26C9F33000EE52262432DC1959CFAFAA95F383A0845E8F47AE682C6BB15C7C69}.
\end{array}                                                 \tag{45.11}
\]

% =============================================================================
% END APPENDED FOURTEENTH-CYCLE V45 MATERIAL
% =============================================================================
% =============================================================================
% BEGIN APPENDED FOURTEENTH-CYCLE V46 MATERIAL
% =============================================================================

\section{Fourteenth-cycle V46: literal source extraction and the
flat-\(K_4\) bosonic Hessian}
\label{sec:v14v46}

\subsection{Source files and claim boundary}

This version returns to the two primary manuscripts:
\[
\begin{array}{p{0.58\linewidth}|r|r}
\textbf{file}&\textbf{lines}&\textbf{bytes}\\ \hline
\texttt{Renewal\_Einstein\_Standard\_Model\_Physical\_Source\_Metric\_
Duhamel\_Ancestry\_and\_Quantum\_Continuum\_Programme\_V30.tex}
 &32301&1262823\\
\texttt{Renewal\_Grand\_Tensor\_Monograph\_V067.tex}
 &68938&2804854.
\end{array}                                                  \tag{46.1}
\]
Their SHA--256 digests are
\[
\begin{aligned}
 &\texttt{E0AE0E0990DE472021A1E90A9C10D41C3F747F28D97A5B1D8F1281852E3A0E69},
                                                               \tag{46.2}\\
 &\texttt{FACD058EC8BB210AD8B296C9D7F4D78A229D08D7A2AE378808700655CB8BC0FE}.
                                                               \tag{46.3}
\end{aligned}
\]

The attached Einstein--SM V30 proves a finite transfer-log
acquisition theorem, a phase-resolved line-source geometry, and a
relation-covariant common-action command compiler.  Its own terminal
status leaves the active common-action transfer, current/stress,
charge, and Einstein cards unpopulated.

The Grand Tensor selects the structural \(K_4\) relation cell,
proves
\[
 B_4B_4^*=4I_{W_4},                                       \tag{46.4}
\]
classifies minimal weak-Higgs and plaquette actions, proves finite
Legendre reciprocity, and gives a conditional weak
Einstein--matter handoff.  It also states explicitly that a physical
metric, action, transport, and regulator control are separate loaded
data.  Therefore (46.4) is a source-incidence Gram, not the reduced
Hessian sought in V44.

\subsection{What the attached transfer theorem actually supplies}

On a physically transported finite Hilbert bundle, the attached V8
theorem starts from
\[
 P_{\theta,\ell}
 =U_\theta^*e^{-\ell H_\theta}U_\theta
 =e^{-\ell\widetilde H_\theta}                            \tag{46.5}
\]
and proves the exact central secant
\[
 -\frac{\log P_{\varepsilon,\ell}
        -\log P_{-\varepsilon,\ell}}
       {2\varepsilon\ell}
 =
 \frac{\widetilde H_\varepsilon-
       \widetilde H_{-\varepsilon}}{2\varepsilon}.         \tag{46.6}
\]
It also proves lag independence, gain removal after scalar
centering, a two-lag incidence test, a logarithm perturbation bound,
and the obstruction to taking a logarithm after a nonreducing
compression.

These are genuine source theorems and agree with the operator
firewalls in V39--V43.  They do not supply the matrices on the active
RPESM regulator.  The attached acquisition card still asks for the
transport \(U_\theta\), full or reducing transfers at
\(\theta=\pm\varepsilon\), a spectral floor and outward radii, a
held-out lag, and a third-derivative bound.  Hence no entry of
\[
 [T_{{\rm phys},h}],\quad
 [T_{{\rm clock},h}],\quad
 \epsilon_{{\rm clock},h}                                 \tag{46.7}
\]
can be assigned a numerical value from (46.5)--(46.6).

\subsection{Two inequivalent visibility notions}

The attached V21 line-source visibility is
\[
 \mathfrak v_{\rm ph}
 =\left|\mathbb E_\mu u\right|,                            \tag{46.8}
\]
where \(u\) is a determinant/Pfaffian phase.  The V42
gauge-projection visibility is
\[
 v_G(\psi)=\langle\psi,P_G\psi\rangle=\|P_G\psi\|^2.       \tag{46.9}
\]

\begin{countertheorem}[Phase and Haar visibility do not bound one another]
\label{cth:v14v46-two-visibilities}
There is no universal positive function \(F\) such that either
\[
 \mathfrak v_{\rm ph}\ge F(v_G)
 \quad\hbox{or}\quad
 v_G\ge F(\mathfrak v_{\rm ph})                            \tag{46.10}
\]
holds for all finite packets.
\end{countertheorem}

\begin{proof}
For the first direction, take a trivial gauge representation, so
\(P_G=I\) and \(v_G=1\).  On a two-point probability space with the
uniform law, take the phase values \(u=(1,-1)\).  Then
\(\mathfrak v_{\rm ph}=0\).

For the reverse direction, take a one-point probability space with
\(u=1\), so \(\mathfrak v_{\rm ph}=1\).  Let \(U(1)\) act on
\({\cal H}=\mathbb C\) by a nontrivial character.  Haar averaging
gives \(P_G=0\), so \(v_G=0\) for the unit vector.  These two examples
rule out both implications.
\end{proof}

The attached visibility Hessian
\[
 -\frac12D_cD_d\log|\mathbb E u|^2
 =\operatorname{Cov}(b_c,b_d)                             \tag{46.11}
\]
is therefore a phase-source metric.  It cannot populate V42's
trivial-isotypic visibility floor without an additional theorem
identifying the two carriers.

\subsection{The literal minimal bosonic action branch}

The Grand Tensor action-rigidity theorem supplies, on its stated
minimal local branch,
\[
\begin{aligned}
 {\mathscr S}_{\rm B}(U,\phi)
 ={}&
 \sum_{e=(x,y)}w_e\|U_e\phi_x-\phi_y\|^2\\
 &+\sum_x\lambda_x(\|\phi_x\|^2-\rho_x^2)^2\\
 &+\sum_f w_f\bigl(n-\operatorname{Re}\operatorname{Tr}U_f\bigr),
                                                               \tag{46.12}
\end{aligned}
\]
where
\[
 w_e>0,\qquad\lambda_x>0,\qquad w_f>0.                    \tag{46.13}
\]
The theorem fixes the functional forms under its covariance,
degree, positivity, and zero-set hypotheses.  It does not give the
values or cutoff dependence of the coefficients in (46.13).

Take the full tetrahedral two-complex: four vertices, six unoriented
edges with one chosen orientation, and four oriented triangular
faces.  Consider a flat homogeneous vacuum
\[
 U_{e,0}=I,\qquad \phi_{x,0}=\phi_0,\qquad
 \|\phi_0\|=\rho>0.                                       \tag{46.14}
\]
Write variations as
\[
 U_e(t)=e^{ta_e}+O(t^2),\qquad
 \phi_x(t)=\phi_0+t\eta_x+O(t^2),                          \tag{46.15}
\]
with \(a_e\in\mathfrak u(n)\).  Define
\[
\begin{aligned}
 (D_{H,0}(a,\eta))_e
 &=a_e\phi_0+\eta_{s(e)}-\eta_{t(e)},\\
 (d_1a)_{xyz}
 &=a_{xy}+a_{yz}+a_{zx},\\
 (R_0\eta)_x
 &=\operatorname{Re}\langle\phi_0,\eta_x\rangle.          \tag{46.16}
\end{aligned}
\]

\begin{theorem}[Exact flat-\(K_4\) minimal bosonic Hessian]
\label{thm:v14v46-bosonic-hessian}
At the vacuum (46.14), the second variation of (46.12) is
\[
 \boxed{
 \begin{aligned}
 D^2{\mathscr S}_{\rm B,0}[(a,\eta)]^2
 ={}&
 2\sum_e w_e\|D_{H,0}(a,\eta)_e\|^2\\
 &+\sum_f w_f\|(d_1a)_f\|_{\rm HS}^2\\
 &+8\sum_x\lambda_x
   \bigl(\operatorname{Re}
          \langle\phi_0,\eta_x\rangle\bigr)^2.
 \end{aligned}}                                           \tag{46.17}
\]
Equivalently, with diagonal weight operators
\(W_E,W_F,\Lambda\),
\[
 D^2{\mathscr S}_{\rm B,0}
 =
 2D_{H,0}^*W_ED_{H,0}
 +d_1^*W_Fd_1
 +8R_0^*\Lambda R_0.                                      \tag{46.18}
\]
\end{theorem}

\begin{proof}
For an edge,
\[
 U_e(t)\phi_{s(e)}(t)-\phi_{t(e)}(t)
 =tD_{H,0}(a,\eta)_e+O(t^2).
\]
The second derivative of its squared norm at zero is therefore
\(2\|D_{H,0}(a,\eta)_e\|^2\).

For a site,
\[
 \|\phi_0+t\eta\|^2-\rho^2
 =2t\operatorname{Re}\langle\phi_0,\eta\rangle+O(t^2),
\]
so the second derivative of the squared radial defect is
\(8(\operatorname{Re}\langle\phi_0,\eta\rangle)^2\).

At a flat face, the first variation of the holonomy is
\((d_1a)_f\).  The class function
\(n-\operatorname{Re}\operatorname{Tr}e^{tX}\) has zero first
derivative and second derivative
\(\|X\|_{\rm HS}^2\) at \(X=0\) for anti-Hermitian \(X\).  Since the
first derivative vanishes, the second-order acceleration of the
holonomy path contributes nothing.  Summing the three types proves
(46.17), and polarization gives (46.18).
\end{proof}

This is a derived symbolic Hessian from a literal source action
shape.  It is not a numerical Hessian of a selected RPESM
configuration.

\subsection{Exact gauge kernel on the weak \(U(2)\) branch}

For vertex gauge parameters
\(\xi=(\xi_x)_{x\in K_4}\in\mathfrak u(2)^4\), define
\[
 \Gamma_0\xi
 =\bigl(d_0\xi,(\xi_x\phi_0)_x\bigr),\qquad
 (d_0\xi)_{xy}=\xi_y-\xi_x.                               \tag{46.19}
\]

\begin{theorem}[The minimal flat bosonic kernel is precisely gauge]
\label{thm:v14v46-gauge-kernel}
Assume the weak fundamental \(U(2)\) action, all coefficients in
(46.13) are positive, and all four triangular faces are retained.
Then
\[
 \boxed{
 \ker D^2{\mathscr S}_{\rm B,0}
 =\operatorname{Ran}\Gamma_0.}                            \tag{46.20}
\]
Consequently the Hessian is strictly positive on the orthogonal
quotient by infinitesimal vertex gauge transformations.
\end{theorem}

\begin{proof}
Gauge covariance gives
\(\operatorname{Ran}\Gamma_0\subset\ker D^2{\mathscr S}_{\rm B,0}\);
this also follows by direct substitution in (46.16).

Conversely, the sum of nonnegative terms in (46.17) can vanish only
if
\[
 d_1a=0,\qquad D_{H,0}(a,\eta)=0,\qquad R_0\eta=0.         \tag{46.21}
\]
The tetrahedron is a contractible simplicial complex, so its first
cohomology vanishes.  Explicitly, choose vertex \(1\), set
\(\xi_1=0\), \(\xi_x=a_{1x}\), and use the triangle
\((1,x,y)\) in \(d_1a=0\); this gives
\[
 a_{xy}=\xi_y-\xi_x=d_0\xi_{xy}.                          \tag{46.22}
\]
The middle equation in (46.21) then says
\[
 \eta_x-\xi_x\phi_0=\zeta
 \quad\hbox{independently of \(x\)}.                       \tag{46.23}
\]
Because \(\xi_x\) is anti-Hermitian, the last equation in (46.21)
reduces to
\[
 \operatorname{Re}\langle\phi_0,\zeta\rangle=0.           \tag{46.24}
\]
The \(U(2)\)-orbit of a nonzero vector is the sphere of the same
norm; hence its infinitesimal orbit is exactly the tangent space
defined by (46.24).  Thus there is
\(\xi_0\in\mathfrak u(2)\) with
\(\zeta=\xi_0\phi_0\).  Replace every \(\xi_x\) by
\(\xi_x+\xi_0\).  Equation (46.22) is unchanged and (46.23) becomes
\(\eta_x=\xi_x\phi_0\).  Hence \((a,\eta)=\Gamma_0\xi\).
\end{proof}

Define the unweighted constraint synthesis
\[
 {\cal C}_0(a,\eta)
 =\bigl(D_{H,0}(a,\eta),d_1a,R_0\eta\bigr)                \tag{46.25}
\]
and let
\(\sigma_0=\sigma_{\min}^{+}({\cal C}_0)>0\).
If
\[
 c_{\rm B}
 =\min\{2\min_ew_e,\min_fw_f,8\min_x\lambda_x\},          \tag{46.26}
\]
then Theorem~\ref{thm:v14v46-gauge-kernel} gives the explicit
symbolic quotient floor
\[
 \boxed{
 D^2{\mathscr S}_{\rm B,0}
 \succeq c_{\rm B}\sigma_0^2 I
 \quad\hbox{on }(\operatorname{Ran}\Gamma_0)^\perp.}       \tag{46.27}
\]
This is the first literal positive reduced Hessian block extracted in
the present cycle.  Its cofinal use still requires lower bounds for
the source coefficients and for \(\sigma_0\) after all metric and
cutoff normalizations.

\subsection{Why the full coupled Hessian is still unavailable}

The Grand Tensor's source-native ADM Hamiltonian has gravitational
kinetic factor
\[
 \frac{\chi^{-1}}{\sqrt q}
 \left(\pi^{ab}\pi_{ab}-\frac12\pi^2\right).              \tag{46.28}
\]

\begin{countertheorem}[The unreduced DeWitt kinetic form is indefinite]
\label{cth:v14v46-dewitt}
At \(q=I\), decompose a symmetric momentum as
\[
 \pi=\pi_0+\frac{\tau}{3}I,\qquad\operatorname{Tr}\pi_0=0.
\]
Then
\[
 \boxed{
 \pi^{ab}\pi_{ab}-\frac12\pi^2
 =\|\pi_0\|_{\rm HS}^2-\frac{\tau^2}{6}.}                 \tag{46.29}
\]
It has five positive shape directions and one negative conformal
direction.
\end{countertheorem}

\begin{proof}
Orthogonality of trace and traceless matrices gives
\[
 \|\pi\|_{\rm HS}^2=\|\pi_0\|_{\rm HS}^2+\tau^2/3.
\]
Subtracting \(\tau^2/2\) gives (46.29).
\end{proof}

Moreover, the Grand Tensor explicitly keeps lapse and shift affine
in the unreduced canonical action.  Their Hessian is therefore zero,
as classified abstractly in V27 and V36.  A positive
lapse--shift block would belong to a declared gauge-fixed or reduced
system.

The following literal rows are absent:

\begin{enumerate}
\item a selected finite gravitational action matrix and its
Hamiltonian-constraint reduction on the \(K_4\) carrier;
\item the metric dependence and cutoff normalization of
\(w_e,w_f,\lambda_x,\rho_x\);
\item the gravity--gauge--Higgs cross derivatives produced by that
dependence;
\item the fermion, Yukawa, ghost, boundary, and counterterm Hessian
blocks on the same background;
\item a projector onto the complete regular constraint quotient; and
\item a two-sided ellipticity window for the resulting coupled
physical Hessian.
\end{enumerate}

Therefore (46.27) cannot be promoted to the full
\(G_{K_4}^{\rm red}\) in V44.

\subsection{Source-population table}

\[
\begin{array}{p{0.31\linewidth}|p{0.19\linewidth}|p{0.39\linewidth}}
\textbf{row}&\textbf{status}&\textbf{evidence}\\ \hline
\(K_4\) incidence/cycle split&literal&
Grand Tensor \(B_4B_4^*=4I\)\\
minimal weak/Higgs/face action shape&literal conditional&
action-rigidity theorem\\
flat minimal bosonic Hessian&derived here&
sum of squares (46.17)\\
weak \(U(2)\) gauge kernel&derived here&
tetrahedral \(H^1=0\)\\
bosonic quotient floor&symbolic&
positive weights and \(\sigma_0\)\\
active coefficient values&missing&
only positivity is stated\\
physical gravitational Hessian&missing&
unreduced DeWitt form indefinite\\
lapse/shift stiffness&forbidden unreduced&
affine multiplier variables\\
common-action transfer matrices&acquisition only&
attached V8 card unpopulated\\
phase-source visibility&theorem only&
no numerical RPESM line card\\
gauge-Haar visibility&missing&
different quantity from phase visibility\\
weak Einstein continuum&conditional&
Grand Tensor handoff hypotheses unpopulated
\end{array}                                                 \tag{46.30}
\]

\begin{decision}[Direction after V46]
\label{dec:v14v46-next}
The next obstruction is the negative trace direction in (46.29), not
the already positive minimal bosonic quotient.  V47 will construct
the finite linearized ADM constraint reduction at a flat \(K_4\)
background.  It will determine whether the Hamiltonian and momentum
constraints remove the DeWitt negative direction, separate the
transverse-traceless physical block from gauge and conformal
directions, and state the exact source rows needed to make that
reduction compatible with the \(K_4\) incidence complex.
\end{decision}

\subsection{Acquisition ledger}

\[
\begin{array}{p{0.30\linewidth}|p{0.20\linewidth}|p{0.39\linewidth}}
\textbf{gate}&\textbf{status after V46}&\textbf{content}\\ \hline
literal source return&completed&
two primary files read with hashes\\
transfer-log theorem&aligned&
proved compiler, inputs unpopulated\\
visibility name collision&resolved&
two counterexamples\\
minimal bosonic Hessian&proved&
exact \(K_4\) sum of squares\\
bosonic gauge quotient&proved&
kernel equals vertex gauge\\
bosonic finite floor&conditional&
positive source weights\\
gravity conformal direction&negative&
unreduced DeWitt signature\\
coupled reduced Hessian&missing&
constraint and cross rows absent\\
linearized ADM reduction&next&
flat \(K_4\) constraint complex
\end{array}                                                 \tag{46.31}
\]

\subsection{Parent-integrity record}

The installed parent was
\[
\begin{array}{c|c}
 \text{file}&
 \texttt{Renewal\_SM\_Einstein\_Continuum\_Research\_Notes\_V45.tex}\\
 \text{lines}&18050\\
 \text{bytes}&589476\\
 \text{SHA--256}&
 \texttt{E2B0B5784B12553E5C14204260E05CA3AFB283E85719B705BE5382282AC53220}.
\end{array}                                                 \tag{46.32}
\]

% =============================================================================
% END APPENDED FOURTEENTH-CYCLE V46 MATERIAL
% =============================================================================
% =============================================================================
% BEGIN APPENDED FOURTEENTH-CYCLE V47 MATERIAL
% =============================================================================

\section{Fourteenth-cycle V47: the finite ADM constraint complex and
the flat-background conformal obstruction}
\label{sec:v14v47}

\subsection{The correct finite linear constraint quotient}

Let \(Q_h\) be a finite real inner-product space of spatial metric
perturbations on a loaded multi-cell regulator.  Use the metric to
identify momenta with another copy of \(Q_h\), and put
\[
 E_h=Q_h\oplus Q_h,\qquad
 \Omega_h((u,p),(v,r))=\langle u,r\rangle-\langle v,p\rangle.
                                                               \tag{47.1}
\]
Let
\[
 C_{H,h}:Q_h\longrightarrow L_h,\qquad
 C_{M,h}:Q_h\longrightarrow S_h                              \tag{47.2}
\]
be the linearized Hamiltonian and momentum constraint operators.
The linear constraint surface and Hamiltonian gauge directions are
\[
 {\cal K}_h
 =\{(u,p):C_{H,h}u=0,\ C_{M,h}p=0\},                       \tag{47.3}
\]
\[
 {\cal G}_h
 =\{(C_{M,h}^*\beta,-C_{H,h}^*N):
       \beta\in S_h,\ N\in L_h\}.                          \tag{47.4}
\]
The two terms in (47.4) are respectively the shift-generated metric
variation and lapse-generated momentum variation.

\begin{theorem}[Finite first-class ADM Hodge quotient]
\label{thm:v14v47-hodge}
Assume the linear first-class complex identity
\[
 \boxed{C_{H,h}C_{M,h}^*=0,}                              \tag{47.5}
\]
equivalently \(C_{M,h}C_{H,h}^*=0\).  Define
\[
 {\cal H}_h
 =\ker C_{H,h}\cap\ker C_{M,h}.                            \tag{47.6}
\]
Then
\[
 {\cal G}_h\subset{\cal K}_h,\qquad
 {\cal K}_h/{\cal G}_h
 \cong{\cal H}_h\oplus{\cal H}_h,                          \tag{47.7}
\]
and the induced symplectic form is the restriction of (47.1) to the
right side.  In particular it is nondegenerate.
\end{theorem}

\begin{proof}
Equation (47.5) and its adjoint show directly that every vector in
(47.4) satisfies (47.3).

For the metric component, the orthogonal representative of a class
in
\[
 \ker C_{H,h}/\operatorname{Ran}C_{M,h}^*
\]
lies in
\[
 \ker C_{H,h}\cap
 (\operatorname{Ran}C_{M,h}^*)^\perp
 =\ker C_{H,h}\cap\ker C_{M,h}
 ={\cal H}_h.                                              \tag{47.8}
\]
For the momentum component, the same argument gives
\[
 \ker C_{M,h}/\operatorname{Ran}C_{H,h}^*
 \cong
 \ker C_{M,h}\cap\ker C_{H,h}
 ={\cal H}_h.                                              \tag{47.9}
\]
This proves the vector-space identification.  If
\((u,p)\in{\cal H}_h\oplus{\cal H}_h\) is symplectically orthogonal
to that space, testing against \((0,r)\) and \((v,0)\) gives
\(u=0\) and \(p=0\).  Hence the reduced form is nondegenerate.
\end{proof}

A relation projector alone is not this quotient.  One must acquire
both differential operators in (47.2), verify (47.5), and control
their ranks.

\subsection{The local \(K_4\) packet does not define (47.2)}

The locked tetrahedral packet supplies
\[
 W_4=\mathbf1_4^\perp,\qquad
 B_4B_4^*=4I_{W_4},\qquad
 \operatorname{Sym}(W_4)\cong\mathbb R^6.                 \tag{47.10}
\]
It identifies one local spatial coefficient module and its edge
cut--cycle split.  The ADM operators in (47.2), however, differentiate
a metric or momentum field between neighboring spatial cells:
\(C_M\) is a discrete divergence and \(C_H\) contains the
linearized spatial scalar curvature.  Their domains require a
global cell adjacency, face orientation, boundary convention, and
metric weights.

\begin{proposition}[One tetrahedron is insufficient for ADM reduction]
\label{prop:v14v47-one-cell}
The data in (47.10) do not determine either \(C_{H,h}\) or
\(C_{M,h}\).  More precisely, the same local space
\(\operatorname{Sym}(W_4)\) and the same internal incidence \(B_4\)
can be embedded into finite spatial regulators with different
linearized constraint ranks and different harmonic spaces
\({\cal H}_h\).
\end{proposition}

\begin{proof}
Attach two copies of the tetrahedral cell either along no face, along
one face, or with periodic identifications.  The internal vertex and
edge incidence in each copy remains (47.10).  The spaces of
cellwise-constant tensor fields have different continuity and
boundary relations in the three assemblies.  Consequently their
discrete divergences have different kernels; the periodic assembly
also retains global constant tensors that a Dirichlet assembly
removes.  Since the linearized curvature uses the same intercell
differences, its kernel changes as well.  Thus no construction from
the internal map \(B_4\) alone can determine both constraint
operators or their common harmonic space.
\end{proof}

This is exactly consistent with the Grand Tensor statement that the
local \(K_4\) selection supplies neither a global atlas nor a metric
and action.

\subsection{Sharp DeWitt compression criterion}

At a Euclidean spatial frame, let
\[
 Q_{\rm cell}=\operatorname{Sym}(\mathbb R^3)
 =Q_{\rm tf}\oplus Q_{\rm tr}.                             \tag{47.11}
\]
Let \(P_{\rm tr}\) be the orthogonal projection onto the trace
subspace, extended cellwise on \(Q_h\).  The DeWitt momentum form is
\[
 {\cal D}=I-\frac32P_{\rm tr},                             \tag{47.12}
\]
because
\[
 \langle p,{\cal D}p\rangle
 =p^{ab}p_{ab}-\frac12(\operatorname{Tr}p)^2.              \tag{47.13}
\]

\begin{theorem}[Sharp reduced DeWitt floor]
\label{thm:v14v47-dewitt-floor}
Let \(P_{{\cal H},h}\) project onto the harmonic space (47.6), and
put
\[
 \alpha_h=\|P_{\rm tr}P_{{\cal H},h}\|.                   \tag{47.14}
\]
Then the exact lower edge of the reduced momentum kinetic form is
\[
 \boxed{
 \lambda_{\min}
 \left(P_{{\cal H},h}{\cal D}P_{{\cal H},h}
       |_{\cal H_h}\right)
 =1-\frac32\alpha_h^2.}                                   \tag{47.15}
\]
It is positive exactly when
\[
 \alpha_h^2<\frac23.                                      \tag{47.16}
\]
\end{theorem}

\begin{proof}
For a unit \(p\in{\cal H}_h\),
\[
 \langle p,{\cal D}p\rangle
 =1-\frac32\|P_{\rm tr}p\|^2.                             \tag{47.17}
\]
The supremum of
\(\|P_{\rm tr}p\|\) over unit \(p\in{\cal H}_h\) is exactly the
operator norm in (47.14).  Taking the infimum of (47.17) proves
(47.15), and (47.16) follows.
\end{proof}

Thus removing the pure trace line is sufficient but not necessary;
the exact issue is the angle of the reduced harmonic space to the
whole trace sector.  This is the ADM specialization of the
constraint-escape machinery established earlier in the notes.

\subsection{Flat zero-momentum linearization fails the positivity test}

The canonical ADM Hamiltonian constraint has the schematic form
\[
 {\cal C}_H(q,\pi)
 =
 \frac1{\sqrt q}
 \left(\pi^{ab}\pi_{ab}-\frac12\pi^2\right)
 -\sqrt q\,(R(q)-2\Lambda)
 +{\cal C}_{H,\rm mat}.                                   \tag{47.18}
\]
At a vacuum with \(\pi_0=0\),
\[
 D_\pi{\cal C}_H(q_0,0)=0.                                \tag{47.19}
\]
Therefore its linearization constrains the metric curvature
variation but does not constrain a momentum variation.

\begin{countertheorem}[The homogeneous conformal momentum survives the
flat linear quotient]
\label{cth:v14v47-flat-conformal}
Assume a periodic or closed finite regulator on which the linearized
divergence and curvature annihilate spatially constant tensors.  Let
\[
 p_{\rm c}=c\,I_3
\]
on every cell, \(c\ne0\).  At a flat zero-momentum background,
\[
 C_{M,h}p_{\rm c}=0,\qquad C_{H,h}p_{\rm c}=0,             \tag{47.20}
\]
so \(p_{\rm c}\in{\cal H}_h\), while
\[
 \langle p_{\rm c},{\cal D}p_{\rm c}\rangle
 =-\frac12\|p_{\rm c}\|^2<0.                              \tag{47.21}
\]
Consequently the linearized constraint quotient does not make the
ADM kinetic Hessian positive.
\end{countertheorem}

\begin{proof}
A discrete divergence made from first differences annihilates a
constant tensor.  The operator \(C_{H,h}\) acts on the configuration
slot; under the identification in Theorem~\ref{thm:v14v47-hodge},
its application to a constant momentum representative also
vanishes because the same second-difference curvature symbol
annihilates constants.  Hence (47.20).  A pure trace unit vector is
an eigenvector of (47.12) with eigenvalue \(-1/2\), proving
(47.21).
\end{proof}

The obstruction is not repaired by the quadratic Hamiltonian
constraint at the level of its tangent space.  The momentum term in
(47.18) first appears at second order around \(\pi_0=0\).

\begin{proposition}[Second-order constraint information is lost by tangent
reduction]
\label{prop:v14v47-second-order}
Let a scalar constraint near the origin have
\[
 C(x)=\ell(x)+Q(x,x)+O(\|x\|^3),                           \tag{47.22}
\]
and let \(v\in\ker\ell\).  A curve
\(x(t)=tv+O(t^2)\) lying in \(C^{-1}(0)\) must satisfy
\[
 Q(v,v)+\ell(x_2)=0                                       \tag{47.23}
\]
at order \(t^2\), where \(x_2\) is its quadratic correction.
The tangent condition \(\ell(v)=0\) alone contains no information
about the sign or solvability of (47.23).
\end{proposition}

\begin{proof}
Insert \(x(t)=tv+t^2x_2+o(t^2)\) into (47.22).  The coefficient of
\(t\) is \(\ell(v)\), and the coefficient of \(t^2\) is
\(\ell(x_2)+Q(v,v)\).
\end{proof}

Applied to (47.18), the negative conformal momentum enters
\(Q(v,v)\).  It must be treated by a nonlinear constraint chart,
a nonzero-momentum background, or a true-clock deparametrization;
it cannot be deleted merely by naming the linear Hamiltonian
constraint.

\subsection{The nonzero mean-curvature entrance}

Let the background momentum be
\[
 \pi_0^{ab}=\kappa q_0^{ab}.                               \tag{47.24}
\]
At \(q_0=I\), the momentum derivative of the kinetic numerator is
\[
 \begin{aligned}
 D_\pi\!
 \left(\pi^{ab}\pi_{ab}-\frac12\pi^2\right)_{\pi_0}
 [\delta\pi]
 &=
 2\pi_0:\delta\pi-\pi_0\,\operatorname{Tr}\delta\pi\\
 &=-\kappa\,\operatorname{Tr}\delta\pi.                   \tag{47.25}
 \end{aligned}
\]
Here the second occurrence of \(\pi_0\) in the first line denotes its
trace \(3\kappa\).

\begin{corollary}[CMC regularizes the trace-momentum derivative]
\label{cor:v14v47-cmc}
If \(\kappa\ne0\), the linear Hamiltonian constraint has a nonzero
derivative in the trace-momentum direction.  At fixed metric
variation, its tangent equation determines
\(\operatorname{Tr}\delta\pi\) uniquely.  Pure momentum tangents with
\(\delta q=0\) are therefore traceless and have positive DeWitt
kinetic form.
\end{corollary}

\begin{proof}
Equation (47.25) is a nonzero scalar multiple of
\(\operatorname{Tr}\delta\pi\).  Solve the scalar tangent constraint
for that trace.  On the resulting pure momentum kernel,
\(P_{\rm tr}\delta\pi=0\), so (47.12) equals the identity.
\end{proof}

This is only a local regularity statement.  It does not construct a
global constant-mean-curvature slice, solve the nonlinear
Lichnerowicz equation, or identify the physical Hamiltonian selected
by the source clock.

\subsection{Three admissible gravitational routes}

The source programme must now choose one of three mathematically
different routes:

\[
\begin{array}{p{0.23\linewidth}|p{0.32\linewidth}|p{0.34\linewidth}}
\textbf{route}&\textbf{required finite object}&\textbf{output}\\ \hline
nonlinear Dirac reduction&
constraint manifold and gauge slice&
reduced symplectic Hamiltonian\\
CMC/York branch&
\(\kappa\ne0\), discrete elliptic solve&
regular trace elimination\\
physical-clock deparametrization&
clock variable and solvable clock momentum&
true positive evolution generator.
\end{array}                                                 \tag{47.26}
\]
The OS transfer branch of V39 may construct the last output directly,
but only after the attached source populates the same-history
transfer and clock card.

\subsection{Exact missing source rows}

A finite ADM reduction card must contain
\[
 \mathfrak A_h^{\rm ADM}=
 \bigl(
 [\mathcal C_{H,h}],[\mathcal C_{M,h}],
 [C_{H,h}],[C_{M,h}],
 \Delta_{{\rm FC},h},
 [P_{{\cal H},h}],\alpha_h,
 [q_{0,h},\pi_{0,h}],
 [\hbox{gauge slice}_h],
 [\hbox{clock solve}_h]
 \bigr),                                                   \tag{47.27}
\]
where
\[
 \Delta_{{\rm FC},h}
 =\|C_{H,h}C_{M,h}^*\|^2                                 \tag{47.28}
\]
is the linear first-class residual.  The source provides the abstract
constraint and HDA requirements, but no literal multi-cell matrices
for (47.27) on the selected \(K_4\)-based regulator.

\begin{decision}[Direction after V47]
\label{dec:v14v47-next}
V48 will develop the CMC/York route because it is the first branch on
which the trace derivative is regular.  It will prove a finite
implicit constraint solve with explicit inverse and remainder
bounds, derive the reduced quadratic form after eliminating the
trace momentum, and expose the discrete elliptic/York operator that
must be acquired from the source.  If that operator has no uniform
floor, the programme will return to the physical-clock transfer
route rather than recycle the flat linear quotient.
\end{decision}

\subsection{Acquisition ledger}

\[
\begin{array}{p{0.30\linewidth}|p{0.20\linewidth}|p{0.39\linewidth}}
\textbf{gate}&\textbf{status after V47}&\textbf{content}\\ \hline
finite ADM linear quotient&proved&
Hodge space (47.6)\\
single-\(K_4\) sufficiency&refuted&
global differences missing\\
reduced DeWitt floor&proved&
sharp trace angle\\
flat conformal removal&refuted&
constant negative mode survives\\
second-order constraint&exposed&
tangent quotient loses it\\
CMC trace derivative&positive route&
nonzero mean curvature\\
literal multi-cell constraint matrices&missing&
source acquisition (47.27)\\
CMC/York implicit solve&next&
quantitative nonlinear reduction
\end{array}                                                 \tag{47.29}
\]

\subsection{Parent-integrity record}

The installed parent was
\[
\begin{array}{c|c}
 \text{file}&
 \texttt{Renewal\_SM\_Einstein\_Continuum\_Research\_Notes\_V46.tex}\\
 \text{lines}&18514\\
 \text{bytes}&605086\\
 \text{SHA--256}&
 \texttt{77622A77D62C0213C4B30EDD931E9E050882CE0F62E0DAA540B2E54895E17B7D}.
\end{array}                                                 \tag{47.30}
\]

% =============================================================================
% END APPENDED FOURTEENTH-CYCLE V47 MATERIAL
% =============================================================================
% =============================================================================
% BEGIN APPENDED FOURTEENTH-CYCLE V48 MATERIAL
% =============================================================================

\section{Fourteenth-cycle V48: quantitative nonzero-mean-curvature
constraint charts and the finite York compiler}
\label{sec:v14v48}

\subsection{A quantitative finite implicit constraint theorem}

Let \(X,Y\) be finite-dimensional Hilbert spaces and let
\[
 F:X\times Y\longrightarrow Y,\qquad F(0,0)=0.            \tag{48.1}
\]
Write
\[
 A=D_yF(0,0),\qquad \sigma_{\min}(A)\ge\mu>0.              \tag{48.2}
\]

\begin{theorem}[Certified implicit constraint chart]
\label{thm:v14v48-ift}
Fix radii \(r_X,r_Y>0\).  Suppose throughout
\(\|x\|\le r_X,\ \|y\|\le r_Y\),
\[
 \|D_yF(x,y)-A\|\le\frac\mu2,                              \tag{48.3}
\]
and, for every \(\|x\|\le r_X\),
\[
 \|A^{-1}F(x,0)\|\le\frac{r_Y}{2}.                         \tag{48.4}
\]
Then for each such \(x\) there is a unique \(y(x)\) in the ball
\(\|y\|\le r_Y\) satisfying
\[
 F(x,y(x))=0.                                              \tag{48.5}
\]
Moreover,
\[
 \|y(x)\|\le2\|A^{-1}F(x,0)\|,                             \tag{48.6}
\]
\[
 \|D_yF(x,y(x))^{-1}\|\le\frac2\mu,                        \tag{48.7}
\]
and
\[
 Dy(x)=-D_yF(x,y(x))^{-1}D_xF(x,y(x)).                    \tag{48.8}
\]
\end{theorem}

\begin{proof}
For fixed \(x\), define
\[
 {\cal T}_x(y)=y-A^{-1}F(x,y).                             \tag{48.9}
\]
By (48.2)--(48.3),
\[
 \|D_y{\cal T}_x(y)\|
 =\|A^{-1}(A-D_yF(x,y))\|\le\frac12.                      \tag{48.10}
\]
Together with (48.4), this shows that \({\cal T}_x\) maps the
closed \(r_Y\)-ball into itself and is a contraction.  Its unique
fixed point is (48.5).  Comparing that fixed point with zero gives
\[
 \|y(x)\|\le\frac{\|{\cal T}_x(0)\|}{1-1/2},
\]
which is (48.6).

Factor
\[
 D_yF=A\left[I+A^{-1}(D_yF-A)\right].
\]
The bracket differs from the identity by at most \(1/2\), so its
inverse has norm at most two.  This proves (48.7).
Differentiating (48.5) gives (48.8).
\end{proof}

All constants are finite and directly auditable.  The load-bearing
quantity is the least singular value \(\mu\); pointwise invertibility
without a cofinal lower bound is not enough.

\subsection{The local trace-momentum branch}

At one Euclidean cell, decompose
\[
 \pi=\pi_{\rm tf}+\frac{\tau}{3}I,\qquad
 \operatorname{Tr}\pi_{\rm tf}=0.                         \tag{48.11}
\]
The DeWitt numerator is
\[
 \pi^{ab}\pi_{ab}-\frac12\pi^2
 =\|\pi_{\rm tf}\|^2-\frac{\tau^2}{6}.                    \tag{48.12}
\]
Collect the curvature, matter, cosmological, traceless-momentum, and
other fixed constraint contributions into a scalar \(A(x)\).  The
local scalar constraint then has the normal form
\[
 F(x,\tau)=A(x)-\frac{\tau^2}{6}=0.                        \tag{48.13}
\]

\begin{proposition}[Exact nonzero-mean-curvature branches]
\label{prop:v14v48-trace-branch}
On a connected set where \(A(x)>0\), (48.13) has exactly two
nonzero branches
\[
 \boxed{\tau_\pm(x)=\pm\sqrt{6A(x)}.}                      \tag{48.14}
\]
Their derivatives are
\[
 D\tau_\pm=\frac{3}{\tau_\pm}DA,                           \tag{48.15}
\]
\[
 D^2\tau_\pm
 =\frac3{\tau_\pm}D^2A
 -\frac9{\tau_\pm^3}\,DA\otimes DA.                       \tag{48.16}
\]
If \(A\ge a_*>0\), the inverse trace-constraint derivative has the
uniform bound
\[
 \left|\left(D_\tau F\right)^{-1}\right|
 =\frac3{|\tau_\pm|}
 \le\sqrt{\frac{3}{2a_*}}.                                \tag{48.17}
\]
\end{proposition}

\begin{proof}
Equation (48.14) is the complete solution of
\(\tau^2=6A\).  Differentiate that identity once and twice:
\[
 2\tau D\tau=6DA,\qquad
 2D\tau\otimes D\tau+2\tau D^2\tau=6D^2A.
\]
Substitution of the first equation into the second gives
(48.15)--(48.16).  Since \(D_\tau F=-\tau/3\), (48.17)
follows.
\end{proof}

The sign in (48.14) is time-orientation data.  Squaring the
constraint or retaining only \(A\) erases it.

\begin{countertheorem}[The CMC chart degenerates at the flat branch point]
\label{cth:v14v48-branch-point}
For \(F(A,\tau)=A-\tau^2/6\), the two solution branches meet at
\((A,\tau)=(0,0)\), where
\[
 D_\tau F=0.                                               \tag{48.18}
\]
No differentiable single-valued solution \(\tau(A)\) exists on a
two-sided neighbourhood of \(A=0\).
\end{countertheorem}

\begin{proof}
For \(A>0\), every real solution has
\(|\tau(A)|=\sqrt{6A}\), whose derivative has magnitude
\(3/|\tau(A)|\to\infty\) as \(A\downarrow0\).  For \(A<0\)
there is no real solution.  Hence no two-sided differentiable chart
exists.
\end{proof}

This is the nonlinear form of the flat-background failure in V47.

\subsection{The Hessian after solving a constraint}

Constraint elimination is not auxiliary stationarity elimination.
The Schur complement from V36 applies when \(D_yS=0\) is solved.
Here one solves \(F=0\).

Let \(\iota(x)=(x,y(x))\), and at the base point put
\[
 B=D_xF,\qquad
 J_{\rm tan}=
 \begin{pmatrix}I\\-A^{-1}B\end{pmatrix}.                 \tag{48.19}
\]

\begin{theorem}[Constrained reduced Hessian]
\label{thm:v14v48-constrained-hessian}
Let \(S:X\times Y\to\mathbb R\) be \(C^2\), and suppose the base
point is constrained stationary:
\[
 DS+DF^*\lambda=0                                         \tag{48.20}
\]
for a multiplier \(\lambda\in Y\).  Then
\[
 \boxed{
 D^2(S\circ\iota)
 =
 J_{\rm tan}^*
 \left(D^2S+D^2\langle\lambda,F\rangle\right)
 J_{\rm tan}.}                                             \tag{48.21}
\]
\end{theorem}

\begin{proof}
Twice differentiating \(F\circ\iota=0\) gives
\[
 DF\,D^2\iota+
 J_{\rm tan}^*(D^2F)J_{\rm tan}=0.                         \tag{48.22}
\]
The second chain rule for \(S\circ\iota\) is
\[
 D^2(S\circ\iota)
 =J_{\rm tan}^*(D^2S)J_{\rm tan}+DS\,D^2\iota.            \tag{48.23}
\]
Use (48.20) in the last term and then (48.22).  This adds
\(J_{\rm tan}^*D^2\langle\lambda,F\rangle J_{\rm tan}\),
proving (48.21).
\end{proof}

Thus even when the trace variable is solved with a uniform inverse,
positivity must be tested on the Hessian of the constrained
Lagrangian, not on the ambient DeWitt block alone.

\begin{countertheorem}[Regular constraint elimination need not yield a
positive Hessian]
\label{cth:v14v48-regular-not-positive}
Let
\[
 F(x,y)=y,\qquad S(x,y)=-x^2+y^2.
\]
Then \(D_yF=1\), the constraint chart is \(y(x)=0\), and the
base point is constrained stationary, but
\[
 D^2(S\circ\iota)=-2<0.                                   \tag{48.24}
\]
\end{countertheorem}

\begin{proof}
All statements follow by differentiation.
\end{proof}

The CMC inverse floor and the physical energy floor are separate
rows.

\subsection{Finite momentum and conformal York stages}

A genuine CMC/York construction does more than solve one independent
trace variable at every cell.  It first solves the momentum
constraint and then a coupled conformal scalar equation.

Let \(C_{M,h}:Q_h\to S_h\) be the finite momentum divergence and put
\[
 \Delta_{L,h}=C_{M,h}C_{M,h}^*.                            \tag{48.25}
\]
For a prescribed momentum source \(j_h\), seek
\[
 p_h=p_{{\rm TT},h}+C_{M,h}^*W_h,\qquad
 p_{{\rm TT},h}\in\ker C_{M,h}.                            \tag{48.26}
\]

\begin{proposition}[Finite momentum-constraint solve]
\label{prop:v14v48-momentum}
The equation \(C_{M,h}p_h=j_h\) is solvable exactly when
\[
 j_h\in\operatorname{Ran}C_{M,h}.                          \tag{48.27}
\]
Its minimum-norm potential is
\[
 W_h=\Delta_{L,h}^\dagger j_h.                             \tag{48.28}
\]
If
\[
 \lambda_{\min}^{+}(\Delta_{L,h})\ge\gamma_M>0,            \tag{48.29}
\]
then
\[
 \|W_h\|\le\gamma_M^{-1}\|j_h\|,\qquad
 \|C_{M,h}^*W_h\|
 \le\gamma_M^{-1/2}\|j_h\|.                               \tag{48.30}
\]
\end{proposition}

\begin{proof}
The range of \(C_{M,h}C_{M,h}^*\) equals the range of
\(C_{M,h}\).  The Moore--Penrose inverse gives the minimum-norm
solution on that range.  The first bound follows from (48.29).
For the second,
\[
 \|C_{M,h}^*W_h\|^2
 =\langle W_h,\Delta_{L,h}W_h\rangle
 =\langle W_h,j_h\rangle
 \le\gamma_M^{-1}\|j_h\|^2.
\]
\end{proof}

After this step, let \(\psi_h\) denote the discrete conformal factor
and let
\[
 {\cal Y}_h(\psi_h;\mathfrak d_h)=0                        \tag{48.31}
\]
be the source-defined finite Lichnerowicz--York equation for the
remaining data \(\mathfrak d_h\).  Its Jacobian is
\[
 {\cal L}_{Y,h}=D_\psi{\cal Y}_h(\psi_{0,h};\mathfrak d_{0,h}).
                                                               \tag{48.32}
\]

\begin{theorem}[Finite York compiler]
\label{thm:v14v48-york}
Suppose:

\begin{enumerate}
\item the range condition (48.27) and floor (48.29) hold;
\item \(\sigma_{\min}({\cal L}_{Y,h})\ge\gamma_H>0\);
\item the nonlinear derivative and base residual satisfy
(48.3)--(48.4) with \(A={\cal L}_{Y,h}\);
\item the reconstructed metric is nondegenerate and the chosen
mean-curvature sign remains separated from zero.
\end{enumerate}

Then the momentum and Hamiltonian constraints have a unique local
finite CMC solution modulo the kernels already removed in
(48.28).  Its derivative is bounded explicitly by (48.7)--(48.8)
and (48.30), and the Hessian of every constrained observable is
given by (48.21).
\end{theorem}

\begin{proof}
Proposition~\ref{prop:v14v48-momentum} solves the linear momentum
stage on its supported range.  Apply
Theorem~\ref{thm:v14v48-ift} to (48.31) for the conformal stage.
The nondegeneracy and sign hypotheses keep the reconstructed point
inside the declared geometric branch.  The derivative and Hessian
statements are the cited formulas.
\end{proof}

The theorem is deliberately independent of a guessed discretization
of the continuum Lichnerowicz equation.  The actual map
\({\cal Y}_h\), boundary condition, weights, and source terms must be
taken from one common finite action.

\subsection{A CMC solution is not yet a physical Hamiltonian}

Solving constraints constructs a reduced phase chart.  Evolution
requires either a gauge-fixed lapse or a canonical clock pair
\((T,P_T)\).  A deparametrized constraint has the oriented local
form
\[
 {\cal C}_{\rm clock}
 =P_T+H_{\rm true}(z)=0,                                   \tag{48.33}
\]
or a multi-branch form from which one orientation has been selected.
Only then does the gauge \(T=t\) produce \(H_{\rm true}\).

The Grand Tensor's positive lapse
\(\mathcal N_h=\beta_h\ell_{0,h}\) calibrates elapsed physical time,
but the inspected source passage does not provide a canonical
\(P_T\), an identity of the form (48.33), or a proof that the CMC
trace solution equals the logarithm of the OS transfer.  Those are
separate incidence rows.

\subsection{Source acquisition card}

The finite nonzero-mean-curvature card is
\[
 \mathfrak Y_h=
 \bigl(
 [C_{M,h}], [\Delta_{L,h}],\gamma_M,
 [j_h],
 [\mathcal Y_h],[\mathcal L_{Y,h}],\gamma_H,
 [r_X,r_Y],
 [\tau_{0,h}],\tau_*,
 [\lambda_h],
 [D^2(S_h+\langle\lambda_h,F_h\rangle)]
 \bigr).                                                   \tag{48.34}
\]
The Grand Tensor supplies abstract Hodge, CMC/ADM, constraint,
positive-screen, and action hypotheses, but it does not give the
literal matrices and nonlinear map in (48.34) on the selected
\(K_4\)-based regulator.

\begin{decision}[Direction after V48]
\label{dec:v14v48-next}
V49 will take the physical-clock route exposed by (48.33).  It will
prove the finite deparametrization theorem, treat the two-sheeted
quadratic clock constraint without losing orientation, and derive
the exact comparison inequality between the resulting classical
true Hamiltonian and the V39/V43 transfer logarithm.  This will
identify whether the CMC branch and the OS branch can meet without
assuming their generators are equal.
\end{decision}

\subsection{Acquisition ledger}

\[
\begin{array}{p{0.30\linewidth}|p{0.20\linewidth}|p{0.39\linewidth}}
\textbf{gate}&\textbf{status after V48}&\textbf{content}\\ \hline
quantitative constraint chart&proved&
inverse and radius bounds\\
nonzero trace branch&proved&
\(\tau_\pm=\pm\sqrt{6A}\)\\
flat branch continuation&refuted&
square-root singularity\\
constrained Hessian&proved&
Lagrangian restriction (48.21)\\
positivity after solve&separate&
explicit counterexample\\
momentum York stage&compiled&
supported vector Laplacian\\
conformal York stage&conditional&
Jacobian floor required\\
literal York matrices&missing&
source card (48.34)\\
physical deparametrization&next&
clock momentum and transfer match
\end{array}                                                 \tag{48.35}
\]

\subsection{Parent-integrity record}

The installed parent was
\[
\begin{array}{c|c}
 \text{file}&
 \texttt{Renewal\_SM\_Einstein\_Continuum\_Research\_Notes\_V47.tex}\\
 \text{lines}&18930\\
 \text{bytes}&619476\\
 \text{SHA--256}&
 \texttt{57FCF11F92AA93E0AA89DADB224F84489CA52EC4BE8047C228B147AF30329696}.
\end{array}                                                 \tag{48.36}
\]

% =============================================================================
% END APPENDED FOURTEENTH-CYCLE V48 MATERIAL
% =============================================================================
% =============================================================================
% BEGIN APPENDED FOURTEENTH-CYCLE V49 MATERIAL
% =============================================================================

\section{Fourteenth-cycle V49: finite physical-clock
deparametrization and calibrated transfer matching}
\label{sec:v14v49}

\subsection{Canonical clock reduction}

Let \((T,P_T)\) be a canonical clock pair and let
\((Z,\omega=d\vartheta)\) be the remaining finite phase space.
Consider the first-order constrained action
\[
 {\mathscr S}
 =\int
 \left(
 P_T\dot T+\vartheta_z(\dot z)-N{\cal C}(T,P_T,z)
 \right)dt.                                                \tag{49.1}
\]

\begin{theorem}[Affine-clock deparametrization]
\label{thm:v14v49-affine-clock}
Suppose
\[
 {\cal C}=P_T+H_{\rm true}(z)                              \tag{49.2}
\]
and impose the oriented gauge \(T=t\).  Then preservation of the
gauge fixes \(N=1\), the constraint gives
\[
 P_T=-H_{\rm true}(z),                                     \tag{49.3}
\]
and (49.1) reduces to
\[
 \boxed{
 {\mathscr S}_{\rm red}
 =\int\bigl(\vartheta_z(\dot z)-H_{\rm true}(z)\bigr)\,dt.}
                                                               \tag{49.4}
\]
Thus \(H_{\rm true}\) generates evolution in the declared clock
variable.
\end{theorem}

\begin{proof}
Hamilton's equation for the clock is
\[
 \dot T=N\partial_{P_T}{\cal C}=N.
\]
The gauge \(\dot T=1\) gives \(N=1\).  Substitute (49.3) and
\(\dot T=1\) in (49.1); the constraint term vanishes and the clock
canonical term becomes \(-H_{\rm true}\), proving (49.4).
\end{proof}

More generally, if
\[
 {\cal C}=A(T,z)P_T+B(T,z),\qquad A\ne0,                   \tag{49.5}
\]
then
\[
 P_T=-B/A,\qquad N=A^{-1},\qquad
 H_{\rm true}=B/A                                         \tag{49.6}
\]
in the gauge \(T=t\).  Hence lapse calibration and the true
Hamiltonian are reciprocal parts of the same clock constraint.

\subsection{Nonlinear and two-sheeted clocks}

\begin{theorem}[Regular nonlinear clock branch]
\label{thm:v14v49-nonlinear-clock}
Let
\[
 {\cal C}(P_T,z)=0,\qquad
 |\partial_{P_T}{\cal C}|\ge\mu_T>0                       \tag{49.7}
\]
on a selected connected branch.  Then locally
\[
 P_T=-H_{\rm true}(z),                                    \tag{49.8}
\]
and in the gauge \(T=t\)
\[
 N=\left(\partial_{P_T}{\cal C}\right)^{-1}.              \tag{49.9}
\]
The derivative of the true Hamiltonian is
\[
 DH_{\rm true}
 =
 \left(\partial_{P_T}{\cal C}\right)^{-1}D_z{\cal C}.      \tag{49.10}
\]
If the hypotheses of Theorem~\ref{thm:v14v48-ift} hold for the
clock momentum, its radius and inverse estimates apply with
\(\mu=\mu_T\).
\end{theorem}

\begin{proof}
The implicit-function theorem gives (49.8).  Hamilton's equation
\(\dot T=N\partial_{P_T}{\cal C}\) gives (49.9).
Differentiating
\({\cal C}(-H_{\rm true}(z),z)=0\) gives
\[
 -\partial_{P_T}{\cal C}\,DH_{\rm true}+D_z{\cal C}=0,
\]
which is (49.10).
\end{proof}

\begin{proposition}[Orientation of a quadratic clock]
\label{prop:v14v49-quadratic-clock}
Let \(H(z)>0\) and
\[
 {\cal C}=P_T^2-H(z)^2.                                   \tag{49.11}
\]
The constraint surface has two disjoint sheets
\[
 P_T=\sigma H(z),\qquad \sigma\in\{+1,-1\}.               \tag{49.12}
\]
In the gauge \(T=t\),
\[
 N=\frac1{2\sigma H(z)},\qquad
 H_{\rm true}^{(\sigma)}=-\sigma H(z).                    \tag{49.13}
\]
Therefore the positive true Hamiltonian \(H\) is the
\(\sigma=-1\) sheet, with the corresponding clock orientation in
the lapse.
\end{proposition}

\begin{proof}
Equations (49.12) solve (49.11).  Since
\(\partial_{P_T}{\cal C}=2P_T\), equation (49.9) gives the lapse.
The reduced clock term is \(P_T\dot T=\sigma H\), which must equal
\(-H_{\rm true}^{(\sigma)}\) in the canonical reduced action.
\end{proof}

A squared clock constraint without a sheet and orientation does not
define one physical Hamiltonian.

\begin{countertheorem}[The quadratic quantum constraint has two frequency
sectors]
\label{cth:v14v49-two-frequency}
Let \(\widehat H\ge0\) be self-adjoint and independent of \(T\).
With \(\widehat P_T=-i\hbar\partial_T\), the equation
\[
 (\widehat P_T^2-\widehat H^2)\Psi=0                       \tag{49.14}
\]
admits both families
\[
 \Psi_+(T)=e^{\,iT\widehat H/\hbar}\psi_+,\qquad
 \Psi_-(T)=e^{-iT\widehat H/\hbar}\psi_-.                  \tag{49.15}
\]
No unique one-parameter evolution follows until one frequency sheet
is selected.
\end{countertheorem}

\begin{proof}
Differentiate each exponential twice.  Both satisfy
\(-\hbar^2\partial_T^2\Psi=\widehat H^2\Psi\), and their first-order
generators have opposite signs.
\end{proof}

The passage from the selected unitary clock evolution to a positive
Euclidean transfer additionally requires the OS/reflection packet.
It is not a consequence of (49.11).

\subsection{Clock changes rescale the generator}

Let \(\widetilde T=f(T)\) with \(f'(T)>0\).  Canonical covariance
requires
\[
 P_T\,dT=P_{\widetilde T}\,d\widetilde T,\qquad
 P_T=f'(T)P_{\widetilde T}.                               \tag{49.16}
\]
For the affine constraint (49.2),
\[
 P_{\widetilde T}+\frac{H_{\rm true}}{f'(T)}=0             \tag{49.17}
\]
after division by \(f'(T)\).  Thus the generator per unit
\(\widetilde T\) is \(H_{\rm true}/f'(T)\).  A clock label without
its calibration does not determine an energy scale.

\subsection{Calibrated transfer-symbol inequality}

Let \(T_{\alpha,\ell}\) be a positive transfer at operational lag
\(\ell_\alpha>0\), and let its physical clock increment be
\(\tau_\alpha>0\).  Put
\[
 \widehat K_\alpha
 =-\ell_\alpha^{-1}\log T_{\alpha,\ell},\qquad
 r_\alpha=\frac{\ell_\alpha}{\tau_\alpha}.                \tag{49.18}
\]
For a normalized projected coherent state \(\psi_{\alpha,z}^P\),
define the energy per unit physical clock
\[
 E_{\alpha}^{T}(z)
 =-\tau_\alpha^{-1}
 \log\langle\psi_{\alpha,z}^P,
             T_{\alpha,\ell}\psi_{\alpha,z}^P\rangle.     \tag{49.19}
\]
Suppose
\[
 m_\alpha I\preceq T_{\alpha,\ell}\preceq I              \tag{49.20}
\]
on the relevant state carrier, and define
\[
 \delta_{{\rm var},\alpha}
 =
 \sup_{z\in K}
 \frac{{\rm Var}_{\psi_{\alpha,z}^P}(T_{\alpha,\ell})}
      {2\ell_\alpha m_\alpha^2},                           \tag{49.21}
\]
\[
 \epsilon_{K,\alpha}
 =
 \sup_{z\in K}
 |\sigma_\alpha^P(\widehat K_\alpha)(z)-K_{\rm cl,\alpha}(z)|,
                                                               \tag{49.22}
\]
\[
 \epsilon_{{\rm clock},\alpha}^{\rm cl}
 =
 \sup_{z\in K}
 |r_\alpha K_{\rm cl,\alpha}(z)-H_{\rm true,\alpha}(z)|.
                                                               \tag{49.23}
\]

\begin{theorem}[Physical-clock transfer handoff]
\label{thm:v14v49-clock-handoff}
Under (49.18)--(49.23),
\[
 \boxed{
 \sup_{z\in K}
 |E_\alpha^T(z)-H_{\rm true,\alpha}(z)|
 \le
 r_\alpha
 \bigl(\epsilon_{K,\alpha}
       +\delta_{{\rm var},\alpha}\bigr)
 +\epsilon_{{\rm clock},\alpha}^{\rm cl}.}                \tag{49.24}
\]
Consequently the transfer and deparametrized classical generators
have one common principal symbol if the right side tends to zero.
\end{theorem}

\begin{proof}
By definition,
\[
 E_\alpha^T
 =r_\alpha
 \left[
 -\ell_\alpha^{-1}
 \log\langle T_{\alpha,\ell}\rangle
 \right].                                                  \tag{49.25}
\]
Theorem~\ref{thm:v14v41-jensen} bounds the difference between the
bracket in (49.25) and
\(\sigma_\alpha^P(\widehat K_\alpha)\) by
\(\delta_{{\rm var},\alpha}\).  Add and subtract
\(r_\alpha K_{\rm cl,\alpha}\) and use
(49.22)--(49.23).
\end{proof}

This separates three independent errors: operator-symbol matching,
state transfer variance, and physical-clock calibration.  A held-out
two-lag identity tests whether one operational generator exists; it
does not by itself determine \(r_\alpha\).

If the transfer logarithm is unbounded as in V43, the same statement
holds on coherent vectors for which the logarithmic expectation is
finite, with the explicit quadratic bound (43.25) replacing
(49.21).

\subsection{Absolute transfer gain and the Einstein source}

The attached V8 source theorem permits acquired transfers known up
to positive gains:
\[
 \overline T_{\theta,\ell}
 =g_{\theta,\ell}T_{\theta,\ell}.                          \tag{49.26}
\]
Trace or Gibbs centering removes the scalar logarithm.  That is
sufficient for a centered Duhamel tangent, but not for a coupled
Einstein energy.

\begin{proposition}[Gain-free logarithms lose the vacuum-energy row]
\label{prop:v14v49-gain}
Let \(g=e^{-\ell E_0}>0\).  Then
\[
 -\ell^{-1}\log(gT)
 =-\ell^{-1}\log T+E_0I.                                  \tag{49.27}
\]
The scalar-centered logarithms of \(gT\) and \(T\) are identical,
and their excitation gaps are identical, but their absolute vacuum
energies differ by \(E_0\).  After division by the declared spatial
volume, this is precisely an additive cosmological source row.
\end{proposition}

\begin{proof}
The scalar \(gI\) commutes with \(T\), so
\[
 \log(gT)=(\log g)I+\log T.
\]
Equation (49.27) follows.  Centering annihilates \(E_0I\), while
differences of eigenvalues also annihilate it.  Coupling the constant
energy density to the spacetime volume changes the metric variation
by a term proportional to the metric, which is the cosmological
source contribution.
\end{proof}

Thus the gain-free V8 acquisition theorem and the zero-point
subtraction in V44 share the same Einstein firewall: an absolute
normalization or a declared counterterm is required.

\begin{countertheorem}[Two-lag centering cannot recover absolute energy]
\label{cth:v14v49-two-lag}
For any self-adjoint \(H\) and any scalar \(E_0\),
\[
 T_\ell=e^{-\ell H},\qquad
 \widetilde T_\ell=e^{-\ell(H+E_0I)}
\]
both obey the exact two-lag centered identity
\[
 \operatorname{ctr}
 (-\ell_1^{-1}\log T_{\ell_1})
 =
 \operatorname{ctr}
 (-\ell_2^{-1}\log T_{\ell_2}),                            \tag{49.28}
\]
and likewise for \(\widetilde T\), but their absolute generators
differ by \(E_0I\).
\end{countertheorem}

\begin{proof}
Every logarithmic generator is respectively \(H\) or
\(H+E_0I\), independently of lag.  Scalar centering removes the
difference.
\end{proof}

\subsection{The classical--quantum clock card}

The finite clock card exposed by this version is
\[
 \mathfrak C_\alpha=
 \bigl(
 [T,P_T], [{\cal C}_{\rm clock}],
 \mu_T,[\hbox{sheet/orientation}],
 [H_{\rm true}],
 \ell_\alpha,\tau_\alpha,r_\alpha,
 [T_{\alpha,\ell}],
 [g_{\alpha,\ell}],E_{0,\alpha},
 \epsilon_K,\delta_{\rm var},
 \epsilon_{\rm clock}^{\rm cl}
 \bigr).                                                   \tag{49.29}
\]
The Grand Tensor supplies a positive lapse calibration
\(\mathcal N_h=\beta_h\ell_{0,h}\) on its conditional ADM packet.
The attached V8 source supplies the transfer-log compiler.  The
inspected sources do not supply one canonical clock momentum,
constraint sheet, absolute transfer gain, and common
\(\ell_\alpha/\tau_\alpha\) calibration tying those two results
together.

\begin{decision}[Direction after V49]
\label{dec:v14v49-next}
V50 is the compulsory companion audit.  It will check whether the
latest Yang--Mills or Navier--Stokes note supplies an absolute
transfer normalization, a physical clock ratio, a nonzero-frequency
trace control, or a CMC/York inverse floor.  V51 will then attack the
coupled vacuum-energy and counterterm consistency condition, because
V44 and V49 show that this scalar row survives every centered
transfer and gap calculation while still sourcing Einstein's
equation.
\end{decision}

\subsection{Acquisition ledger}

\[
\begin{array}{p{0.30\linewidth}|p{0.20\linewidth}|p{0.39\linewidth}}
\textbf{gate}&\textbf{status after V49}&\textbf{content}\\ \hline
affine clock reduction&proved&
true Hamiltonian (49.4)\\
nonlinear clock branch&proved&
implicit derivative and lapse\\
quadratic clock orientation&proved&
two frequency sheets\\
physical-clock transfer bound&proved&
three-error inequality (49.24)\\
gain-free Einstein use&refuted&
vacuum scalar is lost\\
two-lag absolute calibration&refuted&
centering removes \(E_0\)\\
literal joint clock card&missing&
source items (49.29)\\
vacuum/counterterm consistency&after audit&
Einstein scalar row
\end{array}                                                 \tag{49.30}
\]

\subsection{Parent-integrity record}

The installed parent was
\[
\begin{array}{c|c}
 \text{file}&
 \texttt{Renewal\_SM\_Einstein\_Continuum\_Research\_Notes\_V48.tex}\\
 \text{lines}&19352\\
 \text{bytes}&632633\\
 \text{SHA--256}&
 \texttt{890A6E66754E83BDF5398A27CD73F3D51A7102A05FBE1364A745E60CF2744435}.
\end{array}                                                 \tag{49.31}
\]

% =============================================================================
% END APPENDED FOURTEENTH-CYCLE V49 MATERIAL
% =============================================================================
% =============================================================================
% BEGIN APPENDED FOURTEENTH-CYCLE V50 MATERIAL
% =============================================================================

\section{Fourteenth-cycle V50: compulsory companion audit and the
adaptive-atlas boundary}
\label{sec:v14v50}

\subsection{Files inspected}

The fiftieth-version audit was performed read-only on
\[
\begin{array}{p{0.51\linewidth}|r|r}
\textbf{file}&\textbf{lines}&\textbf{bytes}\\ \hline
\texttt{Renewal\_Yang\_Mills\_Mass\_Gap\_Research\_Notes\_V17.tex}
 &3763&138213\\
\texttt{Renewal\_NS\_Stability\_Research\_Notes\_V26.tex}
 &5319&278283.
\end{array}                                                  \tag{50.1}
\]
Their SHA--256 digests are
\[
\begin{aligned}
 &\texttt{739E0192D62B0660EF2C3C0D04C461EB02F9CC8B2DB8D731DBEB32CD01B070CE},
                                                               \tag{50.2}\\
 &\texttt{82C887F8C2C69E4F28FAD7C3DC8DE5B9099CB5A4BF431EBA1FFF3E91935978AD}.
                                                               \tag{50.3}
\end{aligned}
\]
No companion file was modified or removed.

\subsection{Yang--Mills V16--V17}

Yang--Mills V16 applies its centre-relative coefficientwise
calculus at all \(64\) quarter-turn anchors for each of two actions.
It obtains \(128\) exact local patches, then proves that their union
does not cover the transverse torus by the exact uncovered point
\((\pi/4,\pi/4,\pi/4)\).  V17 replaces the algebraic eighth-turn
phase by a nearby rational point of the unit circle, certifies
positive interior patches for both signs and both actions, and
covers that first explicit gap.  It correctly leaves the remaining
global cover open.

The transferable lesson is already represented in V40 and V45:
local strict positivity plus an exact uncovered-point oracle supports
an adaptive finite atlas.  The rational parametrization
\[
 z(r)=\frac{1-r^2+2ri}{1+r^2},\qquad r\in\mathbb Q,        \tag{50.4}
\]
may be useful whenever a source coefficient is Laurent-polynomial in
a unit phase.  It does not supply a clock normalization or an
Einstein energy.

None of the YM V16--V17 matrices acts on the reduced
SM--Einstein phase space.  Their exact floors are therefore not
inserted into the CMC/York, transfer, or vacuum trace cards.

\subsection{Navier--Stokes V26}

The Navier--Stokes note remains unchanged and concerns rank-one
Oseen-kernel derivative estimates.  It provides no physical clock
pair, transfer gain, York Jacobian, symplectic frequency sum, or
metric counterterm.

\subsection{Checkpoint decision}

The audit found no companion value for
\[
 \mu_T,\quad
 r_\alpha=\ell_\alpha/\tau_\alpha,\quad
 E_{0,\alpha},\quad
 \hbar_\alpha\sum_j\nu_{j,\alpha},\quad
 \gamma_M,\quad\gamma_H.                                  \tag{50.5}
\]
The next bottleneck therefore remains the scalar information erased
by centering and gap computations.

\begin{proposition}[Local-atlas progress does not determine a scalar
normalization]
\label{prop:v14v50-atlas-scalar}
Let \(A(x)\) be an operator family and let
\(\lambda(x)I\) be any scalar family.  Every local positivity
certificate for a derivative or commutator family invariant under
\(A\mapsto A+\lambda I\) is identical for \(A\) and
\(A+\lambda I\), while their absolute transfer generators differ by
\(\lambda I\).  Hence completing an atlas for such a quotient target
cannot determine the Einstein vacuum scalar.
\end{proposition}

\begin{proof}
Derivatives with scalar centering, commutators, and spectral gaps
annihilate scalar multiples of the identity.  Functional calculus
gives
\[
 e^{-h(A+\lambda I)}=e^{-h\lambda}e^{-hA},
\]
so the absolute logarithmic generator differs by \(\lambda I\).
The two data types are therefore independent.
\end{proof}

\begin{decision}[Direction after V50]
\label{dec:v14v50-next}
V51 will derive the finite vacuum-energy/counterterm consistency
compiler.  It will express the quadratic zero-point energy as a
symplectic trace, prove its perturbation bound, differentiate the
cell-volume-weighted vacuum functional, and give a summable
\(C^1\) cutoff-cocycle criterion sufficient for a limiting Einstein
source.  It will also prove that subtracting only a base numerical
energy does not cancel the metric derivative of the frequency sum.
\end{decision}

\subsection{Acquisition ledger}

\[
\begin{array}{p{0.30\linewidth}|p{0.20\linewidth}|p{0.39\linewidth}}
\textbf{gate}&\textbf{status after V50}&\textbf{content}\\ \hline
companion checkpoint&acquired&
YM V17 and NS V26 with hashes\\
quarter-anchor atlas&YM progress&
exact but globally incomplete\\
first rational interior patch&YM progress&
one explicit gap covered\\
adaptive exact atlas&structural import&
no numerical SM coefficients\\
clock/York data&absent&
no rows in (50.5)\\
absolute vacuum normalization&absent&
centred targets cannot recover it\\
vacuum/counterterm compiler&next&
metric \(C^1\) cofinal control
\end{array}                                                 \tag{50.6}
\]

The next compulsory companion audit is V55.

\subsection{Parent-integrity record}

The installed parent was
\[
\begin{array}{c|c}
 \text{file}&
 \texttt{Renewal\_SM\_Einstein\_Continuum\_Research\_Notes\_V49.tex}\\
 \text{lines}&19762\\
 \text{bytes}&645493\\
 \text{SHA--256}&
 \texttt{0A72E8C18910C8325F5023B56C31BA55B3D42CDEBCE2CEC19AE406260185B113}.
\end{array}                                                 \tag{50.7}
\]

% =============================================================================
% END APPENDED FOURTEENTH-CYCLE V50 MATERIAL
% =============================================================================
% =============================================================================
% BEGIN APPENDED FOURTEENTH-CYCLE V51 MATERIAL
% =============================================================================

\section{Fourteenth-cycle V51: symplectic vacuum traces,
metric first jets, and the counterterm cocycle}
\label{sec:v14v51}

\subsection{The exact quadratic vacuum trace}

Let \(G=G^{\mathsf T}\succ0\) be the \(2n\times2n\) reduced
quadratic matrix of V44, and put
\[
 M(G)=G^{1/2}\Omega G^{1/2}.                               \tag{51.1}
\]
Its singular values are
\[
 \nu_1,\nu_1,\ldots,\nu_n,\nu_n,                           \tag{51.2}
\]
where the \(\nu_j\) are the Williamson frequencies.

\begin{proposition}[Basis-free zero-point formula]
\label{prop:v14v51-trace}
The Weyl-to-normal vacuum shift is
\[
 \boxed{
 E_0(G,\hbar)
 =\frac\hbar2\sum_{j=1}^n\nu_j
 =\frac\hbar4\operatorname{Tr}|M(G)|.}                    \tag{51.3}
\]
It is independent of the chosen Williamson basis.
\end{proposition}

\begin{proof}
Theorem~\ref{thm:v14v44-williamson} identifies the oscillator
frequencies.  Each contributes \(\hbar\nu_j/2\) in Weyl ordering.
Equation (51.2) gives
\(\operatorname{Tr}|M|=2\sum_j\nu_j\), proving (51.3).
Singular values are invariant under orthogonal changes used in the
Williamson construction.
\end{proof}

\begin{theorem}[Vacuum-trace perturbation bound]
\label{thm:v14v51-perturb}
For two positive matrices \(G,\widetilde G\),
\[
 |E_0(G,\hbar)-E_0(\widetilde G,\hbar)|
 \le\frac\hbar4
 \|M(G)-M(\widetilde G)\|_1.                              \tag{51.4}
\]
If
\[
 aI\preceq G,\widetilde G\preceq bI,                      \tag{51.5}
\]
then
\[
 \boxed{
 |E_0(G,\hbar)-E_0(\widetilde G,\hbar)|
 \le
 \frac{\hbar n}{2}\sqrt{\frac ba}\,
 \|G-\widetilde G\|.}                                     \tag{51.6}
\]
\end{theorem}

\begin{proof}
Mirsky's singular-value inequality gives
\[
 \sum_{k=1}^{2n}
 |s_k(M)-s_k(\widetilde M)|
 \le\|M-\widetilde M\|_1.                                 \tag{51.7}
\]
Apply (51.3).  Under (51.5), the square-root perturbation bound and
the two-product expansion give
\[
 \|M(G)-M(\widetilde G)\|
 \le\sqrt{\frac ba}\,\|G-\widetilde G\|,                  \tag{51.8}
\]
as in V44.  Finally
\(\|A\|_1\le2n\|A\|\), which gives (51.6).
\end{proof}

A fixed-dimensional \(O(\hbar)\) statement becomes
\(O(\hbar n)\) when the regulator adds modes.  The multiplicity is
load bearing.

\subsection{The exact first derivative}

Because \(G\succ0\), \(M(G)\) is invertible.  Let
\[
 M=U|M|                                                    \tag{51.9}
\]
be its polar decomposition.  For a variation \(\dot G\), let \(X\)
be the unique solution of the Sylvester equation
\[
 G^{1/2}X+XG^{1/2}=\dot G.                                \tag{51.10}
\]
Then
\[
 \dot M
 =X\Omega G^{1/2}+G^{1/2}\Omega X.                        \tag{51.11}
\]

\begin{theorem}[Vacuum-frequency first jet]
\label{thm:v14v51-first-jet}
The zero-point functional is differentiable on the positive cone,
and
\[
 \boxed{
 DE_0(G,\hbar)[\dot G]
 =\frac\hbar4\operatorname{Re}
  \operatorname{Tr}(U^*\dot M).}                          \tag{51.12}
\]
Under \(aI\preceq G\preceq bI\),
\[
 \boxed{
 |DE_0(G,\hbar)[\dot G]|
 \le
 \frac{\hbar n}{2}\sqrt{\frac ba}\,\|\dot G\|.}           \tag{51.13}
\]
\end{theorem}

\begin{proof}
The derivative of the matrix square root is the solution of
(51.10), obtained by differentiating
\((G^{1/2})^2=G\).  This gives (51.11).

The trace norm is differentiable at every invertible square matrix.
Its derivative at \(M=U|M|\) is
\[
 D\operatorname{Tr}|M|[\dot M]
 =\operatorname{Re}\operatorname{Tr}(U^*\dot M).          \tag{51.14}
\]
One proof is to use the singular-value decomposition and
differentiate the sum of singular values; repeated positive singular
values cause no ambiguity because their spectral projector trace is
basis independent.  Equations (51.3) and (51.14) give (51.12).

Alternatively, divide (51.6) by the variation parameter and let it
tend to zero.  This yields (51.13).
\end{proof}

The metric derivative of the vacuum energy contains the derivative
of the physical Hessian \(G(q)\).  It is not determined by the
frequency values at one metric.

\subsection{Cell volume and the two vacuum-stress terms}

On a finite spatial regulator with cells \(c\), let
\[
 V_c(q)>0,\qquad G_c(q)\succ0
\]
be the physical cell volume and the reduced quadratic matrix.
Define the raw bosonic vacuum functional
\[
 {\mathscr V}_{0,X}(q)
 =\sum_{c}V_c(q)E_{0,c}(q).                               \tag{51.15}
\]

\begin{proposition}[Exact vacuum first variation]
\label{prop:v14v51-volume}
For every metric variation \(k\),
\[
 \boxed{
 D{\mathscr V}_{0,X}(q)[k]
 =
 \sum_c
 \left(
 E_{0,c}(q)\,DV_c(q)[k]
 +V_c(q)\,DE_{0,c}(q)[k]
 \right).}                                                 \tag{51.16}
\]
The first term is the volume or cosmological contribution.  The
second is the frequency-polarization stress.
\end{proposition}

\begin{proof}
Apply the product rule to every summand in (51.15).
\end{proof}

\begin{countertheorem}[Subtracting a base energy does not renormalize its
metric first jet]
\label{cth:v14v51-base-subtraction}
Fix \(q_0\).  The numerical subtraction
\[
 \widetilde{\mathscr V}(q)
 ={\mathscr V}_{0,X}(q)-{\mathscr V}_{0,X}(q_0)            \tag{51.17}
\]
satisfies \(\widetilde{\mathscr V}(q_0)=0\), but
\[
 D\widetilde{\mathscr V}(q_0)
 =D{\mathscr V}_{0,X}(q_0).                               \tag{51.18}
\]
The covariant volume subtraction
\[
 \widehat{\mathscr V}(q)
 =\sum_cV_c(q)
 \bigl(E_{0,c}(q)-E_{0,c}(q_0)\bigr)                       \tag{51.19}
\]
removes the volume term at \(q_0\), but leaves
\[
 D\widehat{\mathscr V}(q_0)[k]
 =\sum_cV_c(q_0)DE_{0,c}(q_0)[k].                         \tag{51.20}
\]
\end{countertheorem}

\begin{proof}
Equation (51.18) follows because the subtracted number is independent
of \(q\).  Differentiate (51.19) and use the vanishing of each
parenthesis at \(q_0\); the result is (51.20).
\end{proof}

To cancel the vacuum source at an anchor, a counterterm must match
the required metric first jet, not merely the energy value.

\subsection{A cofinal \(C^1\) counterterm theorem}

Let \(K\) be a compact convex metric screen in a fixed
finite-dimensional coordinate chart.  Transport every cutoff
functional to \(K\), and let
\[
 {\mathscr R}_X(q)
 ={\mathscr V}_{0,X}(q)+{\mathscr C}_X(q)                  \tag{51.21}
\]
be the raw vacuum functional plus the declared counterterm.  Define
the adjacent first-jet defect
\[
 \epsilon_X^{(1)}
 =
 \sup_{q\in K}|{\mathscr R}_{X+1}(q)-{\mathscr R}_X(q)|
 +\sup_{q\in K}
 \|D{\mathscr R}_{X+1}(q)-D{\mathscr R}_X(q)\|.            \tag{51.22}
\]

\begin{theorem}[Summable vacuum-counterterm cocycle]
\label{thm:v14v51-cocycle}
If
\[
 \sum_{X=1}^{\infty}\epsilon_X^{(1)}<\infty,               \tag{51.23}
\]
then there is a \(C^1\) functional
\({\mathscr R}_\infty\) on \(K\) such that
\[
 {\mathscr R}_X\longrightarrow{\mathscr R}_\infty,\qquad
 D{\mathscr R}_X\longrightarrow D{\mathscr R}_\infty       \tag{51.24}
\]
uniformly on \(K\).  In particular, the renormalized vacuum
contribution to the finite Einstein residual has a unique limiting
first variation on that screen.
\end{theorem}

\begin{proof}
The first sum in (51.22)--(51.23) makes
\(({\mathscr R}_X)\) uniformly Cauchy.  The second makes
\((D{\mathscr R}_X)\) uniformly Cauchy.  Let their uniform limits be
\({\mathscr R}_\infty\) and \(A_\infty\).

For \(q,q+v\in K\), convexity and the fundamental theorem of calculus
give
\[
 {\mathscr R}_X(q+v)-{\mathscr R}_X(q)
 =\int_0^1D{\mathscr R}_X(q+tv)[v]\,dt.                   \tag{51.25}
\]
Pass uniformly to the limit.  The right side becomes
\(\int_0^1A_\infty(q+tv)[v]dt\), so
\({\mathscr R}_\infty\) is differentiable with derivative
\(A_\infty\).  This proves (51.24).
\end{proof}

A merely convergent sequence of vacuum numbers does not imply
convergence of the stress.  The second term in (51.22) is
independent.

\subsection{A usable trace budget}

Suppose on \(K\)
\[
 a_XI\preceq G_{c,X}(q)\preceq b_XI,\qquad
 \|D_qG_{c,X}\|\le L_{G,X},                               \tag{51.26}
\]
and
\[
 0<V_{c,X}\le V_X,\qquad
 \|D_qV_{c,X}\|\le L_{V,X}.                               \tag{51.27}
\]
If each cell has \(n_X\) oscillator modes, then
\[
 E_{0,c,X}\le\frac{\hbar_Xn_Xb_X}{2},                     \tag{51.28}
\]
and Theorem~\ref{thm:v14v51-first-jet} gives
\[
 \|DE_{0,c,X}\|
 \le\frac{\hbar_Xn_X}{2}
 \sqrt{\frac{b_X}{a_X}}\,L_{G,X}.                         \tag{51.29}
\]
Hence one raw first-jet budget is
\[
 \boxed{
 \|D{\mathscr V}_{0,X}\|
 \le
 \frac{\hbar_Xn_X}{2}
 \sum_c
 \left(
 b_XL_{V,X}
 +V_X\sqrt{\frac{b_X}{a_X}}\,L_{G,X}
 \right).}                                                 \tag{51.30}
\]
A claimed vanishing vacuum source must make the right side tend to
zero or cancel it by a declared counterterm with a controlled
residual.  Bounded frequencies alone do not suffice when the number
of modes or cells grows.

\subsection{Finite Ward consistency}

Let a finite relabeling or gauge group act on the metric screen.
Suppose both \({\mathscr V}_{0,X}\) and
\({\mathscr C}_X\) are invariant:
\[
 {\mathscr R}_X(gq)={\mathscr R}_X(q).                     \tag{51.31}
\]
Then every infinitesimal generator \(\mathcal L_\xi q\) satisfies
\[
 D{\mathscr R}_X(q)[\mathcal L_\xi q]=0.                  \tag{51.32}
\]
Under Theorem~\ref{thm:v14v51-cocycle}, this identity passes to the
limit.  Thus a covariant counterterm preserves the finite
diffeomorphism Ward row.  A subtraction performed only after taking
a spectrum, with no common metric functional, has no such guarantee.

\subsection{Renormalization ambiguity}

\begin{countertheorem}[Centred transfer data do not fix the cosmological
counterterm]
\label{cth:v14v51-renormalization}
Let \({\mathscr R}_X\) satisfy all centered transfer, spectral-gap,
commutator, and gauge-covariance tests.  For any constant \(c\), the
modified functional
\[
 {\mathscr R}_X^{(c)}(q)
 ={\mathscr R}_X(q)+c\sum_cV_c(q)                         \tag{51.33}
\]
has the same centered transfer generators, excitation gaps, and
commutators, but its metric first variation differs by
\[
 c\sum_cDV_c(q).                                          \tag{51.34}
\]
\end{countertheorem}

\begin{proof}
The added term is a scalar energy on each fixed metric carrier, so
centering, eigenvalue differences, and commutators remove it.
Differentiating the common metric functional gives (51.34).
\end{proof}

At least one absolute gravitational renormalization condition is
therefore required.  With more allowed local covariant counterterms,
a determining family of metric and curvature tests is required to
fix all of their coefficients.

\subsection{The complete vacuum card}

For bosons, fermions, ghosts, and declared regulator fields, the
total vacuum card must record their statistics, multiplicities,
boundary conditions, and ordering convention before any cancellation
is taken.  Its finite form is
\[
 \mathfrak V_X^{\rm Ein}=
 \bigl(
 [G_{s,X}(q)], [\nu_{s,j,X}(q)],
 [\epsilon_s], [n_{s,X}],
 [V_{c,X}(q)],
 [{\mathscr C}_X(q)],
 [{\mathscr R}_X(q)],
 \epsilon_X^{(1)},
 [\hbox{renormalization conditions}]
 \bigr),                                                   \tag{51.35}
\]
where \(\epsilon_s\) denotes the declared signed/statistical
contribution.  A cancellation between sectors is valid only on one
common metric, boundary, clock, and cutoff card.

The attached and Grand Tensor sources permit counterterm
contributions in the common action and demand their derivative
control, but the inspected passages do not populate (51.35) for the
active SM--Einstein regulator.

\begin{decision}[Direction after V51]
\label{dec:v14v51-next}
V52 will derive the Ward-consistent first-jet matching theorem for
the complete coupled action.  It will separate cosmological,
Einstein--Hilbert, matter-parameter, and genuinely new metric-stress
directions by a finite determining Gram, and prove the exact residual
that remains invariant under allowed counterterm reparametrization.
This is needed before the V46 bosonic Hessian and V48 constraint
chart can be placed in one renormalized Einstein equation.
\end{decision}

\subsection{Acquisition ledger}

\[
\begin{array}{p{0.30\linewidth}|p{0.20\linewidth}|p{0.39\linewidth}}
\textbf{gate}&\textbf{status after V51}&\textbf{content}\\ \hline
zero-point trace&proved&
basis-free formula (51.3)\\
trace perturbation&proved&
dimension-sensitive bound\\
metric first jet&proved&
polar/Sylvester formula\\
base-value subtraction&refuted&
stress remains\\
cofinal counterterm transport&proved&
summable \(C^1\) cocycle\\
raw trace budget&proved&
mode, cell, and metric factors\\
centred cosmological inference&refuted&
finite counterexample family\\
literal total vacuum card&missing&
sector and counterterm data\\
Ward-consistent source short&next&
determining metric Gram
\end{array}                                                 \tag{51.36}
\]

\subsection{Parent-integrity record}

The installed parent was
\[
\begin{array}{c|c}
 \text{file}&
 \texttt{Renewal\_SM\_Einstein\_Continuum\_Research\_Notes\_V50.tex}\\
 \text{lines}&19914\\
 \text{bytes}&651148\\
 \text{SHA--256}&
 \texttt{9A6261ECB27FAA4A86DEF8345AF16E72AEC1B8B0F832676FFC7541533EB2ED17}.
\end{array}                                                 \tag{51.37}
\]

% =============================================================================
% END APPENDED FOURTEENTH-CYCLE V51 MATERIAL
% =============================================================================
% =============================================================================
% BEGIN APPENDED FOURTEENTH-CYCLE V52 MATERIAL
% =============================================================================

\section{Fourteenth-cycle V52: Ward-compatible counterterm quotients
and the determining Einstein first-jet Gram}
\label{sec:v14v52}

\subsection{Finite first-variation geometry}

Let \(E\) be the finite real space of represented metric variations
on one screen, and let
\[
 G_E:E\longrightarrow E^*
\]
be a faithful positive metric.  Covectors carry the dual norm
\[
 \|r\|_{G_E^{-1}}^2
 =\langle r,G_E^{-1}r\rangle.                             \tag{52.1}
\]
Let
\[
 {\mathscr I}_1,\ldots,{\mathscr I}_m
\]
be the predeclared finite covariant counterterm functionals.  Their
first-jet synthesis at the base metric is
\[
 C:\mathbb R^m\longrightarrow E^*,\qquad
 Ca=\sum_{j=1}^ma_jD{\mathscr I}_j.                        \tag{52.2}
\]
This list may include the finite cosmological volume, the
Einstein--Hilbert/Palatini coefficient, allowed matter-parameter
renormalizations, and admitted boundary terms.  It may not be
enlarged after reading the residual.

Define
\[
 K_C=C^*G_E^{-1}C,\qquad
 P_C=CK_C^\dagger C^*G_E^{-1}.                            \tag{52.3}
\]

\begin{theorem}[Source-minimal counterterm residual]
\label{thm:v14v52-short}
For every raw first variation \(r\in E^*\),
\[
 r_{\rm irr}:=(I-P_C)r                                    \tag{52.4}
\]
is the unique component orthogonal to all allowed counterterm
directions.  Its squared norm is
\[
 \boxed{
 \|r_{\rm irr}\|_{G_E^{-1}}^2
 =
 \langle r,G_E^{-1}r\rangle
 -
 b^*K_C^\dagger b,\qquad
 b=C^*G_E^{-1}r.}                                         \tag{52.5}
\]
The minimum-norm coefficient that removes the allowed component is
\[
 a_*=-K_C^\dagger b,                                      \tag{52.6}
\]
and
\[
 \inf_a\|r+Ca\|_{G_E^{-1}}
 =\|r_{\rm irr}\|_{G_E^{-1}}.                             \tag{52.7}
\]
Moreover,
\[
 (I-P_C)(r+Ca)=r_{\rm irr}                                \tag{52.8}
\]
for every \(a\).  Thus \(r_{\rm irr}\) is invariant under allowed
counterterm reparametrization.
\end{theorem}

\begin{proof}
The operator \(P_C\) is the orthogonal projection onto
\(\operatorname{Ran}C\) in the \(G_E^{-1}\) inner product.  Indeed,
the Moore--Penrose identities give \(P_C^2=P_C\), its range is
\(\operatorname{Ran}C\), and
\[
 G_E^{-1}P_C=P_C^*G_E^{-1}.                               \tag{52.9}
\]
Pythagoras gives (52.5).  Expanding
\[
 \|r+Ca\|_{G_E^{-1}}^2
 =\|r\|_{G_E^{-1}}^2+2\langle a,b\rangle+\langle a,K_Ca\rangle
\]
and completing the supported square gives (52.6)--(52.7).
Finally \(P_CC=C\), which proves (52.8).
\end{proof}

The residual in (52.5) is an obstruction only relative to the
declared counterterm class.  Enlarging that class changes the
physical theory or its renormalization prescription.

\subsection{Ward identities survive exactly when the counterterms do}

Let
\[
 L:\mathfrak d\longrightarrow E                           \tag{52.10}
\]
map an infinitesimal represented relabeling to its metric variation.
A first variation \(r\) is Ward compatible when
\[
 L^*r=0.                                                   \tag{52.11}
\]

\begin{theorem}[Ward-compatible counterterm short]
\label{thm:v14v52-ward}
Suppose
\[
 L^*C=0.                                                   \tag{52.12}
\]
Then
\[
 L^*r_{\rm irr}=L^*r.                                     \tag{52.13}
\]
In particular:

\begin{enumerate}
\item if \(r\) is Ward compatible, so are its removable and
irreducible components;
\item if \(L^*r\ne0\), no allowed covariant counterterm can remove
that anomaly.
\end{enumerate}
\end{theorem}

\begin{proof}
Every vector in \(\operatorname{Ran}P_C=\operatorname{Ran}C\)
is killed by \(L^*\), by (52.12).  Hence
\[
 L^*(I-P_C)r=L^*r.
\]
The two conclusions follow.
\end{proof}

Thus gauge or diffeomorphism failure must be reported before a
counterterm fit.  A noncovariant counterterm can cancel a numerical
residual while violating the finite Ward identity.

\subsection{Cosmological and Einstein directions}

On a common finite coframe/metric action, let
\[
 {\mathscr I}_0(q)=\operatorname{Vol}(q),\qquad
 {\mathscr I}_1(q)={\mathscr S}_{\rm Palatini/EH}(q)       \tag{52.14}
\]
with the boundary completion fixed.  The first two counterterm
columns are
\[
 c_0=D{\mathscr I}_0,\qquad c_1=D{\mathscr I}_1.           \tag{52.15}
\]
Their \(2\times2\) identifying Gram is
\[
 K_{01}=
 \begin{pmatrix}
 \langle c_0,G_E^{-1}c_0\rangle&
 \langle c_0,G_E^{-1}c_1\rangle\\
 \langle c_1,G_E^{-1}c_0\rangle&
 \langle c_1,G_E^{-1}c_1\rangle
 \end{pmatrix}.                                           \tag{52.16}
\]

\begin{proposition}[Finite identifiability of Newton and cosmological
renormalizations]
\label{prop:v14v52-identify}
The two coefficients multiplying (52.14) are separately determined
by the first-variation screen exactly when
\[
 \det K_{01}>0.                                            \tag{52.17}
\]
If \(K_{01}\succeq\kappa I_2\), perturbing the measured first
variation by \(e\) perturbs the fitted two coefficients by at most
\[
 \kappa^{-1/2}\|e\|_{G_E^{-1}}.                           \tag{52.18}
\]
\end{proposition}

\begin{proof}
The Gram is nonsingular exactly when \(c_0,c_1\) are linearly
independent.  In that case the coefficient map has least singular
value \(\sqrt{\lambda_{\min}(K_{01})}\).  Inversion gives (52.18).
\end{proof}

A metric test panel on which \(c_0\) and \(c_1\) become collinear
cannot distinguish a cosmological shift from a gravitational
coupling shift.

\subsection{The determining physical panel}

Let
\[
 {\cal R}
 =(\operatorname{Ran}C)^{\perp_{G_E^{-1}}}
 \cap\ker L^*                                               \tag{52.19}
\]
be the Ward-compatible irreducible covector space.  Let
\[
 A_{\rm det}:{\cal R}\longrightarrow\mathbb R^N           \tag{52.20}
\]
record the first variation on a finite family of protected metric
tests.

\begin{theorem}[Determining-panel certification]
\label{thm:v14v52-determining}
If
\[
 A_{\rm det}^*A_{\rm det}
 \succeq c_{\rm det}I_{\cal R},\qquad c_{\rm det}>0,        \tag{52.21}
\]
then every Ward-compatible raw residual satisfies
\[
 \boxed{
 \|r_{\rm irr}\|_{G_E^{-1}}
 \le c_{\rm det}^{-1/2}
 \|A_{\rm det}r_{\rm irr}\|.}                             \tag{52.22}
\]
In particular, vanishing of the protected determining responses
implies \(r_{\rm irr}=0\).
\end{theorem}

\begin{proof}
Apply (52.21) to \(r_{\rm irr}\in{\cal R}\):
\[
 \|A_{\rm det}r_{\rm irr}\|^2
 \ge c_{\rm det}\|r_{\rm irr}\|_{G_E^{-1}}^2.
\]
Take square roots.
\end{proof}

\begin{countertheorem}[A nondetermining panel can miss an Einstein source]
\label{cth:v14v52-nondetermining}
Let \(E^*=\mathbb R^2\), let there be no counterterm or Ward
direction, and let
\[
 A_{\rm det}(x,y)=x.
\]
Then the nonzero residual \(r=(0,1)\) has zero measured response.
\end{countertheorem}

\begin{proof}
Direct substitution gives \(A_{\rm det}r=0\).
\end{proof}

A list of familiar metric variations is not determining until its
lower frame bound on (52.19) is proved.

\subsection{Pointwise fitting is not a common action}

The short in Theorem~\ref{thm:v14v52-short} is performed at one
base metric.  Applying it independently at every metric can produce
a nonintegrable covector field.

\begin{countertheorem}[Zero pointwise residual with nonzero action curl]
\label{cth:v14v52-curl}
On \(K=\mathbb R^2\) with coordinates \((x,y)\), let the permitted
pointwise covector span be all of \(T^*K\).  The field
\[
 \alpha=-y\,dx+x\,dy                                      \tag{52.23}
\]
has zero pointwise quotient residual, but
\[
 d\alpha=2\,dx\wedge dy\ne0.                              \tag{52.24}
\]
Hence no scalar counterterm functional has \(D{\mathscr C}=\alpha\)
on a simply connected neighbourhood.
\end{countertheorem}

\begin{proof}
The full pointwise span makes the algebraic residual zero.  Direct
exterior differentiation gives (52.24).  Exact one-forms are
closed, so \(\alpha\) is not a functional derivative.
\end{proof}

The admissible construction is to choose actual scalar functionals
\({\mathscr I}_j(q)\) and cutoff-dependent but metric-independent
coefficients \(a_{j,X}\).  Then
\[
 {\mathscr C}_X(q)=\sum_ja_{j,X}{\mathscr I}_j(q)           \tag{52.25}
\]
is exact automatically.  If the coefficients depend on \(q\), the
full derivative contains
\[
 D{\mathscr C}_X
 =\sum_ja_jD{\mathscr I}_j+\sum_j{\mathscr I}_jDa_j,       \tag{52.26}
\]
so retaining only the first sum is not a common-action derivative.
The finite command plaquettes and periods in the attached V30 source
test precisely this integrability issue.

\subsection{First jets do not determine fluctuation Hessians}

\begin{countertheorem}[Invisible first-jet counterterm with arbitrary Hessian]
\label{cth:v14v52-second-jet}
At the base point \(x=0\), the one-dimensional counterterm
\[
 {\mathscr C}_a(x)=a x^2                                  \tag{52.27}
\]
has
\[
 {\mathscr C}_a(0)=0,\qquad D{\mathscr C}_a(0)=0,\qquad
 D^2{\mathscr C}_a(0)=2a.                                 \tag{52.28}
\]
Thus every vacuum value and Einstein first-variation test at the
anchor is identical for all \(a\), while the quadratic fluctuation
Hessian is arbitrary.
\end{countertheorem}

\begin{proof}
Differentiate (52.27).
\end{proof}

Consequently V51--V52 can close the renormalized Einstein first
variation without yet populating the positive physical Hessian
needed by V44.  The allowed counterterm two-jets must be recorded on
the same background.

\subsection{Cofinal transport of the quotient residual}

Let all finite screens be transported to one \(E^*\).  Suppose
\(G_{E,X}^{-1}\to G_E^{-1}\), \(C_X\to C\), and the positive
singular values of \(C_X\) have a common lower bound on one stable
support.  Then Moore--Penrose perturbation on that support gives
\[
 P_{C_X}\longrightarrow P_C.                              \tag{52.29}
\]
If also \(r_X\to r\), then
\[
 (I-P_{C_X})r_X\longrightarrow(I-P_C)r.                   \tag{52.30}
\]
A rank transition in the allowed counterterm bank is therefore a
separate support event; the quotient residual must not be transported
through it with an uncontrolled pseudoinverse.

\subsection{The renormalized first-jet card}

The complete finite card is
\[
 \mathfrak R_X^{(1)}=
 \bigl(
 [E_X,G_{E,X}], [L_X],
 [{\mathscr I}_{j,X}], [C_X], [K_{C,X}],
 [r_X], [r_{{\rm irr},X}],
 [A_{{\rm det},X}],c_{{\rm det},X},
 [\hbox{plaquettes/periods}],
 [\hbox{renormalization conditions}]
 \bigr).                                                   \tag{52.31}
\]
The Grand Tensor supplies the abstract finite Einstein residual and
a determining-panel hypothesis, while the attached source supplies
common-action plaquette machinery.  The active regulator does not
populate the matrices and counterterm jets in (52.31).

\begin{decision}[Direction after V52]
\label{dec:v14v52-next}
V53 will add the second-jet layer.  It will derive the exact
counterterm action on the coupled fluctuation Hessian, quotient that
Hessian by gauge and constraint directions, and prove a
simultaneous first-jet/stability acceptance theorem.  The target is
to prevent a counterterm that repairs the Einstein equation from
silently destroying the V46 bosonic quotient floor or the V44
Williamson positivity needed for the transfer Hamiltonian.
\end{decision}

\subsection{Acquisition ledger}

\[
\begin{array}{p{0.30\linewidth}|p{0.20\linewidth}|p{0.39\linewidth}}
\textbf{gate}&\textbf{status after V52}&\textbf{content}\\ \hline
allowed counterterm quotient&proved&
source-minimal residual\\
Ward compatibility&proved&
anomaly cannot be fitted away\\
Newton/cosmological identifiability&proved&
two-column Gram\\
determining metric panel&proved&
lower-frame criterion\\
pointwise action fit&refuted&
curl counterexample\\
first-jet Hessian inference&refuted&
quadratic counterterm\\
cofinal support transport&conditional&
stable counterterm rank\\
literal first-jet card&missing&
source matrices unpopulated\\
coupled second-jet stability&next&
counterterm Hessian firewall
\end{array}                                                 \tag{52.32}
\]

\subsection{Parent-integrity record}

The installed parent was
\[
\begin{array}{c|c}
 \text{file}&
 \texttt{Renewal\_SM\_Einstein\_Continuum\_Research\_Notes\_V51.tex}\\
 \text{lines}&20353\\
 \text{bytes}&664779\\
 \text{SHA--256}&
 \texttt{4732C82877A2625614363436D8A2F34B4B10A793B9706B7A928AE84CFD50A696}.
\end{array}                                                 \tag{52.33}
\]

% =============================================================================
% END APPENDED FOURTEENTH-CYCLE V52 MATERIAL
% =============================================================================
% =============================================================================
% BEGIN APPENDED FOURTEENTH-CYCLE V53 MATERIAL
% =============================================================================

\section{Fourteenth-cycle V53: simultaneous counterterm first-jet
matching and physical Hessian stability}
\label{sec:v14v53}

\subsection{Assemble before reducing}

Let \(S_a\) be the finite total action with predeclared counterterm
coefficients \(a=(a_1,\ldots,a_m)\):
\[
 S_a=S_{\rm raw}+\sum_{j=1}^ma_j{\mathscr I}_j.            \tag{53.1}
\]
Let \(F=0\) be the complete nonlinear constraint surface and let
\(\iota\) be a local constraint chart with tangent lift
\(J_{\rm tan}=D\iota\).  At a constrained stationary point, choose
\(\lambda_a\) so that
\[
 DS_a+DF^*\lambda_a=0.                                    \tag{53.2}
\]
The physical constrained Hessian is
\[
 \boxed{
 {\mathbb H}_{\rm phys}(a)
 =
 J_{\rm tan}^*
 \left(
 D^2S_a+D^2\langle\lambda_a,F\rangle
 \right)
 J_{\rm tan}.}                                             \tag{53.3}
\]
This is Theorem~\ref{thm:v14v48-constrained-hessian} applied after
the complete action has been assembled.

\begin{proposition}[Individual invariant Hessians need not have gauge
kernels off shell]
\label{prop:v14v53-assemble}
Let a Lie group act by \(z\mapsto g z\), with infinitesimal vector
field \(R_\xi(z)\), and let \(I\) be invariant.  Then
\[
 D^2I(z)[R_\xi(z),v]
 =-DI(z)[D_vR_\xi(z)].                                    \tag{53.4}
\]
Thus \(D^2I(z)R_\xi(z)=0\) follows at a stationary point of \(I\),
but need not hold when \(DI(z)\ne0\).
\end{proposition}

\begin{proof}
Differentiate the infinitesimal invariance identity
\[
 DI(z)[R_\xi(z)]=0
\]
in the direction \(v\).  This gives (53.4).
\end{proof}

A counterterm Hessian may therefore fail to descend by itself even
though the complete stationary action descends.  One must form
(53.3), including the constraint multiplier curvature, before taking
the physical quotient.

\subsection{Affine second-jet model on a fixed chart}

On a fixed base point, constraint chart, and stable multiplier
support, the stationarity equations are linear in the action
coefficients.  When they admit an affine solution
\[
 \lambda_a=\lambda_0+\sum_ja_j\lambda_j,                  \tag{53.5}
\]
equation (53.3) takes the finite affine form
\[
 {\mathbb H}_{\rm phys}(a)
 =H_0+\sum_{j=1}^ma_jB_j,                                 \tag{53.6}
\]
where
\[
\begin{aligned}
 H_0&=
 J_{\rm tan}^*
 \left(D^2S_{\rm raw}
       +D^2\langle\lambda_0,F\rangle\right)J_{\rm tan},\\
 B_j&=
 J_{\rm tan}^*
 \left(D^2{\mathscr I}_j
       +D^2\langle\lambda_j,F\rangle\right)J_{\rm tan}.
                                                               \tag{53.7}
\end{aligned}
\]
If the background or constraint chart changes with \(a\), its
transport and Taylor remainder must be added; (53.6) is not then an
exact global parameterization.

\subsection{The first-jet solution fibre}

Let \(r\in E^*\) and \(C\) be the raw first variation and counterterm
first-jet synthesis of V52.  Exact first-jet cancellation requires
\[
 r+Ca=0.                                                   \tag{53.8}
\]

\begin{theorem}[First-jet solutions and second-jet ambiguity]
\label{thm:v14v53-fibre}
Equation (53.8) is solvable exactly when
\[
 (I-P_C)r=0.                                               \tag{53.9}
\]
When it is solvable, every solution is
\[
 a=a_0+z,\qquad
 a_0=-K_C^\dagger C^*G_E^{-1}r,\qquad z\in\ker C.          \tag{53.10}
\]
All such solutions give the same physical Hessian if and only if
\[
 \boxed{
 \ker C\subset
 \ker{\mathcal B},\qquad
 {\mathcal B}z:=\sum_jz_jB_j.}                            \tag{53.11}
\]
\end{theorem}

\begin{proof}
Theorem~\ref{thm:v14v52-short} gives (53.9) and the
minimum-norm solution \(a_0\).  The complete affine solution set of
a consistent linear equation is \(a_0+\ker C\).  By (53.6), two
solutions separated by \(z\in\ker C\) have Hessian difference
\({\mathcal B}z\).  This difference vanishes for every pair exactly
under (53.11).
\end{proof}

\begin{countertheorem}[Einstein first-jet matching can leave an arbitrary
physical floor]
\label{cth:v14v53-arbitrary-floor}
Take a one-dimensional physical fluctuation space, two counterterm
coefficients, and
\[
 C=(1\ \ 0),\qquad r=-1,\qquad
 H_0=0,\qquad B_1=0,\qquad B_2=1.                         \tag{53.12}
\]
Every coefficient pair \(a=(1,t)\) solves the same first-jet
equation, but
\[
 {\mathbb H}_{\rm phys}(1,t)=t.                           \tag{53.13}
\]
\end{countertheorem}

\begin{proof}
Substitution in (53.8) and (53.6) gives the claims.
\end{proof}

A minimum-norm choice \(t=0\) is a convention; it is not a stability
theorem.

\subsection{Simultaneous finite acceptance as a semidefinite problem}

Let \({\cal A}\subset\mathbb R^m\) be the predeclared closed convex
set allowed by the renormalization conditions.  For a requested
physical floor \(\delta\), define
\[
 {\cal F}_\delta=
 \left\{
 a\in{\cal A}:
 r+Ca=0,\quad
 H_0+\sum_ja_jB_j\succeq\delta I
 \right\}.                                                 \tag{53.14}
\]

\begin{theorem}[First-jet/stability acceptance theorem]
\label{thm:v14v53-sdp}
The set \({\cal F}_\delta\) is convex.  A counterterm prescription
simultaneously satisfies the finite Einstein first variation and the
physical quadratic stability floor \(\delta\) exactly when
\[
 {\cal F}_\delta\ne\varnothing.                            \tag{53.15}
\]
If \({\cal A}\) is compact, the optimal certified floor
\[
 \delta_*=
 \max_{\substack{a\in{\cal A}\\r+Ca=0}}
 \lambda_{\min}
 \left(H_0+\sum_ja_jB_j\right)                            \tag{53.16}
\]
is attained.
\end{theorem}

\begin{proof}
The affine equation and \({\cal A}\) are convex constraints.  The
positive-semidefinite cone is convex and its inverse image under the
affine map (53.6) is convex.  This proves the first assertion;
(53.15) is the definition of simultaneous feasibility.
On a compact feasible coefficient set, the smallest eigenvalue is a
continuous function, so the maximum in (53.16) is attained.
\end{proof}

This is a genuine finite semidefinite feasibility problem.  For
rational source matrices, a proposed rational coefficient and floor
can be checked by exact congruence or an exact \(LDL^*\)
certificate.  Floating feasibility alone is not the installed
source card.

\subsection{A conservative analytic floor}

\begin{corollary}[Counterterm Hessian debit]
\label{cor:v14v53-debit}
Suppose
\[
 H_0\succeq\delta_0I,\qquad \|B_j\|\le b_j.                \tag{53.17}
\]
Then every first-jet solution \(a\) obeys
\[
 \boxed{
 {\mathbb H}_{\rm phys}(a)
 \succeq
 \left(\delta_0-\sum_j|a_j|b_j\right)I.}                  \tag{53.18}
\]
Hence
\[
 \sum_j|a_j|b_j<\delta_0                                  \tag{53.19}
\]
is a sufficient simultaneous stability condition.
\end{corollary}

\begin{proof}
The perturbation in (53.6) has norm at most
\(\sum_j|a_j|b_j\).  Weyl's inequality gives (53.18).
\end{proof}

The bound is conservative because it ignores cancellation between
the \(B_j\).  The semidefinite test (53.14) is exact.

\subsection{Robust intervals and measured matrices}

Suppose submitted matrices obey
\[
 \|\widehat H_0-H_0\|\le\epsilon_0,\qquad
 \|\widehat B_j-B_j\|\le\epsilon_j,                        \tag{53.20}
\]
and a candidate coefficient \(a\) satisfies
\[
 \widehat H_0+\sum_ja_j\widehat B_j
 \succeq\widehat\delta I.                                 \tag{53.21}
\]

\begin{proposition}[Outward stability certificate]
\label{prop:v14v53-outward}
The exact physical Hessian satisfies
\[
 \boxed{
 {\mathbb H}_{\rm phys}(a)
 \succeq
 \left(
 \widehat\delta-\epsilon_0-\sum_j|a_j|\epsilon_j
 \right)I.}                                                \tag{53.22}
\]
\end{proposition}

\begin{proof}
The norm of the difference between the exact and submitted affine
Hessians is at most
\(\epsilon_0+\sum_j|a_j|\epsilon_j\).  Apply Weyl's
inequality.
\end{proof}

The first-jet equation requires its own outward residual and
determining-panel bound.  A common coefficient vector must pass both
the first- and second-jet intervals.

\subsection{Cofinal \(C^2\) transport}

Transport every constrained physical Hessian to one fixed screen.
Let
\[
 \epsilon_X^{(2)}
 =\|{\mathbb H}_{{\rm phys},X+1}
       -{\mathbb H}_{{\rm phys},X}\|.                     \tag{53.23}
\]

\begin{theorem}[Summable second-jet stability]
\label{thm:v14v53-c2}
If
\[
 {\mathbb H}_{{\rm phys},X_0}\succeq\delta_0I,\qquad
 \sum_{X=X_0}^{\infty}\epsilon_X^{(2)}<\delta_0,           \tag{53.24}
\]
then the Hessians converge in operator norm to
\({\mathbb H}_{{\rm phys},\infty}\), and
\[
 \boxed{
 {\mathbb H}_{{\rm phys},\infty}
 \succeq
 \left(
 \delta_0-\sum_{X=X_0}^{\infty}\epsilon_X^{(2)}
 \right)I\succ0.}                                         \tag{53.25}
\]
\end{theorem}

\begin{proof}
Summability makes the operator sequence Cauchy.  Telescoping and
Weyl's inequality give, for every \(Y>X_0\),
\[
 {\mathbb H}_{{\rm phys},Y}
 \succeq
 \left(
 \delta_0-\sum_{X=X_0}^{Y-1}\epsilon_X^{(2)}
 \right)I.
\]
Pass to the operator-norm limit.
\end{proof}

Combined with the \(C^1\) cocycle in V51, this is the finite
\(C^2\) requirement needed to transport both the Einstein equation
and the Gaussian fluctuation Hamiltonian.

\subsection{Constraint and projector stability}

The affine matrices in (53.6) live on the physical quotient.  If the
constraint rank changes, the tangent lift \(J_{\rm tan}\) and
physical screen change.  Before comparing adjacent Hessians one must
supply Kato or polar transport of the constraint projectors and a
positive singular-value floor for the constraint Jacobian.  A small
ambient Hessian difference does not control (53.23) across an
unresolved rank jump.

Likewise the positive bosonic quotient floor (46.27) is only one
diagonal block of the full coupled Hessian.  Gravity, fermion,
Yukawa, ghost, boundary, and counterterm blocks and every cross block
must be assembled before the physical compression.  The earlier
coupled Schur-floor theorem then applies; no diagonal floor licenses
discarding a cross term.

\subsection{The simultaneous jet card}

The finite acquisition packet is
\[
 \mathfrak R_X^{(2)}=
 \bigl(
 [F_X],[J_{{\rm tan},X}],[\lambda_X],
 [r_X],[C_X],
 [H_{0,X}],[B_{j,X}],
 [{\cal A}_X],
 [a_X],\delta_X,
 [\epsilon_X^{(1)}],[\epsilon_X^{(2)}],
 [\hbox{exact PSD certificate}]
 \bigr).                                                   \tag{53.26}
\]
The active sources do not populate this packet.  In particular,
they allow counterterms in the common action but do not give their
complete first and second metric jets on the selected regulator.

\begin{decision}[Direction after V53]
\label{dec:v14v53-next}
V54 will turn (53.14) into a source-facing robust coupled-block
compiler.  It will combine the literal V46 bosonic floor with
gravity, fermion, and cross-block intervals, derive the exact
two-block lower edge, and state a finite infeasibility witness when
no allowed counterterm can make the constrained Hessian positive.
That witness is needed before the V55 companion audit.
\end{decision}

\subsection{Acquisition ledger}

\[
\begin{array}{p{0.30\linewidth}|p{0.20\linewidth}|p{0.39\linewidth}}
\textbf{gate}&\textbf{status after V53}&\textbf{content}\\ \hline
assemble-before-reduce identity&proved&
Noether Hessian correction\\
first-jet solution fibre&proved&
kernel ambiguity (53.10)\\
unique second jet criterion&proved&
\(\ker C\subset\ker\mathcal B\)\\
simultaneous feasibility&proved&
finite LMI (53.14)\\
analytic Hessian debit&proved&
bound (53.18)\\
outward matrix errors&proved&
robust floor (53.22)\\
cofinal \(C^2\) stability&proved&
summable Hessian transport\\
literal simultaneous card&missing&
source jets unpopulated\\
robust coupled-block witness&next&
gravity--matter intervals
\end{array}                                                 \tag{53.27}
\]

\subsection{Parent-integrity record}

The installed parent was
\[
\begin{array}{c|c}
 \text{file}&
 \texttt{Renewal\_SM\_Einstein\_Continuum\_Research\_Notes\_V52.tex}\\
 \text{lines}&20761\\
 \text{bytes}&677603\\
 \text{SHA--256}&
 \texttt{9A557DC085D2A21AA8E2356CA31E2651F23DE8A09B342639FB410604FDC37363}.
\end{array}                                                 \tag{53.28}
\]

% =============================================================================
% END APPENDED FOURTEENTH-CYCLE V53 MATERIAL
% =============================================================================
% =============================================================================
% BEGIN APPENDED FOURTEENTH-CYCLE V54 MATERIAL
% =============================================================================

\section{Fourteenth-cycle V54: sharp coupled bosonic floors and
finite infeasibility witnesses}
\label{sec:v14v54}

\subsection{The physical bosonic block}

After the complete action is assembled and the constraints are
reduced as in V53, split the commuting physical tangent into a
gravitational/clock part and a gauge--Higgs part:
\[
 {\mathbb H}_{\rm B}(a)
 =
 \begin{pmatrix}
 A(a)&X(a)\\
 X(a)^*&D(a)
 \end{pmatrix}.                                           \tag{54.1}
\]
All blocks may include allowed counterterm contributions and depend
affinely on \(a\) on the fixed chart.

The literal minimal \(K_4\) weak-Higgs-gauge result in V46 supplies
the conditional diagonal floor
\[
 D_{\rm min}
 \succeq
 \delta_{\rm HG}I,\qquad
 \delta_{\rm HG}=c_{\rm B}\sigma_0^2,                     \tag{54.2}
\]
before additional bosonic sectors and metric cross derivatives are
inserted.  It supplies no value for \(A\) or \(X\).

\subsection{The sharp two-block lower edge}

\begin{theorem}[Sharp lower edge from diagonal floors and cross norm]
\label{thm:v14v54-two-block}
Suppose
\[
 A\succeq\alpha I,\qquad
 D\succeq\delta I,\qquad
 \|X\|\le\chi.                                             \tag{54.3}
\]
Then
\[
 \boxed{
 \begin{pmatrix}A&X\\X^*&D\end{pmatrix}
 \succeq\lambda_-I,}
 \qquad
 \lambda_-=
 \frac{\alpha+\delta-
 \sqrt{(\alpha-\delta)^2+4\chi^2}}2.                      \tag{54.4}
\]
For \(\alpha,\delta>0\), this lower bound is positive exactly when
\[
 \boxed{\chi^2<\alpha\delta.}                              \tag{54.5}
\]
The bound is optimal given only the three scalars in (54.3).
\end{theorem}

\begin{proof}
For \(u,v\) in the two summands,
\[
 \begin{aligned}
 \left\langle
 \binom uv,
 \begin{pmatrix}A&X\\X^*&D\end{pmatrix}
 \binom uv
 \right\rangle
 &\ge
 \alpha\|u\|^2+\delta\|v\|^2-2\chi\|u\|\|v\|\\
 &=
 \begin{pmatrix}\|u\|\\\|v\|\end{pmatrix}^{\mathsf T}
 \begin{pmatrix}\alpha&-\chi\\-\chi&\delta\end{pmatrix}
 \begin{pmatrix}\|u\|\\\|v\|\end{pmatrix}.
 \end{aligned}                                             \tag{54.6}
\]
The smaller eigenvalue of the scalar \(2\times2\) matrix is
\(\lambda_-\), proving (54.4).  Its determinant is
\(\alpha\delta-\chi^2\), which gives (54.5).
Equality is attained for one-dimensional blocks
\(A=\alpha,D=\delta,X=-\chi\).
\end{proof}

The exact Schur criterion remains
\[
 {\mathbb H}_{\rm B}\succeq0
 \quad\Longleftrightarrow\quad
 D\succ0\ \hbox{ and }\
 A-XD^{-1}X^*\succeq0,                                    \tag{54.7}
\]
with the symmetric alternative interchanging \(A,D\).
Equation (54.4) is the robust scalar version when only outward
block bounds are available.

\subsection{Outward block intervals}

Suppose submitted matrices and radii satisfy
\[
 \|A-\widehat A\|\le\varepsilon_A,\qquad
 \|D-\widehat D\|\le\varepsilon_D,\qquad
 \|X-\widehat X\|\le\varepsilon_X.                         \tag{54.8}
\]
Define
\[
 \underline\alpha
 =\lambda_{\min}(\widehat A)-\varepsilon_A,\qquad
 \underline\delta
 =\lambda_{\min}(\widehat D)-\varepsilon_D,\qquad
 \overline\chi=\|\widehat X\|+\varepsilon_X.              \tag{54.9}
\]

\begin{corollary}[Robust coupled floor]
\label{cor:v14v54-robust}
The exact block (54.1) has lower edge at least
\[
 \boxed{
 \underline\lambda_-=
 \frac{\underline\alpha+\underline\delta-
 \sqrt{(\underline\alpha-\underline\delta)^2+
       4\overline\chi^2}}2.}                              \tag{54.10}
\]
In particular, a positive robust certificate follows from
\[
 \underline\alpha>0,\qquad
 \underline\delta>0,\qquad
 \overline\chi^2<
 \underline\alpha\,\underline\delta.                      \tag{54.11}
\]
\end{corollary}

\begin{proof}
Weyl's inequality and the triangle inequality give (54.3) with the
three values in (54.9).  Apply
Theorem~\ref{thm:v14v54-two-block}.
\end{proof}

The cross term is load bearing.  Positivity of the two submitted
diagonal blocks does not imply positivity of the coupled Hessian.

\begin{countertheorem}[Positive diagonals with an indefinite coupling]
\label{cth:v14v54-cross}
For
\[
 H_\chi=\begin{pmatrix}1&\chi\\\chi&1\end{pmatrix},
\]
both diagonal blocks equal one, but \(H_\chi\) is indefinite when
\(|\chi|>1\).
\end{countertheorem}

\begin{proof}
Its eigenvalues are \(1\pm\chi\).
\end{proof}

\subsection{Counterterm-dependent block test}

On the affine V53 chart, write
\[
 A(a)=A_0+\sum_ja_jA_j,\quad
 D(a)=D_0+\sum_ja_jD_j,\quad
 X(a)=X_0+\sum_ja_jX_j.                                   \tag{54.12}
\]
The exact simultaneous acceptance problem is
\[
 r+Ca=0,\qquad
 \begin{pmatrix}A(a)&X(a)\\X(a)^*&D(a)\end{pmatrix}
 \succeq\delta I,                                         \tag{54.13}
\]
together with the predeclared coefficient restrictions.
It is the same linear matrix inequality as (53.14), now in
source-facing blocks.  The scalar conditions (54.11) are sufficient,
while (54.13) is exact.

\begin{proposition}[A fixed negative direction is an exact obstruction]
\label{prop:v14v54-fixed-witness}
Suppose a nonzero physical vector \(v\) satisfies
\[
 \langle v,B_jv\rangle=0\quad(1\le j\le m),               \tag{54.14}
\]
and
\[
 \langle v,H_0v\rangle<\delta\|v\|^2.                     \tag{54.15}
\]
Then no counterterm coefficient can make
\(H_0+\sum_ja_jB_j\succeq\delta I\).
\end{proposition}

\begin{proof}
The quadratic form on \(v\) is independent of \(a\) by (54.14) and
violates the requested floor by (54.15).
\end{proof}

This is the simplest possible certificate that the declared
counterterm family cannot repair a physical mode.

\subsection{A semidefinite dual obstruction}

Ignore additional coefficient inequalities for the moment and
consider
\[
 Ca=-r,\qquad
 Q(a):=H_0+\sum_ja_jB_j-\delta I\succeq0.                 \tag{54.16}
\]

\begin{theorem}[Checkable dual infeasibility witness]
\label{thm:v14v54-dual}
Suppose there exist a positive semidefinite matrix \(Y\succeq0\) and
a dual vector \(y\) such that
\[
 \operatorname{Tr}(YB_j)+(C^*y)_j=0
 \quad(1\le j\le m),                                      \tag{54.17}
\]
and
\[
 \boxed{
 \operatorname{Tr}\bigl(Y(H_0-\delta I)\bigr)
 +\langle y,r\rangle<0.}                                  \tag{54.18}
\]
Then (54.16) is infeasible.
\end{theorem}

\begin{proof}
If \(a\) satisfied (54.16), then \(Y,Q(a)\succeq0\), so
\[
 \operatorname{Tr}(YQ(a))\ge0.                            \tag{54.19}
\]
Using (54.17) and \(Ca=-r\),
\[
 \begin{aligned}
 \operatorname{Tr}(YQ(a))
 &=
 \operatorname{Tr}\bigl(Y(H_0-\delta I)\bigr)
 +\sum_ja_j\operatorname{Tr}(YB_j)\\
 &=
 \operatorname{Tr}\bigl(Y(H_0-\delta I)\bigr)
 -\langle y,Ca\rangle\\
 &=
 \operatorname{Tr}\bigl(Y(H_0-\delta I)\bigr)
 +\langle y,r\rangle,
 \end{aligned}
\]
contradicting (54.18).
\end{proof}

For rational matrices, rational \(Y,y\), an exact PSD factorization
of \(Y\), and a strict rational inequality in (54.18) give a
reviewable finite obstruction.  The theorem supplies a sufficient
dual certificate; no claim of automatic strong duality is made
without the appropriate closedness or Slater hypotheses.

\subsection{Strict feasibility and perturbation radius}

\begin{proposition}[Robustness of a simultaneous passing point]
\label{prop:v14v54-slater}
Suppose a coefficient \(a_*\) satisfies the exact first-jet equation
and
\[
 {\mathbb H}_{\rm B}(a_*)\succeq\delta_*I,\qquad\delta_*>0.
                                                               \tag{54.20}
\]
If the combined first-jet perturbation is corrected within the same
coefficient support and the resulting coefficient change
\(\Delta a\) obeys
\[
 \sum_j|\Delta a_j|\,\|B_j\|+\epsilon_H<\delta_*,          \tag{54.21}
\]
where \(\epsilon_H\) is the remaining Hessian data error, then the
corrected physical Hessian remains positive with floor at least
\[
 \delta_*-\sum_j|\Delta a_j|\,\|B_j\|-\epsilon_H.          \tag{54.22}
\]
\end{proposition}

\begin{proof}
The total operator perturbation is bounded by the two debits in
(54.21).  Apply Weyl's inequality.
\end{proof}

This is the local Slater margin needed to transport a passing
counterterm prescription between adjacent cutoffs.

\subsection{Fermionic firewall}

The ordinary Loewner order in (54.1) applies to commuting real
bosonic and metric fluctuations.  A Grassmann Hessian is not an
additional positive real diagonal block.  The fermionic rows are:

\begin{enumerate}
\item the self-adjoint one-particle Dirac--Yukawa operator and its
declared filled/empty spectral convention;
\item the reflected crossing and one-sided factorization;
\item the CAR second quantization and charge algebra;
\item determinant/Pfaffian phase, orientation, and counterterm
contribution; and
\item the induced bosonic effective-action first and second metric
jets, if the fermions are integrated out.
\end{enumerate}

The graded nilpotent corrections discussed in V33 may be inverted
over a positive bosonic body, but they do not create an ordinary
Hilbert-space positivity floor.  Fermion stability and the bosonic
Loewner condition must remain typed separately.

\subsection{Application to the current source position}

The current literal information is
\[
 D_{\rm min}\succeq c_{\rm B}\sigma_0^2I                 \tag{54.23}
\]
for the minimal flat weak-Higgs-gauge quotient.  The following
remain unpopulated:
\[
 \underline\alpha,\qquad
 \overline\chi,\qquad
 \epsilon_A,\epsilon_D,\epsilon_X,\qquad
 [A_j,D_j,X_j].                                           \tag{54.24}
\]
The unreduced flat ADM block cannot supply
\(\underline\alpha>0\), by V47.  A passing value must come from a
regular nonlinear/clock reduction and the complete constrained
counterterm Hessian of V53.

\subsection{Robust coupled acquisition card}

\[
 \mathfrak B_X^{\rm coup}=
 \bigl(
 [P_{{\rm phys},X}],
 [A_0,D_0,X_0],
 [A_j,D_j,X_j],
 [\epsilon_A,\epsilon_D,\epsilon_X],
 [r,C,a],
 \delta_{\rm HG},
 [\underline\alpha,\underline\delta,\overline\chi],
 [\hbox{PSD witness or dual obstruction}]
 \bigr).                                                   \tag{54.25}
\]
No source passage inspected in V46 populates this complete card.

\begin{decision}[Direction after V54]
\label{dec:v14v54-next}
V55 is the compulsory companion audit.  It will check whether the
newest Yang--Mills or Navier--Stokes note supplies an adaptive
coverage certificate, a robust block-coupling theorem, or an exact
dual obstruction that changes the source-facing strategy.  V56 will
then return to the physical gravity block and derive a positive
true-Hamiltonian Hessian criterion from a regular clock solve,
because \(\underline\alpha\) in (54.24) is now the first missing
load-bearing scalar.
\end{decision}

\subsection{Acquisition ledger}

\[
\begin{array}{p{0.30\linewidth}|p{0.20\linewidth}|p{0.39\linewidth}}
\textbf{gate}&\textbf{status after V54}&\textbf{content}\\ \hline
two-block lower edge&proved&
sharp formula (54.4)\\
outward coupled floor&proved&
interval certificate (54.10)\\
diagonal-only inference&refuted&
scalar cross example\\
counterterm LMI&source-facing&
equation (54.13)\\
fixed-vector obstruction&proved&
counterterm-blind mode\\
dual infeasibility witness&proved&
finite certificate (54.17)--(54.18)\\
fermionic block typing&fixed&
CAR/Grassmann kept separate\\
literal coupled interval card&missing&
entries (54.24)\\
positive gravity true Hessian&after audit&
regular clock reduction
\end{array}                                                 \tag{54.26}
\]

\subsection{Parent-integrity record}

The installed parent was
\[
\begin{array}{c|c}
 \text{file}&
 \texttt{Renewal\_SM\_Einstein\_Continuum\_Research\_Notes\_V53.tex}\\
 \text{lines}&21163\\
 \text{bytes}&690154\\
 \text{SHA--256}&
 \texttt{E3CFDC6148958E231C49F16887BE420ECD90DF3E797B0207EC124AEFD9DCB1F9}.
\end{array}                                                 \tag{54.27}
\]

% =============================================================================
% END APPENDED FOURTEENTH-CYCLE V54 MATERIAL
% =============================================================================
\section{Fourteenth-cycle V55: compulsory companion audit and direct-certificate import rule}
\label{sec:v14v55}

This is the compulsory fifth-version audit following
Section~\ref{sec:v14v50}.  The audit is read-only: companion statements
are used only after their hypotheses have been translated into the
typed finite Standard Model--Einstein problem.  Numerical constants
belonging to a companion action are not transferred to the present
action.

\subsection{Live companion provenance}

\begin{record}[Files inspected at V55]
\label{rec:v14v55-provenance}
The newest active Yang--Mills note in the shared folder is
\[
 \path{Renewal_Yang_Mills_Mass_Gap_Research_Notes_V19.tex};
\]
it has \(152282\) bytes, \(4140\) lines, and SHA--256
\[
 \texttt{8CF5FEE84D9E0E8AB13D9DDD906EC9DFE8AC500EBBE7B4FBDD73DB49ABE8A593}.
                                                               \tag{55.1}
\]
The newest active Navier--Stokes note is
\[
 \path{Renewal_NS_Stability_Research_Notes_V26.tex};
\]
it has \(278283\) bytes, \(5319\) lines, and SHA--256
\[
 \texttt{82C887F8C2C69E4F28FAD7C3DC8DE5B9099CB5A4BF431EBA1FFF3E91935978AD}.
                                                               \tag{55.2}
\]
The NS digest is unchanged from the V50 audit.  The YM companion has
advanced from V17 to V19.
\end{record}

\begin{assessment}[What is new in the companions]
\label{ass:v14v55-reading}
Yang--Mills V18 constructs exact rational relative-metric patches at
all signed eighth-turn cell centres and proves by the explicit point
\((\pi/8,\pi/8,\pi/8)\) that the resulting atlas is still incomplete.
Yang--Mills V19 then goes upstream: it abandons auxiliary splitting at
two bottlenecks and certifies the direct \(45\)-dimensional Hermitian
pencil on four-dimensional boxes by a coefficientwise relative
Taylor estimate.  This removes a gluing obligation on those boxes,
but it is not a torus cover and supplies no cofinal floor.

Navier--Stokes V26 computes exact rank-one norms for two Oseen-kernel
derivative contractions.  That refinement is specific to a
parabolic kernel estimate.  It supplies neither a Standard Model
constraint operator nor a clock, counterterm, or physical Hessian
coefficient.
\end{assessment}

\subsection{The admissible import}

The useful companion lesson is the direct-certificate architecture.
We state it in the notation needed for the physical bosonic Hessian;
this also fixes exactly what source data would be required to turn
the V54 block criterion into a certified local scalar.

\begin{lemma}[Direct relative box certificate]
\label{lem:v14v55-direct-box}
Let
\[
 H(\theta)=\sum_{\alpha\in{\cal S}}C_\alpha
 e^{i\alpha\cdot\theta},\qquad
 C_{-\alpha}=C_\alpha^*,
                                                               \tag{55.3}
\]
be a finite Hermitian Laurent matrix on \(\mathbb T^d\).  Fix
\(a\in\mathbb T^d\), an invertible matrix \(R\), and rational numbers
\(\beta,\gamma,L_1,L_2>0\) for which
\[
\begin{split}
 R^*H(a)R&\succeq\beta I,\qquad R^*R\preceq\gamma I,\\
 \sum_{j=1}^d\|R^*(\partial_jH)(a)R\|_{\rm F}&\le L_1,\\
 \sum_{\alpha\in{\cal S}}|\alpha|_1^2
       \|R^*C_\alpha R\|_{\rm F}&\le L_2.
\end{split}                                                  \tag{55.4}
\]
If \(d_{\infty,\mathbb T^d}(\theta,a)\le h<\pi\), then
\[
 H(\theta)\succeq
 \frac{\beta-hL_1-\frac12h^2L_2}{\gamma}\,I.
                                                               \tag{55.5}
\]
Consequently the box is positive whenever the numerator in
\((55.5)\) is positive.
\end{lemma}

\begin{proof}
Choose the coordinate displacement \(u\in[-h,h]^d\) from \(a\) to
\(\theta\).  Taylor's formula with integral remainder applied to
\(t\mapsto R^*H(a+tu)R\) gives
\[
 \|R^*(H(\theta)-H(a))R\|_{\rm op}
 \le hL_1+\frac12h^2L_2,                                  \tag{55.6}
\]
because
\[
 |\,\alpha\cdot u\,|^2\le h^2|\alpha|_1^2,\qquad
 \|M\|_{\rm op}\le\|M\|_{\rm F}.                           \tag{55.7}
\]
Thus
\[
 R^*H(\theta)R\succeq
 \left(\beta-hL_1-\frac12h^2L_2\right)I.                  \tag{55.8}
\]
For \(x=Ry\), the second inequality in \((55.4)\) gives
\(\|x\|^2\le\gamma\|y\|^2\), hence
\(\|y\|^2\ge\|x\|^2/\gamma\).  Substitution into \((55.8)\)
proves \((55.5)\).
\end{proof}

\begin{corollary}[Application to the V54 coupled block]
\label{cor:v14v55-block}
Suppose the physical commuting bosonic block of V54 is supplied as a
finite Laurent bank \(H_B(a;\theta)\), after all constraint and gauge
reductions, and suppose Lemma~\ref{lem:v14v55-direct-box} certifies
\[
 H_B(a_0;\theta)\succeq\delta_{\rm dir}I
 \quad\hbox{on a parameter box }Q.                         \tag{55.9}
\]
Then the V54 diagonal/cross decomposition is unnecessary on \(Q\):
\(\delta_{\rm dir}\) itself is a valid coupled floor.  If only the
diagonal gravitational and Higgs--gauge blocks are available, the
sharp V54 two-block floor remains the applicable certificate.
\end{corollary}

\begin{proof}
Equation \((55.9)\) is already the Loewner conclusion sought for the
whole commuting bosonic block.  Introducing a block decomposition
can only replace this direct certificate by a sufficient lower
bound.  In the absence of a populated direct bank, the V54 theorem is
the available blockwise implication.
\end{proof}

\begin{decision}[Audit decision and direction]
\label{dec:v14v55-direction}
No companion numerical value is imported.  The direct relative-box
lemma is imported as a certification method because its proof uses
only finite Hermitian algebra and applies verbatim once a typed SM--
Einstein Laurent bank exists.  The Yang--Mills audit also gives a
warning relevant here: an indefinitely refined list of local
certificates is not a substitute for a finite cover with a uniform
floor.

The next upstream obstruction is the unpopulated gravitational entry
\(\underline\alpha\) in V54.  V56 will derive it from a regular
canonical clock sheet.  The calculation will distinguish the exact
true-Hamiltonian Hessian from the Hessian of the unreduced scalar
constraint, and will isolate a finite curvature debit caused by a
non-affine clock.
\end{decision}

\begin{firewall}[Claim boundary after V55]
\label{fire:v14v55}
The audit imports a proof architecture, not Yang--Mills or
Navier--Stokes constants.  No populated SM--Einstein Laurent bank,
regular clock sheet, positive gravitational true-Hamiltonian
Hessian, cofinal continuum measure, or full SM--Einstein continuum
has yet been constructed.
\end{firewall}

\section*{Version V55 append-only record}
\addcontentsline{toc}{section}{Version V55 append-only record}
The complete V54 source was retained before this continuation.  It
had \(702323\) bytes, \(21559\) lines, and SHA--256
\[
 \texttt{FD5F10AFD5EFCCBF89AFB7DE30DAA3BA3C90A0158A5CE960EF701D8EE7C34259}.
\]
V55 performed the compulsory companion audit and recorded the exact
direct-certificate import rule.  The next compulsory audit is V60.
\section{Fourteenth-cycle V56: regular clock sheets and the true-Hamiltonian Hessian}
\label{sec:v14v56}

Section~\ref{sec:v14v54} isolated the lower gravitational entry
\(\underline\alpha\) as the first missing scalar in the coupled
bosonic floor.  Section~\ref{sec:v14v49} showed that a clock must be
fixed before transfer energy can be compared with a canonical
generator.  We now compute the generator Hessian produced by an
arbitrary regular scalar-clock sheet.  The calculation is finite and
exact; it also identifies the additional debit that is absent for an
affine clock.

\subsection{Regular deparametrization}

Let \(Y\subset\mathbb R^n\) be an affine physical chart obtained after
the other first-class constraints and gauge directions have been
removed.  Let \((T,P)\) be a canonical clock pair and let
\[
 {\cal C}:Y\times I_T\times I_P\longrightarrow\mathbb R
                                                               \tag{56.1}
\]
be a \(C^3\) scalar constraint.  Fix
\((y_0,T_0,P_0)\) with
\[
 {\cal C}(y_0,T_0,P_0)=0,\qquad
 \kappa_0:=\partial_P{\cal C}(y_0,T_0,P_0)>0.              \tag{56.2}
\]
Multiplying the constraint by \(-1\), if needed, fixes the displayed
orientation without changing its zero set.

\begin{theorem}[Exact first and second jets of the true Hamiltonian]
\label{thm:v14v56-implicit}
There are neighbourhoods of \((y_0,T_0)\) and a unique \(C^3\)
function \(H_{\rm true}(y,T)\) such that
\[
 P=-H_{\rm true}(y,T),\qquad
 {\cal C}\bigl(y,T,-H_{\rm true}(y,T)\bigr)=0.             \tag{56.3}
\]
At a point of this sheet put
\[
 \kappa={\cal C}_P,\quad
 A={\cal C}_{yy},\quad b={\cal C}_{yP},\quad
 c={\cal C}_{PP},\quad h=\nabla_yH_{\rm true}.             \tag{56.4}
\]
Then
\[
 h=\frac{{\cal C}_y}{\kappa}                               \tag{56.5}
\]
and
\[
 D_y^2H_{\rm true}
 =
 \frac1{\kappa}
 \left(A-b\otimes h-h\otimes b+c\,h\otimes h\right).       \tag{56.6}
\]
All derivatives of \({\cal C}\) in \((56.4)\)--\((56.6)\) are
evaluated at \((y,T,-H_{\rm true}(y,T))\).
\end{theorem}

\begin{proof}
The implicit-function theorem applies to
\[
 G(y,T,H):={\cal C}(y,T,-H)
\]
because \(G_H=-{\cal C}_P=-\kappa\ne0\).  Differentiating
\(G(y,T,H_{\rm true}(y,T))=0\) in \(y_i\) gives
\[
 0={\cal C}_i-\kappa\,\partial_iH_{\rm true},
\]
which is \((56.5)\).  A second derivative gives
\[
 0={\cal C}_{ij}
   -{\cal C}_{iP}h_j-{\cal C}_{jP}h_i
   +{\cal C}_{PP}h_ih_j-\kappa\,\partial_{ij}H_{\rm true}.
                                                               \tag{56.7}
\]
Solving \((56.7)\) for the last term proves \((56.6)\).
\end{proof}

\begin{remark}[Constraint rescaling]
\label{rem:v14v56-rescaling}
Formula \((56.6)\) is invariant on the constraint sheet under
\({\cal C}\mapsto f{\cal C}\) for any positive \(C^2\) function
\(f\).  Indeed, terms containing a derivative of \(f\) cancel after
\({\cal C}=0\) and \({\cal C}_y=\kappa h\) are used.  Thus the clock
Hessian does not depend on a positive choice of defining function for
the same oriented sheet.
\end{remark}

\begin{proof}
On the sheet,
\[
 (f{\cal C})_P=f\kappa,\qquad
 (f{\cal C})_y=f{\cal C}_y.
\]
Expanding the four second derivatives in the numerator of
\((56.6)\), the terms proportional to \(D f\) occur in opposite
pairs, while the terms proportional to \(D^2f\) carry
\({\cal C}=0\).  The remaining numerator and denominator are both
multiplied by \(f\).
\end{proof}

\subsection{A checkable positive floor}

\begin{theorem}[Clock-sheet lower bound]
\label{thm:v14v56-floor}
On a set \(Q\) in the regular sheet suppose
\[
\begin{gathered}
 0<\underline\kappa\le\kappa\le\overline\kappa,\qquad
 A\succeq\alpha I,\\
 \|b\|_2\le\beta,\qquad c\ge-\gamma,\qquad
 \|h\|_2\le\eta.                                          \tag{56.8}
\end{gathered}
\]
Then
\[
 D_y^2H_{\rm true}\succeq
 \underline\alpha_{\rm clock} I,\qquad
 \underline\alpha_{\rm clock}
 :=
 \frac{\alpha-2\beta\eta-\gamma\eta^2}
      {\overline\kappa}.                                  \tag{56.9}
\]
In particular this is a positive physical gravitational floor if
\[
 \alpha>2\beta\eta+\gamma\eta^2.                          \tag{56.10}
\]
At a stationary point of the true Hamiltonian,
\[
 h=0,\qquad
 D_y^2H_{\rm true}=\frac{A}{\kappa};                       \tag{56.11}
\]
there the nonlinear-clock debit vanishes exactly.
\end{theorem}

\begin{proof}
For a unit vector \(v\), the numerator of \((56.6)\) obeys
\[
\begin{split}
 v^*\!\left(A-b\otimes h-h\otimes b+c\,h\otimes h\right)v
 &\ge \alpha-2|\langle b,v\rangle|\,|\langle h,v\rangle|
        -\gamma|\langle h,v\rangle|^2\\
 &\ge \alpha-2\beta\eta-\gamma\eta^2.                    \tag{56.12}
\end{split}
\]
When the last expression is positive, division by
\(\kappa\le\overline\kappa\) gives \((56.9)\).  If \(h=0\),
formula \((56.6)\) reduces to \((56.11)\).
\end{proof}

\begin{corollary}[Typed input for the V54 coupled floor]
\label{cor:v14v56-v54}
Let \(A_G\) be the restriction of \({\cal C}_{yy}\) to the
gravitational physical coordinates in the chart \(Y\), and suppose
the bounds in Theorem~\ref{thm:v14v56-floor} hold with
\(A_G\succeq\alpha_GI\).  Then V54 may take
\[
 \underline\alpha
 =
 \frac{\alpha_G-2\beta_G\eta_G-\gamma\eta_G^2}
      {\overline\kappa},                                  \tag{56.13}
\]
provided the numerator is positive.  Coupling this entry to the
V46 Higgs--gauge floor
\(\delta_{HG}=c_B\sigma_0^2\) and cross norm \(\chi\) gives
\[
 \lambda_{\rm phys}\ge
 \frac{\underline\alpha+\delta_{HG}
 -\sqrt{(\underline\alpha-\delta_{HG})^2+4\chi^2}}2,       \tag{56.14}
\]
which is positive exactly when
\(\chi^2<\underline\alpha\,\delta_{HG}\).
\end{corollary}

\begin{proof}
The gravitational diagonal estimate is
Theorem~\ref{thm:v14v56-floor}.  Substitute it into the sharp
two-block theorem of Section~\ref{sec:v14v54}.
\end{proof}

\subsection{Affine and quadratic clocks}

\begin{corollary}[Affine clock has no curvature debit]
\label{cor:v14v56-affine}
If
\[
 {\cal C}(y,T,P)=P+K(y,T),                                \tag{56.15}
\]
then
\[
 H_{\rm true}=K,\qquad
 D_y^2H_{\rm true}=D_y^2K.                               \tag{56.16}
\]
Thus \(b=c=0\), \(\kappa=1\), and the clock contributes no correction
to the physical Hessian.
\end{corollary}

\begin{proof}
Equation \((56.3)\) reads \(-H_{\rm true}+K=0\).
The derivative statement follows directly, or from \((56.6)\).
\end{proof}

For a massless scalar clock one instead encounters, after spatial
constraints have been treated, the local algebraic model
\[
 {\cal C}(y,T,P)=C_0(y,T)+\frac{P^2}{2M(y,T)},\qquad M>0.
                                                               \tag{56.17}
\]
On a sheet with \(C_0<0\) and \(P=-H_{\rm true}<0\), set
\[
 R(y,T):=-2M(y,T)C_0(y,T)>0.                              \tag{56.18}
\]

\begin{theorem}[Square-root clock Hessian and its sharp rank-one debit]
\label{thm:v14v56-square-root}
For the sheet just specified,
\[
 H_{\rm true}=\sqrt R,\qquad
 D_y^2H_{\rm true}
 =\frac1{2\sqrt R}D_y^2R
  -\frac1{4R^{3/2}}\,dR\otimes dR.                        \tag{56.19}
\]
If throughout a panel \(Q\)
\[
 0<H_-\le\sqrt R\le H_+,\qquad
 D_y^2R\succeq\rho I,\qquad \|dR\|_2\le g,                \tag{56.20}
\]
then
\[
 D_y^2H_{\rm true}\succeq
 \left(\frac{\rho}{2H_+}-\frac{g^2}{4H_-^3}\right)I.       \tag{56.21}
\]
The negative rank-one term in \((56.19)\) cannot in general be
discarded.
\end{theorem}

\begin{proof}
The constraint equation gives \(P^2=R\); the chosen sheet gives
\(H_{\rm true}=-P=\sqrt R\).  Two differentiations of the square root
give \((56.19)\).  For a unit vector \(v\),
\[
 \frac{D^2R[v,v]}{2\sqrt R}
 -\frac{|dR[v]|^2}{4R^{3/2}}
 \ge \frac{\rho}{2H_+}-\frac{g^2}{4H_-^3},
\]
which proves \((56.21)\).

For necessity of retaining the second term, take the one-dimensional
local function
\[
 R(x)=1+4x+x^2.
\]
At \(x=0\), \(R>0\), \(R''=2>0\), but
\[
 (\sqrt R)''(0)=\frac{2}{2}-\frac{4^2}{4}=-3<0.           \tag{56.22}
\]
Thus positivity of \(D^2R\) alone does not imply positivity of the
true-Hamiltonian Hessian.
\end{proof}

\begin{assessment}[What has been gained]
\label{ass:v14v56-gain}
The gravitational scalar required by V54 is now an explicit
clock-sheet quantity.  For a regular clock it is \((56.13)\); for a
scalar square-root clock it is bounded by \((56.21)\).  These formulas
show two independent failure modes: loss of clock regularity when
\(\kappa\) approaches zero, and loss of convexity from the
clock-gradient rank-one debit.  Both must be measured before any
coupled SM--Einstein positivity claim.

The formulas are local in a specified physical chart.  At a critical
point the Hessian is intrinsic under changes of physical coordinates.
Away from a critical point, a chart-independent statement requires a
specified connection; all floors above refer to the fixed affine
chart \(Y\).
\end{assessment}

\begin{decision}[V57 source target]
\label{dec:v14v56-next}
Populate the square-root-clock inputs from a finite reduced
Einstein--matter cell: compute \(R=-2MC_0\), its gradient and Hessian,
and separate the pure gravitational, matter, and cross blocks before
using \((56.21)\).  If the attached sources do not supply a finite
clocked cell action with these coefficients, prove that provenance
failure explicitly and go upstream to the ADM lapse/shift elimination
that produces \(C_0\).
\end{decision}

\begin{firewall}[Claim boundary after V56]
\label{fire:v14v56}
Theorems~\ref{thm:v14v56-implicit} and
\ref{thm:v14v56-square-root} are exact conditional clock-sheet
calculi.  The required finite Einstein--matter coefficients
\((M,C_0)\) and a regular physical panel have not yet been populated.
No positive numerical gravitational floor or full SM--Einstein
continuum follows.
\end{firewall}

\section*{Version V56 append-only record}
\addcontentsline{toc}{section}{Version V56 append-only record}
The complete V55 source was retained before this continuation.  It
had \(709215\) bytes, \(21733\) lines, and SHA--256
\[
 \texttt{B7FC9B7385C882C4016C5B397427897334142382B83460AF03B3497035F0F9B5}.
\]
V56 derived the exact true-Hamiltonian Hessian on a regular clock
sheet and the quantitative nonlinear-clock curvature debit.
\section{Fourteenth-cycle V57: source ledger and the scalar-clock inertia obstruction}
\label{sec:v14v57}

V56 requested a populated finite clocked Einstein--matter cell.  The
attached manuscripts do not provide that object, but the Grand Tensor
source does provide enough symbolic ADM information to prove a
stronger statement: the usual quadratic scalar-clock solve cannot
supply the positive gravitational diagonal assumed in V54 on a
physical chart that retains both DeWitt signs.

\subsection{Literal source ledger}

\begin{record}[Canonical data present in the attached sources]
\label{rec:v14v57-source}
The Grand Tensor manuscript
\[
 \path{Renewal_Grand_Tensor_Monograph_V067.tex}
\]
has the SHA--256 digest recorded in V46,
\[
 \texttt{FACD058EC8BB210AD8B296C9D7F4D78A229D08D7A2AE378808700655CB8BC0FE}.
                                                               \tag{57.1}
\]
Its theorem labelled \(\texttt{thm:DeWitt}\) derives the cotangent
quadratic form
\[
 {\cal G}_{\rm DW}(\pi,\pi)
 =
 \pi^{ab}\pi_{ab}-\frac12\pi^2                         \tag{57.2}
\]
and records signature \((5,1)\) on the six metric-momentum
coordinates.  Its theorem labelled
\(\texttt{thm:SMST-common-action}\) gives the symbolic gravitational
density
\[
 {\cal H}_{\rm g}
 =
 \frac{\chi^{-1}}{\sqrt q}
 \left(\pi^{ab}\pi_{ab}-\frac12\pi^2\right)
 -\chi\sqrt q\,{\cal R}+\Lambda_H\sqrt q.                 \tag{57.3}
\]
Its finite-action Einstein theorem defines a common finite action,
its stress first variation, Ward identities, and a nonnegative
stationarity residual.  It does not display a canonical clock pair
\((T,P)\), a scalar constraint of the form \((56.17)\), or numerical
finite-cell coefficients for \(M,C_0\).

The attached Einstein--SM V30 source, with the digest recorded in
V46,
\[
 \texttt{E0AE0E0990DE472021A1E90A9C10D41C3F747F28D97A5B1D8F1281852E3A0E69},
                                                               \tag{57.4}
\]
contains physical-duration, resampling-clock, Duhamel, and
transfer-log constructions.  Its clock terminology is operational;
the inspected source contains no declaration of a canonical scalar
clock momentum \(P\) or a finite ADM constraint that can be inserted
into \((56.6)\).
\end{record}

\begin{assessment}[Exact provenance conclusion]
\label{ass:v14v57-provenance}
The available source rows determine the algebraic DeWitt coefficient
and a conditional common-action Hamiltonian.  They do not determine
a clock sheet or its numerical first and second jets.  Therefore the
numbers \(\underline\kappa,\overline\kappa,\alpha,\beta,\gamma,\eta\)
in \((56.8)\) cannot be populated from the attached files.  This is a
source-provenance failure, not a proof that no suitable physical
clock exists.
\end{assessment}

\subsection{The quadratic scalar-clock calculation}

Let \(E\) be a real finite-dimensional gravitational momentum space
with a nondegenerate symmetric form \(G\).  Write
\[
 Q(\pi)=\langle\pi,G\pi\rangle.                            \tag{57.5}
\]
At fixed metric and matter coordinates assume the non-clock part of
the scalar constraint has the ADM kinetic dependence
\[
 C_0(\pi)=\frac{\chi^{-1}}{M}Q(\pi)+V,\qquad
 \chi>0,\quad M>0,                                        \tag{57.6}
\]
and use the quadratic clock
\[
 {\cal C}(\pi,P)=C_0(\pi)+\frac{P^2}{2M}.                  \tag{57.7}
\]
Suppose \(R_0:=-2MV>0\), so both local sheets exist at \(\pi=0\).

\begin{theorem}[Scalar-clock inertia theorem]
\label{thm:v14v57-inertia}
For \(\varepsilon\in\{+1,-1\}\), the two signed true-Hamiltonian
branches are
\[
 H_\varepsilon(\pi)
 =\varepsilon\sqrt{R_0-2\chi^{-1}Q(\pi)}.                 \tag{57.8}
\]
At \(\pi=0\),
\[
 dH_\varepsilon(0)=0,\qquad
 D^2H_\varepsilon(0)
 =-\frac{2\varepsilon\chi^{-1}}{\sqrt{R_0}}\,G.            \tag{57.9}
\]
If \(G\) has inertia \((p,q)\) with \(p,q>0\), then
\(D^2H_+(0)\) has inertia \((q,p)\), while
\(D^2H_-(0)\) has inertia \((p,q)\).  In particular neither clock
sheet has a positive definite full gravitational momentum Hessian.
For the DeWitt form \((57.2)\), the two inertias are respectively
\[
 (1,5)\quad\hbox{and}\quad(5,1).                          \tag{57.10}
\]
\end{theorem}

\begin{proof}
Equations \((57.6)\)--\((57.7)\) give
\[
 P^2=-2MC_0=R_0-2\chi^{-1}Q(\pi).
\]
The two orientations give \((57.8)\).  Because \(Q\) is quadratic,
\[
 dQ(0)=0,\qquad D^2Q(0)=2G.
\]
Applying the square-root formula \((56.19)\), the rank-one derivative
term vanishes at zero and
\[
 D^2H_\varepsilon(0)
 =\frac{\varepsilon}{2\sqrt{R_0}}
   \bigl(-4\chi^{-1}G\bigr),
\]
which is \((57.9)\).  Multiplication of a nondegenerate symmetric
form by a positive scalar preserves inertia and multiplication by a
negative scalar interchanges its positive and negative counts.
This proves all assertions.
\end{proof}

\begin{corollary}[The flat linearized quotient retains both signs]
\label{cor:v14v57-flat-quotient}
On a compact flat background with \(\pi=0\), the kernel of the
linearized ADM constraints contains constant trace momentum and
constant traceless momentum.  Infinitesimal spatial gauge acts
trivially on the background momentum.  Hence the linearized
constraint quotient retains a vector on each side of the DeWitt
cone, and the two clock-sheet Hessians in
Theorem~\ref{thm:v14v57-inertia} remain indefinite after this
linearized reduction.
\end{corollary}

\begin{proof}
At \(\pi=0\), the first variation of the quadratic kinetic part of
the scalar constraint with respect to \(\pi\) is zero.  The
linearized scalar constraint therefore imposes no condition on a
pure momentum variation.  The linearized momentum constraint is the
divergence of the momentum variation, which vanishes for every
constant symmetric tensor.  The infinitesimal spatial-diffeomorphism
variation of \(\pi\) is its Lie derivative; it is zero at the
background \(\pi=0\).

For \(\pi=fq^{-1}\), \(f\ne0\), one has
\[
 Q(\pi)=3f^2-\frac12(3f)^2=-\frac32f^2<0.                 \tag{57.11}
\]
For a nonzero constant traceless \(\pi\),
\[
 Q(\pi)=\pi^{ab}\pi_{ab}>0.                               \tag{57.12}
\]
Both classes survive, proving the claim.
\end{proof}

\begin{proposition}[Robustness under small finite corrections]
\label{prop:v14v57-robust}
Let
\[
 K_0=-\frac{2\varepsilon\chi^{-1}}{\sqrt{R_0}}G
\]
and let
\[
 m_0:=\min\{|\lambda|:\lambda\in\sigma(K_0)\}>0.
\]
If a symmetric correction \(E\) satisfies
\[
 \|E\|_{\rm op}<m_0,                                      \tag{57.13}
\]
then \(K_0+E\) has the same inertia as \(K_0\).  Therefore metric-only
counterterms, which contribute no momentum Hessian, and sufficiently
small gravity--matter kinetic corrections cannot make the
scalar-clock Hessian positive definite.
\end{proposition}

\begin{proof}
Order the eigenvalues increasingly.  Weyl's inequality gives
\[
 |\lambda_j(K_0+E)-\lambda_j(K_0)|\le\|E\|_{\rm op}<m_0.
\]
No eigenvalue can cross zero, so the positive and negative counts are
unchanged.  A counterterm depending only on metric coordinates has
zero second derivative in the momentum block.
\end{proof}

\begin{countertheorem}[The V54 positive-diagonal route fails for this clock]
\label{ctr:v14v57-v54}
For the quadratic scalar clock \((57.7)\) on any physical chart
retaining both DeWitt signs, there is no
\(\underline\alpha>0\) such that
\[
 D_\pi^2H_\varepsilon(0)\succeq\underline\alpha I.        \tag{57.14}
\]
Consequently the positive two-diagonal hypothesis used by the V54
Schur floor cannot be satisfied by this gravitational momentum block.
\end{countertheorem}

\begin{proof}
Theorem~\ref{thm:v14v57-inertia} gives a negative eigenvalue on
either sheet.  Inequality \((57.14)\) would make every eigenvalue
positive, a contradiction.
\end{proof}

\begin{assessment}[Required change of direction]
\label{ass:v14v57-pivot}
The missing scalar in V54 is not merely waiting for numerical
acquisition on the standard quadratic scalar-clock branch.  Its
positivity is false at the flat zero-momentum reference whenever the
physical quotient retains both DeWitt signs.  Continuing to optimize
the V54 cross norm would therefore be circular.

There are four logically distinct exits: choose a reduction that
provably removes every negative direction; use a different clock
whose true generator has a positive physical Hessian; perform a
controlled conformal contour rotation before the positive transfer
construction; or abandon local Hessian positivity and prove
semibounded self-adjoint dynamics by a nonquadratic argument.  The
attached sources currently certify none of these exits.
\end{assessment}

\begin{decision}[V58 attacks the contour branch]
\label{dec:v14v57-next}
Analyze the finite quadratic conformal rotation exactly.  Determine
when a complex linear contour turns the DeWitt form positive, how
the functional measure and reality involution transform, and which
reflection-positivity condition is required to return a self-adjoint
Lorentzian generator.  If the rotation repairs the quadratic form but
breaks the physical involution, record the obstruction and move
upstream to a direct semibounded transfer-form construction.
\end{decision}

\begin{firewall}[Claim boundary after V57]
\label{fire:v14v57}
The source ledger and inertia theorem rule out one positive-Hessian
route; they do not rule out all clocks or all quantum constructions.
No conformal contour, reflection-positive return map, semibounded
physical Hamiltonian, or full SM--Einstein continuum has been proved.
\end{firewall}

\section*{Version V57 append-only record}
\addcontentsline{toc}{section}{Version V57 append-only record}
The complete V56 source was retained before this continuation.  It
had \(719384\) bytes, \(22041\) lines, and SHA--256
\[
 \texttt{2ACB126661CA86E2633786C67C197670B851CA0E63A147DC475DDC0BAB6EB605}.
\]
V57 recorded the canonical source ledger and proved the quadratic
scalar-clock inertia obstruction.
\section{Fourteenth-cycle V58: finite conformal rotation and the reflected-involution test}
\label{sec:v14v58}

The V57 obstruction concerns a real quadratic gravitational momentum
form.  A standard proposed repair is to rotate its negative
directions into an imaginary contour.  At finite regulator this does
make the quadratic form positive.  The missing issue is the physical
adjoint: complex conjugation does not preserve the rotated contour
pointwise and induces an additional sign involution.  We compute that
involution and the exact Gaussian reflection-positivity test.

\subsection{Quadratic contour rotation}

Let \(E\) be a finite-dimensional real inner-product space and let
\(G=G^*\) be invertible with spectral decomposition
\[
 E=E_{\rm p}\oplus E_{\rm n},\qquad
 G=A\oplus(-B),\qquad A\succ0,\quad B\succ0.               \tag{58.1}
\]
Let \(P_{\rm p},P_{\rm n}\) be the two orthogonal projections and put
\[
 L=P_{\rm p}+iP_{\rm n},\qquad
 J=P_{\rm p}-P_{\rm n}.                                   \tag{58.2}
\]

\begin{theorem}[Exact finite conformal rotation]
\label{thm:v14v58-rotation}
The complex contour \({\cal C}=L(E)\subset E_{\mathbb C}\) has the
following properties:
\[
 L^{\mathsf T}GL=A\oplus B=|G|\succ0,                     \tag{58.3}
\]
\[
 \overline{Ly}=LJy,\qquad y\in E,                         \tag{58.4}
\]
and, after compatible orientations are chosen,
\[
 dz=i^{\dim E_{\rm n}}\,dy.                               \tag{58.5}
\]
Consequently, for every entire \(F\) for which the contour
deformation is justified at infinity,
\[
 \int_{\cal C}F(z)e^{-\frac12z^{\mathsf T}Gz}\,dz
 =
 i^{\dim E_{\rm n}}
 \int_E F(Ly)e^{-\frac12\langle y,|G|y\rangle}\,dy.        \tag{58.6}
\]
The rotated weight is positive, but the physical real involution
pulls back to \(J\), not to the identity.
\end{theorem}

\begin{proof}
On \(E_{\rm p}\), \(L=I\), while on \(E_{\rm n}\), \(L=iI\).
Because the bilinear transpose, rather than the Hermitian adjoint,
enters a holomorphic quadratic form,
\[
 i(-B)i=B,
\]
which proves \((58.3)\).  Complex conjugation changes \(i\) to
\(-i\), giving \((58.4)\).  The determinant of \(L\) is
\(i^{\dim E_{\rm n}}\), proving \((58.5)\) and the change-of-variables
formula.  The analytic and decay hypotheses are exactly those needed
to move the real contour to \(L(E)\) without boundary terms.
\end{proof}

\begin{remark}[The Jacobian phase is not the adjoint problem]
\label{rem:v14v58-phase}
For normalized finite expectations, the constant phase in
\((58.6)\) cancels between numerator and partition function whenever
both contour deformations are valid.  The sign \(J\) in \((58.4)\)
does not cancel: it changes the involution used to define real
observables and reflected inner products.
\end{remark}

\subsection{The induced reflection}

Let \(\theta:E\to E\) be the real linear time-reflection involution,
assumed orthogonal and commuting with \(G\).  Then it commutes with
the spectral projections in \((58.2)\).  Let \(E_+^{\rm time}\) be
the future coordinate space and let
\(\iota_+:E_+^{\rm time}\to E\) be its inclusion.

\begin{proposition}[Pulled-back physical reflection]
\label{prop:v14v58-reflection}
If the original antilinear reflection acts on complex configurations
by
\[
 z\longmapsto\theta\overline z,                            \tag{58.7}
\]
then on the real rotated coordinate \(y\) it acts by
\[
 y\longmapsto\Theta_Cy,\qquad
 \Theta_C:=\theta J=J\theta.                              \tag{58.8}
\]
Thus the correct reflected adjoint of a test function is
\[
 (\mathfrak R_CF)(y)
 :=\overline{F(\Theta_Cy)}.                               \tag{58.9}
\]
\end{proposition}

\begin{proof}
Using \((58.4)\) and commutation of \(\theta\) with \(L,J\),
\[
 \theta\overline{Ly}
 =\theta LJy=L\theta Jy=L\Theta_Cy.
\]
Pullback of the original reflected conjugation therefore gives
\((58.9)\).
\end{proof}

\subsection{Finite Gaussian reflection positivity}

Let \(K\succ0\) commute with an orthogonal involution \(\Theta\), and
let \(\mu_K\) be the centered Gaussian probability law with
covariance \(C=K^{-1}\).

\begin{theorem}[Finite Gaussian OS criterion]
\label{thm:v14v58-gaussian-os}
For the polynomial algebra generated by future linear coordinates,
\[
 \int_E\overline{F(\Theta y)}F(y)\,d\mu_K(y)\ge0
 \quad\hbox{for every future polynomial }F               \tag{58.10}
\]
holds if and only if the future cross-covariance
\[
 C_{\rm OS}
 :=
 \iota_+^*\Theta C\,\iota_+                              \tag{58.11}
\]
is positive semidefinite.
\end{theorem}

\begin{proof}
Necessity follows by taking
\(F(y)=\langle u,\iota_+^*y\rangle\):
\[
 \int\overline{F(\Theta y)}F(y)\,d\mu_K(y)
 =\langle u,C_{\rm OS}u\rangle.
\]

For sufficiency, first use exponential vectors
\[
 F_u(y)=\exp\langle\iota_+u,y\rangle.
\]
The Gaussian moment formula shows that their reflected Gram, after
multiplication of each row and column by a positive diagonal
normalization, is
\[
 {\cal K}(u,v)=\exp\langle u,C_{\rm OS}v\rangle.           \tag{58.12}
\]
If \(C_{\rm OS}\succeq0\), then every Hadamard power of its Gram
matrix is positive semidefinite.  The power series of the exponential
therefore makes \((58.12)\) a positive kernel.  Differentiating the
exponential vectors at the origin gives all polynomial vectors, so
their reflected Gram matrices are positive semidefinite as well.
\end{proof}

\begin{corollary}[Exact test after conformal rotation]
\label{cor:v14v58-rotated-os}
For the rotated gravitational Gaussian of
Theorem~\ref{thm:v14v58-rotation}, the physical reflected form is
positive exactly when
\[
 \boxed{
 C_{\rm OS}^{\rm rot}
 =
 \iota_+^*\,\theta J\,|G|^{-1}\iota_+
 \succeq0.}                                               \tag{58.13}
\]
Positivity of the Euclidean quadratic matrix \(|G|\) alone does not
imply \((58.13)\).
\end{corollary}

\begin{proof}
Apply Theorem~\ref{thm:v14v58-gaussian-os} with
\(K=|G|\) and the pulled-back involution
\(\Theta_C=\theta J\) from Proposition~\ref{prop:v14v58-reflection}.
\end{proof}

\begin{counterexample}[Positive rotated action can fail physical reflection positivity]
\label{ctr:v14v58-os}
Let \(E=\mathbb R^2\) represent one scalar coordinate at negative and
positive time, let
\[
 \theta=\begin{pmatrix}0&1\\1&0\end{pmatrix},\qquad
 C=\begin{pmatrix}1&r\\r&1\end{pmatrix},\qquad 0<r<1,
                                                               \tag{58.14}
\]
and rotate this entire scalar sector, so \(J=-I\).  Then
\(C\succ0\), hence \(K=C^{-1}\succ0\), but for the positive-time
inclusion \(\iota_+u=(0,u)\),
\[
 \iota_+^*\theta JC\iota_+=-r<0.                          \tag{58.15}
\]
The positive Gaussian weight therefore violates the physical
reflected inequality already on a linear observable.
\end{counterexample}

\begin{proof}
The eigenvalues of \(C\) are \(1\pm r>0\).  Direct multiplication in
\((58.15)\) gives the negative number \(-r\).
The necessity part of Theorem~\ref{thm:v14v58-gaussian-os} completes
the proof.
\end{proof}

\subsection{Cofinal requirements}

\begin{theorem}[Sufficient cofinal contour packet]
\label{thm:v14v58-cofinal}
At regulator \(n\), suppose a finite gravitational quadratic packet
has matrices \(G_n\), contour maps \(L_n\), sign involutions \(J_n\),
time reflections \(\theta_n\), future inclusions \(\iota_{n,+}\), and
positive rotated covariances \(C_n=|G_n|^{-1}\).  Assume:
\begin{enumerate}
\item the contour deformation identity \((58.6)\) holds on a
      reflection-stable core algebra;
\item \(\theta_nJ_n=J_n\theta_n\) and
      \(C_{{\rm OS},n}:=\iota_{n,+}^*\theta_nJ_nC_n\iota_{n,+}
      \succeq0\);
\item the embeddings intertwine \(J_n,\theta_n\) and the future
      algebras;
\item the reflected seminorms and normalized moments converge on the
      core.
\end{enumerate}
Then the limiting core functional is reflection positive.
\end{theorem}

\begin{proof}
By item 2 and Theorem~\ref{thm:v14v58-gaussian-os}, every finite
reflected Gram matrix is positive semidefinite.  Items 1 and 3
identify these matrices with the physical contour observables and
make them compatible under refinement.  Entrywise convergence in
item 4 preserves positive semidefiniteness of every fixed finite Gram
matrix.  Hence the limiting core functional is reflection positive.
\end{proof}

\begin{assessment}[Contour branch status]
\label{ass:v14v58-status}
The conformal rotation repairs the finite quadratic weight exactly,
but it replaces ordinary time reflection by the twisted involution
\(\theta J\).  The new load-bearing datum is the cross-covariance
matrix \((58.13)\).  Neither attached source supplies a conformal
contour map together with this physical reflected covariance, and
ordinary positivity of a rotated Hessian cannot substitute for it.

The result also clarifies why a modewise rotation cannot be inserted
only at the end of a transfer proof: the negative spectral projector
must be compatible with time reflection, refinement embeddings,
gauge reduction, and the observable adjoint at every regulator.
\end{assessment}

\begin{decision}[V59 attacks semibounded transfer forms]
\label{dec:v14v58-next}
Because the contour branch creates an unpopulated reflected-covariance
obligation, return to the positive transfer operator.  Derive
semibounded self-adjoint dynamics directly from a closed
reflection-positive transfer form, without requiring a positive
classical gravitational Hessian.  Isolate the exact finite and
cofinal hypotheses, and determine which of them the V39--V44 transfer
cards already provide.
\end{decision}

\begin{firewall}[Claim boundary after V58]
\label{fire:v14v58}
A positive finite contour-rotated quadratic action is not by itself a
physical OS theory.  The twisted cross covariance, analytic contour
deformation, compatible refinement, interacting measure, and
Lorentzian return remain unproved.  No full SM--Einstein continuum is
claimed.
\end{firewall}

\section*{Version V58 append-only record}
\addcontentsline{toc}{section}{Version V58 append-only record}
The complete V57 source was retained before this continuation.  It
had \(729265\) bytes, \(22303\) lines, and SHA--256
\[
 \texttt{3BDB3BF401AD73EBEE9B9EDF82A2FD202CCDE2A55775B461708722B45CA29B98}.
\]
V58 proved the finite conformal-rotation identity and the exact
twisted reflection-positivity criterion.
\section{Fourteenth-cycle V59: semibounded dynamics from finite reflected transfer forms}
\label{sec:v14v59}

The classical clock Hessian can be indefinite and a conformal contour
requires an additional reflected-covariance certificate.  Neither
fact prevents a positive transfer operator from defining
semibounded quantum dynamics.  This section gives a finite Gram
compiler for that operator and proves a uniform resolvent comparison
between its exact logarithm and the discrete Dirichlet generator.
The comparison requires no regulator-uniform lower spectral floor.

\subsection{Finite reflected Gram compiler}

Let \(F_1,\ldots,F_N\) be a finite family of future observables and
let \(\Theta\) be the reflected adjoint.  Define
\[
 (G_0)_{ij}:=\omega(\Theta F_i\,F_j),\qquad
 (G_1)_{ij}:=\omega(\Theta F_i\,\tau F_j),                 \tag{59.1}
\]
where \(\tau\) is one positive-time step.

\begin{theorem}[Exact finite transfer criterion]
\label{thm:v14v59-gram}
Assume
\[
 G_0=G_0^*\succ0,\qquad
 0\preceq G_1=G_1^*\preceq G_0.                          \tag{59.2}
\]
On \(\mathbb C^N\) with inner product
\(\langle u,v\rangle_0=u^*G_0v\), define \(T\) by
\[
 \langle u,Tv\rangle_0=u^*G_1v.                           \tag{59.3}
\]
Then \(T\) is a positive self-adjoint contraction.  It is injective
if and only if \(G_1\succ0\).  In an ordinary Euclidean coordinate,
\[
 \widetilde T=G_0^{-1/2}G_1G_0^{-1/2},\qquad
 0\preceq\widetilde T\preceq I.                           \tag{59.4}
\]
\end{theorem}

\begin{proof}
Riesz representation in the positive metric \(G_0\) gives a unique
operator \(T\).  The Hermiticity of \(G_1\) gives self-adjointness in
that metric, and
\[
 \langle u,Tu\rangle_0=u^*G_1u\ge0,\qquad
 \langle u,(I-T)u\rangle_0=u^*(G_0-G_1)u\ge0.
\]
Thus \(0\preceq T\preceq I\).  Conjugation by \(G_0^{1/2}\) gives
\((59.4)\).  Since \(G_0\) is invertible,
\(\ker T=\ker G_1\), proving the final assertion.
\end{proof}

\begin{remark}[Null spaces]
\label{rem:v14v59-null}
If the initial reflected Gram is only semidefinite, first quotient by
\(\ker G_0\).  The one-step form descends exactly when
\(G_1\) annihilates \(\ker G_0\) in both slots.  The theorem then
applies to any basis of the quotient.  This is the finite OS null
quotient, not an additional physical assumption.
\end{remark}

\subsection{Exact logarithm versus the discrete transfer form}

Fix a time step \(a>0\).  For an injective positive contraction
\(T\), functional calculus defines
\[
 H_T:=-a^{-1}\log T,\qquad
 A_T:=a^{-1}(I-T).                                        \tag{59.5}
\]
The first operator can be unbounded in a cofinal limit even though
each finite matrix is bounded.

\begin{theorem}[Floor-free resolvent comparison]
\label{thm:v14v59-resolvent}
Let \(0<a\le1\) and let \(T\) be an injective positive self-adjoint
contraction.  Then
\[
 0\preceq A_T\preceq H_T                                 \tag{59.6}
\]
in quadratic-form order, and
\[
 \left\|(A_T+I)^{-1}-(H_T+I)^{-1}\right\|_{\rm op}
 \le \frac{a}{1-e^{-1}}.                                  \tag{59.7}
\]
No lower bound \(T\succeq mI\) occurs in \((59.7)\).
\end{theorem}

\begin{proof}
For a spectral value \(\lambda=e^{-t}\), \(t\ge0\),
\[
 1-e^{-t}\le t,
\]
which gives \((59.6)\).  The scalar difference of the two resolvents
is nonnegative and equals
\[
 d_a(t)
 =
 \frac{a\{t-1+e^{-t}\}}
 {(a+1-e^{-t})(a+t)}.                                    \tag{59.8}
\]
If \(0\le t\le1\), then
\[
 1-e^{-t}\ge\frac t2,\qquad
 t-1+e^{-t}\le\frac{t^2}{2}.
\]
Therefore
\[
 d_a(t)\le
 \frac{at^2/2}{(a+t/2)(a+t)}\le a.                        \tag{59.9}
\]
If \(t\ge1\), discard the second resolvent and use
\[
 d_a(t)\le\frac{a}{a+1-e^{-t}}
 \le\frac{a}{1-e^{-1}}.                                  \tag{59.10}
\]
Functional calculus and the supremum of
\((59.9)\)--\((59.10)\) prove \((59.7)\).
\end{proof}

\begin{corollary}[Cofinal logarithm from Dirichlet-form convergence]
\label{cor:v14v59-cofinal}
Let \(a_n\downarrow0\) and let \(T_n\) be injective positive
self-adjoint contractions on a common Hilbert space.  Put
\[
 H_n=-a_n^{-1}\log T_n,\qquad
 A_n=a_n^{-1}(I-T_n).                                    \tag{59.11}
\]
If
\[
 (A_n+I)^{-1}\longrightarrow R
 \quad\hbox{strongly},                                   \tag{59.12}
\]
then
\[
 (H_n+I)^{-1}\longrightarrow R
 \quad\hbox{strongly}.                                   \tag{59.13}
\]
If \(R=(H+I)^{-1}\) for a nonnegative self-adjoint \(H\), then both
sequences converge to \(H\) in the strong-resolvent sense.
\end{corollary}

\begin{proof}
Theorem~\ref{thm:v14v59-resolvent} gives
\[
 \|(A_n+I)^{-1}-(H_n+I)^{-1}\|
 \le\frac{a_n}{1-e^{-1}}\longrightarrow0.
\]
Combine this norm convergence with \((59.12)\).  The last assertion
is the definition of strong-resolvent convergence at the real point
\(-1\).
\end{proof}

\begin{theorem}[Closed-form handoff]
\label{thm:v14v59-mosco}
Suppose the quadratic forms
\[
 {\mathfrak a}_n[\psi]
 :=a_n^{-1}\langle\psi,(I-T_n)\psi\rangle                 \tag{59.14}
\]
are placed on a common Hilbert space and Mosco-converge to a densely
defined closed nonnegative form \(\mathfrak a\).  Let \(H\ge0\) be
the self-adjoint operator represented by \(\mathfrak a\).  Then the
exact logarithmic generators \(H_n=-a_n^{-1}\log T_n\) converge to
\(H\) in strong resolvent sense.
\end{theorem}

\begin{proof}
Mosco convergence of closed nonnegative forms is equivalent to
strong convergence of the resolvents of their represented operators.
The operator representing \((59.14)\) is \(A_n\).  Apply
Corollary~\ref{cor:v14v59-cofinal}.
\end{proof}

\subsection{What this route does and does not require}

\begin{proposition}[Semiboundedness without a classical Hessian floor]
\label{prop:v14v59-no-hessian}
Under the finite Gram conditions \((59.2)\) with \(G_1\succ0\), the
operator \(H_T=-a^{-1}\log T\) is self-adjoint and nonnegative,
irrespective of the inertia of any separately proposed classical
action Hessian.  A classical Hessian enters only through an
additional symbol-matching theorem.
\end{proposition}

\begin{proof}
Theorem~\ref{thm:v14v59-gram} places the spectrum of \(T\) in
\((0,1]\).  The real function \(-a^{-1}\log\lambda\) is nonnegative
there, so spectral calculus gives the conclusion.
\end{proof}

\begin{proposition}[The gap remains a separate row]
\label{prop:v14v59-gap}
The hypotheses of Theorem~\ref{thm:v14v59-mosco} do not imply a
positive spectral gap above the ground state.
\end{proposition}

\begin{proof}
On \(\ell^2(\mathbb N_0)\), let
\[
 He_j=\frac1{j+1}e_j
\]
and set \(T_n=e^{-a_nH}\).  Then every \(T_n\) is an injective
positive contraction and its exact logarithmic generator is \(H\);
hence all convergence assertions hold.  But
\(\inf\sigma(H)=0\) and zero is an accumulation point, so there is no
positive gap above zero.
\end{proof}

\begin{assessment}[Relation to the earlier transfer cards]
\label{ass:v14v59-relation}
V39 supplied the finite positive-transfer logarithm and already
distinguished injectivity from a uniform transfer floor.  V43
supplied exact quadratic coherent symbols on a reduced bosonic Fock
space.  V59 adds the missing cofinal bridge: the exact logarithm has
the same strong-resolvent limit as the simpler discrete form
\(a_n^{-1}(I-T_n)\), with a norm-resolvent error bounded only by the
time step.  This permits a semibounded continuum generator even when
the smallest finite transfer eigenvalue tends to zero.

The load-bearing acquisition is now the pair of reflected Grams
\((G_{0,n},G_{1,n})\), their quotient compatibility, and Mosco
compactness of \((59.14)\).  None is equivalent to positivity of the
classical DeWitt Hessian.  A later correspondence theorem must still
show that the resulting generator has the intended Einstein--SM
classical limit.
\end{assessment}

\begin{decision}[V60 compulsory audit and acquisition test]
\label{dec:v14v59-next}
At V60 inspect the live Yang--Mills and Navier--Stokes companions.
Then compare the attached transfer and Grand Tensor records with the
finite Gram packet \((59.1)\): identify which of \(G_0\), \(G_1\),
injectivity, common embeddings, and Mosco bounds are actually
populated.  Continue in V61 with the first missing row, going upstream
to the protected finite path law if the one-step Gram is absent.
\end{decision}

\begin{firewall}[Claim boundary after V59]
\label{fire:v14v59}
The transfer-form route proves a conditional semibounded generator
and a floor-free cofinal comparison.  The required SM--Einstein
reflected Grams and their cofinal convergence have not been
populated, and no spectral gap, interacting continuum measure, or
classical Einstein correspondence follows.
\end{firewall}

\section*{Version V59 append-only record}
\addcontentsline{toc}{section}{Version V59 append-only record}
The complete V58 source was retained before this continuation.  It
had \(739567\) bytes, \(22584\) lines, and SHA--256
\[
 \texttt{1AD41098D4E68CDA311125084968836D1F8210DCC5396D0753ACB7BA3A015394}.
\]
V59 constructed the finite reflected-Gram transfer compiler and
proved the regulator-floor-free resolvent comparison.
\section{Fourteenth-cycle V60: compulsory audit and reflected-transfer acquisition ledger}
\label{sec:v14v60}

This is the compulsory companion audit after V55.  It is followed by
a source-by-source comparison with the V59 reflected-transfer packet.
The comparison distinguishes a proved general compiler from a
populated active Standard Model--Einstein instance.

\subsection{Live companion audit}

\begin{record}[Companion files inspected at V60]
\label{rec:v14v60-companions}
At the audit time the shared folder contained the two nominal active
Yang--Mills files
\[
\begin{split}
 &\path{Renewal_Yang_Mills_Mass_Gap_Research_Notes_V19.tex},\\
 &\path{Renewal_Yang_Mills_Mass_Gap_Research_Notes_V20.tex}.
\end{split}
\]
Each had \(152282\) bytes, \(4140\) lines, and SHA--256
\[
 \texttt{B6C7E17219764E932E3798860D58EA77CEC5586337748EE61381E29FE945CE56}.
                                                               \tag{60.1}
\]
They were byte-for-byte identical, and the nominal V20 file still
ended with the V19 decision to perform its V20 audit and direct-grid
calculation.  Thus no V20 mathematical appendix was present to read.

The active Navier--Stokes file remained
\[
 \path{Renewal_NS_Stability_Research_Notes_V26.tex},
\]
with \(278283\) bytes, \(5319\) lines, and SHA--256
\[
 \texttt{82C887F8C2C69E4F28FAD7C3DC8DE5B9099CB5A4BF431EBA1FFF3E91935978AD}.
                                                               \tag{60.2}
\]
It is unchanged from V55.
\end{record}

\begin{decision}[Companion import decision]
\label{dec:v14v60-companion}
No new companion result is imported.  Yang--Mills V19's direct
relative-box method was already imported at V55.  Navier--Stokes V26
still concerns rank-one Oseen-kernel constants and supplies no
reflected Standard Model--Einstein path law.  The duplicate nominal
YM file is treated only as a provenance fact and is not modified.
\end{decision}

\subsection{Comparison with the attached source programmes}

Let the V59 active packet be
\[
 {\cal P}_{\rm tr}
 =(G_{0,n},G_{1,n},a_n,J_n,\mathfrak a_n)_{n\ge1},        \tag{60.3}
\]
where the first two entries are the zero-step and one-step reflected
Grams, \(a_n\) is the time step, \(J_n\) is the physical Hilbert
embedding, and
\[
 \mathfrak a_n[u]
 =a_n^{-1}\langle u,(I-T_n)u\rangle.                      \tag{60.4}
\]

\begin{theorem}[Exact acquisition ledger]
\label{thm:v14v60-ledger}
The attached Grand Tensor and Einstein--SM sources establish the
following logical ledger for \((60.3)\).
\[
\begin{array}{c|c|p{8.0cm}}
\text{row}&\text{status}&\text{source conclusion}\\ \hline
G_{0,n}&\text{conditional compiler}&
The Grand Tensor finite reflection-Gram theorem reduces reflection
positivity of a declared finite absolute cylinder to positivity of
one explicit Gram.  No numerical active SM--Einstein cylinder weight
is supplied.\\
G_{1,n}&\text{missing}&
The attached Einstein--SM V8 transfer-log acquisition asks for the
positive transfer on a full or independently reducing carrier and
states that this active card has not been populated.\\
\ker G_{1,n}=\{0\}&\text{missing}&
No active one-step matrix is present, hence neither injectivity nor a
supported zero-mode line is determined.\\
J_n&\text{conditional architecture}&
The Grand Tensor Mosco theorem allows canonical isometries of a
strict Hilbert direct system, but the active reflected OS quotients
and their intertwining embeddings are not supplied.\\
\mathfrak a_n\to\mathfrak a&\text{conditional theorem}&
Grand Tensor proves Mosco--resolvent--semigroup equivalence for forms
already placed in that direct system.  It does not prove Mosco
bounds for the active SM--Einstein transfer forms.\\
\text{field measure}&\text{conditional theorem}&
The reflection-positive projective-limit theorem requires finite
absolute weights, summable total-variation defects, and tightness.
Those are hypotheses, not populated active data.\\
\text{Einstein correspondence}&\text{missing}&
The finite action theorem gives Ward, stress, and stationarity
identities, but no theorem identifies the V59 limit generator with
the canonical Einstein--SM Hamiltonian.
\end{array}                                                \tag{60.5}
\]
In particular, no source theorem permits \(G_{1,n}\) or
\(\mathfrak a_n\) to be inferred from the static source Grams,
Duhamel Grams, or expectation-level Einstein residuals.
\end{theorem}

\begin{proof}
The Grand Tensor theorem labelled
\(\texttt{thm:SMST-reflection-Gram}\) starts with a finite nonnegative
absolute weight and returns the finite OS Gram criterion.  Its
projective-limit theorem assumes each such measure is reflection
positive and assumes a summable total-variation defect.  The theorem
labelled \(\texttt{thm:GT-Mosco}\) starts with closed forms on a
strict Hilbert direct system and proves equivalence with embedded
resolvent and semigroup convergence.  These theorems prove the three
conditional compiler entries of \((60.5)\).

The attached Einstein--SM theorem labelled
\(\texttt{thm:V8-transfer-log}\) reconstructs a generator from
positive transfer matrices.  Its acquisition block explicitly lists
those matrices, the physical connection, a reducing carrier, a
spectral floor, and a held-out lag as inputs.  Its integrated theorem
states that these rows have not been populated on the original active
carrier.  Therefore \(G_{1,n}\) and injectivity remain missing.

A static Gram is an equal-time quadratic form.  A one-step Gram is a
joint two-time moment.  Equal one-time marginals admit different
couplings, so the former cannot determine the latter.  Likewise an
expectation-level first-variation identity does not determine a
two-time kernel or a closed-form liminf inequality.  This proves the
last assertion and completes the ledger.
\end{proof}

\begin{counterexample}[A static OS Gram does not determine transfer]
\label{ctr:v14v60-static}
Let the one-time law on \(\{0,1\}\) be uniform.  The two reversible
transition matrices
\[
 P_{\rm id}=\begin{pmatrix}1&0\\0&1\end{pmatrix},\qquad
 P_{\rm mix}=\begin{pmatrix}1/2&1/2\\1/2&1/2\end{pmatrix} \tag{60.6}
\]
have the same static Gram
\[
 G_0=\frac12I,
\]
but their one-step Grams are
\[
 G_1^{\rm id}=\frac12P_{\rm id},\qquad
 G_1^{\rm mix}=\frac12P_{\rm mix}.                        \tag{60.7}
\]
The first transfer is injective and the second has a zero mode.
\end{counterexample}

\begin{proof}
Both chains preserve the uniform law and satisfy detailed balance.
In the point-mass basis the equal-time inner product is
\(\frac12I\), while the joint one-step law is
\(\operatorname{diag}(\pi)P=\frac12P\).  The ranks in \((60.7)\)
are two and one.
\end{proof}

\begin{assessment}[First missing load-bearing row]
\label{ass:v14v60-first-missing}
The earliest missing row is not Mosco convergence.  It is the finite
two-time object \(G_{1,n}\) on the same reflected quotient as
\(G_{0,n}\).  Attempting compactness before that row exists would be
circular, because there is no transfer form to compactify.

The Grand Tensor finite cylinder architecture suggests the correct
upstream object: one protected, reflection-compatible finite path law
whose one-time and adjacent-time marginals are both retained.  The
one-step Gram must be read from that joint law, and a three-time
marginal must test its Markov or semigroup composition rather than
assuming it.
\end{assessment}

\begin{decision}[V61 finite path-law compiler]
\label{dec:v14v60-next}
Construct the exact \(G_0,G_1\) compiler for a finite stationary
reversible path law.  Prove the equivalence between transfer
positivity and a positive adjacent-time joint kernel, identify
injectivity, and use a three-time marginal as a held-out Markov
identity.  Then state the precise finite acquisition card required
from the active SM--Einstein regulator.
\end{decision}

\begin{firewall}[Claim boundary after V60]
\label{fire:v14v60}
The attached sources contain general reflection-positive and Mosco
theorems, but the active reflected one-step SM--Einstein Gram is
absent.  No transfer form, cofinal generator, Einstein
correspondence, or full SM--Einstein continuum has been populated.
\end{firewall}

\section*{Version V60 append-only record}
\addcontentsline{toc}{section}{Version V60 append-only record}
The complete V59 source was retained before this continuation.  It
had \(748678\) bytes, \(22841\) lines, and SHA--256
\[
 \texttt{48F56CEAF2AA3A9C81E841F0ED4C69F9ADF5A839C6E216EA709E02F06C90CC7F}.
\]
V60 performed the compulsory companion audit and located the first
missing active transfer row.
\section{Fourteenth-cycle V61: exact one-step Grams from a finite reversible path law}
\label{sec:v14v61}

V60 found that the active programme lacks the adjacent-time Gram
\(G_1\).  The minimal upstream object that determines it is a
two-time joint law, not another static source metric.  We now give the
exact compiler, its positivity test, and a three-time held-out
identity.

\subsection{Two-time compiler}

Let \(S=\{1,\ldots,m\}\), let
\(\pi_i>0\) with \(\sum_i\pi_i=1\), and let
\(J=(J_{ij})\) be a real matrix satisfying
\[
 J_{ij}\ge0,\qquad J=J^{\mathsf T},\qquad
 \sum_jJ_{ij}=\pi_i.                                      \tag{61.1}
\]
Thus \(J\) is a reflection-symmetric adjacent-time joint law with
one-time marginal \(\pi\).  Put
\[
 D=\operatorname{diag}(\pi_1,\ldots,\pi_m),\qquad
 P=D^{-1}J.                                               \tag{61.2}
\]

\begin{theorem}[Adjacent-law transfer compiler]
\label{thm:v14v61-compiler}
The matrix \(P\) is a stationary reversible Markov transition:
\[
 P\mathbf1=\mathbf1,\qquad \pi^{\mathsf T}P=\pi^{\mathsf T},
 \qquad DP=P^{\mathsf T}D.                               \tag{61.3}
\]
On \(L^2(\pi)\),
\[
 \langle f,g\rangle_\pi=f^*Dg,
\]
the one-time and one-step Grams in the point basis are
\[
 G_0=D,\qquad G_1=J,                                     \tag{61.4}
\]
and the transfer is \(T=P\).  Moreover,
\[
 0\preceq T\preceq I\quad\hbox{on }L^2(\pi)
 \iff J\succeq0.                                         \tag{61.5}
\]
It is injective exactly when \(J\succ0\), and its transfer floor is
\[
 g_*=\lambda_{\min}\!\left(D^{-1/2}JD^{-1/2}\right).      \tag{61.6}
\]
\end{theorem}

\begin{proof}
The row-sum identity in \((61.1)\) gives \(P\mathbf1=\mathbf1\).
Symmetry of \(J=DP\) is detailed balance.  Summing detailed balance in
the first coordinate gives stationarity.

For point observables,
\[
 \mathbb E[\overline{f(X_0)}g(X_0)]=f^*Dg,\qquad
 \mathbb E[\overline{f(X_0)}g(X_1)]=f^*Jg,
\]
which proves \((61.4)\) and identifies \(T=P\).
The whitened transfer is
\[
 \widetilde T=D^{1/2}PD^{-1/2}
 =D^{-1/2}JD^{-1/2}.                                     \tag{61.7}
\]
Thus \(T\succeq0\) exactly when \(J\succeq0\).

For every \(f\),
\[
 f^*(D-J)f
 =\frac12\sum_{i,j}J_{ij}|f_i-f_j|^2\ge0,                \tag{61.8}
\]
using symmetry and the two marginal identities.  Hence \(T\preceq I\).
Since \(D\) is invertible,
\(\ker T=\ker J\).  Formula \((61.6)\) is the least eigenvalue of the
whitened representation \((61.7)\).
\end{proof}

\begin{corollary}[Exact discrete transfer form]
\label{cor:v14v61-dirichlet}
For a physical time step \(a>0\), the V59 form is
\[
 \mathfrak a_{J,a}[f]
 =\frac1a\langle f,(I-T)f\rangle_\pi
 =\frac1{2a}\sum_{i,j}J_{ij}|f_i-f_j|^2.                 \tag{61.9}
\]
If \(J\succ0\), the exact logarithmic generator
\[
 H_{J,a}=-a^{-1}\log P                                  \tag{61.10}
\]
is nonnegative and self-adjoint on \(L^2(\pi)\), and V59 gives
\[
 \left\|(I+\mathfrak A_{J,a})^{-1}
       -(I+H_{J,a})^{-1}\right\|
 \le\frac{a}{1-e^{-1}},                                  \tag{61.11}
\]
where \(\mathfrak A_{J,a}=a^{-1}(I-P)\).
\end{corollary}

\begin{proof}
Equation \((61.9)\) is \((61.8)\) divided by \(a\).
If \(J\succ0\), Theorem~\ref{thm:v14v61-compiler} makes \(P\) an
injective positive contraction.  Apply
Theorem~\ref{thm:v14v59-resolvent}.
\end{proof}

\begin{counterexample}[Reversibility does not imply a positive transfer]
\label{ctr:v14v61-flip}
For the uniform law on two states, let
\[
 P=\begin{pmatrix}0&1\\1&0\end{pmatrix},\qquad
 J=\frac12P.                                              \tag{61.12}
\]
This is a stationary reversible Markov chain and its Dirichlet form
\((61.9)\) is nonnegative.  Nevertheless \(P\) has eigenvalues
\(1,-1\), so \(J\not\succeq0\) and no self-adjoint real principal
logarithm \(-a^{-1}\log P\) exists.
\end{counterexample}

\begin{proof}
All claims follow by direct diagonalization.  The negative eigenvalue
violates the necessary spectral condition for the real logarithm in
\((61.10)\).
\end{proof}

\subsection{Compression and exact coarse graining}

Let \(W:\mathbb C^r\to\mathbb C^m\) synthesize a declared observable
subspace.  Its two Grams are
\[
 G_0^W=W^*DW,\qquad G_1^W=W^*JW.                         \tag{61.13}
\]

\begin{proposition}[When the compressed logarithm is licensed]
\label{prop:v14v61-compression}
Assume \(G_0^W\succ0\).  The compressed transfer form is positive
when \(G_1^W\succeq0\), but its logarithm equals the compression of
the full logarithm only if \(\operatorname{Ran}W\) reduces \(P\) in
the \(D\)-metric.  Equivalently, for the \(D\)-orthogonal projection
\[
 \Pi_W=W(W^*DW)^{-1}W^*D,                                 \tag{61.14}
\]
one must have
\[
 \Pi_WP=P\Pi_W.                                           \tag{61.15}
\]
\end{proposition}

\begin{proof}
Congruence preserves positivity, giving the first statement.
Condition \((61.15)\) makes \(P\) block diagonal relative to
\(\operatorname{Ran}W\) and its \(D\)-orthogonal complement.
Functional calculus then commutes with compression.  Conversely, an
arbitrary nonreducing compression is subject to the V8
no-compressed-log counterexample and is not licensed by its two
compressed Grams alone.
\end{proof}

\begin{proposition}[Reflection-compatible lumping preserves transfer positivity]
\label{prop:v14v61-lumping}
Let \(L:\mathbb C^{\bar m}\to\mathbb C^m\) pull a coarse observable
back to the fine state space.  Define
\[
 \bar D=L^*DL,\qquad \bar J=L^*JL.                        \tag{61.16}
\]
If \(J\succeq0\), then \(\bar J\succeq0\); if
\(D-J\succeq0\), then \(\bar D-\bar J\succeq0\).
Thus exact marginalization or deterministic reflection-compatible
coarse graining preserves the finite transfer inequalities.
\end{proposition}

\begin{proof}
Both assertions are preservation of Loewner order under congruence
by \(L\).
\end{proof}

\subsection{Three-time held-out law}

Let \(K_{ijk}\ge0\) be a stationary three-time law whose two adjacent
marginals equal \(J\):
\[
 \sum_kK_{ijk}=J_{ij},\qquad
 \sum_iK_{ijk}=J_{jk}.                                   \tag{61.17}
\]
Define the Markov reconstruction
\[
 K^{\rm M}_{ijk}:=\frac{J_{ij}J_{jk}}{\pi_j}              \tag{61.18}
\]
and the exact residual
\[
 \Delta_{\rm M}:=\sum_{i,j,k}
 \left|K_{ijk}-K^{\rm M}_{ijk}\right|^2.                  \tag{61.19}
\]

\begin{theorem}[Held-out Markov and lag-two identities]
\label{thm:v14v61-three-time}
One has
\[
 \Delta_{\rm M}=0
 \iff
 K_{ijk}=\pi_iP_{ij}P_{jk}\quad\hbox{for all }i,j,k.       \tag{61.20}
\]
On this branch the endpoint two-step Gram is
\[
 (G_2)_{ik}:=\sum_jK_{ijk}
 =DP^2
 =G_1G_0^{-1}G_1.                                        \tag{61.21}
\]
Thus \((61.19)\) is a full conditional-independence test and
\((61.21)\) is its weaker lag-two matrix consequence.
\end{theorem}

\begin{proof}
Substitute \(J=DP\) into \((61.18)\):
\[
 \frac{(\pi_iP_{ij})(\pi_jP_{jk})}{\pi_j}
 =\pi_iP_{ij}P_{jk}.
\]
A sum of squares vanishes exactly when every summand vanishes, proving
\((61.20)\).  Summing the factorized law over \(j\) gives
\(DP^2\), and
\[
 G_1G_0^{-1}G_1=(DP)D^{-1}(DP)=DP^2.
\]
The endpoint identity alone sums over the intermediate state and
therefore cannot recover every triple entry, so it is weaker.
\end{proof}

\begin{counterexample}[Lag two does not prove the full Markov law]
\label{ctr:v14v61-lag-two}
There exist nonzero arrays \(R_{ijk}\) with all one-coordinate sums
zero:
\[
 \sum_iR_{ijk}=\sum_jR_{ijk}=\sum_kR_{ijk}=0.             \tag{61.22}
\]
For sufficiently small real \(\epsilon\),
\[
 K=K^{\rm M}+\epsilon R                                  \tag{61.23}
\]
remains nonnegative and has the same adjacent and endpoint marginals
as \(K^{\rm M}\), but \(\Delta_{\rm M}>0\).
\end{counterexample}

\begin{proof}
On two states take
\[
 R_{ijk}=(-1)^{i+j+k},
\]
after indexing states by \(0,1\).  Every sum over one index vanishes.
If the Markov reference law has full support, sufficiently small
\(|\epsilon|\) preserves nonnegativity.  All two-coordinate
marginals are unchanged by \((61.22)\), while \(R\ne0\), so
\((61.19)\) is positive.
\end{proof}

\subsection{Finite active acquisition card}

\begin{acquisition}[Minimal reflected path-law card]
\label{acq:v14v61}
At cutoff \(n\), the active SM--Einstein programme must provide:
\begin{enumerate}
\item a finite reflected boundary state set after the gauge and OS
      null quotients, and its faithful marginal \(\pi_n\);
\item the physically normalized adjacent-time joint matrix \(J_n\),
      with exact reflection symmetry and marginal identities;
\item an exact or outward-certified test
      \(J_n\succeq0\), together with either \(J_n\succ0\) or an
      explicitly retained zero-mode support;
\item the physical time step \(a_n\);
\item a held-out three-time law \(K_n\) and the residual
      \(\Delta_{{\rm M},n}\);
\item the refinement pullback and its residuals against
      \(D_n,J_n,K_n\).
\end{enumerate}
For rational finite data, symmetry, marginalization, positivity, and
the Markov residual are decidable by exact arithmetic.  The output is
the transfer \(P_n=D_n^{-1}J_n\), the form \((61.9)\), its exact
logarithm when injective, and the complete finite obstruction witness
otherwise.
\end{acquisition}

\begin{assessment}[Progress and next bottleneck]
\label{ass:v14v61-progress}
V61 reduces the first missing V60 row to one physical adjacent-time
joint law.  It also prevents two shortcuts: a reversible Markov law
can have a negative transfer eigenvalue, and a lag-two identity can
hold while the full three-time law is non-Markov.  The active sources
supply neither \(J_n\) nor \(K_n\), so no numerical transfer has been
created.

Once the path-law card is populated, the next independent problem is
cofinal: prove Mosco compactness of the exact edge form
\((61.9)\) under refinement.  Before that, one can derive sharp
transport bounds showing how errors in \((\pi,J)\) change the
whitened transfer and its resolvent.
\end{assessment}

\begin{decision}[V62 robustness before compactness]
\label{dec:v14v61-next}
Derive a support-safe perturbation theorem for
\[
 \widetilde T=D^{-1/2}JD^{-1/2}
\]
under outward errors in \(\pi\) and \(J\).  Propagate the bound to the
transfer positivity margin, the discrete Dirichlet resolvent, and
the held-out Markov residual.  This is required before approximate
physical path-law data can enter the exact V59 cofinal theorem.
\end{decision}

\begin{firewall}[Claim boundary after V61]
\label{fire:v14v61}
The finite path-law compiler is exact, but no active SM--Einstein
\(\pi_n,J_n,K_n\) has been acquired.  Transfer positivity, injectivity,
Mosco convergence, nontriviality, Einstein correspondence, and the
full continuum remain open.
\end{firewall}

\section*{Version V61 append-only record}
\addcontentsline{toc}{section}{Version V61 append-only record}
The complete V60 source was retained before this continuation.  It
had \(757217\) bytes, \(23045\) lines, and SHA--256
\[
 \texttt{2AF517A0B1F399060E6AE4638959F4BFD26B97259DC5A310097C35381B5C54E2}.
\]
V61 compiled the one-step reflected transfer and its held-out
three-time test from a finite reversible path law.
\section{Fourteenth-cycle V62: support-safe path-law perturbation and resolvent transport}
\label{sec:v14v62}

The exact V61 compiler assumes a faithful marginal and an exact
adjacent joint law.  A physical finite card will generally arrive
with outward error bounds.  Because whitening contains
\(D^{-1/2}\), marginal support must be controlled before any transfer
error is quoted.  This section gives the required support-safe
estimate and propagates it to the exact transfer logarithm.

\subsection{Whitened transfer perturbation}

Let
\[
 D=\operatorname{diag}(\pi),\qquad
 \widehat D=\operatorname{diag}(\widehat\pi),\qquad
 T^\circ=D^{-1/2}JD^{-1/2},\qquad
 \widehat T^\circ=\widehat D^{-1/2}\widehat J\widehat D^{-1/2}.
                                                               \tag{62.1}
\]
The circles indicate Euclidean whitened representatives.

\begin{theorem}[Support-safe whitening bound]
\label{thm:v14v62-whitening}
Assume
\[
 \min_i\{\pi_i,\widehat\pi_i\}\ge p_*>0,\qquad
 \|\widehat\pi-\pi\|_\infty\le\delta_\pi,\qquad
 \|\widehat J-J\|_{\rm op}\le\delta_J,\qquad
 \|J\|_{\rm op}\le j_* .                                  \tag{62.2}
\]
Then
\[
 \boxed{
 \|\widehat T^\circ-T^\circ\|_{\rm op}
 \le
 \varepsilon_T:=
 \frac{\delta_J}{p_*}
 +\frac{j_*\delta_\pi}{p_*^2}.}                           \tag{62.3}
\]
If
\[
 g_*I\preceq T^\circ\preceq I,
\]
then
\[
 (g_*-\varepsilon_T)I
 \preceq\widehat T^\circ
 \preceq(1+\varepsilon_T)I.                               \tag{62.4}
\]
When the submitted pair \((\widehat\pi,\widehat J)\) itself satisfies
the exact joint-law identities of V61, its upper bound improves to
\(\widehat T^\circ\preceq I\).  In particular,
\[
 \varepsilon_T<g_*                                       \tag{62.5}
\]
certifies submitted injectivity.
\end{theorem}

\begin{proof}
On \([p_*,\infty)\), the scalar derivative of \(x^{-1/2}\) has
absolute value at most \(1/(2p_*^{3/2})\).  Since the two matrices are
diagonal,
\[
 \|\widehat D^{-1/2}-D^{-1/2}\|
 \le\frac{\delta_\pi}{2p_*^{3/2}},\qquad
 \|D^{-1/2}\|,\|\widehat D^{-1/2}\|\le p_*^{-1/2}.         \tag{62.6}
\]
Insert and subtract the exact joint matrix:
\[
\begin{split}
 \widehat T^\circ-T^\circ
 ={}&\widehat D^{-1/2}(\widehat J-J)\widehat D^{-1/2}\\
 &+(\widehat D^{-1/2}-D^{-1/2})J\widehat D^{-1/2}\\
 &+D^{-1/2}J(\widehat D^{-1/2}-D^{-1/2}).
                                                               \tag{62.7}
\end{split}
\]
Using \((62.6)\) in \((62.7)\) gives
\[
 \frac{\delta_J}{p_*}
 +2\frac{\delta_\pi}{2p_*^{3/2}}
      j_*p_*^{-1/2},
\]
which is \((62.3)\).  Weyl's inequality gives \((62.4)\).
If the submitted data are an exact symmetric joint law,
Theorem~\ref{thm:v14v61-compiler} gives the sharper upper bound.
The lower edge in \((62.4)\) is positive under \((62.5)\).
\end{proof}

\begin{counterexample}[Why the support floor is load-bearing]
\label{ctr:v14v62-support}
Let
\[
 D_\epsilon=\begin{pmatrix}\epsilon&0\\0&1-\epsilon\end{pmatrix},
 \qquad
 J_\epsilon=D_\epsilon.
\]
Perturb only the first joint entry by an amount of order \(\epsilon\).
The unwhitened perturbation tends to zero, while the corresponding
first whitened entry changes by order one.  Hence no bound of the
form \(\|\Delta T^\circ\|\le C\|\Delta J\|\), with \(C\) independent
of the marginal support floor, is possible.
\end{counterexample}

\begin{proof}
The first whitened entry is \(J_{11}/\epsilon\).
Changing \(J_{11}\) by \(c\epsilon\) changes this entry by \(c\),
although the matrix perturbation has norm \(|c|\epsilon\to0\).
\end{proof}

\subsection{Dirichlet and exact-log resolvents}

Assume from now on that both exact and submitted data obey the V61
joint-law identities and have positive transfer matrices.  For the
same time step \(a\), define
\[
 A=a^{-1}(I-T^\circ),\qquad
 \widehat A=a^{-1}(I-\widehat T^\circ),                   \tag{62.8}
\]
\[
 H=-a^{-1}\log T^\circ,\qquad
 \widehat H=-a^{-1}\log\widehat T^\circ.                  \tag{62.9}
\]

\begin{theorem}[Measured-card resolvent radius]
\label{thm:v14v62-resolvent}
Under the hypotheses above,
\[
 \|(A+I)^{-1}-(\widehat A+I)^{-1}\|
 \le\frac{\varepsilon_T}{a},                              \tag{62.10}
\]
and
\[
 \boxed{
 \|(H+I)^{-1}-(\widehat H+I)^{-1}\|
 \le
 \frac{2a}{1-e^{-1}}+\frac{\varepsilon_T}{a}.}             \tag{62.11}
\]
Thus a cofinal measured path-law approximation is resolvent accurate
provided
\[
 a_n\longrightarrow0,\qquad
 \frac{\varepsilon_{T,n}}{a_n}\longrightarrow0.           \tag{62.12}
\]
\end{theorem}

\begin{proof}
Both \(A,\widehat A\) are nonnegative.  The resolvent identity gives
\[
 (A+I)^{-1}-(\widehat A+I)^{-1}
 =(A+I)^{-1}(\widehat A-A)(\widehat A+I)^{-1}.
\]
The two resolvents have norm at most one, while
\[
 \|\widehat A-A\|
 =a^{-1}\|\widehat T^\circ-T^\circ\|
 \le a^{-1}\varepsilon_T,
\]
which proves \((62.10)\).  Apply V59's floor-free comparison once to
\((A,H)\) and once to \((\widehat A,\widehat H)\), then insert
\((62.10)\) between them.  The triangle inequality gives
\((62.11)\).  Condition \((62.12)\) makes its right side vanish.
\end{proof}

\begin{corollary}[Robust cofinal handoff]
\label{cor:v14v62-cofinal}
Suppose the submitted discrete forms represented by \(\widehat A_n\)
Mosco-converge, after the declared Hilbert embeddings, to a closed
nonnegative form with operator \(H_\infty\).  If \((62.12)\) holds,
then the exact logarithmic generators \(H_n\) and the submitted
logarithmic generators \(\widehat H_n\) both converge to
\(H_\infty\) in the strong-resolvent sense.
\end{corollary}

\begin{proof}
Mosco convergence gives strong convergence of
\((\widehat A_n+I)^{-1}\).  Equation \((62.10)\) transfers that limit
to \(A_n\), and \((62.11)\), together with the individual V59
comparison, transfers it to both logarithmic generators.
\end{proof}

\subsection{Three-time residual under interval data}

Let the state set have \(N\) elements.  Suppose \(K,\widehat K\) and
\(J,\widehat J\) are normalized nonnegative arrays with
\[
 \min_j\{\pi_j,\widehat\pi_j\}\ge p_*,
\]
\[
 \|\widehat J-J\|_{\max}\le\delta_J^\infty,\qquad
 \|\widehat K-K\|_{\max}\le\delta_K^\infty,\qquad
 \|\widehat\pi-\pi\|_\infty\le\delta_\pi.                 \tag{62.13}
\]
Let \(K^{\rm M}\) and \(\widehat K^{\rm M}\) be the Markov
reconstructions \((61.18)\).

\begin{proposition}[Robust Markov-residual radius]
\label{prop:v14v62-markov}
One has
\[
 \|\widehat K^{\rm M}-K^{\rm M}\|_{\max}
 \le
 \varepsilon_{\rm M}:=
 \frac{2\delta_J^\infty}{p_*}
 +\frac{\delta_\pi}{p_*^2},                               \tag{62.14}
\]
and
\[
 \left|
 \sqrt{\widehat\Delta_{\rm M}}-\sqrt{\Delta_{\rm M}}
 \right|
 \le
 N^{3/2}\bigl(\delta_K^\infty+\varepsilon_{\rm M}\bigr).  \tag{62.15}
\]
Consequently
\[
 \sqrt{\Delta_{\rm M}}
 \le
 \sqrt{\widehat\Delta_{\rm M}}
 +N^{3/2}\bigl(\delta_K^\infty+\varepsilon_{\rm M}\bigr)  \tag{62.16}
\]
is a certified outward true residual.
\end{proposition}

\begin{proof}
For \(a,b,p>0\), add and subtract one factor at a time:
\[
 \left|\frac{\widehat a\widehat b}{\widehat p}
             -\frac{ab}{p}\right|
 \le
 \frac{|\widehat a-a|\,|\widehat b|}{\widehat p}
 +\frac{|a|\,|\widehat b-b|}{\widehat p}
 +|ab|\left|\frac1{\widehat p}-\frac1p\right|.            \tag{62.17}
\]
All joint probabilities are at most one, and both denominators are
at least \(p_*\).  Applying \((62.17)\) to
\(J_{ij}J_{jk}/\pi_j\) gives \((62.14)\).

Now
\[
 \sqrt{\Delta_{\rm M}}=\|K-K^{\rm M}\|_{\rm F}.
\]
The reverse triangle inequality and the fact that an \(N^3\)-entry
array has Frobenius norm at most \(N^{3/2}\) times its maximum norm
give \((62.15)\), and \((62.16)\) follows.
\end{proof}

\begin{assessment}[Error scale forced by the continuum]
\label{ass:v14v62-scale}
A fixed absolute accuracy in \(J_n\) is insufficient as
\(a_n\to0\).  The transfer is \(I-a_nH+o(a_n)\), so the physically
relevant difference from the identity is itself of order \(a_n\).
Equation \((62.12)\) requires the whitened transfer uncertainty to be
little-oh of the time step.  This is an acquisition scaling law, not
a technical preference.

The estimates are support-safe but intentionally do not repair
violations of symmetry, marginal consistency, or positivity.  Those
are separate finite residuals.  Projecting noisy data onto a desired
positive Markov cone would be a fit unless the projection radius and
all physical normalization effects were retained.
\end{assessment}

\begin{decision}[V63 cofinal edge-form criterion]
\label{dec:v14v62-next}
Use the exact Dirichlet representation
\[
 \mathfrak a_n[f]
 =\frac1{2a_n}\sum_{i,j}J_{n,ij}|f_i-f_j|^2
\]
to state checkable sufficient conditions for Mosco convergence under
a refining sequence of finite boundary spaces.  Separate liminf,
recovery, tightness, and nontriviality.  Include a counterexample
showing that pointwise convergence of transition probabilities does
not prevent energy escape.
\end{decision}

\begin{firewall}[Claim boundary after V62]
\label{fire:v14v62}
The perturbation theorems quantify a supplied path-law card; they do
not supply one.  The active marginal support, adjacent law, required
little-oh error scale, Mosco bounds, and Einstein correspondence
remain unverified.
\end{firewall}

\section*{Version V62 append-only record}
\addcontentsline{toc}{section}{Version V62 append-only record}
The complete V61 source was retained before this continuation.  It
had \(768402\) bytes, \(23373\) lines, and SHA--256
\[
 \texttt{501F749FC8362E6371BA4671DC6751A45372D35ABFD9F9E7E8D8A61FCA7A42C1}.
\]
V62 propagated support-safe path-law errors through whitening,
Dirichlet resolvents, exact transfer logarithms, and the held-out
Markov residual.
\section{Fourteenth-cycle V63: a conductance-level Mosco criterion and soft-mode escape}
\label{sec:v14v63}

The Grand Tensor source already proves the abstract equivalence between
Mosco, resolvent, and semigroup convergence on its strict Hilbert
direct system.  We do not repeat that theorem.  Instead, we supply
the missing specialization: sufficient conductance conditions under
which the V61 adjacent-law forms satisfy its two Mosco clauses.

\subsection{Conforming reference forms}

Let \(\mathcal H\) be a Hilbert space and let
\[
 V_1\subset V_2\subset\cdots\subset\mathcal H             \tag{63.1}
\]
be finite-dimensional subspaces.  Let \(\mathfrak a\) be a densely
defined closed nonnegative form on \(\mathcal H\), and assume
\[
 \bigcup_nV_n\subset{\cal D}(\mathfrak a)
\quad\hbox{is a form core}.                               \tag{63.2}
\]
Write
\[
 \mathfrak a_n^0=\mathfrak a|_{V_n}.                      \tag{63.3}
\]

\begin{theorem}[Perturbed conforming-form Mosco criterion]
\label{thm:v14v63-conforming}
Let \(\mathfrak a_n\) be nonnegative quadratic forms on \(V_n\).
Suppose \(\varepsilon_n,\eta_n\to0\) and, for all \(u\in V_n\),
\[
 |\mathfrak a_n[u]-\mathfrak a_n^0[u]|
 \le
 \varepsilon_n\mathfrak a_n^0[u]+\eta_n\|u\|^2.           \tag{63.4}
\]
Extend \(\mathfrak a_n\) to \(\mathcal H\) by \(+\infty\) outside
\(V_n\).  Then \(\mathfrak a_n\) Mosco-converges to
\(\mathfrak a\).  Consequently the represented operators converge in
strong resolvent sense, and their semigroups converge strongly,
uniformly on compact time intervals.
\end{theorem}

\begin{proof}
For the lower condition, let \(u_n\rightharpoonup u\) and suppose
along a subsequence that
\(\sup_n\mathfrak a_n[u_n]<\infty\); otherwise the assertion is
immediate.  For large \(n\), \(\varepsilon_n<1/2\), and \((63.4)\)
gives
\[
 (1-\varepsilon_n)\mathfrak a[u_n]
 \le\mathfrak a_n[u_n]+\eta_n\|u_n\|^2.                  \tag{63.5}
\]
Weak convergence bounds \(\|u_n\|\), so
\(\sup_n\mathfrak a[u_n]<\infty\).  A closed nonnegative quadratic
form, extended by \(+\infty\), is weakly lower semicontinuous.
Therefore
\[
 \mathfrak a[u]\le\liminf_n\mathfrak a[u_n]
 \le\liminf_n\mathfrak a_n[u_n],                          \tag{63.6}
\]
where the vanishing errors in \((63.4)\) were used in the last step.

For recovery, take \(u\in{\cal D}(\mathfrak a)\).  The core property
gives \(v_k\in\bigcup_nV_n\) with
\[
 \|v_k-u\|^2+\mathfrak a[v_k-u]\longrightarrow0.          \tag{63.7}
\]
Choose an increasing diagonal \(k(n)\) such that
\(v_{k(n)}\in V_n\), \(v_{k(n)}\to u\) in form norm, and
\[
 \varepsilon_n\mathfrak a[v_{k(n)}]
 +\eta_n\|v_{k(n)}\|^2\longrightarrow0.
\]
Then \((63.4)\) gives
\(\mathfrak a_n[v_{k(n)}]\to\mathfrak a[u]\).
This is the recovery clause.  The final assertions are the inherited
Grand Tensor Mosco theorem.
\end{proof}

\subsection{A checkable edge condition}

At cutoff \(n\), suppose the finite path-law form is written in
unoriented edge notation as
\[
 \mathfrak a_n[u]
 =\frac1{a_n}\sum_{e=\{x,y\}}c_{n,e}|u_x-u_y|^2,          \tag{63.8}
\]
where \(c_{n,e}=J_{n,xy}=J_{n,yx}\ge0\).  Let a conforming reference
form on the same \(V_n\) be
\[
 \mathfrak a_n^0[u]
 =\frac1{a_n}\sum_e c^0_{n,e}|u_x-u_y|^2.                 \tag{63.9}
\]

\begin{proposition}[Relative conductance plus degree-tail test]
\label{prop:v14v63-edges}
Split
\[
 c_{n,e}=c^{\rm main}_{n,e}+r_{n,e}.
\]
Assume
\[
 |c^{\rm main}_{n,e}-c^0_{n,e}|
 \le\varepsilon_n c^0_{n,e}                              \tag{63.10}
\]
and, for the symmetric tail,
\[
 \sup_x\frac1{\pi_{n,x}}
 \sum_y\frac{|r_{n,\{x,y\}}|}{a_n}
 \le\zeta_n.                                              \tag{63.11}
\]
Then
\[
 |\mathfrak a_n[u]-\mathfrak a_n^0[u]|
 \le
 \varepsilon_n\mathfrak a_n^0[u]
 +2\zeta_n\|u\|_{L^2(\pi_n)}^2.                           \tag{63.12}
\]
Hence Theorem~\ref{thm:v14v63-conforming} applies when
\(\varepsilon_n,\zeta_n\to0\) and the reference forms are conforming
restrictions of a closed form with core \(\bigcup_nV_n\).
\end{proposition}

\begin{proof}
The main-edge difference is bounded by
\(\varepsilon_n\mathfrak a_n^0[u]\).  For the tail, use
\[
 |u_x-u_y|^2\le2|u_x|^2+2|u_y|^2.
\]
Symmetry gives
\[
 \frac1{a_n}\sum_{\{x,y\}}|r_{n,\{x,y\}}||u_x-u_y|^2
 \le
 2\sum_x|u_x|^2\sum_y\frac{|r_{n,\{x,y\}}|}{a_n}
 \le2\zeta_n\sum_x\pi_{n,x}|u_x|^2,
\]
which is \((63.12)\).
\end{proof}

\begin{remark}[The time-step scaling reappears]
\label{rem:v14v63-step}
Condition \((63.11)\) divides conductance error by \(a_n\).  Thus an
\(o(1)\) probability error need not be a small form error.  This is
the edgewise version of the V62 requirement
\(\varepsilon_{T,n}=o(a_n)\).
\end{remark}

\subsection{Nontriviality and compactness are separate}

Let \(e_n\) denote the normalized constant vacuum in \(V_n\), assumed
compatible under the embeddings.

\begin{proposition}[Protected nonvacuum witness]
\label{prop:v14v63-nontrivial}
Suppose \(u_n\in V_n\) converges strongly to \(u\in\mathcal H\) and
\[
 \langle u_n,e_n\rangle=0,\qquad
 \|u_n\|\ge v_*>0,\qquad
 \sup_n\mathfrak a_n[u_n]<\infty.                         \tag{63.13}
\]
If the vacuum vectors converge strongly to \(e\), then
\[
 \langle u,e\rangle=0,\qquad \|u\|\ge v_*.
\]
Thus the limiting OS Hilbert space contains a nonvacuum vector of
finite limiting energy along every Mosco-liminf subsequence.
\end{proposition}

\begin{proof}
Strong convergence passes both the inner products and norms to the
limit.  The Mosco lower bound gives
\(\mathfrak a[u]\le\liminf_n\mathfrak a_n[u_n]<\infty\).
\end{proof}

Strong convergence in \((63.13)\) cannot be replaced by mere bounded
energy.  The next example realizes this failure with genuine
positive reversible Markov transitions.

\begin{countertheorem}[Pointwise Markov convergence permits a remote soft mode]
\label{ctr:v14v63-escape}
Let the countable state space be
\(S=\{0,1,2,\ldots\}\) and set
\[
 \pi_0=\frac12,\qquad \pi_k=2^{-(k+1)}\quad(k\ge1).        \tag{63.14}
\]
For the reference reversible star chain, put the conductance between
\(0\) and \(k\) equal to
\[
 c_k=\frac{\pi_k}{4}.                                     \tag{63.15}
\]
For \(n\ge5\), change only the \(n\)-th conductance to
\[
 c_n^{(n)}=\frac{\pi_n}{n},                               \tag{63.16}
\]
retain \((63.15)\) for \(k\ne n\), and define the stochastic matrix
\(P_n\) by
\[
 \pi_iP_n(i,j)=c_{\{i,j\}}^{(n)}
 \quad(i\ne j),
\]
with diagonal entries completing each row.

Then:
\begin{enumerate}
\item every \(P_n\) is reversible and is a positive self-adjoint
      contraction on \(L^2(\pi)\);
\item for every fixed \(i,j\),
      \(P_n(i,j)\to P_0(i,j)\), where \(P_0\) is the reference chain;
\item the spectral gap of \(I-P_n\) above constants tends to zero;
\item normalized almost-invariant vectors concentrate at state \(n\)
      and converge weakly to zero.
\end{enumerate}
Thus pointwise convergence of all fixed transition probabilities
does not prevent soft spectral mass from escaping to new states.
\end{countertheorem}

\begin{proof}
Detailed balance is built into the conductances.  Every noncentral
row has off-diagonal sum at most \(1/4\), and the central row has
off-diagonal sum at most
\[
 \pi_0^{-1}\sum_{k\ge1}c_k=\frac14.
\]
Hence every diagonal entry is at least \(3/4\).  Gershgorin applied to
the self-adjoint whitened matrix places its spectrum in
\([1/2,1]\), proving item 1.

For fixed \(i,j\), the only persistent entry affected away from the
moving state is \(P_n(0,0)\), and its change is bounded by
\[
 \pi_0^{-1}|c_n-c_n^{(n)}|\le\frac{\pi_n}{2}\longrightarrow0.
\]
This proves item 2.

Let
\[
 f_n=\mathbf1_{\{n\}}-\pi_n\mathbf1.
\]
It has mean zero,
\[
 \|f_n\|_{L^2(\pi)}^2=\pi_n(1-\pi_n),
\]
and only the edge \(\{0,n\}\) contributes to its Dirichlet energy:
\[
 \langle f_n,(I-P_n)f_n\rangle_\pi
 =c_n^{(n)}=\frac{\pi_n}{n}.                              \tag{63.17}
\]
Therefore its Rayleigh quotient is
\[
 \frac1{n(1-\pi_n)}\longrightarrow0,                     \tag{63.18}
\]
which proves gap collapse.  After normalization, the whitened vectors
are a moving coordinate vector plus a norm-vanishing constant
correction, hence converge weakly to zero.
\end{proof}

\begin{assessment}[What Mosco convergence still needs]
\label{ass:v14v63-status}
Theorem~\ref{thm:v14v63-conforming} gives a finite, noncircular route
to the V59 cofinal generator: a declared master form, a form-core
refinement, relative control of inherited conductances, and a
vanishing degree tail.  None of these data is present on the active
SM--Einstein path law because that path law itself remains missing.

The star-chain example shows why fixed-cylinder or pointwise
probability convergence is not enough for particle content or a mass
gap.  The remote soft mode is invisible on every fixed matrix entry.
A compact graph-screen or an equivalent tightness estimate is needed
in addition to Mosco convergence if low-energy spectral projections
are to converge.
\end{assessment}

\begin{decision}[V64 total variation versus energy]
\label{dec:v14v63-next}
Compare the Grand Tensor summable total-variation projective-limit
criterion with the energy-form conditions \((63.10)\)--\((63.11)\).
Derive the exact time-step and marginal-support rates under which a
total-variation path-law error controls the Dirichlet form, and give
a counterexample showing that summable measure error alone need not
control generators when \(a_n\) shrinks faster.
\end{decision}

\begin{firewall}[Claim boundary after V63]
\label{fire:v14v63}
The conductance theorem is conditional on a master form and a
form-core refinement.  Neither those inputs nor a compact spectral
screen has been populated for the active SM--Einstein cylinder.
No mass gap or full continuum follows from pointwise transition
convergence.
\end{firewall}

\section*{Version V63 append-only record}
\addcontentsline{toc}{section}{Version V63 append-only record}
The complete V62 source was retained before this continuation.  It
had \(778054\) bytes, \(23669\) lines, and SHA--256
\[
 \texttt{60E62DE54FD7D61929C97942E30650E06EBDB7EFAE2BA6CB78A9B63E5FDA567F}.
\]
V63 specialized Mosco convergence to adjacent-law conductances and
proved a reversible Markov soft-mode escape example.
\section{Fourteenth-cycle V64: total variation does not control a shrinking-step generator}
\label{sec:v14v64}

The Grand Tensor projective-limit theorem uses summable
total-variation defects to construct a compatible
reflection-positive measure.  V61--V63 use the rescaled adjacent law
\(a_n^{-1}J_n\) to construct a generator.  These are different
topologies.  We now give the exact bridge and an explicit
reflection-positive counterexample when its rate fails.

\subsection{Bounded-observable energy control}

Let \(J,\widehat J\) be two symmetric adjacent-time probability
matrices on the same finite state set and let
\[
 \delta_{\rm TV}(J,\widehat J)
 :=\frac12\sum_{i,j}|J_{ij}-\widehat J_{ij}|.              \tag{64.1}
\]
For a common physical time step \(a>0\), write
\[
 \mathfrak a_J[f]
 =\frac1{2a}\sum_{i,j}J_{ij}|f_i-f_j|^2.                 \tag{64.2}
\]

\begin{theorem}[Total-variation to bounded-form estimate]
\label{thm:v14v64-bounded}
For every bounded observable \(f\),
\[
 \boxed{
 |\mathfrak a_J[f]-\mathfrak a_{\widehat J}[f]|
 \le
 \frac{4\|f\|_\infty^2}{a}\,
 \delta_{\rm TV}(J,\widehat J).}                          \tag{64.3}
\]
Consequently total-variation convergence controls the energy of each
uniformly bounded protected observable only when
\[
 \delta_{{\rm TV},n}=o(a_n).                              \tag{64.4}
\]
\end{theorem}

\begin{proof}
Since
\(|f_i-f_j|^2\le4\|f\|_\infty^2\),
\[
\begin{split}
 |\mathfrak a_J[f]-\mathfrak a_{\widehat J}[f]|
 &\le\frac1{2a}\sum_{i,j}|J_{ij}-\widehat J_{ij}|
                     |f_i-f_j|^2\\
 &\le\frac{4\|f\|_\infty^2}{a}\delta_{\rm TV}(J,\widehat J),
\end{split}
\]
which proves \((64.3)\).  Its right side vanishes under \((64.4)\).
\end{proof}

\subsection{Full Hilbert-space control and support}

Let \(\pi,\widehat\pi\) be the marginals of \(J,\widehat J\) and
assume
\[
 \min_i\{\pi_i,\widehat\pi_i\}\ge p_*>0.                  \tag{64.5}
\]

\begin{theorem}[Total-variation to whitened-transfer radius]
\label{thm:v14v64-whitened}
The V62 whitened transfer error obeys
\[
 \boxed{
 \|\widehat T^\circ-T^\circ\|
 \le
 2\delta_{\rm TV}
 \left(\frac1{p_*}+\frac1{p_*^2}\right).}                 \tag{64.6}
\]
If both transfers are positive and injective, their exact logarithmic
resolvents satisfy
\[
\begin{split}
 &\|(I+\widehat H)^{-1}-(I+H)^{-1}\|\\
 &\qquad\le
 \frac{2a}{1-e^{-1}}
 +\frac{2\delta_{\rm TV}}{a}
  \left(\frac1{p_*}+\frac1{p_*^2}\right).                 \tag{64.7}
\end{split}
\]
Thus a sufficient cofinal rate is
\[
 \frac{\delta_{{\rm TV},n}}{a_np_{*,n}^2}\longrightarrow0.
                                                               \tag{64.8}
\]
\end{theorem}

\begin{proof}
Marginalization contracts the entrywise \(l^1\) norm, so
\[
 \|\widehat\pi-\pi\|_\infty\le2\delta_{\rm TV}.            \tag{64.9}
\]
Also
\[
 \|\widehat J-J\|_{\rm op}
 \le\|\widehat J-J\|_{\rm F}
 \le\sum_{i,j}|\widehat J_{ij}-J_{ij}|
 =2\delta_{\rm TV}.                                      \tag{64.10}
\]
A symmetric nonnegative joint matrix has
\(\|J\|_{\rm op}\le\max_i\sum_jJ_{ij}\le1\).
Insert \((64.9)\)--\((64.10)\) into
Theorem~\ref{thm:v14v62-whitening} to obtain \((64.6)\).
Equation \((64.7)\) is then
Theorem~\ref{thm:v14v62-resolvent}.  Since \(p_*\le1\),
condition \((64.8)\) makes the second term vanish; \(a_n\to0\) handles
the first.
\end{proof}

\begin{assessment}[Growing state spaces]
\label{ass:v14v64-support}
On a growing finite boundary space, the least atom \(p_{*,n}\)
usually tends to zero.  Total variation can therefore be small while
whitening amplifies the error on a rare state.  A useful programme
must either prove \((64.8)\), replace atomic whitening by a
support-adapted form estimate such as V63, or declare a protected
screen on which the marginal has a uniform lower frame.  Simply
assuming summable total variation is insufficient for the generator.
\end{assessment}

\subsection{A summable reflection-positive counterexample}

Let the state space be \(\{0,1\}\), let
\(\pi=(1/2,1/2)\), and for \(0<\epsilon<1/2\) define
\[
 P_\epsilon
 =
 \begin{pmatrix}
 1-\epsilon&\epsilon\\
 \epsilon&1-\epsilon
 \end{pmatrix},\qquad
 J_\epsilon=\frac12P_\epsilon.                            \tag{64.11}
\]
The eigenvalues of \(P_\epsilon\) are \(1\) and
\(1-2\epsilon>0\), so each adjacent law is reflection positive and
has an injective positive transfer.

\begin{countertheorem}[Summable measure error with divergent generator]
\label{ctr:v14v64-rate}
Take
\[
 \epsilon_n=2^{-n},\qquad a_n=\epsilon_n^2,\qquad
 J_n=J_{\epsilon_n}.                                     \tag{64.12}
\]
Then
\[
 \sum_n\delta_{\rm TV}(J_{n+1},J_n)<\infty,               \tag{64.13}
\]
and \(J_n\) converges in total variation to the
reflection-positive frozen law \(J_0=\frac12I\).
Nevertheless, on the mean-zero vector \((1,-1)\),
\[
 a_n^{-1}(I-P_{\epsilon_n})
 \quad\hbox{has eigenvalue}\quad
 \frac{2}{\epsilon_n}\longrightarrow\infty,               \tag{64.14}
\]
and the exact logarithmic generator has eigenvalue
\[
 -\frac{\log(1-2\epsilon_n)}{\epsilon_n^2}
 \longrightarrow\infty.                                  \tag{64.15}
\]
The resolvents converge to the projection onto constants, which is
not the resolvent of a densely defined self-adjoint operator on the
full two-dimensional Hilbert space.
\end{countertheorem}

\begin{proof}
Every entry of \(J_\epsilon\) is affine in \(\epsilon\), and direct
summation gives
\[
 \delta_{\rm TV}(J_{\epsilon},J_{\epsilon'})
 =|\epsilon-\epsilon'|.
\]
Hence \((64.13)\) is the convergent telescoping series
\(\sum_n(2^{-n}-2^{-(n+1)})\).

The nonconstant eigenvalue of \(I-P_\epsilon\) is \(2\epsilon\),
which gives \((64.14)\).  The logarithmic eigenvalue is
\(-a_n^{-1}\log(1-2\epsilon_n)\), and
\(-\log(1-x)\ge x\) gives divergence in \((64.15)\).
The generator resolvent tends to one on constants and zero on the
mean-zero line.  A resolvent \((H+I)^{-1}\) of a densely defined
self-adjoint operator is injective, whereas this limiting projection
has a nontrivial kernel.
\end{proof}

\begin{corollary}[Measure continuum and Hamiltonian continuum are independent rows]
\label{cor:v14v64-independent}
The Grand Tensor summable-total-variation theorem may yield a
compatible reflection-positive cylinder measure even when the
shrinking-step Hamiltonians have no densely defined strong-resolvent
limit on the corresponding fixed OS Hilbert space.
\end{corollary}

\begin{proof}
The family \((64.12)\) satisfies the measure hypothesis and consists
of reflection-positive finite laws, while
Countertheorem~\ref{ctr:v14v64-rate} proves failure of the generator
conclusion.
\end{proof}

\begin{assessment}[Correct interface]
\label{ass:v14v64-interface}
The measure row controls bounded cylinder expectations.  The
Hamiltonian row differentiates the transfer in physical time and
therefore magnifies errors by \(a_n^{-1}\).  A complete continuum
packet needs both: projective/tight measure convergence and
form/resolvent convergence at the derivative scale.

There is also a clock-consistency issue.  A fine one-step law with
step \(a_{n+1}\) should be compared with a coarse step \(a_n\) only
after composing the fine transfer the exact number of physical
substeps.  Comparing raw adjacent laws with unequal durations mixes
dynamics with clock calibration.
\end{assessment}

\begin{decision}[V65 audit, then exact time blocking]
\label{dec:v14v64-next}
At the compulsory V65 checkpoint, inspect the live Yang--Mills and
Navier--Stokes notes.  Then continue in V66 with the exact time-block
identity
\[
 P_{a_n}=P_{a_{n+1}}^{\,k_n},\qquad a_n=k_na_{n+1},
\]
its Gram form, and a quantitative defect bound.  This will align
coarse and fine one-step laws before any projective or Mosco
comparison.
\end{decision}

\begin{firewall}[Claim boundary after V64]
\label{fire:v14v64}
Summable total variation alone does not produce a Hamiltonian
continuum.  The active programme still lacks physical adjacent laws,
time-block consistency, derivative-scale error control, and a
classical Einstein correspondence.
\end{firewall}

\section*{Version V64 append-only record}
\addcontentsline{toc}{section}{Version V64 append-only record}
The complete V63 source was retained before this continuation.  It
had \(788292\) bytes, \(23969\) lines, and SHA--256
\[
 \texttt{C3C75645426D7CABDB6E5F7813B67ECC234DCA2A91921C36740FB00285BBF9D3}.
\]
V64 separated total-variation measure convergence from
shrinking-step generator convergence and proved the exact rate
bridge.
\section{Fourteenth-cycle V65: compulsory companion audit and architecture feedback}
\label{sec:v14v65}

This is the compulsory fifth-version audit following V60.  The
companion files were inspected read-only before the time-block
calculation planned for V66.

\begin{record}[Live companion provenance at V65]
\label{rec:v14v65-provenance}
The active Yang--Mills file was
\[
 \path{Renewal_Yang_Mills_Mass_Gap_Research_Notes_V21.tex},
\]
with \(169998\) bytes, \(4625\) lines, and SHA--256
\[
 \texttt{BEF70771D3A0AD4F4FB6C80A08D149E308CA7DE87D6EA4626C467C92C5939BC2}.
                                                               \tag{65.1}
\]
The active Navier--Stokes file remained
\[
 \path{Renewal_NS_Stability_Research_Notes_V26.tex},
\]
with \(278283\) bytes, \(5319\) lines, and SHA--256
\[
 \texttt{82C887F8C2C69E4F28FAD7C3DC8DE5B9099CB5A4BF431EBA1FFF3E91935978AD}.
                                                               \tag{65.2}
\]
\end{record}

\begin{assessment}[Yang--Mills V20--V21]
\label{ass:v14v65-ym}
Yang--Mills V20 certifies the direct four-dimensional quarter-grid
boxes at all \(256\) centres for each action, with one uniform
radius, and proves by an exact point that their union is not a torus
cover.  V21 then applies the V54 two-block architecture to the actual
\(14\oplus31\) parity split.  It proves an exact operator-norm cross
certificate which succeeds at the centre, while the Frobenius cross
bound fails.  Despite that improvement, the sharp scalar block
margin yields a smaller local radius than the direct full-matrix
certificate.  V21 therefore closes that split and returns to a
full-matrix interval enclosure retaining the near-null cancellation.

This is a useful negative transfer to the present programme.  It
supports the V55 rule that a populated direct whole operator should
be certified before introducing diagonal/cross bounds, and it
supports the V57 pivot away from trying to force a positive
gravitational diagonal through a block estimate.  Its matrices and
rational constants act on the Yang--Mills fibre and are not imported.
\end{assessment}

\begin{assessment}[Navier--Stokes status]
\label{ass:v14v65-ns}
Navier--Stokes V26 is unchanged.  Its rank-one Oseen-kernel
calculation has no path-time blocking, reflected transfer, or
Einstein correspondence row.  Nothing is imported from it.
\end{assessment}

\begin{decision}[Audit decision]
\label{dec:v14v65}
The V66 direction is unchanged but sharpened by the Yang--Mills
lesson: time blocking will be performed on the full transfer
operator or full adjacent joint kernel.  It will not be inferred from
separately blocked source sectors unless their common reducing
property has already been proved.  V66 will derive the exact
coarse/fine power identity, a telescoping defect estimate, and the
effect on logarithmic generators.  The next compulsory companion
audit is V70.
\end{decision}

\begin{firewall}[Claim boundary after V65]
\label{fire:v14v65}
The audit changes certificate architecture only.  It supplies no
active SM--Einstein adjacent law, blocked transfer, Mosco limit,
spectral gap, or classical correspondence.
\end{firewall}

\section*{Version V65 append-only record}
\addcontentsline{toc}{section}{Version V65 append-only record}
The complete V64 source was retained before this continuation.  It
had \(796789\) bytes, \(24220\) lines, and SHA--256
\[
 \texttt{93041B3BE4CD278F90CB8E5087C532A4AB298855F7C7AE68DCED3FE25661585E}.
\]
V65 performed the compulsory companion audit and retained the direct
full-transfer direction.  The next audit is V70.
\section{Fourteenth-cycle V66: exact time blocking and a floor-free block defect}
\label{sec:v14v66}

Coarse and fine adjacent laws represent different physical durations.
They can be compared only after the fine transfer has been composed
through the coarse duration.  This section gives the exact identity
and a resolvent-stable defect estimate on the full transfer carrier.

\subsection{Exact powers and adjacent-law Grams}

Let \(D=\operatorname{diag}(\pi)\) be a faithful stationary marginal
and let \(P\) be a positive self-adjoint Markov contraction on
\(L^2(\pi)\).  Put
\[
 J_r:=DP^r,\qquad r\in\mathbb N_0.                        \tag{66.1}
\]

\begin{theorem}[Exact time-block Gram law]
\label{thm:v14v66-gram}
For all \(r,s\ge0\),
\[
 J_0=D,\qquad
 \boxed{J_{r+s}=J_rD^{-1}J_s.}                            \tag{66.2}
\]
In particular,
\[
 J_k=J_1(D^{-1}J_1)^{k-1}=DP^k.                          \tag{66.3}
\]
Every \(J_k\) is symmetric positive semidefinite, has marginal
\(\pi\), and is positive definite when \(P\) is injective.
\end{theorem}

\begin{proof}
Using \(J_r=DP^r\),
\[
 J_rD^{-1}J_s=(DP^r)D^{-1}(DP^s)=DP^{r+s}.
\]
Reversibility makes \(P\) self-adjoint in \(L^2(\pi)\), so
\(D^{1/2}P^kD^{-1/2}\) is Hermitian positive semidefinite.  Congruence
gives \(J_k\succeq0\).  Since \(P^k\mathbf1=\mathbf1\), the row sums
are \(\pi\).  Finally
\(\ker P^k=\ker P\) for a positive self-adjoint \(P\), proving the
last assertion.
\end{proof}

\begin{corollary}[Generator is invariant under exact clock blocking]
\label{cor:v14v66-log}
Let a fine step have duration \(a>0\), let the coarse duration be
\(b=ka\), and suppose \(P\) is injective.  Then
\[
 \boxed{
 -\frac1b\log(P^k)
 =-\frac1a\log P.}                                       \tag{66.4}
\]
\end{corollary}

\begin{proof}
For every spectral value \(\lambda\in(0,1]\),
\[
 -\frac1{ka}\log(\lambda^k)
 =-\frac1a\log\lambda.
\]
Functional calculus proves \((66.4)\).
\end{proof}

\subsection{Different physical carriers}

Let \(\mathcal H_c,\mathcal H_f\) be coarse and fine OS Hilbert spaces,
let \(W:\mathcal H_c\to\mathcal H_f\) be an isometry, and let
\(T_c,T_f\) be injective positive self-adjoint contractions.  Put
\[
 b=ka,\qquad
 \Delta_{\rm blk}:=T_f^kW-WT_c.                           \tag{66.5}
\]

\begin{theorem}[Exact and approximate blocking across an embedding]
\label{thm:v14v66-embedding}
If \(\Delta_{\rm blk}=0\), then
\[
 H_fW=WH_c                                                   \tag{66.6}
\]
on the operator domains, where
\[
 H_f=-a^{-1}\log T_f,\qquad
 H_c=-b^{-1}\log T_c.                                    \tag{66.7}
\]
For arbitrary block defect
\(\|\Delta_{\rm blk}\|\le\delta_{\rm blk}\) and \(0<b\le1\),
\[
\boxed{
 \left\|(H_f+I)^{-1}W-W(H_c+I)^{-1}\right\|
 \le
 \frac{2b}{1-e^{-1}}+\frac{\delta_{\rm blk}}{b}.}          \tag{66.8}
\]
No uniform lower spectral floor for either transfer is required.
\end{theorem}

\begin{proof}
When the defect vanishes,
\(T_f^kW=WT_c\).  Bounded functional calculus first gives the same
identity for every continuous function on \([0,1]\).  Approximating
the logarithm monotonically on \([\epsilon,1]\) and then sending
\(\epsilon\downarrow0\) gives the domain identity
\[
 -b^{-1}\log(T_f^k)W=W(-b^{-1}\log T_c).
\]
The left operator equals \(H_f\) by \((66.4)\), proving \((66.6)\).

For the quantitative assertion define
\[
 A_f^{(k)}=b^{-1}(I-T_f^k),\qquad
 A_c=b^{-1}(I-T_c).                                      \tag{66.9}
\]
The intertwining resolvent identity is
\[
\begin{split}
 &(A_f^{(k)}+I)^{-1}W-W(A_c+I)^{-1}\\
 &\quad=(A_f^{(k)}+I)^{-1}
 \bigl(WA_c-A_f^{(k)}W\bigr)(A_c+I)^{-1}.                \tag{66.10}
\end{split}
\]
Both resolvents have norm at most one, and
\[
 \|WA_c-A_f^{(k)}W\|
 =b^{-1}\|\Delta_{\rm blk}\|
 \le\delta_{\rm blk}/b.                                  \tag{66.11}
\]
Apply the V59 floor-free resolvent comparison at step \(b\) to
\(T_c\) and to \(T_f^k\).  The exact logarithm of the latter, divided
by \(b\), is \(H_f\).  Inserting \((66.10)\) between the two
comparison estimates gives \((66.8)\).
\end{proof}

\begin{corollary}[Cofinal block-consistency rate]
\label{cor:v14v66-rate}
If
\[
 b_n\longrightarrow0,\qquad
 \frac{\delta_{{\rm blk},n}}{b_n}\longrightarrow0,        \tag{66.12}
\]
then the coarse and fine logarithmic resolvents agree asymptotically
through the embeddings.  Hence a strong-resolvent limit on either
side is inherited by the other.
\end{corollary}

\begin{proof}
The right side of \((66.8)\) tends to zero.
\end{proof}

\subsection{A telescoping full-transfer certificate}

\begin{lemma}[Product defect]
\label{lem:v14v66-product}
If \(S_1,\ldots,S_k\) and \(R_1,\ldots,R_k\) are contractions on one
Hilbert space, then
\[
 \left\|S_k\cdots S_1-R_k\cdots R_1\right\|
 \le\sum_{j=1}^k\|S_j-R_j\|.                              \tag{66.13}
\]
\end{lemma}

\begin{proof}
The exact telescoping identity is
\[
 S_k\cdots S_1-R_k\cdots R_1
 =
 \sum_{j=1}^k
 S_k\cdots S_{j+1}(S_j-R_j)R_{j-1}\cdots R_1.
\]
Take norms and use contractivity.
\end{proof}

\begin{corollary}[Inhomogeneous substep card]
\label{cor:v14v66-inhomogeneous}
Suppose the physically observed fine substeps are contractions
\(S_1,\ldots,S_k\), a stationary proposal is \(T_f\), and
\[
 \|S_j-T_f\|\le\epsilon_j,\qquad
 \|T_f^{\rm obs}-S_k\cdots S_1\|\le\epsilon_0.             \tag{66.14}
\]
Then the full observed block obeys
\[
 \|T_f^{\rm obs}-T_f^k\|
 \le\epsilon_0+\sum_{j=1}^k\epsilon_j.                    \tag{66.15}
\]
After insertion of the coarse embedding and observed coarse-transfer
radius, \((66.15)\) supplies \(\delta_{\rm blk}\) in \((66.8)\).
\end{corollary}

\begin{proof}
Apply Lemma~\ref{lem:v14v66-product} with every \(R_j=T_f\), then use
the triangle inequality.
\end{proof}

\subsection{Clock and gain defects}

\begin{proposition}[Clock-ratio error]
\label{prop:v14v66-clock}
Suppose \(T_c=T_f^k\) exactly but the declared coarse duration is
\(b\ne ka\).  Then
\[
 H_c=\frac{ka}{b}\,H_f.                                   \tag{66.16}
\]
Thus a transfer-power identity without the physical duration identity
does not identify the generator.
\end{proposition}

\begin{proof}
Functional calculus gives
\[
 H_c=-b^{-1}\log(T_f^k)
 =-(k/b)\log T_f=(ka/b)H_f.
\]
\end{proof}

\begin{proposition}[Scalar gain and the absolute energy row]
\label{prop:v14v66-gain}
If
\[
 T_c=g\,T_f^k,\qquad g>0,                                 \tag{66.17}
\]
and both sides are positive contractions, then
\[
 H_c=H_f-\frac{\log g}{b}I.                               \tag{66.18}
\]
Trace or Gibbs centering removes this scalar, but the vacuum-energy
and Einstein-source rows do not.
\end{proposition}

\begin{proof}
The scalar commutes with \(T_f\), so
\[
 \log(gT_f^k)=(\log g)I+k\log T_f.
\]
Divide by \(-b=-ka\).  Centering annihilates scalar multiples of the
identity.  V51 showed that the uncentered constant and its metric
first variation contribute to the gravitational source.
\end{proof}

\subsection{Why sectors cannot be blocked independently}

\begin{counterexample}[Diagonal sector data miss blocked cross propagation]
\label{ctr:v14v66-sector}
For \(0<s<r\) and \(r+s<1\), let
\[
 T_\pm=
 \begin{pmatrix}r&\pm s\\ \pm s&r\end{pmatrix}.           \tag{66.19}
\]
Both are injective positive contractions and have the same diagonal
sector compressions.  But
\[
 T_\pm^2=
 \begin{pmatrix}
 r^2+s^2&\pm2rs\\
 \pm2rs&r^2+s^2
 \end{pmatrix},                                          \tag{66.20}
\]
so separately blocked diagonal sectors do not determine the full
blocked transfer.
\end{counterexample}

\begin{proof}
The eigenvalues of both \(T_\pm\) are \(r+s,r-s\in(0,1)\).
Equation \((66.20)\) is direct multiplication.
\end{proof}

\begin{assessment}[Time-block acquisition]
\label{ass:v14v66-acquisition}
A valid coarse/fine card now consists of the full transfers, the
physical durations, one OS-Hilbert isometry, and the defect
\((66.5)\).  If data arrive as path laws, the exact target is the
Gram identity \((66.2)\).  Inhomogeneous substeps may be used through
\((66.15)\).  A comparison of unblocked adjacent laws at unequal
durations is not a dynamical residual.

The floor-free bound \((66.8)\) removes a serious acquisition burden:
small finite transfer eigenvalues may approach zero.  What must be
small is the block defect relative to the coarse duration.
\end{assessment}

\begin{decision}[V67 gauge projection and blocking]
\label{dec:v14v66-next}
Before attempting Einstein correspondence, prove that gauge
projection, OS null quotient, and time blocking commute on the active
carrier.  Derive an approximate commutator bound when the Haar
projector and transfer fail to commute, and show how that defect
propagates to the physical resolvent.  This will prevent a
nonreducing gauge compression from being logarithmically promoted.
\end{decision}

\begin{firewall}[Claim boundary after V66]
\label{fire:v14v66}
The time-block theorems are exact conditional interfaces.  No active
SM--Einstein fine or coarse transfer, duration calibration, Hilbert
embedding, or block defect has been populated.
\end{firewall}

\section*{Version V66 append-only record}
\addcontentsline{toc}{section}{Version V66 append-only record}
The complete V65 source was retained before this continuation.  It
had \(800369\) bytes, \(24305\) lines, and SHA--256
\[
 \texttt{8FC04FE574563F3D9939776FE677FAAAE1B877A44A95DFE55521B730DC31E91F}.
\]
V66 proved exact coarse/fine time blocking and the floor-free
resolvent defect bound.
\section{Fourteenth-cycle V67: Haar reduction, OS null descent, and transfer blocking}
\label{sec:v14v67}

V35--V36 established the finite Haar projector and its relation to
temporal gauge.  V66 now creates a new compatibility question:
whether that projector reduces the full transfer before powers and
logarithms are taken.  We answer it exactly and quantify the
nonreducing defect.

\subsection{Exact gauge reduction}

Let \(\mathcal H\) be a finite OS Hilbert space, let \(T\) be an
injective positive self-adjoint contraction, and let
\[
 Q=\int_GU(g)\,dg                                         \tag{67.1}
\]
be the Haar projector for a unitary representation of a compact
gauge group \(G\).

\begin{theorem}[Exact commuting diagram]
\label{thm:v14v67-exact}
Assume
\[
 [T,U(g)]=0\quad\hbox{for every }g\in G.                  \tag{67.2}
\]
Then \(Q\) reduces \(T\), the physical transfer
\[
 T_{\rm phys}:=QTQ|_{Q\mathcal H}                         \tag{67.3}
\]
is an injective positive self-adjoint contraction, and for every
\(k\ge1\),
\[
 T_{\rm phys}^k=QT^kQ|_{Q\mathcal H}.                     \tag{67.4}
\]
For a step \(a>0\),
\[
 H_{\rm phys}:=-a^{-1}\log T_{\rm phys}
 =Q(-a^{-1}\log T)Q|_{Q\mathcal H}.                       \tag{67.5}
\]
Thus Haar reduction, time blocking, and the transfer logarithm
commute exactly.
\end{theorem}

\begin{proof}
Integrating \((67.2)\) gives \([T,Q]=0\), so
\(Q\mathcal H\) and its orthogonal complement reduce \(T\).
Compression to a reducing subspace preserves positivity,
self-adjointness, and contractivity.  If
\(T_{\rm phys}x=0\) with \(x\in Q\mathcal H\), reduction gives
\(Tx=0\), hence \(x=0\).  Block diagonal multiplication proves
\((67.4)\), and block diagonal functional calculus proves
\((67.5)\).
\end{proof}

\subsection{OS null descent}

Before Hilbert completion, let \({\cal V}_+\) be the future test
space with positive semidefinite OS form and null space
\[
 {\cal N}=\{v:\langle v,v\rangle_{\rm OS}=0\}.             \tag{67.6}
\]

\begin{proposition}[Exact null invariance is required]
\label{prop:v14v67-null}
Suppose linear maps \(T_0,Q_0:{\cal V}_+\to{\cal V}_+\) satisfy
\[
 T_0{\cal N}\subseteq{\cal N},\qquad Q_0{\cal N}\subseteq{\cal N}.
                                                               \tag{67.7}
\]
Then they descend to the OS quotient
\({\cal V}_+/{\cal N}\).  If their quotient maps satisfy the
hypotheses of Theorem~\ref{thm:v14v67-exact}, all identities
\((67.3)\)--\((67.5)\) hold after completion.

If either inclusion in \((67.7)\) fails, the corresponding quotient
operator is not well defined; a small ambient defect cannot replace
exact null descent.
\end{proposition}

\begin{proof}
Two representatives of the same quotient class differ by
\(n\in{\cal N}\).  Their images define the same class exactly when
the image of every such \(n\) again lies in \({\cal N}\).  This proves
necessity and sufficiency of \((67.7)\).  The remaining statements
then follow on the quotient from
Theorem~\ref{thm:v14v67-exact}.
\end{proof}

\subsection{Approximate gauge leakage}

Put
\[
 \ell:=\|(I-Q)TQ\|.                                       \tag{67.8}
\]
For self-adjoint \(T\), this is also the norm of the opposite cross
block.

\begin{lemma}[Haar commutator controls leakage]
\label{lem:v14v67-haar}
If
\[
 \sup_{g\in G}\|[T,U(g)]\|\le\epsilon_G,                  \tag{67.9}
\]
then
\[
 \|[T,Q]\|\le\epsilon_G,\qquad \ell\le\epsilon_G.         \tag{67.10}
\]
\end{lemma}

\begin{proof}
Bochner integration and the triangle inequality give
\[
 [T,Q]=\int_G[T,U(g)]\,dg,\qquad
 \|[T,Q]\|\le\int_G\|[T,U(g)]\|\,dg\le\epsilon_G.
\]
Since
\[
 (I-Q)[T,Q]Q=(I-Q)TQ,
\]
compression gives the second bound.
\end{proof}

Let
\[
 B:=QTQ|_{Q\mathcal H}.                                   \tag{67.11}
\]
Even without commutation, \(B\) is a positive self-adjoint
contraction and is injective because \(T\) is positive and injective.

\begin{theorem}[Gauge leakage under time blocking]
\label{thm:v14v67-block}
For every \(k\ge1\),
\[
 \boxed{
 \|QT^kQ|_{Q\mathcal H}-B^k\|
 \le(k-1)\ell.}                                          \tag{67.12}
\]
Thus separately projecting each fine substep and projecting only
after the full block agree asymptotically only when the accumulated
leakage is controlled.
\end{theorem}

\begin{proof}
Put
\[
 D_k=QT^kQ|_{Q\mathcal H}-B^k.
\]
Then \(D_1=0\), and
\[
\begin{split}
 D_{k+1}
 &=
 \bigl(QT^kQ-B^k\bigr)B
 +QT^k(I-Q)TQ\\
 &=D_kB+QT^k(I-Q)TQ.
\end{split}
\]
Contractivity of \(T,B,Q\) gives
\[
 \|D_{k+1}\|\le\|D_k\|+\ell.
\]
Induction proves \((67.12)\).
\end{proof}

\begin{theorem}[Floor-free physical resolvent radius]
\label{thm:v14v67-resolvent}
Let
\[
 H=-a^{-1}\log T,\qquad
 H_B=-a^{-1}\log B.
\]
Let \(W:Q\mathcal H\hookrightarrow\mathcal H\) be the inclusion.
For \(0<a\le1\),
\[
 \boxed{
 \|(H+I)^{-1}W-W(H_B+I)^{-1}\|
 \le
 \frac{2a}{1-e^{-1}}+\frac{\ell}{a}.}                     \tag{67.13}
\]
For a \(k\)-step block of duration \(b=ka\), the difference between
project-after-block and block-after-project logarithmic resolvents is
bounded by
\[
 \boxed{
 \frac{2b}{1-e^{-1}}+\frac{(k-1)\ell}{b}.}                \tag{67.14}
\]
\end{theorem}

\begin{proof}
The one-step intertwining defect is
\[
 TW-WB=(I-Q)TQ|_{Q\mathcal H},
\]
whose norm is \(\ell\).  Apply
Theorem~\ref{thm:v14v66-embedding} with \(k=1\) to obtain
\((67.13)\).

For the blocked comparison, use Theorem~\ref{thm:v14v67-block} as
the block defect between \(QT^kQ\) and \(B^k\), then apply the V66
floor-free block estimate at duration \(b\).
\end{proof}

\begin{corollary}[Cofinal gauge-reduction rate]
\label{cor:v14v67-rate}
For a fixed blocking factor \(k\), a sufficient condition for
physical and full logarithmic resolvents to agree through the gauge
inclusion is
\[
 a_n\to0,\qquad \frac{\ell_n}{a_n}\to0.                  \tag{67.15}
\]
If \(k=k_n\) varies, the blocked condition is
\[
 b_n\to0,\qquad
 \frac{(k_n-1)\ell_n}{b_n}\to0.                           \tag{67.16}
\]
\end{corollary}

\begin{proof}
Insert the stated rates in \((67.13)\) and \((67.14)\).
\end{proof}

\subsection{A finite warning}

\begin{counterexample}[Projecting before and after blocking differ]
\label{ctr:v14v67-order}
Let
\[
 T=\begin{pmatrix}1/2&1/4\\1/4&1/2\end{pmatrix},\qquad
 Q=\begin{pmatrix}1&0\\0&0\end{pmatrix}.                  \tag{67.17}
\]
Then \(T\) is an injective positive contraction, but
\[
 (QTQ|_{Q\mathcal H})^2=\frac14,\qquad
 QT^2Q|_{Q\mathcal H}=\frac5{16}.                         \tag{67.18}
\]
Hence a nonreducing gauge compression does not commute even with two
time steps, and its logarithm cannot be identified with the
compression of the full generator.
\end{counterexample}

\begin{proof}
The eigenvalues of \(T\) are \(3/4\) and \(1/4\).
Direct multiplication gives \((67.18)\).
\end{proof}

\subsection{Active acquisition card}

\begin{acquisition}[Gauge-compatible transfer card]
\label{acq:v14v67}
At each cutoff the active programme must retain:
\begin{enumerate}
\item the full OS quotient and exact descent of both time translation
      and the compact gauge action through its null space;
\item the full positive transfer \(T_n\) and Haar projector \(Q_n\);
\item either the exact commutators \([T_n,U_n(g)]=0\), verified on a
      finite determining representation set, or an outward leakage
      bound \(\ell_n\);
\item the physical step and blocking factor, with rates
      \((67.15)\) or \((67.16)\);
\item refinement embeddings intertwining the gauge representations
      and projectors.
\end{enumerate}
Only after these rows pass may \(Q_nT_nQ_n\) be logarithmically
promoted as the physical gauge-invariant transfer.
\end{acquisition}

\begin{assessment}[Progress]
\label{ass:v14v67-progress}
V67 closes a compatibility gap between the earlier Haar construction
and the new transfer-form route.  Exact gauge covariance makes every
operation commute.  Approximate covariance is useful only at the
derivative scale \(\ell_n=o(a_n)\), or its blocked analogue.  The OS
null space is stricter: descent through it must be exact because
otherwise the quotient operator depends on the chosen representative.

The active source provides a finite Haar architecture and Ward
identities but no populated full transfer on which the commutator or
leakage can be evaluated.
\end{assessment}

\begin{decision}[V68 physical-form Mosco reduction]
\label{dec:v14v67-next}
Assuming exact gauge reduction at each cutoff, prove that Mosco
convergence of full transfer forms descends to the physical
gauge-invariant forms when the Haar projectors converge strongly and
the recovery sequences can be averaged without energy growth.
Give the corresponding failure witness when projector ranges drift
or acquire hidden directions.
\end{decision}

\begin{firewall}[Claim boundary after V67]
\label{fire:v14v67}
The gauge-transfer compatibility theorem is conditional.  The active
full transfer, exact OS null descent, Haar commutator, leakage rate,
and compatible refinement projectors remain unpopulated.
\end{firewall}

\section*{Version V67 append-only record}
\addcontentsline{toc}{section}{Version V67 append-only record}
The complete V66 source was retained before this continuation.  It
had \(809876\) bytes, \(24611\) lines, and SHA--256
\[
 \texttt{2759912C9F1D0335170A597BDCC632C721A23EB267701D83D0328DD323BCF89B}.
\]
V67 proved exact and approximate compatibility of Haar reduction,
OS null descent, time blocking, and the transfer logarithm.
\section{Fourteenth-cycle V68: Mosco descent to the gauge-invariant OS Hilbert space}
\label{sec:v14v68}

The Grand Tensor Mosco theorem applies to full OS Hilbert spaces.
V67 supplies finite Haar projectors.  We now prove the precise
condition under which the full form limit descends to the physical
gauge-invariant ranges.  Strong convergence of projectors is enough
for the form limit, but not for convergence of low-energy
multiplicity.

\subsection{Physical restriction theorem}

After the canonical embeddings, place all forms on one Hilbert space
\(\mathcal H\).  Let \(\mathfrak a_n\) be closed nonnegative forms
Mosco-converging to a closed nonnegative form \(\mathfrak a\).
Let \(Q_n,Q\) be orthogonal projections with
\[
 Q_n\longrightarrow Q\quad\hbox{strongly}.               \tag{68.1}
\]
Assume
\[
 Q_n{\cal D}(\mathfrak a_n)\subset{\cal D}(\mathfrak a_n)
\]
and that numbers \(\eta_n\downarrow0\) satisfy
\[
 \mathfrak a_n[Q_nu]
 \le\mathfrak a_n[u]+\eta_n\|u\|^2
 \quad(u\in{\cal D}(\mathfrak a_n)).                      \tag{68.2}
\]
Define the physical restrictions
\[
 {\cal D}(\mathfrak a_n^{\rm phys})
 ={\cal D}(\mathfrak a_n)\cap Q_n\mathcal H,\qquad
 \mathfrak a_n^{\rm phys}=\mathfrak a_n
 \ \hbox{on this domain}.                                \tag{68.3}
\]

\begin{theorem}[Mosco descent through convergent gauge projectors]
\label{thm:v14v68-descent}
Suppose the limiting projection preserves the limiting form domain
and its restricted form on \(Q\mathcal H\) is closed.  Then
\[
 \mathfrak a_n^{\rm phys}
 \longrightarrow
 \mathfrak a^{\rm phys}:=
 \mathfrak a|_{{\cal D}(\mathfrak a)\cap Q\mathcal H}     \tag{68.4}
\]
in the Mosco sense.  Consequently the physical resolvents and
semigroups converge strongly through the gauge-invariant
embeddings.
\end{theorem}

\begin{proof}
For the lower condition, let
\(u_n\in Q_n\mathcal H\) converge weakly to \(u\) and suppose the
physical energies have finite liminf.  The full Mosco lower condition
gives
\[
 \mathfrak a[u]\le\liminf_n\mathfrak a_n[u_n].            \tag{68.5}
\]
It remains to prove \(u\in Q\mathcal H\).  For every \(v\),
\[
 \langle u_n,(I-Q)v\rangle
 =\langle u_n,Q_n(I-Q)v\rangle.
\]
The right side tends to zero because \(u_n\) is bounded and
\(Q_n(I-Q)v\to Q(I-Q)v=0\).  Passing to the weak limit gives
\(\langle u,(I-Q)v\rangle=0\), hence \(Qu=u\).
This proves the physical lower clause.

For recovery, take
\(u\in{\cal D}(\mathfrak a)\cap Q\mathcal H\).  Full Mosco
convergence gives \(v_n\to u\) with
\[
 \mathfrak a_n[v_n]\longrightarrow\mathfrak a[u].         \tag{68.6}
\]
Set \(u_n=Q_nv_n\).  Strong convergence of the projections and
\(v_n\to u=Qu\) imply \(u_n\to u\).  By \((68.2)\),
\[
 \limsup_n\mathfrak a_n[u_n]
 \le\mathfrak a[u].
\]
The already proved lower clause applied to this strongly convergent
sequence gives the reverse liminf.  Thus
\(\mathfrak a_n[u_n]\to\mathfrak a[u]\), proving recovery.
The resolvent and semigroup statements follow from the inherited
Mosco equivalence.
\end{proof}

\begin{corollary}[Exact gauge covariance supplies energy contraction]
\label{cor:v14v68-exact}
If the operator \(A_n\) represented by \(\mathfrak a_n\) commutes with
\(Q_n\), then
\[
 \mathfrak a_n[Q_nu]\le\mathfrak a_n[u],                  \tag{68.7}
\]
so Theorem~\ref{thm:v14v68-descent} applies with \(\eta_n=0\).
\end{corollary}

\begin{proof}
Commutation makes the physical and orthogonal ranges reducing, and
\[
 \mathfrak a_n[u]
 =\mathfrak a_n[Q_nu]+\mathfrak a_n[(I-Q_n)u].
\]
The second term is nonnegative.
\end{proof}

\subsection{Approximate gauge covariance at the form scale}

For the transfer form,
\[
 A_n=a_n^{-1}(I-T_n),\qquad
 \ell_n=\|(I-Q_n)T_nQ_n\|.                                \tag{68.8}
\]

\begin{proposition}[Leakage gives the Mosco projection debit]
\label{prop:v14v68-leakage}
For every \(u\in\mathcal H\),
\[
 \mathfrak a_n[Q_nu]
 \le
 \mathfrak a_n[u]+\frac{\ell_n}{a_n}\|u\|^2.              \tag{68.9}
\]
Hence the V67 rate
\[
 \ell_n/a_n\longrightarrow0                              \tag{68.10}
\]
supplies \((68.2)\).
\end{proposition}

\begin{proof}
Write \(u=x+y\) with \(x=Q_nu\) and \(y=(I-Q_n)u\).
Since \(A_n\ge0\),
\[
\begin{split}
 \mathfrak a_n[u]
 &=\mathfrak a_n[x]+\mathfrak a_n[y]
   +2\operatorname{Re}\langle x,A_ny\rangle\\
 &\ge\mathfrak a_n[x]
   -2\|Q_nA_n(I-Q_n)\|\,\|x\|\,\|y\|.
\end{split}
\]
The cross block of \(I\) vanishes and self-adjointness gives
\[
 \|Q_nA_n(I-Q_n)\|=\ell_n/a_n.
\]
Finally \(2\|x\|\|y\|\le\|u\|^2\), which proves \((68.9)\).
\end{proof}

\subsection{Projector convergence witnesses}

\begin{proposition}[Finite source witness for projector drift]
\label{prop:v14v68-drift}
Let \(\{s_j\}_{j\ge1}\) be dense in \(\mathcal H\).  If
\(Q_n\) does not converge strongly to \(Q\), then there are an index
\(j\), a number \(\epsilon>0\), and a subsequence \(n_k\) such that
\[
 \|(Q_{n_k}-Q)s_j\|\ge\epsilon                            \tag{68.11}
\]
for all \(k\).
\end{proposition}

\begin{proof}
If \((Q_n-Q)s_j\to0\) for every vector in a dense family, uniform
boundedness \(\|Q_n-Q\|\le2\) extends convergence to every vector in
\(\mathcal H\).  The contrapositive gives \((68.11)\).
\end{proof}

In the programme, the \(s_j\) can be chosen as the countable primitive
source words of the Grand Tensor direct system.  Thus failure of
strong Haar transport has a finite-word witness.

\subsection{Hidden physical modes and spectral escape}

Mosco descent does not imply norm-resolvent convergence or stability
of physical kernel multiplicity.

\begin{countertheorem}[An extra gauge-invariant ground mode can escape]
\label{ctr:v14v68-hidden}
Let
\[
 \mathcal H=\mathbb Ce_0\oplus\ell^2(\mathbb N),
\]
let \(Q\) project onto \(\mathbb Ce_0\), and let
\[
 Q_n=Q+|e_n\rangle\langle e_n|.                           \tag{68.12}
\]
Define nonnegative diagonal operators
\[
 A_ne_0=A_ne_n=0,\qquad A_ne_j=e_j\ (j\ne n),
                                                               \tag{68.13}
\]
and
\[
 Ae_0=0,\qquad Ae_j=e_j\ (j\ge1).                         \tag{68.14}
\]
Then:
\begin{enumerate}
\item \(Q_n\to Q\) strongly and \([A_n,Q_n]=0\);
\item the full forms of \(A_n\) Mosco-converge to the form of \(A\);
\item the physical restricted forms Mosco-converge as in
      Theorem~\ref{thm:v14v68-descent};
\item every finite physical operator has a two-dimensional kernel,
      while the limiting physical operator has a one-dimensional
      kernel.
\end{enumerate}
The normalized extra ground vectors \(e_n\) converge weakly to zero
and have no strongly convergent subsequence.
\end{countertheorem}

\begin{proof}
Moving coordinate projections converge strongly to zero, giving item
1.  The resolvents satisfy
\[
 (A_n+I)^{-1}-(A+I)^{-1}
 =\frac12|e_n\rangle\langle e_n|,
\]
which converges strongly to zero.  The inherited
Mosco--resolvent equivalence gives item 2, and
Theorem~\ref{thm:v14v68-descent} gives item 3.
The kernel statements follow directly from \((68.13)\)--\((68.14)\).
The standard basis of \(\ell^2\) is weakly null and not norm
precompact.
\end{proof}

\begin{theorem}[Physical compact-screen supplement]
\label{thm:v14v68-screen}
Assume Theorem~\ref{thm:v14v68-descent}.  If the embedded physical
graph unit balls
\[
 \left\{
 (u,A_n^{1/2}u):
 u\in Q_n{\cal D}(\mathfrak a_n),\
 \|u\|^2+\mathfrak a_n[u]\le1
 \right\}                                                 \tag{68.15}
\]
are uniformly approximable by finite-rank projections in
\(\mathcal H\oplus\mathcal H\), then the physical resolvents converge
in norm and isolated physical eigenvalues, including kernel
multiplicity, converge.
\end{theorem}

\begin{proof}
This is the compact graph-screen branch of the inherited Grand Tensor
theorem applied to the physical forms whose Mosco convergence was
proved in Theorem~\ref{thm:v14v68-descent}.  Uniform graph
precompactness gives collective compactness of the resolvents; strong
plus collective compact convergence gives norm convergence, and
Riesz projections give spectral multiplicity convergence.
\end{proof}

\begin{assessment}[What is now structurally closed]
\label{ass:v14v68-status}
Once full transfer forms exist, exact Haar covariance and strong
projector transport are sufficient to pass their Mosco limit to the
gauge-invariant OS Hilbert space.  Approximate covariance remains
admissible at the already identified rate \(\ell_n=o(a_n)\).
Low-energy particle content still requires the separate physical
compact-screen row; V68's escaping ground mode proves this cannot be
deduced from Mosco convergence.

The active programme has formal Haar and source-word architectures,
but no populated transfer forms, projector transport residuals, or
physical graph screens.  The next issue is correspondence: even a
well-defined nonnegative physical generator must be shown to
implement the finite Standard Model--Einstein constraints and
classical brackets.
\end{assessment}

\begin{decision}[V69 generator correspondence residual]
\label{dec:v14v68-next}
Define a finite, basis-independent residual comparing commutators of
the reconstructed physical generator with quantized coordinate
observables against the intended Hamiltonian vector field.  Prove a
Duhamel bound from this residual to finite-time Heisenberg evolution,
and separate first-order correspondence from the Einstein constraint
and stress identities.
\end{decision}

\begin{firewall}[Claim boundary after V68]
\label{fire:v14v68}
Mosco descent is conditional on full forms, strong Haar transport,
and form-scale leakage control.  The physical graph screen and the
Einstein--SM correspondence remain unpopulated; no spectral or
continuum field content is claimed.
\end{firewall}

\section*{Version V68 append-only record}
\addcontentsline{toc}{section}{Version V68 append-only record}
The complete V67 source was retained before this continuation.  It
had \(819353\) bytes, \(24917\) lines, and SHA--256
\[
 \texttt{4AD6A68CF59DCDB98D7A42C7DE76B2AABF3472279C7A357638C80E5958D3A375}.
\]
V68 proved Mosco descent to gauge-invariant OS Hilbert spaces and
isolated the separate physical compact-screen requirement.
\section{Fourteenth-cycle V69: a generator-correspondence residual and finite-time control}
\label{sec:v14v69}

A nonnegative physical transfer generator is not yet the
Einstein--Standard-Model Hamiltonian.  The first correspondence test
is whether its commutators implement the intended Hamiltonian vector
field on a determining finite observable family.  We define that test
basis-independently, prove a finite-time Duhamel bound, and identify
the scalar energy row that commutators cannot see.

\subsection{The finite correspondence packet}

Let \(E\) be a finite-dimensional real coefficient space with metric
\(M\succ0\), let \(\mathcal K\) be a finite physical Hilbert space,
and let
\[
 {\cal Q}:E\longrightarrow{\rm Herm}(\mathcal K)          \tag{69.1}
\]
be a real-linear quantization map for a declared family of reduced
observables.  Let \(L:E\to E\) be the linearized classical
Hamiltonian vector field at the chosen physical background.  For a
self-adjoint transfer generator \(H\), define
\[
 {\cal R}_{H,L}(a)
 :=i[H,{\cal Q}(a)]-{\cal Q}(La).                         \tag{69.2}
\]

\begin{definition}[Basis-independent correspondence residual]
\label{def:v14v69-residual}
Give \({\rm Herm}(\mathcal K)\) the Hilbert--Schmidt metric and put
\[
 \boxed{
 \Delta_{\rm corr}(H;{\cal Q},L,M)
 :=
 \|{\cal R}_{H,L}M^{-1/2}\|_{\rm HS}^2.}                 \tag{69.3}
\]
Here the Hilbert--Schmidt norm includes both the coefficient and
operator indices.
\end{definition}

\begin{proposition}[Coordinate invariance and exact meaning]
\label{prop:v14v69-invariance}
Under an invertible coefficient change \(a=Sa'\), let
\[
 {\cal Q}'={\cal Q}S,\qquad L'=S^{-1}LS,\qquad M'=S^*MS.
                                                               \tag{69.4}
\]
Then
\[
 \Delta_{\rm corr}(H;{\cal Q}',L',M')
 =
 \Delta_{\rm corr}(H;{\cal Q},L,M).                       \tag{69.5}
\]
Moreover,
\[
 \Delta_{\rm corr}=0
 \iff i[H,{\cal Q}(a)]={\cal Q}(La)
 \quad\hbox{for every }a\in E.                            \tag{69.6}
\]
\end{proposition}

\begin{proof}
The transformed residual synthesis is
\({\cal R}'={\cal R}S\).  Since
\[
 S(S^*MS)^{-1}S^*=M^{-1},
\]
the squared Hilbert--Schmidt norm, written as a trace, is unchanged.
A squared norm vanishes exactly when its linear map vanishes, proving
\((69.6)\).
\end{proof}

\subsection{Finite-time linearized evolution}

Let
\[
 \alpha_t^H(A)=e^{itH}Ae^{-itH}.                          \tag{69.7}
\]

\begin{theorem}[Duhamel correspondence bound]
\label{thm:v14v69-duhamel}
For every \(a\in E\) and \(t\ge0\),
\[
\begin{split}
 &\alpha_t^H({\cal Q}(a))-{\cal Q}(e^{tL}a)\\
 &\quad=
 \int_0^t
 \alpha_{t-s}^H\!\left(
 {\cal R}_{H,L}(e^{sL}a)\right)\,ds.                      \tag{69.8}
\end{split}
\]
Consequently
\[
 \boxed{
 \|\alpha_t^H({\cal Q}(a))-{\cal Q}(e^{tL}a)\|_{\rm HS}
 \le
 t e^{t\|L\|_M}
 \sqrt{\Delta_{\rm corr}}\,\|a\|_M.}                     \tag{69.9}
\]
\end{theorem}

\begin{proof}
Differentiate
\[
 F(s)=\alpha_{t-s}^H({\cal Q}(e^{sL}a)).
\]
Using
\(\frac d{ds}\alpha_{t-s}^H(A)=-i[H,\alpha_{t-s}^H(A)]\)
and covariance of the commutator under \(\alpha^H\) gives
\[
 F'(s)
 =-\alpha_{t-s}^H({\cal R}_{H,L}(e^{sL}a)).
\]
Integrating from \(0\) to \(t\) and reversing the endpoints gives
\((69.8)\).  Unitary conjugation preserves the Hilbert--Schmidt norm.
Furthermore,
\[
 \|{\cal R}_{H,L}(b)\|_{\rm HS}
 \le\sqrt{\Delta_{\rm corr}}\|b\|_M,
\qquad
 \|e^{sL}a\|_M\le e^{s\|L\|_M}\|a\|_M.
\]
Insert these bounds in \((69.8)\) and use
\(\int_0^te^{s\|L\|_M}ds\le te^{t\|L\|_M}\).
\end{proof}

\subsection{Identification up to the invisible scalar}

Let \({\cal A}_{\cal Q}\) be the unital complex algebra generated by
\({\cal Q}(E)\).

\begin{theorem}[Exact uniqueness from a determining observable family]
\label{thm:v14v69-uniqueness}
Let \(H,K\) be self-adjoint and suppose
\[
 i[H,{\cal Q}(a)]=i[K,{\cal Q}(a)]
 \quad\hbox{for all }a\in E.                              \tag{69.10}
\]
If
\[
 {\cal A}_{\cal Q}=M_{\dim\mathcal K}(\mathbb C),          \tag{69.11}
\]
then
\[
 H-K=cI                                                  \tag{69.12}
\]
for a real scalar \(c\).  One absolute energy anchor, for example
\[
 \langle\Omega,H\Omega\rangle
 =\langle\Omega,K\Omega\rangle                            \tag{69.13}
\]
for a unit vector \(\Omega\), forces \(c=0\).
\end{theorem}

\begin{proof}
Equation \((69.10)\) says that \(H-K\) commutes with every generator
of \({\cal A}_{\cal Q}\), hence with the full matrix algebra in
\((69.11)\).  Its commutant consists of scalar multiples of the
identity, proving \((69.12)\).  Taking the expectation in
\((69.13)\) gives \(c=0\).
\end{proof}

For a quantitative version, let
\[
 {\rm Herm}_0(\mathcal K)
 =\{X=X^*:\operatorname{Tr}X=0\}
\]
and define the commutator synthesis
\[
 {\cal C}_{\cal Q}(X)(a)=i[X,{\cal Q}(a)].                \tag{69.14}
\]
Let
\[
 \sigma_{\rm ad}
 :=
 \inf_{\substack{X\in{\rm Herm}_0(\mathcal K)\\
                  \|X\|_{\rm HS}=1}}
 \|{\cal C}_{\cal Q}(X)M^{-1/2}\|_{\rm HS}.               \tag{69.15}
\]

\begin{theorem}[Quantitative generator identification]
\label{thm:v14v69-quantitative}
Condition \((69.11)\) is equivalent to \(\sigma_{\rm ad}>0\).
For any self-adjoint \(H,K\),
\[
 \boxed{
 \|\operatorname{ctr}(H-K)\|_{\rm HS}
 \le
 \frac{
 \|({\cal R}_{H,L}-{\cal R}_{K,L})M^{-1/2}\|_{\rm HS}}
 {\sigma_{\rm ad}},}                                     \tag{69.16}
\]
where
\[
 \operatorname{ctr}X=X-\frac{\operatorname{Tr}X}
 {\dim\mathcal K}I.
\]
If \(K\) has exact correspondence residual zero, then
\[
 \|\operatorname{ctr}(H-K)\|_{\rm HS}
 \le\frac{\sqrt{\Delta_{\rm corr}(H)}}{\sigma_{\rm ad}}. \tag{69.17}
\]
\end{theorem}

\begin{proof}
The kernel of \({\cal C}_{\cal Q}\) consists of self-adjoint elements
of the commutant of \({\cal A}_{\cal Q}\).  Its restriction to
\({\rm Herm}_0\) is injective exactly when that commutant contains
only scalars, which is equivalent in finite dimensions to
\((69.11)\).  On the unit sphere of the finite-dimensional traceless
space, injectivity makes the continuous norm attain a positive
minimum, proving \(\sigma_{\rm ad}>0\); the converse is immediate.

Put \(X=\operatorname{ctr}(H-K)\).  Scalar centering does not change
commutators, and
\[
 {\cal C}_{\cal Q}(X)
 ={\cal R}_{H,L}-{\cal R}_{K,L}.
\]
The lower-frame definition \((69.15)\) gives \((69.16)\), and
\((69.17)\) follows when the second residual vanishes.
\end{proof}

\begin{corollary}[Finite unitary-evolution comparison]
\label{cor:v14v69-unitary}
After scalar centering,
\[
 \|\alpha_t^H(A)-\alpha_t^K(A)\|
 \le
 \frac{2|t|\,\|A\|}{\sigma_{\rm ad}}
 \|({\cal R}_{H,L}-{\cal R}_{K,L})M^{-1/2}\|_{\rm HS}.
                                                               \tag{69.18}
\]
\end{corollary}

\begin{proof}
Scalar shifts do not change Heisenberg evolution.  Duhamel's formula
for the two unitary conjugations gives
\[
 \|\alpha_t^H(A)-\alpha_t^K(A)\|
 \le2|t|\|A\|\,\|\operatorname{ctr}(H-K)\|.
\]
Use operator norm at most Hilbert--Schmidt norm and
Theorem~\ref{thm:v14v69-quantitative}.
\end{proof}

\subsection{Independent Einstein and constraint rows}

\begin{proposition}[Commutator correspondence is blind to vacuum energy]
\label{prop:v14v69-scalar}
For every real \(c\),
\[
 \Delta_{\rm corr}(H+cI;{\cal Q},L,M)
 =
 \Delta_{\rm corr}(H;{\cal Q},L,M).                       \tag{69.19}
\]
The transfer changes by the scalar gain
\[
 e^{-a(H+cI)}=e^{-ac}e^{-aH}.                             \tag{69.20}
\]
Therefore the absolute gain and its metric variation remain required
for the cosmological and stress source rows.
\end{proposition}

\begin{proof}
Scalar operators commute with every observable, proving
\((69.19)\).  Equation \((69.20)\) is functional calculus.
The V51 source calculation showed that centering removes the scalar
from dynamics but not from the Einstein variation.
\end{proof}

\begin{assessment}[Typed separation of correspondence tests]
\label{ass:v14v69-separation}
The residual \((69.3)\) tests the linearized Hamiltonian vector field
on a declared reduced observable family.  It does not by itself test:
\begin{enumerate}
\item the finite Gauss, momentum, and scalar constraint algebra;
\item Ward or BRST identities outside the reduced observable algebra;
\item the metric derivative defining stress;
\item the absolute vacuum/counterterm energy;
\item nonlinear correspondence away from the chosen background.
\end{enumerate}
Those rows remain separate.  Conversely, expectation-level Einstein
stationarity does not imply \((69.3)\), because a first variation of
a state does not determine the operator commutators.

The finite lower frame \(\sigma_{\rm ad}\) is the exact
identifiability constant.  If it collapses cofinally, a normalized
traceless operator becomes invisible to all selected commutators.
Adding more evolution estimates without repairing that observable
blind spot would be circular.
\end{assessment}

\begin{acquisition}[Finite generator-correspondence card]
\label{acq:v14v69}
At each cutoff one must supply on the same physical OS carrier:
\begin{enumerate}
\item the reconstructed self-adjoint transfer generator \(H_n\);
\item the reduced observable synthesis \({\cal Q}_n\), coefficient
      metric \(M_n\), and intended finite Hamiltonian linearization
      \(L_n\) obtained from the common action;
\item the residual \((69.3)\) and commutator lower frame
      \(\sigma_{{\rm ad},n}\);
\item one absolute energy and metric-variation anchor;
\item independent constraint-algebra, Ward/BRST, and stress
      residuals.
\end{enumerate}
A regulator-uniform finite-time correspondence estimate requires the
ratio
\(\sqrt{\Delta_{{\rm corr},n}}/\sigma_{{\rm ad},n}\) to vanish on
the protected observable core.
\end{acquisition}

\begin{decision}[V70 audit and constraint closure]
\label{dec:v14v69-next}
At the compulsory V70 checkpoint inspect the live Yang--Mills and
Navier--Stokes notes.  Then continue in V71 by adjoining the finite
constraint-algebra residual to \((69.3)\).  Prove that vanishing
generator correspondence on a determining set does not imply HDA
closure, and give a joint quantitative packet that transports both
dynamics and the first-class algebra.
\end{decision}

\begin{firewall}[Claim boundary after V69]
\label{fire:v14v69}
The correspondence residual is an exact finite compiler.  No active
transfer generator, intended finite Einstein--SM linearization,
commutator lower frame, absolute energy anchor, or cofinal residual
rate has been populated.
\end{firewall}

\section*{Version V69 append-only record}
\addcontentsline{toc}{section}{Version V69 append-only record}
The complete V68 source was retained before this continuation.  It
had \(829467\) bytes, \(25214\) lines, and SHA--256
\[
 \texttt{A44C0D688813E4566BF67D32B1D5D6D81B6B6ACA27D743BBE336D9FA6490993B}.
\]
V69 defined the finite generator-correspondence residual, proved its
Duhamel bound, and identified the scalar energy row invisible to
commutators.
\section{Fourteenth-cycle V70: compulsory companion audit and joint-residual decision}
\label{sec:v14v70}

This is the compulsory audit after V65.  The live companions were
inspected read-only before constructing the finite constraint-algebra
packet.

\begin{record}[Live companion provenance at V70]
\label{rec:v14v70-provenance}
The active Yang--Mills file was
\[
 \path{Renewal_Yang_Mills_Mass_Gap_Research_Notes_V23.tex},
\]
with \(185996\) bytes, \(5070\) lines, and SHA--256
\[
 \texttt{5117E4168F6E208A7B3127586BD089A85059106B418294BEF33DF4FAB88AF4D1}.
                                                               \tag{70.1}
\]
The active Navier--Stokes file remained
\[
 \path{Renewal_NS_Stability_Research_Notes_V26.tex},
\]
with \(278283\) bytes, \(5319\) lines, and SHA--256
\[
 \texttt{82C887F8C2C69E4F28FAD7C3DC8DE5B9099CB5A4BF431EBA1FFF3E91935978AD}.
                                                               \tag{70.2}
\]
\end{record}

\begin{assessment}[Yang--Mills V22--V23]
\label{ass:v14v70-ym}
Yang--Mills V22 aligns a rational \(1\oplus44\) splitting with the
actual soft eigendirection.  This reduces the normalized centre cross
norm by many orders of magnitude and doubles the previously certified
direct radius at one rational phase.  V23 then computes anisotropic
coordinate profiles.  Individual transverse axes reach a much larger
radius, while simultaneous transverse motion remains restricted
because the proof replaces signed mixed coefficients by
triangle-square bounds.  The next YM step therefore retains the
mixed signs in a joint enclosure.

The transferable lesson is about residual organization: a basis or
splitting aligned with the near-null direction can be decisive, but
separate norm bounds can again destroy the cancellation when several
directions act together.  No Yang--Mills matrix, phase, floor, or
radius belongs to the SM--Einstein carrier.
\end{assessment}

\begin{assessment}[Navier--Stokes status]
\label{ass:v14v70-ns}
Navier--Stokes V26 is unchanged and supplies no finite constraint,
transfer, or commutator datum.  Nothing is imported.
\end{assessment}

\begin{decision}[Joint constraint residual]
\label{dec:v14v70}
V71 will not add separately maximized Gauss, momentum, scalar, and
generator errors.  It will assemble the complete typed
constraint-commutator defect as one linear synthesis and retain its
full Gram.  Source whitening will be performed only on the supported
constraint coefficient space.  The resulting smallest singular
value will be the identifiability margin, while every null direction
will be reported explicitly.  The next compulsory audit is V75.
\end{decision}

\begin{firewall}[Claim boundary after V70]
\label{fire:v14v70}
The audit changes only the architecture of the next finite residual.
It supplies no active SM--Einstein constraint matrices, structure
coefficients, transfer generator, or continuum rate.
\end{firewall}

\section*{Version V70 append-only record}
\addcontentsline{toc}{section}{Version V70 append-only record}
The complete V69 source was retained before this continuation.  It
had \(840450\) bytes, \(25555\) lines, and SHA--256
\[
 \texttt{C850F11C686371D965DB68280DF9758EFB290E3CFB5DD88F95E070013BE7EB9F}.
\]
V70 performed the compulsory companion audit and selected a joint
constraint-residual compiler.  The next audit is V75.
\section{Fourteenth-cycle V71: joint generator and first-class constraint residuals}
\label{sec:v14v71}

V69 tests Hamiltonian evolution on a determining observable family.
That condition does not imply that the finite Gauss, momentum, and
scalar constraints close or that their common kernel is preserved.
We now assemble all three defects as one supported operator synthesis.
The full Gram is retained, so mixed constraint directions are tested
without summing separate worst-case bounds.

\subsection{Finite typed constraint system}

Let \(A\) be a finite real constraint coefficient space with metric
\(N\succ0\), and let
\[
 {\cal C}:A\longrightarrow{\rm Herm}(\mathcal K)          \tag{71.1}
\]
be the quantized constraint synthesis on the finite physical carrier
before imposing its common kernel.  Let
\[
 B:\Lambda^2A\longrightarrow A                           \tag{71.2}
\]
be the intended antisymmetric structure map at the chosen finite
background.  Define
\[
 {\cal F}_B(\xi\wedge\eta)
 :=
 -i[{\cal C}(\xi),{\cal C}(\eta)]
 -{\cal C}(B(\xi\wedge\eta)).                             \tag{71.3}
\]
Let \(N_{\wedge}\) be the metric induced by \(N\) on
\(\Lambda^2A\).

\begin{definition}[Constraint-algebra Gram]
\label{def:v14v71-algebra}
The complete closure Gram and its scalar residual are
\[
 G_{\rm alg}
 :={\cal F}_B^*{\cal F}_B,\qquad
 \Delta_{\rm alg}
 :=\|{\cal F}_BN_{\wedge}^{-1/2}\|_{\rm HS}^2.            \tag{71.4}
\]
The adjoint in \(G_{\rm alg}\) uses the Hilbert--Schmidt metric on
operators and the declared coefficient metrics.
\end{definition}

\begin{proposition}[Exact meaning and basis invariance]
\label{prop:v14v71-algebra}
The inertia and spectrum of the whitened Gram
\[
 N_{\wedge}^{-1/2}G_{\rm alg}N_{\wedge}^{-1/2}            \tag{71.5}
\]
are invariant under invertible changes of constraint coordinates.
Moreover,
\[
 \Delta_{\rm alg}=0
 \iff
 -i[{\cal C}(\xi),{\cal C}(\eta)]
 ={\cal C}(B(\xi\wedge\eta))
 \quad\hbox{for all }\xi,\eta.                            \tag{71.6}
\]
\end{proposition}

\begin{proof}
A coordinate change acts on \({\cal F}_B\) by right composition with
the induced invertible map on \(\Lambda^2A\), while
\(N_{\wedge}\) transforms by congruence.  The whitened Gram therefore
changes by unitary equivalence.  A Hilbert--Schmidt squared norm
vanishes exactly when its synthesis vanishes, proving \((71.6)\).
\end{proof}

\subsection{Source-minimal anomaly split}

Define the raw commutator synthesis
\[
 {\cal K}(\xi\wedge\eta)
 :=-i[{\cal C}(\xi),{\cal C}(\eta)].                      \tag{71.7}
\]
Let \(P_{\cal C}\) be the Hilbert--Schmidt orthogonal projection onto
\(\operatorname{Ran}{\cal C}\).

\begin{theorem}[Best structure fit and irreducible anomaly]
\label{thm:v14v71-anomaly}
Among all linear structure maps
\(B:\Lambda^2A\to A\), the minimum Hilbert--Schmidt closure residual
is
\[
 \boxed{
 \Delta_{\rm anom}^{\min}
 =
 \|(I-P_{\cal C}){\cal K}N_{\wedge}^{-1/2}\|_{\rm HS}^2.} \tag{71.8}
\]
One minimum-norm fitted structure map is
\[
 B_{\min}={\cal C}^{\dagger}{\cal K},                     \tag{71.9}
\]
with the Moore--Penrose inverse taken on the declared coefficient and
operator metrics.  For every \(B\),
\[
\begin{split}
 \|({\cal K}-{\cal C}B)N_{\wedge}^{-1/2}\|_{\rm HS}^2
 ={}&\|(I-P_{\cal C}){\cal K}N_{\wedge}^{-1/2}\|_{\rm HS}^2\\
 &+\|(P_{\cal C}{\cal K}-{\cal C}B)
        N_{\wedge}^{-1/2}\|_{\rm HS}^2.                  \tag{71.10}
\end{split}
\]
Thus a commutator component outside the constraint span cannot be
removed by changing structure coefficients.
\end{theorem}

\begin{proof}
Split
\[
 {\cal K}-{\cal C}B
 =(I-P_{\cal C}){\cal K}
 +(P_{\cal C}{\cal K}-{\cal C}B).
\]
The two summands have orthogonal operator ranges, so Pythagoras gives
\((71.10)\).  The first is independent of \(B\).  The Moore--Penrose
solution \((71.9)\) realizes the orthogonal projection of
\({\cal K}\) onto \(\operatorname{Ran}{\cal C}\), proving
\((71.8)\).
\end{proof}

\begin{firewall}[No post hoc HDA fit]
\label{fire:v14v71-fit}
The intended finite HDA structure map must be fixed by the common
action, incidence, lapse/shift convention, and background before its
closure residual is interpreted physically.  Equation \((71.9)\) is
a diagnostic best fit.  It cannot be substituted for the intended
structure map merely because it reduces \(\Delta_{\rm alg}\).
\end{firewall}

\subsection{Constraint preservation by the transfer generator}

Let \(H\) be the reconstructed self-adjoint transfer generator and
let \(D:A\to A\) be the intended linear transport of constraints.
Define
\[
 {\cal E}_{H,D}(\xi)
 :=i[H,{\cal C}(\xi)]-{\cal C}(D\xi),                     \tag{71.11}
\]
\[
 \Delta_{\rm pres}
 :=\|{\cal E}_{H,D}N^{-1/2}\|_{\rm HS}^2.                \tag{71.12}
\]

\begin{theorem}[Finite-time constraint transport]
\label{thm:v14v71-transport}
For every \(\xi\in A\) and \(t\ge0\),
\[
\begin{split}
 &e^{itH}{\cal C}(\xi)e^{-itH}
 -{\cal C}(e^{tD}\xi)\\
 &\quad=
 \int_0^t
 e^{i(t-s)H}
 {\cal E}_{H,D}(e^{sD}\xi)
 e^{-i(t-s)H}\,ds.                                       \tag{71.13}
\end{split}
\]
Hence
\[
 \boxed{
 \|e^{itH}{\cal C}(\xi)e^{-itH}
 -{\cal C}(e^{tD}\xi)\|_{\rm HS}
 \le
 te^{t\|D\|_N}\sqrt{\Delta_{\rm pres}}\,\|\xi\|_N.}       \tag{71.14}
\]
If \(\Delta_{\rm pres}=0\), the common constraint kernel
\[
 \mathcal K_{\rm con}
 :=\bigcap_{\xi\in A}\ker{\cal C}(\xi)                    \tag{71.15}
\]
is invariant under \(H\).
\end{theorem}

\begin{proof}
The Duhamel calculation is the same differentiation as in
Theorem~\ref{thm:v14v69-duhamel}, with \({\cal Q},L\) replaced by
\({\cal C},D\).  This gives \((71.13)\), and unitary invariance of the
Hilbert--Schmidt norm gives \((71.14)\).

If the residual vanishes and
\(\psi\in\mathcal K_{\rm con}\), then
\[
 {\cal C}(\xi)H\psi
 =H{\cal C}(\xi)\psi+[{\cal C}(\xi),H]\psi
 =i{\cal C}(D\xi)\psi=0.
\]
Thus \(H\psi\in\mathcal K_{\rm con}\).
\end{proof}

\subsection{Quantitative leakage from the constraint kernel}

Choose an \(N\)-orthonormal constraint basis
\(\xi_1,\ldots,\xi_r\) and define
\[
 \Gamma:\mathcal K\longrightarrow\mathcal K^{\oplus r},
 \qquad
 \Gamma\psi=({\cal C}(\xi_\alpha)\psi)_{\alpha=1}^r.      \tag{71.16}
\]
Let \(P_{\rm con}\) project onto \(\ker\Gamma\), and put
\[
 \sigma_{\rm con}
 :=
 \inf_{\substack{\psi\perp\ker\Gamma\\\|\psi\|=1}}
 \|\Gamma\psi\|.                                         \tag{71.17}
\]

\begin{theorem}[Constraint-kernel leakage bound]
\label{thm:v14v71-leakage}
If \(\sigma_{\rm con}>0\), then
\[
 \boxed{
 \|(I-P_{\rm con})HP_{\rm con}\|
 \le
 \frac{\sqrt{\Delta_{\rm pres}}}{\sigma_{\rm con}}.}      \tag{71.18}
\]
Thus exact preservation is recovered when
\(\Delta_{\rm pres}=0\), while a cofinal physical reduction requires
the ratio
\[
 \sqrt{\Delta_{{\rm pres},n}}/\sigma_{{\rm con},n}\to0.   \tag{71.19}
\]
\end{theorem}

\begin{proof}
For \(\psi\in\ker\Gamma\), equation \((71.11)\) gives, componentwise,
\[
 {\cal C}(\xi_\alpha)H\psi
 =i{\cal E}_{H,D}(\xi_\alpha)\psi,
\]
because every \({\cal C}(D\xi_\alpha)\psi\) vanishes.  Hence
\[
 \|\Gamma H\psi\|
 \le\sqrt{\Delta_{\rm pres}}\|\psi\|.
\]
By the definition of \(\sigma_{\rm con}\),
\[
 \sigma_{\rm con}\|(I-P_{\rm con})H\psi\|
 \le\|\Gamma(I-P_{\rm con})H\psi\|
 =\|\Gamma H\psi\|.
\]
Take the supremum over unit \(\psi\in\ker\Gamma\).
\end{proof}

\subsection{The complete joint Gram}

Let \(E,M,{\cal Q},L\) be the observable packet of V69.  On the direct
sum
\[
 Z:=E\oplus\Lambda^2A\oplus A                             \tag{71.20}
\]
with metric
\[
 {\cal M}_Z:=M\oplus N_{\wedge}\oplus N,
\]
define the joint synthesis
\[
 {\cal S}_{\rm joint}
 :=
 {\cal R}_{H,L}\oplus{\cal F}_B\oplus{\cal E}_{H,D}.       \tag{71.21}
\]
Its complete Gram is
\[
 G_{\rm joint}:={\cal S}_{\rm joint}^*{\cal S}_{\rm joint}.
                                                               \tag{71.22}
\]

\begin{theorem}[Joint finite correspondence packet]
\label{thm:v14v71-joint}
The following are equivalent:
\[
 G_{\rm joint}=0,                                        \tag{71.23}
\]
\[
 \Delta_{\rm corr}+\Delta_{\rm alg}+\Delta_{\rm pres}=0, \tag{71.24}
\]
and simultaneous exact generator correspondence, first-class closure,
and constraint transport on all declared coefficient directions.
For a nonzero packet, the operator norm
\[
 \|{\cal S}_{\rm joint}{\cal M}_Z^{-1/2}\|^2
 =
 \lambda_{\max}\!\left(
 {\cal M}_Z^{-1/2}G_{\rm joint}{\cal M}_Z^{-1/2}\right)   \tag{71.25}
\]
is the sharp combined-direction residual.  It need not equal the sum
of separately maximized sector norms.
\end{theorem}

\begin{proof}
Every term in \((71.24)\) is nonnegative, so their sum vanishes
exactly when all three syntheses vanish.  A Gram vanishes exactly
when its synthesis vanishes, proving the equivalence.
Equation \((71.25)\) is the singular-value characterization of the
operator norm on the whitened joint coefficient space.
\end{proof}

\begin{counterexample}[Generator correspondence does not imply closure]
\label{ctr:v14v71-independence}
On \(\mathcal K=\mathbb C^2\), take \(H=0\), \(L=0\), and any
observable synthesis \({\cal Q}\).  Then
\(\Delta_{\rm corr}=0\).  Declare two intended abelian constraints
\[
 {\cal C}(e_1)=\sigma_x,\qquad
 {\cal C}(e_2)=\sigma_y,\qquad B=0.
\]
Their closure defect is
\[
 -i[\sigma_x,\sigma_y]=2\sigma_z\ne0,                    \tag{71.26}
\]
so \(\Delta_{\rm alg}>0\).
\end{counterexample}

\begin{proof}
The zero generator commutes with every observable, while the Pauli
commutator gives \((71.26)\).
\end{proof}

\begin{assessment}[Constraint status]
\label{ass:v14v71-status}
V71 provides a finite joint test that neither hides an anomaly by
fitting structure coefficients nor loses a mixed direction by adding
separate worst-case bounds.  The two lower frames have distinct
roles: \(\sigma_{\rm ad}\) identifies the generator modulo scalars,
while \(\sigma_{\rm con}\) detects leakage from the common constraint
kernel.

The Grand Tensor source gives conditional finite HDA, Ward, BRST,
stress, and Einstein residual theorems, but the active transfer
generator and its constraint matrices are not jointly populated.
The joint Gram \((71.22)\) therefore remains an acquisition target,
not a computed numerical certificate.
\end{assessment}

\begin{decision}[V72 stress and Einstein source compatibility]
\label{dec:v14v71-next}
Add the metric-variation row that commutators miss.  Starting from one
finite common action and one absolute transfer normalization, derive
the exact relation between the transfer vacuum-energy variation,
the matter stress covector, and the finite Einstein residual.
Prove which scalar gain changes are harmless for dynamics but
observable by gravity, and formulate the joint second-jet acquisition
needed for cofinal correspondence.
\end{decision}

\begin{firewall}[Claim boundary after V71]
\label{fire:v14v71}
The joint residual is a rigorous finite compiler.  No active
SM--Einstein constraint matrices, intended structure map, common
kernel frame, preservation residual, or cofinal rates have been
populated.
\end{firewall}

\section*{Version V71 append-only record}
\addcontentsline{toc}{section}{Version V71 append-only record}
The complete V70 source was retained before this continuation.  It
had \(843796\) bytes, \(25636\) lines, and SHA--256
\[
 \texttt{D8D9A8E909CEEA1CDBDB2F6A7B95E64D74301A9276FAD262B083E1DD06B5F76C}.
\]
V71 constructed the joint generator, constraint-closure, and
constraint-preservation Gram with quantitative physical leakage.
\section{Fourteenth-cycle V72: absolute transfer energy, vacuum stress, and the Einstein row}
\label{sec:v14v72}

The V69--V71 commutator packets identify dynamics only modulo a
scalar operator.  Gravity detects that scalar through metric
variation.  We now derive the exact finite relation between the
leading transfer eigenvalue, vacuum stress, counterterms, and the
finite Einstein residual.  The second derivative exposes a separate
gap-dependent response term.

\subsection{Absolute transfer normalization}

Let \(q\) vary in a finite affine metric panel \(\mathcal Q\), and let
\(T(q)\) be a \(C^2\) family of injective positive contractions on a
fixed finite OS Hilbert space.  Fix a physical step \(a>0\) and set
\[
 H(q)=-a^{-1}\log T(q).                                   \tag{72.1}
\]
Assume at \(q_0\) that the smallest eigenvalue \(E_0(q_0)\) of
\(H(q_0)\) is simple, with normalized eigenvector \(\Omega_0\).
Equivalently, the largest transfer eigenvalue
\[
 \lambda_0(q)=e^{-aE_0(q)}                               \tag{72.2}
\]
is simple near \(q_0\).

\begin{theorem}[Vacuum-energy first variation from the absolute transfer]
\label{thm:v14v72-first}
For every metric variation \(v\in\mathcal Q\),
\[
 dE_0[v]
 =\langle\Omega_0,dH[v]\Omega_0\rangle                   \tag{72.3}
\]
and
\[
 \boxed{
 dE_0[v]
 =-\frac1{a\lambda_0}
   \langle\Omega_0,dT[v]\Omega_0\rangle.}                 \tag{72.4}
\]
For a slab of physical duration \(\tau>0\), define the vacuum action
and stress covector by
\[
 S_{\rm vac}=\tau E_0,\qquad
 {\mathsf T}_{\rm vac}:=-2D_qS_{\rm vac}.                 \tag{72.5}
\]
Then
\[
 \boxed{
 {\mathsf T}_{\rm vac}[v]
 =\frac{2\tau}{a\lambda_0}
   \langle\Omega_0,dT[v]\Omega_0\rangle.}                 \tag{72.6}
\]
\end{theorem}

\begin{proof}
Differentiate
\(H(q)\Omega(q)=E_0(q)\Omega(q)\) and pair with
\(\Omega_0\).  The eigenvector derivative terms cancel by
self-adjointness and normalization, giving the
Hellmann--Feynman identity \((72.3)\).  Applying the same argument to
\(T(q)\Omega(q)=\lambda_0(q)\Omega(q)\) gives
\[
 d\lambda_0[v]=\langle\Omega_0,dT[v]\Omega_0\rangle.
\]
Differentiate \((72.2)\) to obtain \((72.4)\), and substitute it into
\((72.5)\).
\end{proof}

\subsection{Common effective action and Einstein residual}

Let
\[
 S_{\rm eff}(q,\Phi)
 =
 S_g(q)+S_{\rm mat}(q,\Phi)
 +S_{\rm vac}(q)+C_{\rm ct}(q)                            \tag{72.7}
\]
be one finite common effective action, where \(C_{\rm ct}\) is an
admissible local counterterm already present on the declared action
branch.  Define
\[
 {\mathsf T}_{\rm mat}:=-2D_qS_{\rm mat},\qquad
 {\mathsf T}_{\rm ct}:=-2D_qC_{\rm ct}.                   \tag{72.8}
\]

\begin{theorem}[Finite Einstein row including the transfer vacuum]
\label{thm:v14v72-einstein}
The metric Euler covector of the common effective action is
\[
 \boxed{
 {\cal E}^{\rm eff}_q
 :=2D_qS_{\rm eff}
 =2D_qS_g
 -{\mathsf T}_{\rm mat}
 -{\mathsf T}_{\rm vac}
 -{\mathsf T}_{\rm ct}.}                                 \tag{72.9}
\]
For any faithful metric \(G_{\mathcal Q}\) on metric variations,
\[
 \Delta_{\rm Ein}^{\rm eff}
 :=
 \langle{\cal E}^{\rm eff}_q,
 G_{\mathcal Q}^{-1}{\cal E}^{\rm eff}_q\rangle\ge0       \tag{72.10}
\]
vanishes exactly when the finite effective action is metrically
stationary.
\end{theorem}

\begin{proof}
Differentiate \((72.7)\), use the three definitions of stress, and
multiply by two.  Faithfulness and positivity of
\(G_{\mathcal Q}^{-1}\) make \((72.10)\) zero exactly when
\({\cal E}^{\rm eff}_q=0\).
\end{proof}

\begin{remark}[Relation to the Grand Tensor finite Einstein theorem]
\label{rem:v14v72-grand}
The Grand Tensor theorem
\(\texttt{thm:SMST-finite-action-Einstein}\) gives
\(2D_qS_g-{\mathsf T}_{\rm mat}\) for one finite classical common
action.  Equation \((72.9)\) is its effective transfer extension.
It is valid only when the absolute vacuum and counterterm terms are
part of that same action; an unrelated fitted transfer gain is not a
common-action source.
\end{remark}

\subsection{Scalar gain ambiguity}

Let the observed positive transfer be
\[
 \widehat T(q)=g(q)T(q),\qquad g(q)>0,                    \tag{72.11}
\]
with the product still a contraction on the panel.

\begin{theorem}[Exact gain stress]
\label{thm:v14v72-gain}
The observed generator and ground energy are
\[
 \widehat H=H-a^{-1}\log g\,I,\qquad
 \widehat E_0=E_0-a^{-1}\log g.                           \tag{72.12}
\]
The centered generator and every commutator residual of V69--V71 are
unchanged, while the vacuum stress changes by
\[
 \boxed{
 \widehat{\mathsf T}_{\rm vac}
 -{\mathsf T}_{\rm vac}
 =\frac{2\tau}{a}\,d\log g.}                              \tag{72.13}
\]
\end{theorem}

\begin{proof}
Functional calculus gives \((72.12)\).  Scalar operators commute
with every observable and constraint, so all commutator syntheses are
unchanged.  Apply \(-2\tau d\) to the second identity in
\((72.12)\), obtaining \((72.13)\).
\end{proof}

\begin{corollary}[Determining metric panel identifies the gain first jet]
\label{cor:v14v72-gain-identify}
If \(v_1,\ldots,v_r\) span \(\mathcal Q\), then the stress differences
\[
 s_i:=
 (\widehat{\mathsf T}_{\rm vac}
 -{\mathsf T}_{\rm vac})[v_i]
\]
determine \(d\log g\) uniquely by
\[
 d\log g[v_i]=\frac{a}{2\tau}s_i.                         \tag{72.14}
\]
If the panel does not span, the gain gradient is identifiable only
modulo the annihilator of the tested variations.
\end{corollary}

\begin{proof}
Equation \((72.14)\) is \((72.13)\) evaluated on the spanning family.
A linear covector is determined by its values on a spanning set.
\end{proof}

\begin{firewall}[Gain is either a source or an error]
\label{fire:v14v72-gain}
A metric-dependent gain may be absorbed only if
\(-a^{-1}\log g\) belongs to the predeclared admissible local
counterterm space and satisfies its Ward, locality, and cofinal
cocycle conditions.  Otherwise \((72.13)\) is a physical source
change or a normalization defect.  Centering cannot decide which.
\end{firewall}

\subsection{The vacuum second jet}

Let
\[
 P_0=|\Omega_0\rangle\langle\Omega_0|,\qquad Q_0=I-P_0,
\]
and suppose the ground gap is positive:
\[
 \gamma_0
 :=\lambda_{\min}\left(
 Q_0(H-E_0)Q_0|_{Q_0\mathcal K}\right)>0.                \tag{72.15}
\]
Define the reduced resolvent
\[
 R_0:=
 \left(Q_0(H-E_0)Q_0|_{Q_0\mathcal K}\right)^{-1}.        \tag{72.16}
\]

\begin{theorem}[Exact vacuum-energy Hessian]
\label{thm:v14v72-second}
For metric directions \(v,w\),
\[
\begin{split}
 D^2E_0[v,w]
 ={}&\langle\Omega_0,H''[v,w]\Omega_0\rangle\\
 &-\langle Q_0H'[v]\Omega_0,
       R_0Q_0H'[w]\Omega_0\rangle\\
 &-\langle Q_0H'[w]\Omega_0,
       R_0Q_0H'[v]\Omega_0\rangle.                       \tag{72.17}
\end{split}
\]
In particular,
\[
 D^2E_0[v,v]
 =
 \langle\Omega_0,H''[v,v]\Omega_0\rangle
 -2\langle Q_0H'[v]\Omega_0,
 R_0Q_0H'[v]\Omega_0\rangle,                             \tag{72.18}
\]
and therefore
\[
 D^2E_0[v,v]
 \ge
 \langle\Omega_0,H''[v,v]\Omega_0\rangle
 -\frac{2}{\gamma_0}\|Q_0H'[v]\Omega_0\|^2.              \tag{72.19}
\]
\end{theorem}

\begin{proof}
Choose the differentiable phase of \(\Omega(q)\) so that
\(\langle\Omega_0,\Omega'[v]\rangle=0\).  Projecting the differentiated
eigenvalue equation to \(Q_0\mathcal K\) gives
\[
 \Omega'[v]
 =-R_0Q_0H'[v]\Omega_0.                                  \tag{72.20}
\]
Differentiate the Hellmann--Feynman identity in direction \(w\).
The derivative of \(H'[v]\) gives the first line of \((72.17)\);
the two eigenvector derivatives, using \((72.20)\), give the next two
lines.  Setting \(w=v\) proves \((72.18)\).
Since \(0\prec R_0\preceq\gamma_0^{-1}Q_0\), equation \((72.19)\)
follows.
\end{proof}

\begin{countertheorem}[A first jet does not determine vacuum curvature]
\label{ctr:v14v72-first-not-second}
On \(\mathbb C^2\), let
\[
 H_c(t)=
 \begin{pmatrix}0&t\\t&1\end{pmatrix}
 +ct^2\begin{pmatrix}1&0\\0&0\end{pmatrix}.               \tag{72.21}
\]
All values \(H_c(0)\) and first derivatives \(H_c'(0)\) are identical,
but
\[
 E_{0,c}''(0)=2c-2.                                      \tag{72.22}
\]
Thus transfer value and first tangent data do not determine the
vacuum stress derivative or the effective Hessian.
\end{countertheorem}

\begin{proof}
At zero, the ground vector is \(e_1\), the gap is one,
\(H_c'[0]e_1=e_2\), and
\(\langle e_1,H_c''[0]e_1\rangle=2c\).
Equation \((72.18)\) gives \((72.22)\).
\end{proof}

\subsection{Degenerate ground spaces}

\begin{proposition}[Stress covector obstruction at a degeneracy]
\label{prop:v14v72-degenerate}
If the ground projection \(P_0\) has rank greater than one, the
one-sided directional derivatives of the lowest energy in direction
\(v\) are the extreme eigenvalues of
\[
 P_0H'[v]P_0|_{P_0\mathcal K}.                            \tag{72.23}
\]
A linear vacuum stress covector exists through the degeneracy only if
this compression is scalar on \(P_0\mathcal K\) for every metric
direction in the panel.
\end{proposition}

\begin{proof}
Finite-dimensional degenerate perturbation theory diagonalizes the
first-order perturbation inside the ground eigenspace.  The lowest
branch takes the least eigenvalue for positive parameter and the
greatest eigenvalue with the opposite sign.  These directional
derivatives assemble into one linear functional exactly when all
eigenvalues of every compressed first derivative coincide, which is
the scalar condition.
\end{proof}

\begin{assessment}[Complete finite metric row]
\label{ass:v14v72-status}
The missing scalar in V69 is now fully typed.  Its value comes from
the absolute leading transfer eigenvalue; its first metric derivative
is vacuum stress; its second derivative contains a negative response
term controlled by the ground gap.  Centered transfer-log data
determine none of these absolute quantities.

A complete second-jet Einstein certificate therefore needs the
absolute transfer, a simple or scalar-splitting ground sector, first
and second metric derivatives, the reduced resolvent or gap, the
classical matter stress, and the admissible counterterm jet on one
common action.  The attached sources provide this architecture but
not the active numerical packet.
\end{assessment}

\begin{decision}[V73 cofinal vacuum-stress transport]
\label{dec:v14v72-next}
Derive sufficient cofinal conditions under which
\(E_{0,n}\), \(dE_{0,n}\), and \(D^2E_{0,n}\) converge through the
physical embeddings.  Separate norm-resolvent convergence, isolated
ground-projection convergence, derivative convergence, and a uniform
gap.  Give a soft-mode counterexample showing that energy convergence
alone does not control stress curvature.
\end{decision}

\begin{firewall}[Claim boundary after V72]
\label{fire:v14v72}
The vacuum stress and Einstein formulas are exact finite identities.
No active absolute SM--Einstein transfer, metric derivative, ground
gap, common effective action, or counterterm jet has been populated.
\end{firewall}

\section*{Version V72 append-only record}
\addcontentsline{toc}{section}{Version V72 append-only record}
The complete V71 source was retained before this continuation.  It
had \(855303\) bytes, \(26001\) lines, and SHA--256
\[
 \texttt{DC01D8C8459B8EA67FE52C91F75EF5302AFB42DD9CDD48C615F69F5B2B01B0E8}.
\]
V72 derived the absolute transfer vacuum stress, its finite Einstein
row, the gain ambiguity, and the gap-dependent second jet.
\section{Fourteenth-cycle V73: cofinal vacuum stress and the soft-mode curvature obstruction}
\label{sec:v14v73}

V72 gives exact finite first and second derivatives of the vacuum
energy.  Their cofinal passage is stronger than convergence of the
ground energy itself.  It requires convergence of the isolated
ground projection, the metric derivatives of the generator, and the
reduced resolvent.  We now state and prove that handoff.

\subsection{Isolated ground-sector transport}

After the physical embeddings, work on a common Hilbert space
\(\mathcal H\).  Let \(H_n(q)\) and \(H(q)\) be self-adjoint
families near \(q_0\), bounded below uniformly.  At \(q_0\), suppose
\(H\) has a simple ground eigenvalue \(E_0\) isolated by a circle
\(\Gamma\) from the rest of the spectrum.

\begin{theorem}[Ground energy and projection from norm-resolvent convergence]
\label{thm:v14v73-ground}
Assume
\[
 \sup_{z\in\Gamma}
 \|(H_n-z)^{-1}-(H-z)^{-1}\|
 \longrightarrow0.                                      \tag{73.1}
\]
Then, for all sufficiently large \(n\), \(\Gamma\) encloses exactly
one eigenvalue \(E_{0,n}\) of \(H_n\), counted with multiplicity, and
\[
 E_{0,n}\to E_0,\qquad
 P_{0,n}\to P_0                                           \tag{73.2}
\]
in operator norm, where
\[
 P_{0,n}=-\frac1{2\pi i}\int_\Gamma(H_n-z)^{-1}\,dz.       \tag{73.3}
\]
Normalized ground vectors can be phased so that
\[
 \Omega_{0,n}\to\Omega_0.                                \tag{73.4}
\]
\end{theorem}

\begin{proof}
Equation \((73.1)\) and the Riesz formula give norm convergence of the
spectral projections.  For large \(n\), projections within norm less
than one have the same rank, so \(P_{0,n}\) has rank one.  Spectral
localization inside shrinking contours around \(E_0\), or the
resolvent trace on the one-dimensional Riesz range, gives
\(E_{0,n}\to E_0\).  Rank-one projection convergence permits phases
with
\[
 \Omega_{0,n}=
 \frac{P_{0,n}\Omega_0}{\|P_{0,n}\Omega_0\|}
 \longrightarrow\Omega_0.
\]
\end{proof}

\subsection{First and second metric jets}

Write at \(q_0\)
\[
 H'_{n,v}=DH_n[v],\qquad
 H''_{n,vw}=D^2H_n[v,w],                                  \tag{73.5}
\]
and similarly without \(n\).  Let
\[
 Q_{0,n}=I-P_{0,n},\qquad
 R_{0,n}=
 \left(Q_{0,n}(H_n-E_{0,n})Q_{0,n}\right)^{-1}            \tag{73.6}
\]
on the orthogonal complement.

\begin{theorem}[Cofinal vacuum first- and second-jet theorem]
\label{thm:v14v73-jets}
Assume Theorem~\ref{thm:v14v73-ground} and a uniform ground gap
\[
 \operatorname{dist}\!
 \left(E_{0,n},
 \sigma(H_n)\setminus\{E_{0,n}\}\right)\ge\gamma>0.       \tag{73.7}
\]
Suppose, uniformly for unit metric directions \(v,w\),
\[
 \|H'_{n,v}-H'_v\|\to0,\qquad
 \|H''_{n,vw}-H''_{vw}\|\to0.                             \tag{73.8}
\]
Then
\[
 \|R_{0,n}-R_0\|\to0,                                    \tag{73.9}
\]
after the complementary ranges are identified by the convergent
projections, and
\[
 DE_{0,n}\longrightarrow DE_0,\qquad
 D^2E_{0,n}\longrightarrow D^2E_0                       \tag{73.10}
\]
in the operator norms on the finite metric coefficient space.
Consequently the vacuum stresses
\[
 {\mathsf T}_{{\rm vac},n}=-2\tau_nDE_{0,n}               \tag{73.11}
\]
converge when \(\tau_n\to\tau\).
\end{theorem}

\begin{proof}
The uniform spectral gap and norm-resolvent convergence imply norm
convergence of the bounded continuous reduced-resolvent functional
which equals \((\lambda-E_{0,n})^{-1}\) off the ground spectral
window and zero on it.  Together with \(P_{0,n}\to P_0\), this gives
\((73.9)\).

The first derivative is
\[
 DE_{0,n}[v]
 =\langle\Omega_{0,n},H'_{n,v}\Omega_{0,n}\rangle.
\]
Equations \((73.4)\) and \((73.8)\) give uniform convergence on the
unit sphere of metric directions.  For the second derivative, insert
\(P_{0,n},R_{0,n},H'_{n,v},H''_{n,vw}\) into the exact V72 formula.
Every factor converges in norm and all factors are uniformly bounded
by \((73.7)\)--\((73.8)\).  The two bilinear response terms and the
direct second derivative therefore converge uniformly, proving
\((73.10)\).  Equation \((73.11)\) then gives stress convergence.
\end{proof}

\begin{remark}[Unbounded metric derivatives]
\label{rem:v14v73-unbounded}
Norm convergence in \((73.8)\) can be weakened to graph-bounded
convergence on a common core, provided the ground vectors and their
reduced-resolvent images lie in that core with uniform graph bounds.
The active programme has not populated such domains, so the stronger
finite-card hypothesis is retained here.
\end{remark}

\subsection{Cofinal gain and counterterm jets}

Let
\[
 \widehat T_n(q)=g_n(q)T_n(q),\qquad
 c_n(q):=-a_n^{-1}\log g_n(q).                            \tag{73.12}
\]

\begin{proposition}[Derivative-scale gain criterion]
\label{prop:v14v73-gain}
The scalar contributions to the generator, vacuum stress, and vacuum
curvature converge respectively if
\[
 c_n\to c,\qquad Dc_n\to Dc,\qquad D^2c_n\to D^2c.        \tag{73.13}
\]
Equivalently, the corresponding zeroth, first, and second derivatives
of \(\log g_n\), divided by \(a_n\), must converge.  The condition
\(g_n\to1\) alone implies none of these three statements.
\end{proposition}

\begin{proof}
Equation \((72.12)\) gives
\(\widehat H_n=H_n+c_nI\).  Differentiation gives the three scalar
jets directly.  For the last assertion, take any divergent sequence
\(b_n\) with \(a_nb_n\to0\) and set
\(g_n=e^{-a_nb_n}\).  Then \(g_n\to1\) while
\(c_n=b_n\to\infty\).
\end{proof}

\begin{corollary}[Summable effective-source cocycle]
\label{cor:v14v73-cocycle}
If the shell increments
\[
 d_n:=c_{n+1}-c_n
\]
satisfy
\[
 \sum_n\left(
 \|d_n\|_{C^0}
 +\|Dd_n\|
 +\|D^2d_n\|\right)<\infty,                              \tag{73.14}
\]
then \(c_n\) converges in \(C^2\), and its limiting zeroth, stress,
and curvature contributions are well defined.
\end{corollary}

\begin{proof}
The three telescoping series are absolutely Cauchy in the finite
metric panel.
\end{proof}

\subsection{Soft-mode curvature counterexample}

\begin{countertheorem}[Energy and stress convergence do not control curvature]
\label{ctr:v14v73-soft}
Let \(\gamma_n\downarrow0\) and consider
\[
 H_n(t)=
 \begin{pmatrix}
 1&t\\ t&1+\gamma_n
 \end{pmatrix}.                                          \tag{73.15}
\]
At \(t=0\),
\[
 E_{0,n}(0)=1,\qquad E'_{0,n}(0)=0                       \tag{73.16}
\]
for every \(n\), but
\[
 E''_{0,n}(0)=-\frac2{\gamma_n}\longrightarrow-\infty.   \tag{73.17}
\]
For every sufficiently small \(t\) at fixed \(n\), \(H_n(t)\) is
positive and defines an injective positive transfer
\(e^{-aH_n(t)}\).  Thus convergence of vacuum energy and stress at
one background does not control the Einstein second jet when a soft
mode closes the ground gap.
\end{countertheorem}

\begin{proof}
At zero the ground vector is \(e_1\), the excited energy is
\(1+\gamma_n\), and the derivative maps \(e_1\) to \(e_2\).
The direct second derivative of \(H_n\) is zero.  Formula \((72.18)\)
therefore gives
\[
 E''_{0,n}(0)
 =-2\langle e_2,\gamma_n^{-1}e_2\rangle
 =-2/\gamma_n.
\]
Positivity holds for \(|t|^2<1+\gamma_n\), by the positive
determinant and trace.
\end{proof}

\subsection{Passage of the finite Einstein row}

Let
\[
 {\cal E}^{\rm eff}_{q,n}
 =
 2D_qS_{g,n}
 -{\mathsf T}_{{\rm mat},n}
 -{\mathsf T}_{{\rm vac},n}
 -{\mathsf T}_{{\rm ct},n}.                              \tag{73.18}
\]

\begin{theorem}[Cofinal effective Einstein residual]
\label{thm:v14v73-einstein}
Assume the four covectors on the right of \((73.18)\) converge in one
common finite metric-coordinate system, and faithful metrics
\(G_{\mathcal Q,n}\) converge with uniform upper and lower bounds.
Then
\[
 {\cal E}^{\rm eff}_{q,n}\to{\cal E}^{\rm eff}_{q},
\]
and
\[
 \Delta_{{\rm Ein},n}^{\rm eff}
 =
 \langle{\cal E}^{\rm eff}_{q,n},
 G_{\mathcal Q,n}^{-1}{\cal E}^{\rm eff}_{q,n}\rangle
 \longrightarrow
 \Delta_{\rm Ein}^{\rm eff}.                             \tag{73.19}
\]
In particular, if the finite residuals tend to zero, the limiting
effective Einstein covector vanishes on the declared metric panel.
\end{theorem}

\begin{proof}
Finite-dimensional convergence and the uniform metric bounds imply
convergence of the inverse metrics.  Substitute the four convergent
covectors in \((73.18)\), then pass to the quadratic forms.
\end{proof}

\begin{assessment}[Cofinal source boundary]
\label{ass:v14v73-status}
The absolute transfer row now has a complete cofinal theorem.
Norm-resolvent convergence locates the ground sector; a uniform gap
stabilizes its reduced resolvent; first and second generator
derivatives pass the stress and curvature; and a \(C^2\) scalar-gain
cocycle controls the absolute counterterm source.  Each hypothesis is
independent.

The active source does not supply the absolute transfers, their
metric derivatives, the isolated ground projections, or the gap.
Consequently no vacuum stress has been numerically coupled to the
finite Einstein residual.
\end{assessment}

\begin{decision}[V74 conditional protected-continuum assembly]
\label{dec:v14v73-next}
Assemble the results of V59--V73 into one conditional protected
SM--Einstein continuum theorem.  List the minimal finite rows and
cofinal rates, prove that they yield a gauge-invariant semibounded OS
generator with first-class constraint preservation, finite-time
classical correspondence, and a vanishing effective Einstein
covector on the protected panel.  State separately every property
still absent, especially full Euclidean covariance, locality beyond
the protected algebra, and a mass gap.
\end{decision}

\begin{firewall}[Claim boundary after V73]
\label{fire:v14v73}
The cofinal vacuum-stress theorem is conditional.  The active ground
sector, uniform gap, generator metric derivatives, gain cocycle, and
effective Einstein residual have not been populated.
\end{firewall}

\section*{Version V73 append-only record}
\addcontentsline{toc}{section}{Version V73 append-only record}
The complete V72 source was retained before this continuation.  It
had \(866672\) bytes, \(26345\) lines, and SHA--256
\[
 \texttt{A6CE0E6DD189AB71B454CB0742A395876D7F23D773404A45B420499EB917F4A4}.
\]
V73 proved cofinal passage of vacuum energy, stress, and curvature
and exhibited the ground-gap soft-mode obstruction.
\section{Fourteenth-cycle V74: conditional assembly of the protected SM--Einstein continuum}
\label{sec:v14v74}

The preceding block replaced a failed positive-gravitational-Hessian
route by a reflected-transfer route.  We now assemble its finite and
cofinal interfaces into one theorem.  The word
\emph{protected} is essential: the conclusion concerns the declared
cylinder algebra, physical OS Hilbert direct system, observable core,
constraint family, and metric panel.  It is not yet the full
Einstein--Standard-Model continuum.

\subsection{The complete input packet}

At cutoff \(n\), let the following objects be defined on one common
absolute Euclidean cylinder and its OS quotient.

\begin{enumerate}
\item[\textup{(P1)}] A nonnegative reflection-invariant finite weight
      \(\mu_n\), positive-time algebra \({\cal A}_{n,+}\), and
      reflection Gram \(G^{\rm RP}_n\succeq0\).
\item[\textup{(P2)}] A faithful one-time marginal \(\pi_n\), a
      symmetric adjacent law \(J_n\succeq0\), a three-time law
      \(K_n\), physical step \(a_n\downarrow0\), and the V61
      marginal and held-out Markov identities, with every submitted
      uncertainty satisfying the V62 derivative-scale rate.
\item[\textup{(P3)}] Positive injective transfers
      \(T_n=D_n^{-1}J_n\), exact or asymptotically exact time blocking
      with \(\delta_{{\rm blk},n}/b_n\to0\), and compatible physical
      clock ratios.
\item[\textup{(P4)}] Compact-gauge Haar projectors \(Q_n\) descending
      exactly through the OS null spaces, converging strongly through
      the Hilbert embeddings, with transfer leakage
      \(\ell_n/a_n\to0\).
\item[\textup{(P5)}] Physical transfer forms
      \[
       \mathfrak a_n^{\rm phys}[u]
       =a_n^{-1}\langle u,(I-Q_nT_nQ_n)u\rangle
      \]
      satisfying either the V63 conductance criterion or the abstract
      Grand Tensor Mosco clauses, with a densely defined closed
      nonnegative limiting form \(\mathfrak a_\infty^{\rm phys}\).
\item[\textup{(P6)}] A uniform physical compact graph screen and one
      protected nonvacuum covariance lower bound with the tightness
      and uniform-integrability rows of the Grand Tensor quantum
      cylinder theorem.
\item[\textup{(P7)}] Reduced observable syntheses
      \({\cal Q}_n\), classical linearizations \(L_n\), and the V69
      correspondence data, with
      \[
       \sqrt{\Delta_{{\rm corr},n}}/
       \sigma_{{\rm ad},n}\longrightarrow0,               \tag{74.1}
      \]
      and compatible convergence on a protected operator core.
\item[\textup{(P8)}] Constraint syntheses \({\cal C}_n\), intended
      structure maps \(B_n\), transport maps \(D_n^{\rm con}\), and
      the V71 joint Grams, with
      \[
      \Delta_{{\rm alg},n}\to0,\qquad
      \sqrt{\Delta_{{\rm pres},n}}/
      \sigma_{{\rm con},n}\to0,                           \tag{74.2}
      \]
      plus graph convergence of the protected constraint operators.
\item[\textup{(P9)}] Absolute transfer normalizations, isolated
      simple ground sectors with a uniform gap, convergent first and
      second metric derivatives, a \(C^2\) gain/counterterm cocycle,
      and finite effective Einstein covectors satisfying
      \[
       \Delta_{{\rm Ein},n}^{\rm eff}\longrightarrow0.    \tag{74.3}
      \]
\item[\textup{(P10)}] Reflection-equivariant cutoff maps with
      summable total-variation defects, a common quasilocal
      counterterm tail, compatible determinant/Pfaffian and gauge
      phases, and the Grand Tensor field-tightness row.
\end{enumerate}

Every item is an input condition or a finite residual with an
explicit failure witness.  No item is inferred from the desired
continuum conclusion.

\subsection{Conditional assembly theorem}

\begin{theorem}[Protected reflected SM--Einstein continuum]
\label{thm:v14v74-protected}
If \textup{(P1)}--\textup{(P10)} hold, then after passage to a
subsequence there exist:
\begin{enumerate}
\item a nontrivial reflection-positive probability measure
      \(\mu_\infty\) on the declared protected field space;
\item a gauge-invariant physical OS Hilbert space
      \(\mathcal H_{\rm phys,\infty}\);
\item a nonnegative self-adjoint operator
      \(H_{\rm phys,\infty}\) whose semigroup is the strong limit of
      the blocked physical transfers;
\item a limiting protected constraint kernel invariant under
      \(H_{\rm phys,\infty}\), with the intended first-class algebra
      on the common constraint core;
\item finite-time Heisenberg evolution agreeing with the intended
      linearized Einstein--SM Hamiltonian flow on every protected
      observable, in the limit of \((74.1)\);
\item a limiting vacuum stress and effective metric Euler covector
      satisfying
      \[
       {\cal E}^{\rm eff}_{q,\infty}=0                    \tag{74.4}
      \]
      on the declared determining metric panel.
\end{enumerate}
Isolated finite physical eigenvalues and their multiplicities
converge on every energy window controlled by the compact graph
screen.
\end{theorem}

\begin{proof}
By \textup{(P1)}, each finite cylinder has a positive OS form.
The Grand Tensor reflection-positive projective-limit theorem,
summable total-variation correction, and tightness assumptions in
\textup{(P6)} and \textup{(P10)} produce a subsequential probability
limit which is reflection positive on every protected
positive-time cylinder.  The protected covariance lower bound and
uniform integrability pass a nonzero variance to the limit, proving
nontriviality.

By V61, \textup{(P2)} produces the positive finite transfers.  V62
shows that submitted path-law errors vanish at the resolvent scale.
V66 identifies coarse and fine physical durations and removes the
time-block defect.  V67 descends the transfer through the exact OS
null quotient and Haar reduction; its leakage rate makes the physical
and full resolvents agree asymptotically.  The V68 Mosco descent
applied to \textup{(P5)} therefore yields a closed nonnegative
physical limiting form.  Its representation theorem gives the
nonnegative self-adjoint \(H_{\rm phys,\infty}\), and the inherited
Mosco theorem gives strong semigroup convergence.  The V59
floor-free comparison shows that the exact transfer logarithms have
the same limit.

For constraints, the V71 leakage bound and \((74.2)\) make the
off-kernel part of the generator vanish.  Graph convergence of the
constraint operators passes their common kernels and the joint
closure Gram to the protected core.  Hence the limiting constraint
kernel is invariant and the intended first-class commutators hold
there.

For observables, V69 gives, at every finite \(n\),
\[
 \|\alpha_t^{H_n}({\cal Q}_n(a))
       -{\cal Q}_n(e^{tL_n}a)\|
 \le
 te^{t\|L_n\|}
 \sqrt{\Delta_{{\rm corr},n}}\|a\|.
\]
The identifiability frame in \((74.1)\), the protected-core
convergence, and strong semigroup convergence pass the intended
finite-time flow to the limit.

Finally, V73 passes the isolated ground energy, vacuum stress, and
second metric jet under \textup{(P9)}.  It also passes the finite
effective Einstein covectors and their faithful metrics.  Equation
\((74.3)\) therefore implies \((74.4)\).  The compact graph-screen
branch of V68 upgrades physical resolvents to norm convergence on
controlled windows, so their isolated spectral projections and
multiplicities converge.
\end{proof}

\subsection{Necessity of the separate rows}

\begin{proposition}[No listed convergence row subsumes the others]
\label{prop:v14v74-independence}
The following failures are logically independent:
\begin{enumerate}
\item reflection-positive measure convergence without a Hamiltonian
      limit, as in V64;
\item Mosco generator convergence with escaping low-energy
      multiplicity, as in V68;
\item exact generator correspondence with a failed constraint
      algebra, as in V71;
\item centered dynamical correspondence with an arbitrary vacuum
      stress, as in V72;
\item convergent vacuum energy and stress with divergent curvature,
      as in V73.
\end{enumerate}
Therefore none of \textup{(P5)}--\textup{(P9)} can be deleted from
the stated conclusion.
\end{proposition}

\begin{proof}
Each cited finite counterexample satisfies the earlier property in
its item and violates the later one.  Taking direct sums with fixed
spectator sectors supplies any unrelated rows without changing that
failure.  Hence the logical implications do not hold.
\end{proof}

\subsection{Exact status against the attached sources}

\begin{assessment}[Populated and unpopulated rows]
\label{ass:v14v74-status}
The attached manuscripts prove general compilers for reflection
Grams, projective limits, tightness, Mosco convergence, finite HDA,
Ward/BRST identities, stress, and an expectation-level Einstein
residual.  The present cycle adds the adjacent-law transfer compiler,
floor-free logarithmic convergence, time blocking, gauge/Mosco
descent, generator and constraint correspondence, and absolute
vacuum-stress transport.

For the active Einstein--SM regulator, the following load-bearing
objects remain unpopulated:
\[
\begin{gathered}
 \mu_n,\quad J_n,\quad K_n,\quad T_n,\quad a_n,\quad
 Q_nT_n-T_nQ_n,\quad \mathfrak a_n,\\
 {\cal Q}_n,\quad L_n,\quad {\cal C}_n,\quad B_n,\quad
 D_n^{\rm con},\quad
 E_{0,n},\quad DH_n,\quad D^2H_n,\quad
 {\cal E}^{\rm eff}_{q,n}.
                                                               \tag{74.5}
\end{gathered}
\]
Thus Theorem~\ref{thm:v14v74-protected} is a conditional closure
theorem, not a completed construction.
\end{assessment}

\begin{firewall}[What the protected theorem still does not prove]
\label{fire:v14v74-scope}
Even with \textup{(P1)}--\textup{(P10)}, the theorem as stated does
not prove full Euclidean invariance, local covariance on every
observable, uniqueness of the continuum measure, clustering, a mass
gap, analytic continuation to a complete Lorentzian theory,
asymptotic particle states, numerical Standard Model coupling
matching, or nonlinear Einstein evolution outside the protected
panel.
\end{firewall}

\begin{decision}[V75 audit, then upstream finite weight construction]
\label{dec:v14v74-next}
At V75 perform the compulsory companion audit.  Then go upstream of
the missing adjacent law: in V76 construct a reflection-positive
one-step kernel from a finite common Euclidean action by an exact
boundary factorization
\[
 J(x,y)=Z^{-1}e^{-V(x)/2}K(x,y)e^{-V(y)/2}.
\]
Identify sufficient conditions on the kinetic kernel \(K\), gauge
action, fermion phase, and potential \(V\), and state which active
coefficients remain to be supplied.
\end{decision}

\section*{Version V74 append-only record}
\addcontentsline{toc}{section}{Version V74 append-only record}
The complete V73 source was retained before this continuation.  It
had \(876917\) bytes, \(26651\) lines, and SHA--256
\[
 \texttt{360682B0FA979456D979C35F42FDA18981D3B6D3155131F18F38E7B86F66568D}.
\]
V74 assembled the reflected transfer, gauge, constraint,
correspondence, and Einstein rows into one conditional protected
continuum theorem.
\section{Fourteenth-cycle V75: compulsory companion audit and kernel-positivity decision}
\label{sec:v14v75}

This is the compulsory audit after V70.  It is performed before the
reflection-factorized adjacent-law construction.

\begin{record}[Live companion provenance at V75]
\label{rec:v14v75-provenance}
The active Yang--Mills file was
\[
 \path{Renewal_Yang_Mills_Mass_Gap_Research_Notes_V24.tex},
\]
with \(193362\) bytes, \(5277\) lines, and SHA--256
\[
 \texttt{B37D902F84D9F324597B952F2F671D6A935C4042E0AC422AD1C60B1F1558B032}.
                                                               \tag{75.1}
\]
The active Navier--Stokes file remained
\[
 \path{Renewal_NS_Stability_Research_Notes_V26.tex},
\]
with \(278283\) bytes, \(5319\) lines, and SHA--256
\[
 \texttt{82C887F8C2C69E4F28FAD7C3DC8DE5B9099CB5A4BF431EBA1FFF3E91935978AD}.
                                                               \tag{75.2}
\]
\end{record}

\begin{assessment}[Yang--Mills V24]
\label{ass:v14v75-ym}
Yang--Mills V24 groups the complete transformed Laurent pencil as a
matrix polynomial affine in six sine/cosine variables.  Positivity at
the \(64\) vertices of a rational outer rectangle then proves
positivity throughout that rectangle by matrix convexity.  This
recovers the mixed-coordinate cancellation lost by the V23
triangle-square bound and enlarges the simultaneous transverse
radius to the single-axis scale.  Its negative radius-\(1/4\) witness
belongs only to the outer relaxation and is correctly not promoted
to a physical counterexample.

The admissible import is the multi-affine vertex-hull method.  If the
V76 kinetic boundary kernel depends multi-affinely on finitely many
trigonometric cell variables, exact positivity at the outer
rectangle's vertices will certify the whole parameter cell.  No
Yang--Mills coefficient, phase, transform, or floor is imported.
\end{assessment}

\begin{assessment}[Navier--Stokes status]
\label{ass:v14v75-ns}
Navier--Stokes V26 is unchanged and contains no new reflected-kernel
input.  Nothing is imported.
\end{assessment}

\begin{decision}[V76 direction]
\label{dec:v14v75}
Construct the full adjacent kernel by a reflection factorization, not
by independently fitting its entries.  Prove positivity, marginal
normalization, injectivity, gauge descent, and the transfer formula.
Keep fermionic determinant/Pfaffian orientation as a separate
multiplicative reflected-kernel factor.  When the kinetic kernel is
multi-affine in a compact parameter cell, permit the V24 vertex-hull
method as an exact positivity certificate.  The next compulsory audit
is V80.
\end{decision}

\begin{firewall}[Claim boundary after V75]
\label{fire:v14v75}
The audit imports only a finite matrix-convexity method.  It supplies
no active SM--Einstein boundary kernel, absolute weight, fermion
phase, or continuum datum.
\end{firewall}

\section*{Version V75 append-only record}
\addcontentsline{toc}{section}{Version V75 append-only record}
The complete V74 source was retained before this continuation.  It
had \(887937\) bytes, \(26905\) lines, and SHA--256
\[
 \texttt{879BA5A5F22B8B5B495F32B69E41DB4A40B274DC3B0AA9A24301D3DA345A1389}.
\]
V75 performed the compulsory companion audit and selected the
reflection-factorized kernel construction.  The next audit is V80.
\section{Fourteenth-cycle V76: reflection-factorized adjacent laws from a finite common action}
\label{sec:v14v76}

V74 located the first unpopulated object at the level of an absolute
finite Euclidean cylinder.  V75 directed us upstream from the
adjacent law to a reflection factorization of the common action.
This section constructs the complete V61 path-law card from a
positive kinetic boundary kernel and a real potential.

\subsection{Exact boundary factorization}

Let \(X=\{1,\ldots,m\}\) be a finite boundary configuration set.
Let \(K=K^*\) be a matrix satisfying
\[
 K_{xy}\ge0,\qquad K\succeq0,                             \tag{76.1}
\]
and let \(V:X\to\mathbb R\).  For a physical step \(a>0\), put
\[
 A_V=\operatorname{diag}\bigl(e^{-aV(x)/2}\bigr),\qquad
 \widetilde J=A_VKA_V,                                   \tag{76.2}
\]
\[
 Z=\sum_{x,y}\widetilde J_{xy},\qquad
 J=Z^{-1}\widetilde J,\qquad
 \pi_x=\sum_yJ_{xy}.                                     \tag{76.3}
\]

\begin{theorem}[Reflection-factorized adjacent-law compiler]
\label{thm:v14v76-factor}
Assume \(K\ne0\) and every row of \(K\) has a positive entry.  Then:
\begin{enumerate}
\item \(Z>0\), \(J\) is a symmetric nonnegative probability matrix,
      and \(\pi_x>0\);
\item \(J\succeq0\);
\item
      \[
       P=D_\pi^{-1}J                                     \tag{76.4}
      \]
      is a stationary reversible Markov transfer which is positive
      and self-adjoint on \(L^2(\pi)\);
\item if \(K\succ0\), then \(J\succ0\), \(P\) is injective, and
      \(-a^{-1}\log P\) is a nonnegative self-adjoint generator.
\end{enumerate}
\end{theorem}

\begin{proof}
The diagonal \(A_V\) is positive and invertible.  Entrywise
nonnegativity of \(K\) gives that of \(\widetilde J\).
The row hypothesis and positive diagonal factors give a positive row
sum at every \(x\); hence \(Z>0\) and \(\pi_x>0\).
Symmetry and total normalization are immediate.

Congruence preserves Loewner positivity:
\[
 \widetilde J=A_VKA_V\succeq0,
\]
and division by \(Z>0\) gives \(J\succeq0\).  The V61 compiler then
proves the Markov, reversible, and positive-transfer assertions.
If \(K\succ0\), invertible congruence gives \(J\succ0\), and V61--V59
give the logarithmic generator.
\end{proof}

\begin{corollary}[Boundary action form]
\label{cor:v14v76-action}
Suppose
\[
 K_{xy}=e^{-S_{\rm kin}(x,y)},\qquad
 S_{\rm kin}(x,y)=S_{\rm kin}(y,x),                       \tag{76.5}
\]
and \(K\succeq0\).  Then
\[
 J_{xy}
 =Z^{-1}\exp\left[
 -\frac a2V(x)-S_{\rm kin}(x,y)-\frac a2V(y)\right].      \tag{76.6}
\]
Thus one symmetric finite Euclidean slab action produces both the
absolute adjacent law and its positive transfer.
\end{corollary}

\begin{proof}
Substitute \((76.5)\) into \((76.2)\)--\((76.3)\).
\end{proof}

\begin{firewall}[Pointwise action positivity is insufficient]
\label{fire:v14v76-entrywise}
The facts \(K_{xy}>0\) and \(S_{\rm kin}\in\mathbb R\) do not imply
\(K\succeq0\).  The Grand Tensor reflection-Gram counterexample
\(\left(\begin{smallmatrix}1&2\\2&1\end{smallmatrix}\right)\)
has positive entries and a negative eigenvalue.  The matrix
positivity in \((76.1)\) is a separate load-bearing certificate.
\end{firewall}

\subsection{Canonical kinetic kernels from local moves}

Let \({\cal L}\) be a real symmetric graph Laplacian on \(X\):
\[
 {\cal L}_{xy}\le0\ (x\ne y),\qquad
 {\cal L}\mathbf1=0,\qquad {\cal L}\succeq0.              \tag{76.7}
\]
Set
\[
 K_a=e^{-a{\cal L}}.                                     \tag{76.8}
\]

\begin{theorem}[Graph heat kernel supplies the kinetic factor]
\label{thm:v14v76-heat}
For every \(a>0\),
\[
 K_a=K_a^{\mathsf T}\succ0,\qquad
 (K_a)_{xy}\ge0,\qquad K_a\mathbf1=\mathbf1.              \tag{76.9}
\]
If the move graph is connected, every entry of \(K_a\) is strictly
positive.  Consequently \(K_a\) satisfies all hypotheses of
Theorem~\ref{thm:v14v76-factor}.
\end{theorem}

\begin{proof}
Spectral calculus for \({\cal L}\succeq0\) gives positive eigenvalues
\(e^{-a\lambda}>0\), symmetry, and \(K_a\mathbf1=\mathbf1\).
The matrix \(-{\cal L}\) has nonnegative off-diagonal entries and
generates a finite continuous-time Markov chain; its exponential is
entrywise nonnegative.  If the graph is connected, for every
\(x,y\) some power of the nonnegative off-diagonal move matrix has a
positive \(xy\) entry.  The exponential series then makes
\((K_a)_{xy}>0\).
\end{proof}

Let \(V\) also denote the diagonal multiplication matrix.  The
unnormalized symmetric transfer is the Strang factor
\[
 S_a=e^{-aV/2}e^{-a{\cal L}}e^{-aV/2}.                    \tag{76.10}
\]

\begin{theorem}[Continuum generator of the finite factorization]
\label{thm:v14v76-strang}
Assume \({\cal L}\succeq0\) and \(V\succeq0\).  Then \(S_a\) is an
injective positive contraction and
\[
 a^{-1}(I-S_a)\longrightarrow{\cal L}+V                  \tag{76.11}
\]
in operator norm as \(a\downarrow0\).  Its exact logarithmic
generator satisfies
\[
 -a^{-1}\log S_a\longrightarrow{\cal L}+V                \tag{76.12}
\]
in norm-resolvent sense.
\end{theorem}

\begin{proof}
Each exponential in \((76.10)\) is a self-adjoint contraction, and
\[
 S_a=
 \left(e^{-a{\cal L}/2}e^{-aV/2}\right)^*
 \left(e^{-a{\cal L}/2}e^{-aV/2}\right)\succ0,
\]
so \(S_a\) is positive and injective.  Finite-dimensional Taylor
expansion gives
\[
 e^{-aV/2}=I-\frac a2V+O(a^2),\qquad
 e^{-a{\cal L}}=I-a{\cal L}+O(a^2).
\]
Multiplication yields
\[
 S_a=I-a({\cal L}+V)+O(a^2),
\]
which proves \((76.11)\).  Apply the V59 floor-free resolvent
comparison to \(S_a\): the exact logarithm and
\(a^{-1}(I-S_a)\) have resolvents differing by \(O(a)\).  Norm
convergence in \((76.11)\) then proves \((76.12)\).
\end{proof}

\begin{remark}[Normalization versus the symmetric transfer]
\label{rem:v14v76-normalization}
The Markov transfer \(P=D_\pi^{-1}J\) and the symmetric slab matrix
\(S_a\) use different Hilbert metrics and normalizations.  The former
has constant eigenvector and eigenvalue one; the latter retains the
absolute vacuum amplitude.  Passing from \(S_a\) to \(P\) is a
ground-state or Doob normalization and removes a scalar energy row.
For the Einstein source, both the normalizing eigenvalue and its
metric variation must be retained as in V72.
\end{remark}

\subsection{Potential lower shifts}

If \(V\) is merely bounded below on the finite set, put
\[
 c=\min_xV(x),\qquad \overline V=V-cI\succeq0.            \tag{76.13}
\]

\begin{proposition}[The contraction shift is an absolute source]
\label{prop:v14v76-shift}
The two symmetric transfers satisfy
\[
 S_a(V)=e^{-ac}S_a(\overline V),                          \tag{76.14}
\]
and their logarithmic generators differ by
\[
 -a^{-1}\log S_a(V)
 =
 -a^{-1}\log S_a(\overline V)+cI.                        \tag{76.15}
\]
Thus shifting the finite potential to prove contractivity is allowed
only when the scalar \(c\) and its metric first and second jets are
retained in the V72--V73 Einstein row.
\end{proposition}

\begin{proof}
The scalar \(cI\) commutes with every factor, giving
\(e^{-aV/2}=e^{-ac/2}e^{-a\overline V/2}\).
This proves \((76.14)\), and functional calculus gives
\((76.15)\).
\end{proof}

\subsection{Gauge covariance and exact quotient}

Let \(U(g)\) act on configurations by permutation matrices.

\begin{theorem}[Gauge-compatible boundary factorization]
\label{thm:v14v76-gauge}
If
\[
 U(g)K=KU(g),\qquad V(gx)=V(x)                            \tag{76.16}
\]
for every \(g\), then \(U(g)\) commutes with
\(\widetilde J,J,D_\pi,P\).  The Haar projector therefore reduces the
transfer exactly, and all V67 time-block and logarithm identities
hold on the gauge-invariant range.  Any deterministic
reflection-compatible coarse configuration pullback \(L\) produces
\[
 J_{\rm coarse}=L^*JL\succeq0.                            \tag{76.17}
\]
\end{theorem}

\begin{proof}
The potential condition makes \(A_V\) commute with every \(U(g)\).
Equation \((76.2)\) then gives covariance of \(\widetilde J\), and
normalization preserves it.  Row sums are constant on gauge orbits,
so \(D_\pi\) commutes as well.  The remaining claims follow from V67
and preservation of positivity under congruence.
\end{proof}

\subsection{Fermionic reflected factor}

Suppose a scalarized fermionic boundary factor \(F\) has already been
oriented on the same boundary state space.

\begin{proposition}[Schur-product fermion insertion]
\label{prop:v14v76-fermion}
If
\[
 F=F^*\succeq0,                                          \tag{76.18}
\]
then
\[
 J^{\rm BF}:=J\circ F\succeq0.                            \tag{76.19}
\]
If, in addition, every entry of \(J\circ F\) is real nonnegative and
the total mass is positive, normalization gives another
reflection-positive adjacent probability law.  Gauge covariance is
preserved when \(F\) is gauge invariant.
\end{proposition}

\begin{proof}
The Schur product theorem gives \((76.19)\).  The additional
entrywise and total-mass hypotheses are exactly those needed for an
absolute probability law.  Invariance is preserved entrywise.
\end{proof}

\begin{firewall}[Fermion phase remains separate]
\label{fire:v14v76-fermion}
A determinant or Pfaffian may be gauge covariant and nonzero while
failing either the Hermitian positivity in \((76.18)\) or the
entrywise scalar orientation needed for a probability.  Neither
condition follows from the bosonic factorization.  The determinant
line, zero modes, spectral flow, and cutoff phase holonomy remain
separate Grand Tensor rows.
\end{firewall}

\subsection{Multi-affine parameter certification}

\begin{corollary}[Vertex-hull kinetic certificate]
\label{cor:v14v76-vertex}
Suppose \(K(\theta)\) is Hermitian and affine separately in
\(r\) real variables on a rectangle, and suppose
\[
 K(v)\succeq\kappa I\quad(\kappa>0)                       \tag{76.20}
\]
at every one of its \(2^r\) vertices.  Then
\(K(\theta)\succeq\kappa I\) throughout the rectangle, and
Theorem~\ref{thm:v14v76-factor} gives an injective positive transfer
throughout the parameter cell.
\end{corollary}

\begin{proof}
The Yang--Mills V24 multi-affine hull writes every interior matrix as
a convex combination of its vertex matrices.  The Loewner cone is
convex.  Positive diagonal congruence by \(A_V\) preserves the
strict floor.
\end{proof}

\subsection{Automatic three-time card}

\begin{corollary}[Exact held-out Markov law]
\label{cor:v14v76-three}
The factorized adjacent law produces
\[
 K^{(3)}_{xyz}:=\pi_xP_{xy}P_{yz}.                        \tag{76.21}
\]
It is a normalized nonnegative three-time law with both adjacent
marginals \(J\), and its V61 Markov residual vanishes exactly.
\end{corollary}

\begin{proof}
Sum \((76.21)\) in the last or first variable and use stochasticity
and stationarity.  It equals the V61 Markov reconstruction by
definition.
\end{proof}

\begin{assessment}[What V76 actually populates]
\label{ass:v14v76-status}
V76 turns a positive kinetic kernel and a real finite potential into
the complete symbolic \((\pi,J,K^{(3)},P)\) path-law packet.  It also
gives a canonical sufficient kinetic kernel,
\(e^{-a{\cal L}}\), from any symmetric local-move graph Laplacian.
This is the first construction in the present block that goes
upstream of the missing adjacent law rather than listing it as an
input.

The active SM--Einstein regulator still lacks the concrete finite
boundary configuration set, local-move conductances defining
\({\cal L}\), common potential \(V\), absolute scalar shift,
fermionic reflected factor, and metric derivatives.  The Grand
Tensor action-rigidity theorem supplies allowed functional forms for
some local bosonic terms but not these numerical boundary matrices.
\end{assessment}

\begin{decision}[V77 local-move Laplacian from a typed action]
\label{dec:v14v76-next}
Construct \({\cal L}\) from a finite set of typed gauge, Higgs,
geometric, and admissible fermionic update moves with detailed-balance
conductances.  Prove locality, gauge covariance, connectedness on
each physical sector, and the exact kernel.  Separate the graph
connectivity gap from the continuum mass gap, and state the finite
coefficient card needed from the common action.
\end{decision}

\begin{firewall}[Claim boundary after V76]
\label{fire:v14v76}
The reflection factorization is a conditional finite construction.
No active configuration graph, conductances, common potential,
fermion orientation, or cofinal scaling has been supplied.  Its
finite connectivity or transfer positivity would not by itself prove
a continuum mass gap or the full SM--Einstein continuum.
\end{firewall}

\section*{Version V76 append-only record}
\addcontentsline{toc}{section}{Version V76 append-only record}
The complete V75 source was retained before this continuation.  It
had \(891234\) bytes, \(26986\) lines, and SHA--256
\[
 \texttt{A697F584EA8DAD5456C31DC6687A8603158CFC46FA7FC8FEC8BBF7EE75864BF3}.
\]
V76 constructed the reflection-factorized adjacent law, its graph
heat-kernel realization, gauge descent, and fermionic insertion
criteria.
\section{Fourteenth-cycle V77: typed local-move Laplacian and its exact physical sectors}
\label{sec:v14v77}

V76 reduces the kinetic boundary factor to a finite graph Laplacian.
We now construct that Laplacian from typed local configuration moves.
The construction fixes its exact kernel and gauge covariance and
makes explicit why a finite connectivity gap is not yet a physical
continuum mass gap.

\subsection{Reversible typed moves}

Let \(X\) be a finite configuration set and let
\(\rho_x>0\), \(\sum_x\rho_x=1\), be a declared reference law.
Let \({\cal E}\) be a finite set of unoriented moves
\(e=\{x,y\}\), each carrying:
\[
 \operatorname{type}(e)\in
 \{\mathrm{gauge},\mathrm{Higgs},\mathrm{geometry},
   \mathrm{fermion\ occupation}\},                       \tag{77.1}
\]
a finite support region \(\operatorname{supp}e\), and a conductance
\(c_e>0\).  The fermion-occupation type is included only after the
absolute determinant/Pfaffian line has been scalarized; Grassmann
phases are not Markov conductances.

Define on \(L^2(\rho)\)
\[
 ({\cal L}f)(x)
 :=
 \rho_x^{-1}\sum_{y:\{x,y\}\in{\cal E}}
 c_{\{x,y\}}\bigl(f(x)-f(y)\bigr).                       \tag{77.2}
\]

\begin{theorem}[Exact typed-move Dirichlet form]
\label{thm:v14v77-form}
The operator \({\cal L}\) is self-adjoint and nonnegative on
\(L^2(\rho)\), with
\[
 \boxed{
 \langle f,{\cal L}g\rangle_\rho
 =
 \sum_{\{x,y\}\in{\cal E}}
 c_{\{x,y\}}
 \overline{(f(x)-f(y))}(g(x)-g(y)).}                     \tag{77.3}
\]
Its kernel consists exactly of the functions constant on every
connected component of the move graph:
\[
 \ker{\cal L}
 =
 \{f:f(x)=f(y)\hbox{ whenever }x,y
       \hbox{ lie in the same component}\}.              \tag{77.4}
\]
\end{theorem}

\begin{proof}
Expand the left side of \((77.3)\), pair the two orientations of each
unoriented edge, and use symmetry of the conductance.  Equation
\((77.3)\) gives self-adjointness and
\[
 \langle f,{\cal L}f\rangle_\rho
 =\sum_e c_e|f(x_e)-f(y_e)|^2\ge0.
\]
It vanishes exactly when \(f\) is equal at the endpoints of every
edge.  Transitivity along paths gives constancy on components, and
every componentwise constant function plainly has zero energy.
\end{proof}

\begin{corollary}[Kinetic adjacent law]
\label{cor:v14v77-adjacent}
Let
\[
 P_a=e^{-a{\cal L}},\qquad
 J_a=D_\rho P_a.                                         \tag{77.5}
\]
Then \(P_a\) is a reversible Markov transfer, and
\[
 J_a=J_a^{\mathsf T}\succeq0,\qquad
 (J_a)_{xy}\ge0,\qquad
 \sum_y(J_a)_{xy}=\rho_x.                                \tag{77.6}
\]
It is injective for every \(a>0\).  If the move graph is connected,
\(P_a\) has a one-dimensional fixed space.
\end{corollary}

\begin{proof}
The off-diagonal entries of \(-{\cal L}\) are nonnegative, its rows
sum to zero, and it is self-adjoint in \(L^2(\rho)\).
Thus \(P_a\) is a positivity-preserving stochastic semigroup and is
reversible.  Its whitened representative
\[
 D_\rho^{1/2}P_aD_\rho^{-1/2}
 =e^{-aD_\rho^{1/2}{\cal L}D_\rho^{-1/2}}
\]
is positive definite.  Congruence gives \(J_a\succeq0\);
reversibility gives symmetry, and stochasticity gives the marginal.
The fixed space is \(\ker{\cal L}\), which is one-dimensional on a
connected graph by Theorem~\ref{thm:v14v77-form}.
\end{proof}

\subsection{Locality}

For a function \(O:X\to\mathbb C\), let \(M_O\) be multiplication by
\(O\).

\begin{proposition}[Exact move-local commutator]
\label{prop:v14v77-locality}
The off-diagonal matrix entries of the commutator are
\[
 [\,{\cal L},M_O\,]_{xy}
 =
 -\rho_x^{-1}c_{\{x,y\}}\bigl(O(y)-O(x)\bigr)
 \quad(x\ne y).                                          \tag{77.7}
\]
Hence only moves which change \(O\) contribute.  If \(O\) depends on
a spatial region \(B\) and every move disjoint from \(B\) leaves it
unchanged, then
\[
 [{\cal L},M_O]
 =
 \sum_{e:\operatorname{supp}e\cap B\ne\varnothing}
 [\,{\cal L}_e,M_O\,].                                   \tag{77.8}
\]
\end{proposition}

\begin{proof}
Insert \((77.2)\) into
\({\cal L}(Of)-O{\cal L}f\); the diagonal terms cancel and the
coefficient of \(f(y)\) is \((77.7)\).  If a move does not change any
coordinate read by \(O\), then \(O(x)=O(y)\), so its term vanishes.
\end{proof}

\begin{corollary}[Finite propagation in the move count]
\label{cor:v14v77-propagation}
If every move changes variables within diameter at most \(R\), then
the matrix entry of \({\cal L}^k\) between configurations which
require more than \(k\) allowed moves is zero.  Consequently the
heat kernel has the path-tail bound
\[
 |(e^{-a{\cal L}})_{xy}|
 \le\sum_{k\ge d_{\cal E}(x,y)}
 \frac{a^k\|{\cal L}\|^k}{k!}.                            \tag{77.9}
\]
\end{corollary}

\begin{proof}
A nonzero matrix entry of a product of \(k\) move matrices determines
a path of at most \(k\) moves.  Expand the exponential series and
bound each remaining term in operator norm.
\end{proof}

\subsection{Gauge covariance and physical components}

Let a compact or finite gauge group act on \(X\).  Assume the action
permutes moves and preserves their data:
\[
 \rho_{gx}=\rho_x,\qquad
 c_{\{gx,gy\}}=c_{\{x,y\}},\qquad
 \operatorname{type}(ge)=\operatorname{type}(e).          \tag{77.10}
\]

\begin{theorem}[Gauge covariance and quotient kernel]
\label{thm:v14v77-gauge}
The induced unitary representation \(U(g)\) on \(L^2(\rho)\) commutes
with \({\cal L}\), \(P_a\), and \(J_a\).  Let \(Q\) be the Haar
projector.  The kernel of
\[
 {\cal L}_{\rm phys}:=Q{\cal L}Q|_{Q L^2(\rho)}           \tag{77.11}
\]
consists of gauge-invariant functions constant on each move-graph
component.  Equivalently, its dimension is the number of connected
components of the quotient graph of gauge orbits linked by allowed
moves.
\end{theorem}

\begin{proof}
Changing variables \(x\mapsto gx\) in \((77.2)\) and using
\((77.10)\) gives \(U(g){\cal L}={\cal L}U(g)\).
Functional calculus gives covariance of \(P_a\), and invariance of
\(\rho\) gives covariance of \(J_a\).  The Haar projector therefore
reduces \({\cal L}\).  Intersect \((77.4)\) with the gauge-invariant
functions.  Such a function assigns one value to each gauge orbit
and is constant precisely along paths in the quotient move graph.
\end{proof}

\begin{remark}[Superselection ledger]
\label{rem:v14v77-superselection}
Every disconnected quotient component is an exact conserved sector.
It must be labelled by a protected charge, topology, boundary
condition, or other physical record.  Joining components by an
unobserved artificial move changes the theory; deleting a component
after seeing its spectrum is a post hoc selection.
\end{remark}

\subsection{A finite connectivity floor}

Let a connected component contain \(m\ge2\) states and let
\[
 c_{\min}:=\min_{e\ {\rm in\ component}}c_e,\qquad
 \rho_{\max}:=\max_{x\ {\rm in\ component}}\rho_x.        \tag{77.12}
\]

\begin{theorem}[Elementary component Poincare bound]
\label{thm:v14v77-poincare}
For every \(f\) with component mean zero,
\[
 \langle f,{\cal L}f\rangle_\rho
 \ge
 \frac{c_{\min}}{m(m-1)\rho_{\max}}\,
 \|f\|_{L^2(\rho)}^2.                                    \tag{77.13}
\]
Thus every finite connected component has a positive update gap.
\end{theorem}

\begin{proof}
Choose a spanning tree.  Since the unweighted mean can be subtracted
without changing edge differences, write
\[
 f(x)-\bar f
 =\frac1m\sum_y(f(x)-f(y)).
\]
Every tree path has at most \(m-1\) edges, so Cauchy--Schwarz gives
\[
 |f(x)-\bar f|^2
 \le(m-1)\sum_{e\ {\rm in\ tree}}|d_ef|^2.
\]
Summing in \(x\) gives
\[
 \sum_x|f(x)-\bar f|^2
 \le m(m-1)\sum_{e\ {\rm in\ tree}}|d_ef|^2.             \tag{77.14}
\]
The best constant shift for the \(\rho\)-norm is the \(\rho\)-mean,
so
\[
 \|f\|_{L^2(\rho)}^2
 \le\rho_{\max}\sum_x|f(x)-\bar f|^2.
\]
Finally,
\[
 \langle f,{\cal L}f\rangle_\rho
 \ge c_{\min}\sum_{e\ {\rm in\ tree}}|d_ef|^2.
\]
Combine these inequalities to obtain \((77.13)\).
\end{proof}

\begin{countertheorem}[Finite connectivity gaps can vanish cofinally]
\label{ctr:v14v77-path}
For the uniform nearest-neighbour path on \(m\) states with unit
conductances, the first nonzero unweighted graph-Laplacian eigenvalue
is
\[
 2\left(1-\cos\frac{\pi}{m}\right)\longrightarrow0.       \tag{77.15}
\]
Hence connectedness at every cutoff does not give a regulator-uniform
gap.
\end{countertheorem}

\begin{proof}
The discrete cosine vectors diagonalize the path Laplacian and give
the displayed eigenvalue.  Its convergence to zero follows from
continuity of cosine at zero.
\end{proof}

\subsection{Typed coefficient and provenance card}

\begin{acquisition}[Finite local-move card]
\label{acq:v14v77}
The active common action must supply, before the kinetic kernel is
claimed:
\begin{enumerate}
\item the finite boundary configuration set and faithful reference
      law \(\rho\);
\item every allowed move, its type, spatial support, and positive
      conductance;
\item the gauge action on configurations and the exact covariance
      identities \((77.10)\);
\item the quotient-component ledger and the protected physical label
      of every component;
\item the scaling of conductances with spatial cutoff and physical
      time step;
\item the absolute fermionic orientation before occupation moves are
      assigned positive conductances.
\end{enumerate}
The output is the exact local generator \({\cal L}\), its component
kernel, heat kernel, adjacent law, and finite Poincare bound.
\end{acquisition}

\begin{assessment}[Connectivity gap versus physical mass]
\label{ass:v14v77-status}
The update gap of Theorem~\ref{thm:v14v77-poincare} belongs to the
configuration move process.  It becomes a physical energy statement
only after the V69 correspondence residual identifies the resulting
transfer generator with the common Einstein--SM Hamiltonian on a
determining observable family.  Even then, V77 gives no uniform
cofinal lower bound; the path example proves that finite connectivity
alone cannot supply one.

The Grand Tensor source provides typed local action forms and
representation data, but it does not list an active boundary
configuration graph with numerical conductances and cutoff scaling.
Those are now the first missing inputs upstream of V76.
\end{assessment}

\begin{decision}[V78 bounded-potential comparison]
\label{dec:v14v77-next}
Quantify how the V76 potential tilt changes the reversible kinetic
law and its spectral gap.  Derive a finite comparison in terms of the
oscillation of \(V\), retain the absolute scalar shift for the
Einstein row, and show why a uniform kinetic gap plus an unbounded
potential oscillation gives no uniform physical gap.
\end{decision}

\begin{firewall}[Claim boundary after V77]
\label{fire:v14v77}
V77 constructs a typed finite move generator conditionally on a
supplied conductance ledger.  That ledger, its fermionic
scalarization, and its cofinal scaling are absent.  A finite update
gap is not the Standard Model--Einstein mass gap or the full
continuum.
\end{firewall}

\section*{Version V77 append-only record}
\addcontentsline{toc}{section}{Version V77 append-only record}
The complete V76 source was retained before this continuation.  It
had \(904214\) bytes, \(27351\) lines, and SHA--256
\[
 \texttt{35455D9DD66BE413FAD7AFB22028BA62F90B13565764AE81A67A9ED7E0509597}.
\]
V77 constructed the typed local-move Laplacian, its exact physical
sector kernel, locality, gauge covariance, and finite connectivity
floor.
\section{Fourteenth-cycle V78: bounded-potential comparison and the barrier soft mode}
\label{sec:v14v78}

V77 constructs the reversible kinetic adjacent law.  V76 tilts it by
the diagonal potential factors on the two sides of the reflection
plane.  We now quantify the resulting change of conductances,
marginal, and spectral gap.  The only scale entering the normalized
comparison is the potential oscillation; the absolute lower shift
remains a separate Einstein source.

\subsection{Tilted adjacent law}

Let \(J_0=J_0^{\mathsf T}\succeq0\) be a nonnegative adjacent
probability matrix with faithful marginal
\[
 \rho_x=\sum_y(J_0)_{xy}.
\]
For a real potential \(V\) and step \(a>0\), put
\[
 w_x=e^{-aV(x)/2},\qquad
 \widetilde J_V(x,y)=w_xJ_0(x,y)w_y,                     \tag{78.1}
\]
\[
 Z_V=\sum_{x,y}\widetilde J_V(x,y),\qquad
 J_V=Z_V^{-1}\widetilde J_V,\qquad
 \pi_V(x)=\sum_yJ_V(x,y).                                \tag{78.2}
\]
Let
\[
 \omega_V:=\max_xV(x)-\min_xV(x).                         \tag{78.3}
\]

\begin{theorem}[Exact density and conductance comparison]
\label{thm:v14v78-comparison}
For every \(x,y\),
\[
 e^{-a\omega_V}
 \le\frac{J_V(x,y)}{J_0(x,y)}
 \le e^{a\omega_V}                                       \tag{78.4}
\]
whenever \(J_0(x,y)>0\), and
\[
 \boxed{
 e^{-a\omega_V}\rho_x
 \le\pi_V(x)\le
 e^{a\omega_V}\rho_x.}                                   \tag{78.5}
\]
Consequently, for the Dirichlet forms
\[
 {\cal E}_0[f]
 =\frac1{2a}\sum_{x,y}J_0(x,y)|f_x-f_y|^2,
\]
\[
 {\cal E}_V[f]
 =\frac1{2a}\sum_{x,y}J_V(x,y)|f_x-f_y|^2,               \tag{78.6}
\]
one has
\[
 e^{-a\omega_V}{\cal E}_0[f]
 \le{\cal E}_V[f]
 \le e^{a\omega_V}{\cal E}_0[f].                         \tag{78.7}
\]
\end{theorem}

\begin{proof}
Let \(V_{\min}\) and \(V_{\max}\) be the extrema.  The two endpoint
factors in \((78.1)\) give
\[
 e^{-aV_{\max}}
 \le w_xw_y\le e^{-aV_{\min}}.
\]
Since \(J_0\) has total mass one,
\[
 e^{-aV_{\max}}\le Z_V\le e^{-aV_{\min}}.                \tag{78.8}
\]
Dividing the endpoint bound by \((78.8)\) gives \((78.4)\).
Summing it in \(y\) gives \((78.5)\), and multiplying by squared edge
differences and summing gives \((78.7)\).
\end{proof}

\subsection{Poincare and exact-log gaps}

Define
\[
 \operatorname{Var}_\rho(f)
 :=\inf_{c\in\mathbb C}\sum_x\rho_x|f_x-c|^2,
\]
and similarly for \(\pi_V\).  Let
\[
 \gamma_0
 :=\inf_{f\not\equiv{\rm const}}
 \frac{{\cal E}_0[f]}{\operatorname{Var}_\rho(f)},\qquad
 \gamma_V
 :=\inf_{f\not\equiv{\rm const}}
 \frac{{\cal E}_V[f]}{\operatorname{Var}_{\pi_V}(f)}.     \tag{78.9}
\]

\begin{theorem}[Bounded-potential gap comparison]
\label{thm:v14v78-gap}
One has
\[
 \boxed{
 e^{-2a\omega_V}\gamma_0
 \le\gamma_V
 \le e^{2a\omega_V}\gamma_0.}                            \tag{78.10}
\]
If the tilted transfer \(P_V=D_{\pi_V}^{-1}J_V\) is injective and
positive, then the spectral gap of its exact logarithmic generator
\[
 H_V=-a^{-1}\log P_V
\]
above constants is at least \(\gamma_V\), hence at least
\(e^{-2a\omega_V}\gamma_0\).
\end{theorem}

\begin{proof}
From \((78.5)\), for every scalar \(c\),
\[
 \sum_x\rho_x|f_x-c|^2
 \ge e^{-a\omega_V}
 \sum_x\pi_V(x)|f_x-c|^2.
\]
Taking the infimum gives
\[
 \operatorname{Var}_\rho(f)
 \ge e^{-a\omega_V}\operatorname{Var}_{\pi_V}(f).         \tag{78.11}
\]
Together with the lower half of \((78.7)\) and the base Poincare
inequality, this gives the lower half of \((78.10)\).  Interchanging
the two measures and forms gives the upper half.

If \(\mu\in(0,1]\) is a nonconstant transfer eigenvalue, its
Dirichlet-generator eigenvalue is \(a^{-1}(1-\mu)\), whereas its
exact-log eigenvalue is \(-a^{-1}\log\mu\).  Since
\(-\log\mu\ge1-\mu\), the logarithmic gap dominates the Dirichlet
gap.
\end{proof}

\begin{corollary}[Uniform bounded-oscillation handoff]
\label{cor:v14v78-uniform}
If
\[
 \inf_n\gamma_{0,n}\ge\gamma_*>0,\qquad
 \sup_na_n\omega_{V_n}\le C<\infty,                       \tag{78.12}
\]
then
\[
 \inf_n\gamma_{V,n}\ge e^{-2C}\gamma_*>0.                \tag{78.13}
\]
\end{corollary}

\begin{proof}
Apply the lower half of \((78.10)\) at each cutoff.
\end{proof}

\subsection{A three-state barrier}

The oscillation hypothesis cannot be removed, even when the kinetic
transfer is uniformly positive.

Let
\[
 P_0=
 \begin{pmatrix}
 3/4&1/4&0\\
 1/4&1/2&1/4\\
 0&1/4&3/4
 \end{pmatrix},\qquad
 \rho=(1/3,1/3,1/3),\qquad J_0=\frac13P_0.                \tag{78.14}
\]
The eigenvalues of \(P_0\) are
\[
 1,\quad\frac34,\quad\frac14.                            \tag{78.15}
\]
For \(R>0\), take
\[
 V_R=(0,R,0),\qquad t=e^{-aR/2}.                          \tag{78.16}
\]

\begin{countertheorem}[A high local barrier destroys the update gap]
\label{ctr:v14v78-barrier}
For the tilted law \((78.1)\)--\((78.2)\), the odd vector
\((1,0,-1)\) is a transfer eigenvector with eigenvalue
\[
 \mu_{\rm odd}(t)=\frac3{3+t}.                            \tag{78.17}
\]
Consequently
\[
 \gamma_V
 \le\frac{1-\mu_{\rm odd}(t)}a
 =\frac{t}{a(3+t)}.                                      \tag{78.18}
\]
Choose any \(a_n\downarrow0\) and \(R_n\) so that
\[
 t_n=e^{-a_nR_n/2}=a_n^2.                                \tag{78.19}
\]
Then the kinetic Dirichlet gap is
\[
 \gamma_{0,n}=\frac1{4a_n},
\]
which is uniformly bounded below and in fact diverges, while
\[
 \gamma_{V,n}\le\frac{a_n}{3+a_n^2}\longrightarrow0.      \tag{78.20}
\]
\end{countertheorem}

\begin{proof}
After the diagonal tilt, the unnormalized endpoint row sum is
\[
 \frac13\left(\frac34+\frac14t\right),
\]
while the middle row sum is
\[
 \frac13\left(\frac12t+\frac12t^2\right).
\]
Normalization cancels in transition ratios.  Hence
\[
 P_V(0,1)=P_V(2,1)=\frac{t}{3+t}.                         \tag{78.21}
\]
Symmetry between the two endpoints makes \((1,0,-1)\) an eigenvector,
and its endpoint component is multiplied by
\(1-t/(3+t)=3/(3+t)\), proving \((78.17)\).
The Rayleigh principle gives \((78.18)\).

Equation \((78.15)\) makes the smallest nonzero eigenvalue of
\(a_n^{-1}(I-P_0)\) equal to \(1/(4a_n)\).  Substitution of
\((78.19)\) into \((78.18)\) gives \((78.20)\).
\end{proof}

\begin{assessment}[Barrier meaning]
\label{ass:v14v78-barrier}
The soft mode is the difference between the two low-potential
endpoints.  Communication between them must cross the suppressed
middle state.  Every finite tilted transfer remains positive and
injective, but its gap collapses.  Thus finite connectivity of the
typed move graph and a uniform kinetic gap do not control the
physical tilted gap without a potential-oscillation or comparable
conductance-profile bound.

This is an update-process statement.  It becomes a particle mass
statement only after the V69 correspondence and the V68 compact
spectral screen pass.
\end{assessment}

\subsection{Absolute shift and Einstein source}

Write
\[
 V=\overline V+cI,\qquad
 c=\min_xV(x),\qquad \min\overline V=0.                   \tag{78.22}
\]
The oscillation \(\omega_V\) is unchanged by \(c\).

\begin{proposition}[Gap comparison is blind to the absolute potential]
\label{prop:v14v78-shift}
The normalized adjacent law \(J_V\), its marginal, Markov transfer,
Dirichlet form, and normalized gap depend only on
\(\overline V\).  The unnormalized symmetric transfer is multiplied
by \(e^{-ac}\), and its logarithmic generator gains \(cI\).
Therefore \(c\), \(dc\), and \(D^2c\) remain compulsory V72--V73
Einstein rows although they do not enter \((78.10)\).
\end{proposition}

\begin{proof}
Both endpoint factors in \((78.1)\) acquire the common multiplier
\(e^{-ac/2}\), so \(\widetilde J\) and \(Z\) acquire the same total
factor \(e^{-ac}\), which cancels from \(J_V\).  The statement for
the unnormalized transfer and its logarithm is V76's scalar-shift
identity.
\end{proof}

\subsection{Finite acquisition consequence}

\begin{acquisition}[Potential-profile card]
\label{acq:v14v78}
In addition to the V77 conductance ledger, the active finite common
action must supply:
\begin{enumerate}
\item the real potential value on every retained boundary
      configuration;
\item its absolute minimum and metric first and second jets;
\item an outward bound on \(a_n\omega_{V_n}\), or a sharper direct
      conductance comparison;
\item the tilted marginal support profile;
\item a compact-screen or conductance argument excluding barrier
      soft modes on the physical sector.
\end{enumerate}
The normalized gap certificate is \((78.10)\); the absolute shift is
sent separately to the effective Einstein row.
\end{acquisition}

\begin{assessment}[Current source boundary]
\label{ass:v14v78-status}
The attached sources specify allowed local gauge, Higgs, and
geometric action forms but do not list the potential on a finite
active boundary configuration set, its oscillation scaling, or its
absolute metric jets.  Therefore neither \((78.12)\) nor the barrier
exclusion can presently be evaluated.
\end{assessment}

\begin{decision}[V79 tensorized locality versus interacting barriers]
\label{dec:v14v78-next}
Treat a finite spatial product of local boundary cells.  Prove the
spectral gap of the independent tensorized move generator, then add
bounded finite-range interactions and derive a comparison in terms
of local conditional oscillations rather than the total extensive
oscillation.  This is necessary because \(\operatorname{osc}V\)
normally grows with volume and makes \((78.10)\) useless
thermodynamically.
\end{decision}

\begin{firewall}[Claim boundary after V78]
\label{fire:v14v78}
The bounded-potential comparison is finite and conditional.  The
active local potential, conditional oscillation bounds, fermionic
orientation, thermodynamic limit, physical correspondence, and
continuum mass gap remain unpopulated.
\end{firewall}

\section*{Version V78 append-only record}
\addcontentsline{toc}{section}{Version V78 append-only record}
The complete V77 source was retained before this continuation.  It
had \(915615\) bytes, \(27682\) lines, and SHA--256
\[
 \texttt{73C180E263B3AA8484C0BA02CCD525693CAB89ED37908DD7B62AD2978ECCE45F}.
\]
V78 proved the bounded-potential gap comparison and an explicit
barrier-induced cofinal soft mode.
\section{Fourteenth-cycle V79: tensorized locality and the phase-coexistence row}
\label{sec:v14v79}

V78 controlled one tilted boundary cell by the total potential
oscillation.  On a spatial volume \(\Lambda\), however, an extensive
potential normally has
\(\operatorname{osc}V_\Lambda=O(|\Lambda|)\).  Direct use of
\((78.10)\) then loses all thermodynamic information.  This
continuation replaces that global number by two logically separate
objects:
\[
 \text{a worst one-cell conditional gap}
 \quad\hbox{and}\quad
 \text{an approximate-tensorization constant}.           \tag{79.1}
\]
The first object is local and can be estimated by V78.  The second
detects collective phase bottlenecks and cannot be inferred from
bounded local couplings alone.

\subsection{Exact independent tensorization}

Let \(\Lambda\) be finite.  For \(x\in\Lambda\), let
\((X_x,\mu_x)\) be a finite probability space and let \(A_x\) be a
nonnegative self-adjoint operator on \(L^2(\mu_x)\) satisfying
\[
 A_x{\bf1}=0,\qquad
 \langle g,A_xg\rangle_{\mu_x}
 \ge\gamma_x\operatorname{Var}_{\mu_x}(g),\qquad
 \gamma_x>0.                                             \tag{79.2}
\]
Put
\[
 X_\Lambda=\prod_{x\in\Lambda}X_x,\qquad
 \mu_0=\bigotimes_{x\in\Lambda}\mu_x,\qquad
 A_0=\sum_{x\in\Lambda}A_x^{(x)},                         \tag{79.3}
\]
where \(A_x^{(x)}\) acts only in the \(x\)-th factor.

\begin{theorem}[Exact product gap]
\label{thm:v14v79-product}
The kernel of \(A_0\) is the tensor product of the local kernels.  If
each local kernel consists of constants, then
\[
 \operatorname{gap}(A_0)=\min_{x\in\Lambda}\gamma_x.       \tag{79.4}
\]
For the random-scan normalization
\(|\Lambda|^{-1}A_0\), the gap is instead
\(|\Lambda|^{-1}\min_x\gamma_x\).
\end{theorem}

\begin{proof}
Choose an orthonormal eigenbasis
\(\{e_{x,j}\}_{j\ge0}\) for each finite-dimensional \(A_x\), with
eigenvalues
\[
 0=\lambda_{x,0}<\lambda_{x,1}\le\lambda_{x,2}\le\cdots .
\]
The tensor products
\(\bigotimes_xe_{x,j_x}\) form an orthonormal eigenbasis for \(A_0\),
and their eigenvalues are
\[
 \lambda_{\boldsymbol j}=\sum_{x\in\Lambda}\lambda_{x,j_x}. \tag{79.5}
\]
The zero eigenspace is therefore precisely the tensor product of the
local zero eigenspaces.  When those spaces are one-dimensional, the
smallest positive sum in \((79.5)\) is obtained by exciting one
factor whose first positive eigenvalue is minimal.  This proves
\((79.4)\); scalar rescaling proves the final statement.
\end{proof}

\begin{assessment}[Update normalization]
\label{ass:v14v79-normalization}
The sum normalization in \((79.3)\) is the natural normalization for
a local Hamiltonian: every cell carries its own rate and a
volume-independent product gap is possible.  The random-scan chain
performs only one uniformly selected update per unit time and
therefore has the kinematic factor \(|\Lambda|^{-1}\).  A gap claim
must declare which clock is being used.
\end{assessment}

\subsection{Conditional fibers and approximate tensorization}

Let \(\mu\) now be a strictly positive probability measure on
\(X_\Lambda\).  Write \(\xi=\eta_{\Lambda\setminus\{x\}}\) for an
exterior configuration and let \(\mu_x^\xi\) be the conditional law
of \(\eta_x\) given \(\xi\).  Suppose a nonnegative reversible
one-cell operator \(A_x^\xi\) on \(L^2(\mu_x^\xi)\) is supplied for
every \((x,\xi)\), with
\[
 A_x^\xi{\bf1}=0,\qquad
 \langle g,A_x^\xi g\rangle_{\mu_x^\xi}
 \ge\gamma_x(\xi)\operatorname{Var}_{\mu_x^\xi}(g).       \tag{79.6}
\]
Define the global form
\[
 \mathcal E_\mu(f)
 :=
 \sum_{x\in\Lambda}
 \mathbb E_{\xi\sim\mu_{\Lambda\setminus\{x\}}}
 \left[
  \left\langle f(\,\cdot\,,\xi),
  A_x^\xi f(\,\cdot\,,\xi)\right\rangle_{\mu_x^\xi}
 \right].                                                \tag{79.7}
\]
It is a nonnegative symmetric finite-dimensional form and hence
defines a unique nonnegative self-adjoint generator \(A_\mu\).

\begin{definition}[Approximate tensorization constant]
\label{def:v14v79-at}
The least \(C_{\rm AT}(\mu)\in[1,\infty]\) for which
\[
 \operatorname{Var}_\mu(f)
 \le C_{\rm AT}(\mu)
 \sum_{x\in\Lambda}
 \mathbb E_\mu\!
 \left[
  \operatorname{Var}_{\mu_x^{\eta_{\ne x}}}
  \bigl(f(\,\cdot\,,\eta_{\ne x})\bigr)
 \right]                                                 \tag{79.8}
\]
holds for every \(f\) is the approximate-tensorization constant.
For a product measure, the Efron--Stein martingale decomposition
gives \(C_{\rm AT}=1\).
\end{definition}

\begin{theorem}[Local-to-global conditional gap]
\label{thm:v14v79-local-global}
Set
\[
 \gamma_{\rm fib}:=
 \inf_{\substack{x\in\Lambda\\
                  \xi\in X_{\Lambda\setminus\{x\}}}}
 \gamma_x(\xi).                                          \tag{79.9}
\]
If \(\gamma_{\rm fib}>0\) and \(C_{\rm AT}(\mu)<\infty\), then
\[
 \operatorname{gap}(A_\mu)
 \ge\frac{\gamma_{\rm fib}}{C_{\rm AT}(\mu)}.             \tag{79.10}
\]
Consequently a cofinal family has a uniform sum-clock gap whenever
both
\[
 \inf_n\gamma_{{\rm fib},n}>0,\qquad
 \sup_n C_{\rm AT}(\mu_n)<\infty                         \tag{79.11}
\]
hold.
\end{theorem}

\begin{proof}
Apply \((79.6)\) inside every summand of \((79.7)\):
\[
 \mathcal E_\mu(f)
 \ge\gamma_{\rm fib}
 \sum_x\mathbb E_\mu
 \left[
  \operatorname{Var}_{\mu_x^{\eta_{\ne x}}}
  \bigl(f(\,\cdot\,,\eta_{\ne x})\bigr)
 \right].
\]
Equation \((79.8)\) makes the right-hand side at least
\(\gamma_{\rm fib}C_{\rm AT}^{-1}\operatorname{Var}_\mu(f)\).
The Rayleigh characterization of the first nonzero eigenvalue gives
\((79.10)\).
\end{proof}

\subsection{Conditional oscillations from a finite-range action}

Assume that, after the geometric and gauge constraints have been
restricted to an admissible exterior fiber, the \(x\)-move transfer
is a V78 endpoint tilt of a base reversible transfer.  Let
\(\gamma^0_x(\xi)\) be its base Dirichlet gap and let
\(V_x(s;\xi)\) be the real potential seen by the moving cell.  Define
the dimensionless conditional oscillation
\[
 \Theta_x
 :=
 \sup_\xi
 a\,\left(
       \max_{s\in X_x}V_x(s;\xi)
       -\min_{s\in X_x}V_x(s;\xi)
    \right).                                             \tag{79.12}
\]

\begin{proposition}[Conditional V78 comparison]
\label{prop:v14v79-conditional}
If the tilted fiber marginal in V78 is the conditional law
\(\mu_x^\xi\) used in \((79.7)\), then
\[
 \gamma_x(\xi)
 \ge e^{-2\Theta_x}\gamma_x^0(\xi).                       \tag{79.13}
\]
Thus, with
\[
 \gamma^0_{\rm fib}:=\inf_{x,\xi}\gamma_x^0(\xi),
 \qquad \Theta_*:=\sup_x\Theta_x,
\]
one has
\[
 \operatorname{gap}(A_\mu)
 \ge
 \frac{e^{-2\Theta_*}\gamma^0_{\rm fib}}
      {C_{\rm AT}(\mu)}.                                 \tag{79.14}
\]
\end{proposition}

\begin{proof}
Freeze \(x\) and \(\xi\).  V78, Theorem
\ref{thm:v14v78-comparison}, applied on that finite fiber gives
\[
 \gamma_x(\xi)
 \ge
 \exp\!\left[
  -2a\,\operatorname{osc}_{s}V_x(s;\xi)
 \right]\gamma_x^0(\xi).
\]
Use \((79.12)\), take the infimum over fibers, and apply
Theorem \ref{thm:v14v79-local-global}.
\end{proof}

Suppose the dimensionless interaction is written
\[
 aV_\Lambda(\eta)=\sum_{B\subseteq\Lambda}\Phi_B(\eta_B), \tag{79.15}
\]
and only sets \(B\) containing \(x\) change during an \(x\)-move.
Set
\[
 \operatorname{osc}_x\Phi_B
 :=
 \sup_{\xi}
 \left[
  \max_{s\in X_x}\Phi_B(s,\xi_{B\setminus\{x\}})
  -
  \min_{s\in X_x}\Phi_B(s,\xi_{B\setminus\{x\}})
 \right].                                                \tag{79.16}
\]

\begin{lemma}[Finite-range conditional budget]
\label{lem:v14v79-range}
For every \(x\),
\[
 \Theta_x
 \le\sum_{B\ni x}\operatorname{osc}_x\Phi_B.              \tag{79.17}
\]
In particular, if at most \(D\) interaction terms meet a cell and
each has one-cell oscillation at most \(K\), then
\[
 \Theta_*\le DK,                                         \tag{79.18}
\]
independently of \(|\Lambda|\).  Meanwhile
\(\operatorname{osc}(aV_\Lambda)\) may grow like
\(|\Lambda|\).
\end{lemma}

\begin{proof}
For fixed exterior data and two values \(s,t\in X_x\), terms not
containing \(x\) cancel and
\[
 aV_\Lambda(s,\xi)-aV_\Lambda(t,\xi)
 =
 \sum_{B\ni x}
 \bigl[
  \Phi_B(s,\xi_{B\setminus\{x\}})
  -\Phi_B(t,\xi_{B\setminus\{x\}})
 \bigr].
\]
Take the supremum over \(s,t,\xi\) and use the triangle inequality.
The bounded-incidence conclusion follows immediately.
\end{proof}

\begin{assessment}[What locality does and does not buy]
\label{ass:v14v79-locality}
Equation \((79.18)\) repairs the extensive exponential in V78, but
only the numerator of \((79.14)\).  A collective two-phase
bottleneck appears in \(C_{\rm AT}(\mu)\), even when every cell has
bounded degree and every one-cell conditional problem is perfectly
gapped.
\end{assessment}

\subsection{An explicit bounded-degree phase bottleneck}

The missing collective row is not a technical artifact.  Let
\(G_N=(\Lambda_N,E_N)\) be a sequence of \(d\)-regular graphs with
odd \(N=|\Lambda_N|\) and edge-expansion constant \(h>0\):
\[
 |\partial S|\ge h\min\{|S|,N-|S|\}
 \quad\hbox{for every }S\subseteq\Lambda_N.               \tag{79.19}
\]
On \(X_N=\{-1,+1\}^{\Lambda_N}\), define
\[
 \mu_{N,\beta}(\sigma)
 =
 Z_{N,\beta}^{-1}
 \exp\!\left(\beta\sum_{\{x,y\}\in E_N}\sigma_x\sigma_y\right).
                                                               \tag{79.20}
\]
Let \(E_x\) be conditional expectation in spin \(x\), and take the
sum-clock heat-bath generator
\[
 A_{N,\beta}=\sum_{x\in\Lambda_N}(I-E_x).                 \tag{79.21}
\]
Each conditional operator \(I-E_x\) is an orthogonal projection, so
its nonzero conditional eigenvalue and conditional gap are exactly
one.  Flipping \(x\) changes the exponent in \((79.20)\) by at most
\(2\beta d\), uniformly in \(N\).

\begin{countertheorem}[Uniform local gaps do not exclude a global
two-phase soft mode]
\label{ctr:v14v79-expander}
If
\[
 \beta h>\log2,                                          \tag{79.22}
\]
then
\[
 \operatorname{gap}(A_{N,\beta})
 \le
 2^{N+1}N\,e^{-\beta h(N-1)},                            \tag{79.23}
\]
which tends to zero exponentially.  Consequently
\[
 C_{\rm AT}(\mu_{N,\beta})
 \ge
 \frac{e^{\beta h(N-1)}}{2^{N+1}N},                      \tag{79.24}
\]
although the degree, the conditional interaction oscillation, and
all one-cell heat-bath gaps are uniform.
\end{countertheorem}

\begin{proof}
Let
\[
 M(\sigma)=\sum_x\sigma_x,\qquad
 F(\sigma)={\bf1}_{\{M(\sigma)>0\}}.
\]
Because \(N\) is odd and \(\mu_{N,\beta}\) is invariant under global
spin flip,
\[
 \operatorname{Var}_{\mu_{N,\beta}}(F)=\frac14.           \tag{79.25}
\]
A one-spin fiber contributes to
\(\langle F,(I-E_x)F\rangle\) only when its two endpoints have
magnetizations \(+1\) and \(-1\).  The elementary conditional
variance formula on a two-point fiber bounds that contribution by
the total probability of the two endpoints.  Hence
\[
 \langle F,A_{N,\beta}F\rangle
 \le N\,\mu_{N,\beta}(|M|=1).                             \tag{79.26}
\]

For a configuration with \(|M|=1\), the minority-spin set \(S\) has
cardinality \((N-1)/2\).  Relative to the all-plus configuration,
each disagreeing edge lowers
\(\sum_{\{x,y\}}\sigma_x\sigma_y\) by \(2\).  Expansion gives
\[
 |\partial S|\ge \frac{h(N-1)}2,
\]
so every such configuration has weight at most
\[
 e^{\beta|E_N|}e^{-\beta h(N-1)}.                         \tag{79.27}
\]
There are at most \(2^N\) configurations, while the two constant
configurations give
\[
 Z_{N,\beta}\ge2e^{\beta|E_N|}.
\]
Therefore
\[
 \mu_{N,\beta}(|M|=1)
 \le2^{N-1}e^{-\beta h(N-1)}.                             \tag{79.28}
\]
The Rayleigh quotient from \((79.25)\)--\((79.28)\) is bounded by
the right-hand side of \((79.23)\).

For heat-bath fibers, the right-hand sum in \((79.8)\) equals
\(\langle f,A_{N,\beta}f\rangle\).  Applying \((79.8)\) to \(F\)
shows
\(C_{\rm AT}\ge\operatorname{Var}(F)/
 \langle F,A_{N,\beta}F\rangle\); use \((79.23)\).
\end{proof}

\begin{assessment}[Spatial interpretation]
\label{ass:v14v79-phase}
The expander example is a finite bounded-degree local system and
already proves the logical independence of the collective row.
For a Euclidean lattice, the corresponding obstruction is an
interface free-energy estimate rather than the volume-expansion
estimate \((79.19)\).  The continuum programme must therefore
supply either a uniform approximate-tensorization estimate in the
selected phase or a different conductance argument that controls
interfaces, topological sectors, and gauge-reduced components.
\end{assessment}

\subsection{Physical acquisition ledger}

\begin{acquisition}[Thermodynamic locality card]
\label{acq:v14v79}
For each finite cutoff and admissible reduced exterior fiber, record:
\begin{enumerate}
\item the sum-clock or random-scan normalization;
\item the base one-cell gap
      \(\gamma_x^0(\xi)\);
\item the conditional action oscillation \(\Theta_x\), with the
      incident interaction terms used in \((79.17)\);
\item compatibility of the local tilted marginal with the global
      conditional law;
\item a quantitative \(C_{\rm AT}(\mu_\Lambda)\), or a replacement
      global conductance/profile bound;
\item the selected phase, boundary condition, and any interface or
      topological-sector exclusion;
\item the separate absolute action shift and its metric jets for the
      V72--V73 Einstein source row.
\end{enumerate}
Only after all seven entries are populated does \((79.14)\) yield a
volume-uniform transfer gap.
\end{acquisition}

\begin{assessment}[Current source boundary]
\label{ass:v14v79-source}
The attached grand-tensor and Einstein--SM sources describe local
interaction types, but they do not provide a finite reduced
conditional law, a uniform fiber-gap table, conditional
oscillations, or an approximate-tensorization/interface estimate.
The thermodynamic gap row therefore remains conditional.  The
counterexample above forbids filling it from locality alone.
\end{assessment}

\begin{decision}[V80 five-version audit]
\label{dec:v14v79-next}
Inspect the newest Yang--Mills and Navier--Stokes companion notes as
required by the five-version cadence.  Test specifically whether
their newest contraction, sector, or local-vacuum mechanisms supply
a quantitative substitute for \(C_{\rm AT}\), a phase-selection
criterion, or a reduced conditional boundary law.  Then continue
with the first genuinely shared bottleneck.
\end{decision}

\begin{firewall}[Claim boundary after V79]
\label{fire:v14v79}
Equations \((79.4)\), \((79.10)\), and \((79.14)\) are finite
spectral statements under their displayed hypotheses.  No
thermodynamic SM gap, physical particle mass, interacting OS
positivity, Einstein source limit, or full SM--Einstein continuum is
claimed.
\end{firewall}

\section*{Version V79 append-only record}
\addcontentsline{toc}{section}{Version V79 append-only record}
The complete V78 source was retained before this continuation.  It
had \(925725\) bytes, \(28010\) lines, and SHA--256
\[
 \texttt{61AF8C11D278ED031AD27D3C5EF6123CE94DB302521418C91A09F3C5AC9F78CD}.
\]
V79 replaced the extensive oscillation bound by a conditional
local-to-global criterion and isolated the independent
phase-coexistence obstruction.
\section{Fourteenth-cycle V80: compulsory companion audit and the influence-matrix turn}
\label{sec:v14v80}

This is the compulsory five-version audit following V75.  The
newest active companion files in the shared folder were inspected
read-only before choosing the next SM--Einstein step.

\subsection{Frozen companion provenance}

\begin{record}[V80 companion checkpoint]
\label{rec:v14v80-companions}
The Yang--Mills checkpoint is
\begin{center}
\path{Renewal_Yang_Mills_Mass_Gap_Research_Notes_V25.tex},
\end{center}
with \(203895\) bytes, \(5554\) lines, and SHA--256
\[
 \texttt{7DA88CD644369B2F0C00E31641AA6AAC5E4E7315A493A989383EE0F323B252B0}.
                                                               \tag{80.1}
\]
The Navier--Stokes checkpoint is
\begin{center}
\path{Renewal_NS_Stability_Research_Notes_V26.tex},
\end{center}
with \(278283\) bytes, \(5319\) lines, and SHA--256
\[
 \texttt{82C887F8C2C69E4F28FAD7C3DC8DE5B9099CB5A4BF431EBA1FFF3E91935978AD}.
                                                               \tag{80.2}
\]
These hashes freeze what was read; later replacement of either
active companion does not alter this audit.
\end{record}

\subsection{Yang--Mills V25}

Yang--Mills V25 retains the full multi-affine Wilson matrix rather
than separating coefficient norms.  At one exact rational active
phase it proves, by a \(256\)-vertex hull, a four-coordinate
radius-\(1/8\) positive box with an exact rational Loewner floor.  It
also supplies an exact outer-rectangle negative witness at a larger
active radius and correctly records that this witness lies off the
physical unit circles.  Its next task is symmetry transport and
exact torus coverage.

\begin{assessment}[YM import decision]
\label{ass:v14v80-ym}
The vertex-hull principle was already imported at V75 and placed in
the finite adjacent-law certificate at V76.  V25 improves the
Yang--Mills local symbol atlas, but it supplies none of the four
objects missing from V79:
\[
 \text{a global conditional boundary law},\quad
 C_{\rm AT},\quad
 \text{an interface estimate},\quad
 \text{a cofinal physical graph-screen}.                 \tag{80.3}
\]
Indeed its own firewall lists the compact graph-screen,
norm-resolvent, tightness, and cofinal scaling rows as open.  The
exact Wilson floor may eventually populate a base fiber gap
\(\gamma_x^0(\xi)\), but only after the active finite common action,
constraint reduction, and compatibility of its conditional
marginals have been supplied.  No V25 constant is inserted into
\((79.14)\).
\end{assessment}

\subsection{Navier--Stokes V26}

Navier--Stokes V26 contracts the first two Oseen-kernel derivatives
only against rank-one inputs generated by \(W\otimes W\).  It proves
that the first-derivative norm is unchanged, whereas the
second-derivative rank-one norm is strictly smaller on an
intermediate radial range.  This is a deterministic kernel estimate,
not a boundary probability or transfer theorem.

\begin{assessment}[NS import decision]
\label{ass:v14v80-ns}
No NS theorem supplies a conditional SM law, a transfer gap, or an
approximate-tensorization constant.  One methodological lesson is
relevant: an ambient operator norm can be needlessly large when the
admissible perturbations occupy a constrained cone.  Accordingly,
the next influence calculation should compare only
constraint-compatible one-cell boundary variations.  This changes
the domain over which an influence supremum is taken; it does not
license importing the NS numerical constants.
\end{assessment}

\subsection{Audit verdict}

\begin{theorem}[No companion closes the phase row]
\label{thm:v14v80-verdict}
At the frozen checkpoints \((80.1)\)--\((80.2)\), neither companion
note implies a volume-uniform bound on
\(C_{\rm AT}(\mu_\Lambda)\), selects a unique SM--Einstein phase, or
constructs the compatible reduced conditional law required by
Theorem~\ref{thm:v14v79-local-global}.  Hence the V79
phase-coexistence row remains open.
\end{theorem}

\begin{proof}
Yang--Mills V25 proves a local finite matrix inequality on a momentum
box and explicitly leaves its global and cofinal spectral rows open.
Navier--Stokes V26 proves pointwise deterministic tensor norms.
Neither defines the conditional measures in \((79.8)\), estimates
their interdependence, or proves an interface/conductance inequality.
The types of the conclusions therefore do not match the hypotheses
of Theorem~\ref{thm:v14v79-local-global}.
\end{proof}

The audit nevertheless fixes the next upstream move.  A
Dobrushin-type matrix records how much the conditional law at \(x\)
can change when the exterior datum at \(y\) changes.  A strict
contraction of that matrix is a verifiable sufficient condition for
approximate tensorization.  It is also local enough to accept
Yang--Mills vertex certificates and constraint-adapted variation
sets when the finite common action becomes available.

\begin{decision}[V81 conditional influence compiler]
\label{dec:v14v80-next}
Define total-variation conditional influence coefficients on the
reduced admissible fibers and prove a finite Dobrushin
variance/tensorization theorem with an explicit constant.  Then:
\begin{enumerate}
\item derive computable influence bounds from finite-range action
      oscillations;
\item state the exact constraint-compatibility condition needed when
      not every single-cell replacement is admissible;
\item give a finite counterexample showing that row sums at or above
      one do not prove nonuniqueness but do make this contraction
      compiler inconclusive;
\item feed the resulting constant into \((79.14)\).
\end{enumerate}
This attacks the first unpopulated thermodynamic row rather than
adding further local symbol positivity.
\end{decision}

\begin{firewall}[Claim boundary after V80]
\label{fire:v14v80}
The companion audit is a provenance and applicability statement.
It proves no new SM transfer gap and imports no companion numerical
constant.  The full SM--Einstein continuum remains conditional on
the finite common action, constraint-compatible conditional law,
uniform influence or interface control, OS reconstruction, and the
cofinal source rows.
\end{firewall}

\section*{Version V80 append-only record}
\addcontentsline{toc}{section}{Version V80 append-only record}
The complete V79 source was retained before this continuation.  It
had \(940842\) bytes, \(28465\) lines, and SHA--256
\[
 \texttt{D2C054C5EA31E0F30D3D259BEFF15F707501B032D529E458A31B301044E63BBF}.
\]
V80 performed the compulsory companion audit and selected the
conditional influence matrix as the next upstream bottleneck.
\section{Fourteenth-cycle V81: a constraint-compatible Dobrushin compiler}
\label{sec:v14v81}

V79 isolated \(C_{\rm AT}\) as the collective row not controlled by
one-cell gaps.  V80 found no companion theorem that fills it.  This
continuation gives a finite, self-contained sufficient condition:
strict contraction of a matrix of conditional-law influences.  The
proof is spectral and therefore produces exactly the variance
inequality needed in \((79.14)\).

\subsection{Admissible coordinate geometry}

Let
\[
 \Omega\subseteq\prod_{x\in\Lambda}X_x
\]
be a finite admissible reduced boundary space, and let \(\mu\) be
strictly positive on \(\Omega\).  Two configurations are
\(x\)-adjacent when they agree off \(x\).  Write
\[
 \delta_x f
 :=
 \sup_{\substack{\eta,\zeta\in\Omega\\
                  \eta_{\ne x}=\zeta_{\ne x}}}
 |f(\eta)-f(\zeta)|,                                     \tag{81.1}
\]
with supremum zero when there is no nontrivial \(x\)-edge.

For an exterior datum \(\xi=\eta_{\ne x}\), let
\(S_x(\xi)\subseteq X_x\) be the support of the conditional law
\(\mu_x^\xi\).

\begin{definition}[Connected square-compatible reduction]
\label{def:v14v81-square}
The chosen coordinate or block reduction is called
square-compatible when:
\begin{enumerate}
\item the graph formed by all one-coordinate admissible edges is
      connected;
\item if two admissible configurations differ only at \(y\ne x\),
      their \(x\)-conditional supports coincide;
\item for every value \(s\) in that common support, inserting \(s\)
      into both exteriors gives two admissible configurations that
      still differ only at \(y\).
\end{enumerate}
\end{definition}

The condition is automatic on a full product.  It is not automatic
after Gauss, diffeomorphism, fixed-charge, or topological-sector
constraints.  In that case the coordinate labelled \(x\) must be a
constraint-preserving block move, such as a closed loop or a
balanced star, rather than a forbidden isolated-field replacement.

\subsection{Conditional influence matrix}

For \(x\ne y\), define
\[
 c_{xy}
 :=
 \sup_{\substack{\eta,\zeta\in\Omega\\
                  \eta_{\ne y}=\zeta_{\ne y}}}
 \left\|
   \mu_x^{\eta_{\ne x}}-\mu_x^{\zeta_{\ne x}}
 \right\|_{\rm TV},\qquad c_{xx}=0,                       \tag{81.2}
\]
where the two measures are identified on the common support
guaranteed by Definition~\ref{def:v14v81-square}.  Let
\[
 \alpha:=\max_{x\in\Lambda}\sum_{y\in\Lambda}c_{xy}.      \tag{81.3}
\]
Denote conditional expectation in coordinate \(x\) by \(E_x\), so
\(E_x\) is the orthogonal projection in \(L^2(\mu)\) onto functions
of \(\eta_{\ne x}\).

\begin{lemma}[One-step oscillation transport]
\label{lem:v14v81-transport}
For \(x\ne y\),
\[
 \delta_y(E_xf)
 \le\delta_yf+c_{xy}\delta_xf,                            \tag{81.4}
\]
whereas
\[
 \delta_x(E_xf)=0.                                       \tag{81.5}
\]
Consequently, for the random-scan conditional-expectation operator
\[
 P:=\frac1{|\Lambda|}\sum_{x\in\Lambda}E_x,              \tag{81.6}
\]
the seminorm
\[
 {\cal L}(f):=\sum_{x\in\Lambda}\delta_xf                \tag{81.7}
\]
obeys
\[
 {\cal L}(Pf)
 \le
 \left(1-\frac{1-\alpha}{|\Lambda|}\right){\cal L}(f).
                                                               \tag{81.8}
\]
\end{lemma}

\begin{proof}
Take two admissible exteriors differing only at \(y\).  By square
compatibility, both \(x\)-conditional laws live on the same set
\(S\), and the two functions obtained by inserting a common
\(s\in S\) differ along an admissible \(y\)-edge.  Add and subtract
the expectation of either inserted function under the other
conditional law.  The common-argument difference is at most
\(\delta_yf\).  The change of law is at most its total-variation
distance times the oscillation of \(f\) along the \(x\)-fiber, which
is at most \(\delta_xf\).  This proves \((81.4)\).
Conditional expectation depends only on the exterior, proving
\((81.5)\).

Put \(N=|\Lambda|\).  Equations \((81.4)\)--\((81.6)\) give
\[
 \delta_y(Pf)
 \le\frac{N-1}{N}\delta_yf
    +\frac1N\sum_{x\ne y}c_{xy}\delta_xf.                 \tag{81.9}
\]
Sum over \(y\) and interchange the second pair of sums:
\[
 {\cal L}(Pf)
 \le\frac{N-1}{N}{\cal L}(f)
   +\frac1N\sum_x\delta_xf\sum_yc_{xy}
 \le\left(1-\frac{1-\alpha}{N}\right){\cal L}(f).
\]
\end{proof}

\begin{theorem}[Dobrushin approximate tensorization]
\label{thm:v14v81-dobrushin}
Assume Definition~\ref{def:v14v81-square} and
\[
 \alpha<1.                                               \tag{81.10}
\]
Then the sum-clock heat-bath generator
\[
 A_{\rm HB}:=\sum_{x\in\Lambda}(I-E_x)                  \tag{81.11}
\]
has
\[
 \operatorname{gap}(A_{\rm HB})\ge1-\alpha.              \tag{81.12}
\]
Equivalently,
\[
 \operatorname{Var}_\mu(f)
 \le\frac1{1-\alpha}
 \sum_{x\in\Lambda}
 \mathbb E_\mu
 \left[
  \operatorname{Var}_{\mu_x^{\eta_{\ne x}}}
   \bigl(f(\,\cdot\,,\eta_{\ne x})\bigr)
 \right],                                                \tag{81.13}
\]
so
\[
 C_{\rm AT}(\mu)\le\frac1{1-\alpha}.                     \tag{81.14}
\]
\end{theorem}

\begin{proof}
Each \(E_x\) is an orthogonal projection, hence \(P\) in \((81.6)\)
is self-adjoint, positive semidefinite, and has spectrum in
\([0,1]\).  The connectedness in
Definition~\ref{def:v14v81-square} implies that a function fixed by
every \(E_x\) is constant.  Thus the eigenspace of \(P\) at one is
the constants.

Let \(Pf=\lambda f\) for a nonconstant eigenfunction.  Then
\({\cal L}(f)>0\), and positivity gives \(\lambda\ge0\).  By
\((81.8)\),
\[
 \lambda{\cal L}(f)={\cal L}(Pf)
 \le\left(1-\frac{1-\alpha}{N}\right){\cal L}(f),
\]
so every nonconstant eigenvalue satisfies
\[
 \lambda\le1-\frac{1-\alpha}{N}.                         \tag{81.15}
\]
Since
\[
 A_{\rm HB}=N(I-P),
\]
its first positive eigenvalue is at least \(1-\alpha\), proving
\((81.12)\).

Finally,
\[
 \langle f,(I-E_x)f\rangle_\mu
 =\|(I-E_x)f\|_{L^2(\mu)}^2
 =\mathbb E_\mu[
   \operatorname{Var}_{\mu_x^{\eta_{\ne x}}}(f)],
\]
because \(E_x\) is an orthogonal projection.  The Poincare
inequality equivalent to \((81.12)\) is therefore \((81.13)\).
\end{proof}

\begin{corollary}[Dobrushin completion of the V79 gap row]
\label{cor:v14v81-gap}
Under the hypotheses of
Proposition~\ref{prop:v14v79-conditional} and
Theorem~\ref{thm:v14v81-dobrushin},
\[
 \operatorname{gap}(A_\mu)
 \ge
 (1-\alpha)e^{-2\Theta_*}\gamma^0_{\rm fib}.             \tag{81.16}
\]
A cofinal volume-uniform gap follows if the three displayed factors
have positive uniform lower bounds.
\end{corollary}

\begin{proof}
Insert \((81.14)\) into \((79.14)\).
\end{proof}

\subsection{Compiling influences from the action}

The influence coefficients become finite arithmetic once the
conditional action is known.  The next elementary estimate records
the sharp dependence on a likelihood-ratio oscillation.

\begin{lemma}[Likelihood-ratio oscillation bound]
\label{lem:v14v81-likelihood}
Let \(p,q\) be strictly positive probability laws on the same finite
set and put
\[
 D=\operatorname{osc}\log\frac pq.
\]
Then
\[
 \|p-q\|_{\rm TV}\le\tanh\frac D4.                       \tag{81.17}
\]
\end{lemma}

\begin{proof}
Let \(r=p/q\), \(m=\min r\), and \(M=\max r\).  Then
\(m\le1\le M\), \(M/m\le e^D\), and
\[
 \|p-q\|_{\rm TV}=\frac12\mathbb E_q|r-1|.
\]
Among random variables in \([m,M]\) with mean one, convexity of
\(|r-1|\) makes the endpoint law maximal.  Hence
\[
 \|p-q\|_{\rm TV}
 \le\frac{(1-m)(M-1)}{M-m}.                              \tag{81.18}
\]
For fixed \(R=M/m\), the right-hand side is maximized at
\(m=R^{-1/2}\), where it equals
\[
 \frac{\sqrt R-1}{\sqrt R+1}.
\]
This expression increases with \(R\).  Use \(R\le e^D\) to obtain
\((81.17)\).
\end{proof}

Suppose a common positive reference law \(\nu_x\) exists on each
square-compatible fiber and
\[
 \mu_x^\xi(s)
 =
 \frac{\nu_x(s)e^{-U_x(s;\xi)}}
      {\sum_t\nu_x(t)e^{-U_x(t;\xi)}}.                   \tag{81.19}
\]
Here \(U_x\) is dimensionless; in the V78 convention it contains
\(aV_x\).  Define the mixed action oscillation
\[
 \Delta_{xy}
 :=
 \sup_{\substack{\eta,\zeta\in\Omega\\
                  \eta_{\ne y}=\zeta_{\ne y}}}
 \operatorname{osc}_{s\in S_x}
 \left[
  U_x(s;\eta_{\ne x})-U_x(s;\zeta_{\ne x})
 \right].                                                \tag{81.20}
\]

\begin{proposition}[Action-to-influence compiler]
\label{prop:v14v81-action}
For the conditional laws \((81.19)\),
\[
 c_{xy}\le\tanh\frac{\Delta_{xy}}4.                      \tag{81.21}
\]
Consequently the explicit condition
\[
 \max_x\sum_{y\ne x}\tanh\frac{\Delta_{xy}}4<1           \tag{81.22}
\]
implies \((81.13)\)--\((81.16)\).

If
\[
 U(\eta)=\sum_{B\subseteq\Lambda}\Phi_B(\eta_B),
\]
then only interactions containing both \(x\) and \(y\) enter:
\[
 \Delta_{xy}
 \le
 \sum_{B\ni x,y}\Delta_{xy}(\Phi_B),                     \tag{81.23}
\]
where
\[
 \Delta_{xy}(\Phi_B)
 :=
 \sup_{\substack{s,t\\ \eta,\zeta\ {\rm differ\ at}\ y}}
 \left|
  [\Phi_B(s,\eta)-\Phi_B(t,\eta)]
  -
  [\Phi_B(s,\zeta)-\Phi_B(t,\zeta)]
 \right|.                                                \tag{81.24}
\]
In particular,
\[
 \Delta_{xy}(\Phi_B)\le2\operatorname{osc}_x\Phi_B.       \tag{81.25}
\]
\end{proposition}

\begin{proof}
For two exteriors differing at \(y\), the normalizing constants in
\((81.19)\) add an \(s\)-independent constant to the logarithm of
the likelihood ratio.  Its oscillation is therefore exactly the
oscillation in \((81.20)\).  Apply
Lemma~\ref{lem:v14v81-likelihood} and take the supremum to get
\((81.21)\).  Equation \((81.22)\) gives \(\alpha<1\).

Terms not containing \(x\) cancel from a conditional law.  Among the
remaining terms, those not containing \(y\) cancel between the two
exteriors.  The triangle inequality gives \((81.23)\).  Each of the
two bracketed \(x\)-differences in \((81.24)\) has magnitude at most
\(\operatorname{osc}_x\Phi_B\), proving \((81.25)\).
\end{proof}

\begin{assessment}[High-temperature character]
\label{ass:v14v81-high-temperature}
Finite interaction range makes the matrix in \((81.22)\) sparse, but
does not make its row sums smaller than one.  Condition \((81.22)\)
is a quantitative weak-dependence or single-phase certificate.  It
is not a consequence of locality, and its failure is not a proof of
phase coexistence.
\end{assessment}

\subsection{Failure of the compiler is not failure of the theory}

Consider the ferromagnetic Ising law on the complete graph \(K_4\):
\[
 \mu_\beta(\sigma)
 \propto
 \exp\!\left(\beta\sum_{1\le i<j\le4}\sigma_i\sigma_j\right),
 \qquad \sigma_i\in\{-1,+1\}.                            \tag{81.26}
\]

\begin{proposition}[A finite positive-gap model outside the
Dobrushin region]
\label{prop:v14v81-k4}
For every \(i\ne j\),
\[
 c_{ij}=\tanh\beta,\qquad
 \alpha=3\tanh\beta.                                     \tag{81.27}
\]
Thus \(\alpha\ge1\) once
\(\beta\ge\operatorname{arctanh}(1/3)\).  Nevertheless, for every
finite \(\beta\), the four-spin heat-bath chain is irreducible and
has a strictly positive spectral gap.
\end{proposition}

\begin{proof}
Condition on the three spins other than \(i\).  If their sum is
\(H\), then
\[
 \mu_\beta(\sigma_i=+1\mid\sigma_{\ne i})
 =\frac{1+\tanh(\beta H)}2.                              \tag{81.28}
\]
When spin \(j\) changes, the other two exterior spins can be chosen
with opposite signs, so the two fields are \(-1\) and \(+1\).
The total-variation distance between the resulting Bernoulli laws is
\(\tanh\beta\).  For any other background, monotonicity and the
decrease of the derivative of \(\tanh\) away from zero show that the
distance is no larger.  This proves \((81.27)\).

All conditional probabilities in \((81.28)\) are strictly between
zero and one at finite \(\beta\).  Single-spin heat-bath moves
therefore connect the full hypercube, and the finite reversible
generator has constants as its only zero modes.  Its first positive
eigenvalue is consequently strictly positive.
\end{proof}

The proposition prevents a false contraposition: \(\alpha\ge1\)
means only that Theorem~\ref{thm:v14v81-dobrushin} is silent.
Conversely, the expander family in V79 shows that uniform local
conditional gaps can coexist with a collapsing global gap.  Beyond
the Dobrushin region one needs a sharper block influence, a
phase-restricted tensorization theorem, or an interface/conductance
estimate.

\subsection{Constraint and source acquisition}

\begin{acquisition}[Conditional influence card]
\label{acq:v14v81}
For every finite SM--Einstein regulator, the active source must now
provide:
\begin{enumerate}
\item constraint-preserving coordinate blocks whose move graph is
      connected inside the selected physical sector;
\item the square-compatibility check of
      Definition~\ref{def:v14v81-square}, or a block replacement
      for it;
\item the real positive conditional law \((81.19)\);
\item the exact mixed action table \(\Delta_{xy}\), preferably before
      the crude loss \((81.25)\);
\item a certified row-sum bound \(\alpha<1\), or an identified
      alternative phase/interface theorem;
\item the fiber kinetic gap and conditional potential oscillation
      needed by \((81.16)\);
\item the absolute normalization and its first two metric jets,
      which remain separate Einstein-source data.
\end{enumerate}
A complex fermion determinant or an unsigned Grassmann expression
does not populate item 3; its positive OS realization must be
constructed first.
\end{acquisition}

\begin{assessment}[Current source boundary]
\label{ass:v14v81-source}
The attached sources do not enumerate the reduced admissible
boundary blocks, their conditional supports, or the mixed action
oscillations \((81.20)\).  Therefore \(\alpha\) cannot yet be
evaluated.  The theorem identifies a finite calculation to perform
once the common regulator and physical conditional law are fixed,
but it does not infer weak dependence from the formal locality of
the Lagrangian.
\end{assessment}

\begin{decision}[V82 block influences and sector connectivity]
\label{dec:v14v81-next}
Go upstream to the constraint geometry.  Prove a block version of
Theorem~\ref{thm:v14v81-dobrushin} for overlapping
constraint-preserving updates, identify the additional coverage
constant caused by overlap, and separate three obstructions:
disconnected topological sectors, changing conditional supports,
and an influence norm at least one.  This is the form needed for
Gauss-law-compatible gauge and gravity moves.
\end{decision}

\begin{firewall}[Claim boundary after V81]
\label{fire:v14v81}
The Dobrushin compiler is a finite sufficient theorem.  Its
hypotheses have not been verified for a common regulated
SM--Einstein measure.  It establishes neither phase uniqueness nor a
continuum mass gap when \(\alpha\ge1\), and it does not handle a
complex fermionic weight.  No full SM--Einstein continuum is claimed.
\end{firewall}

\section*{Version V81 append-only record}
\addcontentsline{toc}{section}{Version V81 append-only record}
The complete V80 source was retained before this continuation.  It
had \(947436\) bytes, \(28619\) lines, and SHA--256
\[
 \texttt{0084155F4643167DA67388F2859D018C61D1D65F93799948C6BE1B6E25F92BBC}.
\]
V81 proved an explicit Dobrushin approximate-tensorization theorem,
compiled its influence matrix from mixed local action oscillations,
and isolated the constraint-square requirement.
\section{Fourteenth-cycle V82: block influences, coverage, and sector kernels}
\label{sec:v14v82}

V81's proof uses single-coordinate conditional expectations.  Gauge
and gravitational constraints often make such a replacement
inadmissible: a link or flux cannot change without a compensating
move.  This continuation passes to constraint-preserving blocks and
keeps three obstructions separate:
\[
 \text{sector disconnection},\qquad
 \text{support incompatibility},\qquad
 \text{large conditional influence}.                    \tag{82.1}
\]

\subsection{Exact kernel of a block-move form}

Let \(\Omega\) be finite, let \(\mu(\eta)>0\) for every
\(\eta\in\Omega\), and let \({\cal B}\) be a family of allowed
blocks.  For each \(b\in{\cal B}\), suppose the conditional
\(b\)-fiber carries a reversible form with strictly positive
conductance on every declared \(b\)-move edge.  Let
\({\cal E}=\sum_b{\cal E}_b\), and let \({\mathfrak G}_{\cal B}\) be
the undirected graph obtained by taking the union of those edges.

\begin{proposition}[Block kernel equals the sector algebra]
\label{prop:v14v82-kernel}
A function \(f\) satisfies \({\cal E}(f)=0\) if and only if it is
constant on every connected component of
\({\mathfrak G}_{\cal B}\).  Hence
\[
 \dim\ker A_{\cal B}
 =
 \#\{\text{block-move components}\}.                     \tag{82.2}
\]
If there is more than one component, the Poincare gap relative to
the single global constant subspace is zero.
\end{proposition}

\begin{proof}
Every reversible finite form is a sum of nonnegative terms
\(c(\eta,\zeta)|f(\eta)-f(\zeta)|^2\) over its move edges, with
\(c(\eta,\zeta)>0\).  The total sum vanishes exactly when every edge
difference vanishes.  This is equivalent to constancy on each
connected component.  If there are at least two components, a
nonconstant function that is constant on each component can be
chosen with \(\mu\)-mean zero; it has zero energy and positive
variance.
\end{proof}

Thus a topological or charge sector must be selected before a
unique-vacuum gap theorem is possible, unless the common action
contains an explicitly controlled sector-changing move.

\subsection{Product charts inside one sector}

Fix one connected component \(\Omega_s\).  Suppose a
constraint-solving chart gives a bijection
\[
 \Psi:\prod_{i\in I}Z_i\longrightarrow\Omega_s,           \tag{82.3}
\]
and pull \(\mu\) back to a strictly positive law \(\widehat\mu\) on
the full chart product.  The chart coordinates need not be physical
single fields; they may be loop variables, balanced fluxes, or
gauge-fixed representatives.

For \(S\subseteq I\), write \(E_S\) for conditional expectation given
the coordinates outside \(S\), and set
\[
 D_S(f):=\langle f,(I-E_S)f\rangle_{\widehat\mu}
       =\mathbb E_{\widehat\mu}
         [\operatorname{Var}(f\mid z_{I\setminus S})].    \tag{82.4}
\]
If \(i\in S\), then
\[
 D_i(f)\le D_S(f),                                       \tag{82.5}
\]
because the subspace of functions measurable outside \(S\) is
contained in the subspace of functions measurable outside \(i\);
orthogonal distance to the smaller subspace is larger.

Let a physical update \(b\) vary a chart-coordinate set
\(S_b\subseteq I\).  Assume its conditional form satisfies the
uniform fiber inequality
\[
 {\cal E}_b(f)\ge g_bD_{S_b}(f),\qquad g_b>0.             \tag{82.6}
\]
Weights or physical clock rates are included in \(g_b\).

\begin{definition}[Fractional block coverage]
\label{def:v14v82-coverage}
A fractional allocation is a family
\(\theta_{bi}\ge0\), defined for \(i\in S_b\), such that
\[
 \sum_{i\in S_b}\theta_{bi}\le g_b.                      \tag{82.7}
\]
Its coverage is
\[
 \kappa(\theta)
 :=
 \min_{i\in I}\sum_{b:\,i\in S_b}\theta_{bi},             \tag{82.8}
\]
and the optimal coverage constant is
\[
 \kappa_*:=\sup_\theta\kappa(\theta).                    \tag{82.9}
\]
This is a finite linear programme.
\end{definition}

\begin{theorem}[Overlapping-block gap compiler]
\label{thm:v14v82-block}
Let \(\widehat\alpha\) be the V81 Dobrushin row-sum coefficient of
the pulled-back chart law \(\widehat\mu\).  If
\[
 \widehat\alpha<1,\qquad \kappa_*>0,                      \tag{82.10}
\]
then the physical block generator satisfies
\[
 \operatorname{gap}_{\Omega_s}(A_{\cal B})
 \ge\kappa_*(1-\widehat\alpha).                          \tag{82.11}
\]
The explicit equal-allocation lower bound is
\[
 \kappa_*
 \ge
 \min_{i\in I}
 \sum_{b:\,i\in S_b}\frac{g_b}{|S_b|}.                   \tag{82.12}
\]
If the \(b\)-fiber itself is a V78 tilted transfer with base gap
\(g_b^0\) and conditional oscillation \(\Theta_b\), one may take
\[
 g_b=e^{-2\Theta_b}g_b^0.                                \tag{82.13}
\]
\end{theorem}

\begin{proof}
By \((82.5)\)--\((82.7)\),
\[
 {\cal E}_b(f)
 \ge g_bD_{S_b}(f)
 \ge\sum_{i\in S_b}\theta_{bi}D_i(f).
\]
Summing over blocks gives
\[
 {\cal E}(f)
 \ge\sum_i
 \left(\sum_{b:\,i\in S_b}\theta_{bi}\right)D_i(f)
 \ge\kappa(\theta)\sum_iD_i(f).                          \tag{82.14}
\]
The chart is a full product, so Theorem
\ref{thm:v14v81-dobrushin} applies and yields
\[
 \sum_iD_i(f)\ge(1-\widehat\alpha)
 \operatorname{Var}_{\widehat\mu}(f).
\]
Optimize over \(\theta\) to prove \((82.11)\).  Choosing
\(\theta_{bi}=g_b/|S_b|\) proves \((82.12)\).
Equation \((82.13)\) is the conditional V78 comparison.

The chart bijection is unitary between the two finite \(L^2\)
spaces, so the same gap holds on \(\Omega_s\).
\end{proof}

\begin{assessment}[Meaning of the coverage loss]
\label{ass:v14v82-coverage}
The factor \(\kappa_*\) is not an interaction or phase constant.  It
measures how the declared physical update clocks cover the
independent constraint-solved directions.  It can vanish because a
chart direction is never updated, even when every declared block has
a positive internal gap.  Overlap itself is harmless; uncovered or
arbitrarily slow directions are not.
\end{assessment}

\subsection{Parity constraint: exact failure and repair}

Let
\[
 \Omega_{\rm even}
 :=
 \left\{\eta\in\{0,1\}^N:
        \sum_{j=1}^N\eta_j=0\pmod2\right\}               \tag{82.15}
\]
with the uniform law.

\begin{countertheorem}[Forbidden single-site updates fake zero
influence]
\label{ctr:v14v82-parity}
Every one-site conditional fiber on \(\Omega_{\rm even}\) is a
singleton.  Hence
\[
 E_j=I,\qquad c_{ij}=0,\qquad\alpha=0,                   \tag{82.16}
\]
but the one-site heat-bath generator is identically zero.  Thus
discarding the connectedness and square-compatibility hypotheses of
V81 would turn \(\alpha=0\) into a false positive gap certificate.
\end{countertheorem}

\begin{proof}
Once the other \(N-1\) bits are fixed, parity determines the last
bit.  Therefore conditional resampling changes nothing.  All
one-site forms vanish and their move graph has \(2^{N-1}\) isolated
vertices.
\end{proof}

Now take the pair blocks
\[
 b_j=\{1,j\},\qquad j=2,\ldots,N.                         \tag{82.17}
\]
Conditioned outside \(b_j\), exactly two even-parity states remain,
and the block heat bath resamples them uniformly.

\begin{proposition}[Constraint-solving chart restores the exact gap]
\label{prop:v14v82-parity-repair}
Use chart coordinates
\[
 z_j=\eta_j\quad(2\le j\le N),\qquad
 \eta_1=\sum_{j=2}^Nz_j\pmod2.                           \tag{82.18}
\]
Then the uniform law becomes the product uniform law on
\(\{0,1\}^{N-1}\), the block \(b_j\) updates exactly \(z_j\), and
\[
 A_{\rm pair}=\sum_{j=2}^N(I-E_{b_j})
\]
has
\[
 \operatorname{gap}(A_{\rm pair})=1.                    \tag{82.19}
\]
\end{proposition}

\begin{proof}
Equation \((82.18)\) is a bijection.  Each choice of
\((z_2,\ldots,z_N)\) has the same mass, so the pullback is product
uniform.  Flipping the two bits \(\eta_1,\eta_j\) changes only \(z_j\)
and preserves parity.  The generator is therefore the sum of
\(N-1\) independent two-point heat-bath projections.  Theorem
\ref{thm:v14v79-product} gives the exact gap one.  Equivalently,
Theorem~\ref{thm:v14v82-block} has
\(\widehat\alpha=0\), \(g_{b_j}=1\), and
\(\kappa_*=1\).
\end{proof}

If the same pair blocks act on the full hypercube, parity is
preserved and there are exactly two move components.  Each component
has the sector gap \((82.19)\), while the global Poincare gap relative
to one constant is zero by Proposition~\ref{prop:v14v82-kernel}.
This is the finite prototype of a topological-sector zero mode.

\subsection{Support changes and chart ancestry}

A constraint-solving map can turn a local physical action into a
long-range chart action.  Therefore the influence matrix must be
computed after applying \(\Psi\), not read from the unreduced
Lagrangian.

\begin{proposition}[Ancestry of mixed chart interactions]
\label{prop:v14v82-ancestry}
Let
\[
 U(\eta)=\sum_B\Phi_B(\eta_B),\qquad \eta=\Psi(z).
\]
For chart coordinates \(i\ne j\), the mixed action oscillation
\(\widehat\Delta_{ij}\) of V81 receives a contribution from
\(\Phi_B\) only if the composed observable
\(\Phi_B\circ\Psi\) depends on both \(z_i\) and \(z_j\).
Consequently an unreduced finite-range bound implies a sparse chart
influence matrix only when the constraint-solving ancestry sets
\[
 {\cal A}(B)
 :=
 \{i:\Phi_B\circ\Psi\text{ depends on }z_i\}             \tag{82.20}
\]
have uniformly bounded size and incidence.
\end{proposition}

\begin{proof}
In the mixed difference \((81.24)\), a term independent of either
one of the varied chart coordinates cancels.  Thus only terms with
\(\{i,j\}\subseteq{\cal A}(B)\) survive.  Uniform bounds on the
number and size of the surviving ancestry sets give the asserted
sparsity; without them, solving a global constraint can make a local
term depend on arbitrarily many chart directions.
\end{proof}

When conditional supports change with the exterior, a single product
chart may not exist.  One must then either:
\begin{enumerate}
\item enlarge a block until its conditional support is stable;
\item split the state space into support-stable charts and prove
      quantitative overlap transport;
\item construct a direct conductance comparison without the V81
      common-support step.
\end{enumerate}
Assigning zero influence to a missing move repeats the error in
Countertheorem~\ref{ctr:v14v82-parity}.

\subsection{Sector-resolved acquisition ledger}

\begin{acquisition}[Block and sector card]
\label{acq:v14v82}
For every finite regulator, record:
\begin{enumerate}
\item the physical block-move graph and its connected components;
\item the selected charge, topology, and boundary-condition sector;
\item a constraint-solving product chart \(\Psi\), or an explicit
      support-stable atlas;
\item each physical update's chart support \(S_b\) and conditional
      fiber gap \(g_b\);
\item the fractional coverage linear programme \((82.7)\)--\((82.9)\);
\item the post-reduction ancestry sets \({\cal A}(B)\);
\item the chart mixed-action matrix
      \(\widehat\Delta_{ij}\) and influence row sum
      \(\widehat\alpha\);
\item the absolute sector normalization and its first two metric
      jets for the Einstein-source comparison.
\end{enumerate}
The finite sector certificate is \((82.11)\).  A global unique-vacuum
statement also needs a comparison of sector ground energies or a
controlled sector-changing mechanism.
\end{acquisition}

\begin{assessment}[Current source boundary]
\label{ass:v14v82-source}
The attached sources contain constraint and sector terminology but
do not enumerate a common finite block-move graph, a product chart,
or the ancestry and coverage tables above.  In particular, formal
locality before solving Gauss and gravitational constraints does not
populate the post-reduction influence row.
\end{assessment}

\begin{decision}[V83 sector transfer and vacuum competition]
\label{dec:v14v82-next}
Treat the direct sum over disconnected physical sectors.  Prove the
exact decomposition of the positive transfer and its logarithmic
Hamiltonian, identify the global vacuum multiplicity from the sector
ground energies, and quantify the effect of a small
sector-changing transfer block by a Schur/Feshbach estimate.  This
will decide whether sector restriction, energy separation, or
controlled tunnelling is required before the V72 Einstein vacuum
source can be called unique.
\end{decision}

\begin{firewall}[Claim boundary after V82]
\label{fire:v14v82}
The block theorem is finite and sector-resolved.  Its product chart,
fiber gaps, coverage, and influence hypotheses remain unverified for
a common SM--Einstein regulator.  It does not remove topological
superselection, construct fermionic positivity, compare sector
vacuum energies, or establish a continuum mass gap.
\end{firewall}

\section*{Version V82 append-only record}
\addcontentsline{toc}{section}{Version V82 append-only record}
The complete V81 source was retained before this continuation.  It
had \(962802\) bytes, \(29080\) lines, and SHA--256
\[
 \texttt{B00971BBD3090D7EE1F5A660A7254FB680A6F1E6792F882E89F692309766A7E5}.
\]
V82 proved the overlapping-block coverage theorem, exposed the
exact sector kernel, and repaired a parity constraint by
constraint-preserving pair moves.
\section{Fourteenth-cycle V83: sector transfer, tunnelling, and vacuum-source regularity}
\label{sec:v14v83}

V82 proved that disconnected block-move components produce exact
zero modes before any continuum limit.  A sector restriction can
restore a gap, but the full SM--Einstein vacuum still requires a
decision among sectors.  This continuation keeps the unnormalized
transfer eigenvalues, because separate Markov normalization would
erase precisely the energy offsets that choose the vacuum.

\subsection{Exact direct-sum spectrum}

Let
\[
 {\cal H}=\bigoplus_{s\in{\cal S}}{\cal H}_s,\qquad
 {\mathbb T}_0=\bigoplus_{s\in{\cal S}}{\mathbb T}_s,     \tag{83.1}
\]
where \({\cal S}\) is finite and each
\({\mathbb T}_s\) is a strictly positive self-adjoint finite
transfer.  List its eigenvalues as
\[
 \lambda_{s,0}\ge\lambda_{s,1}\ge\cdots>0,
\]
and define
\[
 E_s=-\frac1a\log\lambda_{s,0},\qquad
 \Delta_s=\frac1a\log\frac{\lambda_{s,0}}{\lambda_{s,1}},
                                                               \tag{83.2}
\]
with \(\Delta_s=+\infty\) for a one-dimensional sector.

\begin{theorem}[Sector vacuum and global gap ledger]
\label{thm:v14v83-direct}
Let
\[
 E_*=\min_sE_s,\qquad
 {\cal S}_*=\{s:E_s=E_*\}.                               \tag{83.3}
\]
Then
\[
 H_0:=-\frac1a\log{\mathbb T}_0
     =\bigoplus_s\left(-\frac1a\log{\mathbb T}_s\right). \tag{83.4}
\]
The global ground multiplicity is the sum of the sector ground
multiplicities over \(s\in{\cal S}_*\).  If every winning sector has
a simple ground state, the spectral gap above the complete global
ground space is
\[
 \Delta_{\rm glob}
 =
 \min\left\{
       \min_{s\in{\cal S}_*}\Delta_s,\,
       \min_{r\notin{\cal S}_*}(E_r-E_*)
     \right\}.                                           \tag{83.5}
\]
The global vacuum is unique exactly when one sector has a strictly
lowest energy and its sector vacuum is unique.
\end{theorem}

\begin{proof}
The functional calculus of a finite direct sum is the direct sum of
the functional calculi, proving \((83.4)\).  The eigenvalues of
\(H_0\) are
\(-a^{-1}\log\lambda_{s,j}\).  Their minimum is \((83.3)\);
counting its occurrences gives the multiplicity.  The first level
above the ground space is either an internal excitation in a winning
sector or the ground level of a losing sector.  Taking the smaller
of these two possibilities gives \((83.5)\).
\end{proof}

\begin{corollary}[Separate stochastic normalization loses sector
selection]
\label{cor:v14v83-normalization}
Replacing every sector transfer by
\(\lambda_{s,0}^{-1}{\mathbb T}_s\) preserves \(\Delta_s\) but sends
every \(E_s\) to zero.  It therefore destroys the second term in
\((83.5)\), the identity of the winning sector, and the absolute
vacuum energy used by the Einstein source row.
\end{corollary}

\begin{proof}
The replacement subtracts \(E_sI\) from the sector Hamiltonian.
Every within-sector energy difference is unchanged, while all sector
ground levels become equal.
\end{proof}

\subsection{A small sector-changing transfer}

Let \({\mathbb T}_0\) now have a simple largest eigenvalue
\(\lambda_0\), normalized eigenvector \(\psi_0\), and second
eigenvalue \(\lambda_1\).  Put
\[
 P=|\psi_0\rangle\langle\psi_0|,\quad Q=I-P,\quad
 \delta_T=\lambda_0-\lambda_1>0.                         \tag{83.6}
\]
Let
\[
 {\mathbb T}={\mathbb T}_0+K,\qquad
 K=K^*,\qquad \varepsilon=\|K\|,                         \tag{83.7}
\]
and assume \({\mathbb T}>0\) and
\[
 2\varepsilon<\delta_T.                                  \tag{83.8}
\]

\begin{theorem}[Stable sector vacuum and logarithmic gap]
\label{thm:v14v83-perturb}
The largest eigenvalue \(\widetilde\lambda_0\) of \({\mathbb T}\) is
simple, and its next eigenvalue \(\widetilde\lambda_1\) satisfies
\[
 \widetilde\lambda_0\ge\lambda_0-\varepsilon,\qquad
 \widetilde\lambda_1\le\lambda_1+\varepsilon.             \tag{83.9}
\]
Consequently the perturbed Hamiltonian
\(H=-a^{-1}\log{\mathbb T}\) has
\[
 \operatorname{gap}(H)
 \ge
 \frac1a\log
 \frac{\lambda_0-\varepsilon}{\lambda_1+\varepsilon}
 >0.                                                      \tag{83.10}
\]
If \(\widetilde\psi_0\) is its normalized transfer vacuum and
\[
 \varepsilon_\perp:=\|QKP\|,
\]
then
\[
 \|Q\widetilde\psi_0\|
 \le
 \frac{\varepsilon_\perp}{\delta_T-2\varepsilon}
 \|P\widetilde\psi_0\|.                                  \tag{83.11}
\]
\end{theorem}

\begin{proof}
The min--max principle gives the two Weyl bounds in \((83.9)\).
Their difference is at least
\(\delta_T-2\varepsilon>0\), so the leading eigenvalue is simple.
Since \(-\log\) reverses the order of positive transfer eigenvalues,
their logarithmic ratio gives \((83.10)\).

Project the eigenvalue equation onto \(Q{\cal H}\):
\[
 (\widetilde\lambda_0-Q{\mathbb T}Q)
 Q\widetilde\psi_0=QKP\widetilde\psi_0.                  \tag{83.12}
\]
On \(Q{\cal H}\),
\[
 Q{\mathbb T}Q\preceq(\lambda_1+\varepsilon)Q,
\]
whereas
\(\widetilde\lambda_0\ge\lambda_0-\varepsilon\).  Thus the
inverse in \((83.12)\) has norm at most
\((\delta_T-2\varepsilon)^{-1}\), proving \((83.11)\).
\end{proof}

The same projection gives the promised Feshbach identity.

\begin{proposition}[Vacuum Feshbach correction]
\label{prop:v14v83-feshbach}
Let \(k_{00}=\langle\psi_0,K\psi_0\rangle\).  Under \((83.8)\),
\[
 \widetilde\lambda_0
 =
 \lambda_0+k_{00}
 +
 \left\langle\psi_0,
 KQ(\widetilde\lambda_0-Q{\mathbb T}Q)^{-1}QK\psi_0
 \right\rangle,                                         \tag{83.13}
\]
and therefore
\[
 0\le
 \widetilde\lambda_0-\lambda_0-k_{00}
 \le
 \frac{\varepsilon_\perp^2}{\delta_T-2\varepsilon}.      \tag{83.14}
\]
\end{proposition}

\begin{proof}
The leading eigenvector has nonzero \(P\)-component; otherwise
\((83.9)\) would contradict the upper bound on
\(Q{\mathbb T}Q\).  Solve \((83.12)\) for its \(Q\)-component and
insert the result into the \(P\)-projected eigenvalue equation.
This gives \((83.13)\).  The resolvent is positive because
\(\widetilde\lambda_0\) lies above the spectrum of
\(Q{\mathbb T}Q\), and its norm has the bound used in
\((83.11)\), proving \((83.14)\).
\end{proof}

\begin{assessment}[What must be uniform]
\label{ass:v14v83-uniform}
Finite-dimensional simplicity is insufficient.  Along a cofinal
family, the usable inputs are a lower bound on the relative
isolation
\[
 \frac{\delta_{T,n}}{a_n\lambda_{0,n}},
\]
a smallness bound on
\(\varepsilon_n/\delta_{T,n}\), and compact control of the perturbed
vacuum.  If the unperturbed sector separation collapses, the
denominators in \((83.11)\)--\((83.14)\) expose the failure rather
than hiding it.
\end{assessment}

\subsection{Exact two-sector tunnelling scale}

Take two one-dimensional sector vacua with equal unperturbed
transfer value \(\lambda>0\), and let the only sector-changing block
be
\[
 {\mathbb T}_{\rm eff}
 =
 \begin{pmatrix}
 \lambda&\kappa\\
 \overline\kappa&\lambda
 \end{pmatrix},
 \qquad |\kappa|<\lambda.                                \tag{83.15}
\]

\begin{proposition}[Finite tunnelling and cofinal softening]
\label{prop:v14v83-tunnel}
The vacuum is simple for every \(\kappa\ne0\), and the two-level
Hamiltonian gap is
\[
 \Delta_{\rm tun}
 =
 \frac1a\log\frac{\lambda+|\kappa|}
                         {\lambda-|\kappa|}.             \tag{83.16}
\]
If \(|\kappa|/\lambda\to0\), then
\[
 \Delta_{\rm tun}
 =
 \frac{2|\kappa|}{a\lambda}
 +O\!\left(\frac{|\kappa|^3}{a\lambda^3}\right).         \tag{83.17}
\]
Thus finite-volume vacuum uniqueness does not give a uniform
continuum gap when \(|\kappa|/(a\lambda)\to0\).
\end{proposition}

\begin{proof}
The transfer eigenvalues are \(\lambda\pm|\kappa|\).
Apply \(-a^{-1}\log\) and subtract.  Expansion of
\(\log[(1+x)/(1-x)]\) at \(x=0\) gives \((83.17)\).
\end{proof}

An interface-suppressed tunnelling amplitude can therefore reproduce
the V79 phase soft mode even though every finite transfer has a
unique Perron vector.  A statement of uniqueness must declare
whether it is sector restriction, absolute energy separation, or a
uniform tunnelling split that survives the limit.

\subsection{Metric jets and sector crossings}

Let \(g\) denote a finite collection of metric parameters, and
suppose \({\mathbb T}(g)\) is \(C^2\) with a simple largest
eigenvalue \(\lambda(g)\) separated from the rest.  Define
\[
 E_{\rm vac}(g)=-a^{-1}\log\lambda(g).                   \tag{83.18}
\]

\begin{proposition}[Transfer-to-Einstein vacuum jets]
\label{prop:v14v83-jets}
For a metric variation \(u\),
\[
 D\lambda[u]
 =
 \langle\psi,D{\mathbb T}[u]\psi\rangle,                 \tag{83.19}
\]
and for variations \(u,v\),
\[
 \begin{split}
 D^2\lambda[u,v]
 ={}&
 \langle\psi,D^2{\mathbb T}[u,v]\psi\rangle\\
 &+
 \left\langle QD{\mathbb T}[u]\psi,
  (\lambda-Q{\mathbb T}Q)^{-1}
  QD{\mathbb T}[v]\psi\right\rangle\\
 &+
 \left\langle QD{\mathbb T}[v]\psi,
  (\lambda-Q{\mathbb T}Q)^{-1}
  QD{\mathbb T}[u]\psi\right\rangle .                   \tag{83.20}
 \end{split}
\]
Consequently
\[
 DE_{\rm vac}
 =-\frac1a\frac{D\lambda}{\lambda},\qquad
 D^2E_{\rm vac}
 =-\frac1a
 \left(
  \frac{D^2\lambda}{\lambda}
  -\frac{D\lambda\otimes D\lambda}{\lambda^2}
 \right).                                                \tag{83.21}
\]
\end{proposition}

\begin{proof}
Differentiate the normalized eigenvalue equation and impose
\(\langle\psi,D\psi\rangle=0\).  Projection onto \(Q{\cal H}\)
gives
\[
 QD\psi[u]
 =
 (\lambda-Q{\mathbb T}Q)^{-1}
 QD{\mathbb T}[u]\psi.
\]
The first derivative is the Hellmann--Feynman identity
\((83.19)\); differentiating it and inserting the displayed
resolvent gives \((83.20)\).  The scalar chain rule applied to
\((83.18)\) gives \((83.21)\).
\end{proof}

The resolvent in \((83.20)\) is the sector-sensitive version of the
gap denominator already isolated in V72--V73.  At an exact
uncoupled crossing,
\[
 E_{\rm vac}(g)=\min_sE_s(g),                             \tag{83.22}
\]
which need not be differentiable.  A small avoided crossing smooths
the first derivative but can make the second derivative diverge.

\begin{example}[Avoided-crossing curvature]
\label{ex:v14v83-crossing}
For the two-level Hamiltonian
\[
 H_{\rm eff}(t)=e_0(t)I+m(t)\sigma_z+\kappa\sigma_x,
 \qquad \kappa\ne0,
\]
the lower energy is
\[
 E_-(t)=e_0(t)-\sqrt{m(t)^2+\kappa^2}.                   \tag{83.23}
\]
If \(m(0)=0\), then
\[
 E_-''(0)
 =
 e_0''(0)-\frac{m'(0)^2}{|\kappa|}.                     \tag{83.24}
\]
Thus a tunnelling amplitude tending to zero can leave finite-volume
first jets well defined while destroying a uniform vacuum second
jet.  This is directly relevant to the Einstein susceptibility row.
\end{example}

\subsection{Sector acquisition consequence}

\begin{acquisition}[Absolute sector-vacuum card]
\label{acq:v14v83}
For every cutoff and metric point, record:
\begin{enumerate}
\item all block-move sectors and their physical interpretation;
\item each unnormalized leading transfer eigenvalue
      \(\lambda_{s,0}\), sector multiplicity, and internal gap;
\item the absolute energy differences \(E_s-E_*\);
\item every sector-changing block, its norm
      \(\varepsilon\), and its off-diagonal norm
      \(\varepsilon_\perp\);
\item the transfer isolation \(\delta_T\) and the Feshbach debit
      \((83.14)\);
\item the cofinal tunnelling scale
      \(|\kappa|/(a\lambda)\) when sectors are nearly degenerate;
\item the first two metric derivatives of the unnormalized transfer
      and a uniform reduced-resolvent bound;
\item the selected route: superselection, uniform energy separation,
      or uniform tunnelling.
\end{enumerate}
Separate normalization of sector Markov kernels is not an acceptable
substitute for items 2--3.
\end{acquisition}

\begin{assessment}[Current source boundary]
\label{ass:v14v83-source}
The attached sources do not provide the unnormalized finite
sector-transfer spectrum, sector-changing amplitudes, or their
metric jets.  Hence neither the winning sector nor the regularity of
the vacuum Einstein source can presently be certified.  Formal
claims of a unique normalized stationary law inside each sector do
not answer this question.
\end{assessment}

\begin{decision}[V84 fermionic reflection block before sector
assembly]
\label{dec:v14v83-next}
Return to the remaining positivity bottleneck.  Formulate a finite
fermionic Gaussian reflection block whose Berezin integration
produces an explicitly positive boundary Gram determinant.  Prove
the Schur-complement criterion, identify exactly where a complex
determinant phase enters, and state the anomaly/orientation data
needed for compatibility across gauge and topological sectors.
Only a positive common transfer can be inserted into the sector
ledger above.
\end{decision}

\begin{firewall}[Claim boundary after V83]
\label{fire:v14v83}
The sector formulas are exact finite spectral statements.  No
sector energies, tunnelling amplitudes, or metric jets have been
computed from the attached SM--Einstein action.  Finite-volume
Perron uniqueness is not promoted to a uniform continuum vacuum or
mass gap, and the fermionic reflection-positive transfer remains
unconstructed.
\end{firewall}

\section*{Version V83 append-only record}
\addcontentsline{toc}{section}{Version V83 append-only record}
The complete V82 source was retained before this continuation.  It
had \(975902\) bytes, \(29439\) lines, and SHA--256
\[
 \texttt{79A44E39F11D5703BDA2D123071E18655C0BCA403E88EF68E520217818E445CF}.
\]
V83 derived the exact sector spectrum, a quantitative
sector-changing Feshbach estimate, the tunnelling scale, and the
sector-sensitive vacuum metric jets.
\section{Fourteenth-cycle V84: fermionic reflection Gram blocks and determinant-line orientation}
\label{sec:v14v84}

The sector analysis in V83 assumes a common positive transfer.
Bosonic positivity does not supply the missing fermionic sign.  This
continuation gives a finite Gaussian criterion that is stronger than
positivity of a partition-function determinant: it proves positivity
of every fermion-number boundary Gram block.

\subsection{Exterior-power positivity}

Let \({\mathfrak h}\) be a finite-dimensional complex one-particle
space and
\[
 {\cal F}_{-}({\mathfrak h})
 =
 \bigoplus_{p=0}^{\dim{\mathfrak h}}\bigwedge^p{\mathfrak h}
                                                               \tag{84.1}
\]
its fermionic Fock space.  For an operator \(K\) on
\({\mathfrak h}\), define
\[
 \Gamma_-(K)
 :=
 \bigoplus_{p=0}^{\dim{\mathfrak h}}\bigwedge^pK,         \tag{84.2}
\]
with \(\bigwedge^0K=1\).

\begin{theorem}[Fermionic Gram-minor theorem]
\label{thm:v14v84-minors}
If \(K\succeq0\), then
\[
 \Gamma_-(K)\succeq0.                                    \tag{84.3}
\]
In an orthonormal occupation basis indexed by ordered subsets
\(I,J\) of equal cardinality,
\[
 \langle e_I,\Gamma_-(K)e_J\rangle
 =\det K[I,J].                                           \tag{84.4}
\]
Hence, for every coefficient family \(c_I\),
\[
 \sum_{I,J}\overline{c_I}\det K[I,J]c_J\ge0.             \tag{84.5}
\]
If \(K=B^*B\), the left-hand side in fixed degree \(p\) is
\[
 \left\|
  \sum_{|I|=p}c_I\,(\bigwedge^pB)e_I
 \right\|^2.                                             \tag{84.6}
\]
\end{theorem}

\begin{proof}
The definition of the exterior power gives \((84.4)\) by the
alternating multilinear expansion of a determinant.  Write
\(K=B^*B\), for example with \(B=K^{1/2}\).  Functoriality of
exterior powers gives
\[
 \bigwedge^pK
 =
 (\bigwedge^pB)^*(\bigwedge^pB),
\]
which proves positivity and \((84.6)\) in every degree.  Their direct
sum gives \((84.3)\)--\((84.5)\).
\end{proof}

\begin{countertheorem}[A positive determinant is not reflection
positivity]
\label{ctr:v14v84-det}
For
\[
 K=-I_2,
\]
one has \(\det K=1>0\), but the one-particle block
\(\bigwedge^1K=-I_2\) is negative.  Therefore positivity of the full
fermion determinant, or cancellation of its phase, does not imply
the boundary Gram inequalities \((84.5)\).
\end{countertheorem}

\begin{proof}
The determinant is the two-particle exterior power, whereas the
one-particle exterior power is \(K\) itself.
\end{proof}

\subsection{Berezin Gaussian Schur block}

Fix an ordering of all Grassmann variables and its associated
Berezin measure.  Let \(A=A^*\succ0\) act on a finite bulk
one-particle space, and let \(B\) map the bulk space to the boundary
space.  With boundary generators
\(\overline\chi_+\) and \(\chi_-\), the standard ordered Gaussian
identity is
\[
 \begin{split}
 &\int d\overline\eta\,d\eta\,
 \exp\!\left(
  -\overline\eta A\eta
  +\overline\chi_+B\eta
  +\overline\eta B^*\chi_-
 \right)\\
 &\qquad
 =
 \det(A)\,
 \exp\!\left(
  \overline\chi_+\,BA^{-1}B^*\,\chi_-
 \right).                                                \tag{84.7}
 \end{split}
\]
Changing the global integration ordering changes both sides by the
same orientation sign; the ordering must remain fixed under gluing.

\begin{theorem}[Finite Gaussian fermion reflection block]
\label{thm:v14v84-gaussian}
Under the hypotheses of \((84.7)\), put
\[
 K=BA^{-1}B^*.                                           \tag{84.8}
\]
Then \(K\succeq0\), \(\det A>0\), and the coefficient matrix of the
boundary exponential in every fermion-number degree is
\[
 \det(A)\,[\det K[I,J]]_{|I|=|J|=p}\succeq0.             \tag{84.9}
\]
After quotienting its null space, it defines a positive fermionic
OS boundary Hilbert space.
\end{theorem}

\begin{proof}
Since
\[
 K=(BA^{-1/2})(BA^{-1/2})^*,
\]
it is positive semidefinite, and positivity of \(A\) gives
\(\det A>0\).  Expanding the Grassmann exponential in the fixed
ordered monomial basis gives the minors in \((84.9)\).  Apply
Theorem~\ref{thm:v14v84-minors} and multiply by the positive scalar
\(\det A\).  A positive semidefinite Gram form descends to a positive
definite inner product on its quotient by the null space.
\end{proof}

For comparison, the ordinary block identity
\[
 \begin{pmatrix}A&C\\ C^*&D\end{pmatrix}\succeq0
 \quad\Longleftrightarrow\quad
 D-C^*A^{-1}C\succeq0                                   \tag{84.10}
\]
holds when \(A\succ0\), and
\[
 \det\begin{pmatrix}A&C\\ C^*&D\end{pmatrix}
 =
 \det A\,\det(D-C^*A^{-1}C).                             \tag{84.11}
\]
Equations \((84.10)\)--\((84.11)\) are useful elimination tests, but
\((84.11)\) alone checks only a top exterior power.
Theorem~\ref{thm:v14v84-gaussian} checks the full Fock Gram matrix.

\subsection{Dynamic bosonic boundary configurations}

Pointwise positivity at one frozen background is still insufficient
when the boundary geometry and gauge field vary.  Let \(Q\) be a
finite set of bosonic boundary configurations.  Suppose one common
positive bulk matrix \(A\) and maps \(B(q)\) are available, and
define the operator-valued kernel
\[
 K(q,q')=B(q)A^{-1}B(q')^*.                              \tag{84.12}
\]

\begin{lemma}[Operator-valued feature factorization]
\label{lem:v14v84-feature}
For arbitrary boundary vectors \(v_q\),
\[
 \sum_{q,q'}\langle v_q,K(q,q')v_{q'}\rangle
 =
 \left\|
  \sum_qA^{-1/2}B(q)^*v_q
 \right\|^2\ge0.                                         \tag{84.13}
\]
For each fermion number \(p\), the matrix
\[
 {\cal K}_p((q,I),(q',J))
 :=
 \det K(q,q')[I,J]                                      \tag{84.14}
\]
is positive semidefinite.
\end{lemma}

\begin{proof}
Equation \((84.13)\) is direct expansion of the square.  For
\((84.14)\), associate to \((q,I)\) the feature vector
\[
 \Phi_{q,I}
 :=
 \left(\bigwedge^p(A^{-1/2}B(q)^*)\right)e_I.
\]
Its inner product with \(\Phi_{q',J}\) is the determinant in
\((84.14)\), by the exterior-power determinant identity.  Thus
\({\cal K}_p\) is a Gram matrix.
\end{proof}

\begin{theorem}[Coupled boson--fermion Schur-product criterion]
\label{thm:v14v84-coupled}
Let \(G(q,q')\) be a positive semidefinite bosonic reflection kernel.
For every \(p\), the coupled boundary matrix
\[
 {\cal J}_p((q,I),(q',J))
 :=
 G(q,q')\,{\cal K}_p((q,I),(q',J))                      \tag{84.15}
\]
is positive semidefinite.  Positive endpoint weights
\(w(q)w(q')\) preserve this conclusion.
\end{theorem}

\begin{proof}
After indexing rows by \((q,I)\), lift \(G\) by repeating its entry
over the occupation indices.  This lifted matrix is positive
semidefinite because it is the tensor product of \(G\) with an
all-ones Gram matrix on the repeated labels.  Equation \((84.15)\)
is its entrywise product with the positive matrix
\({\cal K}_p\).  The Schur product theorem gives positivity.
Endpoint weights act by a diagonal congruence.
\end{proof}

This theorem is a finite sufficient realization of the V76
\emph{fermionic Schur insertion}.  Its load-bearing requirement is
the common feature factorization \((84.12)\).  Separate positivity of
\(K(q,q)\) for every \(q\) does not imply positivity of the
off-diagonal kernel \(K(q,q')\).

\subsection{Faithfulness and the fermion Hamiltonian}

Suppose \(0\prec K\preceq I\), and define the one-particle
Hamiltonian
\[
 h=-\frac1a\log K\succeq0.                               \tag{84.16}
\]

\begin{proposition}[Exact second quantization]
\label{prop:v14v84-second}
One has
\[
 \Gamma_-(K)=e^{-a\,d\Gamma(h)},                         \tag{84.17}
\]
where
\[
 d\Gamma(h)\big|_{\wedge^p{\mathfrak h}}
 =
 \sum_{r=1}^p
 I\wedge\cdots\wedge h\wedge\cdots\wedge I.              \tag{84.18}
\]
If the eigenvalues of \(h\) are
\(0\le\epsilon_1\le\cdots\), the Fock energies are the sums over
occupied distinct modes.  In a sector where the vacuum is the empty
state and \(\epsilon_1>0\), the fermionic excitation gap is
\(\epsilon_1\).
\end{proposition}

\begin{proof}
Diagonalize \(K\).  An occupation vector with occupied set \(I\) has
\(\Gamma_-(K)\)-eigenvalue
\(\prod_{j\in I}e^{-a\epsilon_j}
 =e^{-a\sum_{j\in I}\epsilon_j}\), which is exactly
\((84.17)\)--\((84.18)\).
\end{proof}

If \(K\) has a kernel, the OS quotient removes some occupation
vectors and the logarithm is not a faithful finite Hamiltonian on
the unquotiented space.  If \(K\) has eigenvalue greater than one,
a scalar normalization can make the transfer contractive, but that
scalar contributes to the absolute vacuum energy and must be kept
for V72--V73 and V83.

\subsection{Phase and orientation across charts and sectors}

If \(A\) is non-Hermitian, \(\det A\) may have a complex phase and
\(BA^{-1}B^*\) need not be positive.  Pairing a determinant with its
complex conjugate makes the scalar
\[
 \det A\,\overline{\det A}=|\det A|^2\ge0,               \tag{84.19}
\]
but Countertheorem~\ref{ctr:v14v84-det} shows why scalar phase
cancellation is only a necessary part of the full Gram problem.

There is also a global orientation question.  Let chart or sector
labels be vertices of a connected graph, and suppose overlap
identifications of the determinant line carry phases
\[
 u_{\alpha\beta}\in U(1),\qquad
 u_{\beta\alpha}=\overline{u_{\alpha\beta}}.             \tag{84.20}
\]

\begin{theorem}[Finite determinant-line holonomy test]
\label{thm:v14v84-holonomy}
There exist vertex phases \(v_\alpha\in U(1)\) such that
\[
 u_{\alpha\beta}=v_\alpha\overline{v_\beta}              \tag{84.21}
\]
on every overlap edge if and only if the ordered product of the
\(u\)'s around every graph cycle equals one.  When \((84.21)\)
holds, diagonal rephasing by the \(v_\alpha\)'s removes all overlap
phases without changing any Gram eigenvalue.
\end{theorem}

\begin{proof}
If \((84.21)\) holds, phases telescope around a cycle.  Conversely,
fix a root and a spanning tree.  Define \(v_\alpha\) by multiplying
edge phases along the unique tree path from the root to \(\alpha\),
with the orientation chosen to satisfy \((84.21)\) on tree edges.
For a non-tree edge, that edge plus the two tree paths forms a cycle;
trivial cycle product forces \((84.21)\) there as well.  Diagonal
unitary congruence implements the rephasing and preserves the
spectrum and positivity.
\end{proof}

A nontrivial cycle phase is the finite determinant-line obstruction
that must be ruled out by anomaly cancellation and a compatible
Berezin orientation.  Checking a positive convention separately in
each chart does not check this holonomy.

\subsection{Fermionic acquisition ledger}

\begin{acquisition}[Reflection-positive fermion card]
\label{acq:v14v84}
For every finite cutoff and physical sector, record:
\begin{enumerate}
\item the ordered bulk and boundary Grassmann variables and the
      reflection identification;
\item the eliminated bulk matrix \(A\), with a proof of
      \(A=A^*\succ0\), or a stated alternative positive
      factorization;
\item the boundary maps \(B(q)\) and the common feature formula
      \((84.12)\);
\item all occupation Gram blocks \((84.14)\), including their null
      spaces and smallest retained eigenvalues;
\item the bosonic OS kernel used in the Schur product
      \((84.15)\);
\item the absolute determinant scalar and its first two metric jets;
\item determinant-line transition phases and every independent cycle
      holonomy;
\item the gauge, flavour, chirality, and topological-sector
      orientation data that make the card compatible under gluing.
\end{enumerate}
Only after these entries pass can the fermionic transfer be inserted
into the sector and Dobrushin ledgers.
\end{acquisition}

\begin{assessment}[Current source boundary]
\label{ass:v14v84-source}
The attached sources name fermionic, chiral, and anomaly sectors but
do not provide a finite reflected bulk matrix \(A\), the common
boundary features \(B(q)\), all occupation Gram blocks, or the
determinant-line transition cocycle.  Therefore the theorem is an
exact acquisition target, not a proof that the Standard Model
fermionic measure is reflection positive.
\end{assessment}

\begin{decision}[V85 compulsory audit, then common-card synthesis]
\label{dec:v14v84-next}
At the five-version checkpoint inspect the newest Yang--Mills and
Navier--Stokes notes.  Test whether either supplies a new exact
factorization, constrained-norm device, sector transport, or cofinal
screen that can sharpen \((84.12)\)--\((84.15)\).  Then assemble one
finite common-card theorem combining bosonic reflection positivity,
fermionic occupation blocks, constraint-preserving sector dynamics,
and the absolute Einstein source, with every still-unpopulated datum
visible.
\end{decision}

\begin{firewall}[Claim boundary after V84]
\label{fire:v14v84}
The exterior-power and Gaussian theorems are finite sufficient
criteria.  Their matrices and orientation cocycle have not been
extracted from a regulated chiral Standard Model action.  Determinant
positivity alone is expressly insufficient.  No interacting
fermionic continuum, anomaly cancellation theorem, or full
SM--Einstein continuum is claimed.
\end{firewall}

\section*{Version V84 append-only record}
\addcontentsline{toc}{section}{Version V84 append-only record}
The complete V83 source was retained before this continuation.  It
had \(989446\) bytes, \(29855\) lines, and SHA--256
\[
 \texttt{335B7D1090361A1C991ED8AD7410688D98AAB6A95B072E63CF4A6AECAF406F7F}.
\]
V84 proved the fermionic exterior-power Gram criterion, its Gaussian
and coupled-background factorizations, and the determinant-line
cycle-holonomy test.
\section{Fourteenth-cycle V85: compulsory companion audit and transport-safe synthesis}
\label{sec:v14v85}

This is the compulsory checkpoint after V80.  The newest companion
files were inspected read-only before assembling the common finite
SM--Einstein card.

\subsection{Frozen provenance}

\begin{record}[V85 companion checkpoint]
\label{rec:v14v85-companions}
The Yang--Mills file is
\begin{center}
\path{Renewal_Yang_Mills_Mass_Gap_Research_Notes_V27.tex},
\end{center}
with \(218881\) bytes, \(5962\) lines, and SHA--256
\[
 \texttt{848978B5866A0D969E831640A6A55BAFA6F4FBBFF9ACCB34D1743355D37395F5}.
                                                               \tag{85.1}
\]
The Navier--Stokes file remains
\begin{center}
\path{Renewal_NS_Stability_Research_Notes_V26.tex},
\end{center}
with \(278283\) bytes, \(5319\) lines, and SHA--256
\[
 \texttt{82C887F8C2C69E4F28FAD7C3DC8DE5B9099CB5A4BF431EBA1FFF3E91935978AD}.
                                                               \tag{85.2}
\]
\end{record}

\subsection{New Yang--Mills information}

Yang--Mills V26 directly rebuilds its rational metric and vertex
certificate at three transverse permutations rather than assuming
that the spanning-tree reduction preserves manifest coordinate
symmetry.  It proves three radius-\(1/8\) boxes and an exact
noncoverage witness.  V27 then reverses the elimination order: a
fixed noncommutative splitting handles the complete active momentum
circle, while a six-variable whole-matrix hull proves a transverse
radius-\(1/16\) tube.  It also records that copying such tubes over
the quarter grid cannot cover the full transverse torus.

\begin{assessment}[YM applicability]
\label{ass:v14v85-ym}
No YM V26--V27 theorem constructs the fermionic feature maps
\(B(q)\), an interacting positive measure, or a cofinal compact
screen.  Their exact floors therefore do not populate V84.

Two structural lessons do apply:
\begin{enumerate}
\item after constraint reduction, a physical permutation or sector
      symmetry must be supported by an explicit unitary/congruence
      intertwiner or by an independent certificate; it cannot be
      inferred from the unreduced variables;
\item elimination should preserve the entire active matrix block
      before any outer hull or norm bound is taken.
\end{enumerate}
The first lesson strengthens the determinant-line transport row in
V84.  The second agrees with the common-feature and exterior-power
factorizations \((84.12)\)--\((84.15)\).  Neither imports a numerical
constant.
\end{assessment}

\subsection{Navier--Stokes status}

Navier--Stokes V26 is unchanged.  Its exact rank-one contraction
continues to support only the methodological instruction already
recorded in V80: take influence and response suprema over the actual
constraint-compatible variation set.  It supplies no probabilistic
or fermionic datum.

\begin{theorem}[V85 audit verdict]
\label{thm:v14v85-verdict}
At checkpoints \((85.1)\)--\((85.2)\), neither companion supplies a
common positive SM--Einstein transfer, a determinant-line
trivialization, a uniform influence bound, or a cofinal source
limit.  The only new import is a required validation row for every
symmetry or sector transport after constraint reduction.
\end{theorem}

\begin{proof}
The Yang--Mills conclusions are finite Wilson-pencil inequalities on
proper subsets of momentum space and explicitly leave the global
and cofinal rows open.  The Navier--Stokes conclusion is a
deterministic Oseen-kernel norm.  Neither has the types required by
V81--V84.  Yang--Mills V26 directly demonstrates why unreduced
symmetry does not by itself certify the reduced representation,
which proves the stated methodological addition.
\end{proof}

\begin{decision}[V86 common finite-card theorem]
\label{dec:v14v85-next}
Assemble one theorem with typed hypotheses for:
\begin{enumerate}
\item bosonic reflection-factorized adjacent kernels;
\item fermionic exterior-power Gram blocks;
\item constraint-preserving sector charts and block coverage;
\item a Dobrushin or alternative phase estimate;
\item unnormalized sector transfer energies and metric jets;
\item explicit reduced symmetry intertwiners and determinant-line
      holonomy.
\end{enumerate}
Prove exactly which finite conclusions follow when the rows are
populated, and construct a dependency graph showing which source
objects must be acquired before the cofinal OS and Einstein limits
can even be posed.  This will prevent a local symbol floor from
being reused as a missing global input.
\end{decision}

\begin{firewall}[Claim boundary after V85]
\label{fire:v14v85}
The audit changes the acquisition discipline, not the proved
physical scope.  The full common regulator, chiral fermion
orientation, thermodynamic phase control, OS limit, and Einstein
source convergence remain unconstructed.
\end{firewall}

\section*{Version V85 append-only record}
\addcontentsline{toc}{section}{Version V85 append-only record}
The complete V84 source was retained before this continuation.  It
had \(1002981\) bytes, \(30241\) lines, and SHA--256
\[
 \texttt{71EB45A9A6DCE6986D668CD7AEB74BFF13C097993E3FF24DECF68E7689B5F171}.
\]
V85 performed the compulsory companion audit and added explicit
post-reduction symmetry transport to the common-card requirements.
\section{Fourteenth-cycle V86: the common finite SM--Einstein card}
\label{sec:v14v86}

V76--V85 produced separate finite compilers for reflection kernels,
local dynamics, phase control, sectors, fermions, and vacuum metric
jets.  This continuation states their exact composition.  The
purpose is not to rename the missing continuum theorem, but to
prevent a datum of one type from being used in place of another.

\subsection{Typed finite data}

Fix a regulator \(n\), Euclidean time step \(a_n>0\), a reduced
metric parameter \(g\), and a physical sector \(s\).  The finite
common card consists of the following rows.

\begin{enumerate}
\item[\textup{(F1)}] A finite common bosonic--fermionic action and
reflection identify a positive semidefinite boundary Gram matrix
\({\cal J}_{n,s}(g)\).  Its fermionic blocks have the feature form
\((84.12)\)--\((84.15)\), with a fixed Berezin orientation and
trivial determinant-line cycle holonomy.

\item[\textup{(F2)}] Time gluing is compatible with the OS null
space and produces, on the quotient
\({\cal H}_{n,s}(g)\), a strictly positive self-adjoint transfer
\({\mathbb T}_{n,s}(g)\).  Its exact Hamiltonian is
\[
 H_{n,s}(g)=-a_n^{-1}\log{\mathbb T}_{n,s}(g).            \tag{86.1}
\]

\item[\textup{(F3)}] The constraint-preserving block form
\({\cal E}_{n,s}^{\rm upd}\) has a constraint-solving product chart,
fiber bounds \(g_b\), fractional coverage \(\kappa_{n,s}\), and
chart influence row sum \(\alpha_{n,s}<1\).  Thus
\[
 \gamma_{n,s}^{\rm upd}
 :=
 \kappa_{n,s}(1-\alpha_{n,s})>0.                         \tag{86.2}
\]

\item[\textup{(F4)}] There is a unitary physical-correspondence map
\[
 U_{n,s}:L^2_0(\widehat\mu_{n,s})\oplus\mathbb C{\bf1}
 \longrightarrow{\cal H}_{n,s},                         \tag{86.3}
\]
sending \({\bf1}\) to the normalized transfer vacuum
\(\psi_{n,s}\), and for every \(f\perp{\bf1}\),
\[
 \left|
 \langle U_{n,s}f,(H_{n,s}-E_{n,s})U_{n,s}f\rangle
 -
 {\cal E}_{n,s}^{\rm upd}(f)
 \right|
 \le r_{n,s}\|f\|^2,                                    \tag{86.4}
\]
where
\[
 E_{n,s}=-a_n^{-1}\log\lambda_{n,s,0},\qquad
 r_{n,s}<\gamma_{n,s}^{\rm upd}.                         \tag{86.5}
\]

\item[\textup{(F5)}] Every claimed gauge, coordinate, reflection, or
sector symmetry after reduction has an explicit unitary or
congruence intertwiner preserving the action, null space, transfer,
and determinant-line orientation.  Where no such intertwiner is
proved, the target chart is independently certified.

\item[\textup{(F6)}] The leading unnormalized sector eigenvalues are
retained.  One sector \(s_*\) has a simple vacuum and satisfies
\[
 \delta_{n}^{\rm sec}
 :=
 \min_{r\ne s_*}(E_{n,r}-E_{n,s_*})>0.                  \tag{86.6}
\]
If sector-changing blocks are present, they instead satisfy the V83
isolation and Feshbach hypotheses.

\item[\textup{(F7)}] The functions
\({\mathbb T}_{n,s_*}(g)\), the reduced gravitational action
\(S^{\rm grav}_{n}(g)\), and the counterterm action \(C_n(g)\) are
\(C^2\).  At a candidate reduced Euclidean metric \(g_n^*\),
\[
 D\Gamma_n(g_n^*)=0,\qquad
 D^2\Gamma_n(g_n^*)\succeq m_nI,\qquad m_n>0,            \tag{86.7}
\]
where
\[
 \Gamma_n(g)
 :=
 S^{\rm grav}_{n}(g)+C_n(g)+E_{n,s_*}(g).               \tag{86.8}
\]
\end{enumerate}

The correspondence row (F4) is deliberately explicit.  A fast
boundary update chain and the physical time-transfer Hamiltonian are
different operators until \((86.3)\)--\((86.4)\) are proved.

\subsection{Finite composition theorem}

\begin{theorem}[Common finite-card theorem]
\label{thm:v14v86-card}
If (F1)--(F7) hold at one regulator, then:
\begin{enumerate}
\item the OS quotient is a positive Hilbert space and
      \({\mathbb T}_{n,s}\) defines a self-adjoint finite Hamiltonian
      in every certified sector;
\item the physical Hamiltonian gap in sector \(s\) obeys
\[
 \gamma_{n,s}^{\rm phys}
 \ge
 \gamma_{n,s}^{\rm upd}-r_{n,s}
 =
 \kappa_{n,s}(1-\alpha_{n,s})-r_{n,s}>0;                \tag{86.9}
\]
\item when sectors do not mix, \(s_*\) is the unique global vacuum
      sector and
\[
 \gamma_n^{\rm glob}
 \ge
 \min\{\gamma_{n,s_*}^{\rm phys},
        \delta_n^{\rm sec}\};                            \tag{86.10}
\]
\item the vacuum first and second metric jets exist and are given by
      \((83.19)\)--\((83.21)\);
\item after shrinking to a convex neighbourhood of \(g_n^*\),
      \(\Gamma_n\) is strictly convex there and \(g_n^*\) is its
      unique critical point in that neighbourhood.
\end{enumerate}
If a sector-changing transfer block is present, conclusions 2 and
4 hold with the perturbed vacuum, while conclusion 3 is replaced by
the gap bound \((83.10)\) and the source jets use the perturbed
reduced resolvent.
\end{theorem}

\begin{proof}
Row (F1), Theorem~\ref{thm:v14v84-coupled}, and the OS null quotient
give the positive boundary Hilbert space.  Row (F2) supplies the
strictly positive self-adjoint time transfer, so the functional
calculus defines \((86.1)\).

By Theorems \ref{thm:v14v81-dobrushin} and
\ref{thm:v14v82-block}, row (F3) gives
\[
 {\cal E}_{n,s}^{\rm upd}(f)
 \ge\gamma_{n,s}^{\rm upd}\|f\|^2
 \quad(f\perp{\bf1}).                                    \tag{86.11}
\]
Apply the lower half of \((86.4)\):
\[
 \langle U_{n,s}f,(H_{n,s}-E_{n,s})U_{n,s}f\rangle
 \ge
 (\gamma_{n,s}^{\rm upd}-r_{n,s})\|f\|^2.
\]
Since \(U_{n,s}\) is onto and maps the constant to the vacuum, the
Rayleigh principle gives \((86.9)\).

Row (F6) and Theorem~\ref{thm:v14v83-direct} make the unique global
ground state lie in \(s_*\).  The first excitation is either internal
to that sector or the ground state of another sector, yielding
\((86.10)\).  In the mixing case use
Theorem~\ref{thm:v14v83-perturb}.

The simple isolated transfer eigenvalue in (F6), together with the
\(C^2\) dependence in (F7), permits the differentiation in
Proposition~\ref{prop:v14v83-jets}.  Finally, continuity of
\(D^2\Gamma_n\) and \((86.7)\) give a convex neighbourhood on which
\(D^2\Gamma_n\succeq(m_n/2)I\).  Integration of the Hessian along
line segments gives strict convexity, so there is at most one
critical point there; \(g_n^*\) is one by \((86.7)\).
\end{proof}

\begin{corollary}[Fully explicit weak-dependence lower bound]
\label{cor:v14v86-explicit}
If each physical block \(b\) has a V78 base gap \(g_b^0\) and
conditional oscillation \(\Theta_b\), and if a fractional allocation
\(\theta_{bi}\) satisfies
\[
 \sum_{i\in S_b}\theta_{bi}
 \le e^{-2\Theta_b}g_b^0,
\]
then \((86.9)\) holds with
\[
 \gamma_{n,s}^{\rm phys}
 \ge
 \left[
  \min_i\sum_{b:\,i\in S_b}\theta_{bi}
 \right]
 \left[
  1-\max_i\sum_{j\ne i}
       \tanh\frac{\widehat\Delta_{ij}}4
 \right]
 -r_{n,s}.                                               \tag{86.12}
\]
Every quantity in \((86.12)\) is finite arithmetic once the reduced
conditional action and correspondence matrix are supplied.
\end{corollary}

\begin{proof}
Use \((82.13)\), Proposition
\ref{prop:v14v81-action}, Definition
\ref{def:v14v82-coverage}, and Theorem
\ref{thm:v14v86-card}.
\end{proof}

\subsection{The dependency graph}

The logical order of the finite construction is
\[
\begin{array}{c}
 \text{common regulated action plus reflection}\\
 \downarrow\\
 \text{bosonic Gram factor and fermionic common features}\\
 \downarrow\\
 \text{OS null quotient plus time gluing}\\
 \downarrow\\
 {\mathbb T}_{n,s}>0\quad\longrightarrow\quad
 H_{n,s}=-a_n^{-1}\log{\mathbb T}_{n,s}
\end{array}
\qquad
\begin{array}{c}
 \text{constraint equations plus sector choice}\\
 \downarrow\\
 \text{product chart or support-stable atlas}\\
 \downarrow\\
 \text{fiber gaps, coverage, mixed influences}\\
 \downarrow\\
 \gamma_{n,s}^{\rm upd}\\
 \downarrow\ {\rm correspondence}\\
 \gamma_{n,s}^{\rm phys}.
\end{array}                                              \tag{86.13}
\]
The absolute branch is
\[
 \begin{array}{c}
 \text{unnormalized sector transfers}\\
 \downarrow\\
 \text{sector energies, isolation, and metric jets}\\
 \downarrow\\
 \text{unique vacuum or controlled tunnelling}\\
 \downarrow\\
 E_{\rm vac}(g)\ \text{in}\ 
 S_{\rm grav}(g)+C(g)+E_{\rm vac}(g).
 \end{array}                                             \tag{86.14}
\]
Reduced symmetry intertwiners and determinant-line orientation gate
every horizontal identification in \((86.13)\)--\((86.14)\).

\begin{proposition}[Typed non-substitution]
\label{prop:v14v86-types}
A local bosonic symbol floor can populate part of a fiber bound
\(g_b^0\), but by itself implies none of:
\[
 {\cal J}\succeq0,\quad
 \alpha<1,\quad
 r<\gamma^{\rm upd},\quad
 \delta^{\rm sec}>0,\quad
 D^2\Gamma\succeq mI.                                   \tag{86.15}
\]
Likewise, a positive full fermion determinant does not imply the
first item in \((86.15)\).
\end{proposition}

\begin{proof}
The V79 expander family has exact one-cell heat-bath gaps but a
collapsing global gap, so a local floor does not imply
\(\alpha<1\) or a physical gap.  The physical transfer can be chosen
independently of an auxiliary update chain, so no update floor
implies the correspondence residual.  Adding sector-dependent
scalar multiples changes \(\delta^{\rm sec}\) and the vacuum metric
jets without changing normalized local dynamics.  Gravitational
counterterms change \(D^2\Gamma\) without changing the matter symbol.
Finally Countertheorem~\ref{ctr:v14v84-det} gives a positive
determinant with a negative one-particle Gram block.
\end{proof}

\subsection{Cofinal handoff}

Theorem~\ref{thm:v14v86-card} is finite.  A continuum theorem needs
uniform and convergent versions of different rows, not just a
sequence of finite positive numbers.

\begin{acquisition}[Common-card cofinal columns]
\label{acq:v14v86-cofinal}
For a cofinal family, add:
\begin{enumerate}
\item cylinder consistency or tightness of the common reflected
      measures and stability of the OS null quotients;
\item a common limiting physical Hilbert-space identification;
\item norm-resolvent or Mosco convergence of the physical
      Hamiltonians near the isolated vacuum;
\item positive lower bounds on \((86.9)\) and \((86.10)\);
\item compactness of the vacuum projections and exclusion of hidden
      soft modes;
\item convergence of the absolute sector energies after declared
      counterterms;
\item \(C^2\) convergence of the vacuum metric jets and a uniform
      reduced Einstein Hessian floor;
\item stable anomaly cancellation, determinant-line orientation, and
      reduced symmetry intertwiners;
\item recovery of the declared local relativistic SM field algebra
      and Einstein covariance in the limit.
\end{enumerate}
\end{acquisition}

\begin{assessment}[Current source boundary]
\label{ass:v14v86-source}
No attached source supplies all rows (F1)--(F7) for one common
regulator.  In particular, the fermionic common features, the
constraint-solved conditional law, the update-to-physical
correspondence residual, unnormalized sector energies, and the
finite effective Einstein critical point are not jointly populated.
The finite theorem therefore cannot yet be instantiated.
\end{assessment}

\begin{decision}[V87 cofinal common-card theorem]
\label{dec:v14v86-next}
Prove a precise cofinal handoff from finite common cards.  Use
norm-resolvent convergence and convergence of rank-one vacuum
projections to preserve a unique spectral gap, and use locally
uniform \(C^2\) convergence of the effective actions plus a uniform
Hessian floor to preserve the reduced Einstein solution.  Include
counterexamples showing why strong resolvent convergence and
pointwise convergence of metric first jets are insufficient.
\end{decision}

\begin{firewall}[Claim boundary after V86]
\label{fire:v14v86}
The common-card theorem is conditional on typed finite inputs that
have not been extracted from one regulated Standard Model--Einstein
action.  It proves a finite conclusion only.  No cofinal measure,
continuum field algebra, anomaly-free chiral transfer, or Einstein
continuum solution has been constructed.
\end{firewall}

\section*{Version V86 append-only record}
\addcontentsline{toc}{section}{Version V86 append-only record}
The complete V85 source was retained before this continuation.  It
had \(1008270\) bytes, \(30368\) lines, and SHA--256
\[
 \texttt{3906FD1C5A29C52DAF2D9D6BC830E8FA3FD42B782346B03E3922C3D5E82FC795}.
\]
V86 assembled the typed finite common-card theorem, its explicit
gap formula, metric closure, and dependency graph.
\section{Fourteenth-cycle V87: cofinal preservation of the vacuum gap and Einstein solution}
\label{sec:v14v87}

V86 composes the finite rows.  A continuum conclusion requires a
topology strong enough to prevent the vacuum and the metric
susceptibility from escaping.  This continuation proves two cofinal
handoffs on fixed identified spaces:
\[
 \text{norm resolvent for the Hamiltonian},\qquad
 C^2_{\rm loc}\text{ for the reduced effective action}. \tag{87.1}
\]

\subsection{Norm-resolvent preservation of a simple gapped vacuum}

Let \({\cal H}\) be a Hilbert space after the finite OS spaces have
been transported by the declared common physical identifications.
Let \(H_n\ge0\) and \(H\ge0\) be self-adjoint, and assume
\[
 \left\|(H_n+I)^{-1}-(H+I)^{-1}\right\|\longrightarrow0. \tag{87.2}
\]
Suppose every \(H_n\) has a normalized unique vacuum \(\psi_n\) and
a uniform gap
\[
 H_n\psi_n=0,\qquad
 H_n\big|_{\psi_n^\perp}\succeq\gamma I,\qquad
 \gamma>0.                                               \tag{87.3}
\]

\begin{theorem}[Norm-resolvent vacuum compactness and gap]
\label{thm:v14v87-norm}
Under \((87.2)\)--\((87.3)\), the limit \(H\) has a unique normalized
vacuum up to phase and
\[
 H\big|_{\ker(H)^\perp}\succeq\gamma I.                  \tag{87.4}
\]
Writing
\[
 P_n=|\psi_n\rangle\langle\psi_n|,\qquad
 P={\bf1}_{\{0\}}(H),
\]
one has
\[
 \|P_n-P\|\longrightarrow0,\qquad \operatorname{rank}P=1. \tag{87.5}
\]
Thus the vacuum projections cannot escape weakly to zero.
\end{theorem}

\begin{proof}
Put
\[
 R_n=(H_n+I)^{-1},\qquad R=(H+I)^{-1},\qquad
 q=(1+\gamma)^{-1}<1.                                   \tag{87.6}
\]
By \((87.3)\),
\[
 R_nP_n=P_n,\qquad
 \|R_n(I-P_n)\|\le q,                                    \tag{87.7}
\]
so
\[
 \sigma(R_n)\subset[0,q]\cup\{1\},                       \tag{87.8}
\]
and the spectral projection of \(R_n\) at \(1\) is the rank-one
projection \(P_n\).

Choose a circle \({\cal C}\) centred at \(1\) with radius strictly
smaller than \((1-q)/2\).  Norm convergence \(R_n\to R\) implies
uniform resolvent convergence on \({\cal C}\) for all large \(n\).
The Riesz projections
\[
 \frac1{2\pi i}\int_{\cal C}(z-R_n)^{-1}\,dz
\]
therefore converge in norm to the corresponding projection for
\(R\).  They are \(P_n\), so their limit \(P\) satisfies
\(\|P_n-P\|\to0\).  Once this norm is less than one, two orthogonal
projections have the same finite rank; hence \(\operatorname{rank}P=1\).

The spectral inclusion theorem for norm-convergent bounded
self-adjoint operators gives
\[
 \sigma(R)\subset[0,q]\cup\{1\}.                         \tag{87.9}
\]
By spectral mapping for \(R=(H+I)^{-1}\), the eigenvalue one
corresponds exactly to \(H=0\), and every other finite spectral value
\(E\) of \(H\) obeys
\((1+E)^{-1}\le q\), hence \(E\ge\gamma\).  This proves
\((87.4)\)--\((87.5)\).
\end{proof}

\begin{corollary}[Vacuum expectations]
\label{cor:v14v87-expect}
Choose phases so that \(\psi_n\to\psi\), where
\(P=|\psi\rangle\langle\psi|\).  If bounded observables satisfy
\(\|A_n-A\|\to0\), then
\[
 \langle\psi_n,A_n\psi_n\rangle
 \longrightarrow
 \langle\psi,A\psi\rangle.                               \tag{87.10}
\]
\end{corollary}

\begin{proof}
Norm convergence of rank-one projections permits a phase choice
with norm convergence of their unit vectors.  Add and subtract the
three terms obtained by replacing \(A_n\) with \(A\) and each
\(\psi_n\) with \(\psi\).
\end{proof}

\subsection{Why strong resolvent convergence is insufficient}

\begin{countertheorem}[A unique vacuum escapes under strong
resolvent convergence]
\label{ctr:v14v87-strong}
On \(\ell^2(\mathbb N)\), let \(P_n\) project onto the \(n\)-th basis
vector and set
\[
 H_n=I-P_n.                                               \tag{87.11}
\]
Every \(H_n\) has a unique zero-energy vacuum and gap one, but
\[
 H_n\longrightarrow I
\]
strongly and in strong resolvent sense.  The limit has no
zero-energy state.  Moreover
\[
 P_n\longrightarrow0
\]
strongly but not in norm.
\end{countertheorem}

\begin{proof}
For every \(x\in\ell^2\), the coordinate
\(\langle e_n,x\rangle\) tends to zero, hence \(P_nx\to0\) and
\(H_nx\to x\).  Also
\[
 (H_n+I)^{-1}=\frac12I+\frac12P_n
 \longrightarrow\frac12I=(I+I)^{-1}
\]
strongly.  Each \(H_n\) has spectrum \(\{0,1\}\), whereas the limit
has spectrum \(\{1\}\).  Finally \(\|P_n\|=1\) for every \(n\).
\end{proof}

This is the abstract version of the hidden-soft-mode warning in
V63 and V68.  A uniform finite gap and strong convergence alone do
not retain the vacuum.

\subsection{Cofinal sector selection}

Let \(E_{n,s}^{\rm ren}(g)\) be the renormalized absolute sector
ground energies, using counterterms whose scalar parts are common
across sectors unless a physical sector-dependent term is explicitly
declared.

\begin{proposition}[Uniform sector winner persists]
\label{prop:v14v87-sector}
Suppose on a metric neighbourhood \(K\),
\[
 \sup_{g\in K}
 |E_{n,s}^{\rm ren}(g)-E_s^{\rm ren}(g)|
 \longrightarrow0                                      \tag{87.12}
\]
uniformly in the relevant finite sector set, and one sector \(s_*\)
satisfies
\[
 \inf_{n}\inf_{g\in K}\min_{r\ne s_*}
 \left(E_{n,r}^{\rm ren}(g)-E_{n,s_*}^{\rm ren}(g)\right)
 \ge\delta_{\rm sec}>0.                                  \tag{87.13}
\]
Then
\[
 E_r^{\rm ren}(g)-E_{s_*}^{\rm ren}(g)
 \ge\delta_{\rm sec}
\]
for all \(g\in K\) and \(r\ne s_*\).  In particular, no sector
crossing appears in the limit.
\end{proposition}

\begin{proof}
Pass to the limit in the inequality \((87.13)\) using uniform
convergence \((87.12)\).
\end{proof}

If the number of sectors grows with the cutoff, uniformity over a
fixed list is unavailable.  One then needs tightness in sector
space:
\[
 \inf_{\substack{s\ {\rm outside}\\{\rm a\ fixed\ finite\ set}}}
 \bigl(E_{n,s}^{\rm ren}-E_{n,s_*}^{\rm ren}\bigr)
 \longrightarrow+\infty
\]
or another explicit exclusion of new low-energy sectors.

\subsection{Uniform \(C^2\) passage for the Einstein equation}

Let \(U\subset\mathbb R^d\) be an open reduced Euclidean metric
chart and \(K\Subset U\) a nonempty compact convex set.  Let
\(\Gamma_n,\Gamma\in C^2(U)\) satisfy
\[
 \|\Gamma_n-\Gamma\|_{C^2(K)}\longrightarrow0.           \tag{87.14}
\]
Assume
\[
 D^2\Gamma_n(g)\succeq mI
 \quad(g\in K,\ n\ge n_0),\qquad m>0,                    \tag{87.15}
\]
and for every \(n\ge n_0\) there is \(g_n^*\in K\) with
\[
 D\Gamma_n(g_n^*)=0.                                     \tag{87.16}
\]

\begin{theorem}[Cofinal reduced Einstein solution]
\label{thm:v14v87-einstein}
There is a unique \(g^*\in K\) satisfying
\[
 D\Gamma(g^*)=0,                                         \tag{87.17}
\]
and
\[
 g_n^*\longrightarrow g^*,\qquad
 D^2\Gamma(g^*)\succeq mI.                               \tag{87.18}
\]
The inverse reduced susceptibility is uniformly bounded:
\[
 \|(D^2\Gamma_n(g_n^*))^{-1}\|\le m^{-1},\qquad
 \|(D^2\Gamma(g^*))^{-1}\|\le m^{-1}.                    \tag{87.19}
\]
\end{theorem}

\begin{proof}
Compactness gives a convergent subsequence
\(g_{n_k}^*\to g^*\in K\).  Uniform convergence of the first
derivatives in \((87.14)\), together with \((87.16)\), gives
\(D\Gamma(g^*)=0\).  Uniform convergence of the Hessians and
\((87.15)\) gives
\[
 D^2\Gamma(g)\succeq mI\quad(g\in K).                    \tag{87.20}
\]
For \(g,h\in K\), integration along the line segment gives
\[
 \langle D\Gamma(g)-D\Gamma(h),g-h\rangle
 \ge m\|g-h\|^2.                                         \tag{87.21}
\]
Thus \(\Gamma\) has at most one critical point in \(K\).  Every
convergent subsequence of \(g_n^*\) has that same limit, so the whole
sequence converges.  Equation \((87.19)\) follows from the spectral
lower bounds.
\end{proof}

The theorem applies to
\[
 \Gamma_n
 =
 S^{\rm grav}_n+C_n+E_{n,s_*}^{\rm ren}
\]
only if all three summands converge in the same declared
renormalization scheme.  Convergence of the normalized matter
transfer does not supply convergence of the scalar energy or its
metric jets.

\subsection{Pointwise first jets do not control susceptibility}

\begin{countertheorem}[A moving metric spike]
\label{ctr:v14v87-spike}
There exist smooth \(\Gamma_n:\mathbb R\to\mathbb R\) and critical
points \(x_n=1/n\) such that
\[
 \Gamma_n'(x)\longrightarrow x
 \quad\hbox{for every fixed }x,                          \tag{87.22}
\]
while
\[
 \Gamma_n''(x_n)\longrightarrow-\infty.                 \tag{87.23}
\]
Thus pointwise convergence of first metric jets, even with convergent
critical points, does not control the second jet or the inverse
Einstein response.
\end{countertheorem}

\begin{proof}
Choose \(\phi\in C_c^\infty(\mathbb R)\) with
\[
 \phi(0)=1,\qquad\phi'(0)=1,
\]
and define
\[
 \Gamma_n'(x)
 =
 x-\frac1n\phi\!\left(n^2(x-1/n)\right).                 \tag{87.24}
\]
Integrate once to define \(\Gamma_n\).  At \(x_n=1/n\),
\(\Gamma_n'(x_n)=0\).  For every fixed \(x\), the argument of
\(\phi\) eventually leaves its compact support, so
\((87.22)\) holds.  Differentiation gives
\[
 \Gamma_n''(x_n)=1-n\phi'(0)=1-n,
\]
which proves \((87.23)\).
\end{proof}

A second elementary failure is loss of metric compactness:
\[
 \Gamma_n(x)=\left(x-1+\frac1n\right)^2
 \quad\hbox{on }(0,1).
\]
Its critical point tends to the excluded boundary \(x=1\), although
the derivatives converge pointwise.  A compact interior metric card
is therefore independent of jet convergence.

\subsection{Cofinal composition}

\begin{theorem}[Conditional cofinal common-card handoff]
\label{thm:v14v87-cofinal}
Assume:
\begin{enumerate}
\item one common reflected continuum measure and physical Hilbert
      identification have already been constructed;
\item the finite global Hamiltonians, after the declared absolute
      energy subtraction, satisfy \((87.2)\)--\((87.3)\);
\item the sector energies satisfy
      \((87.12)\)--\((87.13)\), including exclusion of new sectors;
\item the reduced effective actions satisfy
      \((87.14)\)--\((87.16)\);
\item the local observable algebra, covariance, determinant-line
      orientation, and anomaly cancellation pass compatibly to the
      same limit.
\end{enumerate}
Then the limit has a unique gapped vacuum in sector \(s_*\), the
vacuum expectations of norm-convergent bounded observables converge,
and the reduced Euclidean Einstein equation has the unique stable
solution \(g^*\) of Theorem~\ref{thm:v14v87-einstein}.
\end{theorem}

\begin{proof}
Apply Theorem~\ref{thm:v14v87-norm}, Corollary
\ref{cor:v14v87-expect}, Proposition
\ref{prop:v14v87-sector}, and Theorem
\ref{thm:v14v87-einstein}.  Item 5 ensures that these limits belong
to one physical theory rather than unrelated operator and metric
limits.
\end{proof}

\begin{assessment}[What this theorem does not construct]
\label{ass:v14v87-scope}
Item 1 of Theorem~\ref{thm:v14v87-cofinal} is a major independent
construction: tightness, cylinder consistency, reflection
positivity, the OS null quotient, and time translations must all
survive.  Norm-resolvent convergence cannot be asserted before the
operators have been embedded into a common limiting Hilbert space.
Likewise, the theorem preserves a relativistic field algebra only
because item 5 assumes its compatible convergence.
\end{assessment}

\begin{acquisition}[Cofinal spectral--metric card]
\label{acq:v14v87}
The next source extraction must provide:
\begin{enumerate}
\item explicit embeddings of the finite OS Hilbert spaces;
\item a norm bound for the resolvent difference in \((87.2)\);
\item a uniform finite physical gap from \((86.12)\);
\item convergence or a Riesz-contour certificate for the vacuum
      projections;
\item renormalized absolute sector energies with a uniform winner
      margin;
\item a compact reduced metric set \(K\);
\item uniform \(C^2(K)\) bounds and convergence of the complete
      effective actions;
\item a uniform reduced Hessian floor.
\end{enumerate}
\end{acquisition}

\begin{assessment}[Current source boundary]
\label{ass:v14v87-source}
The attached sources do not provide a common OS Hilbert embedding,
a norm-resolvent estimate, a uniformly isolated sector vacuum, or
uniform \(C^2\) convergence of the complete effective action.
Therefore Theorem~\ref{thm:v14v87-cofinal} remains conditional.
\end{assessment}

\begin{decision}[V88 reflected measure and OS-space limit]
\label{dec:v14v87-next}
Go upstream to item 1.  Formulate a projective family of finite
reflected probability laws, prove that cylinder-level reflection
positivity survives weak/projective convergence, construct the
pre-Hilbert OS quotient on cylinder observables, and state the
closability and time-translation conditions needed before
norm-resolvent convergence can be posed.  Include a counterexample
showing that convergence of normalized finite partition functions
alone determines no limiting measure.
\end{decision}

\begin{firewall}[Claim boundary after V87]
\label{fire:v14v87}
The cofinal theorem preserves, but does not construct, the common
measure, Hilbert space, field algebra, or covariance.  Its spectral
and metric convergence hypotheses have not been verified for the
attached programme.  No full SM--Einstein continuum is claimed.
\end{firewall}

\section*{Version V87 append-only record}
\addcontentsline{toc}{section}{Version V87 append-only record}
The complete V86 source was retained before this continuation.  It
had \(1020671\) bytes, \(30711\) lines, and SHA--256
\[
 \texttt{66DADCE5D3042B938332DF4961ED38D63EFA52141D0D4B4D60E1FFF0FF41D42A}.
\]
V87 proved norm-resolvent preservation of the unique vacuum and gap,
uniform \(C^2\) preservation of the reduced Einstein solution, and
sharp counterexamples to weaker convergence modes.
\section{Fourteenth-cycle V88: projective reflected laws and the limiting OS space}
\label{sec:v14v88}

V87 assumed a common limiting Hilbert space.  This continuation
constructs the cylinder-level candidate from compatible finite laws
and shows exactly what remains before the finite Hamiltonians can be
placed on it.

\subsection{Projective probability law}

Let
\[
 X_1\xleftarrow{r_{21}}X_2\xleftarrow{r_{32}}X_3
 \xleftarrow{}\cdots                                    \tag{88.1}
\]
be nonempty finite configuration spaces with surjective restriction
maps satisfying
\(r_{kn}=r_{mn}\circ r_{km}\) for \(k\ge m\ge n\).
Let \(\mu_n\) be probability laws satisfying exact consistency
\[
 (r_{mn})_\#\mu_m=\mu_n\qquad(m\ge n).                   \tag{88.2}
\]
Define the projective configuration space
\[
 X_\infty
 :=
 \{(x_n)_n:r_{mn}(x_m)=x_n\text{ for }m\ge n\}.          \tag{88.3}
\]

\begin{theorem}[Finite projective-law construction]
\label{thm:v14v88-projective}
There is a unique probability law \(\mu_\infty\) on the sigma
algebra generated by the coordinate cylinders such that
\[
 (\pi_n)_\#\mu_\infty=\mu_n                              \tag{88.4}
\]
for every \(n\).  For every cylinder function
\(F=f\circ\pi_n\),
\[
 \int_{X_\infty}F\,d\mu_\infty
 =
 \sum_{x\in X_n}f(x)\mu_n(x).                            \tag{88.5}
\]
\end{theorem}

\begin{proof}
Give every finite \(X_n\) the discrete topology.  The product
\(\prod_nX_n\) is compact, and \(X_\infty\) is a closed subset.
Surjectivity makes it nonempty by the finite-intersection property.
Define the value of a cylinder event
\(\pi_n^{-1}(A)\) to be \(\mu_n(A)\).  If the same event is
represented at a finer level, \((88.2)\) gives the same value.
Finite additivity follows at one common refinement level.

The resulting positive normalized functional on cylinder functions
is bounded by the uniform norm.  Cylinder functions form a
self-adjoint unital algebra separating points of the compact space
\(X_\infty\), hence are uniformly dense in
\(C(X_\infty)\).  The functional extends uniquely and positively to
\(C(X_\infty)\), and the Riesz representation theorem gives
\(\mu_\infty\).  Equation \((88.4)\) proves uniqueness on the
cylinder sigma algebra.
\end{proof}

Exact consistency can be obtained after a subsequence when only
finite-marginal convergence is known.

\begin{proposition}[Diagonal marginal handoff]
\label{prop:v14v88-diagonal}
Suppose \(\widetilde\mu_m\) is a law on \(X_m\), without assuming
\((88.2)\).  If for every fixed \(n\),
\[
 (r_{mn})_\#\widetilde\mu_m\longrightarrow\mu_n
 \quad(m\to\infty),                                     \tag{88.6}
\]
then the limiting laws satisfy \((88.2)\) and define
\(\mu_\infty\) by Theorem~\ref{thm:v14v88-projective}.
If convergence has not been established, compactness of every finite
probability simplex permits a diagonal subsequence on which all
fixed marginals converge.
\end{proposition}

\begin{proof}
For \(k\ge n\), continuity of the finite pushforward map and
\(r_{mn}=r_{kn}\circ r_{mk}\) give
\[
 (r_{kn})_\#\mu_k
 =
 \lim_{m\to\infty}
 (r_{kn})_\#(r_{mk})_\#\widetilde\mu_m
 =
 \lim_{m\to\infty}(r_{mn})_\#\widetilde\mu_m
 =\mu_n.
\]
For the final statement, successively extract convergent
subsequences for the first, second, and later marginals, then take
the diagonal subsequence.
\end{proof}

\subsection{Reflection positivity survives on cylinders}

Assume reflections
\(\vartheta_n:X_n\to X_n\) commute with restriction, and let
\({\cal A}_{n,+}\) be a positive-time function algebra whose pullback
to finer levels remains positive-time.  Define the anti-linear
reflection
\[
 (\Theta_nf)(x)=\overline{f(\vartheta_nx)}.              \tag{88.7}
\]
Suppose each finite law is reflection positive:
\[
 \int_{X_n}(\Theta_nf)f\,d\mu_n\ge0
 \qquad(f\in{\cal A}_{n,+}).                             \tag{88.8}
\]

\begin{theorem}[Projective preservation of reflection positivity]
\label{thm:v14v88-rp}
Let \({\cal A}_{\infty,+}\) be the union of the pulled-back
positive-time cylinder algebras.  Then
\[
 \langle F,G\rangle_{\rm OS}
 :=
 \int_{X_\infty}(\Theta F)G\,d\mu_\infty                \tag{88.9}
\]
is a positive semidefinite sesquilinear form on
\({\cal A}_{\infty,+}\).  If \(F=f\circ\pi_n\), its OS norm equals
the finite norm:
\[
 \|F\|_{\rm OS}^2
 =
 \int_{X_n}(\Theta_nf)f\,d\mu_n.                         \tag{88.10}
\]
\end{theorem}

\begin{proof}
Any finite family of cylinder functions has a common refinement
level \(m\).  Reflection compatibility and \((88.5)\) identify its
OS Gram matrix with the finite matrix computed using \(\mu_m\).
Equation \((88.8)\), applied to arbitrary linear combinations, makes
that matrix positive semidefinite.  This proves \((88.9)\);
\((88.10)\) is the one-function case.
\end{proof}

Let
\[
 {\cal N}_{\rm OS}
 :=
 \{F:\|F\|_{\rm OS}=0\}.                                 \tag{88.11}
\]
Cauchy--Schwarz for a positive semidefinite sesquilinear form shows
that every element of \({\cal N}_{\rm OS}\) is orthogonal to every
cylinder observable.

\begin{corollary}[Cylinder OS Hilbert space]
\label{cor:v14v88-hilbert}
The quotient
\[
 {\cal D}_{\rm OS}
 =
 {\cal A}_{\infty,+}/{\cal N}_{\rm OS}                  \tag{88.12}
\]
is a pre-Hilbert space.  Its completion
\({\cal H}_{\rm OS}\) is the limiting cylinder OS Hilbert space.
Under exact consistency, every finite OS quotient maps isometrically
to \({\cal H}_{\rm OS}\), and the union of their cylinder images is
dense.
\end{corollary}

\begin{proof}
The null space is a linear radical by Cauchy--Schwarz, so
\((88.9)\) descends to a positive definite form on the quotient.
Equation \((88.10)\) proves that pullback preserves norms and null
vectors.  Density holds by the definition of the completion of the
cylinder quotient.
\end{proof}

\subsection{Asymptotic laws and null-space stability}

For merely convergent finite marginals, cylinder Gram matrices
converge, but finite quotient dimensions need not stabilize.  Fix a
finite cylinder subspace with basis \(F_1,\ldots,F_d\), and let
\[
 G_n=[\langle F_i,F_j\rangle_{{\rm OS},n}],\qquad
 G=[\langle F_i,F_j\rangle_{\rm OS}].                    \tag{88.13}
\]

\begin{proposition}[Finite-cylinder null stability]
\label{prop:v14v88-null}
Assume
\[
 \|G_n-G\|\to0,\qquad
 \sigma(G_n)\subset\{0\}\cup[\eta,\infty)                \tag{88.14}
\]
for one \(\eta>0\).  Then the orthogonal projections onto
\(\ker G_n\) converge in norm to the projection onto \(\ker G\);
their ranks are eventually constant.  The positive square-root
whitenings on the retained ranges provide uniformly conditioned
finite-cylinder Hilbert identifications.
\end{proposition}

\begin{proof}
Use a circle about zero of radius less than \(\eta/2\).
Norm convergence and the Riesz projection argument of V87 give norm
convergence of the enclosed spectral projections.  Projections at
distance less than one have equal finite rank.  On the orthogonal
complements, every positive eigenvalue is at least \(\eta/2\) for
large \(n\), so \(G_n^{1/2}\) and its inverse there have uniform
bounds and converge by finite-dimensional functional calculus.
\end{proof}

\begin{countertheorem}[OS quotient dimension can collapse]
\label{ctr:v14v88-null}
The positive Gram matrices
\[
 G_n=\begin{pmatrix}1&0\\0&1/n\end{pmatrix}
 \longrightarrow
 G=\begin{pmatrix}1&0\\0&0\end{pmatrix}                 \tag{88.15}
\]
have zero-dimensional finite null spaces but a one-dimensional
limiting null space.  Thus convergence of cylinder expectations
alone does not supply the Hilbert identifications needed in
\((87.2)\).
\end{countertheorem}

\begin{proof}
The spectra are displayed in \((88.15)\); the missing uniform
positive Gram floor is \(1/n\).
\end{proof}

\subsection{Time translation on the quotient}

Let \(\tau_t\) act on positive-time cylinder observables for
\(t\ge0\).  Assume:
\[
 \tau_{t+u}=\tau_t\tau_u,\qquad \tau_0=I,                 \tag{88.16}
\]
\[
 \langle F,\tau_tG\rangle_{\rm OS}
 =
 \langle\tau_tF,G\rangle_{\rm OS},                       \tag{88.17}
\]
and
\[
 \|\tau_tF\|_{\rm OS}\le\|F\|_{\rm OS}.                  \tag{88.18}
\]
Assume also OS-norm continuity at zero on the cylinder domain:
\[
 \|\tau_tF-F\|_{\rm OS}\longrightarrow0
 \quad(t\downarrow0).                                   \tag{88.19}
\]

\begin{theorem}[OS time semigroup and Hamiltonian]
\label{thm:v14v88-time}
Equations \((88.16)\)--\((88.19)\) make
\[
 T(t)[F]=[\tau_tF]                                       \tag{88.20}
\]
a well-defined strongly continuous self-adjoint contraction
semigroup on \({\cal H}_{\rm OS}\).  There is a unique nonnegative
self-adjoint operator \(H\) such that
\[
 T(t)=e^{-tH}.                                           \tag{88.21}
\]
For a discrete time step \(a\), a single transfer \(T(a)\) determines
a finite exact logarithmic Hamiltonian on its injective range, but
strong continuity of a full semigroup is an additional continuum
row.
\end{theorem}

\begin{proof}
The contraction \((88.18)\) maps OS-null vectors to OS-null vectors,
so \((88.20)\) is well defined and extends by continuity to the
completion.  Equations \((88.16)\)--\((88.18)\) give a self-adjoint
contraction semigroup there.  Continuity \((88.19)\) on a dense
domain, together with uniform boundedness, gives strong continuity
on the completion.  The spectral theorem for strongly continuous
self-adjoint contraction semigroups gives the unique generator
\(H\ge0\).
\end{proof}

The contraction, symmetry, and continuity hypotheses must be derived
from the common Euclidean law and its time gluing.  Reflection
positivity by itself supplies the inner product, not all three.

\subsection{Graded cylinder functional}

Fermionic variables need no fictitious positive Grassmann measure.
Let \({\mathfrak A}_n\) be finite-dimensional graded cylinder
algebras with compatible inclusions and normalized linear
functionals \(\omega_n\).  If the functionals are consistent and
\[
 \omega_n((\Theta F)F)\ge0
 \quad(F\in{\mathfrak A}_{n,+}),                         \tag{88.22}
\]
then the same common-refinement proof defines a positive OS form on
the algebraic direct limit.  V84 supplies a sufficient finite
criterion for \((88.22)\) in every occupation sector.  The
determinant-line orientations and symmetry intertwiners must commute
with the inclusions; otherwise the functionals do not define one
graded projective system.

\subsection{Partition functions do not determine a measure}

\begin{countertheorem}[Normalized scalar data have no marginal
content]
\label{ctr:v14v88-partition}
On the fixed two-point space \(X=\{0,1\}\), let
\[
 \mu_{2k}=(3/4,1/4),\qquad
 \mu_{2k+1}=(1/4,3/4).                                  \tag{88.23}
\]
Every normalized partition function equals one, and an unnormalized
representation can be chosen with the same partition function at
every \(n\), but the laws have no limit.
\end{countertheorem}

\begin{proof}
The mass of the cylinder event \(\{0\}\) alternates between
\(3/4\) and \(1/4\).  Scalar normalization records only the sum of
the two weights and cannot recover this marginal.
\end{proof}

The same example can be tensored with any fixed reflection-positive
system.  Convergence of free-energy scalars therefore does not imply
convergence of Schwinger functions, OS Gram matrices, or vacuum
states.

\subsection{Projective acquisition ledger}

\begin{acquisition}[Reflected cylinder card]
\label{acq:v14v88}
For the common regulated programme, record:
\begin{enumerate}
\item finite configuration or graded cylinder algebras and all
      restriction/inclusion maps;
\item normalized finite laws or functionals and every fixed-cylinder
      marginal;
\item exact consistency, or quantitative convergence as in
      \((88.6)\);
\item compatible reflections and positive-time subalgebras;
\item finite OS Gram matrices, their null projections, and retained
      Gram floors;
\item compatible time translations, contraction, symmetry, and
      continuity estimates;
\item determinant-line orientations and reduced symmetry
      intertwiners under refinement;
\item a proof that the cylinder limit carries the intended local
      gauge-invariant SM observables and metric dependence.
\end{enumerate}
Only after these rows pass does
Theorem~\ref{thm:v14v88-time} supply the common Hilbert space and
Hamiltonian on which V87 can be applied.
\end{acquisition}

\begin{assessment}[Current source boundary]
\label{ass:v14v88-source}
The attached sources do not exhibit a projectively consistent common
family of reflected SM--Einstein laws or graded functionals.  They
also do not give uniform positive Gram floors controlling the finite
OS null spaces.  Thus the common Hilbert space assumed in V87
remains unconstructed.
\end{assessment}

\begin{decision}[V89 local observable net and Ward compatibility]
\label{dec:v14v88-next}
Assuming the cylinder functional, construct its local
gauge-invariant observable net.  Prove isotony and reflection
compatibility from the refinement maps, formulate the exact
finite Ward/BRST descent condition that makes gauge-exact insertions
OS-null, and state the Euclidean covariance and locality limits
needed before the OS Hamiltonian can be identified with a
relativistic Standard Model rather than an abstract gapped
semigroup.
\end{decision}

\begin{firewall}[Claim boundary after V88]
\label{fire:v14v88}
The projective and OS theorems are conditional constructions from
consistent finite laws.  No such common SM--Einstein law has been
extracted from the attached action, and no continuum covariance,
locality, anomaly cancellation, or physical field algebra has been
proved.
\end{firewall}

\section*{Version V88 append-only record}
\addcontentsline{toc}{section}{Version V88 append-only record}
The complete V87 source was retained before this continuation.  It
had \(1034324\) bytes, \(31119\) lines, and SHA--256
\[
 \texttt{AF9B9431C37E62330E368075FC492181BFC62F0DFF9C1BFEE565FDF31874689B}.
\]
V88 constructed the projective reflected cylinder law and OS space,
proved reflection positivity passes to the limit, and isolated
finite-cylinder null stability and time-semigroup conditions.
\section{Fourteenth-cycle V89: the physical local net and Ward--OS descent}
\label{sec:v14v89}

V88 constructs an OS Hilbert space from a consistent reflected
cylinder functional.  That space is not yet a Standard Model
Hilbert space: one must identify local gauge-invariant observables,
remove gauge-exact directions, and pass covariance and locality
through the same refinement system.

\subsection{The cylinder local net}

Let \({\cal R}\) be a directed family of bounded Euclidean regions.
For every \(O\in{\cal R}\), let \({\mathfrak A}(O)\) be the graded
cylinder algebra generated by fields supported in \(O\).  Assume the
refinement maps identify generators consistently.

\begin{proposition}[Algebraic isotony and graded locality]
\label{prop:v14v89-net}
If
\[
 O_1\subseteq O_2
 \quad\Longrightarrow\quad
 {\mathfrak A}(O_1)\subseteq{\mathfrak A}(O_2),           \tag{89.1}
\]
at every finite refinement, then the algebraic direct limit is an
isotone net.  If homogeneous elements with disjoint supports obey
\[
 AB=(-1)^{|A||B|}BA,                                     \tag{89.2}
\]
at every common refinement, then \((89.2)\) holds in the direct
limit.
\end{proposition}

\begin{proof}
Any finite family of cylinder elements appears at one common
refinement.  Inclusion and the graded commutation identity can
therefore be checked there, and the compatibility maps preserve both
relations.
\end{proof}

To represent a positive-time local element \(A\) by left
multiplication on the OS quotient, one also needs
\[
 \|AF\|_{\rm OS}\le C_A\|F\|_{\rm OS}
 \quad(F\in{\cal A}_{\infty,+}).                         \tag{89.3}
\]
Then
\[
 \pi_{\rm OS}(A)[F]=[AF]                                 \tag{89.4}
\]
is well defined and bounded.  Algebraic locality alone does not prove
\((89.3)\); the multiplier bound is a separate reconstruction row.

\subsection{Finite BRST complex}

Let \(s\) be an odd graded derivation on the cylinder algebra:
\[
 s(AB)=(sA)B+(-1)^{|A|}A(sB),\qquad s^2=0.              \tag{89.5}
\]
Assume it preserves refinement and the positive-time algebra.  Let
the reflection satisfy
\[
 \Theta s=\epsilon\,s\Theta,\qquad \epsilon\in\{-1,+1\}. \tag{89.6}
\]
The physical cylinder cohomology in a region is
\[
 {\mathfrak A}_{\rm phys}(O)
 =
 \frac{\ker(s:{\mathfrak A}(O)\to{\mathfrak A}(O))}
      {\operatorname{im}(s:{\mathfrak A}(O)\to
                              {\mathfrak A}(O))}.        \tag{89.7}
\]
Because \(s\) is a derivation, the product of closed elements is
closed and the product with an exact element is exact; hence
\((89.7)\) is an algebra.

The load-bearing identity is the finite Ward condition
\[
 \omega(sX)=0                                            \tag{89.8}
\]
for every cylinder element \(X\) on which the reflected products
below are defined.

\begin{theorem}[Ward identity sends BRST-exact states to the OS null
space]
\label{thm:v14v89-ward}
Assume reflection positivity, \((89.5)\)--\((89.8)\), and let
\(B,sB\) belong to the positive-time algebra.  Then
\[
 \|sB\|_{\rm OS}^2=0.                                    \tag{89.9}
\]
Consequently the map
\[
 [F]_{\rm BRST}\longmapsto[F]_{\rm OS}
 \quad(sF=0)                                             \tag{89.10}
\]
is well defined from positive-time BRST cohomology into the OS
Hilbert space, modulo any additional OS-null closed classes.
\end{theorem}

\begin{proof}
Using \((89.6)\) and the graded Leibniz rule,
\[
 \begin{split}
 \|sB\|_{\rm OS}^2
 &=\omega((\Theta sB)sB)\\
 &=\epsilon\,\omega((s\Theta B)sB)\\
 &=\epsilon\,\omega\!\left(
     s\bigl((\Theta B)sB\bigr)
   \right),
                                                               \tag{89.11}
 \end{split}
\]
because the second Leibniz term contains \(s^2B=0\).
Equation \((89.8)\) makes \((89.11)\) zero.  Reflection positivity
and Cauchy--Schwarz then make \(sB\) orthogonal to every OS vector,
so changing a closed representative by \(sB\) does not change its OS
class.
\end{proof}

\begin{corollary}[Local physical isotony]
\label{cor:v14v89-physical-net}
If the refinement and region inclusions commute with \(s\), then the
cohomology algebras \({\mathfrak A}_{\rm phys}(O)\) form an isotone
net, and Theorem~\ref{thm:v14v89-ward} represents that net on the
physical OS quotient whenever the multiplier bounds \((89.3)\) hold.
\end{corollary}

\begin{proof}
Inclusions map closed elements to closed elements and exact elements
to exact elements, so they induce cohomology maps.  The OS
representation is independent of the representative by
Theorem~\ref{thm:v14v89-ward}.
\end{proof}

\subsection{Nilpotency without the Ward identity is insufficient}

\begin{countertheorem}[An exact state need not be OS-null]
\label{ctr:v14v89-ward}
Let
\[
 {\mathfrak A}=\mathbb C[\varepsilon]/(\varepsilon^2),
\]
with \(\varepsilon\) odd, reflection given by conjugating
coefficients and fixing \(\varepsilon\), and
\[
 \omega(a+b\varepsilon)=a.
\]
Define the odd derivation
\[
 s(1)=0,\qquad s(\varepsilon)=1.                         \tag{89.12}
\]
Then \(s^2=0\), the OS form is positive semidefinite, but
\[
 \omega(s\varepsilon)=1,\qquad
 \|s\varepsilon\|_{\rm OS}^2=\|1\|_{\rm OS}^2=1.         \tag{89.13}
\]
\end{countertheorem}

\begin{proof}
For \(F=a+b\varepsilon\),
\[
 \omega((\Theta F)F)=|a|^2\ge0.
\]
The graded Leibniz rule is compatible with
\(\varepsilon^2=0\) because
\(s(\varepsilon^2)=\varepsilon-\varepsilon=0\).
Equation \((89.13)\) is direct.  Thus the Ward identity, rather than
nilpotency alone, is what forces exact directions into the OS null
space.
\end{proof}

A finite anomaly defect is therefore an OS defect.  It cannot be
postponed until after taking the quotient, because the quotient itself
depends on \((89.9)\).

\subsection{Compact gauge averaging}

When a finite or compact gauge group \(G\) acts by algebra
automorphisms \(\alpha_g\) satisfying
\[
 \omega\circ\alpha_g=\omega,\qquad
 \Theta\alpha_g=\alpha_g\Theta,                          \tag{89.14}
\]
there is a direct invariant-sector construction.

\begin{theorem}[Gauge-average projection]
\label{thm:v14v89-haar}
The maps
\[
 U_g[F]=[\alpha_gF]                                      \tag{89.15}
\]
form a unitary representation on the OS Hilbert space, and
\[
 P_G=\int_GU_g\,dg                                       \tag{89.16}
\]
is the orthogonal projection onto the gauge-invariant OS subspace.
For a finite group, the integral is the normalized sum.
\end{theorem}

\begin{proof}
Equation \((89.14)\) preserves the OS inner product and its null
space, so \((89.15)\) is unitary.  Haar invariance gives
\(P_G^*=P_G\), \(P_G^2=P_G\), and
\(U_hP_G=P_G\) for every \(h\).  Its range is therefore exactly the
invariant subspace.
\end{proof}

Gauge averaging and BRST cohomology are not interchangeable without
a comparison theorem.  The finite card must state which construction
is used and prove its compatibility with time translation, sector
decomposition, and refinement.

\subsection{Euclidean covariance in the cylinder limit}

Let a Euclidean transformation \(g\) act on supports and by algebra
automorphisms \(\alpha_g\).  Suppose that for every fixed cylinder
observable \(F\), finite representatives satisfy
\[
 \omega_n(\alpha_{n,g}F_n)-\omega_n(F_n)\longrightarrow0, \tag{89.17}
\]
and the transformed representatives are compatible with the
limiting cylinder maps.

\begin{proposition}[Covariance passes through cylinder expectations]
\label{prop:v14v89-covariance}
Under convergence of the relevant fixed-cylinder marginals,
\[
 \omega_\infty(\alpha_gF)=\omega_\infty(F).              \tag{89.18}
\]
If \(g\) preserves the positive half-space and commutes appropriately
with reflection, it induces an isometry on the OS quotient.
\end{proposition}

\begin{proof}
Pass to the limit in \((89.17)\) using fixed-cylinder expectation
convergence.  The same calculation with
\((\Theta F)G\) shows preservation of the OS inner product.
\end{proof}

Exact lattice symmetry at every cutoff usually covers only a
discrete subgroup.  Recovery of continuous translations and
rotations requires estimates uniform as the mesh tends to zero.
The reduced-intertwiner rule from V85 applies: a geometric symmetry
before gauge fixing must be accompanied by its action on the
constraint-solving chart, determinant line, and OS null space.

\subsection{From Euclidean graded locality to a relativistic net}

Equation \((89.2)\) provides Euclidean support locality in the
cylinder algebra.  To identify the OS reconstruction with a
relativistic local theory one further needs:
\begin{enumerate}
\item strong continuity of the reconstructed Euclidean group action;
\item the spectrum condition for the continued translation
      generators;
\item analytic continuation of a sufficiently rich set of
      Schwinger functions;
\item permutation or graded permutation symmetry and regularity at
      noncoincident points;
\item proof that Euclidean disjoint-support relations continue to
      spacelike graded commutativity;
\item a local stress tensor or metric-response distribution
      compatible with the Ward identities.
\end{enumerate}
These are reconstruction hypotheses, not consequences of a
Hamiltonian gap.

\begin{proposition}[Gap and locality are logically independent]
\label{prop:v14v89-independent}
A positive self-adjoint semigroup with a unique gapped vacuum need
not carry any specified local net, and an isotone graded-local
cylinder net need not have a gapped time generator.
\end{proposition}

\begin{proof}
Any finite-dimensional positive matrix semigroup with a simple
largest transfer eigenvalue supplies the first object without
spatial support data.  Conversely, a free massless Gaussian
cylinder net has the standard local algebraic support relations but
its one-particle energies approach zero with momentum.  Thus neither
set of axioms contains the other.
\end{proof}

\subsection{Ward and local-net acquisition ledger}

\begin{acquisition}[Physical observable card]
\label{acq:v14v89}
For each refinement and its limit, record:
\begin{enumerate}
\item the local graded cylinder algebras and support maps;
\item the constraint-solving and gauge action on every generator;
\item either the BRST differential \(s\) with \(s^2=0\), reflection
      compatibility, and the exact Ward identity, or the compact
      gauge-average projector;
\item all finite anomaly defects and determinant-line holonomies;
\item multiplier bounds \((89.3)\) for the reconstructed local
      operators;
\item time-translation invariance of the physical subspace;
\item quantitative restoration of continuous Euclidean covariance;
\item the Schwinger-function regularity and continuation estimates
      used to obtain spacelike graded locality;
\item the local metric-response insertion and its compatibility with
      the Einstein-source jets.
\end{enumerate}
\end{acquisition}

\begin{assessment}[Current source boundary]
\label{ass:v14v89-source}
The attached sources contain formal gauge, BRST, and Ward structures,
but do not provide one projectively consistent graded functional for
which \((89.8)\), determinant-line triviality, multiplier bounds,
continuous covariance restoration, and OS analytic continuation are
all proved.  Therefore no physical continuum observable net has yet
been constructed.
\end{assessment}

\begin{decision}[V90 compulsory audit]
\label{dec:v14v89-next}
Inspect the newest Yang--Mills and Navier--Stokes companion notes.
Look specifically for an exact symmetry intertwiner, whole-block
factorization, constrained multiplier estimate, or cofinal
compactness mechanism that changes the local-net or Ward acquisition
order.  Then continue in V91 with the stress-tensor insertion and
diffeomorphism Ward identity needed to couple the limiting
observable net to Einstein geometry.
\end{decision}

\begin{firewall}[Claim boundary after V89]
\label{fire:v14v89}
The Ward--OS descent theorem is conditional on an exact finite Ward
identity for the same reflected functional.  The projective
functional, anomaly-free chiral BRST complex, multiplier bounds,
continuous covariance, analytic continuation, and local stress
tensor remain unconstructed.  No full SM--Einstein continuum is
claimed.
\end{firewall}

\section*{Version V89 append-only record}
\addcontentsline{toc}{section}{Version V89 append-only record}
The complete V88 source was retained before this continuation.  It
had \(1048472\) bytes, \(31513\) lines, and SHA--256
\[
 \texttt{892E87B5FB31949F9A275BC04107DA846531B199851D0E7DD55ABD2DA852B534}.
\]
V89 constructed the algebraic physical local net, proved exact
Ward--BRST descent into the OS null space, and separated covariance
and relativistic locality from the spectral-gap row.
\section{Fourteenth-cycle V90: compulsory audit and constraint-locus hulls}
\label{sec:v14v90}

The newest companion files were inspected read-only at the required
five-version checkpoint.

\begin{record}[V90 companion checkpoint]
\label{rec:v14v90-companions}
The Yang--Mills checkpoint is
\begin{center}
\path{Renewal_Yang_Mills_Mass_Gap_Research_Notes_V28.tex},
\end{center}
with \(225902\) bytes, \(6151\) lines, and SHA--256
\[
 \texttt{022E47EA7AE296E6141C589F1524F9F20D647D430CB45009A57438E0B224E6EB}.
                                                               \tag{90.1}
\]
The Navier--Stokes checkpoint remains
\begin{center}
\path{Renewal_NS_Stability_Research_Notes_V26.tex},
\end{center}
with \(278283\) bytes, \(5319\) lines, and SHA--256
\[
 \texttt{82C887F8C2C69E4F28FAD7C3DC8DE5B9099CB5A4BF431EBA1FFF3E91935978AD}.
                                                               \tag{90.2}
\]
\end{record}

Yang--Mills V28 replaces a rectangular relaxation in
\((\cos u,\sin u)\) by an exact rational polygon containing the
physical circular arc.  Because the matrix pencil is separately
affine in each coordinate pair, positivity at the product vertices
proves positivity on the curved physical locus.  This doubles its
certified transverse radius without changing the splitting.  The
note correctly retains the failure of full-torus coverage and all
cofinal rows.

\begin{assessment}[V28 import]
\label{ass:v14v90-ym}
The rational-arc result supplies no stress tensor, Ward identity,
fermion orientation, common measure, or compact OS screen.  Its
specific Wilson floors are not imported.

The transferable method is sharper: when admissible metric, gauge,
or phase variables lie on an algebraic or trigonometric constraint
locus, first construct a rational outer polytope adapted to that
locus and only then apply a whole-matrix vertex certificate.
An oversized coordinate box may contain nonphysical negative
corners and produce a false rejection.  Every hull must include an
exact proof of containment of the physical locus.
\end{assessment}

Navier--Stokes V26 is unchanged and adds no local-net or
stress-response input.  Its constrained-input norm lesson is already
recorded at V80 and V85.

\begin{theorem}[V90 audit verdict]
\label{thm:v14v90-verdict}
Neither companion closes any hypothesis of V88--V89.  The V91
stress-response compiler should, however, preserve the physical
metric-variation locus before taking matrix hulls or operator norms,
and should distinguish a failure of an outer relaxation from a
failure on an admissible metric path.
\end{theorem}

\begin{proof}
Yang--Mills V28 is a finite matrix positivity theorem on one proper
momentum tube and explicitly leaves measure, OS, compactness, and
mass-gap questions open.  Navier--Stokes V26 is a deterministic
kernel-norm theorem.  Their conclusion types do not match the
projective, Ward, or metric-response rows.  The arc containment and
multi-affine convexity proof is independent of the Wilson
coefficients and therefore yields the stated certification rule.
\end{proof}

\begin{decision}[V91 stress insertion and diffeomorphism Ward row]
\label{dec:v14v90-next}
For a differentiable reflected cylinder functional
\(\omega_g\), define the local metric derivative as a stress-tensor
insertion.  Prove:
\begin{enumerate}
\item the derivative of normalized expectations, including the
      connected subtraction;
\item the second derivative and contact term;
\item the diffeomorphism Ward identity from change of variables and
      its anomaly/Jacobian defect;
\item descent of the stress insertion through BRST and OS quotients;
\item the exact relation between the integrated stress jets and the
      vacuum-energy derivatives in V83.
\end{enumerate}
Keep admissible metric variations on their physical constraint locus
and use a rational outer hull only when its containment is proved.
\end{decision}

\begin{firewall}[Claim boundary after V90]
\label{fire:v14v90}
The audit imports a certification method, not a physical theorem or
constant.  The common reflected functional, anomaly-free Ward
identity, local stress distribution, and cofinal Einstein coupling
remain unconstructed.
\end{firewall}

\section*{Version V90 append-only record}
\addcontentsline{toc}{section}{Version V90 append-only record}
The complete V89 source was retained before this continuation.  It
had \(1061171\) bytes, \(31861\) lines, and SHA--256
\[
 \texttt{384B10AD7E789B1F918A0F3B48D5A4AEF2172EC3FC5D6CCAF1EB8F447D38E2C8}.
\]
V90 performed the companion audit and imported only the
physical-constraint-locus rational-hull method.
\section{Fourteenth-cycle V91: stress insertions, diffeomorphism Ward defect, and vacuum jets}
\label{sec:v14v91}

V89 constructed the physical local-net requirements, while V83 and
V87 used derivatives of an absolute vacuum energy.  This
continuation connects the two.  At finite cutoff the connection is
an exact differentiation identity; its continuum use additionally
requires uniform integrability and renormalization.

\subsection{Normalized first and second metric derivatives}

Let \(Q\) be a finite configuration space, let \(\nu\) be a fixed
positive reference measure, and let the complete real effective
finite action be
\[
 {\cal I}_g(q).
\]
It includes the bosonic action, the positive fermionic determinant
scalar after V84, gauge-fixing Jacobians, and declared counterterms.
Define
\[
 Z(g)=\sum_{q\in Q}e^{-{\cal I}_g(q)}\nu(q),\qquad
 \langle F\rangle_g
 =
 Z(g)^{-1}\sum_qF_g(q)e^{-{\cal I}_g(q)}\nu(q).           \tag{91.1}
\]
For metric directions \(u,v\), put
\[
 X_u=D_g{\cal I}_g[u],\qquad
 X_{uv}=D_g^2{\cal I}_g[u,v].                            \tag{91.2}
\]
These are the integrated first stress insertion and its contact
term in the present sign convention.

\begin{theorem}[Normalized stress differentiation]
\label{thm:v14v91-diff}
For every differentiable observable \(F_g\),
\[
 D_g\langle F\rangle[u]
 =
 \langle D_gF[u]\rangle
 -
 \operatorname{Cov}(F,X_u),                             \tag{91.3}
\]
where
\[
 \operatorname{Cov}(A,B)
 =
 \langle AB\rangle-\langle A\rangle\langle B\rangle.
\]
For
\[
 W(g)=-\log Z(g),
\]
one has
\[
 DW[u]=\langle X_u\rangle,                               \tag{91.4}
\]
and
\[
 D^2W[u,v]
 =
 \langle X_{uv}\rangle
 -
 \operatorname{Cov}(X_u,X_v).                           \tag{91.5}
\]
Thus the second metric jet is the contact term minus the connected
stress covariance.
\end{theorem}

\begin{proof}
Differentiate the numerator and denominator in \((91.1)\).
The derivative of the weight is
\(-X_u e^{-{\cal I}_g}\), and the quotient rule gives
\((91.3)\).  Since
\[
 DZ[u]=-Z\langle X_u\rangle,
\]
equation \((91.4)\) follows.  Apply \((91.3)\) to
\(F=X_u\), including \(D_gX_u[v]=X_{uv}\), to obtain
\((91.5)\).
\end{proof}

\begin{assessment}[Contact terms are compulsory]
\label{ass:v14v91-contact}
The covariance term alone is not the Einstein susceptibility.
Metric dependence of the kinetic operator, measure, gauge fixing,
fermionic Schur block, and counterterms contributes to
\(X_{uv}\).  Omitting it changes both the sign and the value of the
finite Hessian.
\end{assessment}

If the reference measure also depends on \(g\), choose one fixed
dominating measure and include the logarithm of its density in
\({\cal I}_g\).  This makes the Jacobian contribution explicit rather
than silently dropping it.

\subsection{Finite diffeomorphism Ward identity}

Let \(V_\xi\) be the infinitesimal field-space change generated by a
discrete or finite-dimensional approximation to a vector field
\(\xi\).  Write
\[
 \delta_\xi F=V_\xi F,\qquad
 {\rm div}_\nu V_\xi
 =
 \left.\frac d{dt}\right|_{0}
 \log\frac{(\Phi_t)_\#\nu}{\nu},                         \tag{91.6}
\]
with the sign fixed by the following change-of-variables identity.
Define
\[
 A_\xi
 :=
 {\rm div}_\nu V_\xi-\delta_\xi{\cal I}_g.               \tag{91.7}
\]

\begin{lemma}[Finite change-of-variables Ward identity]
\label{lem:v14v91-change}
For every \(F\),
\[
 \langle\delta_\xi F\rangle
 +
 \langle FA_\xi\rangle=0,\qquad
 \langle A_\xi\rangle=0.                                \tag{91.8}
\]
\end{lemma}

\begin{proof}
Perform the bijective change of variables \(q\mapsto\Phi_tq\) in
the unnormalized finite sum or integral and differentiate at zero.
The derivative of \(F\) is \(\delta_\xi F\), while the density
derivative is the sum of the Jacobian divergence and
\(-\delta_\xi{\cal I}_g\).  This gives the first identity.  Set
\(F=1\) for the second.
\end{proof}

Let
\[
 X_\xi:=D_g{\cal I}_g[{\cal L}_\xi g]                    \tag{91.9}
\]
be the action response to the induced metric variation.  Define the
combined diffeomorphism defect
\[
 {\mathfrak a}_\xi
 :=
 X_\xi+A_\xi
 =
 D_g{\cal I}_g[{\cal L}_\xi g]
 +{\rm div}_\nu V_\xi-\delta_\xi{\cal I}_g.              \tag{91.10}
\]
It contains regulator noncovariance, a Jacobian anomaly, and any
orientation defect in the fermion determinant line.

\begin{theorem}[Normalized diffeomorphism Ward defect]
\label{thm:v14v91-diffeo}
For a metric-independent observable \(F\),
\[
 D_g\langle F\rangle[{\cal L}_\xi g]
 +
 \langle\delta_\xi F\rangle
 =
 -\operatorname{Cov}(F,{\mathfrak a}_\xi).              \tag{91.11}
\]
In particular, if
\[
 {\mathfrak a}_\xi=0                                     \tag{91.12}
\]
as a finite cylinder insertion, then the exact diffeomorphism Ward
identity holds:
\[
 D_g\langle F\rangle[{\cal L}_\xi g]
 +\langle\delta_\xi F\rangle=0.                          \tag{91.13}
\]
\end{theorem}

\begin{proof}
Equations \((91.7)\) and \((91.10)\) give
\(A_\xi=-X_\xi+{\mathfrak a}_\xi\).
The second identity in \((91.8)\) implies
\(\langle X_\xi\rangle=\langle{\mathfrak a}_\xi\rangle\).
Using \((91.3)\) with \(D_gF=0\),
\[
 D_g\langle F\rangle[{\cal L}_\xi g]
 =-\langle FX_\xi\rangle
   +\langle F\rangle\langle X_\xi\rangle.
\]
The first identity in \((91.8)\) gives
\[
 \langle\delta_\xi F\rangle
 =\langle FX_\xi\rangle-\langle F{\mathfrak a}_\xi\rangle.
\]
Adding and substituting the equality of the defect and response
means proves \((91.11)\).
\end{proof}

When
\[
 X_\xi
 =
 \frac12\int T^{\mu\nu}{\cal L}_\xi g_{\mu\nu},
                                                               \tag{91.14}
\]
integration by parts converts \((91.13)\), for invariant insertions,
into the distributional stress conservation law.  Equation
\((91.11)\) shows that the actual conclusion in the presence of a
regulator defect is an anomalous divergence with insertion
\({\mathfrak a}_\xi\).

\subsection{BRST and OS compatibility of the stress insertion}

Let \(s\) be the V89 BRST differential, independent of the metric
chart coordinates, and assume
\[
 \omega_g(sY)=0                                          \tag{91.15}
\]
for every \(g\) in the chart and every required reflected cylinder
product \(Y\).

\begin{proposition}[Differentiated Ward identity]
\label{prop:v14v91-brst}
If \(Y\) is metric independent, then
\[
 \operatorname{Cov}_g(sY,X_u)=0.                         \tag{91.16}
\]
More generally,
\[
 D_g\langle sY\rangle[u]
 =
 \langle s(D_gY[u])\rangle-\operatorname{Cov}(sY,X_u)=0.
                                                               \tag{91.17}
\]
For every positive-time \(B\), the BRST-exact vector \(sB\) remains
OS-null throughout the metric chart, and differentiating its zero
OS norm produces no physical stress matrix element.
\end{proposition}

\begin{proof}
Differentiate the normalized Ward identity
\(\langle sY\rangle_g=0\) and use
Theorem~\ref{thm:v14v91-diff}; commutation of \(s\) with the metric
derivative gives \((91.17)\), and metric independence gives
\((91.16)\).  Apply this to
\(Y=(\Theta B)sB\).  The V89 proof makes
\(\langle sY\rangle\) the OS norm of \(sB\), up to its fixed sign.
It vanishes for every \(g\), so its differentiated insertion also
vanishes on the physical quotient.
\end{proof}

This result controls BRST-exact null vectors.  A general
metric-dependent OS null space can still rotate or change dimension;
the Gram-floor and null-projection stability of V88 remain necessary.

\subsection{Long-cylinder transfer identity}

Let \({\mathbb T}(g)>0\) be a \(d\)-dimensional finite transfer with
eigenvalues
\[
 \lambda_0(g)>\lambda_1(g)\ge\cdots>0
\]
on a metric neighbourhood, and let
\[
 Z_N(g)=\operatorname{Tr}{\mathbb T}(g)^N,\qquad
 W_N(g)=-\log Z_N(g).                                    \tag{91.18}
\]
Assume a uniform ratio
\[
 \sup_g\frac{\lambda_1(g)}{\lambda_0(g)}\le\rho<1        \tag{91.19}
\]
and uniform \(C^2\) bounds on the eigenvalue data in that
neighbourhood.

\begin{theorem}[Integrated stress converges to the vacuum-energy
jets]
\label{thm:v14v91-cylinder}
With
\[
 E_0(g)=-a^{-1}\log\lambda_0(g),
\]
one has locally uniformly through two metric derivatives
\[
 \frac1{Na}W_N(g)\longrightarrow E_0(g).                 \tag{91.20}
\]
In particular,
\[
 \frac1{Na}DW_N[u]\longrightarrow DE_0[u],\qquad
 \frac1{Na}D^2W_N[u,v]\longrightarrow D^2E_0[u,v].       \tag{91.21}
\]
Combining \((91.4)\)--\((91.5)\) with \((91.21)\) identifies the
per-unit-time integrated stress mean and connected
stress/contact combination with the transfer-vacuum jets
\((83.21)\).
\end{theorem}

\begin{proof}
Factor
\[
 Z_N=\lambda_0^N(1+R_N),\qquad
 R_N=\sum_{j=1}^{d-1}\left(\frac{\lambda_j}{\lambda_0}\right)^N.
                                                               \tag{91.22}
\]
Then
\[
 |R_N|\le(d-1)\rho^N.
\]
The assumed \(C^2\) bounds and the uniform separation imply, after
differentiating \((91.22)\), bounds of the form
\[
 |D^kR_N|\le C_kN^k\rho^{N-k},
 \qquad k=0,1,2,                                        \tag{91.23}
\]
with a harmless adjustment of the constant for small \(N\).
Therefore
\[
 \frac1{Na}W_N
 =
 -\frac1a\log\lambda_0
 -\frac1{Na}\log(1+R_N),
\]
and the second term together with its first two derivatives tends to
zero locally uniformly.  This proves \((91.20)\)--\((91.21)\).
The finite differentiation identities then give the stress
interpretation.
\end{proof}

Separate normalization of \({\mathbb T}\) to top eigenvalue one
removes the first term in \((91.22)\) and therefore erases the
Einstein source.  The absolute determinant, bosonic normalization,
and counterterm scalar must remain in the long-cylinder ledger.

\subsection{Admissible metric loci}

The reduced metric variations entering \(X_u\) and \(X_{uv}\) must
satisfy the chosen gauge, signature, constraint, and clock
conditions.  Let \(g=g(z)\) parameterize that physical finite locus.
All Hessian and Ward tests are to be pulled back to the \(z\)
variables before taking a norm or a convex hull.

\begin{proposition}[Constraint-locus vertex certification]
\label{prop:v14v91-hull}
Suppose a pulled-back Hermitian response matrix \(M(z)\) is
multi-affine in variables \(p_j(z)\), and rational polytopes
\({\cal P}_j\) contain the exact physical images of the
\(p_j\)'s.  If
\[
 M(v_1,\ldots,v_k)\succeq mI
\]
at every product vertex
\(v_j\in\operatorname{Vert}({\cal P}_j)\), then the same floor holds
on the physical metric locus.
\end{proposition}

\begin{proof}
The exact containment puts each physical \(p_j\) in its polytope.
Iterated affine expansion writes \(M\) as a convex combination of
the product-vertex matrices.  The positive Loewner cone is convex.
\end{proof}

A negative vertex outside the physical locus rejects only the chosen
outer hull.  It is not a negative physical susceptibility, exactly
as in the Yang--Mills V28 audit.

\subsection{Stress-source acquisition ledger}

\begin{acquisition}[Local-to-vacuum stress card]
\label{acq:v14v91}
For every cutoff and metric chart, record:
\begin{enumerate}
\item the fixed dominating reference measure and the complete real
      action \({\cal I}_g\);
\item \(X_u\) and every contact insertion \(X_{uv}\), including
      fermion, gauge-fixing, measure, and counterterm terms;
\item the field-space diffeomorphism, its reference divergence, and
      the defect \({\mathfrak a}_\xi\);
\item exact BRST/Ward compatibility for reflected products;
\item the local stress distribution whose integral is \(X_u\);
\item the unnormalized transfer and the long-cylinder identity
      \((91.22)\);
\item a uniform leading-eigenvalue ratio and \(C^2\) derivative
      bounds;
\item the physical reduced metric parameterization and any rational
      outer polytopes used for certification;
\item uniform integrability and distributional convergence of the
      local stress and contact insertions;
\item agreement of the limiting local insertion with the derivative
      of the renormalized absolute vacuum energy.
\end{enumerate}
\end{acquisition}

\begin{assessment}[Current source boundary]
\label{ass:v14v91-source}
The attached sources do not provide one positive finite density
\(\exp(-{\cal I}_g)\) containing all chiral fermionic, gauge,
geometric, Jacobian, and counterterm factors, nor its local stress
and contact insertions.  The diffeomorphism defect and
long-cylinder \(C^2\) bounds are also absent.  Therefore the formal
Einstein source is not yet identified with a limiting OS stress
operator.
\end{assessment}

\begin{decision}[V92 distributional stress limit and Bianchi
compatibility]
\label{dec:v14v91-next}
Prove a uniform-integrability criterion that passes the local stress,
contact term, and diffeomorphism Ward identity from cylinder
functionals to distributions.  Then show that an anomaly-free
limiting Ward identity makes the renormalized source divergence-free
and hence compatible with the reduced Einstein Bianchi identity.
Include a concentration counterexample in which integrated vacuum
energy converges but local stress escapes into a delta singularity.
\end{decision}

\begin{firewall}[Claim boundary after V91]
\label{fire:v14v91}
The stress and Ward identities are exact finite formulas under a
real positive common density.  That density, its anomaly-free
diffeomorphism action, its local stress distribution, and the
uniform long-cylinder bounds have not been constructed for the
regulated Standard Model--Einstein programme.  No continuum
Einstein source is claimed.
\end{firewall}

\section*{Version V91 append-only record}
\addcontentsline{toc}{section}{Version V91 append-only record}
The complete V90 source was retained before this continuation.  It
had \(1065784\) bytes, \(31972\) lines, and SHA--256
\[
 \texttt{B16632A424A02F06CD29B803FA0E28A3B424F9CB8A85B225736CBF6FE227B147}.
\]
V91 proved the normalized stress and contact formulas, the finite
diffeomorphism Ward defect, BRST--OS stress descent, and the
long-cylinder equality with vacuum-energy jets.
\section{Fourteenth-cycle V92: distributional stress, Bianchi compatibility, and concentration}
\label{sec:v14v92}

V91 gives exact finite stress identities.  A local continuum source
requires more than convergence of its integral: the stress and
contact insertions must converge as distributions, their moments
must be uniformly integrable, and the diffeomorphism defect must
vanish in the same topology.

\subsection{Moment criterion for stress distributions}

Let \(M\) be a fixed smooth compact manifold or a fixed relatively
compact coordinate region.  Let
\({\cal D}_2(M)\) denote smooth compactly supported symmetric
two-tensor test fields.  Use the projective law of V88 to regard all
finite stress insertions as random variables on one probability
space.  For \(h,k\in{\cal D}_2(M)\), write
\[
 T_n(h),\qquad C_n(h,k)                                  \tag{92.1}
\]
for the first stress insertion and the second contact insertion.
They are linear and bilinear in their test arguments.

\begin{theorem}[Uniform-integrability passage of stress jets]
\label{thm:v14v92-ui}
Assume for some \(\delta>0\):
\begin{enumerate}
\item for every finite family of test tensors, the corresponding
      \(T_n(h)\)'s converge jointly in probability to variables
      \(T(h)\);
\item for every \(h\),
\[
 \sup_n\mathbb E|T_n(h)|^{2+\delta}<\infty;              \tag{92.2}
\]
\item \(C_n(h,k)\to C(h,k)\) in probability and
\[
 \sup_n\mathbb E|C_n(h,k)|^{1+\delta}<\infty;            \tag{92.3}
\]
\item there are a continuous test-function seminorm \(p\) and a
      constant \(K\) such that
\[
 \sup_n\|T_n(h)\|_{L^{2+\delta}}\le Kp(h),\qquad
 \sup_n\|C_n(h,k)\|_{L^{1+\delta}}
 \le Kp(h)p(k).                                          \tag{92.4}
\]
\end{enumerate}
Then
\[
 \mathbb ET_n(h)\to\mathbb ET(h),                        \tag{92.5}
\]
\[
 \operatorname{Cov}(T_n(h),T_n(k))
 \to
 \operatorname{Cov}(T(h),T(k)),                         \tag{92.6}
\]
and
\[
 \mathbb EC_n(h,k)\to\mathbb EC(h,k).                   \tag{92.7}
\]
The limits
\[
 \tau(h):=\mathbb ET(h),\qquad
 {\cal K}(h,k)
 :=
 \mathbb EC(h,k)-\operatorname{Cov}(T(h),T(k))           \tag{92.8}
\]
are respectively a distribution and a continuous bilinear
distributional second jet.
\end{theorem}

\begin{proof}
An \(L^{1+\delta}\) bound implies uniform integrability.  Convergence
in probability plus uniform integrability gives convergence in
\(L^1\), proving \((92.5)\) and \((92.7)\).  Hölder's inequality and
\((92.2)\) make the products
\(T_n(h)T_n(k)\) uniformly integrable.  Joint convergence in
probability and another truncation argument therefore give
convergence of their expectations; together with \((92.5)\), this
proves \((92.6)\).

Linearity and bilinearity pass to the limits.  The bounds in
\((92.4)\) pass to the limit and give continuity in the
test-function topology.  H\"older's inequality bounds the limiting
covariance by a constant times \(p(h)p(k)\).  Equation \((92.8)\)
therefore gives the stated distributional objects.
\end{proof}

The theorem deliberately requires more integrability for the first
stress than for the contact term because the susceptibility contains
a product of two first insertions.

\subsection{Anomaly-free diffeomorphism limit}

Let \(\xi\) be a smooth compactly supported vector field.  At cutoff
\(n\), let
\[
 X_{n,\xi}=T_n({\cal L}_\xi g),\qquad
 {\mathfrak a}_{n,\xi}
\]
be the metric response and the V91 combined diffeomorphism defect.
The finite change-of-variables identity gives
\[
 D W_n(g)[{\cal L}_\xi g]
 =
 \mathbb E X_{n,\xi}
 =
 \mathbb E{\mathfrak a}_{n,\xi}.                         \tag{92.9}
\]

\begin{theorem}[Ward limit implies distributional conservation]
\label{thm:v14v92-conservation}
Assume Theorem~\ref{thm:v14v92-ui} and
\[
 |\mathbb E{\mathfrak a}_{n,\xi}|
 \le\varepsilon_n q(\xi),\qquad
 \varepsilon_n\to0,                                     \tag{92.10}
\]
for a continuous vector-field seminorm \(q\).  Then
\[
 \tau({\cal L}_\xi g)=0                                  \tag{92.11}
\]
for every compactly supported \(\xi\).  Equivalently, if
\[
 \tau(h)=\frac12\int_M\tau^{\mu\nu}h_{\mu\nu},
\]
then
\[
 \nabla_\mu\tau^{\mu\nu}=0                               \tag{92.12}
\]
in the distributional sense.

If, more generally, the defect expectations converge to a
distribution \({\mathfrak a}(\xi)\), then
\[
 \tau({\cal L}_\xi g)={\mathfrak a}(\xi),                \tag{92.13}
\]
which is the anomalous divergence equation.
\end{theorem}

\begin{proof}
Use \((92.9)\), pass to the stress limit by \((92.5)\), and apply
\((92.10)\) to obtain \((92.11)\).  Since
\[
 ({\cal L}_\xi g)_{\mu\nu}
 =\nabla_\mu\xi_\nu+\nabla_\nu\xi_\mu,
\]
distributional integration by parts identifies \((92.11)\) with
\((92.12)\).  The same passage with a nonzero defect limit proves
\((92.13)\).
\end{proof}

For Ward identities with local insertions \(F\), the required
condition is stronger:
\[
 \operatorname{Cov}(F,{\mathfrak a}_{n,\xi})\to0.        \tag{92.14}
\]
Vanishing of only the mean defect proves vacuum stress conservation,
not the full local diffeomorphism Ward algebra.

\subsection{Bianchi compatibility and the reduced Hessian}

Let
\[
 \Gamma(g)=S_{\rm grav}(g)+C(g)+W_{\rm ren}(g),           \tag{92.15}
\]
where the gravitational and counterterm parts are
diffeomorphism-invariant and the matter source satisfies
\((92.11)\).  Then
\[
 D\Gamma(g)[{\cal L}_\xi g]=0.                           \tag{92.16}
\]

\begin{theorem}[Bianchi compatibility of the limiting source]
\label{thm:v14v92-bianchi}
The Euler--Lagrange tensor of \((92.15)\) is distributionally
divergence-free.  At a critical metric \(g_*\), its Hessian
\(L=D^2\Gamma(g_*)\) annihilates infinitesimal diffeomorphism
directions:
\[
 L({\cal L}_\xi g_*,h)=0
 \quad\hbox{for every }h.                                \tag{92.17}
\]
Consequently a strictly positive Einstein Hessian is possible only
on a gauge-fixed slice or quotient.  For a self-adjoint linearized
equation
\[
 Lh=f,                                                   \tag{92.18}
\]
the Ward condition that \(f\) annihilate
\({\cal L}_\xi g_*\) is the finite-codimension solvability condition
against the gauge kernel.
\end{theorem}

\begin{proof}
Equation \((92.16)\) for every compactly supported \(\xi\) is the
weak Noether identity, hence the Euler--Lagrange tensor has zero
covariant divergence.  Differentiate \((92.16)\) in an arbitrary
metric direction \(h\):
\[
 D^2\Gamma(g_*)[h,{\cal L}_\xi g_*]
 +
 D\Gamma(g_*)[D_h({\cal L}_\xi g)]=0.                   \tag{92.19}
\]
The second term vanishes because \(g_*\) is critical, proving
\((92.17)\).

For self-adjoint \(L\), the range is orthogonal to its kernel in the
finite reduced approximation, and the closed-range analogue gives
the same necessary condition in the continuum.  Thus
\((92.18)\) requires \(f\) to annihilate all gauge directions.
Gauge fixing or quotienting removes those directions; without it,
\((92.17)\) forbids a positive definite Hessian.
\end{proof}

This theorem explains why V54, V56, V86, and V87 consistently use a
reduced metric chart.  A claimed unreduced positive Einstein
Hessian would contradict exact diffeomorphism invariance.

\subsection{Concentration invisible to the integrated energy}

\begin{countertheorem}[A finite energy derivative can converge to a
singular stress]
\label{ctr:v14v92-delta}
Let \(M=(-1,1)\), choose a nonnegative
\(\varphi\in C_c^\infty((-1,1))\) with
\(\int\varphi=1\), and set
\[
 t_n(x)=n\varphi(nx).                                    \tag{92.20}
\]
Then
\[
 \int_Mt_n(x)\,dx=1                                     \tag{92.21}
\]
for all large \(n\), but
\[
 t_n\longrightarrow\delta_0                             \tag{92.22}
\]
as distributions and
\[
 \|t_n\|_{L^p}=n^{1-1/p}\|\varphi\|_{L^p}\longrightarrow\infty
 \quad(p>1).                                             \tag{92.23}
\]
Thus convergence of an integrated vacuum-energy derivative does not
give a regular local stress or the moment bounds in
Theorem~\ref{thm:v14v92-ui}.
\end{countertheorem}

\begin{proof}
Equation \((92.21)\) follows by \(y=nx\).  For a test function \(h\),
\[
 \int t_nh
 =
 \int\varphi(y)h(y/n)\,dy\longrightarrow h(0),
\]
which proves \((92.22)\).  The same substitution gives
\((92.23)\).
\end{proof}

On a noncompact space, translating a fixed unit-mass bump to spatial
infinity gives a second failure: all local distributional tests tend
to zero while the total integral stays one.  Local regularity and
spatial tightness are therefore independent source rows.

\subsection{Contact-term closure}

The second Einstein jet in \((92.8)\) is meaningful only if its two
parts use one renormalization prescription.  A local counterterm can
change \(\mathbb EC(h,k)\) without changing separated-point stress
covariances, whereas a source normalization can shift the first jet.
The finite card must record:
\[
 \text{bare contact}
 +\text{Jacobian contact}
 +\text{fermion contact}
 +\text{counterterm contact}.                            \tag{92.24}
\]

\begin{proposition}[Second-jet Ward compatibility]
\label{prop:v14v92-second-ward}
Suppose \(\Gamma\) is \(C^2\) and satisfies \((92.16)\).  At a
critical metric,
\[
 {\cal K}(h,{\cal L}_\xi g_*)=0                          \tag{92.25}
\]
for every test variation \(h\), where \({\cal K}=D^2W_{\rm ren}\)
is combined with the gravitational and counterterm Hessians in
\(D^2\Gamma\).  Therefore a matter covariance calculation that
violates \((92.25)\) must be completed by its contact, Jacobian, and
counterterm terms before it can enter the Einstein Hessian.
\end{proposition}

\begin{proof}
Equation \((92.25)\), with the stated complete Hessian, is exactly
the differentiated invariance identity \((92.19)\).  Splitting the
matter second derivative by \((91.5)\) gives the required
contact-minus-covariance combination.
\end{proof}

\subsection{Distributional source ledger}

\begin{acquisition}[Stress-limit and Bianchi card]
\label{acq:v14v92}
For the cofinal common regulator, record:
\begin{enumerate}
\item the lattice-to-continuum pairing defining \(T_n(h)\) and
      \(C_n(h,k)\);
\item joint convergence in probability for finite test families;
\item the \(L^{2+\delta}\) stress and \(L^{1+\delta}\) contact
      bounds;
\item uniform distributional seminorm bounds and spatial tightness;
\item the mean and insertion-level diffeomorphism defect estimates
      \((92.10)\), \((92.14)\);
\item the complete contact decomposition \((92.24)\);
\item the gauge-fixed metric slice and the projector onto the
      complement of diffeomorphism directions;
\item the reduced Hessian floor and compatibility
      \((92.25)\);
\item the identification of the limiting stress with the local net
      of V89 and the vacuum-energy jets of V91.
\end{enumerate}
\end{acquisition}

\begin{assessment}[Current source boundary]
\label{ass:v14v92-source}
The attached sources do not provide the uniform local moment bounds,
spatial tightness, anomaly estimates, or complete contact
renormalization required above.  They therefore do not yet produce
a conserved distributional Standard Model source for the Einstein
equation.
\end{assessment}

\begin{decision}[V93 literal source audit of the stress and
regulator rows]
\label{dec:v14v92-next}
Return to the attached grand-tensor and Einstein--SM manuscripts.
Locate their exact finite action, measure, stress, counterterm,
Ward/BRST, reflection, and regulator statements.  Classify each as
a populated formula, a conditional assumption, or a programme
requirement against the V84--V92 ledgers.  Do not infer the moment,
positivity, or cofinal estimates from formal covariance.  Use the
first genuinely populated formula to choose the next proof attack.
\end{decision}

\begin{firewall}[Claim boundary after V92]
\label{fire:v14v92}
The distributional theorem is conditional on convergence and moment
bounds not supplied by the current sources.  Bianchi compatibility
is a necessary consistency result, not existence of a solution.
No conserved renormalized SM stress tensor or full SM--Einstein
continuum has been constructed.
\end{firewall}

\section*{Version V92 append-only record}
\addcontentsline{toc}{section}{Version V92 append-only record}
The complete V91 source was retained before this continuation.  It
had \(1079758\) bytes, \(32404\) lines, and SHA--256
\[
 \texttt{639CE94E22A7C8CE44FFC8438D443CEB717FAE610070A51EE8FC44C08D2647C8}.
\]
V92 proved the uniform-integrability passage of stress jets,
distributional conservation from the Ward limit, Bianchi
compatibility, and the unavoidable diffeomorphism kernel of the
unreduced Einstein Hessian.
\section{Fourteenth-cycle V93: literal source audit and correction of the regulator status}
\label{sec:v14v93}

This continuation rereads the attached manuscripts as mathematical
sources.  Their acquisition directives are treated as statements of
their own programme, not as instructions governing this note.  The
audit corrects one overbroad statement in V88: the Grand Tensor
manuscript does construct an existential exactly projective reflected
control regulator.  It does not construct the physical
SM--Einstein continuum.

\subsection{Frozen source provenance}

\begin{record}[Audited manuscripts]
\label{rec:v14v93-sources}
The Grand Tensor source is
\begin{center}
\path{C:/Users/googo/Downloads/Renewal_Grand_Tensor_Monograph_V067.tex},
\end{center}
with SHA--256
\[
 \texttt{FACD058EC8BB210AD8B296C9D7F4D78A229D08D7A2AE378808700655CB8BC0FE}.
                                                               \tag{93.1}
\]
The attached Einstein--SM source is
\begin{center}
\path{attached_sources/Renewal_Einstein_Standard_Model_Physical_Source_Metric_Duhamel_Ancestry_and_Quantum_Continuum_Programme_V30.tex},
\end{center}
with SHA--256
\[
 \texttt{E0AE0E0990DE472021A1E90A9C10D41C3F747F28D97A5B1D8F1281852E3A0E69}.
                                                               \tag{93.2}
\]
\end{record}

\subsection{What is literally populated}

In the Grand Tensor subsection
\emph{A projective full-coupling reflected Renewal regulator}, the
definition labelled \texttt{def:RPESM-cutoff} specifies:
\[
 G_{\rm SM}=
 \frac{SU(3)\times SU(2)\times U(1)}{\mathbb Z_6},        \tag{93.3}
\]
a finite-range reflection-even bosonic action
\[
 S_{B,n}
 =
 S_{g,n}+S_{{\rm YM},n}+S_{H,n}+S_{{\rm gf},n}
 +S_{{\rm med},n}+S_{{\rm ct},n},                        \tag{93.4}
\]
and the positive control-modulus cylinder
\[
 d\Lambda_n^+
 =
 e^{-S_{B,n}}\|\sigma_n^+\|^2|p_n^+|^2\,d\lambda_n.
                                                               \tag{93.5}
\]
The actual complex determinant and real Pfaffian sections remain
marked line data; the positive sections
\(\sigma_n^+,p_n^+\) are explicitly introduced as massive control
sections.

The theorem labelled \texttt{thm:RPESM-reflected-slab} proves, for a
nonempty open coefficient set, finite normalizability and
\(L^{1+\eta}\) occurrence of fixed-support polynomial action, score,
current, stress, zero-mode, and line-norm writers.  Its reversible
generator is
\[
 L_n=\rho(R_n-I)+L_n^{\rm loc},\qquad \rho>0,             \tag{93.6}
\]
and its uniform centered floor comes from
\[
 -L_n\succeq\rho I.                                      \tag{93.7}
\]
The reflected stationary path law satisfies the exact conditional
square identity
\[
 \mathbb E_n[\overline{F(\Theta_n\omega)}F(\omega)]
 =
 \int|\mathbb E_n[F\mid\omega_0=x]|^2\,d\mu_n(x)\ge0.    \tag{93.8}
\]

The theorem labelled \texttt{thm:RPESM-projective-transport} defines
a normalized conditional detail kernel and proves
\[
 (\pi_{n+1/n})_\#\mu_{n+1}=\mu_n,                        \tag{93.9}
\]
as well as exact intertwining of the old Markov skeleton and
multiplicative determinant/Pfaffian line transport.  Its shell
normalization term is retained in the first and second action
writers.

\begin{theorem}[Corrected projective-regulator status]
\label{thm:v14v93-correction}
Equations \((93.5)\), \((93.8)\), and \((93.9)\) populate the
following rows for the named existential control regulator:
\begin{enumerate}
\item a common positive finite law at every cutoff;
\item exact projective consistency of its static marginals;
\item a reflected stationary Markov path law for bounded
      positive-time writers;
\item finite occurrence of the declared score, stress, contact,
      line, source, and held-out writer grammar.
\end{enumerate}
Accordingly, the claim in Assessment
\ref{ass:v14v88-source} that the attached sources exhibit no
projectively consistent reflected family is superseded by this
theorem.
\end{theorem}

\begin{proof}
The four conclusions are the literal contents of the definition and
theorems cited above.  The distinction between an existential control
regulator and the physical target is retained in their stated
hypotheses and firewall.
\end{proof}

\subsection{What population does not evaluate}

The same Grand Tensor firewall
\texttt{fire:RPESM-scope} explicitly states that exact projectivity
does not prove regulator removal, uniform locality, nuclearity,
split, clustering, a mass gap, full determinant/Pfaffian holonomy,
observed couplings and running, Lorentzian reconstruction,
scattering completeness, or the interacting SM--Einstein
continuum.

The Einstein--SM source sharpens the finite stopping line:
\begin{enumerate}
\item its V12 theorem
      \texttt{thm:V12-RPESM-occurrence} confirms common event and
      boundary-record occurrence, while
      \texttt{ass:V12-population-scope} says the first selected
      numerical time-split crossing matrix, relation residual, cone
      margin, line-transport margin, and cofinal energy window are
      not displayed;
\item V23 proves the generic positive joint line--crossing Gram
\[
 \begin{pmatrix}M&\Gamma\\\Gamma^*&N\end{pmatrix}\succeq0,
 \qquad
 D_{w\times}
 =N-\Gamma^*M^\dagger\Gamma\succeq0,                    \tag{93.10}
\]
      but does not evaluate the selected physical
      \((M,\Gamma,N)\);
\item V26 proves a conditional Gram-measure martingale telescope and
      a selected summable transport route for bounded square
      coordinates, but records the base numerical matrices, their
      supported floors, and physical line sewing as a remaining
      finite anchor;
\item V30 proves Kato transport of constant-rank relation supports,
      exact polar alignment, and the common-action plaquette
      criterion
\[
 r_i(\theta)r_j(\theta+he_i)
 =
 r_j(\theta)r_i(\theta+he_j),                            \tag{93.11}
\]
      while its verdict leaves the shrinking physical command square,
      returned-word Grams, direct responses, supported physical
      transport, line sewing, and derivative-scaled cutoff rectangle
      unpopulated.
\end{enumerate}

\begin{assessment}[Typed source matrix]
\label{ass:v14v93-matrix}
The literal status against V84--V92 is:
\begin{center}
\small
\begin{tabular}{p{0.29\textwidth}|p{0.22\textwidth}|p{0.39\textwidth}}
\textbf{Row}&\textbf{Status}&\textbf{Meaning}\\ \hline
positive finite control law&populated&Equation (93.5), using positive
massive control sections rather than the selected chiral target\\
projective static law&populated&Exact conditional-kernel construction
(93.9)\\
finite reflected path&populated&Markov conditional-square identity
(93.8)\\
finite writer occurrence&populated structurally&Stress, Ward, line,
and held-out objects exist as writers; their physical values and
uniform estimates are not thereby proved\\
fermion occupation Gram&unpopulated&No selected V84 common feature
matrix and all-degree physical Fock Gram ledger\\
selected chiral crossing&unpopulated&V23 is a compiler; V12 records
the physical crossing card as missing\\
local physical gap&unpopulated&The displayed floor (93.7) is supplied
by the global refresh in (93.6)\\
physical correspondence&unpopulated&No bound identifies the refresh
generator with the relativistic OS Hamiltonian\\
local stress limit&unpopulated&Finite occurrence and an
\(L^{1+\eta}\) statement do not give the uniform
\(L^{2+\delta}\), contact, tightness, and anomaly estimates of V92\\
Einstein continuum&open&The source firewalls expressly retain
regulator removal and the interacting continuum
\end{tabular}
\end{center}
\end{assessment}

\subsection{Imported exact tools without repetition}

The following source theorems are now adopted rather than reproved:
\begin{enumerate}
\item the projective conditional-shell construction
      \((93.9)\);
\item Kato normal transport of a constant-rank physical relation
      support and its polar finite chord;
\item the common-action plaquette equivalence \((93.11)\);
\item the primitive-complete positive Gram short and its conditional
      martingale Pythagoras;
\item the joint line--crossing Gram \((93.10)\).
\end{enumerate}
They complement V81--V92.  None supplies the missing numerical
physical cards merely by being a valid abstract theorem.

\subsection{The first active obstruction: global refresh}

The explicit source gap \((93.7)\) is generated by
\(\rho(R_n-I)\), where \(R_n\) replaces every state by the full
stationary mean.  This update is global.  It can connect distinct
phases and topological components even when all local moves fail to
do so.  It therefore cannot be inserted into \((86.9)\) as a
physical local mass gap without an independent correspondence
theorem.

\begin{proposition}[Source audit decision]
\label{prop:v14v93-decision}
The first literally populated formula that changes the current proof
direction is \((93.6)\), not another formal stress or Ward identity.
The next task is to separate the refresh floor from the local
generator spectrum and determine the exact condition under which
refresh may vanish without closing the local gap.
\end{proposition}

\begin{proof}
All later continuum conclusions require a physical Hamiltonian gap.
The only explicit uniform gap in the named projective regulator is
\((93.7)\), and its proof uses the refresh term.  The selected
fermionic crossing and physical correspondence cards remain absent.
Thus the provenance of this floor must be resolved before it can
support V86 or V87.
\end{proof}

\begin{decision}[V94 exact refresh removal]
\label{dec:v14v93-next}
For
\[
 A_{\rho,n}:=-L_n
 =A_n^{\rm loc}+\rho(I-R_n),
\]
prove the exact spectral shift on the centered space, characterize
its kernel and locality, and show by example how refresh hides
sector and phase zero modes.  Derive a necessary and sufficient
uniform condition for sending \(\rho_n\to0\) while retaining a gap.
Then return the genuinely local part \(A_n^{\rm loc}\) to the
V81--V82 influence and coverage compilers.
\end{decision}

\begin{firewall}[Claim boundary after V93]
\label{fire:v14v93}
The source correction recognizes a genuine existential projective
reflected regulator.  It does not identify that control regulator
with nature, the selected chiral Standard Model transfer, or a local
relativistic Hamiltonian.  Writer occurrence is not a moment bound,
and a global-refresh gap is not a particle mass gap.
\end{firewall}

\section*{Version V93 append-only record}
\addcontentsline{toc}{section}{Version V93 append-only record}
The complete V92 source was retained before this continuation.  It
had \(1092413\) bytes, \(32767\) lines, and SHA--256
\[
 \texttt{86D6CEE30421C40D72D257FD81F169D1ECE89CB926216D597986234755570857}.
\]
V93 performed a literal source audit, corrected the projective
regulator status, imported the existing exact projective, Kato,
plaquette, and Gram tools, and isolated the global refresh as the
first active physical-gap obstruction.

\section{Fourteenth-cycle V94: exact removal of the global refresh}
\label{sec:v14v94}

V93 found that the only explicit uniform spectral floor in the
literal projective regulator is supplied by a stationary global
refresh.  This continuation separates that artificial floor from
the local generator exactly.

\subsection{Centered spectral shift}

Let \((\Omega,\mu)\) be a probability space for which the following
operators are well defined on \(L^2(\mu)\).  Let
\[
 Rf=\int f\,d\mu,\qquad Q=I-R,                            \tag{94.1}
\]
so \(R\) is the orthogonal projection onto constants.  Let
\(A_{\rm loc}\ge0\) be self-adjoint and satisfy
\[
 A_{\rm loc}{\bf1}=0.                                    \tag{94.2}
\]
The source convention \(L_{\rm loc}\le0\) is
\(A_{\rm loc}=-L_{\rm loc}\).  Define
\[
 A_\rho=A_{\rm loc}+\rho Q,\qquad \rho>0.                \tag{94.3}
\]

\begin{theorem}[Exact refresh spectral shift]
\label{thm:v14v94-shift}
One has
\[
 RA_{\rm loc}=A_{\rm loc}R=0,\qquad
 QA_{\rm loc}=A_{\rm loc}Q=A_{\rm loc}.                  \tag{94.4}
\]
Consequently
\[
 A_\rho\big|_{\mathbb C{\bf1}}=0,\qquad
 A_\rho\big|_{QL^2}
 =
 A_{\rm loc}\big|_{QL^2}+\rho I.                        \tag{94.5}
\]
If
\[
 \gamma_{\rm loc}:=
 \inf\sigma(A_{\rm loc}|_{QL^2}),                        \tag{94.6}
\]
where zero is allowed, then
\[
 \boxed{\operatorname{gap}(A_\rho)=\rho+\gamma_{\rm loc}.}
                                                               \tag{94.7}
\]
In particular, a claimed lower bound \(\rho\) for \(A_\rho\) gives
no positive lower bound for \(\gamma_{\rm loc}\).
\end{theorem}

\begin{proof}
Self-adjointness and \((94.2)\) give, for every \(f\),
\[
 \langle{\bf1},A_{\rm loc}f\rangle
 =
 \langle A_{\rm loc}{\bf1},f\rangle=0,
\]
so \(RA_{\rm loc}=0\); also \(A_{\rm loc}R=0\).
This proves \((94.4)\) and the reducing decomposition
\((94.5)\).  Adding a scalar \(\rho I\) shifts the entire spectrum on
the centered subspace by \(\rho\), proving \((94.7)\).
\end{proof}

\begin{corollary}[Necessary and sufficient refresh-removal condition]
\label{cor:v14v94-remove}
For a cofinal family with \(\rho_n\downarrow0\),
\[
 \inf_n\operatorname{gap}(A_{\rho_n,n})>0
\]
after discarding finitely many indices if and only if
\[
 \liminf_{n\to\infty}\gamma_{{\rm loc},n}>0.             \tag{94.8}
\]
More generally,
\[
 \liminf_n\operatorname{gap}(A_{\rho_n,n})
 =
 \liminf_n(\rho_n+\gamma_{{\rm loc},n}).                 \tag{94.9}
\]
\end{corollary}

\begin{proof}
Equation \((94.9)\) is \((94.7)\) term by term.  When
\(\rho_n\to0\), its right-hand side has the same lower limit as
\(\gamma_{{\rm loc},n}\).
\end{proof}

\subsection{Exact semigroup mixture}

Because \(Q\) commutes with \(A_{\rm loc}\),
\[
 e^{-tA_\rho}
 =
 R+e^{-\rho t}e^{-tA_{\rm loc}}Q.                       \tag{94.10}
\]
Equivalently, with \(P_t^{\rm loc}=e^{-tA_{\rm loc}}\),
\[
 \boxed{
 e^{-tA_\rho}
 =
 e^{-\rho t}P_t^{\rm loc}
 +(1-e^{-\rho t})R.}                                     \tag{94.11}
\]

\begin{proposition}[Poisson reset interpretation]
\label{prop:v14v94-reset}
Equation \((94.11)\) is the transition semigroup that follows the
local dynamics until an independent rate-\(\rho\) event replaces the
entire state by an independent sample from \(\mu\).  Its off-diagonal
generator rate contains
\[
 q_{\rm ref}(\eta,\zeta)=\rho\,\mu(\zeta)                \tag{94.12}
\]
for arbitrary distinct configurations.  On a spatial product with
more than one cell, this permits one jump to change variables at
arbitrarily separated sites and is not a bounded-range block update.
\end{proposition}

\begin{proof}
The probability of no refresh by time \(t\) is \(e^{-\rho t}\);
stationarity makes the law after the first refresh equal to \(\mu\)
at time \(t\), giving \((94.11)\).  Expanding
\(\rho(R-I)\) in configuration coordinates gives
\((94.12)\).  It is positive for every target with positive
\(\mu\)-mass, regardless of the spatial diameter of the change.
\end{proof}

Algebraic preservation of a local observable space by addition of a
constant does not repair this move-level nonlocality.  The physical
question concerns propagation and the support of transitions, not
only the support label attached to the resulting function.

\subsection{Sector and phase zero modes hidden by refresh}

Let
\[
 \Omega=\Omega_1\sqcup\cdots\sqcup\Omega_k
\]
be the connected components of the local move graph, with positive
\(\mu\)-mass.  Proposition~\ref{prop:v14v82-kernel} gives
\[
 \ker A_{\rm loc}
 =
 \{f:f\text{ is constant on every }\Omega_j\}.           \tag{94.13}
\]

\begin{proposition}[Refresh lifts every relative-sector zero mode]
\label{prop:v14v94-sector}
The \(k-1\) centered functions in \(\ker A_{\rm loc}\) are
\(A_\rho\)-eigenvectors with eigenvalue \(\rho\).  Thus
\[
 k>1\quad\Longrightarrow\quad
 \gamma_{\rm loc}=0,\qquad
 \operatorname{gap}(A_\rho)=\rho.                        \tag{94.14}
\]
The refresh produces a unique stationary vector by directly
jumping between sectors; it does not prove that the local physical
dynamics connects them.
\end{proposition}

\begin{proof}
Every centered sector-constant function lies in \(Q L^2\) and is
annihilated by \(A_{\rm loc}\).  Equation \((94.5)\) therefore gives
eigenvalue \(\rho\).  The exact gap follows from \((94.7)\).
\end{proof}

The same conclusion applies to a local phase bottleneck whose gap
tends to zero: fixed \(\rho\) masks the soft mode by shifting it to
energy at least \(\rho\).  Sending the reset rate to zero exposes the
original collapse.

\subsection{Correspondence debit removes the artificial floor}

The V86 physical correspondence compares the update form with the
exact OS Hamiltonian.  Suppose, in the ideal local case, the physical
Hamiltonian after subtracting its ground energy is
\(H_{\rm phys}=A_{\rm loc}\) on the same centered space.

\begin{theorem}[Refresh cancels in the physical correspondence bound]
\label{thm:v14v94-correspondence}
The exact form difference is
\[
 A_\rho-H_{\rm phys}=\rho Q,\qquad
 \|A_\rho-H_{\rm phys}\|=\rho.                           \tag{94.15}
\]
No residual \(r<\rho\) can satisfy \((86.4)\) for all centered unit
vectors.  Using the sharp residual \(r=\rho\), the V86 lower bound
becomes
\[
 \operatorname{gap}(H_{\rm phys})
 \ge
 \operatorname{gap}(A_\rho)-\rho
 =
 \gamma_{\rm loc}.                                      \tag{94.16}
\]
Thus the refresh contributes no physical gap unless the target
Hamiltonian itself contains and justifies that global reset.
\end{theorem}

\begin{proof}
On \(QL^2\), the difference in \((94.15)\) equals \(\rho I\), so its
operator and quadratic-form norm is exactly \(\rho\).  Insert
\((94.7)\) and \(r=\rho\) into \((86.9)\) to obtain
\((94.16)\).
\end{proof}

More generally, if
\[
 \|H_{\rm phys}-A_{\rm loc}\|\le r_{\rm loc},
\]
then
\[
 \|H_{\rm phys}-A_\rho\|
 \le r_{\rm loc}+\rho,
\]
and the common-card estimate gives only
\[
 \operatorname{gap}(H_{\rm phys})
 \ge\gamma_{\rm loc}-r_{\rm loc}.                       \tag{94.17}
\]
Again the reset cancels.

\subsection{Projective compatibility}

Let
\[
 J_n:L^2(\mu_n)\longrightarrow L^2(\mu_{n+1}),\qquad
 J_nf=f\circ\pi_{n+1/n},                                 \tag{94.18}
\]
be the isometry supplied by exact projectivity \((93.9)\).  Then
\[
 R_{n+1}J_n=J_nR_n,\qquad Q_{n+1}J_n=J_nQ_n.             \tag{94.19}
\]

\begin{proposition}[Refresh transport is easy but not local]
\label{prop:v14v94-projective}
If
\[
 A_{{\rm loc},n+1}J_n=J_nA_{{\rm loc},n},
\]
then equal refresh rates give
\[
 A_{\rho,n+1}J_n=J_nA_{\rho,n}.                          \tag{94.20}
\]
For unequal rates,
\[
 \|(A_{\rho_{n+1},n+1}J_n
      -J_nA_{\rho_n,n})f\|
 =
 |\rho_{n+1}-\rho_n|\,\|Q_nf\|.                         \tag{94.21}
\]
Exact cutoff transport of the refresh therefore supplies no
information about spatial locality or physical correspondence.
\end{proposition}

\begin{proof}
Equation \((94.19)\) follows by integrating a pulled-back observable
against the projective law.  Substitute it and the assumed local
intertwining into \((94.3)\); the difference is
\((\rho_{n+1}-\rho_n)J_nQ_n\), proving both statements.
\end{proof}

\subsection{The local recovery route}

For the local block generator of V82, the exact source regulator may
populate the finite conditional laws.  The refresh-free lower bound
must instead come from
\[
 \boxed{
 \gamma_{{\rm loc},n}
 \ge
 \kappa_{n,s}
 \left(
  1-\max_i\sum_{j\ne i}
      \tanh\frac{\widehat\Delta_{ij,n}}4
 \right),}                                              \tag{94.22}
\]
or from a phase/interface conductance theorem when the row sum is at
least one.  The sector move graph must be restricted or connected
before \(\kappa\) and the influence matrix are evaluated.

\begin{acquisition}[Refresh-removal card]
\label{acq:v14v94}
For the literal RPESM regulator, acquire:
\begin{enumerate}
\item the declared physical local block family inside each sector;
\item the local conditional fiber gaps and fractional coverage;
\item the post-constraint mixed action oscillations and influence
      row sum;
\item a sector or phase interface estimate outside the Dobrushin
      region;
\item the refresh-free local gap
      \(\gamma_{{\rm loc},n}\);
\item the physical OS Hamiltonian and a correspondence residual to
      \(A_{{\rm loc},n}\), with the refresh excluded;
\item a locality or finite-propagation estimate for that Hamiltonian;
\item the cofinal choice \(\rho_n\to0\) and verification of
      \((94.8)\).
\end{enumerate}
The static law, its normalized stress, and its partition function do
not depend on the artificial rate \(\rho\), while time correlations
and the Hamiltonian do.  The refresh therefore cannot be used to
supply the absolute Einstein vacuum energy either.
\end{acquisition}

\begin{assessment}[Current literal status]
\label{ass:v14v94-source}
The Grand Tensor source defines \(L_n^{\rm loc}\) only as a finite
sum of local conditional-expectation generators and proves its
reversibility.  Its displayed uniform floor uses the refresh term.
No numerical or symbolic local influence matrix, coverage constant,
interface estimate, refresh-free gap, or OS-Hamiltonian
correspondence residual is supplied.
\end{assessment}

\begin{decision}[V95 compulsory companion audit]
\label{dec:v14v94-next}
Inspect the newest Yang--Mills and Navier--Stokes notes.  Test whether
their current local symbol, cover, or constrained-kernel estimates
can populate a refresh-free fiber or coverage row for the literal
RPESM regulator.  Continue in V96 with the first such compatible
row, or, if none exists, derive the sharp finite-volume local
conditional generator directly from the projective shell kernel
\((93.9)\).
\end{decision}

\begin{firewall}[Claim boundary after V94]
\label{fire:v14v94}
The exact source regulator remains a valid projective reflected
control construction.  V94 proves that its global-refresh floor is
not a local physical mass gap and disappears from an honest
correspondence estimate.  No refresh-free uniform gap or relativistic
Hamiltonian has been proved.
\end{firewall}

\section*{Version V94 append-only record}
\addcontentsline{toc}{section}{Version V94 append-only record}
The complete V93 source was retained before this continuation.  It
had \(1103350\) bytes, \(33045\) lines, and SHA--256
\[
 \texttt{4F5EA5CA6367BD7E2ADDFBD71DB8B60A85B9A4C696CE2F6D66A41EB3DAA784F2}.
\]
V94 proved the exact centered spectral shift from the source's global
refresh, its sector and locality obstruction, and its cancellation
from the physical correspondence bound.
\section{Fourteenth-cycle V95: compulsory audit and the absence of a refresh-free companion gap}
\label{sec:v14v95}

The newest companion files were inspected read-only after V94.

\begin{record}[V95 companion checkpoint]
\label{rec:v14v95-companions}
The Yang--Mills checkpoint is
\begin{center}
\path{Renewal_Yang_Mills_Mass_Gap_Research_Notes_V29.tex},
\end{center}
with \(232987\) bytes, \(6340\) lines, and SHA--256
\[
 \texttt{85AEAC9219611F19F1C2F98A2F92C4253B7F32F7120F17183099CCA237B0F0B8}.
                                                               \tag{95.1}
\]
The Navier--Stokes checkpoint remains
\begin{center}
\path{Renewal_NS_Stability_Research_Notes_V26.tex},
\end{center}
with \(278283\) bytes, \(5319\) lines, and SHA--256
\[
 \texttt{82C887F8C2C69E4F28FAD7C3DC8DE5B9099CB5A4BF431EBA1FFF3E91935978AD}.
                                                               \tag{95.2}
\]
\end{record}

Yang--Mills V29 proves an exact negative direction for the fixed V10
\(90\times90\) certificate block at a physical rational
trigonometric phase inside radius \(1/4\).  It carefully does not
infer negativity of the direct Wilson pencil: the obstruction closes
only the route using one fixed splitting.  Its surviving direction
is a local atlas of independently feasible splittings with overlap,
or a direct active-circle factorization if the local block itself is
infeasible.

\begin{assessment}[YM applicability to refresh removal]
\label{ass:v14v95-ym}
The V28 radius-\(1/8\) tube remains a valid local symbol
certificate, and V29 supplies a genuine obstruction to extending
its fixed splitting.  Neither theorem:
\begin{enumerate}
\item covers the physical momentum torus;
\item defines the conditional shell law of the RPESM regulator;
\item gives a block-move coverage or influence matrix;
\item proves a refresh-free Markov or OS Hamiltonian gap.
\end{enumerate}
Thus no YM numerical floor can yet be inserted into
\(\gamma_{{\rm loc},n}\).  The transferable rule is to allow
chart-dependent local splittings and certify their overlaps, rather
than force one splitting past a proved physical obstruction.
\end{assessment}

Navier--Stokes V26 remains a deterministic rank-one Oseen-kernel
estimate and supplies no conditional generator or shell gap.

\begin{theorem}[V95 audit verdict]
\label{thm:v14v95-verdict}
At checkpoints \((95.1)\)--\((95.2)\), neither companion populates a
refresh-free fiber gap, fractional coverage, Dobrushin row, interface
bound, or update-to-physical correspondence for the literal
projective RPESM regulator.
\end{theorem}

\begin{proof}
The Yang--Mills results are local finite symbol inequalities on
proper momentum subsets and explicitly leave the interacting
measure, compact screen, continuum response, and Hamiltonian gap
open.  The Navier--Stokes result concerns a different deterministic
kernel.  Their conclusion types do not match the local-generator
rows of V94.
\end{proof}

\begin{decision}[V96 projective-shell generator]
\label{dec:v14v95-next}
Use the exact conditional detail kernel in the source formula
\((93.9)\).  Write the fine Markov skeleton as a conditional
expectation sandwich between the coarse and fine \(L^2\) spaces.
Determine its complete spectrum, its refresh-free Dirichlet gap, and
its kernel.  Then distinguish:
\begin{enumerate}
\item inheritance of the coarse spectral gap;
\item zero eigenvalues of the one-step transfer on new-detail
      directions;
\item whether the conditional detail update is spatially local or
      refreshes an entire shell;
\item what modification can make the transfer faithful without
      importing the global refresh of V94.
\end{enumerate}
\end{decision}

\begin{firewall}[Claim boundary after V95]
\label{fire:v14v95}
The audit imports no companion constant.  The YM physical
fixed-splitting obstruction is respected, and the literal RPESM
local gap remains uncomputed.
\end{firewall}

\section*{Version V95 append-only record}
\addcontentsline{toc}{section}{Version V95 append-only record}
The complete V94 source was retained before this continuation.  It
had \(1114934\) bytes, \(33394\) lines, and SHA--256
\[
 \texttt{F8350B77D7F4A4981FD0FB4A5D20791D0736326DE3D37CEA67067DDDE02D6E0D}.
\]
V95 performed the compulsory companion audit and found no
refresh-free local generator input, so the next step turns directly
to the literal projective shell kernel.

\section{Fourteenth-cycle V96: exact shell-sandwich spectrum and the scalar-laziness obstruction}
\label{sec:v14v96}

V95 found no companion estimate that could replace the global
refresh.  We therefore return to the literal conditional shell law
of the Grand Tensor source.  The source formulas \((93.9)\) and
\({\rm(RP.13)}\) say
\[
 \mu_{n+1}(\dd x,\dd y)
   =\mu_n(\dd x)K_{n+1/n}(\dd y\mid x),
 \qquad
 P_{n+1}((x,y),\dd(x',y'))
   =P_n(x,\dd x')K_{n+1/n}(\dd y'\mid x').
                                                               \tag{96.1}
\]
The following calculation extracts exactly what this construction
does and does not prove about a refresh-free gap.

Put
\[
 {\cal H}_n=L^2(\mu_n),\qquad
 {\cal H}_{n+1}=L^2(\mu_{n+1}),
\]
and define
\[
 (J_nf)(x,y)=f(x),\qquad
 (J_n^*F)(x)=\int F(x,y)K_{n+1/n}(\dd y\mid x).
                                                               \tag{96.2}
\]

\begin{lemma}[Conditional lift and orthogonal detail space]
\label{lem:v14v96-lift}
The map \(J_n:{\cal H}_n\to{\cal H}_{n+1}\) is an isometry,
\(J_n^*\) is the conditional-expectation adjoint, and
\[
 J_n^*J_n=I,\qquad
 \Pi_n:=J_nJ_n^*
\]
is the orthogonal projection onto the closed coarse-observable
space \({\cal C}_n:=\operatorname{Ran}J_n\).  Moreover
\[
 {\cal H}_{n+1}
  ={\cal C}_n\oplus{\cal D}_{n+1},
 \qquad
 {\cal D}_{n+1}:=\ker J_n^*
  =\{F:\mathbb E[F\mid x]=0\}.
                                                               \tag{96.3}
\]
\end{lemma}

\begin{proof}
Disintegration in \((96.1)\) gives
\[
 \|J_nf\|_{{\cal H}_{n+1}}^2
 =\int |f(x)|^2\mu_n(\dd x)
       \int K_{n+1/n}(\dd y\mid x)
 =\|f\|_{{\cal H}_n}^2.
\]
The same calculation with \(\overline fF\) proves the displayed
formula for \(J_n^*\).  Hence \(J_n^*J_n=I\).  The standard
isometry identities make \(J_nJ_n^*\) the orthogonal projection
onto \(\operatorname{Ran}J_n\), whose orthogonal complement is
\(\ker J_n^*\).
\end{proof}

\begin{theorem}[Complete spectrum of the projective shell skeleton]
\label{thm:v14v96-spectrum}
Let \(P_n\) be the reversible Markov contraction on
\({\cal H}_n\), and let \(P_{n+1}\) be the source skeleton
\((96.1)\).  Then
\[
 \boxed{P_{n+1}=J_nP_nJ_n^*.}                                \tag{96.4}
\]
Relative to \((96.3)\),
\[
 P_{n+1}\simeq
 \begin{pmatrix}P_n&0\\0&0\end{pmatrix}.
                                                               \tag{96.5}
\]
Consequently
\[
 \ker P_{n+1}
   =J_n(\ker P_n)\oplus{\cal D}_{n+1}.                         \tag{96.6}
\]
If \({\cal D}_{n+1}\ne\{0\}\), then
\[
 \sigma(P_{n+1})=\sigma(P_n)\cup\{0\},                         \tag{96.7}
\]
with the old spectral multiplicities retained and with an
additional zero sector equal to the conditional-detail space.
\end{theorem}

\begin{proof}
For \(F\in{\cal H}_{n+1}\), the kernel in \((96.1)\) gives
\[
 \begin{split}
 (P_{n+1}F)(x,y)
 &=\int P_n(x,\dd x')\int
       F(x',y')K_{n+1/n}(\dd y'\mid x')\\
 &=(J_nP_nJ_n^*F)(x,y).
 \end{split}
\]
Thus \({\cal C}_n\) reduces \(P_{n+1}\), its restriction there is
unitarily equivalent to \(P_n\), and the operator vanishes on
\({\cal D}_{n+1}\).  Equations \((96.5)\)--\((96.7)\) follow from
the spectrum and kernel of an orthogonal direct sum.
\end{proof}

\begin{corollary}[Exact refresh-free Dirichlet-gap inheritance]
\label{cor:v14v96-dirichlet-gap}
Fix a time step \(a_n>0\) and define
\[
 A_n^{\rm D}=a_n^{-1}(I-P_n),\qquad
 A_{n+1}^{\rm D}=a_n^{-1}(I-P_{n+1}).
\]
Let \(\gamma_n^{\rm D}\) and \(\gamma_{n+1}^{\rm D}\) be their
spectral gaps above the constant state.  Then
\[
 \boxed{\gamma_{n+1}^{\rm D}
   =\min\{\gamma_n^{\rm D},a_n^{-1}\}.}                        \tag{96.8}
\]
If \(P_n\ge0\) and the old centered space is nonzero, then
\(\gamma_n^{\rm D}\le a_n^{-1}\), so
\[
 \gamma_{n+1}^{\rm D}=\gamma_n^{\rm D}.                        \tag{96.9}
\]
This conclusion uses no global refresh.

\end{corollary}

\begin{proof}
The constant state lies in \({\cal C}_n\).  On the centered part of
\({\cal C}_n\), equation \((96.5)\) gives exactly the old
Dirichlet spectrum.  On \({\cal D}_{n+1}\), it gives the single
Dirichlet value \(a_n^{-1}\).  Taking the lower edge proves
\((96.8)\).  If \(P_n\ge0\), every nonconstant eigenvalue or
spectral value \(\lambda\) obeys \(0\le\lambda\le1\), so the old
Dirichlet lower edge is at most \(a_n^{-1}\).
\end{proof}

\begin{assessment}[What the inherited gap means]
\label{ass:v14v96-gap-meaning}
Equation \((96.9)\) is a genuine refresh-free finite-volume
statement, but it creates no new uniform constant: the fine gap is
exactly the unproved old gap.  Iteration transports a lower bound
only after one base-scale bound has been supplied with the physical
volume and clock dependence required by the continuum limit.
Furthermore the transition in \((96.1)\) redraws the entire detail
shell after the coarse move.  If the shell contains arbitrarily
many spatial cells and the conditional law has full support, then
the generator \(P_{n+1}-I\) has direct transitions between detail
configurations differing at arbitrarily many separated cells.
It is therefore a shell-global update, not a bounded-range
physical generator.
\end{assessment}

The zero block in \((96.5)\) causes a different obstruction for the
exact transfer Hamiltonian.

\begin{proposition}[Dirichlet gap does not imply a finite logarithmic Hamiltonian]
\label{prop:v14v96-log-kernel}
Assume \({\cal D}_{n+1}\ne\{0\}\) and \(P_n\ge0\).  There is no
densely defined self-adjoint operator \(H_{n+1}\), finite on a dense
domain containing \({\cal D}_{n+1}\), such that
\[
 P_{n+1}=e^{-a_nH_{n+1}}.                                    \tag{96.10}
\]
Spectral calculus assigns infinite energy to every nonzero detail
vector.  Thus the finite Dirichlet value \(a_n^{-1}\) in
\((96.8)\) must not be reported as a finite energy of
\(-a_n^{-1}\log P_{n+1}\).
\end{proposition}

\begin{proof}
For every self-adjoint \(H\), the positive function
\(e^{-a_n t}\) is strictly positive at every finite
\(t\in\mathbb R\).  Its spectral integral therefore gives
\(\ker e^{-a_nH}=\{0\}\): if \(e^{-a_nH}u=0\), then
\[
 0=\|e^{-a_nH}u\|^2
   =\int_{\mathbb R}e^{-2a_nt}\,\dd\langle E_H(t)u,u\rangle
\]
forces the spectral measure of \(u\) to vanish.  Equation
\((96.6)\) gives the contradictory inclusion
\({\cal D}_{n+1}\subset\ker P_{n+1}\).
\end{proof}

One might try to remove the kernel by scalar laziness.  For
\(0<\eta<1\), set
\[
 P_{n+1}^{(\eta)}
   =\eta I+(1-\eta)J_nP_nJ_n^*.                              \tag{96.11}
\]
When \(P_n\ge0\), this is a positive, injective, reversible Markov
operator.  It has eigenvalue \(\eta\) on
\({\cal D}_{n+1}\), while an old eigenvalue \(\lambda\) becomes
\(\eta+(1-\eta)\lambda\).

\begin{theorem}[No scalar-laziness repair on one continuum clock]
\label{thm:v14v96-laziness-no-go}
Suppose \(a\downarrow0\) and an old physical mode satisfies
\[
 \lambda(a)=1-aE+o(a),\qquad 0<E<\infty.                      \tag{96.12}
\]
For any \(0<\eta(a)<1\), define its coarse and detail logarithmic
energies under \((96.11)\) by
\[
 E_{\rm old}^{(\eta)}(a)
 =-a^{-1}\log\!\left(\eta(a)+(1-\eta(a))\lambda(a)\right),
 \qquad
 E_{\rm det}^{(\eta)}(a)=-a^{-1}\log\eta(a).                  \tag{96.13}
\]
It is impossible that both
\[
 E_{\rm old}^{(\eta)}(a)\longrightarrow cE
 \quad\hbox{for some }c>0,
 \qquad
 E_{\rm det}^{(\eta)}(a)\longrightarrow M
 \quad\hbox{for some }0<M<\infty.                            \tag{96.14}
\]
More precisely:
\begin{enumerate}
\item if \(\eta(a)\to\eta_0<1\), then
      \(E_{\rm old}^{(\eta)}(a)\to(1-\eta_0)E\), whereas
      \(E_{\rm det}^{(\eta)}(a)\to+\infty\);
\item if \(E_{\rm det}^{(\eta)}(a)\to M<\infty\), then
      \(\eta(a)=1-aM+o(a)\) and
      \(E_{\rm old}^{(\eta)}(a)\to0\).
\end{enumerate}
Scalar laziness also changes the exact old-clock intertwining to
\[
 P_{n+1}^{(\eta)}J_n
  =J_n\bigl(\eta I+(1-\eta)P_n\bigr),                         \tag{96.15}
\]
so it is not a faithful lift of \(P_n\) unless \(\eta=0\).
\end{theorem}

\begin{proof}
Equation \((96.12)\) gives
\[
 \eta+(1-\eta)\lambda
  =1-(1-\eta)aE+o\!\left(a(1-\eta)\right).
\]
If \(\eta(a)\to\eta_0<1\), the expansion
\(-\log(1-z)=z+o(z)\) proves the first old-mode limit, while
\(-a^{-1}\log\eta(a)\) diverges along every subsequence on which
\(\eta\) is bounded away from one.  If the detail energy tends to
\(M<\infty\), then
\(\log\eta(a)=-aM+o(a)\), hence
\(\eta(a)=1-aM+o(a)\).  Substitution yields
\[
 \eta+(1-\eta)\lambda
   =1-a^2ME+o(a^2),
\]
and the old energy tends to zero.  Finally \((96.15)\) follows
from \(J_n^*J_n=I\).
\end{proof}

\begin{decision}[V97 faithful detail dynamics]
\label{dec:v14v96-next}
Replace the scalar repair by a genuine operator acting on the
detail directions.  Put the old and detail dynamics on one fine
Hilbert space and study a symmetric split transfer
\[
 T_{n+1}(a)
   =e^{-aB_{n+1}/2}\,\widetilde T_n(a)\,
      e^{-aB_{n+1}/2},                                      \tag{96.16}
\]
where \(B_{n+1}\) is a sum of constraint-preserving local detail
generators and \(\widetilde T_n\) is a faithful lift of the old
transfer rather than the projection sandwich.  State the precise
intertwining condition, prove positivity and injectivity, and
derive a gap or an explicit commutator debit.  If a faithful lift
cannot preserve the interacting conditional law, record that
obstruction before introducing any estimate.
\end{decision}

\begin{firewall}[Claim boundary after V96]
\label{fire:v14v96}
The exact projective skeleton inherits the old Dirichlet gap, but
does not prove a uniform physical mass gap, bounded-range locality,
or a finite logarithmic Hamiltonian on the new-detail space.
Scalar laziness cannot repair all three defects on the same
continuum clock.
\end{firewall}

\section*{Version V96 append-only record}
\addcontentsline{toc}{section}{Version V96 append-only record}
The complete V95 source was retained before this continuation.  It
had \(1119321\) bytes, \(33502\) lines, and SHA--256
\[
 \texttt{EACAB02D2A963DFC7265EFF63543C4B9C21D42C99EA470EC1AA5C6DEA3D0F5A7}.
\]
V96 proved the exact conditional-lift block spectrum, the
refresh-free Dirichlet-gap inheritance formula, the logarithmic
detail-kernel obstruction, and the one-clock scalar-laziness
no-go.


\section{Fourteenth-cycle V97: faithful conditional transport and a split-transfer gap}
\label{sec:v14v97}

V96 showed that the source sandwich \(J_nP_nJ_n^*\) obtains exact
projectivity by erasing every centered detail observable in one
step.  A finite logarithmic Hamiltonian requires a faithful
transport of those observables.  Such a transport is not determined
by the conditional law alone.  We isolate the missing datum before
using it.

\begin{definition}[Conditional-fiber trivialization]
\label{def:v14v97-trivialization}
Write \(K_x=K_{n+1/n}(\,\cdot\mid x)\).  A conditional-fiber
trivialization is a probability space \((U,\nu)\) and a measurable
family of measure-space isomorphisms
\[
 \phi_x:(U,\nu)\longrightarrow(Y_x,K_x)                       \tag{97.1}
\]
defined modulo null sets.  It is reflection covariant if there is a
measure-preserving involution \(\vartheta_U\) such that
\[
 \phi_{\vartheta x}\circ\vartheta_U
   =\vartheta_Y\circ\phi_x\quad\hbox{almost everywhere}.       \tag{97.2}
\]
The induced lattice unitary
\[
 {\cal W}:{\cal H}_{n+1}\longrightarrow
 L^2(\mu_n)\mathbin{\widehat\otimes}L^2(\nu),
 \qquad
 ({\cal W}F)(x,u)=F(x,\phi_x(u)),                              \tag{97.3}
\]
preserves constants, complex conjugation, and pointwise
positivity.
\end{definition}

This is an additional structure row.  Disintegration provides the
conditional measures \(K_x\), but it does not canonically select
the coherent maps \(\phi_x\), their reflection covariance, or their
spatial locality.

\begin{theorem}[Faithful old-clock lift]
\label{thm:v14v97-faithful-lift}
Let \(T_n=e^{-aH_n}\) be an injective, self-adjoint,
positivity-preserving Markov transfer on \(L^2(\mu_n)\), where
\(H_n\ge0\), \(H_n1=0\).  Given \((97.1)\), define
\[
 \widetilde T_n
 ={\cal W}^{-1}(T_n\otimes I){\cal W},
 \qquad
 X_n={\cal W}^{-1}(H_n\otimes I){\cal W}.                     \tag{97.4}
\]
Then \(\widetilde T_n=e^{-aX_n}\) is injective, self-adjoint,
positivity preserving, Markov, and
\[
 \widetilde T_nJ_n=J_nT_n,\qquad
 J_n^*\widetilde T_n=T_nJ_n^*.                               \tag{97.5}
\]
If \(H_n\) has simple kernel and gap \(\gamma_n>0\), then
\[
 \ker X_n={\cal W}^{-1}
 \bigl(\operatorname{span}\{1\}\otimes L^2(\nu)\bigr),         \tag{97.6}
\]
and \(X_n\) has gap \(\gamma_n\) off this detail-degenerate
kernel.
\end{theorem}

\begin{proof}
Under \({\cal W}\), the coarse lift \(J_nf\) is \(f\otimes1\).
All operator assertions follow by unitary conjugacy from
\(T_n\otimes I\).  In particular,
\[
 (T_n\otimes I)(f\otimes1)=(T_nf)\otimes1,
\]
which gives the first identity in \((97.5)\); taking adjoints gives
the second.  Equation \((97.6)\) and the gap statement follow from
the tensor-product spectral resolution.
\end{proof}

The degeneracy in \((97.6)\) is removed by actual detail dynamics,
not by scalar laziness.  Let \(B\ge0\) be a self-adjoint Markov
generator on \(L^2(\nu)\), with
\[
 \ker B=\operatorname{span}\{1\},\qquad
 B\ge\beta(I-E_\nu),\qquad \beta>0,                           \tag{97.7}
\]
where \(E_\nu\) projects onto constants.

\begin{corollary}[Exact commuting completion]
\label{cor:v14v97-commuting}
Define
\[
 Y_n={\cal W}^{-1}(I\otimes B){\cal W},\qquad
 H_{n+1}^{\rm prod}=X_n+Y_n.                                 \tag{97.8}
\]
Then \(X_n\) and \(Y_n\) strongly commute,
\[
 e^{-aH_{n+1}^{\rm prod}}
 ={\cal W}^{-1}(e^{-aH_n}\otimes e^{-aB}){\cal W},            \tag{97.9}
\]
and
\[
 \ker H_{n+1}^{\rm prod}=\operatorname{span}\{1\},\qquad
 \boxed{\operatorname{gap}(H_{n+1}^{\rm prod})
        =\min\{\gamma_n,\beta\}.}                             \tag{97.10}
\]
The transfer is injective and has the exact old-clock
intertwining
\[
 e^{-aH_{n+1}^{\rm prod}}J_n
   =J_ne^{-aH_n}.                                             \tag{97.11}
\]
\end{corollary}

\begin{proof}
Equations \((97.8)\)--\((97.9)\) are the tensor sum and tensor
semigroup identities.  Its spectral energies are all sums
\(E+\varepsilon\), with only \(E=\varepsilon=0\) at the constant
state.  The least positive sum is
\(\min\{\gamma_n,\beta\}\).  Since \(e^{-aB}1=1\), equation
\((97.11)\) follows.
\end{proof}

The product construction is a mathematical completion, but its
physical locality depends on \({\cal W}\).  A generator local in
the reference coordinate \(u\) can become nonlocal in the physical
detail \(y=\phi_x(u)\).  We therefore also give a bound for a
physical detail generator which need not commute with the faithful
old lift.

Let \({\cal H}\) be one fine \(L^2\) space, let \(X,Y\ge0\) be
self-adjoint Markov generators, and suppose their kernel
projections are \(E_X,E_Y\).  Let \(E_0\) project onto constants.
Assume
\[
 \ker X\cap\ker Y=\operatorname{Ran}E_0,                      \tag{97.12}
\]
\[
 X\ge\gamma(I-E_X),\qquad
 Y\ge\beta(I-E_Y),                                            \tag{97.13}
\]
and define the centered Friedrichs-angle constant
\[
 c=\bigl\|(E_X-E_0)(E_Y-E_0)\bigr\|,\qquad 0\le c<1.           \tag{97.14}
\]

\begin{lemma}[Two-kernel angle coercivity]
\label{lem:v14v97-angle}
Under \((97.12)\)--\((97.14)\),
\[
 X+Y\ge\Gamma(I-E_0),
 \qquad
 \boxed{\Gamma=\min\{\gamma,\beta\}(1-c).}                    \tag{97.15}
\]
\end{lemma}

\begin{proof}
For centered \(f\), equations \((97.13)\) give
\[
 \langle f,(X+Y)f\rangle
 \ge\min\{\gamma,\beta\}
 \left(\|(I-E_X)f\|^2+\|(I-E_Y)f\|^2\right).                  \tag{97.16}
\]
On the centered space the two kernel projections are
\(P=E_X-E_0\) and \(Q=E_Y-E_0\).  The standard two-projection
decomposition, obtained by diagonalizing the positive contraction
\(PQP\) on \(\operatorname{Ran}P\), gives
\[
 \|P+Q\|\le1+\|PQ\|=1+c.
\]
Therefore
\[
 \|(I-P)f\|^2+\|(I-Q)f\|^2
 =\langle f,(2I-P-Q)f\rangle
 \ge(1-c)\|f\|^2.
\]
Together with \((97.16)\), this proves \((97.15)\).
\end{proof}

For a bounded finite-carrier approximation, define the symmetric
split transfer
\[
 S(a)=e^{-aY/2}e^{-aX}e^{-aY/2}.                              \tag{97.17}
\]
It is self-adjoint, operator-positive, positivity preserving,
Markov, and injective.

\begin{theorem}[Explicit split-transfer commutator debit]
\label{thm:v14v97-split}
Assume in addition that \(X,Y\) are bounded, and put
\[
 C=\|[X,Y]\|,\qquad
 q(a)=e^{-a\Gamma}+a^2C.                                     \tag{97.18}
\]
Then
\[
 \|S(a)-e^{-a(X+Y)}\|\le a^2C.                               \tag{97.19}
\]
If \(q(a)<1\), the exact logarithmic split Hamiltonian
\[
 H_{\rm split}(a)=-a^{-1}\log S(a)
\]
has simple vacuum and satisfies
\[
 \boxed{\operatorname{gap}H_{\rm split}(a)
 \ge-a^{-1}\log q(a).}                                       \tag{97.20}
\]
Its Dirichlet gap satisfies
\[
 \operatorname{gap}\bigl(a^{-1}(I-S(a))\bigr)
 \ge\frac{1-e^{-a\Gamma}-a^2C}{a}.                           \tag{97.21}
\]
Thus the noncommuting physical split loses an explicit
\(O(aC)\) amount from the continuous-time lower bound.
\end{theorem}

\begin{proof}
For nonnegative \(X,Y\), all semigroups below are contractions.
The double Duhamel identity gives, for \(s,t\ge0\),
\[
 \|e^{-sY}e^{-tX}-e^{-tX}e^{-sY}\|
 \le st\|[X,Y]\|.                                             \tag{97.22}
\]
Hence moving the left half-step in \(S(a)\) past \(e^{-aX}\)
costs at most \(a^2C/2\).

For the Lie product, set
\[
 F(s)=e^{-(a-s)(X+Y)}e^{-sX}e^{-sY}.
\]
Differentiation and one Duhamel expansion of
\(Ye^{-sX}-e^{-sX}Y\) give
\[
 \|F'(s)\|\le sC.
\]
Integration from \(0\) to \(a\) yields
\[
 \|e^{-aX}e^{-aY}-e^{-a(X+Y)}\|\le a^2C/2.
\]
The triangle inequality proves \((97.19)\).

By Lemma~\ref{lem:v14v97-angle}, the restriction of
\(e^{-a(X+Y)}\) to the centered space has norm at most
\(e^{-a\Gamma}\).  Equation \((97.19)\) therefore bounds the
centered norm of \(S(a)\) by \(q(a)\).  Since \(S(a)\) is positive
and fixes constants, \(q(a)<1\) makes the vacuum simple and bounds
every nonconstant transfer eigenvalue above by \(q(a)\).
Monotonicity of \(-a^{-1}\log t\) proves \((97.20)\), and the same
spectral comparison for \(a^{-1}(I-S(a))\) proves \((97.21)\).
\end{proof}

\begin{assessment}[Rows still needed for the RPESM shell]
\label{ass:v14v97-missing}
The source conditional density gives \(K_x\), but the following
are not yet populated:
\begin{enumerate}
\item a reflection-covariant, cutoff-compatible fiber
      trivialization \((97.1)\);
\item locality or a uniform quasilocal tail bound for
      \(x,u\mapsto\phi_x(u)\);
\item a constraint-preserving physical detail generator with
      uniform \(\beta\);
\item a uniform angle \(c<1\) and a clock-scaled commutator bound
      \(a\|[X,Y]\|\to0\).
\end{enumerate}
Without these rows, \((97.10)\) is a conditional product theorem
and \((97.20)\) is a finite-regulator compiler, not a proof of the
RPESM continuum gap.
\end{assessment}

\begin{decision}[V98 locality of conditional transport]
\label{dec:v14v97-next}
Go upstream to the conditional density
\[
 K_x(\dd y)=Z(x)^{-1}e^{-V(x,y)}\kappa_x(\dd y).
\]
Construct a canonical transport from the reference kernel
\(\kappa_x\) to \(K_x\), and determine how a change of one old
cell propagates through that transport.  Prove a finite influence
or quasilocality bound under explicit mixed-Hessian assumptions.
If the compact gauge, Higgs, or constraint factors invalidate
those assumptions, isolate the first invalid factor rather than
claiming a local fiber trivialization.
\end{decision}

\begin{firewall}[Claim boundary after V97]
\label{fire:v14v97}
Faithful projective transport and a finite logarithmic Hamiltonian
are now constructed under an explicit conditional-fiber
trivialization.  The required covariant local trivialization and
its uniform constants have not been derived for the literal source
density.
\end{firewall}

\section*{Version V97 append-only record}
\addcontentsline{toc}{section}{Version V97 append-only record}
The complete V96 source was retained before this continuation.  It
had \(1129571\) bytes, \(33801\) lines, and SHA--256
\[
 \texttt{EAA9E848875B07AB20C8C2564332230E109EC62BA26C690686FF86005DF9ED4F}.
\]
V97 replaced projection-based detail erasure by a faithful
conditional-fiber lift, proved the exact tensor-sum gap, and derived
a two-kernel angle bound with an explicit split-transfer
commutator debit.


\section{Fourteenth-cycle V98: conditional Moser transport, finite influence, and the compact/nonconvex obstruction}
\label{sec:v14v98}

V97 reduced faithful lifting to a cutoff-compatible family of
conditional-fiber isomorphisms.  We first construct such a family
on a Euclidean uniformly convex shell and quantify its dependence
on the old field.  We then test each hypothesis against the literal
RPESM source class.

Let \(x\in{\cal X}\subset\mathbb R^p\), \(y\in\mathbb R^m\), and
\[
 K_x(\dd y)=p_x(y)\dd y,\qquad
 p_x(y)=Z(x)^{-1}e^{-W(x,y)}.                                \tag{98.1}
\]
Write
\[
 {\cal L}_x=\Delta_y-\nabla_yW(x,\cdot)\cdot\nabla_y.
\]

\begin{theorem}[Conditional Moser transport with a sharp energy estimate]
\label{thm:v14v98-moser}
Suppose \(W\) is \(C^3\), has sufficient growth to justify the
integrations below, and obeys the uniform conditional convexity
bound
\[
 \nabla_y^2W(x,y)\succeq m_*I,\qquad m_*>0.                  \tag{98.2}
\]
Let \(x_t\), \(0\le t\le1\), be a \(C^1\) path and set
\[
 G_t(y)=\dot x_t\cdot\nabla_xW(x_t,y),\qquad
 g_t(y)=-G_t(y)+\mathbb E_{K_{x_t}}G_t.                      \tag{98.3}
\]
There is a unique mean-zero weak solution \(\psi_t\) of
\[
 -{\cal L}_{x_t}\psi_t=g_t.                                  \tag{98.4}
\]
Its velocity \(v_t=\nabla_y\psi_t\) solves
\[
 \partial_tp_{x_t}+\operatorname{div}_y(p_{x_t}v_t)=0         \tag{98.5}
\]
and satisfies
\[
 \boxed{
 \|v_t\|_{L^2(K_{x_t})}
 \le \frac{\|\nabla_yG_t\|_{L^\infty}}{m_*}.}                 \tag{98.6}
\]
If \(v_t\) has the regularity needed for its flow
\(\Phi_{s,t}\), then
\[
 (\Phi_{s,t})_*K_{x_s}=K_{x_t}.                              \tag{98.7}
\]
Thus, after fixing a base point \(x_\circ\) and the straight path
from \(x_\circ\) to \(x\), the map
\[
 \phi_x=\Phi_{0,1}:(\mathbb R^m,K_{x_\circ})
                 \longrightarrow(\mathbb R^m,K_x)            \tag{98.8}
\]
is an almost-everywhere invertible conditional-fiber
trivialization.
\end{theorem}

\begin{proof}
The Brascamp--Lieb inequality implied by \((98.2)\) is
\[
 \operatorname{Var}_{K_x}(f)
 \le\int\nabla f^{\mathsf T}
       (\nabla_y^2W)^{-1}\nabla f\,\dd K_x
 \le m_*^{-1}\int|\nabla f|^2\,\dd K_x.                      \tag{98.9}
\]
Lax--Milgram on the mean-zero weighted Sobolev space gives the
unique solution of \((98.4)\).  Testing by \(\psi_t\), using
\((98.9)\), and applying Cauchy--Schwarz gives
\[
 \|\nabla\psi_t\|_2^2
 =\langle\psi_t,g_t\rangle
 \le m_*^{-1/2}\|\nabla\psi_t\|_2\|g_t\|_2.                  \tag{98.10}
\]
A second application of Brascamp--Lieb to \(G_t\) yields
\[
 \|g_t\|_2^2=\operatorname{Var}_{K_{x_t}}(G_t)
 \le m_*^{-1}\|\nabla_yG_t\|_\infty^2.
\]
This proves \((98.6)\).

Differentiating the normalized density gives
\[
 \partial_t\log p_{x_t}
 =-G_t+\mathbb E_{K_{x_t}}G_t=g_t.
\]
Since
\[
 \operatorname{div}(p_{x_t}\nabla\psi_t)
 =p_{x_t}{\cal L}_{x_t}\psi_t=-p_{x_t}g_t,
\]
equation \((98.5)\) follows.  The characteristic formula for the
continuity equation proves \((98.7)\).  A regular ODE flow has the
inverse \(\Phi_{t,s}\), which proves \((98.8)\).
\end{proof}

\begin{corollary}[Finite old-to-detail influence]
\label{cor:v14v98-influence}
Under the assumptions of
Theorem~\ref{thm:v14v98-moser},
\[
 W_2(K_{x_0},K_{x_1})
 \le\int_0^1
 \frac{\|\nabla_y(\dot x_t\cdot\nabla_xW(x_t,\cdot))\|_\infty}
      {m_*}\,\dd t.                                          \tag{98.11}
\]
In particular, if only old block \(x_i\) changes by \(\delta\) and
\[
 \sup_{x,y}\|\nabla_y\partial_{x_i}W(x,y)\|\le L_i,
\]
then
\[
 \boxed{W_2(K_x,K_{x+\delta e_i})
        \le (L_i/m_*)|\delta|.}                              \tag{98.12}
\]
\end{corollary}

\begin{proof}
Use the coupling
\((Y_0,\Phi_{0,1}(Y_0))\) with \(Y_0\sim K_{x_0}\).
Minkowski's integral inequality and \((98.6)\) give
\[
 \left(\mathbb E|Y_1-Y_0|^2\right)^{1/2}
 \le\int_0^1\|v_t\|_{L^2(K_{x_t})}\,\dd t,
\]
which proves both statements.
\end{proof}

\begin{corollary}[Reflection covariance on the convex branch]
\label{cor:v14v98-reflection}
Suppose \(\vartheta_X,\vartheta_Y\) are orthogonal involutions,
\(x_\circ=\vartheta_Xx_\circ\), and
\[
 W(\vartheta_Xx,\vartheta_Yy)=W(x,y).                         \tag{98.13}
\]
If the path assigned to \(\vartheta_Xx\) is the reflected path
assigned to \(x\), then the transport \((98.8)\) obeys
\[
 \phi_{\vartheta_Xx}\circ\vartheta_Y
 =\vartheta_Y\circ\phi_x.                                    \tag{98.14}
\]
\end{corollary}

\begin{proof}
Pulling \((98.4)\) back by the two involutions gives the same
mean-zero weak equation.  Uniqueness makes the velocities
equivariant, and uniqueness of their flows proves \((98.14)\).
\end{proof}

Equations \((98.8)\) and \((98.14)\) populate the measure and
reflection rows of Definition~\ref{def:v14v97-trivialization} on
the stated convex branch.  Equation \((98.12)\) is a quantitative
influence estimate.  It is not yet a spatial quasilocality
estimate: the inverse weighted elliptic operator in \((98.4)\) can
spread a locally supported parameter derivative throughout the
detail shell.  Exponential off-diagonal decay would require an
additional uniform massive resolvent estimate for
\(-{\cal L}_x\).

\begin{proposition}[Why quartic control is insufficient]
\label{prop:v14v98-quartic}
Coercive quartic integrability does not imply \((98.2)\).  For
example,
\[
 W(y)=\lambda y^4-by^2,\qquad \lambda,b>0,                   \tag{98.15}
\]
defines a smooth probability law with moments of every order, but
\[
 W''(0)=-2b<0.
\]
Hence a source hypothesis saying only that positive quartic Higgs
growth dominates all negative lower-order terms cannot supply a
uniform conditional convexity constant.
\end{proposition}

\begin{proof}
The positive leading coefficient makes \(e^{-W}\) integrable
against every polynomial, while direct differentiation gives the
negative Hessian at the origin.
\end{proof}

\begin{proposition}[Compact gauge obstruction to global strict convexity]
\label{prop:v14v98-compact}
Let \(M\) be a compact connected boundaryless Riemannian manifold
of positive dimension.  No \(C^2\) function \(w:M\to\mathbb R\)
can satisfy
\[
 \operatorname{Hess}w\succeq m_*g,\qquad m_*>0,               \tag{98.16}
\]
globally.  In particular, the compact
\(G_{\rm SM}\)-valued link sector cannot be covered by the
Euclidean theorem through one globally strongly convex chart.
\end{proposition}

\begin{proof}
Taking the metric trace in \((98.16)\) gives
\(\Delta w\ge m_*\dim M>0\).  Integration over \(M\) contradicts
\(\int_M\Delta w\,\dd{\rm vol}=0\).
\end{proof}

\begin{assessment}[Literal RPESM applicability]
\label{ass:v14v98-source}
The source cutoff definition supplies:
\begin{enumerate}
\item compact gauge, scheduler, route, and boundary variables;
\item an unbounded Higgs variable with a positive quartic leading
      term;
\item domination of negative terms sufficient for normalization
      and polynomial moments;
\item a finite-range interacting shell density.
\end{enumerate}
It does not state \((98.2)\), a mixed derivative bound \(L_i\), a
smooth positive density on one Euclidean fiber, or an off-diagonal
resolvent estimate.  Propositions
\ref{prop:v14v98-quartic}--\ref{prop:v14v98-compact} show that the
first two source sectors cannot be placed under one global
strong-convexity argument.  Therefore the V98 transport theorem is
valid but cannot yet be applied to the full literal shell.
\end{assessment}

There is also a composition issue.  The straight-path Moser map
from a fixed base is canonical only after the base and coordinate
path have been selected.  For three cutoffs it need not obey
\[
 \phi_{n+2/n}=\phi_{n+2/n+1}\circ\phi_{n+1/n};                \tag{98.17}
\]
the failure is the holonomy of the Poisson velocity connection.
Exact cutoff compatibility requires either zero holonomy or the
retention of this transport as an additional cocycle.

\begin{decision}[V99 bounded-perturbation local dynamics]
\label{dec:v14v98-next}
Do not force a global convex chart onto the compact gauge and
nonconvex Higgs sectors.  Instead construct the physical detail
generator directly:
\begin{enumerate}
\item choose constraint-preserving compact-group heat-bath or
      diffusion blocks;
\item use a coercive reference Higgs block with a proved local
      Poincare gap;
\item compare the interacting conditional blocks by bounded
      oscillation where available;
\item combine the blocks by the V82 coverage and V81 influence
      compiler.
\end{enumerate}
State exactly which oscillation constants grow with the unbounded
old field.  If they are not uniform, replace a false global bound
by a high-probability good-set and Lyapunov decomposition.
\end{decision}

\begin{firewall}[Claim boundary after V98]
\label{fire:v14v98}
A reflection-covariant faithful transport with finite influence is
proved on a uniformly convex Euclidean shell.  The literal RPESM
shell contains compact and potentially nonconvex factors outside
that theorem, and no uniform quasilocal transport has been claimed.
\end{firewall}

\section*{Version V98 append-only record}
\addcontentsline{toc}{section}{Version V98 append-only record}
The complete V97 source was retained before this continuation.  It
had \(1139804\) bytes, \(34107\) lines, and SHA--256
\[
 \texttt{1537256A932A3F904788E9AA089933BC58BCF88F9A8E72BCA3A9E0BECB375EC2}.
\]
V98 constructed the conditional Moser trivialization on the
uniformly convex branch, proved its influence and reflection
bounds, and isolated the compact-topology, nonconvex-quartic,
spatial-resolvent, and cutoff-holonomy obstructions.


\section{Fourteenth-cycle V99: exact heat-bath fiber gaps and failure of uniform total-variation influence}
\label{sec:v14v99}

The source generator already contains a finite sum of local
conditional-expectation updates.  This gives an exact local
spectral statement without any convexity or bounded-oscillation
hypothesis.

Let \(\Omega_s\) be one physical constraint sector.  For a
constraint-preserving block \(B\), let
\[
 {\cal E}_Bf
   =\mathbb E_{\mu_s}[f\mid{\cal F}_{B^c}]
\]
be conditional expectation with respect to the variables outside
\(B\), and let \(r_B>0\).  Define
\[
 L_s^{\rm hb}
   =\sum_{B\in{\mathfrak B}}r_B({\cal E}_B-I),\qquad
 A_s^{\rm hb}=-L_s^{\rm hb}.                                 \tag{99.1}
\]

\begin{theorem}[Exact local projection gap]
\label{thm:v14v99-fiber}
Each \({\cal E}_B\) is an orthogonal projection on
\(L^2(\mu_s)\), and
\[
 \langle f,(I-{\cal E}_B)f\rangle
 =\mathbb E_{\mu_s}
   \operatorname{Var}_{\mu_s}(f\mid{\cal F}_{B^c}).           \tag{99.2}
\]
Consequently the nonzero spectrum of \(I-{\cal E}_B\) is the
singleton \(\{1\}\).  The rate-\(r_B\) block update has exact
conditional fiber gap \(r_B\), independently of whether the
conditional law is compact, nonconvex, or unbounded.
Moreover
\[
 \langle f,A_s^{\rm hb}f\rangle
 =\sum_{B\in{\mathfrak B}}r_B
   \mathbb E_{\mu_s}
   \operatorname{Var}_{\mu_s}(f\mid{\cal F}_{B^c}).           \tag{99.3}
\]
\end{theorem}

\begin{proof}
Conditional expectation is the orthogonal projection onto
\(L^2({\cal F}_{B^c},\mu_s)\).  Pythagoras gives
\[
 \|f\|_2^2-\|{\cal E}_Bf\|_2^2
 =\|f-{\cal E}_Bf\|_2^2,
\]
which is \((99.2)\).  Every orthogonal projection has spectrum
contained in \(\{0,1\}\); on the orthogonal complement of its
range, \(I-{\cal E}_B\) is the identity.  Summing the block forms
proves \((99.3)\).
\end{proof}

\begin{corollary}[Kernel and sector criterion]
\label{cor:v14v99-kernel}
The refresh-free kernel is exactly
\[
 \ker A_s^{\rm hb}
 =\bigcap_{B\in{\mathfrak B}}
   L^2({\cal F}_{B^c},\mu_s).                                \tag{99.4}
\]
It consists only of constants precisely when the common
block-invariant sigma algebra is trivial modulo \(\mu_s\).
If a physical superselection label is unchanged by every block,
then functions of that label remain in the kernel until the
analysis is restricted to a selected sector.
\end{corollary}

\begin{proof}
Every summand in \((99.3)\) is nonnegative.  Their sum vanishes
exactly when every conditional variance vanishes, which is
equivalent to membership in every displayed sigma algebra.
\end{proof}

Thus the missing number is not a block gap.  It is an approximate
tensorization constant comparing total variance with the sum in
\((99.3)\).  The V81--V82 compiler gives the following direct
specialization.

\begin{corollary}[Refresh-free assembly under a populated influence row]
\label{cor:v14v99-assembly}
Suppose the constraint-preserving block cover has fractional
coverage \(\kappa_*>0\), and its post-constraint influence
constant is \(\widehat\alpha<1\) in the sense of V82.  Then
\[
 \boxed{\operatorname{gap}(A_s^{\rm hb})
 \ge r_*\kappa_*(1-\widehat\alpha),\qquad
 r_*=\min_Br_B.}                                             \tag{99.5}
\]
No global refresh enters this estimate.
\end{corollary}

\begin{proof}
Theorem~\ref{thm:v14v99-fiber} supplies the unit fiber-gap row
required by the V82 overlapping-block theorem.  Multiplication by
the smallest update rate and insertion of the fractional coverage
gives \((99.5)\).
\end{proof}

We now test the total-variation influence route on the unbounded
Higgs sector.  For a block conditional density
\[
 K_B(\dd u\mid z)
 =Z_B(z)^{-1}e^{-U_B(u;z)}\nu_B(\dd u),
\]
define, when outside configurations \(z,z'\) differ only in block
\(C\),
\[
 \Delta_{BC}
 =\sup_{u,u',z\sim_Cz'}
 \left|
 \bigl[U_B(u;z)-U_B(u;z')\bigr]
 -\bigl[U_B(u';z)-U_B(u';z')\bigr]
 \right|.                                                    \tag{99.6}
\]
The likelihood-ratio lemma of V81 gives
\[
 c_{BC}
 :=\sup_{z\sim_Cz'}\|K_B(\cdot\mid z)-K_B(\cdot\mid z')\|_{\rm TV}
 \le\tanh(\Delta_{BC}/4)                                     \tag{99.7}
\]
when \(\Delta_{BC}<\infty\).

\begin{proposition}[An unbounded quartic channel has maximal TV influence]
\label{prop:v14v99-tv-one}
Let
\[
 K_z(\dd h)
 =Z(z)^{-1}\exp\{-\lambda h^4-Jzh\}\,\dd h,
 \qquad \lambda>0,\quad J\ne0.                               \tag{99.8}
\]
Then
\[
 \sup_{z,z'\in\mathbb R}
 \|K_z-K_{z'}\|_{\rm TV}=1.                                  \tag{99.9}
\]
In particular, the uniform Dobrushin coefficient of this
nearest-neighbor quartic channel is one, although every
single-block heat-bath gap remains exactly one.
\end{proposition}

\begin{proof}
After changing signs, take \(J>0\).  The unique minimizer of
\(\lambda h^4+Jzh\) is
\[
 h_z=-\left(\frac{Jz}{4\lambda}\right)^{1/3}
 \quad(z>0),                                                  \tag{99.10}
\]
and \(h_{-z}=-h_z\).  Put \(a_z=|h_z|\) and rescale
\(h=a_zs\).  Since \(Jz=4\lambda a_z^3\),
\[
 \lambda h^4+Jzh
 =\lambda a_z^4(s^4+4s).
\]
The function \(s^4+4s\) has its strict minimum at \(s=-1\).
Laplace concentration therefore gives, for every fixed
\(0<\varepsilon<1\),
\[
 K_z\{|h+a_z|<\varepsilon a_z\}\longrightarrow1,
 \qquad
 K_{-z}\{|h+a_z|<\varepsilon a_z\}\longrightarrow0.
\]
The total-variation distance is at least the difference of these
two probabilities, hence tends to one.  Since total variation is
never larger than one, \((99.9)\) follows.
\end{proof}

The same mechanism occurs whenever an unbounded neighboring field
enters the conditional Higgs potential through a nonzero linear
channel.  Quartic confinement normalizes the conditional law, but
does not make conditionals with arbitrarily separated external
fields overlap uniformly.

On the bounded set \(|h|,|z|\le R\), the interaction \(Jzh\)
instead gives
\[
 \Delta_{BC}(R)\le4|J|R^2,\qquad
 c_{BC}(R)\le\tanh(|J|R^2).                                  \tag{99.11}
\]
For an interaction graph of outside degree \(D\), the elementary
row bound is
\[
 \alpha_R\le D\tanh(|J|R^2),                                 \tag{99.12}
\]
which ceases to be contractive as \(R\) grows whenever \(D>1\).
Thus expanding compact truncations do not produce the required
uniform full-space constant.

\begin{lemma}[High probability of a good set does not control the gap]
\label{lem:v14v99-goodset-no-gap}
For every \(0<\varepsilon<1\) and every \(\delta>0\), there is a
reversible two-state continuous-time chain with stationary mass
\(1-\varepsilon\) on a good state and spectral gap exactly
\(\delta\).
\end{lemma}

\begin{proof}
Let the rates from good to bad and bad to good be
\[
 q_{gb}=\varepsilon\delta,\qquad
 q_{bg}=(1-\varepsilon)\delta.
\]
Detailed balance holds for stationary weights
\((1-\varepsilon,\varepsilon)\), and the sole nonzero eigenvalue
of \(-L\) is \(q_{gb}+q_{bg}=\delta\).
\end{proof}

Lemma~\ref{lem:v14v99-goodset-no-gap} prevents an illegitimate
upgrade of \((99.11)\) from a high-probability truncation to a
global Poincare inequality.  A return or tail-conductance estimate
is an independent row.

\begin{assessment}[Exact state after V99]
\label{ass:v14v99-state}
For the literal source local generator:
\begin{enumerate}
\item the local conditional fiber gaps are exactly known;
\item the refresh-free global kernel is exactly characterized;
\item a sectorwise global gap follows from the already proved
      coverage/influence compiler;
\item ordinary uniform total-variation influence can fail
      maximally in an allowed unbounded quartic channel;
\item truncation requires an additional Lyapunov or return
      estimate and cannot be justified by tail mass alone.
\end{enumerate}
The next analytic route should use a distance adapted to quartic
confinement, in which far external fields can have controlled
transport even when their probability laws have TV distance near
one.
\end{assessment}

\begin{decision}[V100 audit and next-cycle handoff]
\label{dec:v14v99-next}
Perform the compulsory YM/NS checkpoint.  Test specifically for:
\begin{enumerate}
\item a weighted Wasserstein influence matrix;
\item a Lyapunov drift or return inequality;
\item a constraint-preserving block cover;
\item a local-to-global gap theorem compatible with unbounded
      polynomial fields.
\end{enumerate}
Then close the fourteenth cycle at V100.  Freeze it as
\(\texttt{\_f14}\), start the next compact V1 with a major-results
summary and context, and continue there with a quartic-adapted
Wasserstein or Lyapunov compiler rather than the failed uniform-TV
route.
\end{decision}

\begin{firewall}[Claim boundary after V99]
\label{fire:v14v99}
The source's refresh-free block gaps are proved exactly, but no
uniform global RPESM gap follows until coverage, sector
irreducibility, and a full-space dependence or return estimate
are populated.  Total-variation Dobrushin contraction is not
assumed.
\end{firewall}

\section*{Version V99 append-only record}
\addcontentsline{toc}{section}{Version V99 append-only record}
The complete V98 source was retained before this continuation.  It
had \(1149395\) bytes, \(34381\) lines, and SHA--256
\[
 \texttt{7259FFCF7B0AF783D5E57F950655E615D70B863184F8544320C08E5624B9ED25}.
\]
V99 extracted the exact unit local heat-bath gaps, specialized the
refresh-free assembly theorem, proved maximal total-variation
influence in an unbounded quartic channel, and isolated the
independent tail-return row.


\section{Fourteenth-cycle V100: compulsory companion audit and cycle closure}
\label{sec:v14v100}

This checkpoint tests the weighted-transport and tail-return rows
identified in V99 before the active lineage is frozen.

\begin{record}[V100 companion checkpoint]
\label{rec:v14v100-companions}
The newest Yang--Mills note is
\begin{center}
\path{Renewal_Yang_Mills_Mass_Gap_Research_Notes_V30.tex},
\end{center}
with \(246165\) bytes, \(6697\) lines, and SHA--256
\[
 \texttt{C42E60C54FB85A218B5FC7AE91C1930DE0A797218F73A493D039356C09ACFAB6}.
                                                               \tag{100.1}
\]
The Navier--Stokes note remains
\begin{center}
\path{Renewal_NS_Stability_Research_Notes_V26.tex},
\end{center}
with \(278283\) bytes, \(5319\) lines, and SHA--256
\[
 \texttt{82C887F8C2C69E4F28FAD7C3DC8DE5B9099CB5A4BF431EBA1FFF3E91935978AD}.
                                                               \tag{100.2}
\]
\end{record}

Yang--Mills V30 constructs a second exact rational splitting of the
four-thirds Wilson pencil at the physical point which invalidated
the old fixed splitting.  It proves a full-active-circle tube of
transverse radius \(99/800\), and proves that this tube overlaps the
earlier radius-\(1/8\) V28 tube.  This is an exact two-chart
pencil atlas across the V29 fixed-splitting obstruction.

\begin{assessment}[Precise YM transfer]
\label{ass:v14v100-ym}
The transferable principle is:
\begin{quote}
use independently certified local charts and verify their overlap
on the physical locus; do not require a single global
factorization when a physical obstruction to that factorization is
already known.
\end{quote}
For the SM shell, this supports a finite-atlas approach to compact
gauge variables.  The V30 charts themselves live in transverse
momentum space and certify a Wilson matrix pencil.  They do not
supply conditional gauge-coordinate charts, transition Jacobians,
a heat-bath influence matrix, a Lyapunov function, or a Markov
gap.  None of the numerical radii or matrix floors in YM V30 is
therefore imported into \((99.5)\).
\end{assessment}

Navier--Stokes V26 is unchanged.  Its rank-one contractions of
first and second derivatives of the Oseen kernel are deterministic
kernel estimates, not weighted Wasserstein contractions for an
unbounded Gibbs sampler.

\begin{theorem}[V100 companion verdict]
\label{thm:v14v100-verdict}
At checkpoints \((100.1)\)--\((100.2)\), neither companion
populates any of the four requested rows:
\begin{enumerate}
\item a quartic-adapted weighted Wasserstein influence matrix;
\item a Lyapunov drift or uniform tail-return inequality;
\item a constraint-preserving RPESM block cover;
\item a local-to-global gap theorem for the literal unbounded
      polynomial field law.
\end{enumerate}
The YM atlas method may be reused only after new gauge-fiber charts
and overlap transports have been constructed on the RPESM state
space.
\end{theorem}

\begin{proof}
The complete new YM V30 section contains rational
finite-dimensional Wilson-pencil inequalities and their overlap
geometry; it explicitly leaves the interacting measure,
reflection-positive continuum, clustering, and Hamiltonian gap
open.  NS V26 contains no conditional probability law or Markov
generator.  Their object types do not match any of the four rows.
\end{proof}

\begin{theorem}[Fourteenth-cycle spectral and locality ledger]
\label{thm:v14v100-ledger}
The exact conclusions acquired in V93--V99 are:
\begin{enumerate}
\item the Grand Tensor source contains an existential exactly
      projective positive RPESM regulator, while physical chiral
      selection, refresh-free local gap, correspondence, uniform
      stress bounds, and the continuum limit remain unpopulated;
\item the source global refresh shifts the centered local spectrum
      by exactly \(\rho\), lifts relative-sector zero modes, and
      cancels dollar-for-dollar from an honest comparison with the
      refresh-free physical generator;
\item the exact shell skeleton has block form
      \(P_{n+1}\simeq P_n\oplus0\), hence inherits the old
      Dirichlet gap but annihilates the new conditional-detail
      space and has no finite logarithmic Hamiltonian there;
\item scalar laziness cannot keep both a nonzero finite old-mode
      energy and a nonzero finite detail energy on one continuum
      clock;
\item a conditional-fiber isomorphism gives an injective faithful
      old transfer; adding detail dynamics yields the exact product
      gap \(\min\{\gamma,\beta\}\), and a noncommuting symmetric
      split obeys the angle and commutator bound \((97.20)\);
\item conditional Moser transport supplies such isomorphisms and a
      finite influence bound on uniformly convex Euclidean fibers,
      but compact gauge topology, allowed nonconvex quartic Higgs
      terms, elliptic nonlocality, and cutoff holonomy block a
      direct global application;
\item every source local conditional-expectation update has exact
      fiber gap one, so global assembly is the remaining problem;
\item an allowed unbounded quartic linear channel has uniform
      total-variation influence exactly one, and high probability
      of a bounded good set alone cannot control the gap.
\end{enumerate}
Each item is a proved theorem or obstruction with the claim
boundary stated above; together they do not prove the full
SM--Einstein continuum.
\end{theorem}

\begin{proof}
Items one through eight are respectively the content of the source
audit V93, refresh calculation V94, shell spectral theorem V96,
laziness no-go V96, faithful split construction V97, conditional
transport analysis V98, and heat-bath/TV results V99.  The
append-only records preserve their full hypotheses and proofs.
\end{proof}

\begin{decision}[Next compact cycle]
\label{dec:v14v100-next}
Freeze this V100 as
\(\texttt{Renewal\_SM\_Einstein\_Continuum\_Research\_Notes\_V100\_f14.tex}\).
Begin a new compact V1 with:
\begin{enumerate}
\item context and a major-results summary;
\item the exact V100 archive hash;
\item a weighted Wasserstein comparison theorem for quartic
      single-block conditionals;
\item a path-coupling or coarse-Ricci compiler for the block
      heat-bath chain;
\item a separate Lyapunov return row for any region where the
      Wasserstein contraction is not uniform.
\end{enumerate}
The new metric must remain finite when total variation is one and
must respect compact gauge charts and physical constraints.
\end{decision}

\begin{firewall}[Claim boundary at the fourteenth-cycle close]
\label{fire:v14v100}
The finite regulator, exact shell, faithful-transfer, and local
heat-bath mechanisms are now separated sharply.  A refresh-free
uniform physical gap, cutoff-compatible local transport,
continuum Hamiltonian, common stress-source limit, and coupled
Einstein solution remain open.
\end{firewall}

\section*{Version V100 append-only record}
\addcontentsline{toc}{section}{Version V100 append-only record}
The complete V99 source was retained before this continuation.  It
had \(1158869\) bytes, \(34655\) lines, and SHA--256
\[
 \texttt{9AD245C7AC25A5AA459EA4D41DBB4543441262ED2F5E5B8F58F252C47A9311E1}.
\]
V100 completed the mandatory companion audit, imported only the
finite-atlas method, and closed the fourteenth cycle with an exact
spectral/locality ledger and a weighted-transport handoff.

\end{document}
